\documentclass[11pt,a4paper]{report}
\usepackage{fontspec,amsmath,amssymb,amsthm,microtype,xcolor}
\setmainfont{lmroman10-regular.otf}[BoldFont=lmroman10-bold.otf,ItalicFont=lmroman10-italic.otf,BoldItalicFont=lmroman10-bolditalic.otf]
\usepackage[most]{tcolorbox}
\usepackage{mathrsfs}
\providecommand{\lt}{<}
\providecommand{\gt}{>}
\setmonofont{DejaVuSansMono.ttf}[BoldFont=DejaVuSansMono-Bold.ttf,ItalicFont=DejaVuSansMono-Oblique.ttf,BoldItalicFont=DejaVuSansMono-BoldOblique.ttf,Scale=0.86]
\newfontfamily\leanfallback{STIXTwoMath-Regular.otf}[Scale=0.86]
\definecolor{leanrule}{gray}{0.62}
\definecolor{leanback}{gray}{0.962}
\definecolor{leantitle}{gray}{0.22}
\definecolor{leanmuted}{gray}{0.38}
\newtcolorbox{leanbox}[1]{enhanced,breakable,colback=leanback,colframe=leanrule,
  boxrule=0.4pt,arc=0.6mm,left=2.4mm,right=2.4mm,top=1.8mm,bottom=1.8mm,
  before skip=0.8em,after skip=1em,fonttitle=\small\bfseries,coltitle=leantitle,
  colbacktitle=leanback,titlerule=0.3pt,title={#1}}
\newlength{\leancharwidth}
\newenvironment{leancode}{\par\smallskip\begingroup\ttfamily\footnotesize\raggedright
  \settowidth{\leancharwidth}{M}\setlength{\parskip}{0pt}\setlength{\parindent}{0pt}
  \frenchspacing}{\par\endgroup\smallskip}
\newcommand{\leanline}[2]{\hangindent=\dimexpr#1\leancharwidth+4\leancharwidth\relax
  \hangafter=1\noindent\hspace*{#1\leancharwidth}#2\par}
\newcommand{\leansource}[1]{\par{\footnotesize\color{leanmuted}#1}\par}
\newcommand{\leanlabel}[1]{\par\smallskip\noindent{\small\bfseries #1}\par}
\usepackage[margin=25mm]{geometry}
\usepackage[unicode,hidelinks]{hyperref}
\hypersetup{pdfauthor={LangevinLean Project},bookmarksdepth=2}
\setlength{\emergencystretch}{3em}
\setlength{\parskip}{0.4em}
\setlength{\parindent}{0pt}
\setcounter{tocdepth}{1}
\makeatletter
\renewcommand*\l@section{\@dottedtocline{1}{1.5em}{3.8em}}
\renewcommand*\@pnumwidth{2.5em}
\renewcommand*\@tocrmarg{3.5em}
\makeatother
\setcounter{secnumdepth}{2}
\allowdisplaybreaks[2]
\raggedbottom
\title{\texorpdfstring{Wasserstein Mixing Time of the Unadjusted Langevin Algorithm}{Wasserstein Mixing Time of the Unadjusted Langevin Algorithm}\\[0.6em]\large Mathematical proofs}
\author{LangevinLean Project}
\date{Draft report LL-11\\[0.4em]\small Content revision e7b1521ab333}
\begin{document}
\pagenumbering{roman}
\maketitle
\chapter*{About this document}
\addcontentsline{toc}{chapter}{About this document}
Chapter 1 states the assumptions and the main results. Chapter 2 proves them, following the argument from the hypotheses to the conclusion in order. Chapter 3 proves the supporting lemmas, each before its first use. The proofs are written for readers with an undergraduate background in analysis, linear algebra and probability; definitions are recalled where they are first used, and each inequality is justified.

\section*{Mathematics and its Lean formalization}
Every result is followed by a grey box headed \emph{Lean}. The box reproduces the exact statement of the corresponding Lean declaration, read from the source file and line range named in the box, and then explains how the Lean proof follows the mathematical argument: which hypothesis encodes which assumption, and which formal step carries out which step of the proof. The mathematics can be read without the boxes; the boxes connect each claim to the statement that Lean has checked.

\section*{Checking the proofs}
The results are verified by Lean 4 with the Mathlib library, at the versions pinned in the repository \url{https://github.com/odakyildiz/LangevinLean}. After installing the Lean toolchain manager elan and cloning the repository, run from its root:
\begin{leancode}
\leanline{0}{lake\ exe\ cache\ get}
\leanline{0}{lake\ build\ WassersteinMixing}
\leanline{0}{lake\ env\ lean\ WassersteinMixing/Check.lean}
\end{leancode}
The second command checks every proof in the library. The third prints the axioms on which the main theorems depend; they are only the standard axioms of propositional extensionality, choice and quotient soundness.

Prepared by the LangevinLean Project with AI assistance, these notes form part of an ongoing effort to formalize Langevin algorithms comprehensively. This is draft work and may contain mistakes.

\section*{Scope of the checked statements}


\textbf{Theorem 1 and Corollary 1.1 are proved for the actual Gibbs target, independent-Gaussian ULA laws and constructed invariant measure, without additional realization assumptions.}



Theorem 1 and Corollary 1.1 are proved under the original \(C^2\) Hessian assumptions for the actual Gibbs target, independent-Gaussian ULA laws and constructed invariant measure. The proof constructs the stationary Brownian dynamics, reversible Markov semigroup and bounded Poisson corrector directly at \(C^2\) regularity. No eigenbasis, Markov realization, invariant law, comparison estimate or smoothing hypothesis is added to the final statements.

\section*{Literature and scope}
\begin{enumerate}
\item Francesco Pedrotti and Peter A. Whalley (2026). \emph{Wasserstein mixing time of the unadjusted Langevin algorithm}. arXiv preprint arXiv:2608.02430, version 1 (3 August 2026). \url{https://arxiv.org/abs/2608.02430v1}.
\end{enumerate}
Results from the earlier reports \cite{LangevinLean2026ULASC,LangevinLean2026IPLA,LangevinLean2026Kinetic} are used below. Each application cites the section containing the required result.

\tableofcontents
\clearpage
\pagenumbering{arabic}
\chapter{Assumptions and main results}
\section*{Assumptions and notation from the source}
\subsection*{\texorpdfstring{Assumption 1.1 — Strong log-concavity and log-smoothness}{Assumption 1.1 — Strong log-concavity and log-smoothness}}


Source: Francesco Pedrotti and Peter A. Whalley (2026), Section 1, Assumption 1.1.



Let \(V\in C^2(\mathbb R^d)\) and let \(0<\alpha\le\beta\).
                 For every \(x\in\mathbb R^d\),


                 \[\alpha I_d\preccurlyeq\nabla^2V(x)\preccurlyeq\beta I_d,
                 \qquad\kappa=\frac\beta\alpha.\]


The target is \(\pi(dx)=Z^{-1}e^{-V(x)}dx\).
                 The source step restriction is \(0<h\le1/\beta\).
                 The quadratic-form bounds on the actual Hessian derive the
                 strong-convexity and gradient-Lipschitz contract used internally.

\subsection*{\texorpdfstring{Equations (1)–(2) — Diffusion and Gaussian ULA recursion}{Equations (1)–(2) — Diffusion and Gaussian ULA recursion}}


Source: Francesco Pedrotti and Peter A. Whalley (2026), Section 1, equations (1)–(2), and the definition of quadratic transport.

\[dX_t=-\nabla V(X_t)\,dt+\sqrt2\,dB_t,
                 \qquad\widehat X_{k+1}=\widehat X_k-h\nabla V(\widehat X_k)
                 +\sqrt2(B_{(k+1)h}-B_{kh}).\]


The source defines \(W_2\) on
                 \(\mathcal P_2(\mathbb R^d)\), so the finite-time and accuracy
                 statements take a probability initialization with finite second moment.
                 The biased law \(\widehat\pi\) is constructed from an iid Gaussian past,
                 with finite second moment, actual ULA invariance and transport attraction.

\subsection*{\texorpdfstring{Assumption 2.1 — Temporary regularity in the source proof}{Assumption 2.1 — Temporary regularity in the source proof}}


Source: Francesco Pedrotti and Peter A. Whalley (2026), Section 2, Assumption 2.1; removed in Section 2.2.



The source first assumes that \(V\) is smooth and
                 every partial derivative has at most polynomial growth, then
                 removes this temporary assumption by Gaussian convolution.
                 The formal proof instead constructs the actual reversible
                 semigroup and its corrector directly under \(C^2\).
                 This temporary source-proof hypothesis is therefore not used
                 by the final statements, and no smoothing comparison is required.

\section{\texorpdfstring{Theorem 1 — The invariant bias bound}{Theorem 1 — The invariant bias bound}}


Source: Francesco Pedrotti and Peter A. Whalley (2026), Section 1, Theorem 1, equation (3).



\textbf{Paper statement proved.}



Let \(V\in C^2(\mathbb R^d)\), \(0<\alpha\le\beta\), and
\(\alpha I_d\preccurlyeq\nabla^2V\preccurlyeq\beta I_d\).
For \(0<h\le1/\beta\), set \(\kappa=\beta/\alpha\).
The Gibbs target is \(\pi(dx)=Z^{-1}e^{-V(x)}dx\), and
\(\widehat\pi\) is the constructed invariant law of the actual independent-Gaussian ULA transition. Then


\[W_2(\pi,\widehat\pi)\le6h\sqrt{\kappa\beta d}.\]


The invariant law, stationary Langevin process, reversible Markov semigroup,
Poisson corrector and synchronous comparison are all constructed from these
source assumptions. No spectral, Markov, invariant-law or comparison estimate
is an additional hypothesis of this endpoint.

\begin{leanbox}{Lean: \texttt{bias\_\allowbreak{}bound}}
\leansource{\texttt{bias\_\allowbreak{}bound} \quad WassersteinMixing/\allowbreak{}Theorem.lean, lines 21–27}
\begin{leancode}
\leanline{0}{theorem\ bias\_bound\ (P\ :\ Parameters)\ (V\ :\ E\ P.dimension\ →\ ℝ)}
\leanline{4}{(hA\ :\ HessianAssumptions\ P\ V)\ :}
\leanline{4}{ULA.wasserstein2Real\ (ULA.gibbsMeasure\ V)}
\leanline{6}{(invariantLaw\ P\ V\ hA.toPotentialAssumptions)\ ≤}
\leanline{8}{6\ *\ P.step\ *\ Real.sqrt\ (P.condition\ *\ P.beta\ *\ P.dimension)}
\end{leancode}
\end{leanbox}

{\small\color{gray}Supporting proof: Section~\ref{proof-module-WassersteinMixing-Theorem-lean}.\par}
\section{\texorpdfstring{Corollary 1.1 — Finite-time Gaussian ULA convergence}{Corollary 1.1 — Finite-time Gaussian ULA convergence}}


Source: Francesco Pedrotti and Peter A. Whalley (2026), Section 1, Corollary 1.1, equation (4).



\textbf{Paper statement proved.}



Under the same source Hessian assumptions and step restriction,
let \(\mu\in\mathcal P_2(\mathbb R^d)\), and let \(\mu_k\) be the law after
\(k\) actual independent-Gaussian ULA updates. For every \(k\ge0\),


\[W_2(\mu_k,\pi)\le(1-\alpha h)^kW_2(\mu,\widehat\pi)
+6h\sqrt{\kappa\beta d}.\]


The measure \(\widehat\pi\) is the constructed invariant law. The exact
contraction factor and the transport triangle inequality are proved for the
same Gaussian ULA laws appearing in this statement.

\begin{leanbox}{Lean: \texttt{convergence\_\allowbreak{}bound}}
\leansource{\texttt{convergence\_\allowbreak{}bound} \quad WassersteinMixing/\allowbreak{}Theorem.lean, lines 31–39}
\begin{leancode}
\leanline{0}{theorem\ convergence\_bound\ (P\ :\ Parameters)\ (V\ :\ E\ P.dimension\ →\ ℝ)}
\leanline{4}{(hA\ :\ HessianAssumptions\ P\ V)\ (μ\ :\ Measure\ (E\ P.dimension))\ [IsProbabilityMeasure\ μ]}
\leanline{4}{(hμ\ :\ MemLp\ (id\ :\ E\ P.dimension\ →\ E\ P.dimension)\ 2\ μ)\ (n\ :\ ℕ)\ :}
\leanline{4}{ULA.wasserstein2Real\ ((ULA.ulaLawStep\ V\ P.step)\textasciicircum{}[n]\ μ)\ (ULA.gibbsMeasure\ V)\ ≤}
\leanline{6}{P.contraction\ \textasciicircum{}\ n\ *\ ULA.wasserstein2Real\ μ}
\leanline{8}{(invariantLaw\ P\ V\ hA.toPotentialAssumptions)\ +}
\leanline{10}{6\ *\ P.step\ *\ Real.sqrt\ (P.condition\ *\ P.beta\ *\ P.dimension)}
\end{leancode}
\end{leanbox}

{\small\color{gray}Supporting proof: Section~\ref{proof-module-WassersteinMixing-Theorem-lean}.\par}
\section{\texorpdfstring{Corollary 1.1 — The exact iteration complexity}{Corollary 1.1 — The exact iteration complexity}}


Source: Francesco Pedrotti and Peter A. Whalley (2026), Section 1, Corollary 1.1, equation (5).



\textbf{Paper statement proved.}



Under the source Hessian assumptions, let
\(\mu\in\mathcal P_2(\mathbb R^d)\), \(d>0\) and \(\varepsilon>0\).
Choose


\[h=\frac1\beta\min\!\left(1,\frac{\varepsilon}{12\sqrt d}\right).\]


Every natural number \(k\) satisfying


\[k\ge\kappa\max\!\left(1,\frac{12\sqrt d}{\varepsilon}\right)
\log\!\left(1+\frac{2\sqrt\alpha\,W_2(\mu,\widehat\pi)}{\varepsilon}\right)\]


gives \(\sqrt\alpha\,W_2(\mu_k,\pi)\le\varepsilon\).
The step, logarithmic count and constants are the literal formulas in
Corollary 1.1. The paired statement uses the actual Gaussian laws and the
constructed invariant measure.

\begin{leanbox}{Lean: \texttt{accuracy\_\allowbreak{}bound}}
\leansource{\texttt{accuracy\_\allowbreak{}bound} \quad WassersteinMixing/\allowbreak{}Theorem.lean, lines 44–54}
\begin{leancode}
\leanline{0}{theorem\ accuracy\_bound\ (P\ :\ Parameters)\ (V\ :\ E\ P.dimension\ →\ ℝ)}
\leanline{4}{(hA\ :\ HessianAssumptions\ P\ V)\ (μ\ :\ Measure\ (E\ P.dimension))\ [IsProbabilityMeasure\ μ]}
\leanline{4}{(hμ\ :\ MemLp\ (id\ :\ E\ P.dimension\ →\ E\ P.dimension)\ 2\ μ)\ (ε\ :\ ℝ)\ (hε\ :\ 0\ <\ ε)}
\leanline{4}{(hstep\ :\ P.step\ =\ paperStep\ P.beta\ P.dimension\ ε)\ (n\ :\ ℕ)}
\leanline{4}{(hn\ :\ iterationPrefactor\ P.alpha\ P.beta\ P.dimension\ ε\ *}
\leanline{6}{Real.log\ (1\ +\ 2\ *\ Real.sqrt\ P.alpha\ *}
\leanline{8}{ULA.wasserstein2Real\ μ\ (invariantLaw\ P\ V\ hA.toPotentialAssumptions)\ /\ ε)\ ≤\ (n\ :\ ℝ))\ :}
\leanline{4}{Real.sqrt\ P.alpha\ *}
\leanline{6}{ULA.wasserstein2Real\ ((ULA.ulaLawStep\ V\ P.step)\textasciicircum{}[n]\ μ)\ (ULA.gibbsMeasure\ V)\ ≤\ ε}
\end{leancode}
\end{leanbox}

{\small\color{gray}Supporting proof: Section~\ref{proof-module-WassersteinMixing-Theorem-lean}.\par}
\chapter{Mathematical argument}
\section{\texorpdfstring{Deriving the potential contract from the source Hessian bounds}{Deriving the potential contract from the source Hessian bounds}}


The source assumes precisely \(V\in C^2\) and
    \(\alpha I_d\preccurlyeq\nabla^2V\preccurlyeq\beta I_d\).
    Symmetry and the quadratic-form bounds give
    \(\|\nabla^2V(x)\|\le\beta\). The mean-value theorem makes
    \(\nabla V\) globally \(\beta\)-Lipschitz. Integrating the lower Hessian
    inequality along line segments gives the strong supporting-plane
    inequality, hence \(\alpha\)-strong convexity.




Thus the first-order contract used by the existing Langevin library is
    a consequence of the paper's original assumptions. The final proof works
    directly under \(C^2\); it does not require the source proof's temporary
    smoothness and polynomial-growth assumption or a smoothing limit.

\begin{leanbox}{Lean: \texttt{HessianAssumptions.\allowbreak{}toPotentialAssumptions}}
\leansource{\texttt{HessianAssumptions.\allowbreak{}toPotentialAssumptions} \quad WassersteinMixing/\allowbreak{}CurvatureAssumptions.lean, lines 138–142}
\begin{leancode}
\leanline{0}{theorem\ HessianAssumptions.toPotentialAssumptions\ (P\ :\ Parameters)\ (V\ :\ H\ →\ ℝ)}
\leanline{4}{(hV\ :\ HessianAssumptions\ P\ V)\ :\ PotentialAssumptions\ P\ V}
\end{leancode}
\leanlabel{How the Lean proof follows the argument}
\texttt{Hessian\allowbreak{}Assumptions} records the quadratic forms of the
    genuine derivative of the gradient. Its
    \texttt{to\allowbreak{}Potential\allowbreak{}Assumptions} theorem derives every field of the
    internal potential contract.
\end{leanbox}
\section{\texorpdfstring{Constructing the actual stationary Brownian comparison}{Constructing the actual stationary Brownian comparison}}


Start the diffusion at an actual Gibbs-distributed random vector
    independent of a canonical Brownian path. Picard iteration in the
    complete \(L^2\) space constructs the global solution of


    \[X_t-X_0=-\int_0^t\nabla V(X_s)\,ds+\sqrt2\,B_t.\]


Uniqueness identifies its successive segments with the reference
    process used by the synchronous Gaussian ULA chain. The Gibbs law is
    stationary. The independent Gaussian increments have the exact laws
    required by the original ULA recursion; their pushforward chain laws are
    the actual iterates of that transition.

\begin{leanbox}{Lean: \texttt{state\_\allowbreak{}integral}, \texttt{state\_\allowbreak{}law}, \texttt{actualChain\_\allowbreak{}law}}
\leansource{\texttt{state\_\allowbreak{}integral} \quad WassersteinMixing/\allowbreak{}SemigroupGlobal.lean, lines 133–150}
\begin{leancode}
\leanline{0}{theorem\ state\_integral\ (P\ :\ Parameters)\ (V\ :\ E\ P.dimension\ →\ ℝ)}
\leanline{4}{(hV\ :\ PotentialAssumptions\ P\ V)\ (t\ :\ ℝ)\ (ht\ :\ 0\ ≤\ t)\ :}
\leanline{4}{letI\ :=\ actualProbability\_isProbability\ P\ V\ hV}
\leanline{4}{state\ P\ V\ hV\ t\ -\ initial\ P\ V\ hV\ =}
\leanline{6}{-(∫\ s\ in\ 0..t,\ ULA.lpDrift\ (actualProbability\ P\ V)\ (gradient\ V)}
\leanline{8}{hV.gradientLipschitz\ (state\ P\ V\ hV\ s))\ +}
\leanline{6}{ULA.brownianLangevinNoise\ (actualProbability\ P\ V)\ (brownian\ P)}
\leanline{8}{(brownian\_preBrownian\ P\ V\ hV)\ 0\ t}
\end{leancode}
\leansource{\texttt{state\_\allowbreak{}law} \quad WassersteinMixing/\allowbreak{}SemigroupGlobal.lean, lines 206–217}
\begin{leancode}
\leanline{0}{theorem\ state\_law\ (P\ :\ Parameters)\ (V\ :\ E\ P.dimension\ →\ ℝ)}
\leanline{4}{(hV\ :\ PotentialAssumptions\ P\ V)\ (t\ :\ ℝ)\ :}
\leanline{4}{(actualProbability\ P\ V).map\ (state\ P\ V\ hV\ t)\ =\ ULA.gibbsMeasure\ V}
\end{leancode}
\leansource{\texttt{actualChain\_\allowbreak{}law} \quad WassersteinMixing/\allowbreak{}ActualComparison.lean, lines 193–212}
\begin{leancode}
\leanline{0}{theorem\ actualChain\_law\ (P\ :\ Parameters)\ (V\ :\ E\ P.dimension\ →\ ℝ)}
\leanline{4}{(hV\ :\ PotentialAssumptions\ P\ V)\ (n\ :\ ℕ)\ :}
\leanline{4}{(actualProbability\ P\ V).map\ (actualChain\ P\ V\ hV\ n)\ =}
\leanline{6}{(ULA.ulaLawStep\ V\ P.step)\textasciicircum{}[n]\ (ULA.gibbsMeasure\ V)}
\end{leancode}
\leanlabel{How the Lean proof follows the argument}
\texttt{Canonical\allowbreak{}Semigroup.\allowbreak{}state} is the constructed global
    solution. \texttt{state\_\allowbreak{}integral}, \texttt{state\_\allowbreak{}law} and
    \texttt{state\_\allowbreak{}grid} identify its equation, law and grid states.
    \texttt{actual\allowbreak{}Chain\_\allowbreak{}law} identifies the comparison chain with the
    original Gaussian ULA laws.
\end{leanbox}
\section{\texorpdfstring{Proving the actual reversible Markov semigroup}{Proving the actual reversible Markov semigroup}}


Evaluating an observable at \(X_t\) is an isometry from
    \(L^2(\pi)\) into the common probability space. Compress evaluation at
    \(t\) against evaluation at zero to define \(P_t\).
    Future Brownian increments are independent of the complete
    initial/\allowbreak{}Brownian past. Finite shifted Euler maps and their genuine
    \(L^2\) limits prove


    \[\mathbb E[f(X_{a+t})\mid\mathcal F_a]=P_tf(X_a),
    \qquad P_{s+t}=P_sP_t.\]


For reversibility, realize MALA with the same Brownian Gaussian
    proposals and independent uniform accept/\allowbreak{}reject variables. Its anchored
    finite-step laws are exactly reversible. The proved rejection-distance
    estimate and Euler convergence give convergence in probability to the
    actual diffusion; stationarity and density upgrade every observable to
    strong \(L^2\) convergence. Self-adjointness passes to \(P_t\).
    Its contraction and strong continuity follow from actual evaluation
    isometries and state continuity.

\begin{leanbox}{Lean: \texttt{actualEvaluation\_\allowbreak{}markov\_\allowbreak{}pastAt}, \texttt{op\_\allowbreak{}isSelfAdjoint}, \texttt{semigroup}}
\leansource{\texttt{actualEvaluation\_\allowbreak{}markov\_\allowbreak{}pastAt} \quad WassersteinMixing/\allowbreak{}ActualMarkov.lean, lines 120–145}
\begin{leancode}
\leanline{0}{theorem\ actualEvaluation\_markov\_pastAt\ (a\ :\ ℝ≥0)\ (t\ :\ ℝ)\ (ht\ :\ 0\ ≤\ t)}
\leanline{4}{(f\ :\ Lp\ H\ 2\ (ULA.gibbsMeasure\ V))\ :}
\leanline{4}{(condExpL2\ H\ ℝ\ (ULA.lpLangevinPast\_le\ (actualProbability\ P\ V)}
\leanline{6}{(CanonicalSemigroup.initial\ P\ V\ hV)\ (actualBrownianNoise\ P\ V\ hV)\ (a\ :\ ℝ))}
\leanline{6}{(CanonicalSemigroup.evaluation\ ({\leanfallback 𝕜}\ :=\ ℝ)\ P\ V\ hV\ (a\ +\ t)\ f)\ :\ Lp\ H\ 2\ (actualProbability\ P\ V))\ =}
\leanline{6}{CanonicalSemigroup.evaluation\ ({\leanfallback 𝕜}\ :=\ ℝ)\ P\ V\ hV\ a}
\leanline{8}{(CanonicalSemigroup.op\ ({\leanfallback 𝕜}\ :=\ ℝ)\ P\ V\ hV\ t\ f)}
\end{leancode}
\leansource{\texttt{op\_\allowbreak{}isSelfAdjoint} \quad WassersteinMixing/\allowbreak{}SemigroupReversible.lean, lines 70–87}
\begin{leancode}
\leanline{0}{theorem\ op\_isSelfAdjoint\ (P\ :\ Parameters)\ (V\ :\ E\ P.dimension\ →\ ℝ)}
\leanline{4}{(hV\ :\ PotentialAssumptions\ P\ V)\ (t\ :\ ℝ)\ (ht\ :\ 0\ ≤\ t)\ :}
\leanline{4}{IsSelfAdjoint\ (op\ ({\leanfallback 𝕜}\ :=\ {\leanfallback 𝕜})\ (F\ :=\ F)\ P\ V\ hV\ t)}
\end{leancode}
\leansource{\texttt{semigroup} \quad WassersteinMixing/\allowbreak{}SemigroupInstance.lean, lines 55–62}
\begin{leancode}
\leanline{0}{def\ semigroup\ (P\ :\ Parameters)\ (V\ :\ E\ P.dimension\ →\ ℝ)\ (hV\ :\ PotentialAssumptions\ P\ V)\ :}
\leanline{4}{ReversibleSemigroup\ (Lp\ F\ 2\ (ULA.gibbsMeasure\ V))\ where}
\leanline{2}{op\ :=\ op\ P\ V\ hV}
\leanline{2}{zero\ :=\ op\_zero\ P\ V\ hV}
\leanline{2}{add\ :=\ fun\ \{s\ t\}\ hs\ ht\ ↦\ op\_add\ P\ V\ hV\ s\ t\ hs\ ht}
\leanline{2}{selfAdjoint\ :=\ fun\ \{t\}\ ht\ ↦\ op\_isSelfAdjoint\ P\ V\ hV\ t\ ht}
\leanline{2}{norm\_le\ :=\ fun\ \{t\}\ \_\ f\ ↦\ op\_norm\_le\ P\ V\ hV\ t\ f}
\leanline{2}{continuous\ :=\ fun\ f\ ↦\ continuousOn\_op\ P\ V\ hV\ f}
\end{leancode}
\leanlabel{How the Lean proof follows the argument}
\texttt{actual\allowbreak{}Evaluation\_\allowbreak{}markov\_\allowbreak{}past\allowbreak{}At} proves the conditional
    identity. \texttt{op\_\allowbreak{}is\allowbreak{}Self\allowbreak{}Adjoint} uses actual reversible MALA
    approximants. \texttt{Canonical\allowbreak{}Semigroup.\allowbreak{}semigroup} assembles all
    proved fields; no abstract Markov or spectral realization is supplied.
\end{leanbox}
\section{\texorpdfstring{Constructing the bounded Poisson corrector without an eigenbasis}{Constructing the bounded Poisson corrector without an eigenbasis}}


On complex Gibbs observables define
    \(A_hf=\int_0^h(f-P_sf)\,ds\) and
    \(G_N=\sum_{n<N}P_{nh}A_h\).
    Self-adjointness and \(P_t=P_{t/2}^2\) imply
    \(0\preccurlyeq P_t\preccurlyeq P_s\preccurlyeq I\) for
    \(0\le s\le t\). Integration and telescoping give


    \[0\preccurlyeq G_N\preccurlyeq h(I-P_{Nh})\preccurlyeq hI.\]


The increasing bounded operators converge strongly to a genuine
    bounded operator \(G\). Passing the finite telescoping identity to the
    limit gives \((I-P_h)G=A_h\), while
    \(\|G\|\le h\). Commutation and positivity of
    \((h^2I-G^2)(I-P_h)\) give the energy domination


    \[\langle Gf,(I-P_h)Gf\rangle
      \le h^2\langle f,(I-P_h)f\rangle.\]


Pointwise real-to-complex embedding commutes with actual evaluation
    and adjunction. Thus the constructed corrector preserves real
    observables. The extracted real semigroup operator is the actual real
    transition. Set \(g=G\nabla V\).

\begin{leanbox}{Lean: \texttt{corrector\_\allowbreak{}poisson}, \texttt{corrector\_\allowbreak{}energy\_\allowbreak{}le}, \texttt{actualPoissonProcess}}
\leansource{\texttt{corrector\_\allowbreak{}poisson} \quad WassersteinMixing/\allowbreak{}SemigroupPoisson.lean, lines 355–367}
\begin{leancode}
\leanline{0}{theorem\ corrector\_poisson\ (h\ :\ ℝ)\ (hh\ :\ 0\ ≤\ h)\ (x\ :\ H)\ :}
\leanline{4}{S.corrector\ h\ hh\ x\ -\ S.op\ h\ (S.corrector\ h\ hh\ x)\ =\ S.defect\ h\ x}
\end{leancode}
\leansource{\texttt{corrector\_\allowbreak{}energy\_\allowbreak{}le} \quad WassersteinMixing/\allowbreak{}SemigroupPoisson.lean, lines 392–402}
\begin{leancode}
\leanline{0}{theorem\ corrector\_energy\_le\ (h\ :\ ℝ)\ (hh\ :\ 0\ ≤\ h)\ \{s\ :\ ℝ\}\ (hs\ :\ 0\ ≤\ s)\ (x\ :\ H)\ :}
\leanline{4}{inner\ ℝ\ (S.corrector\ h\ hh\ x)}
\leanline{6}{(S.corrector\ h\ hh\ x\ -\ S.op\ s\ (S.corrector\ h\ hh\ x))\ ≤}
\leanline{6}{h\ \textasciicircum{}\ 2\ *\ inner\ ℝ\ x\ (x\ -\ S.op\ s\ x)}
\end{leancode}
\leansource{\texttt{actualPoissonProcess} \quad WassersteinMixing/\allowbreak{}ActualPoisson.lean, lines 81–88}
\begin{leancode}
\leanline{0}{def\ actualPoissonProcess\ (P\ :\ Parameters)\ (V\ :\ E\ P.dimension\ →\ ℝ)}
\leanline{4}{(hV\ :\ PotentialAssumptions\ P\ V)\ :}
\leanline{4}{PoissonProcess\ (Lp\ (E\ P.dimension)\ 2\ (ULA.gibbsMeasure\ V))\ (E\ P.dimension)}
\leanline{6}{(actualProbability\ P\ V)\ P.step\ P.step\_pos\ :=}
\leanline{2}{semigroupPoissonProcess\ (actualProbability\ P\ V)\ (actualObservableSemigroup\ P\ V\ hV)}
\leanline{4}{(actualRealEmbedding\ P\ V)\ (actualObservableSemigroup\_preservesReal\ P\ V\ hV)}
\leanline{4}{P.step\ P.step\_pos\ (actualFiltration\ P\ V\ hV)\ (actualGridEvaluation\ P\ V\ hV)}
\leanline{4}{(actualGridEvaluation\_adapted\ P\ V\ hV)\ (actualGridEvaluation\_markov\ P\ V\ hV)}
\end{leancode}
\leanlabel{How the Lean proof follows the argument}
\texttt{corrector\_\allowbreak{}norm\_\allowbreak{}le},
    \texttt{corrector\_\allowbreak{}poisson} and \texttt{corrector\_\allowbreak{}energy\_\allowbreak{}le}
    prove the required conclusions from the constructed semigroup.
    \texttt{actual\allowbreak{}Poisson\allowbreak{}Process} transfers them to the actual adapted
    grid evaluations.
\end{leanbox}
\section{\texorpdfstring{Deriving the actual score increment and defect bounds}{Deriving the actual score increment and defect bounds}}


Gibbs integration gives
    \(\|\nabla V\|_{L^2(\pi)}^2\le\beta d\).
    The actual integral equation and Brownian second moment imply


    \[\|X_{kh+s}-X_{kh}\|_{L^2}^2
      \le d(2s^2\beta+4s).\]


Gradient Lipschitzness and \(s\beta\le1\) then give


    \[\|\nabla V(X_{kh+s})-\nabla V(X_{kh})\|_{L^2}^2
      \le6s\beta^2d.\]


The true local defect is
    \(D_k=\int_0^h(\nabla V(X_{kh})-\nabla V(X_{kh+s}))\,ds\).
    Integral Cauchy–Schwarz gives
    \(\|D_k\|_{L^2}^2\le3h^3\beta^2d\).
    The corrector bound yields
    \(\|g(X_{kh})\|_{L^2}^2\le h^2\beta d\), and its energy bound
    gives a corrector-increment budget \(6h^3\beta^2d\).

\begin{leanbox}{Lean: \texttt{actualObservableScore\_\allowbreak{}energy}, \texttt{actualScoreIncrement\_\allowbreak{}energy}, \texttt{actualDefect\_\allowbreak{}energy}}
\leansource{\texttt{actualObservableScore\_\allowbreak{}energy} \quad WassersteinMixing/\allowbreak{}ActualScoreObservable.lean, lines 64–69}
\begin{leancode}
\leanline{0}{theorem\ actualObservableScore\_energy\ :}
\leanline{4}{‖actualObservableScore\ P\ V\ hV‖\ \textasciicircum{}\ 2\ ≤\ P.beta\ *\ P.dimension}
\end{leancode}
\leansource{\texttt{actualScoreIncrement\_\allowbreak{}energy} \quad WassersteinMixing/\allowbreak{}ActualLocalError.lean, lines 117–126}
\begin{leancode}
\leanline{0}{theorem\ actualScoreIncrement\_energy\ (n\ :\ ℕ)\ (s\ :\ ℝ)\ (hs\ :\ s\ ∈\ Icc\ 0\ P.step)\ :}
\leanline{4}{‖actualScoreIncrement\ P\ V\ hV\ n\ s‖\ \textasciicircum{}\ 2\ ≤\ 6\ *\ s\ *\ P.beta\ \textasciicircum{}\ 2\ *\ P.dimension}
\end{leancode}
\leansource{\texttt{actualDefect\_\allowbreak{}energy} \quad WassersteinMixing/\allowbreak{}ActualLocalError.lean, lines 154–158}
\begin{leancode}
\leanline{0}{theorem\ actualDefect\_energy\ (n\ :\ ℕ)\ :}
\leanline{4}{‖actualDefect\ P\ V\ hV\ n‖\ \textasciicircum{}\ 2\ ≤\ 3\ *\ P.step\ \textasciicircum{}\ 3\ *\ P.beta\ \textasciicircum{}\ 2\ *\ P.dimension}
\end{leancode}
\leanlabel{How the Lean proof follows the argument}
\texttt{actual\allowbreak{}Observable\allowbreak{}Score\_\allowbreak{}energy},
    \texttt{actual\allowbreak{}Score\allowbreak{}Increment\_\allowbreak{}energy} and
    \texttt{actual\allowbreak{}Defect\_\allowbreak{}energy} derive these constants for the genuine
    diffusion. They are not inputs to the final theorem.
\end{leanbox}
\section{\texorpdfstring{Identifying and centering the actual martingale residual}{Identifying and centering the actual martingale residual}}


The proved conditional Markov identity and commutation with
    the Bochner time integral identify the actual defect mean:


    \[\mathbb E[D_k\mid\mathcal F_{kh}]
      =\int_0^h(I-P_s)\nabla V(X_{kh})\,ds
      =(I-P_h)g(X_{kh}).\]


Hence
    \(M_{k+1}=D_k-[g(X_{kh})-g(X_{(k+1)h})]\) is conditionally
    centered. The separate budgets \(3h^3\beta^2d\) and
    \(6h^3\beta^2d\), together with
    \(\|u-v\|^2\le2\|u\|^2+2\|v\|^2\), give


    \[\|M_{k+1}\|_{L^2}^2\le18h^3\beta^2d.\]


The predictable shifted drift belongs to the actual Brownian past,
    so conditional orthogonality removes its cross term with \(M_{k+1}\).

\begin{leanbox}{Lean: \texttt{actualDefect\_\allowbreak{}conditionalMean}, \texttt{martingaleAt\_\allowbreak{}energy\_\allowbreak{}le}, \texttt{shifted\_\allowbreak{}energy\_\allowbreak{}step}}
\leansource{\texttt{actualDefect\_\allowbreak{}conditionalMean} \quad WassersteinMixing/\allowbreak{}ActualDefectMean.lean, lines 81–92}
\begin{leancode}
\leanline{0}{theorem\ actualDefect\_conditionalMean\ (n\ :\ ℕ)\ :}
\leanline{4}{(condExpL2\ (E\ P.dimension)\ ℝ\ (actualPast\_le\ P\ V\ hV\ n)}
\leanline{6}{(actualDefect\ P\ V\ hV\ n)\ :\ Lp\ (E\ P.dimension)\ 2\ (actualProbability\ P\ V))\ =}
\leanline{6}{actualGridEvaluation\ P\ V\ hV\ n}
\leanline{8}{(∫\ s\ in\ 0..P.step,\ actualObservableScore\ P\ V\ hV\ -}
\leanline{10}{CanonicalSemigroup.op\ ({\leanfallback 𝕜}\ :=\ ℝ)\ P\ V\ hV\ s\ (actualObservableScore\ P\ V\ hV))}
\end{leancode}
\leansource{\texttt{martingaleAt\_\allowbreak{}energy\_\allowbreak{}le} \quad WassersteinMixing/\allowbreak{}PoissonProcess.lean, lines 148–156}
\begin{leancode}
\leanline{0}{theorem\ martingaleAt\_energy\_le\ (f\ :\ H)\ (D\ :\ ℕ\ →\ Lp\ E\ 2\ Q)\ (β\ d\ :\ ℝ)\ (n\ :\ ℕ)}
\leanline{4}{(hdefect\ :\ ‖D\ n‖\ \textasciicircum{}\ 2\ ≤\ 3\ *\ h\ \textasciicircum{}\ 3\ *\ β\ \textasciicircum{}\ 2\ *\ d)}
\leanline{4}{(hincrement\ :\ ‖S.eval\ (n\ +\ 1)\ f\ -\ S.eval\ n\ f‖\ \textasciicircum{}\ 2\ ≤\ 6\ *\ h\ *\ β\ \textasciicircum{}\ 2\ *\ d)\ :}
\leanline{4}{‖S.martingaleAt\ f\ D\ n‖\ \textasciicircum{}\ 2\ ≤\ 18\ *\ h\ \textasciicircum{}\ 3\ *\ β\ \textasciicircum{}\ 2\ *\ d}
\end{leancode}
\leansource{\texttt{shifted\_\allowbreak{}energy\_\allowbreak{}step} \quad WassersteinMixing/\allowbreak{}PoissonComparison.lean, lines 134–148}
\begin{leancode}
\leanline{0}{theorem\ shifted\_energy\_step\ (n\ :\ ℕ)\ :}
\leanline{4}{‖C.error\ (n\ +\ 1)\ +\ S.correctorAt\ score\ (n\ +\ 1)‖\ \textasciicircum{}\ 2\ ≤}
\leanline{6}{P.contraction\ *\ ‖C.error\ n\ +\ S.correctorAt\ score\ n‖\ \textasciicircum{}\ 2\ +\ P.energyForcing}
\end{leancode}
\leanlabel{How the Lean proof follows the argument}
\texttt{actual\allowbreak{}Defect\_\allowbreak{}conditional\allowbreak{}Mean} identifies the true
    integral defect. The Poisson process derives martingale centering and
    its energy bound; the synchronous comparison then derives the shifted
    energy recurrence.
\end{leanbox}
\section{\texorpdfstring{Using the genuine averaged Hessian and shifted energy}{Using the genuine averaged Hessian and shifted energy}}


The literal segment average of the Hessian between the actual
    diffusion and ULA states is symmetric, lies between \(\alpha I\) and
    \(\beta I\), and maps their difference to the true gradient difference.
    Writing \(Z_k=X_{kh}-\widehat X_k\) and \(U_k=Z_k+g(X_{kh})\),
    the exact synchronous updates give


    \[U_{k+1}=(I-hA_k)U_k+hA_kg(X_{kh})+M_{k+1}.\]


The positive contraction inequality yields


    \[\|(I-hA_k)u+hA_kv\|^2
      \le(1-\alpha h)\|u\|^2+h\beta\|v\|^2.\]


After conditional orthogonality, the forcing constant is
    \(1+18=19\):


    \[\|U_{k+1}\|_{L^2}^2
      \le(1-\alpha h)\|U_k\|_{L^2}^2+19h^3\beta^2d.\]


Starting the two processes together gives the uniform floor
    \(19h^2\kappa\beta d\). The triangle inequality, with
    \(\sqrt{19}\le5\), gives
    \(\|Z_k\|_{L^2}\le6h\sqrt{\kappa\beta d}\).

\begin{leanbox}{Lean: \texttt{actual\_\allowbreak{}curvature\_\allowbreak{}error\_\allowbreak{}step}, \texttt{curvature\_\allowbreak{}weighted\_\allowbreak{}interpolation}, \texttt{shifted\_\allowbreak{}norm\_\allowbreak{}le\_\allowbreak{}bias}}
\leansource{\texttt{actual\_\allowbreak{}curvature\_\allowbreak{}error\_\allowbreak{}step} \quad WassersteinMixing/\allowbreak{}ActualCurvature.lean, lines 25–37}
\begin{leancode}
\leanline{0}{theorem\ actual\_curvature\_error\_step\ (n\ :\ ℕ)\ :}
\leanline{4}{actualReference\ P\ V\ hV\ (n\ +\ 1)\ -\ actualChain\ P\ V\ hV\ (n\ +\ 1)\ =}
\leanline{6}{(actualReference\ P\ V\ hV\ n\ -\ actualChain\ P\ V\ hV\ n)\ -}
\leanline{8}{P.step\ •\ actualCurvature\ P\ V\ hV\ n\ (actualReference\ P\ V\ hV\ n\ -\ actualChain\ P\ V\ hV\ n)\ +}
\leanline{10}{actualDefect\ P\ V\ hV\ n}
\end{leancode}
\leansource{\texttt{curvature\_\allowbreak{}weighted\_\allowbreak{}interpolation} \quad WassersteinMixing/\allowbreak{}ShiftedEnergy.lean, lines 72–93}
\begin{leancode}
\leanline{0}{theorem\ curvature\_weighted\_interpolation\ (P\ :\ Parameters)\ (A\ :\ H\ →L[ℝ]\ H)}
\leanline{4}{(hsym\ :\ ∀\ x\ y,\ inner\ ℝ\ (A\ x)\ y\ =\ inner\ ℝ\ x\ (A\ y))}
\leanline{4}{(hlower\ :\ ∀\ x,\ P.alpha\ *\ ‖x‖\ \textasciicircum{}\ 2\ ≤\ inner\ ℝ\ x\ (A\ x))}
\leanline{4}{(hupper\ :\ ∀\ x,\ inner\ ℝ\ x\ (A\ x)\ ≤\ P.beta\ *\ ‖x‖\ \textasciicircum{}\ 2)\ (u\ v\ :\ H)\ :}
\leanline{4}{‖u\ -\ P.step\ •\ A\ u\ +\ P.step\ •\ A\ v‖\ \textasciicircum{}\ 2\ ≤}
\leanline{6}{P.contraction\ *\ ‖u‖\ \textasciicircum{}\ 2\ +\ P.step\ *\ P.beta\ *\ ‖v‖\ \textasciicircum{}\ 2}
\end{leancode}
\leansource{\texttt{shifted\_\allowbreak{}norm\_\allowbreak{}le\_\allowbreak{}bias} \quad WassersteinMixing/\allowbreak{}Recurrence.lean, lines 88–121}
\begin{leancode}
\leanline{0}{theorem\ shifted\_norm\_le\_bias\ \{H\ :\ Type*\}\ [NormedAddCommGroup\ H]}
\leanline{4}{(P\ :\ Parameters)\ (z\ g\ :\ H)}
\leanline{4}{(hu\ :\ ‖z\ +\ g‖\ \textasciicircum{}\ 2\ ≤\ P.energyFloor)}
\leanline{4}{(hg\ :\ ‖g‖\ \textasciicircum{}\ 2\ ≤\ P.correctorBudget)\ :\ ‖z‖\ ≤\ P.bias}
\end{leancode}
\leanlabel{How the Lean proof follows the argument}
\texttt{actual\_\allowbreak{}curvature\_\allowbreak{}error\_\allowbreak{}step} uses the genuine
    averaged-Hessian secant identity.
    \texttt{curvature\_\allowbreak{}weighted\_\allowbreak{}interpolation} proves the deterministic
    inequality. \texttt{shifted\_\allowbreak{}energy\_\allowbreak{}uniform} and
    \texttt{shifted\_\allowbreak{}norm\_\allowbreak{}le\_\allowbreak{}bias} derive the floor and literal
    coefficient six.
\end{leanbox}
\section{\texorpdfstring{Constructing the actual Gaussian ULA invariant law}{Constructing the actual Gaussian ULA invariant law}}


Let \(F(x)=x-h\nabla V(x)\), whose Lipschitz factor is
    \(q=1-\alpha h<1\). On an iid Gaussian past
    \(\omega=(\omega_0,\omega_1,\ldots)\), let \(\tau\) drop its first entry.
    The tail preserves the product law and is independent of \(\omega_0\).
    On the complete \(L^2\) space define


    \[\Phi(X)(\omega)=F(X(\tau\omega))+\sqrt{2h}\,\omega_0.\]


This is a contraction with factor \(q\), so Banach's theorem
    constructs a square-integrable fixed point. Its law \(\widehat\pi\)
    is a probability measure with finite second moment. The actual
    recursion and head/\allowbreak{}tail independence prove invariance under the
    original independent-Gaussian ULA transition. The proved two-law
    contraction gives
    \(W_2(T^k\mu,\widehat\pi)\le q^kW_2(\mu,\widehat\pi)\to0\).

\begin{leanbox}{Lean: \texttt{invariantLaw\_\allowbreak{}invariant}, \texttt{invariantLaw\_\allowbreak{}attraction}}
\leansource{\texttt{invariantLaw\_\allowbreak{}invariant} \quad WassersteinMixing/\allowbreak{}Invariant.lean, lines 236–242}
\begin{leancode}
\leanline{0}{theorem\ invariantLaw\_invariant\ (P\ :\ Parameters)\ (V\ :\ E\ →\ ℝ)}
\leanline{4}{(hV\ :\ PotentialAssumptions\ P\ V)\ :}
\leanline{4}{ULA.ulaLawStep\ V\ P.step\ (invariantLaw\ P\ V\ hV)\ =\ invariantLaw\ P\ V\ hV}
\end{leancode}
\leansource{\texttt{invariantLaw\_\allowbreak{}attraction} \quad WassersteinMixing/\allowbreak{}Invariant.lean, lines 256–265}
\begin{leancode}
\leanline{0}{theorem\ invariantLaw\_attraction\ (P\ :\ Parameters)\ (V\ :\ E\ →\ ℝ)}
\leanline{4}{(hV\ :\ PotentialAssumptions\ P\ V)\ (μ\ :\ Measure\ E)\ [IsProbabilityMeasure\ μ]}
\leanline{4}{(hμ\ :\ MemLp\ (id\ :\ E\ →\ E)\ 2\ μ)\ :}
\leanline{4}{Filter.Tendsto\ (fun\ k\ :\ ℕ\ ↦}
\leanline{6}{ULA.wasserstein2Real\ ((ULA.ulaLawStep\ V\ P.step)\textasciicircum{}[k]\ μ)\ (invariantLaw\ P\ V\ hV))}
\leanline{6}{Filter.atTop\ ({\leanfallback 𝓝}\ 0)}
\end{leancode}
\leanlabel{How the Lean proof follows the argument}
\texttt{invariant\allowbreak{}Law} specializes the genuine random
    fixed-point construction. \texttt{invariant\allowbreak{}Law\_\allowbreak{}invariant} and
    \texttt{invariant\allowbreak{}Law\_\allowbreak{}attraction} prove the exact stationary and
    limiting properties needed in Theorem 1.
\end{leanbox}
\section{\texorpdfstring{Passing the actual stationary comparison to transport}{Passing the actual stationary comparison to transport}}


The synchronous trajectories give actual couplings of
    \(\pi\) and each stationary-start Gaussian ULA law. Their uniform
    \(L^2\) bound therefore bounds the actual quadratic transport
    infimum. To prove its triangle inequality, disintegrate two couplings
    with the same middle marginal and construct a genuine triple law.
    Minkowski bounds its endpoint cost; taking both coupling infima gives


    \[W_2(\mu,\xi)\le W_2(\mu,\nu)+W_2(\nu,\xi).\]


Apply this with the stationary-start ULA law between \(\pi\) and
    \(\widehat\pi\), then use the proved attraction to the constructed
    invariant law. Theorem 1 follows. Combining it with actual ULA
    contraction gives the finite-time inequality of Corollary 1.1.
    Second moments prove finiteness before distances are converted to
    real numbers.

\begin{leanbox}{Lean: \texttt{stationary\_\allowbreak{}transport\_\allowbreak{}bound}, \texttt{wasserstein2Real\_\allowbreak{}triangle}, \texttt{ula\_\allowbreak{}iterate\_\allowbreak{}wasserstein2Real\_\allowbreak{}le}}
\leansource{\texttt{stationary\_\allowbreak{}transport\_\allowbreak{}bound} \quad WassersteinMixing/\allowbreak{}PoissonComparison.lean, lines 163–167}
\begin{leancode}
\leanline{0}{theorem\ stationary\_transport\_bound\ [BorelSpace\ E]\ [SecondCountableTopology\ E]\ (n\ :\ ℕ)\ :}
\leanline{4}{ULA.wasserstein2Real\ (Q.map\ (C.reference\ 0))\ (Q.map\ (C.chain\ n))\ ≤\ P.bias}
\end{leancode}
\leansource{\texttt{wasserstein2Real\_\allowbreak{}triangle} \quad WassersteinMixing/\allowbreak{}TransportGeometry.lean, lines 175–181}
\begin{leancode}
\leanline{0}{theorem\ wasserstein2Real\_triangle\ (μ\ ν\ ξ\ :\ Measure\ H)}
\leanline{4}{[IsProbabilityMeasure\ μ]\ [IsProbabilityMeasure\ ν]\ [IsProbabilityMeasure\ ξ]}
\leanline{4}{(hμν\ :\ ULA.wasserstein2\ μ\ ν\ ≠\ ⊤)\ (hνξ\ :\ ULA.wasserstein2\ ν\ ξ\ ≠\ ⊤)\ :}
\leanline{4}{ULA.wasserstein2Real\ μ\ ξ\ ≤\ ULA.wasserstein2Real\ μ\ ν\ +\ ULA.wasserstein2Real\ ν\ ξ}
\end{leancode}
\leansource{\texttt{ula\_\allowbreak{}iterate\_\allowbreak{}wasserstein2Real\_\allowbreak{}le} \quad WassersteinMixing/\allowbreak{}Transport.lean, lines 109–121}
\begin{leancode}
\leanline{0}{theorem\ ula\_iterate\_wasserstein2Real\_le\ (P\ :\ Parameters)\ (V\ :\ H\ →\ ℝ)}
\leanline{4}{(hV\ :\ PotentialAssumptions\ P\ V)\ (μ\ ν\ :\ Measure\ H)}
\leanline{4}{[IsProbabilityMeasure\ μ]\ [IsProbabilityMeasure\ ν]}
\leanline{4}{(hfinite\ :\ ULA.wasserstein2\ μ\ ν\ ≠\ ⊤)\ (k\ :\ ℕ)\ :}
\leanline{4}{ULA.wasserstein2Real\ ((ULA.ulaLawStep\ V\ P.step)\textasciicircum{}[k]\ μ)}
\leanline{6}{((ULA.ulaLawStep\ V\ P.step)\textasciicircum{}[k]\ ν)\ ≤}
\leanline{6}{P.contraction\ \textasciicircum{}\ k\ *\ ULA.wasserstein2Real\ μ\ ν}
\end{leancode}
\leanlabel{How the Lean proof follows the argument}
\texttt{stationary\_\allowbreak{}transport\_\allowbreak{}bound} uses the actual
    synchronous coupling. \texttt{glue\allowbreak{}Couplings} constructs the
    triple law and \texttt{wasserstein2\allowbreak{}Real\_\allowbreak{}triangle} proves the
    finite-distance triangle inequality. The final theorem instantiates
    the comparison from the source Hessian assumptions.
\end{leanbox}
\section{\texorpdfstring{Recovering the exact scale-invariant iteration count}{Recovering the exact scale-invariant iteration count}}


For \(d>0\) and \(\varepsilon>0\), write


    \[m=\max\left(1,\frac{12\sqrt d}{\varepsilon}\right),\qquad
    h=\frac1{\beta m},\qquad K=\frac\beta\alpha m.\]


This is the source minimum formula and gives \(\alpha hK=1\).
    Since
    \(\sqrt\alpha\sqrt{(\beta/\alpha)\beta d}=\beta\sqrt d\),
    the scaled bias is at most \(\varepsilon/2\).
    For \(t=\alpha h\in(0,1]\), use \((1-t)^k\le e^{-tk}\).
    The displayed source iteration restriction makes the scaled transient
    at most \(\varepsilon/2\). Adding the two budgets proves the
    literal accuracy conclusion, including zero initial error.

\begin{leanbox}{Lean: \texttt{paperStep\_\allowbreak{}eq\_\allowbreak{}inv}, \texttt{paper\_\allowbreak{}complexity\_\allowbreak{}bound}, \texttt{accuracy\_\allowbreak{}bound}}
\leansource{\texttt{paperStep\_\allowbreak{}eq\_\allowbreak{}inv} \quad WassersteinMixing/\allowbreak{}Complexity.lean, lines 27–40}
\begin{leancode}
\leanline{0}{theorem\ paperStep\_eq\_inv\ (β\ d\ ε\ :\ ℝ)\ (hβ\ :\ 0\ <\ β)\ (hd\ :\ 0\ <\ d)\ (hε\ :\ 0\ <\ ε)\ :}
\leanline{4}{paperStep\ β\ d\ ε\ =\ 1\ /\ (β\ *\ max\ 1\ (12\ *\ Real.sqrt\ d\ /\ ε))}
\end{leancode}
\leansource{\texttt{paper\_\allowbreak{}complexity\_\allowbreak{}bound} \quad WassersteinMixing/\allowbreak{}Complexity.lean, lines 118–142}
\begin{leancode}
\leanline{0}{theorem\ paper\_complexity\_bound\ (α\ β\ d\ R\ ε\ :\ ℝ)\ (hα\ :\ 0\ <\ α)\ (hαβ\ :\ α\ ≤\ β)}
\leanline{4}{(hd\ :\ 0\ <\ d)\ (hR\ :\ 0\ ≤\ R)\ (hε\ :\ 0\ <\ ε)\ (k\ :\ ℕ)}
\leanline{4}{(hk\ :\ iterationPrefactor\ α\ β\ d\ ε\ *}
\leanline{6}{Real.log\ (1\ +\ 2\ *\ Real.sqrt\ α\ *\ R\ /\ ε)\ ≤\ (k\ :\ ℝ))\ :}
\leanline{4}{Real.sqrt\ α\ *\ ((1\ -\ α\ *\ paperStep\ β\ d\ ε)\ \textasciicircum{}\ k\ *\ R\ +}
\leanline{6}{6\ *\ paperStep\ β\ d\ ε\ *\ Real.sqrt\ ((β\ /\ α)\ *\ β\ *\ d))\ ≤\ ε}
\end{leancode}
\leansource{\texttt{accuracy\_\allowbreak{}bound} \quad WassersteinMixing/\allowbreak{}Theorem.lean, lines 44–54}
\begin{leancode}
\leanline{0}{theorem\ accuracy\_bound\ (P\ :\ Parameters)\ (V\ :\ E\ P.dimension\ →\ ℝ)}
\leanline{4}{(hA\ :\ HessianAssumptions\ P\ V)\ (μ\ :\ Measure\ (E\ P.dimension))\ [IsProbabilityMeasure\ μ]}
\leanline{4}{(hμ\ :\ MemLp\ (id\ :\ E\ P.dimension\ →\ E\ P.dimension)\ 2\ μ)\ (ε\ :\ ℝ)\ (hε\ :\ 0\ <\ ε)}
\leanline{4}{(hstep\ :\ P.step\ =\ paperStep\ P.beta\ P.dimension\ ε)\ (n\ :\ ℕ)}
\leanline{4}{(hn\ :\ iterationPrefactor\ P.alpha\ P.beta\ P.dimension\ ε\ *}
\leanline{6}{Real.log\ (1\ +\ 2\ *\ Real.sqrt\ P.alpha\ *}
\leanline{8}{ULA.wasserstein2Real\ μ\ (invariantLaw\ P\ V\ hA.toPotentialAssumptions)\ /\ ε)\ ≤\ (n\ :\ ℝ))\ :}
\leanline{4}{Real.sqrt\ P.alpha\ *}
\leanline{6}{ULA.wasserstein2Real\ ((ULA.ulaLawStep\ V\ P.step)\textasciicircum{}[n]\ μ)\ (ULA.gibbsMeasure\ V)\ ≤\ ε}
\end{leancode}
\leanlabel{How the Lean proof follows the argument}
\texttt{paper\allowbreak{}Step\_\allowbreak{}eq\_\allowbreak{}inv} identifies the two step
    formulas. \texttt{paper\_\allowbreak{}complexity\_\allowbreak{}bound} proves their exact
    scalar consequence. \texttt{accuracy\_\allowbreak{}bound} applies it to the
    actual Gaussian chain and constructed invariant law.
\end{leanbox}
\chapter{Supporting proofs}
The following sections give the definitions and prove the lemmas used in the convergence arguments.

\section{\texorpdfstring{Actual stationary-pair transition operators}{Actual stationary-pair transition operators}}
\label{proof-module-WassersteinMixing-Semigroup-lean}


For a stationary pair law \(J\) with both marginals \(\mu\), pullback by the first and second coordinates gives actual isometries \(U_0,U_1:L^2(\mu;F)\to L^2(J;F)\). The adjoint compression \(P=U_0^*U_1\) is a genuine contraction. Reversal symmetry of the pair law makes it self-adjoint.

\subsection{\texorpdfstring{Constructing the contraction from the two actual marginals}{Constructing the contraction from the two actual marginals}}
\label{report-step-WassersteinMixing-Semigroup-lean-1}


\textbf{Claim.} The coordinate evaluations preserve norms and

\[\langle Pf,g\rangle=\langle U_1f,U_0g\rangle,\qquad\|Pf\|\le\|f\|.\]

\textbf{Proof.} Each marginal equality makes the corresponding coordinate map measure-preserving, hence its pullback is an isometry. Define the transition by adjoint compression. The adjoint identity proves the pairing formula. The first evaluation and its adjoint have norm at most one, while the second evaluation preserves norms, giving the pointwise and operator contraction bounds.

\begin{leanbox}{Lean: \texttt{evalFst}, \texttt{evalSnd}, \texttt{transition} and 3 more}
\leansource{\texttt{evalFst} \quad WassersteinMixing/\allowbreak{}Semigroup.lean, lines 28–29}
\begin{leancode}
\leanline{0}{def\ evalFst\ :\ Lp\ F\ 2\ μ\ →ₗᵢ[{\leanfallback 𝕜}]\ Lp\ F\ 2\ J\ :=}
\leanline{2}{Lp.compMeasurePreservingₗᵢ\ {\leanfallback 𝕜}\ Prod.fst\ ⟨measurable\_fst,\ hfst⟩}
\end{leancode}
\leansource{\texttt{evalSnd} \quad WassersteinMixing/\allowbreak{}Semigroup.lean, lines 32–33}
\begin{leancode}
\leanline{0}{def\ evalSnd\ :\ Lp\ F\ 2\ μ\ →ₗᵢ[{\leanfallback 𝕜}]\ Lp\ F\ 2\ J\ :=}
\leanline{2}{Lp.compMeasurePreservingₗᵢ\ {\leanfallback 𝕜}\ Prod.snd\ ⟨measurable\_snd,\ hsnd⟩}
\end{leancode}
\leansource{\texttt{transition} \quad WassersteinMixing/\allowbreak{}Semigroup.lean, lines 37–39}
\begin{leancode}
\leanline{0}{def\ transition\ :\ Lp\ F\ 2\ μ\ →L[{\leanfallback 𝕜}]\ Lp\ F\ 2\ μ\ :=}
\leanline{2}{(evalFst\ ({\leanfallback 𝕜}\ :=\ {\leanfallback 𝕜})\ μ\ J\ hfst).toContinuousLinearMap.adjoint\ ∘L}
\leanline{4}{(evalSnd\ ({\leanfallback 𝕜}\ :=\ {\leanfallback 𝕜})\ μ\ J\ hsnd).toContinuousLinearMap}
\end{leancode}
\leansource{\texttt{transition\_\allowbreak{}inner} \quad WassersteinMixing/\allowbreak{}Semigroup.lean, lines 42–46}
\begin{leancode}
\leanline{0}{theorem\ transition\_inner\ (f\ g\ :\ Lp\ F\ 2\ μ)\ :}
\leanline{4}{inner\ {\leanfallback 𝕜}\ (transition\ ({\leanfallback 𝕜}\ :=\ {\leanfallback 𝕜})\ μ\ J\ hfst\ hsnd\ f)\ g\ =}
\leanline{6}{inner\ {\leanfallback 𝕜}\ (evalSnd\ ({\leanfallback 𝕜}\ :=\ {\leanfallback 𝕜})\ μ\ J\ hsnd\ f)\ (evalFst\ ({\leanfallback 𝕜}\ :=\ {\leanfallback 𝕜})\ μ\ J\ hfst\ g)}
\end{leancode}
\leansource{\texttt{transition\_\allowbreak{}norm\_\allowbreak{}le} \quad WassersteinMixing/\allowbreak{}Semigroup.lean, lines 49–59}
\begin{leancode}
\leanline{0}{theorem\ transition\_norm\_le\ (f\ :\ Lp\ F\ 2\ μ)\ :}
\leanline{4}{‖transition\ ({\leanfallback 𝕜}\ :=\ {\leanfallback 𝕜})\ μ\ J\ hfst\ hsnd\ f‖\ ≤\ ‖f‖}
\end{leancode}
\leansource{\texttt{transition\_\allowbreak{}opNorm\_\allowbreak{}le} \quad WassersteinMixing/\allowbreak{}Semigroup.lean, lines 62–65}
\begin{leancode}
\leanline{0}{theorem\ transition\_opNorm\_le\ :\ ‖transition\ ({\leanfallback 𝕜}\ :=\ {\leanfallback 𝕜})\ (F\ :=\ F)\ μ\ J\ hfst\ hsnd‖\ ≤\ 1}
\end{leancode}
\leanlabel{How the Lean proof follows the argument}
\texttt{eval\allowbreak{}Fst} and \texttt{eval\allowbreak{}Snd} are actual measure-preserving pullbacks. \texttt{transition} is their continuous adjoint composition. The correlation and norm estimates use the defining adjoint identity and the evaluation isometries.
\end{leanbox}
\subsection{\texorpdfstring{Reversal symmetry proves self-adjointness}{Reversal symmetry proves self-adjointness}}
\label{report-step-WassersteinMixing-Semigroup-lean-2}


\textbf{Claim.} If swapping the two coordinates preserves \(J\), then \(\langle Pf,g\rangle=\langle f,Pg\rangle\), so \(P\) is self-adjoint.



\textbf{Proof.} Pullback by the actual swap map is an isometry on the joint space and exchanges the two evaluation maps. Apply preservation of the inner product by this isometry to the correlation identity, then use the adjoint identity in the other direction. This is the symmetric pairing characterization of a bounded self-adjoint operator.

\begin{leanbox}{Lean: \texttt{reverse}, \texttt{reverse\_\allowbreak{}evalFst}, \texttt{reverse\_\allowbreak{}evalSnd} and 2 more}
\leansource{\texttt{reverse} \quad WassersteinMixing/\allowbreak{}Semigroup.lean, lines 70–71}
\begin{leancode}
\leanline{0}{def\ reverse\ :\ Lp\ F\ 2\ J\ →ₗᵢ[{\leanfallback 𝕜}]\ Lp\ F\ 2\ J\ :=}
\leanline{2}{Lp.compMeasurePreservingₗᵢ\ {\leanfallback 𝕜}\ Prod.swap\ ⟨measurable\_swap,\ hswap⟩}
\end{leancode}
\leansource{\texttt{reverse\_\allowbreak{}evalFst} \quad WassersteinMixing/\allowbreak{}Semigroup.lean, lines 74–86}
\begin{leancode}
\leanline{0}{theorem\ reverse\_evalFst\ (f\ :\ Lp\ F\ 2\ μ)\ :}
\leanline{4}{reverse\ ({\leanfallback 𝕜}\ :=\ {\leanfallback 𝕜})\ J\ hswap\ (evalFst\ ({\leanfallback 𝕜}\ :=\ {\leanfallback 𝕜})\ μ\ J\ hfst\ f)\ =\ evalSnd\ ({\leanfallback 𝕜}\ :=\ {\leanfallback 𝕜})\ μ\ J\ hsnd\ f}
\end{leancode}
\leansource{\texttt{reverse\_\allowbreak{}evalSnd} \quad WassersteinMixing/\allowbreak{}Semigroup.lean, lines 89–103}
\begin{leancode}
\leanline{0}{theorem\ reverse\_evalSnd\ (f\ :\ Lp\ F\ 2\ μ)\ :}
\leanline{4}{reverse\ ({\leanfallback 𝕜}\ :=\ {\leanfallback 𝕜})\ J\ hswap\ (evalSnd\ ({\leanfallback 𝕜}\ :=\ {\leanfallback 𝕜})\ μ\ J\ hsnd\ f)\ =\ evalFst\ ({\leanfallback 𝕜}\ :=\ {\leanfallback 𝕜})\ μ\ J\ hfst\ f}
\end{leancode}
\leansource{\texttt{transition\_\allowbreak{}symmetric} \quad WassersteinMixing/\allowbreak{}Semigroup.lean, lines 105–117}
\begin{leancode}
\leanline{0}{theorem\ transition\_symmetric\ (f\ g\ :\ Lp\ F\ 2\ μ)\ :}
\leanline{4}{inner\ {\leanfallback 𝕜}\ (transition\ ({\leanfallback 𝕜}\ :=\ {\leanfallback 𝕜})\ μ\ J\ hfst\ hsnd\ f)\ g\ =}
\leanline{6}{inner\ {\leanfallback 𝕜}\ f\ (transition\ ({\leanfallback 𝕜}\ :=\ {\leanfallback 𝕜})\ μ\ J\ hfst\ hsnd\ g)}
\end{leancode}
\leansource{\texttt{transition\_\allowbreak{}isSelfAdjoint} \quad WassersteinMixing/\allowbreak{}Semigroup.lean, lines 119–121}
\begin{leancode}
\leanline{0}{theorem\ transition\_isSelfAdjoint\ :\ IsSelfAdjoint\ (transition\ ({\leanfallback 𝕜}\ :=\ {\leanfallback 𝕜})\ (F\ :=\ F)\ μ\ J\ hfst\ hsnd)}
\end{leancode}
\leanlabel{How the Lean proof follows the argument}
\texttt{reverse\_\allowbreak{}eval\allowbreak{}Fst} and \texttt{reverse\_\allowbreak{}eval\allowbreak{}Snd} compare almost-everywhere pullback representatives. \texttt{transition\_\allowbreak{}symmetric} uses the actual reversal isometry and \texttt{transition\_\allowbreak{}is\allowbreak{}Self\allowbreak{}Adjoint} applies the symmetric-operator characterization.
\end{leanbox}
\subsection{\texorpdfstring{Compression on any common observation space}{Compression on any common observation space}}
\label{report-step-WassersteinMixing-Semigroup-lean-3}


\textbf{Claim.} The same construction works for any two isometries into a common Hilbert space: \(P=U_0^*U_1\) has the same pairing formula and norm bound. For equal observations, \(U_0^*U_0=I\).



\textbf{Proof.} The contraction argument uses only the two isometry properties. At equal observations, the adjoint composition of an isometry with itself is exactly the identity.

\begin{leanbox}{Lean: \texttt{compression}, \texttt{compression\_\allowbreak{}inner}, \texttt{compression\_\allowbreak{}norm\_\allowbreak{}le} and 1 more}
\leansource{\texttt{compression} \quad WassersteinMixing/\allowbreak{}Semigroup.lean, lines 132–133}
\begin{leancode}
\leanline{0}{def\ compression\ (initial\ current\ :\ A\ →ₗᵢ[{\leanfallback 𝕜}]\ B)\ :\ A\ →L[{\leanfallback 𝕜}]\ A\ :=}
\leanline{2}{initial.toContinuousLinearMap.adjoint\ ∘L\ current.toContinuousLinearMap}
\end{leancode}
\leansource{\texttt{compression\_\allowbreak{}inner} \quad WassersteinMixing/\allowbreak{}Semigroup.lean, lines 135–138}
\begin{leancode}
\leanline{0}{theorem\ compression\_inner\ (initial\ current\ :\ A\ →ₗᵢ[{\leanfallback 𝕜}]\ B)\ (f\ g\ :\ A)\ :}
\leanline{4}{inner\ {\leanfallback 𝕜}\ (compression\ initial\ current\ f)\ g\ =\ inner\ {\leanfallback 𝕜}\ (current\ f)\ (initial\ g)}
\end{leancode}
\leansource{\texttt{compression\_\allowbreak{}norm\_\allowbreak{}le} \quad WassersteinMixing/\allowbreak{}Semigroup.lean, lines 141–150}
\begin{leancode}
\leanline{0}{theorem\ compression\_norm\_le\ (initial\ current\ :\ A\ →ₗᵢ[{\leanfallback 𝕜}]\ B)\ (f\ :\ A)\ :}
\leanline{4}{‖compression\ initial\ current\ f‖\ ≤\ ‖f‖}
\end{leancode}
\leansource{\texttt{compression\_\allowbreak{}self} \quad WassersteinMixing/\allowbreak{}Semigroup.lean, lines 153–155}
\begin{leancode}
\leanline{0}{theorem\ compression\_self\ (initial\ :\ A\ →ₗᵢ[{\leanfallback 𝕜}]\ B)\ :}
\leanline{4}{compression\ initial\ initial\ =\ ContinuousLinearMap.id\ {\leanfallback 𝕜}\ A}
\end{leancode}
\leanlabel{How the Lean proof follows the argument}
\texttt{compression} abstracts only the isometric evaluation algebra. Its correlation and contraction lemmas are used by the actual Brownian observations; \texttt{compression\_\allowbreak{}self} supplies their exact time-zero identity.
\end{leanbox}
\section{\texorpdfstring{Continuity of actual stationary observable evaluations}{Continuity of actual stationary observable evaluations}}
\label{proof-module-WassersteinMixing-SemigroupContinuity-lean}


If stationary random states with common law \(\mu\) converge in measure, their evaluation maps converge strongly on every \(L^2(\mu)\) observable. The proof first treats bounded continuous functions and then uses their density and the common evaluation isometry.

\subsection{\texorpdfstring{The actual evaluation and its square norm}{The actual evaluation and its square norm}}
\label{report-step-WassersteinMixing-SemigroupContinuity-lean-1}


\textbf{Claim.} A state \(X\) with law \(\mu\) gives an evaluation isometry \(f\mapsto f(X)\). For bounded continuous \(\phi\),

\[\|\phi(X)-\phi(Y)\|_2^2=\int\|\phi(X(\omega))-\phi(Y(\omega))\|^2\,dQ.\]

\textbf{Proof.} Evaluation is pullback by a measure-preserving state map. Compare the exact almost-everywhere representatives and use the integral square-norm formula. The auxiliary implication from vanishing mean norm to convergence in measure follows from the integral Markov bound.

\begin{leanbox}{Lean: \texttt{tendstoInMeasure\_\allowbreak{}of\_\allowbreak{}integral\_\allowbreak{}norm\_\allowbreak{}sub}, \texttt{evaluation}, \texttt{lp\_\allowbreak{}norm\_\allowbreak{}sq\_\allowbreak{}eq\_\allowbreak{}integral} and 2 more}
\leansource{\texttt{tendstoInMeasure\_\allowbreak{}of\_\allowbreak{}integral\_\allowbreak{}norm\_\allowbreak{}sub} \quad WassersteinMixing/\allowbreak{}SemigroupContinuity.lean, lines 30–43}
\begin{leancode}
\leanline{0}{theorem\ tendstoInMeasure\_of\_integral\_norm\_sub\ \{I\ :\ Type*\}\ \{l\ :\ Filter\ I\}}
\leanline{4}{(X\ :\ I\ →\ Ω\ →\ E)\ (Y\ :\ Ω\ →\ E)}
\leanline{4}{(hint\ :\ ∀\ i,\ Integrable\ (fun\ ω\ ↦\ X\ i\ ω\ -\ Y\ ω)\ Q)}
\leanline{4}{(hlim\ :\ Tendsto\ (fun\ i\ ↦\ ∫\ ω,\ ‖X\ i\ ω\ -\ Y\ ω‖\ ∂Q)\ l\ ({\leanfallback 𝓝}\ 0))\ :}
\leanline{4}{TendstoInMeasure\ Q\ X\ l\ Y}
\end{leancode}
\leansource{\texttt{evaluation} \quad WassersteinMixing/\allowbreak{}SemigroupContinuity.lean, lines 46–49}
\begin{leancode}
\leanline{0}{def\ evaluation\ (X\ :\ Lp\ E\ 2\ Q)\ (hX\ :\ Q.map\ X\ =\ μ)\ :\ Lp\ F\ 2\ μ\ →ₗᵢ[{\leanfallback 𝕜}]\ Lp\ F\ 2\ Q\ :=}
\leanline{2}{Lp.compMeasurePreservingₗᵢ\ {\leanfallback 𝕜}\ X\ ⟨(Lp.stronglyMeasurable\ X).measurable,\ hX⟩}
\leanline{0}{\strut}
\leanline{0}{include\ {\leanfallback 𝕜}\ in}
\end{leancode}
\leansource{\texttt{lp\_\allowbreak{}norm\_\allowbreak{}sq\_\allowbreak{}eq\_\allowbreak{}integral} \quad WassersteinMixing/\allowbreak{}SemigroupContinuity.lean, lines 52–58}
\begin{leancode}
\leanline{0}{theorem\ lp\_norm\_sq\_eq\_integral\ (Z\ :\ Lp\ F\ 2\ Q)\ :}
\leanline{4}{‖Z‖\ \textasciicircum{}\ 2\ =\ ∫\ ω,\ ‖Z\ ω‖\ \textasciicircum{}\ 2\ ∂Q}
\end{leancode}
\leansource{\texttt{evaluation\_\allowbreak{}bcf\_\allowbreak{}coe} \quad WassersteinMixing/\allowbreak{}SemigroupContinuity.lean, lines 63–69}
\begin{leancode}
\leanline{0}{private\ theorem\ evaluation\_bcf\_coe\ (X\ :\ Lp\ E\ 2\ Q)\ (hX\ :\ Q.map\ X\ =\ μ)}
\leanline{4}{(φ\ :\ E\ →ᵇ\ F)\ :}
\leanline{4}{evaluation\ ({\leanfallback 𝕜}\ :=\ {\leanfallback 𝕜})\ μ\ Q\ X\ hX\ (BoundedContinuousFunction.toLp\ 2\ μ\ {\leanfallback 𝕜}\ φ)}
\leanline{6}{=ᵐ[Q]\ fun\ ω\ ↦\ φ\ (X\ ω)}
\end{leancode}
\leansource{\texttt{evaluation\_\allowbreak{}bcf\_\allowbreak{}norm\_\allowbreak{}sq} \quad WassersteinMixing/\allowbreak{}SemigroupContinuity.lean, lines 72–84}
\begin{leancode}
\leanline{0}{private\ theorem\ evaluation\_bcf\_norm\_sq\ (X\ Y\ :\ Lp\ E\ 2\ Q)}
\leanline{4}{(hX\ :\ Q.map\ X\ =\ μ)\ (hY\ :\ Q.map\ Y\ =\ μ)\ (φ\ :\ E\ →ᵇ\ F)\ :}
\leanline{4}{‖evaluation\ ({\leanfallback 𝕜}\ :=\ {\leanfallback 𝕜})\ μ\ Q\ X\ hX\ (BoundedContinuousFunction.toLp\ 2\ μ\ {\leanfallback 𝕜}\ φ)\ -}
\leanline{6}{evaluation\ ({\leanfallback 𝕜}\ :=\ {\leanfallback 𝕜})\ μ\ Q\ Y\ hY\ (BoundedContinuousFunction.toLp\ 2\ μ\ {\leanfallback 𝕜}\ φ)‖\ \textasciicircum{}\ 2\ =}
\leanline{6}{∫\ ω,\ ‖φ\ (X\ ω)\ -\ φ\ (Y\ ω)‖\ \textasciicircum{}\ 2\ ∂Q}
\end{leancode}
\leanlabel{How the Lean proof follows the argument}
\texttt{evaluation} is actual measure-preserving composition. \texttt{lp\_\allowbreak{}norm\_\allowbreak{}sq\_\allowbreak{}eq\_\allowbreak{}integral} takes the real part of the integrable inner-product formula. The bounded-continuous representative lemmas identify the norm of the actual evaluation difference with the displayed integral.
\end{leanbox}
\subsection{\texorpdfstring{Extending convergence from bounded continuous functions to all observables}{Extending convergence from bounded continuous functions to all observables}}
\label{report-step-WassersteinMixing-SemigroupContinuity-lean-2}


\textbf{Claim.} If \(X_i\to X\) in measure and all states have law \(\mu\), then \(\phi(X_i)\to\phi(X)\) in \(L^2\) for bounded continuous \(\phi\), and the same is true for every square-integrable observable.



\textbf{Proof.} Composition with a continuous bounded function preserves convergence in measure. The squared difference has a uniform constant bound, so dominated convergence in measure gives its integral tending to zero. For arbitrary \(f\), approximate it in \(L^2(\mu)\) by a bounded continuous function. All evaluation maps have norm one, so both approximation errors have the same uniform bound; the middle bounded-continuous evaluation difference tends to zero. Norm convergence of states implies convergence in measure, giving the final specializations.

\begin{leanbox}{Lean: \texttt{evaluation\_\allowbreak{}bcf\_\allowbreak{}tendsto\_\allowbreak{}inMeasure}, \texttt{evaluation\_\allowbreak{}tendsto\_\allowbreak{}inMeasure}, \texttt{evaluation\_\allowbreak{}bcf\_\allowbreak{}tendsto} and 1 more}
\leansource{\texttt{evaluation\_\allowbreak{}bcf\_\allowbreak{}tendsto\_\allowbreak{}inMeasure} \quad WassersteinMixing/\allowbreak{}SemigroupContinuity.lean, lines 89–127}
\begin{leancode}
\leanline{0}{theorem\ evaluation\_bcf\_tendsto\_inMeasure\ \{I\ :\ Type*\}\ \{l\ :\ Filter\ I\}\ [l.IsCountablyGenerated]}
\leanline{4}{(X\ :\ I\ →\ Lp\ E\ 2\ Q)\ (Y\ :\ Lp\ E\ 2\ Q)\ (hX\ :\ ∀\ i,\ Q.map\ (X\ i)\ =\ μ)}
\leanline{4}{(hY\ :\ Q.map\ Y\ =\ μ)}
\leanline{4}{(hlim\ :\ TendstoInMeasure\ Q\ (fun\ i\ ω\ ↦\ X\ i\ ω)\ l\ Y)\ (φ\ :\ E\ →ᵇ\ F)\ :}
\leanline{4}{Tendsto\ (fun\ i\ ↦\ evaluation\ ({\leanfallback 𝕜}\ :=\ {\leanfallback 𝕜})\ μ\ Q\ (X\ i)\ (hX\ i)}
\leanline{6}{(BoundedContinuousFunction.toLp\ 2\ μ\ {\leanfallback 𝕜}\ φ))\ l}
\leanline{6}{({\leanfallback 𝓝}\ (evaluation\ ({\leanfallback 𝕜}\ :=\ {\leanfallback 𝕜})\ μ\ Q\ Y\ hY\ (BoundedContinuousFunction.toLp\ 2\ μ\ {\leanfallback 𝕜}\ φ)))}
\end{leancode}
\leansource{\texttt{evaluation\_\allowbreak{}tendsto\_\allowbreak{}inMeasure} \quad WassersteinMixing/\allowbreak{}SemigroupContinuity.lean, lines 132–162}
\begin{leancode}
\leanline{0}{theorem\ evaluation\_tendsto\_inMeasure\ [NormedSpace\ ℝ\ F]\ [μ.WeaklyRegular]\ \{I\ :\ Type*\}\ \{l\ :\ Filter\ I\}}
\leanline{4}{[l.IsCountablyGenerated]\ (X\ :\ I\ →\ Lp\ E\ 2\ Q)\ (Y\ :\ Lp\ E\ 2\ Q)}
\leanline{4}{(hX\ :\ ∀\ i,\ Q.map\ (X\ i)\ =\ μ)\ (hY\ :\ Q.map\ Y\ =\ μ)}
\leanline{4}{(hlim\ :\ TendstoInMeasure\ Q\ (fun\ i\ ω\ ↦\ X\ i\ ω)\ l\ Y)\ (f\ :\ Lp\ F\ 2\ μ)\ :}
\leanline{4}{Tendsto\ (fun\ i\ ↦\ evaluation\ ({\leanfallback 𝕜}\ :=\ {\leanfallback 𝕜})\ μ\ Q\ (X\ i)\ (hX\ i)\ f)\ l}
\leanline{6}{({\leanfallback 𝓝}\ (evaluation\ ({\leanfallback 𝕜}\ :=\ {\leanfallback 𝕜})\ μ\ Q\ Y\ hY\ f))}
\end{leancode}
\leansource{\texttt{evaluation\_\allowbreak{}bcf\_\allowbreak{}tendsto} \quad WassersteinMixing/\allowbreak{}SemigroupContinuity.lean, lines 166–172}
\begin{leancode}
\leanline{0}{theorem\ evaluation\_bcf\_tendsto\ \{I\ :\ Type*\}\ \{l\ :\ Filter\ I\}\ [l.IsCountablyGenerated]}
\leanline{4}{(X\ :\ I\ →\ Lp\ E\ 2\ Q)\ (Y\ :\ Lp\ E\ 2\ Q)\ (hX\ :\ ∀\ i,\ Q.map\ (X\ i)\ =\ μ)}
\leanline{4}{(hY\ :\ Q.map\ Y\ =\ μ)\ (hlim\ :\ Tendsto\ X\ l\ ({\leanfallback 𝓝}\ Y))\ (φ\ :\ E\ →ᵇ\ F)\ :}
\leanline{4}{Tendsto\ (fun\ i\ ↦\ evaluation\ ({\leanfallback 𝕜}\ :=\ {\leanfallback 𝕜})\ μ\ Q\ (X\ i)\ (hX\ i)}
\leanline{6}{(BoundedContinuousFunction.toLp\ 2\ μ\ {\leanfallback 𝕜}\ φ))\ l}
\leanline{6}{({\leanfallback 𝓝}\ (evaluation\ ({\leanfallback 𝕜}\ :=\ {\leanfallback 𝕜})\ μ\ Q\ Y\ hY\ (BoundedContinuousFunction.toLp\ 2\ μ\ {\leanfallback 𝕜}\ φ)))}
\end{leancode}
\leansource{\texttt{evaluation\_\allowbreak{}tendsto} \quad WassersteinMixing/\allowbreak{}SemigroupContinuity.lean, lines 176–182}
\begin{leancode}
\leanline{0}{theorem\ evaluation\_tendsto\ [NormedSpace\ ℝ\ F]\ [μ.WeaklyRegular]\ \{I\ :\ Type*\}\ \{l\ :\ Filter\ I\}}
\leanline{4}{[l.IsCountablyGenerated]\ (X\ :\ I\ →\ Lp\ E\ 2\ Q)\ (Y\ :\ Lp\ E\ 2\ Q)}
\leanline{4}{(hX\ :\ ∀\ i,\ Q.map\ (X\ i)\ =\ μ)\ (hY\ :\ Q.map\ Y\ =\ μ)}
\leanline{4}{(hlim\ :\ Tendsto\ X\ l\ ({\leanfallback 𝓝}\ Y))\ (f\ :\ Lp\ F\ 2\ μ)\ :}
\leanline{4}{Tendsto\ (fun\ i\ ↦\ evaluation\ ({\leanfallback 𝕜}\ :=\ {\leanfallback 𝕜})\ μ\ Q\ (X\ i)\ (hX\ i)\ f)\ l}
\leanline{6}{({\leanfallback 𝓝}\ (evaluation\ ({\leanfallback 𝕜}\ :=\ {\leanfallback 𝕜})\ μ\ Q\ Y\ hY\ f))}
\end{leancode}
\leanlabel{How the Lean proof follows the argument}
The bounded-continuous theorem uses its exact square-norm integral and bounded convergence in measure. \texttt{evaluation\_\allowbreak{}tendsto\_\allowbreak{}in\allowbreak{}Measure} uses density and the common isometry bound. The last two theorems convert norm convergence of states to convergence in measure before applying these results.
\end{leanbox}
\section{\texorpdfstring{Curvature, step size and the potential contract}{Curvature, step size and the potential contract}}
\label{proof-module-WassersteinMixing-Parameters-lean}


Let \(0<\alpha\le\beta\), let \(d\) be a nonnegative integer and let \(h>0\) satisfy \(h\beta\le1\). Define \(\kappa=\beta/\alpha\), \(q=1-\alpha h\) and \(B=6h\sqrt{\kappa\beta d}\). These are numerical parameters; no invariant law or spectral representation is included in their definition.



Earlier results used here: \cite[Section~3.10]{LangevinLean2026ULASC}.
\subsection{\texorpdfstring{The numerical parameters and their elementary bounds}{The numerical parameters and their elementary bounds}}
\label{report-step-WassersteinMixing-Parameters-lean-1}


\textbf{Claim.} Under the displayed restrictions, \(\beta>0\), \(\kappa\ge1\), \(0\le q<1\), \(h\le1/\beta\), and \(B\ge0\).



\textbf{Proof.} Since \(\beta\ge\alpha>0\), division by \(\alpha\) preserves \(\beta\ge\alpha\) and gives \(\kappa\ge1\). Multiplication by \(h\ge0\) gives \(\alpha h\le\beta h\le1\), so \(q\ge0\). The strict inequality \(\alpha h>0\) gives \(q<1\). Dividing \(h\beta\le1\) by \(\beta>0\) gives the step restriction. Finally \(h>0\) and the nonnegativity of the square root imply \(B\ge0\).

\begin{leanbox}{Lean: \texttt{Parameters}, \texttt{Parameters.\allowbreak{}condition}, \texttt{Parameters.\allowbreak{}contraction} and 7 more}
\leansource{\texttt{Parameters} \quad WassersteinMixing/\allowbreak{}Parameters.lean, lines 19–27}
\begin{leancode}
\leanline{0}{structure\ Parameters\ where}
\leanline{2}{alpha\ :\ ℝ}
\leanline{2}{beta\ :\ ℝ}
\leanline{2}{step\ :\ ℝ}
\leanline{2}{dimension\ :\ ℕ}
\leanline{2}{alpha\_pos\ :\ 0\ <\ alpha}
\leanline{2}{alpha\_le\_beta\ :\ alpha\ ≤\ beta}
\leanline{2}{step\_pos\ :\ 0\ <\ step}
\leanline{2}{step\_beta\_le\ :\ step\ *\ beta\ ≤\ 1}
\end{leancode}
\leansource{\texttt{Parameters.\allowbreak{}condition} \quad WassersteinMixing/\allowbreak{}Parameters.lean, lines 30–30}
\begin{leancode}
\leanline{0}{def\ Parameters.condition\ (P\ :\ Parameters)\ :\ ℝ\ :=\ P.beta\ /\ P.alpha}
\end{leancode}
\leansource{\texttt{Parameters.\allowbreak{}contraction} \quad WassersteinMixing/\allowbreak{}Parameters.lean, lines 33–33}
\begin{leancode}
\leanline{0}{def\ Parameters.contraction\ (P\ :\ Parameters)\ :\ ℝ\ :=\ 1\ -\ P.alpha\ *\ P.step}
\end{leancode}
\leansource{\texttt{Parameters.\allowbreak{}bias} \quad WassersteinMixing/\allowbreak{}Parameters.lean, lines 36–37}
\begin{leancode}
\leanline{0}{def\ Parameters.bias\ (P\ :\ Parameters)\ :\ ℝ\ :=}
\leanline{2}{6\ *\ P.step\ *\ Real.sqrt\ (P.condition\ *\ P.beta\ *\ P.dimension)}
\end{leancode}
\leansource{\texttt{Parameters.\allowbreak{}beta\_\allowbreak{}pos} \quad WassersteinMixing/\allowbreak{}Parameters.lean, lines 39–40}
\begin{leancode}
\leanline{0}{theorem\ Parameters.beta\_pos\ (P\ :\ Parameters)\ :\ 0\ <\ P.beta}
\end{leancode}
\leansource{\texttt{Parameters.\allowbreak{}condition\_\allowbreak{}ge\_\allowbreak{}one} \quad WassersteinMixing/\allowbreak{}Parameters.lean, lines 42–43}
\begin{leancode}
\leanline{0}{theorem\ Parameters.condition\_ge\_one\ (P\ :\ Parameters)\ :\ 1\ ≤\ P.condition}
\end{leancode}
\leansource{\texttt{Parameters.\allowbreak{}contraction\_\allowbreak{}nonneg} \quad WassersteinMixing/\allowbreak{}Parameters.lean, lines 45–48}
\begin{leancode}
\leanline{0}{theorem\ Parameters.contraction\_nonneg\ (P\ :\ Parameters)\ :\ 0\ ≤\ P.contraction}
\end{leancode}
\leansource{\texttt{Parameters.\allowbreak{}contraction\_\allowbreak{}lt\_\allowbreak{}one} \quad WassersteinMixing/\allowbreak{}Parameters.lean, lines 50–52}
\begin{leancode}
\leanline{0}{theorem\ Parameters.contraction\_lt\_one\ (P\ :\ Parameters)\ :\ P.contraction\ <\ 1}
\end{leancode}
\leansource{\texttt{Parameters.\allowbreak{}step\_\allowbreak{}le\_\allowbreak{}inv\_\allowbreak{}beta} \quad WassersteinMixing/\allowbreak{}Parameters.lean, lines 54–55}
\begin{leancode}
\leanline{0}{theorem\ Parameters.step\_le\_inv\_beta\ (P\ :\ Parameters)\ :\ P.step\ ≤\ 1\ /\ P.beta}
\end{leancode}
\leansource{\texttt{Parameters.\allowbreak{}bias\_\allowbreak{}nonneg} \quad WassersteinMixing/\allowbreak{}Parameters.lean, lines 57–59}
\begin{leancode}
\leanline{0}{theorem\ Parameters.bias\_nonneg\ (P\ :\ Parameters)\ :\ 0\ ≤\ P.bias}
\end{leancode}
\leanlabel{How the Lean proof follows the argument}
\texttt{Parameters} records all numerical hypotheses. \texttt{condition}, \texttt{contraction} and \texttt{bias} define their literal expressions. The remaining declarations divide only by a proved positive curvature constant; the contraction inequalities use the supplied step restriction and positive step.
\end{leanbox}
\subsection{\texorpdfstring{The first-order potential contract and the original Hessian bounds}{The first-order potential contract and the original Hessian bounds}}
\label{report-step-WassersteinMixing-Parameters-lean-2}


\textbf{Definition.} On a real inner product space, the potential contract requires \(V\in C^2\), \(\alpha\)-strong convexity and a \(\beta\)-Lipschitz gradient:

\[V(y)\ge V(x)+\langle\nabla V(x),y-x\rangle+\frac\alpha2\|y-x\|^2,\qquad\|\nabla V(y)-\nabla V(x)\|\le\beta\|y-x\|.\]

The original source Hessian bounds imply this contract with exactly the same constants, as proved in the Hessian-assumption argument. The contract contains no invariant law, semigroup, spectral basis or comparison conclusion; those objects and properties are constructed separately.

\begin{leanbox}{Lean: \texttt{PotentialAssumptions}}
\leansource{\texttt{PotentialAssumptions} \quad WassersteinMixing/\allowbreak{}Parameters.lean, lines 63–67}
\begin{leancode}
\leanline{0}{structure\ PotentialAssumptions\ \{H\ :\ Type*\}\ [NormedAddCommGroup\ H]}
\leanline{4}{[InnerProductSpace\ ℝ\ H]\ [CompleteSpace\ H]\ (P\ :\ Parameters)\ (V\ :\ H\ →\ ℝ)\ :\ Prop\ where}
\leanline{2}{regularity\ :\ ContDiff\ ℝ\ 2\ V}
\leanline{2}{strongConvexity\ :\ StrongConvexOn\ Set.univ\ P.alpha\ V}
\leanline{2}{gradientLipschitz\ :\ LipschitzWith\ ⟨P.beta,\ P.beta\_pos.le⟩\ (gradient\ V)}
\end{leancode}
\leanlabel{How the Lean proof follows the argument}
\texttt{Potential\allowbreak{}Assumptions} is a proposition-valued structure containing \texttt{Cont\allowbreak{}Diff \ensuremath{\mathbb R} 2 V}, \texttt{Strong\allowbreak{}Convex\allowbreak{}On} and a \texttt{Lipschitz\allowbreak{}With} gradient condition. None of its fields asserts an invariant measure, a generator realization or a spectral basis.
\end{leanbox}
\section{\texorpdfstring{A genuine integral estimate for the local defect}{A genuine integral estimate for the local defect}}
\label{proof-module-WassersteinMixing-LocalError-lean}


Let \(H\) be a complete real normed space and \(J:[0,h]\to H\) a square-integrable increment. The local defect is the actual Bochner integral \(D=\int_0^hJ(s)\,ds\). Time Cauchy–Schwarz yields the source constant \(3\) from the primitive square-increment constant \(6\).

\subsection{\texorpdfstring{Scalar Cauchy–Schwarz through the square-integrable Hilbert space}{Scalar Cauchy–Schwarz through the square-integrable Hilbert space}}
\label{report-step-WassersteinMixing-LocalError-lean-1}


\textbf{Claim.} If \(f,g\) are square-integrable real functions on a measure space, then

\[\int fg\le\left(\int f^2\right)^{1/2}\left(\int g^2\right)^{1/2}.\]


\textbf{Proof.} The squared norm of the equivalence class of \(f\) in \(L^2\) is \(\int f^2\), and likewise for \(g\). Their real inner product is \(\int fg\). Apply the Hilbert-space bound for the inner product and replace each norm by the square root of its squared integral. The identifications hold almost everywhere and therefore preserve the integrals.

\begin{leanbox}{Lean: \texttt{toLp\_\allowbreak{}norm\_\allowbreak{}sq}, \texttt{real\_\allowbreak{}cauchy\_\allowbreak{}schwarz}}
\leansource{\texttt{toLp\_\allowbreak{}norm\_\allowbreak{}sq} \quad WassersteinMixing/\allowbreak{}LocalError.lean, lines 19–25}
\begin{leancode}
\leanline{0}{private\ theorem\ toLp\_norm\_sq\ \{Ω\ :\ Type*\}\ [MeasurableSpace\ Ω]\ (Q\ :\ Measure\ Ω)}
\leanline{4}{(f\ :\ Ω\ →\ ℝ)\ (hf\ :\ MemLp\ f\ 2\ Q)\ :}
\leanline{4}{‖hf.toLp\ f‖\ \textasciicircum{}\ 2\ =\ ∫\ ω,\ f\ ω\ \textasciicircum{}\ 2\ ∂Q}
\end{leancode}
\leansource{\texttt{real\_\allowbreak{}cauchy\_\allowbreak{}schwarz} \quad WassersteinMixing/\allowbreak{}LocalError.lean, lines 27–40}
\begin{leancode}
\leanline{0}{private\ theorem\ real\_cauchy\_schwarz\ \{Ω\ :\ Type*\}\ [MeasurableSpace\ Ω]\ (Q\ :\ Measure\ Ω)}
\leanline{4}{(f\ g\ :\ Ω\ →\ ℝ)\ (hf\ :\ MemLp\ f\ 2\ Q)\ (hg\ :\ MemLp\ g\ 2\ Q)\ :}
\leanline{4}{(∫\ ω,\ f\ ω\ *\ g\ ω\ ∂Q)\ ≤}
\leanline{6}{Real.sqrt\ (∫\ ω,\ f\ ω\ \textasciicircum{}\ 2\ ∂Q)\ *\ Real.sqrt\ (∫\ ω,\ g\ ω\ \textasciicircum{}\ 2\ ∂Q)}
\end{leancode}
\leanlabel{How the Lean proof follows the argument}
\texttt{to\allowbreak{}Lp\_\allowbreak{}norm\_\allowbreak{}sq} uses the real \texttt{L2.\allowbreak{}inner\_\allowbreak{}def} and the almost-everywhere representative equality. \texttt{real\_\allowbreak{}cauchy\_\allowbreak{}schwarz} identifies the integral product with the genuine inner product and applies \texttt{real\_\allowbreak{}inner\_\allowbreak{}le\_\allowbreak{}norm}.
\end{leanbox}
\subsection{\texorpdfstring{Time Cauchy–Schwarz for the Bochner integral}{Time Cauchy–Schwarz for the Bochner integral}}
\label{report-step-WassersteinMixing-LocalError-lean-2}


\textbf{Claim.} For \(h\ge0\) and square-integrable \(J\) on \((0,h]\),

\[\left\|\int_0^hJ(s)\,ds\right\|^2\le h\int_0^h\|J(s)\|^2\,ds.\]


\textbf{Proof.} The norm of the Bochner integral is at most \(\int_0^h\|J(s)\|\,ds\). Apply scalar Cauchy–Schwarz to \(\|J\|\) and the constant function \(1\). The latter has squared integral \(h\). Both resulting sides are nonnegative, so squaring gives the stated inequality. This also defines the local defect \(D\) as the genuine integral, with all square-integrability hypotheses retained.

\begin{leanbox}{Lean: \texttt{time\_\allowbreak{}integral\_\allowbreak{}norm\_\allowbreak{}sq}, \texttt{stationaryLocalDefect}}
\leansource{\texttt{time\_\allowbreak{}integral\_\allowbreak{}norm\_\allowbreak{}sq} \quad WassersteinMixing/\allowbreak{}LocalError.lean, lines 45–63}
\begin{leancode}
\leanline{0}{theorem\ time\_integral\_norm\_sq\ (increment\ :\ ℝ\ →\ H)\ (h\ :\ ℝ)\ (hh\ :\ 0\ ≤\ h)}
\leanline{4}{(hi\ :\ MemLp\ increment\ 2\ (volume.restrict\ (Ioc\ 0\ h)))\ :}
\leanline{4}{‖∫\ s\ in\ (0\ :\ ℝ)..h,\ increment\ s‖\ \textasciicircum{}\ 2\ ≤}
\leanline{6}{h\ *\ ∫\ s\ in\ (0\ :\ ℝ)..h,\ ‖increment\ s‖\ \textasciicircum{}\ 2}
\end{leancode}
\leansource{\texttt{stationaryLocalDefect} \quad WassersteinMixing/\allowbreak{}LocalError.lean, lines 66–67}
\begin{leancode}
\leanline{0}{def\ stationaryLocalDefect\ (increment\ :\ ℝ\ →\ H)\ (h\ :\ ℝ)\ :\ H\ :=}
\leanline{2}{∫\ s\ in\ (0\ :\ ℝ)..h,\ increment\ s}
\end{leancode}
\leanlabel{How the Lean proof follows the argument}
\texttt{time\_\allowbreak{}integral\_\allowbreak{}norm\_\allowbreak{}sq} combines \texttt{interval\allowbreak{}Integral.\allowbreak{}norm\_\allowbreak{}integral\_\allowbreak{}le\_\allowbreak{}integral\_\allowbreak{}norm} with the preceding scalar bound for the restricted interval measure. It evaluates the measure of the interval and uses the nonnegative square-root identities before squaring. \texttt{stationary\allowbreak{}Local\allowbreak{}Defect} is the Bochner interval integral itself.
\end{leanbox}
\subsection{\texorpdfstring{The exact local defect constant and the displacement polynomial}{The exact local defect constant and the displacement polynomial}}
\label{report-step-WassersteinMixing-LocalError-lean-3}


\textbf{Claim.} If \(\|J(s)\|^2\le6s\beta^2d\) on \((0,h]\), then

\[\|D\|^2\le3h^3\beta^2d.\]


\textbf{Proof.} Square integrability makes \(\|J\|^2\) integrable. Integrate its upper bound and use the preceding time estimate:

\[\|D\|^2\le h\int_0^h6s\beta^2d\,ds
=h\cdot3h^2\beta^2d.\]


For the auxiliary source displacement polynomial, \(0\le s\le h\), \(\beta\ge0\) and \(h\beta\le1\) imply \(s\beta\le1\), hence \(2s^2\beta+4s\le6s\). This last simplification supplies no independent probabilistic displacement estimate.

\begin{leanbox}{Lean: \texttt{stationaryLocalDefect\_\allowbreak{}energy\_\allowbreak{}le}, \texttt{displacement\_\allowbreak{}polynomial\_\allowbreak{}le}}
\leansource{\texttt{stationaryLocalDefect\_\allowbreak{}energy\_\allowbreak{}le} \quad WassersteinMixing/\allowbreak{}LocalError.lean, lines 70–89}
\begin{leancode}
\leanline{0}{theorem\ stationaryLocalDefect\_energy\_le\ (increment\ :\ ℝ\ →\ H)\ (h\ β\ d\ :\ ℝ)}
\leanline{4}{(hh\ :\ 0\ ≤\ h)\ (hi\ :\ MemLp\ increment\ 2\ (volume.restrict\ (Ioc\ 0\ h)))}
\leanline{4}{(hbound\ :\ ∀\ s\ ∈\ Ioc\ 0\ h,\ ‖increment\ s‖\ \textasciicircum{}\ 2\ ≤\ 6\ *\ s\ *\ β\ \textasciicircum{}\ 2\ *\ d)\ :}
\leanline{4}{‖stationaryLocalDefect\ increment\ h‖\ \textasciicircum{}\ 2\ ≤\ 3\ *\ h\ \textasciicircum{}\ 3\ *\ β\ \textasciicircum{}\ 2\ *\ d}
\end{leancode}
\leansource{\texttt{displacement\_\allowbreak{}polynomial\_\allowbreak{}le} \quad WassersteinMixing/\allowbreak{}LocalError.lean, lines 92–96}
\begin{leancode}
\leanline{0}{theorem\ displacement\_polynomial\_le\ (s\ h\ β\ :\ ℝ)\ (hs\ :\ 0\ ≤\ s)\ (hsh\ :\ s\ ≤\ h)}
\leanline{4}{(hβ\ :\ 0\ ≤\ β)\ (hstep\ :\ h\ *\ β\ ≤\ 1)\ :}
\leanline{4}{2\ *\ s\ \textasciicircum{}\ 2\ *\ β\ +\ 4\ *\ s\ ≤\ 6\ *\ s}
\end{leancode}
\leanlabel{How the Lean proof follows the argument}
\texttt{stationary\allowbreak{}Local\allowbreak{}Defect\_\allowbreak{}energy\_\allowbreak{}le} proves integrability of the squared norm from \texttt{Mem\allowbreak{}Lp}, applies monotonicity of the interval integral and evaluates the linear polynomial exactly. \texttt{displacement\_\allowbreak{}polynomial\_\allowbreak{}le} isolates the arithmetic consequence of the stable-step restriction.
\end{leanbox}
\section{\texorpdfstring{Constructing the actual synchronous Gibbs comparison}{Constructing the actual synchronous Gibbs comparison}}
\label{proof-module-WassersteinMixing-ActualComparison-lean}


A diagonal Gibbs initialization and an independent canonical Brownian path provide one probability space for the stationary Langevin diffusion and its Gaussian Euler chain. Both random processes, their common initial state and their actual laws are constructed.



Earlier results used here: \cite[Section~3.45]{LangevinLean2026ULASC}.
\subsection{\texorpdfstring{The concrete probability space and initial moments}{The concrete probability space and initial moments}}
\label{report-step-WassersteinMixing-ActualComparison-lean-1}


\textbf{Construction.} Couple a Gibbs variable with itself on the diagonal, then take its independent product with the coordinate Brownian source. The diagonal coupling has zero transport cost. Strong convexity gives a finite, positive Gibbs partition function and a finite second moment, so both the target and extended law are probability measures and the initialization belongs to \(L^2\).

\begin{leanbox}{Lean: \texttt{ActualSpace}, \texttt{diagonalGibbsCoupling}, \texttt{diagonalGibbsCoupling\_\allowbreak{}cost} and 4 more}
\leansource{\texttt{ActualSpace} \quad WassersteinMixing/\allowbreak{}ActualComparison.lean, lines 19–20}
\begin{leancode}
\leanline{0}{abbrev\ ActualSpace\ (P\ :\ Parameters)\ :=}
\leanline{2}{(E\ P.dimension\ ×\ E\ P.dimension)\ ×\ ULA.CanonicalBrownian.VectorSample\ (Fin\ P.dimension)}
\end{leancode}
\leansource{\texttt{diagonalGibbsCoupling} \quad WassersteinMixing/\allowbreak{}ActualComparison.lean, lines 23–27}
\begin{leancode}
\leanline{0}{def\ diagonalGibbsCoupling\ (P\ :\ Parameters)\ (V\ :\ E\ P.dimension\ →\ ℝ)\ :}
\leanline{4}{ULA.Coupling\ (ULA.gibbsMeasure\ V)\ (ULA.gibbsMeasure\ V)\ where}
\leanline{2}{joint\ :=\ (ULA.gibbsMeasure\ V).map\ (fun\ x\ ↦\ (x,\ x))}
\leanline{2}{fst\ :=\ by\ rw\ [Measure.map\_map\ measurable\_fst\ (by\ fun\_prop)];\ exact\ Measure.map\_id}
\leanline{2}{snd\ :=\ by\ rw\ [Measure.map\_map\ measurable\_snd\ (by\ fun\_prop)];\ exact\ Measure.map\_id}
\end{leancode}
\leansource{\texttt{diagonalGibbsCoupling\_\allowbreak{}cost} \quad WassersteinMixing/\allowbreak{}ActualComparison.lean, lines 29–33}
\begin{leancode}
\leanline{0}{theorem\ diagonalGibbsCoupling\_cost\ (P\ :\ Parameters)\ (V\ :\ E\ P.dimension\ →\ ℝ)\ :}
\leanline{4}{ULA.couplingCost\ (diagonalGibbsCoupling\ P\ V).joint\ =\ 0}
\end{leancode}
\leansource{\texttt{actualProbability} \quad WassersteinMixing/\allowbreak{}ActualComparison.lean, lines 36–38}
\begin{leancode}
\leanline{0}{def\ actualProbability\ (P\ :\ Parameters)\ (V\ :\ E\ P.dimension\ →\ ℝ)\ :\ Measure\ (ActualSpace\ P)\ :=}
\leanline{2}{ULA.initialExtension\ (diagonalGibbsCoupling\ P\ V)}
\leanline{4}{(ULA.CanonicalBrownian.vectorProbability\ (Fin\ P.dimension))}
\end{leancode}
\leansource{\texttt{actualTarget\_\allowbreak{}isProbability} \quad WassersteinMixing/\allowbreak{}ActualComparison.lean, lines 41–44}
\begin{leancode}
\leanline{0}{theorem\ actualTarget\_isProbability\ (P\ :\ Parameters)\ (V\ :\ E\ P.dimension\ →\ ℝ)}
\leanline{4}{(hV\ :\ PotentialAssumptions\ P\ V)\ :\ IsProbabilityMeasure\ (ULA.gibbsMeasure\ V)}
\end{leancode}
\leansource{\texttt{actualProbability\_\allowbreak{}isProbability} \quad WassersteinMixing/\allowbreak{}ActualComparison.lean, lines 47–50}
\begin{leancode}
\leanline{0}{theorem\ actualProbability\_isProbability\ (P\ :\ Parameters)\ (V\ :\ E\ P.dimension\ →\ ℝ)}
\leanline{4}{(hV\ :\ PotentialAssumptions\ P\ V)\ :\ IsProbabilityMeasure\ (actualProbability\ P\ V)}
\end{leancode}
\leansource{\texttt{actualTarget\_\allowbreak{}memLp} \quad WassersteinMixing/\allowbreak{}ActualComparison.lean, lines 53–57}
\begin{leancode}
\leanline{0}{theorem\ actualTarget\_memLp\ (P\ :\ Parameters)\ (V\ :\ E\ P.dimension\ →\ ℝ)}
\leanline{4}{(hV\ :\ PotentialAssumptions\ P\ V)\ :\ MemLp\ (id\ :\ E\ P.dimension\ →\ E\ P.dimension)\ 2}
\leanline{6}{(ULA.gibbsMeasure\ V)}
\end{leancode}
\leanlabel{How the Lean proof follows the argument}
\texttt{diagonal\allowbreak{}Gibbs\allowbreak{}Coupling} uses the exact diagonal pushforward, whose cost is computed as zero. \texttt{actual\allowbreak{}Probability} is its independent Brownian extension. The probability and moment lemmas instantiate the shared Gibbs theorems from the actual potential assumptions.
\end{leanbox}


Earlier results used here: \cite[Section~3.17.1]{LangevinLean2026ULASC}; \cite[Section~3.17.2]{LangevinLean2026ULASC}; \cite[Section~3.22.2]{LangevinLean2026ULASC}; \cite[Section~3.22.3]{LangevinLean2026ULASC}; \cite[Section~3.26.13]{LangevinLean2026ULASC}; \cite[Section~3.26.15]{LangevinLean2026ULASC}; \cite[Section~3.26.3]{LangevinLean2026ULASC}; \cite[Section~3.8.16]{LangevinLean2026ULASC}; \cite[Section~3.8.17]{LangevinLean2026ULASC}.
\subsection{\texorpdfstring{The genuine stationary diffusion}{The genuine stationary diffusion}}
\label{report-step-WassersteinMixing-ActualComparison-lean-2}


\textbf{Construction.} Normalize each Brownian increment to obtain independent standard Gaussian Euler innovations. Solve the Langevin integral equation with the Gibbs initialization and these Brownian paths. The reference segments have the Gibbs law at every interior time, and adjacent segment endpoints agree. In particular the grid reference \(Y_n=Y_{nh}\) is stationary.

\begin{leanbox}{Lean: \texttt{actualNoise}, \texttt{actualReferenceData}, \texttt{actualSegment} and 3 more}
\leansource{\texttt{actualNoise} \quad WassersteinMixing/\allowbreak{}ActualComparison.lean, lines 60–67}
\begin{leancode}
\leanline{0}{def\ actualNoise\ (P\ :\ Parameters)\ (V\ :\ E\ P.dimension\ →\ ℝ)}
\leanline{4}{(hV\ :\ PotentialAssumptions\ P\ V)\ :\ ℕ\ →\ Lp\ (E\ P.dimension)\ 2\ (actualProbability\ P\ V)\ :=\ by}
\leanline{2}{letI\ :=\ actualTarget\_isProbability\ P\ V\ hV}
\leanline{2}{exact\ ULA.Hilbert.productNoise}
\leanline{4}{(ULA.CanonicalBrownian.vectorProbability\ (Fin\ P.dimension))\ (diagonalGibbsCoupling\ P\ V)}
\leanline{4}{(ULA.Hilbert.canonicalEulerNoise\ P.step)}
\leanline{4}{(ULA.Hilbert.sourceEulerNoise\_memLp\ \_\ ULA.CanonicalBrownian.coordinate}
\leanline{6}{ULA.CanonicalBrownian.coordinate\_isPreBrownian\ P.step)}
\end{leancode}
\leansource{\texttt{actualReferenceData} \quad WassersteinMixing/\allowbreak{}ActualComparison.lean, lines 71–84}
\begin{leancode}
\leanline{0}{def\ actualReferenceData\ (P\ :\ Parameters)\ (V\ :\ E\ P.dimension\ →\ ℝ)}
\leanline{4}{(hV\ :\ PotentialAssumptions\ P\ V)\ :}
\leanline{4}{letI\ :=\ actualProbability\_isProbability\ P\ V\ hV}
\leanline{4}{ULA.Hilbert.L2Reference\ (actualProbability\ P\ V)\ (Fintype.card\ (Fin\ P.dimension))\ V\ P.step}
\leanline{6}{⟨P.beta,\ P.beta\_pos.le⟩\ hV.gradientLipschitz\ (actualNoise\ P\ V\ hV)}
\leanline{6}{(ULA.gibbsMeasure\ V)\ :=\ by}
\leanline{2}{letI\ :=\ actualTarget\_isProbability\ P\ V\ hV}
\leanline{2}{letI\ :=\ actualProbability\_isProbability\ P\ V\ hV}
\leanline{2}{have\ hc\ :\ ULA.couplingCost\ (diagonalGibbsCoupling\ P\ V).joint\ ≠\ ⊤\ :=\ by}
\leanline{4}{rw\ [diagonalGibbsCoupling\_cost];\ exact\ ENNReal.zero\_ne\_top}
\leanline{2}{exact\ (ULA.Hilbert.canonicalProductReference\ V\ P.alpha\ P.step}
\leanline{4}{⟨P.beta,\ P.beta\_pos.le⟩\ P.alpha\_pos\ (hV.regularity.differentiable\ (by\ norm\_num))}
\leanline{4}{hV.strongConvexity\ hV.gradientLipschitz\ P.step\_pos}
\leanline{4}{(μ\ :=\ ULA.gibbsMeasure\ V)).reference\ (diagonalGibbsCoupling\ P\ V)\ hc}
\end{leancode}
\leansource{\texttt{actualSegment} \quad WassersteinMixing/\allowbreak{}ActualComparison.lean, lines 87–91}
\begin{leancode}
\leanline{0}{def\ actualSegment\ (P\ :\ Parameters)\ (V\ :\ E\ P.dimension\ →\ ℝ)}
\leanline{4}{(hV\ :\ PotentialAssumptions\ P\ V)\ (n\ :\ ℕ)\ (s\ :\ ℝ)\ :}
\leanline{4}{Lp\ (E\ P.dimension)\ 2\ (actualProbability\ P\ V)\ :=\ by}
\leanline{2}{letI\ :=\ actualProbability\_isProbability\ P\ V\ hV}
\leanline{2}{exact\ (actualReferenceData\ P\ V\ hV).segment\ n\ s}
\end{leancode}
\leansource{\texttt{actualReference} \quad WassersteinMixing/\allowbreak{}ActualComparison.lean, lines 94–96}
\begin{leancode}
\leanline{0}{def\ actualReference\ (P\ :\ Parameters)\ (V\ :\ E\ P.dimension\ →\ ℝ)}
\leanline{4}{(hV\ :\ PotentialAssumptions\ P\ V)\ (n\ :\ ℕ)\ :\ Lp\ (E\ P.dimension)\ 2\ (actualProbability\ P\ V)\ :=}
\leanline{2}{actualSegment\ P\ V\ hV\ n\ 0}
\end{leancode}
\leansource{\texttt{actualReference\_\allowbreak{}stationary} \quad WassersteinMixing/\allowbreak{}ActualComparison.lean, lines 99–103}
\begin{leancode}
\leanline{0}{theorem\ actualReference\_stationary\ (P\ :\ Parameters)\ (V\ :\ E\ P.dimension\ →\ ℝ)}
\leanline{4}{(hV\ :\ PotentialAssumptions\ P\ V)\ (n\ :\ ℕ)\ :}
\leanline{4}{(actualProbability\ P\ V).map\ (actualReference\ P\ V\ hV\ n)\ =\ ULA.gibbsMeasure\ V}
\end{leancode}
\leansource{\texttt{actualSegment\_\allowbreak{}stationary} \quad WassersteinMixing/\allowbreak{}ActualComparison.lean, lines 107–155}
\begin{leancode}
\leanline{0}{theorem\ actualSegment\_stationary\ (P\ :\ Parameters)\ (V\ :\ E\ P.dimension\ →\ ℝ)}
\leanline{4}{(hV\ :\ PotentialAssumptions\ P\ V)\ (n\ :\ ℕ)\ (t\ :\ ℝ)\ (ht\ :\ t\ ∈\ Icc\ 0\ P.step)\ :}
\leanline{4}{(actualProbability\ P\ V).map\ (actualSegment\ P\ V\ hV\ n\ t)\ =\ ULA.gibbsMeasure\ V}
\end{leancode}
\leanlabel{How the Lean proof follows the argument}
\texttt{actual\allowbreak{}Reference\allowbreak{}Data} instantiates the proved Brownian reference construction. \texttt{actual\allowbreak{}Segment} and \texttt{actual\allowbreak{}Reference} select its actual segments and grid knots. Their stationarity follows from the constructed reference, rather than a new stationarity hypothesis.
\end{leanbox}


Earlier results used here: \cite[Section~3.17.2]{LangevinLean2026ULASC}; \cite[Section~3.19.3]{LangevinLean2026ULASC}; \cite[Section~3.22.20]{LangevinLean2026ULASC}; \cite[Section~3.22.22]{LangevinLean2026ULASC}; \cite[Section~3.23.5]{LangevinLean2026ULASC}; \cite[Section~3.25.6]{LangevinLean2026ULASC}; \cite[Section~3.26.3]{LangevinLean2026ULASC}; \cite[Section~3.29.4]{LangevinLean2026ULASC}; \cite[Section~3.29.7]{LangevinLean2026ULASC}; \cite[Section~3.32.12]{LangevinLean2026ULASC}; \cite[Section~3.32.13]{LangevinLean2026ULASC}; \cite[Section~3.32.15]{LangevinLean2026ULASC}; \cite[Section~3.32.17]{LangevinLean2026ULASC}; \cite[Section~3.32.5]{LangevinLean2026ULASC}; \cite[Section~3.42.7]{LangevinLean2026ULASC}; \cite[Section~3.44.5]{LangevinLean2026ULASC}; \cite[Section~3.45.3]{LangevinLean2026ULASC}; \cite[Section~3.45.4]{LangevinLean2026ULASC}; \cite[Section~3.8.17]{LangevinLean2026ULASC}; \cite[Section~3.8.19]{LangevinLean2026ULASC}; \cite[Section~3.8.20]{LangevinLean2026ULASC}; \cite[Section~3.8.22]{LangevinLean2026ULASC}.
\subsection{\texorpdfstring{The actual Gaussian Euler chain}{The actual Gaussian Euler chain}}
\label{report-step-WassersteinMixing-ActualComparison-lean-3}


\textbf{Claim.} On the same probability space, define \(X_{n+1}=X_n-h\nabla V(X_n)+\sqrt{2h}\,\xi_n\), with the same initialization and Brownian innovations. Then \(X_0=Y_0\) and

\[\operatorname{Law}(X_n)=K_h^n\pi,\qquad \pi(dx)\propto e^{-V(x)}\,dx.\]

\textbf{Proof.} The common diagonal initialization has zero \(L^2\) difference. The actual independent-Gaussian recursion has the Markov law obtained by applying the Gaussian ULA law step at each iteration; induction identifies its pushforward with the displayed iterate.

\begin{leanbox}{Lean: \texttt{actualInitial}, \texttt{actualChain}, \texttt{actual\_\allowbreak{}initial\_\allowbreak{}equal} and 1 more}
\leansource{\texttt{actualInitial} \quad WassersteinMixing/\allowbreak{}ActualComparison.lean, lines 158–164}
\begin{leancode}
\leanline{0}{def\ actualInitial\ (P\ :\ Parameters)\ (V\ :\ E\ P.dimension\ →\ ℝ)}
\leanline{4}{(hV\ :\ PotentialAssumptions\ P\ V)\ :\ Lp\ (E\ P.dimension)\ 2\ (actualProbability\ P\ V)\ :=\ by}
\leanline{2}{letI\ :=\ actualTarget\_isProbability\ P\ V\ hV}
\leanline{2}{have\ hc\ :\ ULA.couplingCost\ (diagonalGibbsCoupling\ P\ V).joint\ ≠\ ⊤\ :=\ by}
\leanline{4}{rw\ [diagonalGibbsCoupling\_cost];\ exact\ ENNReal.zero\_ne\_top}
\leanline{2}{exact\ ULA.initialExtension\_lpInitial\ (diagonalGibbsCoupling\ P\ V)}
\leanline{4}{(ULA.CanonicalBrownian.vectorProbability\ (Fin\ P.dimension))\ hc\ (actualTarget\_memLp\ P\ V\ hV)}
\end{leancode}
\leansource{\texttt{actualChain} \quad WassersteinMixing/\allowbreak{}ActualComparison.lean, lines 167–171}
\begin{leancode}
\leanline{0}{def\ actualChain\ (P\ :\ Parameters)\ (V\ :\ E\ P.dimension\ →\ ℝ)}
\leanline{4}{(hV\ :\ PotentialAssumptions\ P\ V)\ (n\ :\ ℕ)\ :\ Lp\ (E\ P.dimension)\ 2\ (actualProbability\ P\ V)\ :=\ by}
\leanline{2}{letI\ :=\ actualProbability\_isProbability\ P\ V\ hV}
\leanline{2}{exact\ ULA.Hilbert.lpChain\ (actualProbability\ P\ V)\ V\ P.step\ hV.gradientLipschitz}
\leanline{4}{(actualInitial\ P\ V\ hV)\ (actualNoise\ P\ V\ hV)\ n}
\end{leancode}
\leansource{\texttt{actual\_\allowbreak{}initial\_\allowbreak{}equal} \quad WassersteinMixing/\allowbreak{}ActualComparison.lean, lines 174–189}
\begin{leancode}
\leanline{0}{theorem\ actual\_initial\_equal\ (P\ :\ Parameters)\ (V\ :\ E\ P.dimension\ →\ ℝ)}
\leanline{4}{(hV\ :\ PotentialAssumptions\ P\ V)\ :\ actualReference\ P\ V\ hV\ 0\ =\ actualChain\ P\ V\ hV\ 0}
\end{leancode}
\leansource{\texttt{actualChain\_\allowbreak{}law} \quad WassersteinMixing/\allowbreak{}ActualComparison.lean, lines 193–212}
\begin{leancode}
\leanline{0}{theorem\ actualChain\_law\ (P\ :\ Parameters)\ (V\ :\ E\ P.dimension\ →\ ℝ)}
\leanline{4}{(hV\ :\ PotentialAssumptions\ P\ V)\ (n\ :\ ℕ)\ :}
\leanline{4}{(actualProbability\ P\ V).map\ (actualChain\ P\ V\ hV\ n)\ =}
\leanline{6}{(ULA.ulaLawStep\ V\ P.step)\textasciicircum{}[n]\ (ULA.gibbsMeasure\ V)}
\end{leancode}
\leanlabel{How the Lean proof follows the argument}
\texttt{actual\allowbreak{}Chain} is the genuine noise-driven Euler recursion. \texttt{actual\_\allowbreak{}initial\_\allowbreak{}equal} uses the diagonal coupling. \texttt{actual\allowbreak{}Chain\_\allowbreak{}law} invokes the exact law theorem of that recursion with the actual Gaussian innovations.
\end{leanbox}


Earlier results used here: \cite[Section~3.17.2]{LangevinLean2026ULASC}; \cite[Section~3.19.1]{LangevinLean2026ULASC}; \cite[Section~3.22.19]{LangevinLean2026ULASC}; \cite[Section~3.22.23]{LangevinLean2026ULASC}; \cite[Section~3.23.9]{LangevinLean2026ULASC}; \cite[Section~3.26.3]{LangevinLean2026ULASC}; \cite[Section~3.32.5]{LangevinLean2026ULASC}; \cite[Section~3.36.7]{LangevinLean2026ULASC}; \cite[Section~3.45.1]{LangevinLean2026ULASC}; \cite[Section~3.45.3]{LangevinLean2026ULASC}; \cite[Section~3.45.4]{LangevinLean2026ULASC}; \cite[Section~3.8.17]{LangevinLean2026ULASC}; \cite[Section~3.8.19]{LangevinLean2026ULASC}; \cite[Section~3.8.20]{LangevinLean2026ULASC}; \cite[Section~3.8.22]{LangevinLean2026ULASC}.
\section{\texorpdfstring{The actual Gibbs diffusion evaluations and contractions}{The actual Gibbs diffusion evaluations and contractions}}
\label{proof-module-WassersteinMixing-SemigroupActual-lean}


The canonical Brownian Langevin solution, constructed from the potential, is stationary for the Gibbs law. Evaluating observables on its grid and interior states therefore gives actual isometries and actual contractive transition compressions, with strong continuity along each segment.

\subsection{\texorpdfstring{Actual operators and their true two-time correlations}{Actual operators and their true two-time correlations}}
\label{report-step-WassersteinMixing-SemigroupActual-lean-1}


\textbf{Construction.} Let \(U_nf=f(Y_{nh})\) and \(U_{n,s}f=f(Y_{nh+s})\). Their Gibbs marginal laws make them isometries. Compress against time-zero evaluation to obtain the actual transition. Its norm is at most one, its value at zero is the identity, and

\[\langle P_nf,g\rangle=\int\langle f(Y_{nh}),g(Y_0)\rangle\,dQ.\]

\textbf{Proof.} Apply the general compression identities to the actual stationary evaluation maps. The representative theorem identifies the abstract pairing with the displayed actual two-time integral.

\begin{leanbox}{Lean: \texttt{actualEvaluation}, \texttt{actualSegmentEvaluation}, \texttt{actualTransition} and 7 more}
\leansource{\texttt{actualEvaluation} \quad WassersteinMixing/\allowbreak{}SemigroupActual.lean, lines 23–27}
\begin{leancode}
\leanline{0}{def\ actualEvaluation\ (P\ :\ Parameters)\ (V\ :\ E\ P.dimension\ →\ ℝ)}
\leanline{4}{(hV\ :\ PotentialAssumptions\ P\ V)\ (n\ :\ ℕ)\ :}
\leanline{4}{Lp\ F\ 2\ (ULA.gibbsMeasure\ V)\ →ₗᵢ[{\leanfallback 𝕜}]\ Lp\ F\ 2\ (actualProbability\ P\ V)\ :=}
\leanline{2}{Lp.compMeasurePreservingₗᵢ\ {\leanfallback 𝕜}\ (actualReference\ P\ V\ hV\ n)}
\leanline{4}{⟨(Lp.stronglyMeasurable\ \_).measurable,\ actualReference\_stationary\ P\ V\ hV\ n⟩}
\end{leancode}
\leansource{\texttt{actualSegmentEvaluation} \quad WassersteinMixing/\allowbreak{}SemigroupActual.lean, lines 30–34}
\begin{leancode}
\leanline{0}{def\ actualSegmentEvaluation\ (P\ :\ Parameters)\ (V\ :\ E\ P.dimension\ →\ ℝ)}
\leanline{4}{(hV\ :\ PotentialAssumptions\ P\ V)\ (n\ :\ ℕ)\ (t\ :\ ℝ)\ (ht\ :\ t\ ∈\ Set.Icc\ 0\ P.step)\ :}
\leanline{4}{Lp\ F\ 2\ (ULA.gibbsMeasure\ V)\ →ₗᵢ[{\leanfallback 𝕜}]\ Lp\ F\ 2\ (actualProbability\ P\ V)\ :=}
\leanline{2}{Lp.compMeasurePreservingₗᵢ\ {\leanfallback 𝕜}\ (actualSegment\ P\ V\ hV\ n\ t)}
\leanline{4}{⟨(Lp.stronglyMeasurable\ \_).measurable,\ actualSegment\_stationary\ P\ V\ hV\ n\ t\ ht⟩}
\end{leancode}
\leansource{\texttt{actualTransition} \quad WassersteinMixing/\allowbreak{}SemigroupActual.lean, lines 37–40}
\begin{leancode}
\leanline{0}{def\ actualTransition\ (P\ :\ Parameters)\ (V\ :\ E\ P.dimension\ →\ ℝ)}
\leanline{4}{(hV\ :\ PotentialAssumptions\ P\ V)\ (n\ :\ ℕ)\ :}
\leanline{4}{Lp\ F\ 2\ (ULA.gibbsMeasure\ V)\ →L[{\leanfallback 𝕜}]\ Lp\ F\ 2\ (ULA.gibbsMeasure\ V)\ :=}
\leanline{2}{ObservableTransition.compression\ (actualEvaluation\ ({\leanfallback 𝕜}\ :=\ {\leanfallback 𝕜})\ P\ V\ hV\ 0)\ (actualEvaluation\ ({\leanfallback 𝕜}\ :=\ {\leanfallback 𝕜})\ P\ V\ hV\ n)}
\end{leancode}
\leansource{\texttt{actualSegmentTransition} \quad WassersteinMixing/\allowbreak{}SemigroupActual.lean, lines 43–47}
\begin{leancode}
\leanline{0}{def\ actualSegmentTransition\ (P\ :\ Parameters)\ (V\ :\ E\ P.dimension\ →\ ℝ)}
\leanline{4}{(hV\ :\ PotentialAssumptions\ P\ V)\ (t\ :\ ℝ)\ (ht\ :\ t\ ∈\ Set.Icc\ 0\ P.step)\ :}
\leanline{4}{Lp\ F\ 2\ (ULA.gibbsMeasure\ V)\ →L[{\leanfallback 𝕜}]\ Lp\ F\ 2\ (ULA.gibbsMeasure\ V)\ :=}
\leanline{2}{ObservableTransition.compression\ (actualEvaluation\ ({\leanfallback 𝕜}\ :=\ {\leanfallback 𝕜})\ P\ V\ hV\ 0)}
\leanline{4}{(actualSegmentEvaluation\ ({\leanfallback 𝕜}\ :=\ {\leanfallback 𝕜})\ P\ V\ hV\ 0\ t\ ht)}
\end{leancode}
\leansource{\texttt{actualTransition\_\allowbreak{}norm\_\allowbreak{}le} \quad WassersteinMixing/\allowbreak{}SemigroupActual.lean, lines 50–53}
\begin{leancode}
\leanline{0}{theorem\ actualTransition\_norm\_le\ (P\ :\ Parameters)\ (V\ :\ E\ P.dimension\ →\ ℝ)}
\leanline{4}{(hV\ :\ PotentialAssumptions\ P\ V)\ (n\ :\ ℕ)\ (f\ :\ Lp\ F\ 2\ (ULA.gibbsMeasure\ V))\ :}
\leanline{4}{‖actualTransition\ ({\leanfallback 𝕜}\ :=\ {\leanfallback 𝕜})\ P\ V\ hV\ n\ f‖\ ≤\ ‖f‖}
\end{leancode}
\leansource{\texttt{actualSegmentTransition\_\allowbreak{}norm\_\allowbreak{}le} \quad WassersteinMixing/\allowbreak{}SemigroupActual.lean, lines 56–60}
\begin{leancode}
\leanline{0}{theorem\ actualSegmentTransition\_norm\_le\ (P\ :\ Parameters)\ (V\ :\ E\ P.dimension\ →\ ℝ)}
\leanline{4}{(hV\ :\ PotentialAssumptions\ P\ V)\ (t\ :\ ℝ)\ (ht\ :\ t\ ∈\ Set.Icc\ 0\ P.step)}
\leanline{4}{(f\ :\ Lp\ F\ 2\ (ULA.gibbsMeasure\ V))\ :}
\leanline{4}{‖actualSegmentTransition\ ({\leanfallback 𝕜}\ :=\ {\leanfallback 𝕜})\ P\ V\ hV\ t\ ht\ f‖\ ≤\ ‖f‖}
\end{leancode}
\leansource{\texttt{actualTransition\_\allowbreak{}zero} \quad WassersteinMixing/\allowbreak{}SemigroupActual.lean, lines 63–66}
\begin{leancode}
\leanline{0}{@[simp]\ theorem\ actualTransition\_zero\ (P\ :\ Parameters)\ (V\ :\ E\ P.dimension\ →\ ℝ)}
\leanline{4}{(hV\ :\ PotentialAssumptions\ P\ V)\ :}
\leanline{4}{actualTransition\ ({\leanfallback 𝕜}\ :=\ {\leanfallback 𝕜})\ (F\ :=\ F)\ P\ V\ hV\ 0\ =\ ContinuousLinearMap.id\ {\leanfallback 𝕜}\ \_}
\end{leancode}
\leansource{\texttt{actualTransition\_\allowbreak{}inner} \quad WassersteinMixing/\allowbreak{}SemigroupActual.lean, lines 69–73}
\begin{leancode}
\leanline{0}{theorem\ actualTransition\_inner\ (P\ :\ Parameters)\ (V\ :\ E\ P.dimension\ →\ ℝ)}
\leanline{4}{(hV\ :\ PotentialAssumptions\ P\ V)\ (n\ :\ ℕ)\ (f\ g\ :\ Lp\ F\ 2\ (ULA.gibbsMeasure\ V))\ :}
\leanline{4}{inner\ {\leanfallback 𝕜}\ (actualTransition\ ({\leanfallback 𝕜}\ :=\ {\leanfallback 𝕜})\ P\ V\ hV\ n\ f)\ g\ =}
\leanline{6}{inner\ {\leanfallback 𝕜}\ (actualEvaluation\ ({\leanfallback 𝕜}\ :=\ {\leanfallback 𝕜})\ P\ V\ hV\ n\ f)\ (actualEvaluation\ ({\leanfallback 𝕜}\ :=\ {\leanfallback 𝕜})\ P\ V\ hV\ 0\ g)}
\end{leancode}
\leansource{\texttt{actualEvaluation\_\allowbreak{}coe} \quad WassersteinMixing/\allowbreak{}SemigroupActual.lean, lines 77–81}
\begin{leancode}
\leanline{0}{theorem\ actualEvaluation\_coe\ (P\ :\ Parameters)\ (V\ :\ E\ P.dimension\ →\ ℝ)}
\leanline{4}{(hV\ :\ PotentialAssumptions\ P\ V)\ (n\ :\ ℕ)\ (f\ :\ Lp\ F\ 2\ (ULA.gibbsMeasure\ V))\ :}
\leanline{4}{actualEvaluation\ ({\leanfallback 𝕜}\ :=\ {\leanfallback 𝕜})\ P\ V\ hV\ n\ f\ =ᵐ[actualProbability\ P\ V]}
\leanline{6}{fun\ ω\ ↦\ f\ (actualReference\ P\ V\ hV\ n\ ω)}
\end{leancode}
\leansource{\texttt{actualTransition\_\allowbreak{}correlation} \quad WassersteinMixing/\allowbreak{}SemigroupActual.lean, lines 84–93}
\begin{leancode}
\leanline{0}{theorem\ actualTransition\_correlation\ (P\ :\ Parameters)\ (V\ :\ E\ P.dimension\ →\ ℝ)}
\leanline{4}{(hV\ :\ PotentialAssumptions\ P\ V)\ (n\ :\ ℕ)\ (f\ g\ :\ Lp\ F\ 2\ (ULA.gibbsMeasure\ V))\ :}
\leanline{4}{inner\ {\leanfallback 𝕜}\ (actualTransition\ ({\leanfallback 𝕜}\ :=\ {\leanfallback 𝕜})\ P\ V\ hV\ n\ f)\ g\ =}
\leanline{6}{∫\ ω,\ inner\ {\leanfallback 𝕜}\ (f\ (actualReference\ P\ V\ hV\ n\ ω))}
\leanline{8}{(g\ (actualReference\ P\ V\ hV\ 0\ ω))\ ∂actualProbability\ P\ V}
\end{leancode}
\leanlabel{How the Lean proof follows the argument}
The definitions use the actual reference state and its proved stationary law. Norm and zero identities specialize the compression results. \texttt{actual\allowbreak{}Evaluation\_\allowbreak{}coe} and \texttt{actual\allowbreak{}Transition\_\allowbreak{}correlation} compare exact pullback representatives.
\end{leanbox}


Earlier results used here: \cite[Section~3.26.3]{LangevinLean2026ULASC}.
\subsection{\texorpdfstring{Strong continuity from the actual continuous reference path}{Strong continuity from the actual continuous reference path}}
\label{report-step-WassersteinMixing-SemigroupActual-lean-2}


\textbf{Claim.} Every square-integrable target observable has norm-continuous evaluation along an actual reference segment; its compressed transition orbit is also continuous.



\textbf{Proof.} The reference segment is continuous in \(L^2\) and has the same Gibbs law throughout the interval. The stationary-observable continuity theorem applies to every observable, not just bounded continuous ones. The fixed initial evaluation adjoint is bounded and hence preserves this convergence after compression.

\begin{leanbox}{Lean: \texttt{continuous\_\allowbreak{}actualSegmentEvaluation}, \texttt{continuous\_\allowbreak{}actualSegmentTransition}}
\leansource{\texttt{continuous\_\allowbreak{}actualSegmentEvaluation} \quad WassersteinMixing/\allowbreak{}SemigroupActual.lean, lines 98–115}
\begin{leancode}
\leanline{0}{theorem\ continuous\_actualSegmentEvaluation\ [NormedSpace\ ℝ\ F]}
\leanline{4}{(P\ :\ Parameters)\ (V\ :\ E\ P.dimension\ →\ ℝ)\ (hV\ :\ PotentialAssumptions\ P\ V)}
\leanline{4}{(n\ :\ ℕ)\ (f\ :\ Lp\ F\ 2\ (ULA.gibbsMeasure\ V))\ :}
\leanline{4}{Continuous\ (fun\ t\ :\ Set.Icc\ (0\ :\ ℝ)\ P.step\ ↦}
\leanline{6}{actualSegmentEvaluation\ ({\leanfallback 𝕜}\ :=\ {\leanfallback 𝕜})\ P\ V\ hV\ n\ t\ t.property\ f)}
\end{leancode}
\leansource{\texttt{continuous\_\allowbreak{}actualSegmentTransition} \quad WassersteinMixing/\allowbreak{}SemigroupActual.lean, lines 118–124}
\begin{leancode}
\leanline{0}{theorem\ continuous\_actualSegmentTransition\ [NormedSpace\ ℝ\ F]}
\leanline{4}{(P\ :\ Parameters)\ (V\ :\ E\ P.dimension\ →\ ℝ)\ (hV\ :\ PotentialAssumptions\ P\ V)}
\leanline{4}{(f\ :\ Lp\ F\ 2\ (ULA.gibbsMeasure\ V))\ :}
\leanline{4}{Continuous\ (fun\ t\ :\ Set.Icc\ (0\ :\ ℝ)\ P.step\ ↦}
\leanline{6}{actualSegmentTransition\ ({\leanfallback 𝕜}\ :=\ {\leanfallback 𝕜})\ P\ V\ hV\ t\ t.property\ f)}
\end{leancode}
\leanlabel{How the Lean proof follows the argument}
\texttt{continuous\_\allowbreak{}actual\allowbreak{}Segment\allowbreak{}Evaluation} applies \texttt{Stationary\allowbreak{}Observables.\allowbreak{}evaluation\_\allowbreak{}tendsto} to the constructed continuous segment. \texttt{continuous\_\allowbreak{}actual\allowbreak{}Segment\allowbreak{}Transition} composes with the continuous adjoint of time-zero evaluation.
\end{leanbox}


Earlier results used here: \cite[Section~3.26.3]{LangevinLean2026ULASC}.
\section{\texorpdfstring{A global actual Gibbs diffusion and its continuous operators}{A global actual Gibbs diffusion and its continuous operators}}
\label{proof-module-WassersteinMixing-SemigroupGlobal-lean}


Finite-horizon Langevin solutions driven by the same actual Brownian source agree on overlaps by uniqueness. Their endpoints define a global canonical path. It has the Gibbs law at every nonnegative time and yields strongly continuous actual observable contractions.

\subsection{\texorpdfstring{Actual Brownian input and finite-horizon solutions}{Actual Brownian input and finite-horizon solutions}}
\label{report-step-WassersteinMixing-SemigroupGlobal-lean-1}


\textbf{Construction.} Use the coordinate Brownian source on the independent Gibbs extension and the actual Gibbs initial variable. For each finite horizon, Picard existence gives a continuous \(L^2\) solution of

\[Y_t-Y_0=-\int_0^t\nabla V(Y_s)\,ds+\sqrt2 B_t.\]

The initial law is the actual normalized Gibbs measure; Brownian coordinates retain their actual pre-Brownian laws on the product space.

\begin{leanbox}{Lean: \texttt{brownian}, \texttt{brownian\_\allowbreak{}preBrownian}, \texttt{initial} and 3 more}
\leansource{\texttt{brownian} \quad WassersteinMixing/\allowbreak{}SemigroupGlobal.lean, lines 20–21}
\begin{leancode}
\leanline{0}{def\ brownian\ (P\ :\ Parameters)\ (i\ :\ Fin\ P.dimension)\ (t\ :\ ℝ≥0)\ (ω\ :\ ActualSpace\ P)\ :\ ℝ\ :=}
\leanline{2}{ULA.CanonicalBrownian.coordinate\ i\ t\ ω.2}
\end{leancode}
\leansource{\texttt{brownian\_\allowbreak{}preBrownian} \quad WassersteinMixing/\allowbreak{}SemigroupGlobal.lean, lines 23–29}
\begin{leancode}
\leanline{0}{theorem\ brownian\_preBrownian\ (P\ :\ Parameters)\ (V\ :\ E\ P.dimension\ →\ ℝ)}
\leanline{4}{(hV\ :\ PotentialAssumptions\ P\ V)\ :}
\leanline{4}{∀\ i,\ IsPreBrownianReal\ (brownian\ P\ i)\ (actualProbability\ P\ V)}
\end{leancode}
\leansource{\texttt{initial} \quad WassersteinMixing/\allowbreak{}SemigroupGlobal.lean, lines 32–36}
\begin{leancode}
\leanline{0}{def\ initial\ (P\ :\ Parameters)\ (V\ :\ E\ P.dimension\ →\ ℝ)\ (hV\ :\ PotentialAssumptions\ P\ V)\ :}
\leanline{4}{Lp\ (E\ P.dimension)\ 2\ (actualProbability\ P\ V)\ :=\ by}
\leanline{2}{let\ :=\ actualTarget\_isProbability\ P\ V\ hV}
\leanline{2}{exact\ ULA.initialExtension\_lpReference\ (diagonalGibbsCoupling\ P\ V)}
\leanline{4}{(ULA.CanonicalBrownian.vectorProbability\ (Fin\ P.dimension))\ (actualTarget\_memLp\ P\ V\ hV)}
\end{leancode}
\leansource{\texttt{initial\_\allowbreak{}law} \quad WassersteinMixing/\allowbreak{}SemigroupGlobal.lean, lines 38–43}
\begin{leancode}
\leanline{0}{theorem\ initial\_law\ (P\ :\ Parameters)\ (V\ :\ E\ P.dimension\ →\ ℝ)}
\leanline{4}{(hV\ :\ PotentialAssumptions\ P\ V)\ :}
\leanline{4}{(actualProbability\ P\ V).map\ (initial\ P\ V\ hV)\ =\ ULA.gibbsMeasure\ V}
\end{leancode}
\leansource{\texttt{solutionOn} \quad WassersteinMixing/\allowbreak{}SemigroupGlobal.lean, lines 46–51}
\begin{leancode}
\leanline{0}{def\ solutionOn\ (P\ :\ Parameters)\ (V\ :\ E\ P.dimension\ →\ ℝ)\ (hV\ :\ PotentialAssumptions\ P\ V)}
\leanline{4}{(T\ :\ ℝ)\ (hT\ :\ 0\ ≤\ T)\ :\ ℝ\ →\ Lp\ (E\ P.dimension)\ 2\ (actualProbability\ P\ V)\ :=\ by}
\leanline{2}{let\ :=\ actualProbability\_isProbability\ P\ V\ hV}
\leanline{2}{exact\ Classical.choose\ (ULA.exists\_lp\_langevin\_integral\_of\_brownian\ (actualProbability\ P\ V)}
\leanline{4}{V\ ⟨P.beta,\ P.beta\_pos.le⟩\ hV.gradientLipschitz\ (initial\ P\ V\ hV)}
\leanline{4}{(brownian\ P)\ (brownian\_preBrownian\ P\ V\ hV)\ 0\ T\ le\_rfl\ hT)}
\end{leancode}
\leansource{\texttt{solutionOn\_\allowbreak{}spec} \quad WassersteinMixing/\allowbreak{}SemigroupGlobal.lean, lines 53–66}
\begin{leancode}
\leanline{0}{theorem\ solutionOn\_spec\ (P\ :\ Parameters)\ (V\ :\ E\ P.dimension\ →\ ℝ)}
\leanline{4}{(hV\ :\ PotentialAssumptions\ P\ V)\ (T\ :\ ℝ)\ (hT\ :\ 0\ ≤\ T)\ :}
\leanline{4}{letI\ :=\ actualProbability\_isProbability\ P\ V\ hV}
\leanline{4}{Continuous\ (solutionOn\ P\ V\ hV\ T\ hT)\ ∧\ solutionOn\ P\ V\ hV\ T\ hT\ 0\ =\ initial\ P\ V\ hV\ ∧}
\leanline{4}{(∀\ t\ ∈\ Icc\ 0\ T,\ solutionOn\ P\ V\ hV\ T\ hT\ t\ -\ initial\ P\ V\ hV\ =}
\leanline{6}{-(∫\ s\ in\ 0..t,\ ULA.lpDrift\ (actualProbability\ P\ V)\ (gradient\ V)}
\leanline{8}{hV.gradientLipschitz\ (solutionOn\ P\ V\ hV\ T\ hT\ s))\ +}
\leanline{6}{ULA.brownianLangevinNoise\ (actualProbability\ P\ V)\ (brownian\ P)}
\leanline{8}{(brownian\_preBrownian\ P\ V\ hV)\ 0\ t)}
\end{leancode}
\leanlabel{How the Lean proof follows the argument}
\texttt{brownian} and \texttt{initial} select the actual source and initialization. \texttt{solution\allowbreak{}On} chooses a solution only from the proved existence theorem. \texttt{solution\allowbreak{}On\_\allowbreak{}spec} retains continuity, the exact initial value and the actual integral equation.
\end{leanbox}


Earlier results used here: \cite[Section~3.13.2]{LangevinLean2026ULASC}; \cite[Section~3.22.20]{LangevinLean2026ULASC}; \cite[Section~3.22.22]{LangevinLean2026ULASC}; \cite[Section~3.25.1]{LangevinLean2026ULASC}; \cite[Section~3.25.6]{LangevinLean2026ULASC}; \cite[Section~3.26.3]{LangevinLean2026ULASC}; \cite[Section~3.32.12]{LangevinLean2026ULASC}; \cite[Section~3.8.17]{LangevinLean2026ULASC}; \cite[Section~3.8.19]{LangevinLean2026ULASC}; \cite[Section~3.8.20]{LangevinLean2026ULASC}.
\subsection{\texorpdfstring{Consistency, continuity and identification with the reference path}{Consistency, continuity and identification with the reference path}}
\label{report-step-WassersteinMixing-SemigroupGlobal-lean-2}


\textbf{Claim.} Define the global endpoint at \(t\ge0\) using horizon \(t\), and hold it constant before zero. It agrees with every longer-horizon solution, is continuous, satisfies the same integral equation and coincides with the actual comparison segment and grid reference.



\textbf{Proof.} Uniqueness of solutions with the same initial state and Brownian input identifies overlapping finite horizons. Near any time choose one fixed larger horizon; the global path equals that continuous solution locally, including the clipped-zero presentation. Substituting consistency into the finite-horizon integral equation gives the global one. The actual segmented reference solves that same equation, so uniqueness gives its segment and grid identification.

\begin{leanbox}{Lean: \texttt{state}, \texttt{state\_\allowbreak{}eq\_\allowbreak{}solutionOn}, \texttt{state\_\allowbreak{}zero} and 6 more}
\leansource{\texttt{state} \quad WassersteinMixing/\allowbreak{}SemigroupGlobal.lean, lines 70–72}
\begin{leancode}
\leanline{0}{def\ state\ (P\ :\ Parameters)\ (V\ :\ E\ P.dimension\ →\ ℝ)\ (hV\ :\ PotentialAssumptions\ P\ V)}
\leanline{4}{(t\ :\ ℝ)\ :\ Lp\ (E\ P.dimension)\ 2\ (actualProbability\ P\ V)\ :=}
\leanline{2}{solutionOn\ P\ V\ hV\ (max\ t\ 0)\ (le\_max\_right\ \_\ \_)\ (max\ t\ 0)}
\end{leancode}
\leansource{\texttt{state\_\allowbreak{}eq\_\allowbreak{}solutionOn} \quad WassersteinMixing/\allowbreak{}SemigroupGlobal.lean, lines 75–88}
\begin{leancode}
\leanline{0}{theorem\ state\_eq\_solutionOn\ (P\ :\ Parameters)\ (V\ :\ E\ P.dimension\ →\ ℝ)}
\leanline{4}{(hV\ :\ PotentialAssumptions\ P\ V)\ (T\ :\ ℝ)\ (hT\ :\ 0\ ≤\ T)}
\leanline{4}{(t\ :\ ℝ)\ (ht\ :\ 0\ ≤\ t)\ (htT\ :\ t\ ≤\ T)\ :}
\leanline{4}{state\ P\ V\ hV\ t\ =\ solutionOn\ P\ V\ hV\ T\ hT\ t}
\end{leancode}
\leansource{\texttt{state\_\allowbreak{}zero} \quad WassersteinMixing/\allowbreak{}SemigroupGlobal.lean, lines 90–92}
\begin{leancode}
\leanline{0}{@[simp]\ theorem\ state\_zero\ (P\ :\ Parameters)\ (V\ :\ E\ P.dimension\ →\ ℝ)}
\leanline{4}{(hV\ :\ PotentialAssumptions\ P\ V)\ :\ state\ P\ V\ hV\ 0\ =\ initial\ P\ V\ hV}
\end{leancode}
\leansource{\texttt{state\_\allowbreak{}eq\_\allowbreak{}solutionOn\_\allowbreak{}max} \quad WassersteinMixing/\allowbreak{}SemigroupGlobal.lean, lines 94–99}
\begin{leancode}
\leanline{0}{theorem\ state\_eq\_solutionOn\_max\ (P\ :\ Parameters)\ (V\ :\ E\ P.dimension\ →\ ℝ)}
\leanline{4}{(hV\ :\ PotentialAssumptions\ P\ V)\ (T\ :\ ℝ)\ (hT\ :\ 0\ ≤\ T)}
\leanline{4}{(t\ :\ ℝ)\ (htT\ :\ max\ t\ 0\ ≤\ T)\ :}
\leanline{4}{state\ P\ V\ hV\ t\ =\ solutionOn\ P\ V\ hV\ T\ hT\ (max\ t\ 0)}
\end{leancode}
\leansource{\texttt{continuous\_\allowbreak{}state} \quad WassersteinMixing/\allowbreak{}SemigroupGlobal.lean, lines 103–115}
\begin{leancode}
\leanline{0}{theorem\ continuous\_state\ (P\ :\ Parameters)\ (V\ :\ E\ P.dimension\ →\ ℝ)}
\leanline{4}{(hV\ :\ PotentialAssumptions\ P\ V)\ :\ Continuous\ (state\ P\ V\ hV)}
\end{leancode}
\leansource{\texttt{continuousOn\_\allowbreak{}state} \quad WassersteinMixing/\allowbreak{}SemigroupGlobal.lean, lines 118–130}
\begin{leancode}
\leanline{0}{theorem\ continuousOn\_state\ (P\ :\ Parameters)\ (V\ :\ E\ P.dimension\ →\ ℝ)}
\leanline{4}{(hV\ :\ PotentialAssumptions\ P\ V)\ :\ ContinuousOn\ (state\ P\ V\ hV)\ (Ici\ 0)}
\end{leancode}
\leansource{\texttt{state\_\allowbreak{}integral} \quad WassersteinMixing/\allowbreak{}SemigroupGlobal.lean, lines 133–150}
\begin{leancode}
\leanline{0}{theorem\ state\_integral\ (P\ :\ Parameters)\ (V\ :\ E\ P.dimension\ →\ ℝ)}
\leanline{4}{(hV\ :\ PotentialAssumptions\ P\ V)\ (t\ :\ ℝ)\ (ht\ :\ 0\ ≤\ t)\ :}
\leanline{4}{letI\ :=\ actualProbability\_isProbability\ P\ V\ hV}
\leanline{4}{state\ P\ V\ hV\ t\ -\ initial\ P\ V\ hV\ =}
\leanline{6}{-(∫\ s\ in\ 0..t,\ ULA.lpDrift\ (actualProbability\ P\ V)\ (gradient\ V)}
\leanline{8}{hV.gradientLipschitz\ (state\ P\ V\ hV\ s))\ +}
\leanline{6}{ULA.brownianLangevinNoise\ (actualProbability\ P\ V)\ (brownian\ P)}
\leanline{8}{(brownian\_preBrownian\ P\ V\ hV)\ 0\ t}
\end{leancode}
\leansource{\texttt{state\_\allowbreak{}eq\_\allowbreak{}actualSegment} \quad WassersteinMixing/\allowbreak{}SemigroupGlobal.lean, lines 153–169}
\begin{leancode}
\leanline{0}{theorem\ state\_eq\_actualSegment\ (P\ :\ Parameters)\ (V\ :\ E\ P.dimension\ →\ ℝ)}
\leanline{4}{(hV\ :\ PotentialAssumptions\ P\ V)\ (n\ :\ ℕ)\ (s\ :\ ℝ)\ (hs\ :\ s\ ∈\ Icc\ 0\ P.step)\ :}
\leanline{4}{state\ P\ V\ hV\ ((n\ :\ ℝ)\ *\ P.step\ +\ s)\ =\ actualSegment\ P\ V\ hV\ n\ s}
\end{leancode}
\leansource{\texttt{state\_\allowbreak{}grid} \quad WassersteinMixing/\allowbreak{}SemigroupGlobal.lean, lines 172–176}
\begin{leancode}
\leanline{0}{theorem\ state\_grid\ (P\ :\ Parameters)\ (V\ :\ E\ P.dimension\ →\ ℝ)}
\leanline{4}{(hV\ :\ PotentialAssumptions\ P\ V)\ (n\ :\ ℕ)\ :}
\leanline{4}{state\ P\ V\ hV\ ((n\ :\ ℝ)\ *\ P.step)\ =\ actualReference\ P\ V\ hV\ n}
\end{leancode}
\leanlabel{How the Lean proof follows the argument}
\texttt{state\_\allowbreak{}eq\_\allowbreak{}solution\allowbreak{}On} proves horizon independence by integral-solution uniqueness. The continuity lemmas use local equality with a longer fixed solution. \texttt{state\_\allowbreak{}integral}, \texttt{state\_\allowbreak{}eq\_\allowbreak{}actual\allowbreak{}Segment} and \texttt{state\_\allowbreak{}grid} carry that exact identification to the global and comparison representations.
\end{leanbox}


Earlier results used here: \cite[Section~3.13.2]{LangevinLean2026ULASC}; \cite[Section~3.25.1]{LangevinLean2026ULASC}; \cite[Section~3.43.4]{LangevinLean2026ULASC}; \cite[Section~3.44.5]{LangevinLean2026ULASC}.
\subsection{\texorpdfstring{Gibbs stationarity and the global observable contraction}{Gibbs stationarity and the global observable contraction}}
\label{report-step-WassersteinMixing-SemigroupGlobal-lean-3}


\textbf{Claim.} Every global state has law \(\pi\). Thus \(U_tf=f(Y_t)\) is an isometry, \(P_t=U_0^*U_t\) is a contraction, \(P_0=I\), and every orbit is strongly continuous for \(t\ge0\).



\textbf{Proof.} The actual Euler innovations have the independent coordinate and initialization laws required by the proved Brownian Gibbs-stationarity theorem. Apply it to the global finite-horizon solution to obtain the marginal law. Actual evaluation and compression then give the operator and norm identities. Continuity of the global state together with stationarity implies strong continuity of every observable evaluation; the fixed adjoint preserves it.

\begin{leanbox}{Lean: \texttt{scalar\_\allowbreak{}independent}, \texttt{initial\_\allowbreak{}independent}, \texttt{state\_\allowbreak{}law} and 5 more}
\leansource{\texttt{scalar\_\allowbreak{}independent} \quad WassersteinMixing/\allowbreak{}SemigroupGlobal.lean, lines 178–189}
\begin{leancode}
\leanline{0}{private\ theorem\ scalar\_independent\ (P\ :\ Parameters)\ (V\ :\ E\ P.dimension\ →\ ℝ)}
\leanline{4}{(hV\ :\ PotentialAssumptions\ P\ V)\ (r\ :\ ℝ)\ (hr\ :\ 0\ <\ r)\ :}
\leanline{4}{iIndepFun\ (fun\ p\ :\ ℕ\ ×\ Fin\ P.dimension\ ↦\ fun\ ω\ ↦}
\leanline{6}{ULA.Hilbert.brownianEulerNoise\ (actualProbability\ P\ V)\ (brownian\ P)}
\leanline{8}{(brownian\_preBrownian\ P\ V\ hV)\ r\ p.1\ ω\ p.2)\ (actualProbability\ P\ V)}
\end{leancode}
\leansource{\texttt{initial\_\allowbreak{}independent} \quad WassersteinMixing/\allowbreak{}SemigroupGlobal.lean, lines 191–202}
\begin{leancode}
\leanline{0}{private\ theorem\ initial\_independent\ (P\ :\ Parameters)\ (V\ :\ E\ P.dimension\ →\ ℝ)}
\leanline{4}{(hV\ :\ PotentialAssumptions\ P\ V)\ (r\ :\ ℝ)\ (\_hr\ :\ 0\ <\ r)\ :}
\leanline{4}{IndepFun\ (initial\ P\ V\ hV)\ (fun\ ω\ k\ ↦}
\leanline{6}{ULA.Hilbert.brownianEulerNoise\ (actualProbability\ P\ V)\ (brownian\ P)}
\leanline{8}{(brownian\_preBrownian\ P\ V\ hV)\ r\ k\ ω)\ (actualProbability\ P\ V)}
\end{leancode}
\leansource{\texttt{state\_\allowbreak{}law} \quad WassersteinMixing/\allowbreak{}SemigroupGlobal.lean, lines 206–217}
\begin{leancode}
\leanline{0}{theorem\ state\_law\ (P\ :\ Parameters)\ (V\ :\ E\ P.dimension\ →\ ℝ)}
\leanline{4}{(hV\ :\ PotentialAssumptions\ P\ V)\ (t\ :\ ℝ)\ :}
\leanline{4}{(actualProbability\ P\ V).map\ (state\ P\ V\ hV\ t)\ =\ ULA.gibbsMeasure\ V}
\end{leancode}
\leansource{\texttt{evaluation} \quad WassersteinMixing/\allowbreak{}SemigroupGlobal.lean, lines 223–226}
\begin{leancode}
\leanline{0}{def\ evaluation\ (P\ :\ Parameters)\ (V\ :\ E\ P.dimension\ →\ ℝ)\ (hV\ :\ PotentialAssumptions\ P\ V)}
\leanline{4}{(t\ :\ ℝ)\ :\ Lp\ F\ 2\ (ULA.gibbsMeasure\ V)\ →ₗᵢ[{\leanfallback 𝕜}]\ Lp\ F\ 2\ (actualProbability\ P\ V)\ :=}
\leanline{2}{StationaryObservables.evaluation\ (ULA.gibbsMeasure\ V)\ (actualProbability\ P\ V)}
\leanline{4}{(state\ P\ V\ hV\ t)\ (state\_law\ P\ V\ hV\ t)}
\end{leancode}
\leansource{\texttt{op} \quad WassersteinMixing/\allowbreak{}SemigroupGlobal.lean, lines 229–231}
\begin{leancode}
\leanline{0}{def\ op\ (P\ :\ Parameters)\ (V\ :\ E\ P.dimension\ →\ ℝ)\ (hV\ :\ PotentialAssumptions\ P\ V)\ (t\ :\ ℝ)\ :}
\leanline{4}{Lp\ F\ 2\ (ULA.gibbsMeasure\ V)\ →L[{\leanfallback 𝕜}]\ Lp\ F\ 2\ (ULA.gibbsMeasure\ V)\ :=}
\leanline{2}{ObservableTransition.compression\ (evaluation\ P\ V\ hV\ 0)\ (evaluation\ P\ V\ hV\ t)}
\end{leancode}
\leansource{\texttt{op\_\allowbreak{}zero} \quad WassersteinMixing/\allowbreak{}SemigroupGlobal.lean, lines 233–235}
\begin{leancode}
\leanline{0}{@[simp]\ theorem\ op\_zero\ (P\ :\ Parameters)\ (V\ :\ E\ P.dimension\ →\ ℝ)}
\leanline{4}{(hV\ :\ PotentialAssumptions\ P\ V)\ :\ op\ ({\leanfallback 𝕜}\ :=\ {\leanfallback 𝕜})\ (F\ :=\ F)\ P\ V\ hV\ 0\ =\ 1}
\end{leancode}
\leansource{\texttt{op\_\allowbreak{}norm\_\allowbreak{}le} \quad WassersteinMixing/\allowbreak{}SemigroupGlobal.lean, lines 237–240}
\begin{leancode}
\leanline{0}{theorem\ op\_norm\_le\ (P\ :\ Parameters)\ (V\ :\ E\ P.dimension\ →\ ℝ)}
\leanline{4}{(hV\ :\ PotentialAssumptions\ P\ V)\ (t\ :\ ℝ)\ (f\ :\ Lp\ F\ 2\ (ULA.gibbsMeasure\ V))\ :}
\leanline{4}{‖op\ ({\leanfallback 𝕜}\ :=\ {\leanfallback 𝕜})\ P\ V\ hV\ t\ f‖\ ≤\ ‖f‖}
\end{leancode}
\leansource{\texttt{continuousOn\_\allowbreak{}op} \quad WassersteinMixing/\allowbreak{}SemigroupGlobal.lean, lines 244–254}
\begin{leancode}
\leanline{0}{theorem\ continuousOn\_op\ [NormedSpace\ ℝ\ F]\ (P\ :\ Parameters)\ (V\ :\ E\ P.dimension\ →\ ℝ)}
\leanline{4}{(hV\ :\ PotentialAssumptions\ P\ V)\ (f\ :\ Lp\ F\ 2\ (ULA.gibbsMeasure\ V))\ :}
\leanline{4}{ContinuousOn\ (fun\ t\ ↦\ op\ ({\leanfallback 𝕜}\ :=\ {\leanfallback 𝕜})\ P\ V\ hV\ t\ f)\ (Ici\ 0)}
\end{leancode}
\leanlabel{How the Lean proof follows the argument}
The two independence helpers specialize the actual Brownian product source. \texttt{state\_\allowbreak{}law} invokes the proved Langevin Gibbs-stationarity result. \texttt{evaluation} and \texttt{op} are actual pullback and compression constructions; \texttt{continuous\allowbreak{}On\_\allowbreak{}op} uses the stationary-observable continuity theorem.
\end{leanbox}


Earlier results used here: \cite[Section~3.26.3]{LangevinLean2026ULASC}; \cite[Section~3.29.7]{LangevinLean2026ULASC}; \cite[Section~3.32.13]{LangevinLean2026ULASC}; \cite[Section~3.32.15]{LangevinLean2026ULASC}; \cite[Section~3.32.17]{LangevinLean2026ULASC}; \cite[Section~3.42.7]{LangevinLean2026ULASC}; \cite[Section~3.8.17]{LangevinLean2026ULASC}; \cite[Section~3.8.19]{LangevinLean2026ULASC}; \cite[Section~3.8.20]{LangevinLean2026ULASC}; \cite[Section~3.8.22]{LangevinLean2026ULASC}.
\section{\texorpdfstring{The actual score increments and local defect constants}{The actual score increments and local defect constants}}
\label{proof-module-WassersteinMixing-ActualLocalError-lean}


All quantities in this section come from the constructed stationary Langevin path. The estimates \(6s\beta^2d\) for score increments and \(3h^3\beta^2d\) for the integrated defect are derived, and subtracting the actual diffusion and Euler updates gives the exact synchronous error equation.

\subsection{\texorpdfstring{Deriving the primitive score-increment estimate}{Deriving the primitive score-increment estimate}}
\label{report-step-WassersteinMixing-ActualLocalError-lean-1}


\textbf{Claim.} On a stationary reference segment,

\[\mathbb E\|Y_{nh+s}-Y_{nh}\|^2\le d(2\beta s^2+4s),\qquad \mathbb E\|\nabla V(Y_{nh+s})-\nabla V(Y_{nh})\|^2\le6s\beta^2d\]

for \(0\le s\le h\), \(h\beta\le1\).



\textbf{Proof.} The integral equation splits the displacement into integrated drift and Brownian noise. Squared triangle and time Cauchy–Schwarz, the stationary score moment and exact Brownian variance give the displacement polynomial. Gradient Lipschitzness multiplies it by \(\beta^2\). Since \(\beta s\le1\), the polynomial is at most \(6s\). Continuity follows from the continuous reference segment and Lipschitz drift map.

\begin{leanbox}{Lean: \texttt{referenceScoreIncrement}, \texttt{reference\_\allowbreak{}displacement\_\allowbreak{}energy}, \texttt{reference\_\allowbreak{}score\_\allowbreak{}increment\_\allowbreak{}polynomial} and 2 more}
\leansource{\texttt{referenceScoreIncrement} \quad WassersteinMixing/\allowbreak{}ActualLocalError.lean, lines 26–28}
\begin{leancode}
\leanline{0}{def\ referenceScoreIncrement\ (n\ :\ ℕ)\ (s\ :\ ℝ)\ :\ Lp\ H\ 2\ Q\ :=}
\leanline{2}{ULA.lpDrift\ Q\ (gradient\ V)\ hLip\ (R.segment\ n\ 0)\ -}
\leanline{4}{ULA.lpDrift\ Q\ (gradient\ V)\ hLip\ (R.segment\ n\ s)}
\end{leancode}
\leansource{\texttt{reference\_\allowbreak{}displacement\_\allowbreak{}energy} \quad WassersteinMixing/\allowbreak{}ActualLocalError.lean, lines 32–45}
\begin{leancode}
\leanline{0}{theorem\ reference\_displacement\_energy\ (n\ :\ ℕ)\ (s\ :\ ℝ)\ (hs\ :\ s\ ∈\ Icc\ 0\ h)\ :}
\leanline{4}{‖R.segment\ n\ s\ -\ R.segment\ n\ 0‖\ \textasciicircum{}\ 2\ ≤}
\leanline{6}{(d\ :\ ℝ)\ *\ (2\ *\ s\ \textasciicircum{}\ 2\ *\ (β\ :\ ℝ)\ +\ 4\ *\ s)}
\end{leancode}
\leansource{\texttt{reference\_\allowbreak{}score\_\allowbreak{}increment\_\allowbreak{}polynomial} \quad WassersteinMixing/\allowbreak{}ActualLocalError.lean, lines 49–59}
\begin{leancode}
\leanline{0}{theorem\ reference\_score\_increment\_polynomial\ (n\ :\ ℕ)\ (s\ :\ ℝ)\ (hs\ :\ s\ ∈\ Icc\ 0\ h)\ :}
\leanline{4}{‖referenceScoreIncrement\ Q\ d\ V\ h\ β\ hLip\ noise\ π\ R\ n\ s‖\ \textasciicircum{}\ 2\ ≤}
\leanline{6}{(β\ :\ ℝ)\ \textasciicircum{}\ 2\ *\ d\ *\ (2\ *\ s\ \textasciicircum{}\ 2\ *\ (β\ :\ ℝ)\ +\ 4\ *\ s)}
\end{leancode}
\leansource{\texttt{reference\_\allowbreak{}score\_\allowbreak{}increment\_\allowbreak{}energy} \quad WassersteinMixing/\allowbreak{}ActualLocalError.lean, lines 62–70}
\begin{leancode}
\leanline{0}{theorem\ reference\_score\_increment\_energy\ (hstep\ :\ h\ *\ (β\ :\ ℝ)\ ≤\ 1)}
\leanline{4}{(n\ :\ ℕ)\ (s\ :\ ℝ)\ (hs\ :\ s\ ∈\ Icc\ 0\ h)\ :}
\leanline{4}{‖referenceScoreIncrement\ Q\ d\ V\ h\ β\ hLip\ noise\ π\ R\ n\ s‖\ \textasciicircum{}\ 2\ ≤}
\leanline{6}{6\ *\ s\ *\ (β\ :\ ℝ)\ \textasciicircum{}\ 2\ *\ d}
\end{leancode}
\leansource{\texttt{referenceScoreIncrement\_\allowbreak{}continuous} \quad WassersteinMixing/\allowbreak{}ActualLocalError.lean, lines 74–77}
\begin{leancode}
\leanline{0}{theorem\ referenceScoreIncrement\_continuous\ (n\ :\ ℕ)\ :}
\leanline{4}{Continuous\ (referenceScoreIncrement\ Q\ d\ V\ h\ β\ hLip\ noise\ π\ R\ n)}
\end{leancode}
\leanlabel{How the Lean proof follows the argument}
The generic reference lemmas use the actual reference integral equation and noise moment. \texttt{reference\_\allowbreak{}score\_\allowbreak{}increment\_\allowbreak{}energy} applies \texttt{displacement\_\allowbreak{}polynomial\_\allowbreak{}le} to the proved polynomial estimate. The increment is defined by the true Lipschitz gradient applied to the reference states.
\end{leanbox}


Earlier results used here: \cite[Section~3.11.6]{LangevinLean2026ULASC}; \cite[Section~3.13.2]{LangevinLean2026ULASC}; \cite[Section~3.13.5]{LangevinLean2026ULASC}.
\subsection{\texorpdfstring{Actual scores and the constant-three defect}{Actual scores and the constant-three defect}}
\label{report-step-WassersteinMixing-ActualLocalError-lean-2}


\textbf{Construction and bounds.} Set \(F_n=\nabla V(Y_{nh})\), \(I_n(s)=F_n-\nabla V(Y_{nh+s})\) and \(D_n=\int_0^h I_n(s)\,ds\). Then

\[\|F_n\|_2^2\le\beta d,\quad\|I_n(s)\|_2^2\le6s\beta^2d,\quad\|D_n\|_2^2\le3h^3\beta^2d.\]

\textbf{Proof.} Gibbs stationarity supplies the score moment. Specializing the preceding increment estimate gives the second bound. The time curve is continuous and hence square integrable on the compact interval. Applying time Cauchy–Schwarz and integrating \(6s\beta^2d\) gives the last constant. The endpoint increment is exactly \(F_n-F_{n+1}\).

\begin{leanbox}{Lean: \texttt{actualScore}, \texttt{actualScoreIncrement}, \texttt{actualDefect} and 6 more}
\leansource{\texttt{actualScore} \quad WassersteinMixing/\allowbreak{}ActualLocalError.lean, lines 85–88}
\begin{leancode}
\leanline{0}{def\ actualScore\ (n\ :\ ℕ)\ :\ Lp\ (E\ P.dimension)\ 2\ (actualProbability\ P\ V)\ :=\ by}
\leanline{2}{letI\ :=\ actualProbability\_isProbability\ P\ V\ hV}
\leanline{2}{exact\ ULA.lpDrift\ (actualProbability\ P\ V)\ (gradient\ V)\ hV.gradientLipschitz}
\leanline{4}{(actualReference\ P\ V\ hV\ n)}
\end{leancode}
\leansource{\texttt{actualScoreIncrement} \quad WassersteinMixing/\allowbreak{}ActualLocalError.lean, lines 91–95}
\begin{leancode}
\leanline{0}{def\ actualScoreIncrement\ (n\ :\ ℕ)\ (s\ :\ ℝ)\ :\ Lp\ (E\ P.dimension)\ 2\ (actualProbability\ P\ V)\ :=\ by}
\leanline{2}{letI\ :=\ actualProbability\_isProbability\ P\ V\ hV}
\leanline{2}{exact\ actualScore\ P\ V\ hV\ n\ -}
\leanline{4}{ULA.lpDrift\ (actualProbability\ P\ V)\ (gradient\ V)\ hV.gradientLipschitz}
\leanline{6}{(actualSegment\ P\ V\ hV\ n\ s)}
\end{leancode}
\leansource{\texttt{actualDefect} \quad WassersteinMixing/\allowbreak{}ActualLocalError.lean, lines 98–99}
\begin{leancode}
\leanline{0}{def\ actualDefect\ (n\ :\ ℕ)\ :\ Lp\ (E\ P.dimension)\ 2\ (actualProbability\ P\ V)\ :=}
\leanline{2}{stationaryLocalDefect\ (actualScoreIncrement\ P\ V\ hV\ n)\ P.step}
\end{leancode}
\leansource{\texttt{actualScore\_\allowbreak{}energy} \quad WassersteinMixing/\allowbreak{}ActualLocalError.lean, lines 102–114}
\begin{leancode}
\leanline{0}{theorem\ actualScore\_energy\ (n\ :\ ℕ)\ :}
\leanline{4}{‖actualScore\ P\ V\ hV\ n‖\ \textasciicircum{}\ 2\ ≤\ P.beta\ *\ P.dimension}
\end{leancode}
\leansource{\texttt{actualScoreIncrement\_\allowbreak{}energy} \quad WassersteinMixing/\allowbreak{}ActualLocalError.lean, lines 117–126}
\begin{leancode}
\leanline{0}{theorem\ actualScoreIncrement\_energy\ (n\ :\ ℕ)\ (s\ :\ ℝ)\ (hs\ :\ s\ ∈\ Icc\ 0\ P.step)\ :}
\leanline{4}{‖actualScoreIncrement\ P\ V\ hV\ n\ s‖\ \textasciicircum{}\ 2\ ≤\ 6\ *\ s\ *\ P.beta\ \textasciicircum{}\ 2\ *\ P.dimension}
\end{leancode}
\leansource{\texttt{actualScoreIncrement\_\allowbreak{}continuous} \quad WassersteinMixing/\allowbreak{}ActualLocalError.lean, lines 129–134}
\begin{leancode}
\leanline{0}{theorem\ actualScoreIncrement\_continuous\ (n\ :\ ℕ)\ :\ Continuous\ (actualScoreIncrement\ P\ V\ hV\ n)}
\end{leancode}
\leansource{\texttt{actualScoreIncrement\_\allowbreak{}memLp} \quad WassersteinMixing/\allowbreak{}ActualLocalError.lean, lines 137–143}
\begin{leancode}
\leanline{0}{theorem\ actualScoreIncrement\_memLp\ (n\ :\ ℕ)\ :}
\leanline{4}{MemLp\ (actualScoreIncrement\ P\ V\ hV\ n)\ 2\ (volume.restrict\ (Ioc\ 0\ P.step))}
\end{leancode}
\leansource{\texttt{actualScoreIncrement\_\allowbreak{}endpoint} \quad WassersteinMixing/\allowbreak{}ActualLocalError.lean, lines 146–151}
\begin{leancode}
\leanline{0}{theorem\ actualScoreIncrement\_endpoint\ (n\ :\ ℕ)\ :}
\leanline{4}{actualScoreIncrement\ P\ V\ hV\ n\ P.step\ =\ actualScore\ P\ V\ hV\ n\ -\ actualScore\ P\ V\ hV\ (n\ +\ 1)}
\end{leancode}
\leansource{\texttt{actualDefect\_\allowbreak{}energy} \quad WassersteinMixing/\allowbreak{}ActualLocalError.lean, lines 154–158}
\begin{leancode}
\leanline{0}{theorem\ actualDefect\_energy\ (n\ :\ ℕ)\ :}
\leanline{4}{‖actualDefect\ P\ V\ hV\ n‖\ \textasciicircum{}\ 2\ ≤\ 3\ *\ P.step\ \textasciicircum{}\ 3\ *\ P.beta\ \textasciicircum{}\ 2\ *\ P.dimension}
\end{leancode}
\leanlabel{How the Lean proof follows the argument}
\texttt{actual\allowbreak{}Score}, \texttt{actual\allowbreak{}Score\allowbreak{}Increment} and \texttt{actual\allowbreak{}Defect} use the constructed diffusion. Their moment, continuity and endpoint theorems discharge the generic local-defect estimate directly; none of these bounds is a primitive comparison field.
\end{leanbox}


Earlier results used here: \cite[Section~3.13.2]{LangevinLean2026ULASC}; \cite[Section~3.26.3]{LangevinLean2026ULASC}; \cite[Section~3.28.2]{LangevinLean2026ULASC}.
\subsection{\texorpdfstring{Subtracting the true updates}{Subtracting the true updates}}
\label{report-step-WassersteinMixing-ActualLocalError-lean-3}


\textbf{Claim.} The reference endpoint satisfies

\[Y_{(n+1)h}=Y_{nh}-hF_n+\sqrt{2h}\,\xi_n+D_n.\]

\textbf{Proof.} Integrating the definition of \(I_n\) gives \(D_n=hF_n-\int_0^h\nabla V(Y_{nh+s})\,ds\). Substitute this into the actual Langevin integral equation and use its exact Brownian endpoint. Subtracting the Euler update cancels the shared innovation and gives

\[e_{n+1}=e_n-h\bigl(\nabla V(Y_{nh})-\nabla V(X_n)\bigr)+D_n,\qquad e_n=Y_{nh}-X_n.\]\begin{leanbox}{Lean: \texttt{actualDefect\_\allowbreak{}eq}, \texttt{actualReference\_\allowbreak{}step}, \texttt{actual\_\allowbreak{}error\_\allowbreak{}step}}
\leansource{\texttt{actualDefect\_\allowbreak{}eq} \quad WassersteinMixing/\allowbreak{}ActualLocalError.lean, lines 161–175}
\begin{leancode}
\leanline{0}{theorem\ actualDefect\_eq\ (n\ :\ ℕ)\ :}
\leanline{4}{letI\ :=\ actualProbability\_isProbability\ P\ V\ hV}
\leanline{4}{actualDefect\ P\ V\ hV\ n\ =\ P.step\ •\ actualScore\ P\ V\ hV\ n\ -}
\leanline{6}{∫\ s\ in\ (0\ :\ ℝ)..P.step,}
\leanline{8}{ULA.lpDrift\ (actualProbability\ P\ V)\ (gradient\ V)\ hV.gradientLipschitz}
\leanline{10}{(actualSegment\ P\ V\ hV\ n\ s)}
\end{leancode}
\leansource{\texttt{actualReference\_\allowbreak{}step} \quad WassersteinMixing/\allowbreak{}ActualLocalError.lean, lines 179–193}
\begin{leancode}
\leanline{0}{theorem\ actualReference\_step\ (n\ :\ ℕ)\ :}
\leanline{4}{actualReference\ P\ V\ hV\ (n\ +\ 1)\ =\ actualReference\ P\ V\ hV\ n\ -}
\leanline{6}{P.step\ •\ actualScore\ P\ V\ hV\ n\ +\ Real.sqrt\ (2\ *\ P.step)\ •\ actualNoise\ P\ V\ hV\ n\ +}
\leanline{8}{actualDefect\ P\ V\ hV\ n}
\end{leancode}
\leansource{\texttt{actual\_\allowbreak{}error\_\allowbreak{}step} \quad WassersteinMixing/\allowbreak{}ActualLocalError.lean, lines 197–209}
\begin{leancode}
\leanline{0}{theorem\ actual\_error\_step\ (n\ :\ ℕ)\ :}
\leanline{4}{letI\ :=\ actualProbability\_isProbability\ P\ V\ hV}
\leanline{4}{actualReference\ P\ V\ hV\ (n\ +\ 1)\ -\ actualChain\ P\ V\ hV\ (n\ +\ 1)\ =}
\leanline{6}{(actualReference\ P\ V\ hV\ n\ -\ actualChain\ P\ V\ hV\ n)\ -}
\leanline{8}{P.step\ •\ (actualScore\ P\ V\ hV\ n\ -}
\leanline{10}{ULA.lpDrift\ (actualProbability\ P\ V)\ (gradient\ V)\ hV.gradientLipschitz}
\leanline{12}{(actualChain\ P\ V\ hV\ n))\ +\ actualDefect\ P\ V\ hV\ n}
\end{leancode}
\leanlabel{How the Lean proof follows the argument}
\texttt{actual\allowbreak{}Defect\_\allowbreak{}eq} uses integrability before subtracting interval integrals. \texttt{actual\allowbreak{}Reference\_\allowbreak{}step} uses the proved reference equation and endpoint noise identity. \texttt{actual\_\allowbreak{}error\_\allowbreak{}step} is then algebraic subtraction of the two actual updates.
\end{leanbox}


Earlier results used here: \cite[Section~3.13.2]{LangevinLean2026ULASC}; \cite[Section~3.13.5]{LangevinLean2026ULASC}.
\section{\texorpdfstring{Curvature, shifted errors and genuine conditional centering}{Curvature, shifted errors and genuine conditional centering}}
\label{proof-module-WassersteinMixing-ShiftedEnergy-lean}


This section derives the deterministic weighted interpolation inequality and uses the actual conditional-expectation projection in \(L^2\) to remove the martingale cross term. These general inequalities take an exact error recursion, curvature bounds and conditional centering as hypotheses. The actual averaged-Hessian construction, conditional defect mean and realized stationary comparison prove all those hypotheses for the genuine Langevin and ULA processes.

\subsection{\texorpdfstring{Positive contractions and weighted interpolation}{Positive contractions and weighted interpolation}}
\label{report-step-WassersteinMixing-ShiftedEnergy-lean-1}


\textbf{Claim.} Let \(T\) be a self-adjoint linear operator on a real inner product space with \(0\le\langle x,Tx\rangle\le\|x\|^2\). Then \(\|Tx\|^2\le\langle x,Tx\rangle\). If additionally \(a\|x\|^2\le\langle x,Tx\rangle\le b\|x\|^2\), then

\[\|u-Tu+Tv\|^2\le(1-a)\|u\|^2+b\|v\|^2.\]


\textbf{Proof.} The quadratic function \(F(x)=\langle x,Tx\rangle/2\) has gradient \(Tx\). Positivity and the upper quadratic bound give its convex lower model and unit-smooth upper model. The resulting gradient inequality yields \(\|Tx\|^2\le\langle x,Tx\rangle\). Expanding \(\|u+T(v-u)\|^2\), applying this inequality to \(v-u\), and using symmetry gives the sharper upper bound \(\|u\|^2-\langle u,Tu\rangle+\langle v,Tv\rangle\). The two quadratic-form bounds finish the proof.




For \(T=hA\), where \(A\) is self-adjoint with \(\alpha I\preceq A\preceq\beta I\), the restriction \(h\beta\le1\) supplies \(0\preceq T\preceq I\). Taking \(a=\alpha h\), \(b=\beta h\) gives

\[\|u-hAu+hAv\|^2\le(1-\alpha h)\|u\|^2+h\beta\|v\|^2.\]\begin{leanbox}{Lean: \texttt{positive\_\allowbreak{}contraction\_\allowbreak{}cocoercive}, \texttt{weighted\_\allowbreak{}interpolation}, \texttt{curvature\_\allowbreak{}weighted\_\allowbreak{}interpolation}}
\leansource{\texttt{positive\_\allowbreak{}contraction\_\allowbreak{}cocoercive} \quad WassersteinMixing/\allowbreak{}ShiftedEnergy.lean, lines 23–50}
\begin{leancode}
\leanline{0}{theorem\ positive\_contraction\_cocoercive\ (T\ :\ H\ →L[ℝ]\ H)}
\leanline{4}{(hsym\ :\ ∀\ x\ y,\ inner\ ℝ\ (T\ x)\ y\ =\ inner\ ℝ\ x\ (T\ y))}
\leanline{4}{(hpos\ :\ ∀\ x,\ 0\ ≤\ inner\ ℝ\ x\ (T\ x))}
\leanline{4}{(hupper\ :\ ∀\ x,\ inner\ ℝ\ x\ (T\ x)\ ≤\ ‖x‖\ \textasciicircum{}\ 2)\ (x\ :\ H)\ :}
\leanline{4}{‖T\ x‖\ \textasciicircum{}\ 2\ ≤\ inner\ ℝ\ x\ (T\ x)}
\end{leancode}
\leansource{\texttt{weighted\_\allowbreak{}interpolation} \quad WassersteinMixing/\allowbreak{}ShiftedEnergy.lean, lines 53–69}
\begin{leancode}
\leanline{0}{theorem\ weighted\_interpolation\ (T\ :\ H\ →L[ℝ]\ H)\ (a\ b\ :\ ℝ)}
\leanline{4}{(hsym\ :\ ∀\ x\ y,\ inner\ ℝ\ (T\ x)\ y\ =\ inner\ ℝ\ x\ (T\ y))}
\leanline{4}{(hpos\ :\ ∀\ x,\ 0\ ≤\ inner\ ℝ\ x\ (T\ x))}
\leanline{4}{(hupper\ :\ ∀\ x,\ inner\ ℝ\ x\ (T\ x)\ ≤\ ‖x‖\ \textasciicircum{}\ 2)}
\leanline{4}{(hlower\ :\ ∀\ x,\ a\ *\ ‖x‖\ \textasciicircum{}\ 2\ ≤\ inner\ ℝ\ x\ (T\ x))}
\leanline{4}{(hbound\ :\ ∀\ x,\ inner\ ℝ\ x\ (T\ x)\ ≤\ b\ *\ ‖x‖\ \textasciicircum{}\ 2)\ (u\ v\ :\ H)\ :}
\leanline{4}{‖u\ -\ T\ u\ +\ T\ v‖\ \textasciicircum{}\ 2\ ≤\ (1\ -\ a)\ *\ ‖u‖\ \textasciicircum{}\ 2\ +\ b\ *\ ‖v‖\ \textasciicircum{}\ 2}
\end{leancode}
\leansource{\texttt{curvature\_\allowbreak{}weighted\_\allowbreak{}interpolation} \quad WassersteinMixing/\allowbreak{}ShiftedEnergy.lean, lines 72–93}
\begin{leancode}
\leanline{0}{theorem\ curvature\_weighted\_interpolation\ (P\ :\ Parameters)\ (A\ :\ H\ →L[ℝ]\ H)}
\leanline{4}{(hsym\ :\ ∀\ x\ y,\ inner\ ℝ\ (A\ x)\ y\ =\ inner\ ℝ\ x\ (A\ y))}
\leanline{4}{(hlower\ :\ ∀\ x,\ P.alpha\ *\ ‖x‖\ \textasciicircum{}\ 2\ ≤\ inner\ ℝ\ x\ (A\ x))}
\leanline{4}{(hupper\ :\ ∀\ x,\ inner\ ℝ\ x\ (A\ x)\ ≤\ P.beta\ *\ ‖x‖\ \textasciicircum{}\ 2)\ (u\ v\ :\ H)\ :}
\leanline{4}{‖u\ -\ P.step\ •\ A\ u\ +\ P.step\ •\ A\ v‖\ \textasciicircum{}\ 2\ ≤}
\leanline{6}{P.contraction\ *\ ‖u‖\ \textasciicircum{}\ 2\ +\ P.step\ *\ P.beta\ *\ ‖v‖\ \textasciicircum{}\ 2}
\end{leancode}
\leanlabel{How the Lean proof follows the argument}
\texttt{positive\_\allowbreak{}contraction\_\allowbreak{}cocoercive} derives the gradient inequality from explicit quadratic models using the shared ULA result. \texttt{weighted\_\allowbreak{}interpolation} expands the squared norm and applies the two quadratic-form estimates. \texttt{curvature\_\allowbreak{}weighted\_\allowbreak{}interpolation} constructs the scaled operator and checks all required signs from \texttt{Parameters}.
\end{leanbox}


Earlier results used here: \cite[Section~3.10.1]{LangevinLean2026ULASC}.
\subsection{\texorpdfstring{The shifted identity and the martingale residual constant}{The shifted identity and the martingale residual constant}}
\label{report-step-WassersteinMixing-ShiftedEnergy-lean-2}


\textbf{Claim.} If \(z^+=z-hAz+D\) and \(D=g-g^++M\), then

\[z^++g^+=(z+g)-hA(z+g)+hAg+M.\]


This follows by substituting the two equalities and cancelling the added terms. If \(\langle U,M\rangle=0\), expanding the square gives \(\|U+M\|^2=\|U\|^2+\|M\|^2\).




For any real \(c\), if \(\|D\|^2\le3c\) and \(\|J\|^2\le6c\), then

\[\|D-J\|^2\le18c.\]


Indeed the triangle inequality and \(2ab\le a^2+b^2\) give \(\|D-J\|^2\le2\|D\|^2+2\|J\|^2\le18c\). This is the source's residual constant, derived from its two separate estimates.

\begin{leanbox}{Lean: \texttt{shifted\_\allowbreak{}error\_\allowbreak{}identity}, \texttt{centered\_\allowbreak{}energy\_\allowbreak{}identity}, \texttt{martingale\_\allowbreak{}residual\_\allowbreak{}bound}}
\leansource{\texttt{shifted\_\allowbreak{}error\_\allowbreak{}identity} \quad WassersteinMixing/\allowbreak{}ShiftedEnergy.lean, lines 96–102}
\begin{leancode}
\leanline{0}{theorem\ shifted\_error\_identity\ (A\ :\ H\ →L[ℝ]\ H)\ (h\ :\ ℝ)}
\leanline{4}{(z\ znext\ g\ gnext\ defect\ martingale\ :\ H)}
\leanline{4}{(hstep\ :\ znext\ =\ z\ -\ h\ •\ A\ z\ +\ defect)}
\leanline{4}{(hdefect\ :\ defect\ =\ g\ -\ gnext\ +\ martingale)\ :}
\leanline{4}{znext\ +\ gnext\ =\ (z\ +\ g)\ -\ h\ •\ A\ (z\ +\ g)\ +\ h\ •\ A\ g\ +\ martingale}
\end{leancode}
\leansource{\texttt{centered\_\allowbreak{}energy\_\allowbreak{}identity} \quad WassersteinMixing/\allowbreak{}ShiftedEnergy.lean, lines 105–109}
\begin{leancode}
\leanline{0}{theorem\ centered\_energy\_identity\ (predictable\ martingale\ :\ H)}
\leanline{4}{(horthogonal\ :\ inner\ ℝ\ predictable\ martingale\ =\ 0)\ :}
\leanline{4}{‖predictable\ +\ martingale‖\ \textasciicircum{}\ 2\ =\ ‖predictable‖\ \textasciicircum{}\ 2\ +\ ‖martingale‖\ \textasciicircum{}\ 2}
\end{leancode}
\leansource{\texttt{martingale\_\allowbreak{}residual\_\allowbreak{}bound} \quad WassersteinMixing/\allowbreak{}ShiftedEnergy.lean, lines 113–119}
\begin{leancode}
\leanline{0}{theorem\ martingale\_residual\_bound\ (D\ increment\ :\ H)\ (c\ :\ ℝ)}
\leanline{4}{(hD\ :\ ‖D‖\ \textasciicircum{}\ 2\ ≤\ 3\ *\ c)\ (hg\ :\ ‖increment‖\ \textasciicircum{}\ 2\ ≤\ 6\ *\ c)\ :}
\leanline{4}{‖D\ -\ increment‖\ \textasciicircum{}\ 2\ ≤\ 18\ *\ c}
\end{leancode}
\leanlabel{How the Lean proof follows the argument}
\texttt{shifted\_\allowbreak{}error\_\allowbreak{}identity} proves the vector identity by linearity and module arithmetic. \texttt{centered\_\allowbreak{}energy\_\allowbreak{}identity} expands the actual Hilbert norm. \texttt{martingale\_\allowbreak{}residual\_\allowbreak{}bound} squares the triangle inequality and combines it with the nonnegative square of the difference of the two norms.
\end{leanbox}
\subsection{\texorpdfstring{Conditional expectation supplies the orthogonality}{Conditional expectation supplies the orthogonality}}
\label{report-step-WassersteinMixing-ShiftedEnergy-lean-3}


\textbf{Claim.} Let \(U,M\in L^2(Q;H)\), let \(\mathcal G\) be a sub-\(\sigma\)-algebra, suppose \(U\) is \(\mathcal G\)-measurable and \(\mathbb E[M\mid\mathcal G]=0\). Then \(\langle U,M\rangle_{L^2}=0\). In particular, if \(\|U\|_{L^2}^2\le qE+\ell\), then

\[\|U+M\|_{L^2}^2\le qE+\ell+\|M\|_{L^2}^2.\]


\textbf{Proof.} Conditional expectation is the orthogonal projection onto the closed subspace of \(\mathcal G\)-measurable square-integrable vectors. Hence \(\langle M,U\rangle=\langle\mathbb E[M\mid\mathcal G],U\rangle=0\). Symmetry of the real inner product gives the claimed orientation. Substitute this equality into the squared-norm expansion and apply the predictable energy bound.

\begin{leanbox}{Lean: \texttt{conditional\_\allowbreak{}centering\_\allowbreak{}orthogonal}, \texttt{conditional\_\allowbreak{}shifted\_\allowbreak{}energy}}
\leansource{\texttt{conditional\_\allowbreak{}centering\_\allowbreak{}orthogonal} \quad WassersteinMixing/\allowbreak{}ShiftedEnergy.lean, lines 130–135}
\begin{leancode}
\leanline{0}{theorem\ conditional\_centering\_orthogonal\ \{m\ :\ MeasurableSpace\ Ω\}\ (hm\ :\ m\ ≤\ m0)}
\leanline{4}{(M\ U\ :\ Lp\ H\ 2\ Q)\ (hU\ :\ AEStronglyMeasurable[m]\ U\ Q)}
\leanline{4}{(hM\ :\ (condExpL2\ H\ ℝ\ hm\ M\ :\ Lp\ H\ 2\ Q)\ =\ 0)\ :\ inner\ ℝ\ U\ M\ =\ 0}
\end{leancode}
\leansource{\texttt{conditional\_\allowbreak{}shifted\_\allowbreak{}energy} \quad WassersteinMixing/\allowbreak{}ShiftedEnergy.lean, lines 139–146}
\begin{leancode}
\leanline{0}{theorem\ conditional\_shifted\_energy\ \{m\ :\ MeasurableSpace\ Ω\}\ (hm\ :\ m\ ≤\ m0)}
\leanline{4}{(U\ M\ :\ Lp\ H\ 2\ Q)\ (hU\ :\ AEStronglyMeasurable[m]\ U\ Q)}
\leanline{4}{(hM\ :\ (condExpL2\ H\ ℝ\ hm\ M\ :\ Lp\ H\ 2\ Q)\ =\ 0)}
\leanline{4}{(q\ current\ forcing\ :\ ℝ)\ (hpredictable\ :\ ‖U‖\ \textasciicircum{}\ 2\ ≤\ q\ *\ current\ +\ forcing)\ :}
\leanline{4}{‖U\ +\ M‖\ \textasciicircum{}\ 2\ ≤\ q\ *\ current\ +\ forcing\ +\ ‖M‖\ \textasciicircum{}\ 2}
\end{leancode}
\leanlabel{How the Lean proof follows the argument}
\texttt{conditional\_\allowbreak{}centering\_\allowbreak{}orthogonal} uses \texttt{inner\_\allowbreak{}cond\allowbreak{}Exp\allowbreak{}L2\_\allowbreak{}eq\_\allowbreak{}inner\_\allowbreak{}fun} for the actual conditional-expectation projection. \texttt{conditional\_\allowbreak{}shifted\_\allowbreak{}energy} applies this orthogonality and the exact energy identity before using the predictable estimate. It introduces no new axiom or substitute conditional-expectation operation.
\end{leanbox}
\section{\texorpdfstring{The genuine Hessian and its segment average}{The genuine Hessian and its segment average}}
\label{proof-module-WassersteinMixing-Curvature-lean}


The curvature operator is derived from the actual potential. For \(V\in C^2(H)\), set \(A(x)=D(\nabla V)(x)\) and \(\overline A(x,y)=\int_0^1 A(y+t(x-y))\,dt\). The construction proves symmetry, sharp curvature bounds, continuity and the exact gradient secant identity.

\subsection{\texorpdfstring{The actual second derivative}{The actual second derivative}}
\label{report-step-WassersteinMixing-Curvature-lean-1}


\textbf{Claim.} The gradient is continuously differentiable, \(\langle A(x)u,v\rangle=D^2V(x)[u,v]\), and \(A(x)\) is symmetric and continuous in \(x\).



\textbf{Proof.} Identify the gradient with the inverse Riesz image of the first derivative. Differentiation through this continuous linear identification gives the displayed bilinear form. Equality of mixed second derivatives gives symmetry. Continuous second differentiability gives continuity of the derivative of the gradient.

\begin{leanbox}{Lean: \texttt{contDiff\_\allowbreak{}gradient\_\allowbreak{}of\_\allowbreak{}two}, \texttt{hessianOperator}, \texttt{inner\_\allowbreak{}hessianOperator} and 2 more}
\leansource{\texttt{contDiff\_\allowbreak{}gradient\_\allowbreak{}of\_\allowbreak{}two} \quad WassersteinMixing/\allowbreak{}Curvature.lean, lines 24–27}
\begin{leancode}
\leanline{0}{theorem\ contDiff\_gradient\_of\_two\ (V\ :\ H\ →\ ℝ)\ (hV\ :\ ContDiff\ ℝ\ 2\ V)\ :}
\leanline{4}{ContDiff\ ℝ\ 1\ (gradient\ V)}
\end{leancode}
\leansource{\texttt{hessianOperator} \quad WassersteinMixing/\allowbreak{}Curvature.lean, lines 30–30}
\begin{leancode}
\leanline{0}{def\ hessianOperator\ (V\ :\ H\ →\ ℝ)\ (x\ :\ H)\ :\ H\ →L[ℝ]\ H\ :=\ fderiv\ ℝ\ (gradient\ V)\ x}
\end{leancode}
\leansource{\texttt{inner\_\allowbreak{}hessianOperator} \quad WassersteinMixing/\allowbreak{}Curvature.lean, lines 33–43}
\begin{leancode}
\leanline{0}{theorem\ inner\_hessianOperator\ (V\ :\ H\ →\ ℝ)\ (hV\ :\ ContDiff\ ℝ\ 2\ V)\ (x\ u\ v\ :\ H)\ :}
\leanline{4}{inner\ ℝ\ (hessianOperator\ V\ x\ u)\ v\ =\ fderiv\ ℝ\ (fderiv\ ℝ\ V)\ x\ u\ v}
\end{leancode}
\leansource{\texttt{hessianOperator\_\allowbreak{}symmetric} \quad WassersteinMixing/\allowbreak{}Curvature.lean, lines 46–49}
\begin{leancode}
\leanline{0}{theorem\ hessianOperator\_symmetric\ (V\ :\ H\ →\ ℝ)\ (hV\ :\ ContDiff\ ℝ\ 2\ V)\ (x\ u\ v\ :\ H)\ :}
\leanline{4}{inner\ ℝ\ (hessianOperator\ V\ x\ u)\ v\ =\ inner\ ℝ\ u\ (hessianOperator\ V\ x\ v)}
\end{leancode}
\leansource{\texttt{hessianOperator\_\allowbreak{}continuous} \quad WassersteinMixing/\allowbreak{}Curvature.lean, lines 52–54}
\begin{leancode}
\leanline{0}{theorem\ hessianOperator\_continuous\ (V\ :\ H\ →\ ℝ)\ (hV\ :\ ContDiff\ ℝ\ 2\ V)\ :}
\leanline{4}{Continuous\ (hessianOperator\ V)}
\end{leancode}
\leanlabel{How the Lean proof follows the argument}
\texttt{cont\allowbreak{}Diff\_\allowbreak{}gradient\_\allowbreak{}of\_\allowbreak{}two} differentiates the Riesz representation. \texttt{hessian\allowbreak{}Operator} is the actual Fréchet derivative of the gradient. \texttt{inner\_\allowbreak{}hessian\allowbreak{}Operator} identifies its form with the second derivative; symmetry follows from the second-derivative symmetry theorem.
\end{leanbox}
\subsection{\texorpdfstring{Pointwise curvature bounds}{Pointwise curvature bounds}}
\label{report-step-WassersteinMixing-Curvature-lean-2}


\textbf{Claim.} Under the potential contract, \(\|A(x)\|\le\beta\) and

\[\alpha\|v\|^2\le\langle v,A(x)v\rangle\le\beta\|v\|^2.\]

\textbf{Proof.} A derivative of a \(\beta\)-Lipschitz map has norm at most \(\beta\), giving the upper estimate by Cauchy–Schwarz. Strong convexity makes the gradient \(\alpha\)-strongly monotone. Apply that inequality to \(x+tv\) and \(x\), divide by the positive displacement parameter, and pass to the directional derivative as \(t\downarrow0\) to obtain the lower bound.

\begin{leanbox}{Lean: \texttt{hessianOperator\_\allowbreak{}norm\_\allowbreak{}le}, \texttt{hessianOperator\_\allowbreak{}lower}, \texttt{hessianOperator\_\allowbreak{}upper}}
\leansource{\texttt{hessianOperator\_\allowbreak{}norm\_\allowbreak{}le} \quad WassersteinMixing/\allowbreak{}Curvature.lean, lines 57–59}
\begin{leancode}
\leanline{0}{theorem\ hessianOperator\_norm\_le\ (P\ :\ Parameters)\ (V\ :\ H\ →\ ℝ)}
\leanline{4}{(hV\ :\ PotentialAssumptions\ P\ V)\ (x\ :\ H)\ :\ ‖hessianOperator\ V\ x‖\ ≤\ P.beta}
\end{leancode}
\leansource{\texttt{hessianOperator\_\allowbreak{}lower} \quad WassersteinMixing/\allowbreak{}Curvature.lean, lines 62–85}
\begin{leancode}
\leanline{0}{theorem\ hessianOperator\_lower\ (P\ :\ Parameters)\ (V\ :\ H\ →\ ℝ)}
\leanline{4}{(hV\ :\ PotentialAssumptions\ P\ V)\ (x\ v\ :\ H)\ :}
\leanline{4}{P.alpha\ *\ ‖v‖\ \textasciicircum{}\ 2\ ≤\ inner\ ℝ\ v\ (hessianOperator\ V\ x\ v)}
\end{leancode}
\leansource{\texttt{hessianOperator\_\allowbreak{}upper} \quad WassersteinMixing/\allowbreak{}Curvature.lean, lines 88–94}
\begin{leancode}
\leanline{0}{theorem\ hessianOperator\_upper\ (P\ :\ Parameters)\ (V\ :\ H\ →\ ℝ)}
\leanline{4}{(hV\ :\ PotentialAssumptions\ P\ V)\ (x\ v\ :\ H)\ :}
\leanline{4}{inner\ ℝ\ v\ (hessianOperator\ V\ x\ v)\ ≤\ P.beta\ *\ ‖v‖\ \textasciicircum{}\ 2}
\end{leancode}
\leanlabel{How the Lean proof follows the argument}
\texttt{hessian\allowbreak{}Operator\_\allowbreak{}norm\_\allowbreak{}le} uses the derivative bound of a Lipschitz map. \texttt{hessian\allowbreak{}Operator\_\allowbreak{}lower} takes a one-sided slope limit of the strong-gradient-monotonicity inequality. \texttt{hessian\allowbreak{}Operator\_\allowbreak{}upper} combines the norm estimate with the real inner-product bound.
\end{leanbox}


Earlier results used here: \cite[Section~3.2.3]{LangevinLean2026ULASC}.
\subsection{\texorpdfstring{The exact averaged secant operator}{The exact averaged secant operator}}
\label{report-step-WassersteinMixing-Curvature-lean-3}


\textbf{Claim.} The actual operator integral exists and satisfies

\[\overline A(x,y)(x-y)=\nabla V(x)-\nabla V(y),\qquad \alpha I\preceq\overline A(x,y)\preceq\beta I.\]

\textbf{Proof.} Its integrand is continuous on a compact interval. Evaluation and inner-product functionals commute with the Bochner integral. Integrating the pointwise symmetry and form bounds therefore preserves them, and the norm bound is preserved because the interval has length one. The derivative of the gradient along the segment is \(A(y+t(x-y))(x-y)\); the fundamental theorem gives the secant identity. Joint continuity in the endpoints follows by dominated convergence, with the uniform bound \(\beta\).

\begin{leanbox}{Lean: \texttt{averagedHessian}, \texttt{averagedHessian\_\allowbreak{}curve\_\allowbreak{}continuous}, \texttt{averagedHessian\_\allowbreak{}apply} and 7 more}
\leansource{\texttt{averagedHessian} \quad WassersteinMixing/\allowbreak{}Curvature.lean, lines 97–98}
\begin{leancode}
\leanline{0}{def\ averagedHessian\ (V\ :\ H\ →\ ℝ)\ (x\ y\ :\ H)\ :\ H\ →L[ℝ]\ H\ :=}
\leanline{2}{∫\ t\ in\ (0\ :\ ℝ)..1,\ hessianOperator\ V\ (y\ +\ t\ •\ (x\ -\ y))}
\end{leancode}
\leansource{\texttt{averagedHessian\_\allowbreak{}curve\_\allowbreak{}continuous} \quad WassersteinMixing/\allowbreak{}Curvature.lean, lines 100–102}
\begin{leancode}
\leanline{0}{theorem\ averagedHessian\_curve\_continuous\ (V\ :\ H\ →\ ℝ)\ (hV\ :\ ContDiff\ ℝ\ 2\ V)\ (x\ y\ :\ H)\ :}
\leanline{4}{Continuous\ (fun\ t\ :\ ℝ\ ↦\ hessianOperator\ V\ (y\ +\ t\ •\ (x\ -\ y)))}
\end{leancode}
\leansource{\texttt{averagedHessian\_\allowbreak{}apply} \quad WassersteinMixing/\allowbreak{}Curvature.lean, lines 105–108}
\begin{leancode}
\leanline{0}{theorem\ averagedHessian\_apply\ (V\ :\ H\ →\ ℝ)\ (hV\ :\ ContDiff\ ℝ\ 2\ V)\ (x\ y\ v\ :\ H)\ :}
\leanline{4}{averagedHessian\ V\ x\ y\ v\ =\ ∫\ t\ in\ (0\ :\ ℝ)..1,\ hessianOperator\ V\ (y\ +\ t\ •\ (x\ -\ y))\ v}
\end{leancode}
\leansource{\texttt{averagedHessian\_\allowbreak{}inner} \quad WassersteinMixing/\allowbreak{}Curvature.lean, lines 111–116}
\begin{leancode}
\leanline{0}{theorem\ averagedHessian\_inner\ (V\ :\ H\ →\ ℝ)\ (hV\ :\ ContDiff\ ℝ\ 2\ V)\ (x\ y\ u\ v\ :\ H)\ :}
\leanline{4}{inner\ ℝ\ u\ (averagedHessian\ V\ x\ y\ v)\ =}
\leanline{6}{∫\ t\ in\ (0\ :\ ℝ)..1,\ inner\ ℝ\ u\ (hessianOperator\ V\ (y\ +\ t\ •\ (x\ -\ y))\ v)}
\end{leancode}
\leansource{\texttt{averagedHessian\_\allowbreak{}symmetric} \quad WassersteinMixing/\allowbreak{}Curvature.lean, lines 119–125}
\begin{leancode}
\leanline{0}{theorem\ averagedHessian\_symmetric\ (V\ :\ H\ →\ ℝ)\ (hV\ :\ ContDiff\ ℝ\ 2\ V)\ (x\ y\ u\ v\ :\ H)\ :}
\leanline{4}{inner\ ℝ\ (averagedHessian\ V\ x\ y\ u)\ v\ =\ inner\ ℝ\ u\ (averagedHessian\ V\ x\ y\ v)}
\end{leancode}
\leansource{\texttt{averagedHessian\_\allowbreak{}norm\_\allowbreak{}le} \quad WassersteinMixing/\allowbreak{}Curvature.lean, lines 128–133}
\begin{leancode}
\leanline{0}{theorem\ averagedHessian\_norm\_le\ (P\ :\ Parameters)\ (V\ :\ H\ →\ ℝ)}
\leanline{4}{(hV\ :\ PotentialAssumptions\ P\ V)\ (x\ y\ :\ H)\ :\ ‖averagedHessian\ V\ x\ y‖\ ≤\ P.beta}
\end{leancode}
\leansource{\texttt{averagedHessian\_\allowbreak{}lower} \quad WassersteinMixing/\allowbreak{}Curvature.lean, lines 136–147}
\begin{leancode}
\leanline{0}{theorem\ averagedHessian\_lower\ (P\ :\ Parameters)\ (V\ :\ H\ →\ ℝ)}
\leanline{4}{(hV\ :\ PotentialAssumptions\ P\ V)\ (x\ y\ v\ :\ H)\ :}
\leanline{4}{P.alpha\ *\ ‖v‖\ \textasciicircum{}\ 2\ ≤\ inner\ ℝ\ v\ (averagedHessian\ V\ x\ y\ v)}
\end{leancode}
\leansource{\texttt{averagedHessian\_\allowbreak{}upper} \quad WassersteinMixing/\allowbreak{}Curvature.lean, lines 150–156}
\begin{leancode}
\leanline{0}{theorem\ averagedHessian\_upper\ (P\ :\ Parameters)\ (V\ :\ H\ →\ ℝ)}
\leanline{4}{(hV\ :\ PotentialAssumptions\ P\ V)\ (x\ y\ v\ :\ H)\ :}
\leanline{4}{inner\ ℝ\ v\ (averagedHessian\ V\ x\ y\ v)\ ≤\ P.beta\ *\ ‖v‖\ \textasciicircum{}\ 2}
\end{leancode}
\leansource{\texttt{averagedHessian\_\allowbreak{}secant} \quad WassersteinMixing/\allowbreak{}Curvature.lean, lines 160–170}
\begin{leancode}
\leanline{0}{theorem\ averagedHessian\_secant\ (V\ :\ H\ →\ ℝ)\ (hV\ :\ ContDiff\ ℝ\ 2\ V)\ (x\ y\ :\ H)\ :}
\leanline{4}{averagedHessian\ V\ x\ y\ (x\ -\ y)\ =\ gradient\ V\ x\ -\ gradient\ V\ y}
\end{leancode}
\leansource{\texttt{averagedHessian\_\allowbreak{}continuous} \quad WassersteinMixing/\allowbreak{}Curvature.lean, lines 174–189}
\begin{leancode}
\leanline{0}{theorem\ averagedHessian\_continuous\ (P\ :\ Parameters)\ (V\ :\ H\ →\ ℝ)}
\leanline{4}{(hV\ :\ PotentialAssumptions\ P\ V)\ :}
\leanline{4}{Continuous\ (fun\ p\ :\ H\ ×\ H\ ↦\ averagedHessian\ V\ p.1\ p.2)}
\end{leancode}
\leanlabel{How the Lean proof follows the argument}
\texttt{averaged\allowbreak{}Hessian} is an operator-valued interval integral. The application and pairing lemmas justify moving bounded linear functionals inside it. The secant theorem differentiates the actual gradient curve. Joint continuity uses the proved uniform Hessian bound as a dominated-convergence envelope.
\end{leanbox}
\section{\texorpdfstring{The random averaged Hessian on square-integrable vectors}{The random averaged Hessian on square-integrable vectors}}
\label{proof-module-WassersteinMixing-CurvatureLp-lean}


For actual random states \(X,Y\in L^2(Q;H)\), the field \(\omega\mapsto\overline A(X(\omega),Y(\omega))\) is measurable and bounded by \(\beta\). Multiplication by this field defines a bounded symmetric operator on \(L^2(Q;H)\), with the original curvature constants.



Earlier results used here: \cite[Section~3.13]{LangevinLean2026ULASC}.
\subsection{\texorpdfstring{Constructing multiplication by the actual field}{Constructing multiplication by the actual field}}
\label{report-step-WassersteinMixing-CurvatureLp-lean-1}


\textbf{Construction.} Set \((\mathcal A U)(\omega)=\overline A(X(\omega),Y(\omega))U(\omega)\). Joint continuity of the averaged Hessian gives measurable coefficients; its uniform operator bound gives an essentially bounded field. The estimate \(\|\mathcal A U(\omega)\|\le\beta\|U(\omega)\|\) proves square integrability and a bounded linear action on equivalence classes.

\begin{leanbox}{Lean: \texttt{averagedHessianField}, \texttt{averagedHessianField\_\allowbreak{}memLp}, \texttt{randomCurvature} and 1 more}
\leansource{\texttt{averagedHessianField} \quad WassersteinMixing/\allowbreak{}CurvatureLp.lean, lines 26–29}
\begin{leancode}
\leanline{0}{def\ averagedHessianField\ (X\ Y\ :\ Lp\ H\ 2\ Q)\ (ω\ :\ Ω)\ :\ H\ →L[ℝ]\ H\ :=}
\leanline{2}{averagedHessian\ V\ (X\ ω)\ (Y\ ω)}
\leanline{0}{\strut}
\leanline{0}{include\ P\ hV\ in}
\end{leancode}
\leansource{\texttt{averagedHessianField\_\allowbreak{}memLp} \quad WassersteinMixing/\allowbreak{}CurvatureLp.lean, lines 31–36}
\begin{leancode}
\leanline{0}{theorem\ averagedHessianField\_memLp\ (X\ Y\ :\ Lp\ H\ 2\ Q)\ :}
\leanline{4}{MemLp\ (averagedHessianField\ Q\ V\ X\ Y)\ ⊤\ Q}
\end{leancode}
\leansource{\texttt{randomCurvature} \quad WassersteinMixing/\allowbreak{}CurvatureLp.lean, lines 39–42}
\begin{leancode}
\leanline{0}{def\ randomCurvature\ (X\ Y\ :\ Lp\ H\ 2\ Q)\ :\ Lp\ H\ 2\ Q\ →L[ℝ]\ Lp\ H\ 2\ Q\ :=}
\leanline{2}{(ContinuousLinearMap.id\ ℝ\ (H\ →L[ℝ]\ H)).holderL\ Q\ ⊤\ 2\ 2}
\leanline{4}{((averagedHessianField\_memLp\ Q\ P\ V\ hV\ X\ Y).toLp}
\leanline{6}{(averagedHessianField\ Q\ V\ X\ Y))}
\end{leancode}
\leansource{\texttt{randomCurvature\_\allowbreak{}coe} \quad WassersteinMixing/\allowbreak{}CurvatureLp.lean, lines 45–53}
\begin{leancode}
\leanline{0}{theorem\ randomCurvature\_coe\ (X\ Y\ U\ :\ Lp\ H\ 2\ Q)\ :}
\leanline{4}{randomCurvature\ Q\ P\ V\ hV\ X\ Y\ U\ =ᵐ[Q]\ fun\ ω\ ↦}
\leanline{6}{averagedHessian\ V\ (X\ ω)\ (Y\ ω)\ (U\ ω)}
\end{leancode}
\leanlabel{How the Lean proof follows the argument}
\texttt{averaged\allowbreak{}Hessian\allowbreak{}Field\_\allowbreak{}mem\allowbreak{}Lp} establishes membership at exponent infinity. \texttt{random\allowbreak{}Curvature} uses bounded bilinear evaluation and the Hölder multiplication construction. \texttt{random\allowbreak{}Curvature\_\allowbreak{}coe} proves the exact almost-everywhere pointwise representative.
\end{leanbox}
\subsection{\texorpdfstring{Integrated forms, exact secants and adaptation}{Integrated forms, exact secants and adaptation}}
\label{report-step-WassersteinMixing-CurvatureLp-lean-2}


\textbf{Claim.} The constructed operator satisfies \(\|\mathcal A\|\le\beta\), is symmetric, and obeys

\[\alpha\|U\|_2^2\le\langle U,\mathcal A U\rangle_{L^2}\le\beta\|U\|_2^2,\qquad \mathcal A(X-Y)=\nabla V(X)-\nabla V(Y).\]

\textbf{Proof.} Integrate the pointwise norm and quadratic-form estimates using the exact representative. Pointwise Hessian symmetry gives symmetry of the integral pairing; the secant identity holds almost everywhere. If \(X,Y,U\) are measurable with respect to a sub-sigma-algebra, continuity of the field and evaluation makes their product measurable there too.

\begin{leanbox}{Lean: \texttt{randomCurvature\_\allowbreak{}norm\_\allowbreak{}le}, \texttt{randomCurvature\_\allowbreak{}symmetric}, \texttt{randomCurvature\_\allowbreak{}lower} and 3 more}
\leansource{\texttt{randomCurvature\_\allowbreak{}norm\_\allowbreak{}le} \quad WassersteinMixing/\allowbreak{}CurvatureLp.lean, lines 56–62}
\begin{leancode}
\leanline{0}{theorem\ randomCurvature\_norm\_le\ (X\ Y\ U\ :\ Lp\ H\ 2\ Q)\ :}
\leanline{4}{‖randomCurvature\ Q\ P\ V\ hV\ X\ Y\ U‖\ ≤\ P.beta\ *\ ‖U‖}
\end{leancode}
\leansource{\texttt{randomCurvature\_\allowbreak{}symmetric} \quad WassersteinMixing/\allowbreak{}CurvatureLp.lean, lines 65–73}
\begin{leancode}
\leanline{0}{theorem\ randomCurvature\_symmetric\ (X\ Y\ U\ W\ :\ Lp\ H\ 2\ Q)\ :}
\leanline{4}{inner\ ℝ\ (randomCurvature\ Q\ P\ V\ hV\ X\ Y\ U)\ W\ =}
\leanline{6}{inner\ ℝ\ U\ (randomCurvature\ Q\ P\ V\ hV\ X\ Y\ W)}
\end{leancode}
\leansource{\texttt{randomCurvature\_\allowbreak{}lower} \quad WassersteinMixing/\allowbreak{}CurvatureLp.lean, lines 76–86}
\begin{leancode}
\leanline{0}{theorem\ randomCurvature\_lower\ (X\ Y\ U\ :\ Lp\ H\ 2\ Q)\ :}
\leanline{4}{P.alpha\ *\ ‖U‖\ \textasciicircum{}\ 2\ ≤\ inner\ ℝ\ U\ (randomCurvature\ Q\ P\ V\ hV\ X\ Y\ U)}
\end{leancode}
\leansource{\texttt{randomCurvature\_\allowbreak{}upper} \quad WassersteinMixing/\allowbreak{}CurvatureLp.lean, lines 89–93}
\begin{leancode}
\leanline{0}{theorem\ randomCurvature\_upper\ (X\ Y\ U\ :\ Lp\ H\ 2\ Q)\ :}
\leanline{4}{inner\ ℝ\ U\ (randomCurvature\ Q\ P\ V\ hV\ X\ Y\ U)\ ≤\ P.beta\ *\ ‖U‖\ \textasciicircum{}\ 2}
\end{leancode}
\leansource{\texttt{randomCurvature\_\allowbreak{}secant} \quad WassersteinMixing/\allowbreak{}CurvatureLp.lean, lines 96–109}
\begin{leancode}
\leanline{0}{theorem\ randomCurvature\_secant\ [IsFiniteMeasure\ Q]\ (X\ Y\ :\ Lp\ H\ 2\ Q)\ :}
\leanline{4}{randomCurvature\ Q\ P\ V\ hV\ X\ Y\ (X\ -\ Y)\ =}
\leanline{6}{ULA.lpDrift\ Q\ (gradient\ V)\ hV.gradientLipschitz\ X\ -}
\leanline{8}{ULA.lpDrift\ Q\ (gradient\ V)\ hV.gradientLipschitz\ Y}
\end{leancode}
\leansource{\texttt{randomCurvature\_\allowbreak{}adapted} \quad WassersteinMixing/\allowbreak{}CurvatureLp.lean, lines 113–120}
\begin{leancode}
\leanline{0}{theorem\ randomCurvature\_adapted\ (m\ :\ MeasurableSpace\ Ω)\ (X\ Y\ U\ :\ Lp\ H\ 2\ Q)}
\leanline{4}{(hX\ :\ AEStronglyMeasurable[m]\ X\ Q)\ (hY\ :\ AEStronglyMeasurable[m]\ Y\ Q)}
\leanline{4}{(hU\ :\ AEStronglyMeasurable[m]\ U\ Q)\ :}
\leanline{4}{AEStronglyMeasurable[m]\ (randomCurvature\ (m0\ :=\ m0)\ Q\ P\ V\ hV\ X\ Y\ U)\ Q}
\end{leancode}
\leanlabel{How the Lean proof follows the argument}
The norm and pairing lemmas use the pointwise representative and the integral expression for the real square norm. \texttt{random\allowbreak{}Curvature\_\allowbreak{}secant} invokes the actual segment secant identity. \texttt{random\allowbreak{}Curvature\_\allowbreak{}adapted} proves measurability under the supplied sub-sigma-algebra, without an adaptation field on the operator.
\end{leanbox}


Earlier results used here: \cite[Section~3.13.2]{LangevinLean2026ULASC}; \cite[Section~3.13.3]{LangevinLean2026ULASC}.
\section{\texorpdfstring{The actual synchronous curvature operator}{The actual synchronous curvature operator}}
\label{proof-module-WassersteinMixing-ActualCurvature-lean}


Evaluate the previously constructed random segment Hessian between the actual reference \(Y_{nh}\) and Euler state \(X_n\). Its secant identity turns the true gradient difference in the error update into an actual symmetric bounded operator, retaining the original curvature bounds.

\subsection{\texorpdfstring{Identifying the true error operator}{Identifying the true error operator}}
\label{report-step-WassersteinMixing-ActualCurvature-lean-1}


\textbf{Claim.} With \(A_n(\omega)=\overline A(Y_{nh}(\omega),X_n(\omega))\), the exact error recursion is

\[e_{n+1}=e_n-hA_ne_n+D_n.\]

\textbf{Proof.} The actual random-curvature construction gives a bounded operator on the common \(L^2\) space. Its secant identity is precisely \(A_ne_n=\nabla V(Y_{nh})-\nabla V(X_n)\). Substitution into the previously derived synchronous update gives the displayed equality.

\begin{leanbox}{Lean: \texttt{actualCurvature}, \texttt{actual\_\allowbreak{}curvature\_\allowbreak{}error\_\allowbreak{}step}}
\leansource{\texttt{actualCurvature} \quad WassersteinMixing/\allowbreak{}ActualCurvature.lean, lines 18–22}
\begin{leancode}
\leanline{0}{def\ actualCurvature\ (n\ :\ ℕ)\ :}
\leanline{4}{Lp\ (E\ P.dimension)\ 2\ (actualProbability\ P\ V)\ →L[ℝ]}
\leanline{6}{Lp\ (E\ P.dimension)\ 2\ (actualProbability\ P\ V)\ :=}
\leanline{2}{randomCurvature\ (actualProbability\ P\ V)\ P\ V\ hV}
\leanline{4}{(actualReference\ P\ V\ hV\ n)\ (actualChain\ P\ V\ hV\ n)}
\end{leancode}
\leansource{\texttt{actual\_\allowbreak{}curvature\_\allowbreak{}error\_\allowbreak{}step} \quad WassersteinMixing/\allowbreak{}ActualCurvature.lean, lines 25–37}
\begin{leancode}
\leanline{0}{theorem\ actual\_curvature\_error\_step\ (n\ :\ ℕ)\ :}
\leanline{4}{actualReference\ P\ V\ hV\ (n\ +\ 1)\ -\ actualChain\ P\ V\ hV\ (n\ +\ 1)\ =}
\leanline{6}{(actualReference\ P\ V\ hV\ n\ -\ actualChain\ P\ V\ hV\ n)\ -}
\leanline{8}{P.step\ •\ actualCurvature\ P\ V\ hV\ n\ (actualReference\ P\ V\ hV\ n\ -\ actualChain\ P\ V\ hV\ n)\ +}
\leanline{10}{actualDefect\ P\ V\ hV\ n}
\end{leancode}
\leanlabel{How the Lean proof follows the argument}
\texttt{actual\allowbreak{}Curvature} specializes \texttt{random\allowbreak{}Curvature} to the actual two states. \texttt{actual\_\allowbreak{}curvature\_\allowbreak{}error\_\allowbreak{}step} rewrites the genuine gradient difference using \texttt{random\allowbreak{}Curvature\_\allowbreak{}secant}.
\end{leanbox}


Earlier results used here: \cite[Section~3.13.2]{LangevinLean2026ULASC}.
\subsection{\texorpdfstring{Curvature and measurable-past properties}{Curvature and measurable-past properties}}
\label{report-step-WassersteinMixing-ActualCurvature-lean-2}


\textbf{Claim.} The actual operator is symmetric, \(\alpha\|U\|_2^2\le\langle U,A_nU\rangle\le\beta\|U\|_2^2\), and it preserves measurability with respect to any sub-sigma-algebra that contains both current states.



\textbf{Proof.} These are the integrated properties of the true averaged Hessian. For adaptation, the field depends continuously on the two current states; evaluating it on a measurable vector therefore remains measurable with respect to the same information.

\begin{leanbox}{Lean: \texttt{actualCurvature\_\allowbreak{}symmetric}, \texttt{actualCurvature\_\allowbreak{}lower}, \texttt{actualCurvature\_\allowbreak{}upper} and 1 more}
\leansource{\texttt{actualCurvature\_\allowbreak{}symmetric} \quad WassersteinMixing/\allowbreak{}ActualCurvature.lean, lines 39–42}
\begin{leancode}
\leanline{0}{theorem\ actualCurvature\_symmetric\ (n\ :\ ℕ)}
\leanline{4}{(u\ v\ :\ Lp\ (E\ P.dimension)\ 2\ (actualProbability\ P\ V))\ :}
\leanline{4}{inner\ ℝ\ (actualCurvature\ P\ V\ hV\ n\ u)\ v\ =\ inner\ ℝ\ u\ (actualCurvature\ P\ V\ hV\ n\ v)}
\end{leancode}
\leansource{\texttt{actualCurvature\_\allowbreak{}lower} \quad WassersteinMixing/\allowbreak{}ActualCurvature.lean, lines 44–46}
\begin{leancode}
\leanline{0}{theorem\ actualCurvature\_lower\ (n\ :\ ℕ)\ (u\ :\ Lp\ (E\ P.dimension)\ 2\ (actualProbability\ P\ V))\ :}
\leanline{4}{P.alpha\ *\ ‖u‖\ \textasciicircum{}\ 2\ ≤\ inner\ ℝ\ u\ (actualCurvature\ P\ V\ hV\ n\ u)}
\end{leancode}
\leansource{\texttt{actualCurvature\_\allowbreak{}upper} \quad WassersteinMixing/\allowbreak{}ActualCurvature.lean, lines 48–50}
\begin{leancode}
\leanline{0}{theorem\ actualCurvature\_upper\ (n\ :\ ℕ)\ (u\ :\ Lp\ (E\ P.dimension)\ 2\ (actualProbability\ P\ V))\ :}
\leanline{4}{inner\ ℝ\ u\ (actualCurvature\ P\ V\ hV\ n\ u)\ ≤\ P.beta\ *\ ‖u‖\ \textasciicircum{}\ 2}
\end{leancode}
\leansource{\texttt{actualCurvature\_\allowbreak{}adapted} \quad WassersteinMixing/\allowbreak{}ActualCurvature.lean, lines 54–61}
\begin{leancode}
\leanline{0}{theorem\ actualCurvature\_adapted\ (m\ :\ MeasurableSpace\ (ActualSpace\ P))\ (n\ :\ ℕ)}
\leanline{4}{(u\ :\ Lp\ (E\ P.dimension)\ 2\ (actualProbability\ P\ V))}
\leanline{4}{(hX\ :\ AEStronglyMeasurable[m]\ (actualReference\ P\ V\ hV\ n)\ (actualProbability\ P\ V))}
\leanline{4}{(hY\ :\ AEStronglyMeasurable[m]\ (actualChain\ P\ V\ hV\ n)\ (actualProbability\ P\ V))}
\leanline{4}{(hu\ :\ AEStronglyMeasurable[m]\ u\ (actualProbability\ P\ V))\ :}
\leanline{4}{AEStronglyMeasurable[m]\ (actualCurvature\ P\ V\ hV\ n\ u)\ (actualProbability\ P\ V)}
\end{leancode}
\leanlabel{How the Lean proof follows the argument}
The symmetry and form estimates specialize the proved random-curvature theorems. \texttt{actual\allowbreak{}Curvature\_\allowbreak{}adapted} applies their sub-sigma-algebra measurability result to the actual reference and chain.
\end{leanbox}
\section{\texorpdfstring{Adaptation to the genuine Brownian past}{Adaptation to the genuine Brownian past}}
\label{proof-module-WassersteinMixing-ActualAdapted-lean}


The filtration at step \(n\) contains the actual initialization and all additive Brownian input up to \(nh\). Adaptation of the diffusion is obtained from its integral construction; adaptation of the Euler chain follows from its recursion. The actual curvature consequently preserves this same past.



Earlier results used here: \cite[Section~3.54]{LangevinLean2026ULASC}.
\subsection{\texorpdfstring{Identifying the common initialization and path}{Identifying the common initialization and path}}
\label{report-step-WassersteinMixing-ActualAdapted-lean-1}


\textbf{Claim.} The Euler and global-diffusion initial variables agree, and the actual grid reference equals the endpoint of any canonical solution whose horizon reaches that grid time.



\textbf{Proof.} The diagonal initial coupling has zero squared difference, so its two \(L^2\) variables agree. Uniqueness of the Langevin integral equation identifies the segmented reference with the canonical finite-horizon solution. The additive Brownian input is the actual noise measured from time zero.

\begin{leanbox}{Lean: \texttt{actualBrownianNoise}, \texttt{actualInitial\_\allowbreak{}eq\_\allowbreak{}semigroupInitial}, \texttt{actualReference\_\allowbreak{}eq\_\allowbreak{}solutionOn}}
\leansource{\texttt{actualBrownianNoise} \quad WassersteinMixing/\allowbreak{}ActualAdapted.lean, lines 24–26}
\begin{leancode}
\leanline{0}{def\ actualBrownianNoise\ (t\ :\ ℝ)\ :\ Lp\ (E\ P.dimension)\ 2\ (actualProbability\ P\ V)\ :=}
\leanline{2}{ULA.brownianLangevinNoise\ (actualProbability\ P\ V)\ (CanonicalSemigroup.brownian\ P)}
\leanline{4}{(CanonicalSemigroup.brownian\_preBrownian\ P\ V\ hV)\ 0\ t}
\end{leancode}
\leansource{\texttt{actualInitial\_\allowbreak{}eq\_\allowbreak{}semigroupInitial} \quad WassersteinMixing/\allowbreak{}ActualAdapted.lean, lines 29–40}
\begin{leancode}
\leanline{0}{theorem\ actualInitial\_eq\_semigroupInitial\ :}
\leanline{4}{actualInitial\ P\ V\ hV\ =\ CanonicalSemigroup.initial\ P\ V\ hV}
\end{leancode}
\leansource{\texttt{actualReference\_\allowbreak{}eq\_\allowbreak{}solutionOn} \quad WassersteinMixing/\allowbreak{}ActualAdapted.lean, lines 44–58}
\begin{leancode}
\leanline{0}{theorem\ actualReference\_eq\_solutionOn\ (n\ :\ ℕ)\ (T\ :\ ℝ)\ (hT\ :\ 0\ ≤\ T)}
\leanline{4}{(hn\ :\ (n\ :\ ℝ)\ *\ P.step\ ≤\ T)\ :}
\leanline{4}{actualReference\ P\ V\ hV\ n\ =\ CanonicalSemigroup.solutionOn\ P\ V\ hV\ T\ hT\ ((n\ :\ ℝ)\ *\ P.step)}
\end{leancode}
\leanlabel{How the Lean proof follows the argument}
\texttt{actual\allowbreak{}Initial\_\allowbreak{}eq\_\allowbreak{}semigroup\allowbreak{}Initial} uses the exact diagonal coupling cost. \texttt{actual\allowbreak{}Reference\_\allowbreak{}eq\_\allowbreak{}solution\allowbreak{}On} invokes uniqueness for the actual Brownian integral equation, so no path-identification premise is added.
\end{leanbox}


Earlier results used here: \cite[Section~3.17.2]{LangevinLean2026ULASC}; \cite[Section~3.22.23]{LangevinLean2026ULASC}; \cite[Section~3.25.1]{LangevinLean2026ULASC}; \cite[Section~3.29.4]{LangevinLean2026ULASC}; \cite[Section~3.29.5]{LangevinLean2026ULASC}; \cite[Section~3.44.4]{LangevinLean2026ULASC}; \cite[Section~3.8.17]{LangevinLean2026ULASC}.
\subsection{\texorpdfstring{Constructing the increasing filtration}{Constructing the increasing filtration}}
\label{report-step-WassersteinMixing-ActualAdapted-lean-2}


\textbf{Construction.} Let \(\mathcal F_n=\sigma(Y_0,B_s:0\le s\le nh)\), interpreted with the actual additive Brownian representatives. Inclusion of time intervals gives \(\mathcal F_n\subseteq\mathcal F_{n+1}\). The constructed reference belongs to \(L^2(\mathcal F_n)\).



\textbf{Proof.} The initial variable and each permitted Brownian evaluation generate the stated sub-sigma-algebra. The Picard solution remains in the closed subspace of random variables measurable with respect to that past; uniqueness identifies it with the reference endpoint.

\begin{leanbox}{Lean: \texttt{actualPast}, \texttt{actualPast\_\allowbreak{}le}, \texttt{actualPast\_\allowbreak{}mono} and 2 more}
\leansource{\texttt{actualPast} \quad WassersteinMixing/\allowbreak{}ActualAdapted.lean, lines 61–63}
\begin{leancode}
\leanline{0}{@[instance\_reducible]\ def\ actualPast\ (n\ :\ ℕ)\ :\ MeasurableSpace\ (ActualSpace\ P)\ :=}
\leanline{2}{ULA.lpLangevinPast\ (actualProbability\ P\ V)\ (CanonicalSemigroup.initial\ P\ V\ hV)}
\leanline{4}{(actualBrownianNoise\ P\ V\ hV)\ ((n\ :\ ℝ)\ *\ P.step)}
\end{leancode}
\leansource{\texttt{actualPast\_\allowbreak{}le} \quad WassersteinMixing/\allowbreak{}ActualAdapted.lean, lines 65–66}
\begin{leancode}
\leanline{0}{theorem\ actualPast\_le\ (n\ :\ ℕ)\ :\ actualPast\ P\ V\ hV\ n\ ≤\ (inferInstance\ :\ MeasurableSpace\ (ActualSpace\ P))}
\end{leancode}
\leansource{\texttt{actualPast\_\allowbreak{}mono} \quad WassersteinMixing/\allowbreak{}ActualAdapted.lean, lines 68–77}
\begin{leancode}
\leanline{0}{theorem\ actualPast\_mono\ :\ Monotone\ (actualPast\ P\ V\ hV)}
\end{leancode}
\leansource{\texttt{actualFiltration} \quad WassersteinMixing/\allowbreak{}ActualAdapted.lean, lines 80–83}
\begin{leancode}
\leanline{0}{def\ actualFiltration\ :\ Filtration\ ℕ\ (inferInstance\ :\ MeasurableSpace\ (ActualSpace\ P))\ where}
\leanline{2}{seq\ :=\ actualPast\ P\ V\ hV}
\leanline{2}{mono'\ :=\ actualPast\_mono\ P\ V\ hV}
\leanline{2}{le'\ :=\ actualPast\_le\ P\ V\ hV}
\end{leancode}
\leansource{\texttt{actualReference\_\allowbreak{}adapted} \quad WassersteinMixing/\allowbreak{}ActualAdapted.lean, lines 86–98}
\begin{leancode}
\leanline{0}{theorem\ actualReference\_adapted\ (n\ :\ ℕ)\ :}
\leanline{4}{AEStronglyMeasurable[actualPast\ P\ V\ hV\ n]\ (actualReference\ P\ V\ hV\ n)\ (actualProbability\ P\ V)}
\end{leancode}
\leanlabel{How the Lean proof follows the argument}
\texttt{actual\allowbreak{}Past} specializes the natural-past construction. \texttt{actual\allowbreak{}Past\_\allowbreak{}mono} proves inclusion of every generating time evaluation and \texttt{actual\allowbreak{}Filtration} packages it. \texttt{actual\allowbreak{}Reference\_\allowbreak{}adapted} uses the proved natural-past theorem for the actual integral solution.
\end{leanbox}


Earlier results used here: \cite[Section~3.25.2]{LangevinLean2026ULASC}; \cite[Section~3.25.3]{LangevinLean2026ULASC}; \cite[Section~3.54.4]{LangevinLean2026ULASC}; \cite[Section~3.54.5]{LangevinLean2026ULASC}; \cite[Section~3.54.8]{LangevinLean2026ULASC}.
\subsection{\texorpdfstring{Adaptation of innovations, chain and curvature}{Adaptation of innovations, chain and curvature}}
\label{report-step-WassersteinMixing-ActualAdapted-lean-3}


\textbf{Claim.} The normalized innovation is the difference of the actual Brownian inputs at \((n+1)h\) and \(nh\), divided by \(\sqrt{2h}\), hence is measurable by \(\mathcal F_{n+1}\). The Euler state and synchronous error are \(\mathcal F_n\)-measurable, and \(A_n\) maps \(L^2(\mathcal F_n)\) into itself.



\textbf{Proof.} Use the exact increment identity. Induct through the Euler update using continuity of the gradient and the increasing filtration. Subtracting adapted states gives the adapted error. The curvature-preservation theorem then applies because both of its endpoint states are adapted.

\begin{leanbox}{Lean: \texttt{actualNoise\_\allowbreak{}eq\_\allowbreak{}global\_\allowbreak{}increment}, \texttt{actualNoise\_\allowbreak{}adapted}, \texttt{actualChain\_\allowbreak{}adapted} and 2 more}
\leansource{\texttt{actualNoise\_\allowbreak{}eq\_\allowbreak{}global\_\allowbreak{}increment} \quad WassersteinMixing/\allowbreak{}ActualAdapted.lean, lines 102–117}
\begin{leancode}
\leanline{0}{theorem\ actualNoise\_eq\_global\_increment\ (n\ :\ ℕ)\ :}
\leanline{4}{actualNoise\ P\ V\ hV\ n\ =\ (Real.sqrt\ (2\ *\ P.step))⁻¹\ •}
\leanline{6}{(actualBrownianNoise\ P\ V\ hV\ (((n\ +\ 1\ :\ ℕ)\ :\ ℝ)\ *\ P.step)\ -}
\leanline{8}{actualBrownianNoise\ P\ V\ hV\ ((n\ :\ ℝ)\ *\ P.step))}
\end{leancode}
\leansource{\texttt{actualNoise\_\allowbreak{}adapted} \quad WassersteinMixing/\allowbreak{}ActualAdapted.lean, lines 120–140}
\begin{leancode}
\leanline{0}{theorem\ actualNoise\_adapted\ (n\ :\ ℕ)\ :}
\leanline{4}{AEStronglyMeasurable[actualPast\ P\ V\ hV\ (n\ +\ 1)]\ (actualNoise\ P\ V\ hV\ n)\ (actualProbability\ P\ V)}
\end{leancode}
\leansource{\texttt{actualChain\_\allowbreak{}adapted} \quad WassersteinMixing/\allowbreak{}ActualAdapted.lean, lines 144–168}
\begin{leancode}
\leanline{0}{theorem\ actualChain\_adapted\ (n\ :\ ℕ)\ :}
\leanline{4}{AEStronglyMeasurable[actualPast\ P\ V\ hV\ n]\ (actualChain\ P\ V\ hV\ n)\ (actualProbability\ P\ V)}
\end{leancode}
\leansource{\texttt{actualError\_\allowbreak{}adapted} \quad WassersteinMixing/\allowbreak{}ActualAdapted.lean, lines 171–179}
\begin{leancode}
\leanline{0}{theorem\ actualError\_adapted\ (n\ :\ ℕ)\ :}
\leanline{4}{AEStronglyMeasurable[actualPast\ P\ V\ hV\ n]}
\leanline{6}{(actualReference\ P\ V\ hV\ n\ -\ actualChain\ P\ V\ hV\ n)\ (actualProbability\ P\ V)}
\end{leancode}
\leansource{\texttt{actualCurvature\_\allowbreak{}preserves\_\allowbreak{}past} \quad WassersteinMixing/\allowbreak{}ActualAdapted.lean, lines 182–187}
\begin{leancode}
\leanline{0}{theorem\ actualCurvature\_preserves\_past\ (n\ :\ ℕ)}
\leanline{4}{(u\ :\ Lp\ (E\ P.dimension)\ 2\ (actualProbability\ P\ V))}
\leanline{4}{(hu\ :\ AEStronglyMeasurable[actualPast\ P\ V\ hV\ n]\ u\ (actualProbability\ P\ V))\ :}
\leanline{4}{AEStronglyMeasurable[actualPast\ P\ V\ hV\ n]\ (actualCurvature\ P\ V\ hV\ n\ u)\ (actualProbability\ P\ V)}
\end{leancode}
\leanlabel{How the Lean proof follows the argument}
The proof works inside the closed measurable submodule \texttt{lp\allowbreak{}Meas}, whose closure under sums and scalar multiplication handles the recursion. \texttt{actual\allowbreak{}Curvature\_\allowbreak{}preserves\_\allowbreak{}past} supplies the derived adaptation of both actual endpoint states to the actual random-curvature theorem.
\end{leanbox}


Earlier results used here: \cite[Section~3.13.2]{LangevinLean2026ULASC}; \cite[Section~3.25.1]{LangevinLean2026ULASC}; \cite[Section~3.32.13]{LangevinLean2026ULASC}; \cite[Section~3.44.1]{LangevinLean2026ULASC}; \cite[Section~3.54.2]{LangevinLean2026ULASC}; \cite[Section~3.54.6]{LangevinLean2026ULASC}; \cite[Section~3.54.7]{LangevinLean2026ULASC}.
\section{\texorpdfstring{The semigroup law from the actual conditional Markov identity}{The semigroup law from the actual conditional Markov identity}}
\label{proof-module-WassersteinMixing-SemigroupComposition-lean}


The conditional Markov identity on the actual initial/\allowbreak{}Brownian past implies the composition law for the canonical real observable contractions. Its proof uses actual adaptedness of time-zero evaluation and the defining adjoint pairings.

\subsection{\texorpdfstring{Conditioning the actual two-time pairing}{Conditioning the actual two-time pairing}}
\label{report-step-WassersteinMixing-SemigroupComposition-lean-1}


\textbf{Claim.} If \(\mathbb E[U_{a+t}f\mid\mathcal F_a]=U_aP_tf\), then \(P_{a+t}=P_aP_t\).



\textbf{Proof.} Time-zero evaluation is measurable with respect to the actual past because that past contains the initialization. Pair the conditional identity with \(U_0g\). The conditional-expectation projection preserves its pairing with this adapted vector. The adjoint compression identities identify the resulting sides with \(\langle P_{a+t}f,g\rangle\) and \(\langle P_aP_tf,g\rangle\). Equality against all test observables proves the operator identity.

\begin{leanbox}{Lean: \texttt{evaluation\_\allowbreak{}zero\_\allowbreak{}adapted\_\allowbreak{}pastAt}, \texttt{op\_\allowbreak{}add\_\allowbreak{}of\_\allowbreak{}markov}}
\leansource{\texttt{evaluation\_\allowbreak{}zero\_\allowbreak{}adapted\_\allowbreak{}pastAt} \quad WassersteinMixing/\allowbreak{}SemigroupComposition.lean, lines 24–39}
\begin{leancode}
\leanline{0}{theorem\ evaluation\_zero\_adapted\_pastAt\ (P\ :\ Parameters)\ (V\ :\ E\ P.dimension\ →\ ℝ)}
\leanline{4}{(hV\ :\ PotentialAssumptions\ P\ V)\ (a\ :\ ℝ)\ (f\ :\ Lp\ F\ 2\ (ULA.gibbsMeasure\ V))\ :}
\leanline{4}{AEStronglyMeasurable[ULA.lpLangevinPast\ (actualProbability\ P\ V)}
\leanline{6}{(initial\ P\ V\ hV)\ (actualBrownianNoise\ P\ V\ hV)\ a]}
\leanline{6}{(evaluation\ ({\leanfallback 𝕜}\ :=\ ℝ)\ P\ V\ hV\ 0\ f)\ (actualProbability\ P\ V)}
\end{leancode}
\leansource{\texttt{op\_\allowbreak{}add\_\allowbreak{}of\_\allowbreak{}markov} \quad WassersteinMixing/\allowbreak{}SemigroupComposition.lean, lines 44–66}
\begin{leancode}
\leanline{0}{theorem\ op\_add\_of\_markov\ (P\ :\ Parameters)\ (V\ :\ E\ P.dimension\ →\ ℝ)}
\leanline{4}{(hV\ :\ PotentialAssumptions\ P\ V)\ (a\ t\ :\ ℝ)}
\leanline{4}{(hmarkov\ :\ ∀\ f\ :\ Lp\ F\ 2\ (ULA.gibbsMeasure\ V),}
\leanline{6}{(condExpL2\ F\ ℝ\ (ULA.lpLangevinPast\_le\ (actualProbability\ P\ V)}
\leanline{8}{(initial\ P\ V\ hV)\ (actualBrownianNoise\ P\ V\ hV)\ a)}
\leanline{8}{(evaluation\ ({\leanfallback 𝕜}\ :=\ ℝ)\ P\ V\ hV\ (a\ +\ t)\ f)\ :\ Lp\ F\ 2\ (actualProbability\ P\ V))\ =}
\leanline{8}{evaluation\ ({\leanfallback 𝕜}\ :=\ ℝ)\ P\ V\ hV\ a\ (op\ ({\leanfallback 𝕜}\ :=\ ℝ)\ P\ V\ hV\ t\ f))\ :}
\leanline{4}{op\ ({\leanfallback 𝕜}\ :=\ ℝ)\ (F\ :=\ F)\ P\ V\ hV\ (a\ +\ t)\ =\ op\ P\ V\ hV\ a\ ∘L\ op\ P\ V\ hV\ t}
\end{leancode}
\leanlabel{How the Lean proof follows the argument}
\texttt{evaluation\_\allowbreak{}zero\_\allowbreak{}adapted\_\allowbreak{}past\allowbreak{}At} uses the actual initial-state generator of the past. \texttt{op\_\allowbreak{}add\_\allowbreak{}of\_\allowbreak{}markov} applies the real conditional-expectation pairing theorem and the exact evaluation compression identities, then uses inner-product extensionality.
\end{leanbox}


Earlier results used here: \cite[Section~3.26.3]{LangevinLean2026ULASC}; \cite[Section~3.54.4]{LangevinLean2026ULASC}; \cite[Section~3.54.5]{LangevinLean2026ULASC}; \cite[Section~3.54.6]{LangevinLean2026ULASC}.
\section{\texorpdfstring{Independence of the entire Brownian future and past}{Independence of the entire Brownian future and past}}
\label{proof-module-WassersteinMixing-ActualMarkovNoise-lean}


At a deterministic cut time \(a\ge0\), all future coordinate increments \((B^i_{a+t}-B^i_a)_{i,t}\) are independent of all past coordinates \((B^i_{t\wedge a})_{i,t}\). Adjoining the independent initial state preserves this whole-path independence.

\subsection{\texorpdfstring{Regrouping independent pairs into independent families}{Regrouping independent pairs into independent families}}
\label{report-step-WassersteinMixing-ActualMarkovNoise-lean-1}


\textbf{Claim.} Suppose the rows \((X_i,Y_i)\) are mutually independent and, within each row, \(X_i\) and \(Y_i\) are independent. Then the entire families \((X_i)_i\) and \((Y_i)_i\) are independent.



\textbf{Proof.} Regard each row as a two-coordinate independent family. Independence of the rows gives independence of the family indexed by the product of the row set and the two-point coordinate set. The two collections occupy disjoint index sets, so their generated sigma-algebras are independent. This argument does not require countability of the row set.

\begin{leanbox}{Lean: \texttt{iIndep\_\allowbreak{}bool\_\allowbreak{}of\_\allowbreak{}indep}, \texttt{indepFun\_\allowbreak{}families\_\allowbreak{}of\_\allowbreak{}independent\_\allowbreak{}pairs}}
\leansource{\texttt{iIndep\_\allowbreak{}bool\_\allowbreak{}of\_\allowbreak{}indep} \quad WassersteinMixing/\allowbreak{}ActualMarkovNoise.lean, lines 20–36}
\begin{leancode}
\leanline{0}{private\ theorem\ iIndep\_bool\_of\_indep\ \{Ω\ :\ Type*\}\ [MeasurableSpace\ Ω]}
\leanline{4}{(Q\ :\ Measure\ Ω)\ [IsProbabilityMeasure\ Q]\ (m\ :\ Bool\ →\ MeasurableSpace\ Ω)}
\leanline{4}{(hi\ :\ Indep\ (m\ false)\ (m\ true)\ Q)\ :\ iIndep\ m\ Q}
\end{leancode}
\leansource{\texttt{indepFun\_\allowbreak{}families\_\allowbreak{}of\_\allowbreak{}independent\_\allowbreak{}pairs} \quad WassersteinMixing/\allowbreak{}ActualMarkovNoise.lean, lines 41–83}
\begin{leancode}
\leanline{0}{theorem\ indepFun\_families\_of\_independent\_pairs\ \{Ω\ I\ A\ :\ Type*\}}
\leanline{4}{[MeasurableSpace\ Ω]\ [MeasurableSpace\ A]}
\leanline{4}{(Q\ :\ Measure\ Ω)\ [IsProbabilityMeasure\ Q]\ (X\ Y\ :\ I\ →\ Ω\ →\ A)}
\leanline{4}{(hX\ :\ ∀\ i,\ Measurable\ (X\ i))\ (hY\ :\ ∀\ i,\ Measurable\ (Y\ i))}
\leanline{4}{(hrows\ :\ iIndepFun\ (fun\ i\ ω\ ↦\ (X\ i\ ω,\ Y\ i\ ω))\ Q)}
\leanline{4}{(hpair\ :\ ∀\ i,\ IndepFun\ (X\ i)\ (Y\ i)\ Q)\ :}
\leanline{4}{IndepFun\ (fun\ ω\ i\ ↦\ X\ i\ ω)\ (fun\ ω\ i\ ↦\ Y\ i\ ω)\ Q}
\end{leancode}
\leanlabel{How the Lean proof follows the argument}
The Boolean helper expands all finite subsets of a two-point index. \texttt{indep\allowbreak{}Fun\_\allowbreak{}families\_\allowbreak{}of\_\allowbreak{}independent\_\allowbreak{}pairs} combines row and within-row independence by uncurry, then applies independence of sigma-algebras generated by disjoint index subsets.
\end{leanbox}
\subsection{\texorpdfstring{Whole-path independence on the actual product source}{Whole-path independence on the actual product source}}
\label{report-step-WassersteinMixing-ActualMarkovNoise-lean-2}


\textbf{Claim.} The actual Brownian future at \(a\) is independent of the initial state together with the entire Brownian past.



\textbf{Proof.} For each coordinate, independent Brownian increments give independence of the shifted future path and the path restricted to \([0,a]\). Compose the latter with time clipping. Coordinate paths are mutually independent, so the preceding regrouping result gives vector whole-path independence. Finally take the product with the independent initialization: integrating rectangle events over the initial coordinate leaves the same future marginal and factorizes their masses.

\begin{leanbox}{Lean: \texttt{canonicalBrownian\_\allowbreak{}future\_\allowbreak{}past\_\allowbreak{}independent}, \texttt{indepFun\_\allowbreak{}future\_\allowbreak{}initial\_\allowbreak{}past}, \texttt{actualBrownian\_\allowbreak{}raw\_\allowbreak{}future\_\allowbreak{}initial\_\allowbreak{}past\_\allowbreak{}independent}}
\leansource{\texttt{canonicalBrownian\_\allowbreak{}future\_\allowbreak{}past\_\allowbreak{}independent} \quad WassersteinMixing/\allowbreak{}ActualMarkovNoise.lean, lines 87–115}
\begin{leancode}
\leanline{0}{theorem\ canonicalBrownian\_future\_past\_independent\ (I\ :\ Type*)\ (a\ :\ ℝ≥0)\ :}
\leanline{4}{IndepFun}
\leanline{6}{(fun\ (ω\ :\ ULA.CanonicalBrownian.VectorSample\ I)\ i\ t\ ↦}
\leanline{8}{ULA.CanonicalBrownian.coordinate\ i\ (a\ +\ t)\ ω\ -\ ULA.CanonicalBrownian.coordinate\ i\ a\ ω)}
\leanline{6}{(fun\ (ω\ :\ ULA.CanonicalBrownian.VectorSample\ I)\ i\ t\ ↦}
\leanline{8}{ULA.CanonicalBrownian.coordinate\ i\ (min\ t\ a)\ ω)}
\leanline{6}{(ULA.CanonicalBrownian.vectorProbability\ I)}
\end{leancode}
\leansource{\texttt{indepFun\_\allowbreak{}future\_\allowbreak{}initial\_\allowbreak{}past} \quad WassersteinMixing/\allowbreak{}ActualMarkovNoise.lean, lines 119–139}
\begin{leancode}
\leanline{0}{theorem\ indepFun\_future\_initial\_past\ \{Ω₀\ Ω\ A\ B\ :\ Type*\}}
\leanline{4}{[MeasurableSpace\ Ω₀]\ [MeasurableSpace\ Ω]\ [MeasurableSpace\ A]\ [MeasurableSpace\ B]}
\leanline{4}{(μ\ :\ Measure\ Ω₀)\ (ν\ :\ Measure\ Ω)\ [IsProbabilityMeasure\ μ]\ [IsProbabilityMeasure\ ν]}
\leanline{4}{(X\ :\ Ω\ →\ A)\ (Y\ :\ Ω\ →\ B)\ (hX\ :\ Measurable\ X)\ (hY\ :\ Measurable\ Y)}
\leanline{4}{(hi\ :\ IndepFun\ X\ Y\ ν)\ :}
\leanline{4}{IndepFun\ (fun\ ω\ :\ Ω₀\ ×\ Ω\ ↦\ X\ ω.2)\ (fun\ ω\ :\ Ω₀\ ×\ Ω\ ↦\ (ω.1,\ Y\ ω.2))\ (μ.prod\ ν)}
\end{leancode}
\leansource{\texttt{actualBrownian\_\allowbreak{}raw\_\allowbreak{}future\_\allowbreak{}initial\_\allowbreak{}past\_\allowbreak{}independent} \quad WassersteinMixing/\allowbreak{}ActualMarkovNoise.lean, lines 142–159}
\begin{leancode}
\leanline{0}{theorem\ actualBrownian\_raw\_future\_initial\_past\_independent}
\leanline{4}{(P\ :\ Parameters)\ (V\ :\ E\ P.dimension\ →\ ℝ)\ (hV\ :\ PotentialAssumptions\ P\ V)\ (a\ :\ ℝ≥0)\ :}
\leanline{4}{IndepFun}
\leanline{6}{(fun\ (ω\ :\ ActualSpace\ P)\ i\ t\ ↦}
\leanline{8}{CanonicalSemigroup.brownian\ P\ i\ (a\ +\ t)\ ω\ -\ CanonicalSemigroup.brownian\ P\ i\ a\ ω)}
\leanline{6}{(fun\ (ω\ :\ ActualSpace\ P)\ ↦\ (ω.1,\ fun\ i\ t\ ↦\ CanonicalSemigroup.brownian\ P\ i\ (min\ t\ a)\ ω))}
\leanline{6}{(actualProbability\ P\ V)}
\end{leancode}
\leanlabel{How the Lean proof follows the argument}
\texttt{canonical\allowbreak{}Brownian\_\allowbreak{}future\_\allowbreak{}past\_\allowbreak{}independent} combines the actual coordinate pre-Brownian shift theorem with coordinate-path independence. \texttt{indep\allowbreak{}Fun\_\allowbreak{}future\_\allowbreak{}initial\_\allowbreak{}past} proves the product extension directly by product integration. The final theorem specializes these facts to \texttt{actual\allowbreak{}Probability}.
\end{leanbox}


Earlier results used here: \cite[Section~3.22.1]{LangevinLean2026ULASC}; \cite[Section~3.8.11]{LangevinLean2026ULASC}; \cite[Section~3.8.16]{LangevinLean2026ULASC}; \cite[Section~3.8.17]{LangevinLean2026ULASC}; \cite[Section~3.8.19]{LangevinLean2026ULASC}; \cite[Section~3.8.20]{LangevinLean2026ULASC}; \cite[Section~3.8.21]{LangevinLean2026ULASC}.
\section{\texorpdfstring{Conditional averaging over an independent path input}{Conditional averaging over an independent path input}}
\label{proof-module-WassersteinMixing-ConditionalIndependentAverage-lean}


Let \(X\) be measurable from a past sigma-algebra \(\mathcal F\), and let \(Z\) be independent of all of \(\mathcal F\). For an integrable measurable vector function \(F\), conditioning \(F(X,Z)\) on the past integrates only the independent input. The input space may be a whole path space.

\subsection{\texorpdfstring{Factorization after restriction to a past event}{Factorization after restriction to a past event}}
\label{report-step-WassersteinMixing-ConditionalIndependentAverage-lean-1}


\textbf{Claim.} For every \(A\in\mathcal F\), the joint law under the restricted measure satisfies

\[\operatorname{Law}_{Q|_A}(X,Z)=\operatorname{Law}_{Q|_A}(X)\otimes\operatorname{Law}_{Q}(Z).\]

\textbf{Proof.} Evaluate both sides on a measurable rectangle. Its preimage is the intersection of a past event, obtained from \(X\) and \(A\), with an event determined by \(Z\). Independence factorizes its mass. Equality on all rectangles determines the product measure.

\begin{leanbox}{Lean: \texttt{independent\_\allowbreak{}input\_\allowbreak{}restrict\_\allowbreak{}joint\_\allowbreak{}map}}
\leansource{\texttt{independent\_\allowbreak{}input\_\allowbreak{}restrict\_\allowbreak{}joint\_\allowbreak{}map} \quad WassersteinMixing/\allowbreak{}ConditionalIndependentAverage.lean, lines 19–35}
\begin{leancode}
\leanline{0}{theorem\ independent\_input\_restrict\_joint\_map\ \{Ω\ A\ B\ :\ Type*\}}
\leanline{4}{(m\ :\ MeasurableSpace\ Ω)\ [m0\ :\ MeasurableSpace\ Ω]\ [MeasurableSpace\ A]\ [MeasurableSpace\ B]}
\leanline{4}{(Q\ :\ Measure\ Ω)\ [IsProbabilityMeasure\ Q]\ (hm\ :\ m\ ≤\ m0)}
\leanline{4}{(X\ :\ Ω\ →\ A)\ (Y\ :\ Ω\ →\ B)\ (hX\ :\ Measurable[m]\ X)\ (hY\ :\ Measurable\ Y)}
\leanline{4}{(hi\ :\ Indep\ ((inferInstance\ :\ MeasurableSpace\ B).comap\ Y)\ m\ Q)}
\leanline{4}{(s\ :\ Set\ Ω)\ (hs\ :\ MeasurableSet[m]\ s)\ :}
\leanline{4}{(Q.restrict\ s).map\ (fun\ ω\ ↦\ (X\ ω,\ Y\ ω))\ =}
\leanline{6}{((Q.restrict\ s).map\ X).prod\ (Q.map\ Y)}
\end{leancode}
\leanlabel{How the Lean proof follows the argument}
\texttt{independent\_\allowbreak{}input\_\allowbreak{}restrict\_\allowbreak{}joint\_\allowbreak{}map} uses product-measure extensionality and the definition of independence on the two preimage events. The independent input retains its unrestricted marginal.
\end{leanbox}
\subsection{\texorpdfstring{The genuine conditional mean depends only on the current state}{The genuine conditional mean depends only on the current state}}
\label{report-step-WassersteinMixing-ConditionalIndependentAverage-lean-2}


\textbf{Claim.}

\[\mathbb E[F(X,Z)\mid\mathcal F]=\int F(X,z)\,\operatorname{Law}(Z)(dz).\]

\textbf{Proof.} Product integrability gives integrability of the averaged function and its measurable dependence on \(X\). On every past event, the restricted pair-law factorization and Fubini identify its integral with the integral of \(F(X,Z)\). These are the defining properties of conditional expectation. If \(X\) is only almost surely past-measurable, replace it by a measurable version, apply the result and transport back through almost-sure equality. The conditional mean is then measurable from \(X\) alone.

\begin{leanbox}{Lean: \texttt{condExp\_\allowbreak{}independent\_\allowbreak{}input\_\allowbreak{}average}, \texttt{condExp\_\allowbreak{}independent\_\allowbreak{}input\_\allowbreak{}state\_\allowbreak{}measurable}}
\leansource{\texttt{condExp\_\allowbreak{}independent\_\allowbreak{}input\_\allowbreak{}average} \quad WassersteinMixing/\allowbreak{}ConditionalIndependentAverage.lean, lines 39–80}
\begin{leancode}
\leanline{0}{theorem\ condExp\_independent\_input\_average\ \{Ω\ A\ B\ H\ :\ Type*\}}
\leanline{4}{(m\ :\ MeasurableSpace\ Ω)\ [m0\ :\ MeasurableSpace\ Ω]\ [MeasurableSpace\ A]\ [MeasurableSpace\ B]}
\leanline{4}{[NormedAddCommGroup\ H]\ [NormedSpace\ ℝ\ H]\ [CompleteSpace\ H]}
\leanline{4}{(Q\ :\ Measure\ Ω)\ [IsProbabilityMeasure\ Q]\ (hm\ :\ m\ ≤\ m0)}
\leanline{4}{(X\ :\ Ω\ →\ A)\ (Y\ :\ Ω\ →\ B)\ (hX\ :\ Measurable[m]\ X)\ (hY\ :\ Measurable\ Y)}
\leanline{4}{(hi\ :\ Indep\ ((inferInstance\ :\ MeasurableSpace\ B).comap\ Y)\ m\ Q)}
\leanline{4}{(F\ :\ A\ ×\ B\ →\ H)\ (hF\ :\ StronglyMeasurable\ F)}
\leanline{4}{(hFi\ :\ Integrable\ (fun\ ω\ ↦\ F\ (X\ ω,\ Y\ ω))\ Q)\ :}
\leanline{4}{Q[(fun\ ω\ ↦\ F\ (X\ ω,\ Y\ ω))\ |\ m]\ =ᵐ[Q]}
\leanline{6}{fun\ ω\ ↦\ ∫\ y,\ F\ (X\ ω,\ y)\ ∂Q.map\ Y}
\end{leancode}
\leansource{\texttt{condExp\_\allowbreak{}independent\_\allowbreak{}input\_\allowbreak{}state\_\allowbreak{}measurable} \quad WassersteinMixing/\allowbreak{}ConditionalIndependentAverage.lean, lines 85–112}
\begin{leancode}
\leanline{0}{theorem\ condExp\_independent\_input\_state\_measurable\ \{Ω\ A\ B\ H\ :\ Type*\}}
\leanline{4}{(m\ :\ MeasurableSpace\ Ω)\ [m0\ :\ MeasurableSpace\ Ω]}
\leanline{4}{[NormedAddCommGroup\ A]\ [mA\ :\ MeasurableSpace\ A]\ [BorelSpace\ A]}
\leanline{4}{[MeasurableSpace\ B]\ [NormedAddCommGroup\ H]\ [NormedSpace\ ℝ\ H]\ [CompleteSpace\ H]}
\leanline{4}{(Q\ :\ Measure\ Ω)\ [IsProbabilityMeasure\ Q]\ (hm\ :\ m\ ≤\ m0)}
\leanline{4}{(X\ :\ Ω\ →\ A)\ (Y\ :\ Ω\ →\ B)\ (hX\ :\ AEStronglyMeasurable[m]\ X\ Q)\ (hY\ :\ Measurable\ Y)}
\leanline{4}{(hi\ :\ Indep\ ((inferInstance\ :\ MeasurableSpace\ B).comap\ Y)\ m\ Q)}
\leanline{4}{(F\ :\ A\ ×\ B\ →\ H)\ (hF\ :\ StronglyMeasurable\ F)}
\leanline{4}{(hFi\ :\ Integrable\ (fun\ ω\ ↦\ F\ (X\ ω,\ Y\ ω))\ Q)\ :}
\leanline{4}{AEStronglyMeasurable[mA.comap\ X]\ (Q[(fun\ ω\ ↦\ F\ (X\ ω,\ Y\ ω))\ |\ m])\ Q}
\end{leancode}
\leanlabel{How the Lean proof follows the argument}
\texttt{cond\allowbreak{}Exp\_\allowbreak{}independent\_\allowbreak{}input\_\allowbreak{}average} applies the set-integral characterization of conditional expectation. The final theorem uses the measurable version furnished by \texttt{AEStrongly\allowbreak{}Measurable} and proves measurability for the sigma-algebra generated by the original state.
\end{leanbox}
\section{\texorpdfstring{Strong convergence of bounded observables without stationary approximants}{Strong convergence of bounded observables without stationary approximants}}
\label{proof-module-WassersteinMixing-BoundedComposition-lean}


Euler approximation states need not have the target law. A bounded continuous observable nevertheless gives a genuine square-integrable random variable at each state, and convergence in measure of those states implies strong convergence of the bounded observable compositions.

\subsection{\texorpdfstring{Actual composition, subsequences and dominated convergence}{Actual composition, subsequences and dominated convergence}}
\label{report-step-WassersteinMixing-BoundedComposition-lean-1}


\textbf{Claim.} For a bounded continuous \(\phi\), if \(X_i\to X\) in measure on a probability space, then \(\phi(X_i)\to\phi(X)\) in \(L^2\), even when their marginal laws differ.



\textbf{Proof.} Boundedness gives square integrability and an exact representative for every composition. Every sequence approaching the filter has a subsequence along which state convergence holds almost surely. Continuity gives almost-sure convergence of the compositions. Their squared difference is bounded by the integrable constant \(4\|\phi\|_\infty^2\), so dominated convergence forces its integral to zero along that subsequence. The subsequence criterion gives the full strong limit. Strong state convergence is a special case because it implies convergence in measure.

\begin{leanbox}{Lean: \texttt{boundedComposition}, \texttt{boundedComposition\_\allowbreak{}coe}, \texttt{boundedComposition\_\allowbreak{}sub\_\allowbreak{}norm\_\allowbreak{}sq} and 2 more}
\leansource{\texttt{boundedComposition} \quad WassersteinMixing/\allowbreak{}BoundedComposition.lean, lines 23–25}
\begin{leancode}
\leanline{0}{def\ boundedComposition\ (X\ :\ Lp\ E\ 2\ Q)\ (φ\ :\ E\ →ᵇ\ H)\ :\ Lp\ H\ 2\ Q\ :=}
\leanline{2}{(MemLp.of\_bound\ (φ.continuous.comp\_aestronglyMeasurable\ (Lp.aestronglyMeasurable\ X))}
\leanline{4}{‖φ‖\ (Eventually.of\_forall\ fun\ ω\ ↦\ φ.norm\_coe\_le\_norm\ (X\ ω))).toLp\ (fun\ ω\ ↦\ φ\ (X\ ω))}
\end{leancode}
\leansource{\texttt{boundedComposition\_\allowbreak{}coe} \quad WassersteinMixing/\allowbreak{}BoundedComposition.lean, lines 27–31}
\begin{leancode}
\leanline{0}{theorem\ boundedComposition\_coe\ (X\ :\ Lp\ E\ 2\ Q)\ (φ\ :\ E\ →ᵇ\ H)\ :}
\leanline{4}{boundedComposition\ Q\ X\ φ\ =ᵐ[Q]\ fun\ ω\ ↦\ φ\ (X\ ω)}
\end{leancode}
\leansource{\texttt{boundedComposition\_\allowbreak{}sub\_\allowbreak{}norm\_\allowbreak{}sq} \quad WassersteinMixing/\allowbreak{}BoundedComposition.lean, lines 32–41}
\begin{leancode}
\leanline{0}{private\ theorem\ boundedComposition\_sub\_norm\_sq\ (X\ Y\ :\ Lp\ E\ 2\ Q)\ (φ\ :\ E\ →ᵇ\ H)\ :}
\leanline{4}{‖boundedComposition\ Q\ X\ φ\ -\ boundedComposition\ Q\ Y\ φ‖\ \textasciicircum{}\ 2\ =}
\leanline{6}{∫\ ω,\ ‖φ\ (X\ ω)\ -\ φ\ (Y\ ω)‖\ \textasciicircum{}\ 2\ ∂Q}
\end{leancode}
\leansource{\texttt{boundedComposition\_\allowbreak{}tendsto\_\allowbreak{}inMeasure} \quad WassersteinMixing/\allowbreak{}BoundedComposition.lean, lines 44–79}
\begin{leancode}
\leanline{0}{theorem\ boundedComposition\_tendsto\_inMeasure\ \{I\ :\ Type*\}\ \{l\ :\ Filter\ I\}}
\leanline{4}{[l.IsCountablyGenerated]\ (X\ :\ I\ →\ Lp\ E\ 2\ Q)\ (Y\ :\ Lp\ E\ 2\ Q)}
\leanline{4}{(hlim\ :\ TendstoInMeasure\ Q\ (fun\ i\ ω\ ↦\ X\ i\ ω)\ l\ Y)\ (φ\ :\ E\ →ᵇ\ H)\ :}
\leanline{4}{Tendsto\ (fun\ i\ ↦\ boundedComposition\ Q\ (X\ i)\ φ)\ l\ ({\leanfallback 𝓝}\ (boundedComposition\ Q\ Y\ φ))}
\end{leancode}
\leansource{\texttt{boundedComposition\_\allowbreak{}tendsto} \quad WassersteinMixing/\allowbreak{}BoundedComposition.lean, lines 81–84}
\begin{leancode}
\leanline{0}{theorem\ boundedComposition\_tendsto\ \{I\ :\ Type*\}\ \{l\ :\ Filter\ I\}\ [l.IsCountablyGenerated]}
\leanline{4}{(X\ :\ I\ →\ Lp\ E\ 2\ Q)\ (Y\ :\ Lp\ E\ 2\ Q)\ (hlim\ :\ Tendsto\ X\ l\ ({\leanfallback 𝓝}\ Y))\ (φ\ :\ E\ →ᵇ\ H)\ :}
\leanline{4}{Tendsto\ (fun\ i\ ↦\ boundedComposition\ Q\ (X\ i)\ φ)\ l\ ({\leanfallback 𝓝}\ (boundedComposition\ Q\ Y\ φ))}
\end{leancode}
\leanlabel{How the Lean proof follows the argument}
\texttt{bounded\allowbreak{}Composition} constructs the actual \(L^2\) class from its uniform bound. The private norm identity uses its exact representative. \texttt{bounded\allowbreak{}Composition\_\allowbreak{}tendsto\_\allowbreak{}in\allowbreak{}Measure} applies subsequence extraction and dominated convergence; the final theorem converts strong state convergence to convergence in measure.
\end{leanbox}
\section{\texorpdfstring{Closedness of the conditional state-measurability property}{Closedness of the conditional state-measurability property}}
\label{proof-module-WassersteinMixing-ConditionalMarkovClosure-lean}


Conditional expectation of a measurable function of a current state and an independent future input depends only on the state. Expressed inside the true \(L^2\) space, this property is closed under strong limits and can therefore pass from actual Euler approximants to the diffusion.

\subsection{\texorpdfstring{Independent-input averaging and strong-limit closure}{Independent-input averaging and strong-limit closure}}
\label{report-step-WassersteinMixing-ConditionalMarkovClosure-lean-1}


\textbf{Claim.} If \(Z=F(X,W)\) is square integrable, \(X\) is past-measurable and \(W\) is independent of the whole past, then \(\mathbb E[Z\mid\mathcal F]\) is measurable from \(X\) alone. If \(Z_i\to Z\) strongly and all such conditional means are state-measurable, the limit conditional mean is state-measurable too.



\textbf{Proof.} Square integrability implies integrability on the probability space. The independent-input averaging theorem identifies ordinary conditional expectation as an average over the independent input; the actual \(L^2\) projection has the same almost-everywhere representative. For the limit, conditional expectation is a bounded linear projection, hence continuous. The square-integrable random variables measurable from the state form a closed subspace, so the strong limit remains in it.

\begin{leanbox}{Lean: \texttt{condExpL2\_\allowbreak{}independent\_\allowbreak{}input\_\allowbreak{}state\_\allowbreak{}measurable}, \texttt{condExpL2\_\allowbreak{}state\_\allowbreak{}measurable\_\allowbreak{}of\_\allowbreak{}tendsto}}
\leansource{\texttt{condExpL2\_\allowbreak{}independent\_\allowbreak{}input\_\allowbreak{}state\_\allowbreak{}measurable} \quad WassersteinMixing/\allowbreak{}ConditionalMarkovClosure.lean, lines 26–39}
\begin{leancode}
\leanline{0}{theorem\ condExpL2\_independent\_input\_state\_measurable}
\leanline{4}{(X\ :\ Ω\ →\ A)\ (Y\ :\ Ω\ →\ B)\ (hX\ :\ AEStronglyMeasurable[m]\ X\ Q)\ (hY\ :\ Measurable\ Y)}
\leanline{4}{(hi\ :\ Indep\ ((inferInstance\ :\ MeasurableSpace\ B).comap\ Y)\ m\ Q)}
\leanline{4}{(Z\ :\ Lp\ H\ 2\ Q)\ (F\ :\ A\ ×\ B\ →\ H)\ (hF\ :\ StronglyMeasurable\ F)}
\leanline{4}{(hZ\ :\ Z\ =ᵐ[Q]\ fun\ ω\ ↦\ F\ (X\ ω,\ Y\ ω))\ :}
\leanline{4}{AEStronglyMeasurable[mA.comap\ X]\ (condExpL2\ H\ ℝ\ hm\ Z\ :\ Lp\ H\ 2\ Q)\ Q}
\end{leancode}
\leansource{\texttt{condExpL2\_\allowbreak{}state\_\allowbreak{}measurable\_\allowbreak{}of\_\allowbreak{}tendsto} \quad WassersteinMixing/\allowbreak{}ConditionalMarkovClosure.lean, lines 43–54}
\begin{leancode}
\leanline{0}{theorem\ condExpL2\_state\_measurable\_of\_tendsto\ \{I\ :\ Type*\}\ \{l\ :\ Filter\ I\}\ [l.NeBot]}
\leanline{4}{(X\ :\ Ω\ →\ A)\ (hXa\ :\ Measurable\ X)}
\leanline{4}{(Z\ :\ I\ →\ Lp\ H\ 2\ Q)\ (W\ :\ Lp\ H\ 2\ Q)\ (hlim\ :\ Tendsto\ Z\ l\ ({\leanfallback 𝓝}\ W))}
\leanline{4}{(hZ\ :\ ∀\ i,\ AEStronglyMeasurable[mA.comap\ X]\ (condExpL2\ H\ ℝ\ hm\ (Z\ i)\ :\ Lp\ H\ 2\ Q)\ Q)\ :}
\leanline{4}{AEStronglyMeasurable[mA.comap\ X]\ (condExpL2\ H\ ℝ\ hm\ W\ :\ Lp\ H\ 2\ Q)\ Q}
\end{leancode}
\leanlabel{How the Lean proof follows the argument}
The first theorem identifies \texttt{cond\allowbreak{}Exp\allowbreak{}L2} with ordinary conditional expectation and applies the actual averaging theorem. The second uses continuity of that projection and the closed measurable subspace \texttt{is\allowbreak{}Closed\_\allowbreak{}aestrongly\allowbreak{}Measurable}.
\end{leanbox}
\section{\texorpdfstring{The actual future sigma-algebra and the natural past}{The actual future sigma-algebra and the natural past}}
\label{proof-module-WassersteinMixing-ActualMarkovPast-lean}


The actual \(L^2\) Brownian noise has an explicit coordinate representative. This identifies the natural past sigma-algebra with information contained in the initial state and raw Brownian past, allowing whole-path independence to pass to the actual conditioning filtration.

\subsection{\texorpdfstring{Representatives and future information}{Representatives and future information}}
\label{report-step-WassersteinMixing-ActualMarkovPast-lean-1}


\textbf{Construction.} For \(t\ge0\), the additive noise equals \(\sqrt2 B_t\) almost surely. Define the future path at \(a\) by \(W_a(t)=B_{a+t}-B_a\), and let its generated sigma-algebra be the future information. Each coordinate evaluation is measurable, hence the whole path and its generated information are measurable on the original probability space.

\begin{leanbox}{Lean: \texttt{actualBrownianNoise\_\allowbreak{}coe}, \texttt{actualBrownianFuture}, \texttt{actualBrownianFuture\_\allowbreak{}measurable} and 2 more}
\leansource{\texttt{actualBrownianNoise\_\allowbreak{}coe} \quad WassersteinMixing/\allowbreak{}ActualMarkovPast.lean, lines 21–30}
\begin{leancode}
\leanline{0}{theorem\ actualBrownianNoise\_coe\ (t\ :\ ℝ)\ (ht\ :\ 0\ ≤\ t)\ :}
\leanline{4}{actualBrownianNoise\ P\ V\ hV\ t\ =ᵐ[actualProbability\ P\ V]}
\leanline{6}{fun\ ω\ ↦\ Real.sqrt\ 2\ •\ ULA.brownianCoordinates\ (CanonicalSemigroup.brownian\ P)\ ⟨t,\ ht⟩\ ω}
\end{leancode}
\leansource{\texttt{actualBrownianFuture} \quad WassersteinMixing/\allowbreak{}ActualMarkovPast.lean, lines 33–34}
\begin{leancode}
\leanline{0}{def\ actualBrownianFuture\ (a\ :\ ℝ≥0)\ (ω\ :\ ActualSpace\ P)\ :\ Fin\ P.dimension\ →\ ℝ≥0\ →\ ℝ\ :=}
\leanline{2}{fun\ i\ t\ ↦\ CanonicalSemigroup.brownian\ P\ i\ (a\ +\ t)\ ω\ -\ CanonicalSemigroup.brownian\ P\ i\ a\ ω}
\end{leancode}
\leansource{\texttt{actualBrownianFuture\_\allowbreak{}measurable} \quad WassersteinMixing/\allowbreak{}ActualMarkovPast.lean, lines 36–39}
\begin{leancode}
\leanline{0}{theorem\ actualBrownianFuture\_measurable\ (a\ :\ ℝ≥0)\ :\ Measurable\ (actualBrownianFuture\ P\ a)}
\end{leancode}
\leansource{\texttt{actualFutureInformation} \quad WassersteinMixing/\allowbreak{}ActualMarkovPast.lean, lines 42–43}
\begin{leancode}
\leanline{0}{@[instance\_reducible]\ def\ actualFutureInformation\ (a\ :\ ℝ≥0)\ :\ MeasurableSpace\ (ActualSpace\ P)\ :=}
\leanline{2}{(inferInstance\ :\ MeasurableSpace\ (Fin\ P.dimension\ →\ ℝ≥0\ →\ ℝ)).comap\ (actualBrownianFuture\ P\ a)}
\end{leancode}
\leansource{\texttt{actualFutureInformation\_\allowbreak{}le} \quad WassersteinMixing/\allowbreak{}ActualMarkovPast.lean, lines 45–47}
\begin{leancode}
\leanline{0}{theorem\ actualFutureInformation\_le\ (a\ :\ ℝ≥0)\ :}
\leanline{4}{actualFutureInformation\ P\ a\ ≤\ (inferInstance\ :\ MeasurableSpace\ (ActualSpace\ P))}
\end{leancode}
\leanlabel{How the Lean proof follows the argument}
\texttt{actual\allowbreak{}Brownian\allowbreak{}Noise\_\allowbreak{}coe} identifies the actual noise representative with the coordinate Brownian vector. \texttt{actual\allowbreak{}Brownian\allowbreak{}Future} uses raw coordinate increments, and its measurability is proved coordinatewise. \texttt{actual\allowbreak{}Future\allowbreak{}Information} is the corresponding comap sigma-algebra.
\end{leanbox}


Earlier results used here: \cite[Section~3.24.15]{LangevinLean2026ULASC}; \cite[Section~3.24.21]{LangevinLean2026ULASC}; \cite[Section~3.24.23]{LangevinLean2026ULASC}; \cite[Section~3.24.8]{LangevinLean2026ULASC}; \cite[Section~3.25.1]{LangevinLean2026ULASC}; \cite[Section~3.8.11]{LangevinLean2026ULASC}.
\subsection{\texorpdfstring{Independence from the actual conditioning past}{Independence from the actual conditioning past}}
\label{report-step-WassersteinMixing-ActualMarkovPast-lean-2}


\textbf{Claim.} The future information at \(a\) is independent of \(\sigma(Y_0,\sqrt2 B_s:0\le s\le a)\), and therefore independent of the grid filtration at \(a=nh\).



\textbf{Proof.} Whole-path independence has already been proved for the raw source. Replace each additive-noise evaluation by its almost-sure coordinate representative and the initialization by its actual initial-coordinate representative. Every generator of the natural past is measurable from the raw past and initialization. Independence descends to this generated sub-sigma-algebra; the grid statement is its specialization.

\begin{leanbox}{Lean: \texttt{actualBrownian\_\allowbreak{}future\_\allowbreak{}independent\_\allowbreak{}pastAt}, \texttt{actualBrownian\_\allowbreak{}future\_\allowbreak{}independent\_\allowbreak{}past}}
\leansource{\texttt{actualBrownian\_\allowbreak{}future\_\allowbreak{}independent\_\allowbreak{}pastAt} \quad WassersteinMixing/\allowbreak{}ActualMarkovPast.lean, lines 51–97}
\begin{leancode}
\leanline{0}{theorem\ actualBrownian\_future\_independent\_pastAt\ (a\ :\ ℝ≥0)\ :}
\leanline{4}{Indep\ (actualFutureInformation\ P\ a)}
\leanline{6}{(ULA.lpLangevinPast\ (actualProbability\ P\ V)\ (CanonicalSemigroup.initial\ P\ V\ hV)}
\leanline{8}{(actualBrownianNoise\ P\ V\ hV)\ (a\ :\ ℝ))\ (actualProbability\ P\ V)}
\end{leancode}
\leansource{\texttt{actualBrownian\_\allowbreak{}future\_\allowbreak{}independent\_\allowbreak{}past} \quad WassersteinMixing/\allowbreak{}ActualMarkovPast.lean, lines 100–104}
\begin{leancode}
\leanline{0}{theorem\ actualBrownian\_future\_independent\_past\ (n\ :\ ℕ)\ :}
\leanline{4}{Indep\ (actualFutureInformation\ P\ ⟨(n\ :\ ℝ)\ *\ P.step,}
\leanline{6}{mul\_nonneg\ (Nat.cast\_nonneg\ n)\ P.step\_pos.le⟩)}
\leanline{6}{(actualPast\ P\ V\ hV\ n)\ (actualProbability\ P\ V)}
\end{leancode}
\leanlabel{How the Lean proof follows the argument}
The arbitrary-time theorem compares the generating comap sigma-algebras using their almost-everywhere representatives. The grid theorem substitutes the nonnegative time \texttt{n * P.\allowbreak{}step} and unfolds \texttt{actual\allowbreak{}Past}.
\end{leanbox}


Earlier results used here: \cite[Section~3.22.17]{LangevinLean2026ULASC}; \cite[Section~3.24.8]{LangevinLean2026ULASC}; \cite[Section~3.54.4]{LangevinLean2026ULASC}; \cite[Section~3.8.17]{LangevinLean2026ULASC}.
\section{\texorpdfstring{Stationary independent inputs for the actual future dynamics}{Stationary independent inputs for the actual future dynamics}}
\label{proof-module-WassersteinMixing-ActualMarkovInput-lean}


At each deterministic time, the actual state has the Gibbs law and is measurable from the past. Its pair with the whole independent future Brownian path therefore has a product law, and this pair law is independent of the cut time.

\subsection{\texorpdfstring{The state and future form an actual product input}{The state and future form an actual product input}}
\label{report-step-WassersteinMixing-ActualMarkovInput-lean-1}


\textbf{Claim.} At \(a\ge0\), the state \(Y_a\) is past-measurable, is independent of \(W_a\), and

\[\operatorname{Law}(Y_a,W_a)=\pi\otimes\operatorname{Law}(W_a).\]

\textbf{Proof.} Apply natural-past adaptation to a finite horizon ending at \(a\), then identify the endpoint with the global solution. Independence of the whole future from the past implies independence from this measurable state. The pair-law factorization follows from independence and the already proved Gibbs stationarity.

\begin{leanbox}{Lean: \texttt{actualState\_\allowbreak{}adapted\_\allowbreak{}pastAt}, \texttt{actualState\_\allowbreak{}future\_\allowbreak{}independent}, \texttt{actualState\_\allowbreak{}future\_\allowbreak{}joint\_\allowbreak{}law}}
\leansource{\texttt{actualState\_\allowbreak{}adapted\_\allowbreak{}pastAt} \quad WassersteinMixing/\allowbreak{}ActualMarkovInput.lean, lines 19–29}
\begin{leancode}
\leanline{0}{theorem\ actualState\_adapted\_pastAt\ (a\ :\ ℝ≥0)\ :}
\leanline{4}{AEStronglyMeasurable[ULA.lpLangevinPast\ (actualProbability\ P\ V)}
\leanline{6}{(CanonicalSemigroup.initial\ P\ V\ hV)\ (actualBrownianNoise\ P\ V\ hV)\ (a\ :\ ℝ)]}
\leanline{6}{(CanonicalSemigroup.state\ P\ V\ hV\ a)\ (actualProbability\ P\ V)}
\end{leancode}
\leansource{\texttt{actualState\_\allowbreak{}future\_\allowbreak{}independent} \quad WassersteinMixing/\allowbreak{}ActualMarkovInput.lean, lines 32–41}
\begin{leancode}
\leanline{0}{theorem\ actualState\_future\_independent\ (a\ :\ ℝ≥0)\ :}
\leanline{4}{IndepFun\ (CanonicalSemigroup.state\ P\ V\ hV\ a)\ (actualBrownianFuture\ P\ a)\ (actualProbability\ P\ V)}
\end{leancode}
\leansource{\texttt{actualState\_\allowbreak{}future\_\allowbreak{}joint\_\allowbreak{}law} \quad WassersteinMixing/\allowbreak{}ActualMarkovInput.lean, lines 44–51}
\begin{leancode}
\leanline{0}{theorem\ actualState\_future\_joint\_law\ (a\ :\ ℝ≥0)\ :}
\leanline{4}{(actualProbability\ P\ V).map\ (fun\ ω\ ↦}
\leanline{6}{(CanonicalSemigroup.state\ P\ V\ hV\ a\ ω,\ actualBrownianFuture\ P\ a\ ω))\ =}
\leanline{6}{(ULA.gibbsMeasure\ V).prod\ ((actualProbability\ P\ V).map\ (actualBrownianFuture\ P\ a))}
\end{leancode}
\leanlabel{How the Lean proof follows the argument}
\texttt{actual\allowbreak{}State\_\allowbreak{}adapted\_\allowbreak{}past\allowbreak{}At} applies the actual integral-solution adaptation theorem. \texttt{actual\allowbreak{}State\_\allowbreak{}future\_\allowbreak{}independent} descends future/\allowbreak{}past independence to the state. \texttt{actual\allowbreak{}State\_\allowbreak{}future\_\allowbreak{}joint\_\allowbreak{}law} uses the exact product pushforward identity.
\end{leanbox}


Earlier results used here: \cite[Section~3.25.2]{LangevinLean2026ULASC}; \cite[Section~3.25.3]{LangevinLean2026ULASC}; \cite[Section~3.26.3]{LangevinLean2026ULASC}; \cite[Section~3.54.4]{LangevinLean2026ULASC}; \cite[Section~3.54.8]{LangevinLean2026ULASC}.
\subsection{\texorpdfstring{Equality of whole future-input laws}{Equality of whole future-input laws}}
\label{report-step-WassersteinMixing-ActualMarkovInput-lean-2}


\textbf{Claim.} The shifted Brownian path has the same law at any two cut times, so \(\operatorname{Law}(Y_a,W_a)=\operatorname{Law}(Y_b,W_b)\).



\textbf{Proof.} Finite collections of shifted Brownian coordinates have the same centered Gaussian law, with covariance determined solely by their elapsed times. Equality on finite coordinate cylinders gives equality of the full path laws. The actual vector source is a product of independent coordinate sources; equality therefore transfers to the vector path. Substitute this equality into the preceding product-input factorization.

\begin{leanbox}{Lean: \texttt{preBrownian\_\allowbreak{}shift\_\allowbreak{}path\_\allowbreak{}law}, \texttt{actualBrownianFuture\_\allowbreak{}law\_\allowbreak{}eq}, \texttt{actualState\_\allowbreak{}future\_\allowbreak{}joint\_\allowbreak{}law\_\allowbreak{}eq}}
\leansource{\texttt{preBrownian\_\allowbreak{}shift\_\allowbreak{}path\_\allowbreak{}law} \quad WassersteinMixing/\allowbreak{}ActualMarkovInput.lean, lines 53–63}
\begin{leancode}
\leanline{0}{private\ theorem\ preBrownian\_shift\_path\_law\ \{Ω\ :\ Type*\}\ [MeasurableSpace\ Ω]}
\leanline{4}{(Q\ :\ Measure\ Ω)\ [IsProbabilityMeasure\ Q]\ (B\ :\ ℝ≥0\ →\ Ω\ →\ ℝ)}
\leanline{4}{(hB\ :\ IsPreBrownianReal\ B\ Q)\ (hBmeas\ :\ ∀\ t,\ Measurable\ (B\ t))\ (a\ b\ :\ ℝ≥0)\ :}
\leanline{4}{Q.map\ (fun\ ω\ t\ ↦\ B\ (a\ +\ t)\ ω\ -\ B\ a\ ω)\ =\ Q.map\ (fun\ ω\ t\ ↦\ B\ (b\ +\ t)\ ω\ -\ B\ b\ ω)}
\end{leancode}
\leansource{\texttt{actualBrownianFuture\_\allowbreak{}law\_\allowbreak{}eq} \quad WassersteinMixing/\allowbreak{}ActualMarkovInput.lean, lines 66–98}
\begin{leancode}
\leanline{0}{theorem\ actualBrownianFuture\_law\_eq\ (a\ b\ :\ ℝ≥0)\ :}
\leanline{4}{(actualProbability\ P\ V).map\ (actualBrownianFuture\ P\ a)\ =}
\leanline{6}{(actualProbability\ P\ V).map\ (actualBrownianFuture\ P\ b)}
\end{leancode}
\leansource{\texttt{actualState\_\allowbreak{}future\_\allowbreak{}joint\_\allowbreak{}law\_\allowbreak{}eq} \quad WassersteinMixing/\allowbreak{}ActualMarkovInput.lean, lines 101–107}
\begin{leancode}
\leanline{0}{theorem\ actualState\_future\_joint\_law\_eq\ (a\ b\ :\ ℝ≥0)\ :}
\leanline{4}{(actualProbability\ P\ V).map\ (fun\ ω\ ↦}
\leanline{6}{(CanonicalSemigroup.state\ P\ V\ hV\ a\ ω,\ actualBrownianFuture\ P\ a\ ω))\ =}
\leanline{4}{(actualProbability\ P\ V).map\ (fun\ ω\ ↦}
\leanline{6}{(CanonicalSemigroup.state\ P\ V\ hV\ b\ ω,\ actualBrownianFuture\ P\ b\ ω))}
\end{leancode}
\leanlabel{How the Lean proof follows the argument}
\texttt{pre\allowbreak{}Brownian\_\allowbreak{}shift\_\allowbreak{}path\_\allowbreak{}law} proves equality of path measures from finite-dimensional Gaussian data. \texttt{actual\allowbreak{}Brownian\allowbreak{}Future\_\allowbreak{}law\_\allowbreak{}eq} applies it to the actual coordinate product. The final theorem rewrites both exact product input laws.
\end{leanbox}


Earlier results used here: \cite[Section~3.22.1]{LangevinLean2026ULASC}; \cite[Section~3.8.11]{LangevinLean2026ULASC}; \cite[Section~3.8.16]{LangevinLean2026ULASC}; \cite[Section~3.8.17]{LangevinLean2026ULASC}; \cite[Section~3.8.19]{LangevinLean2026ULASC}; \cite[Section~3.8.20]{LangevinLean2026ULASC}; \cite[Section~3.8.21]{LangevinLean2026ULASC}.
\section{\texorpdfstring{Euler convergence with a moving step and endpoint}{Euler convergence with a moving step and endpoint}}
\label{proof-module-WassersteinMixing-SemigroupEuler-lean}


For a Lipschitz drift and continuous additive input, Euler approximations converge to the actual integral solution when the step tends to zero and the elapsed grid time tends to the desired time. The step and iteration count may vary independently along any filter subject to a fixed finite horizon bound.



Earlier results used here: \cite[Section~3.16]{LangevinLean2026ULASC}.
\subsection{\texorpdfstring{Uniform finite-horizon control gives the moving-time limit}{Uniform finite-horizon control gives the moving-time limit}}
\label{report-step-WassersteinMixing-SemigroupEuler-lean-1}


\textbf{Claim.} If \(r_i\to0\), \(n_ir_i\le T\) and \(n_ir_i\to t\), then the actual additive Euler endpoint converges to \(Y_t\).



\textbf{Proof.} The deterministic finite-horizon Euler estimate bounds the endpoint error relative to \(Y_{n_ir_i}\) by a constant times the modulus of continuity of the drift along the solution over intervals of length \(r_i\). Uniform continuity on the compact horizon makes that modulus tend to zero. Continuity of \(Y\) makes \(Y_{n_ir_i}\to Y_t\). The triangle inequality combines these limits. The case of a zero Lipschitz constant is included in the same finite-horizon estimate.

\begin{leanbox}{Lean: \texttt{additiveEulerChain\_\allowbreak{}variable\_\allowbreak{}tendsto}}
\leansource{\texttt{additiveEulerChain\_\allowbreak{}variable\_\allowbreak{}tendsto} \quad WassersteinMixing/\allowbreak{}SemigroupEuler.lean, lines 22–70}
\begin{leancode}
\leanline{0}{theorem\ additiveEulerChain\_variable\_tendsto}
\leanline{4}{(g\ :\ H\ →\ H)\ (L\ :\ ℝ≥0)\ (hLip\ :\ LipschitzWith\ L\ g)}
\leanline{4}{(x\ :\ H)\ (B\ Y\ :\ ℝ\ →\ H)\ (hB0\ :\ B\ 0\ =\ 0)\ (hY\ :\ Continuous\ Y)}
\leanline{4}{(T\ :\ ℝ)\ (hT\ :\ 0\ ≤\ T)}
\leanline{4}{(hYe\ :\ ∀\ t\ ∈\ Icc\ 0\ T,\ Y\ t\ -\ x\ =\ -(∫\ s\ in\ 0..t,\ g\ (Y\ s))\ +\ B\ t)}
\leanline{4}{(step\ :\ I\ →\ ℝ)\ (count\ :\ I\ →\ ℕ)\ (hstep\ :\ ∀\ i,\ 0\ ≤\ step\ i)}
\leanline{4}{(hmesh\ :\ Tendsto\ step\ l\ ({\leanfallback 𝓝}\ 0))\ (hbound\ :\ ∀\ i,\ (count\ i\ :\ ℝ)\ *\ step\ i\ ≤\ T)}
\leanline{4}{(t\ :\ ℝ)\ (htime\ :\ Tendsto\ (fun\ i\ ↦\ (count\ i\ :\ ℝ)\ *\ step\ i)\ l\ ({\leanfallback 𝓝}\ t))\ :}
\leanline{4}{Tendsto\ (fun\ i\ ↦\ ULA.additiveEulerChain\ g\ (step\ i)\ x\ B\ (count\ i))\ l\ ({\leanfallback 𝓝}\ (Y\ t))}
\end{leancode}
\leanlabel{How the Lean proof follows the argument}
\texttt{additive\allowbreak{}Euler\allowbreak{}Chain\_\allowbreak{}variable\_\allowbreak{}tendsto} applies the proved uniform additive Euler error estimate and compact-horizon uniform continuity. It separates the Euler error at the actual elapsed grid time from the moving-endpoint continuity error, then combines their limits.
\end{leanbox}


Earlier results used here: \cite[Section~3.16.1]{LangevinLean2026ULASC}; \cite[Section~3.16.7]{LangevinLean2026ULASC}.
\section{\texorpdfstring{Actual shifted Langevin Euler convergence}{Actual shifted Langevin Euler convergence}}
\label{proof-module-WassersteinMixing-SemigroupEulerActual-lean}


Subtracting the actual integral equation at two times yields the shifted equation started from the actual current state. Applying deterministic additive Euler convergence to this genuine \(L^2\) equation gives both fixed-terminal and moving-grid limits.

\subsection{\texorpdfstring{Restarting the actual integral equation}{Restarting the actual integral equation}}
\label{report-step-WassersteinMixing-SemigroupEulerActual-lean-1}


\textbf{Claim.} For \(a,t\ge0\),

\[Y_{a+t}-Y_a=-\int_0^t\nabla V(Y_{a+s})\,ds+\sqrt2(B_{a+t}-B_a).\]

Euler recursion started from \(Y_a\) with these actual future increments converges in \(L^2\) to \(Y_{a+t}\), including along variable steps and counts whose elapsed time tends to \(t\).



\textbf{Proof.} Subtract the global integral equations at \(a+t\) and \(a\). Additivity and translation of the integrals give the shifted drift integral, while the exact Brownian increment theorem gives the shifted noise. The actual gradient map on \(L^2\) is Lipschitz and the global path is continuous, so the proved additive Euler convergence theorem applies.

\begin{leanbox}{Lean: \texttt{shifted\_\allowbreak{}state\_\allowbreak{}integral}, \texttt{shiftedEuler\_\allowbreak{}tendsto}, \texttt{shiftedEuler\_\allowbreak{}variable\_\allowbreak{}tendsto}}
\leansource{\texttt{shifted\_\allowbreak{}state\_\allowbreak{}integral} \quad WassersteinMixing/\allowbreak{}SemigroupEulerActual.lean, lines 20–33}
\begin{leancode}
\leanline{0}{theorem\ shifted\_state\_integral\ (P\ :\ Parameters)\ (V\ :\ E\ P.dimension\ →\ ℝ)}
\leanline{4}{(hV\ :\ PotentialAssumptions\ P\ V)\ (a\ T\ :\ ℝ)\ (ha\ :\ 0\ ≤\ a)\ (hT\ :\ 0\ ≤\ T)\ :}
\leanline{4}{letI\ :=\ actualProbability\_isProbability\ P\ V\ hV}
\leanline{4}{∀\ t\ ∈\ Icc\ 0\ T,\ state\ P\ V\ hV\ (a\ +\ t)\ -\ state\ P\ V\ hV\ a\ =}
\leanline{6}{-(∫\ s\ in\ 0..t,\ ULA.lpDrift\ (actualProbability\ P\ V)\ (gradient\ V)}
\leanline{8}{hV.gradientLipschitz\ (state\ P\ V\ hV\ (a\ +\ s)))\ +}
\leanline{6}{ULA.brownianLangevinNoise\ (actualProbability\ P\ V)\ (brownian\ P)}
\leanline{8}{(brownian\_preBrownian\ P\ V\ hV)\ a\ t}
\end{leancode}
\leansource{\texttt{shiftedEuler\_\allowbreak{}tendsto} \quad WassersteinMixing/\allowbreak{}SemigroupEulerActual.lean, lines 36–51}
\begin{leancode}
\leanline{0}{theorem\ shiftedEuler\_tendsto\ (P\ :\ Parameters)\ (V\ :\ E\ P.dimension\ →\ ℝ)}
\leanline{4}{(hV\ :\ PotentialAssumptions\ P\ V)\ (a\ t\ :\ ℝ)\ (ha\ :\ 0\ ≤\ a)\ (ht\ :\ 0\ ≤\ t)\ :}
\leanline{4}{letI\ :=\ actualProbability\_isProbability\ P\ V\ hV}
\leanline{4}{Tendsto\ (fun\ n\ :\ ℕ\ ↦\ ULA.additiveEulerChain}
\leanline{6}{(ULA.lpDrift\ (actualProbability\ P\ V)\ (gradient\ V)\ hV.gradientLipschitz)}
\leanline{6}{(t\ /\ ((n\ +\ 1\ :\ ℕ)\ :\ ℝ))\ (state\ P\ V\ hV\ a)}
\leanline{6}{(ULA.brownianLangevinNoise\ (actualProbability\ P\ V)\ (brownian\ P)}
\leanline{8}{(brownian\_preBrownian\ P\ V\ hV)\ a)\ (n\ +\ 1))\ atTop\ ({\leanfallback 𝓝}\ (state\ P\ V\ hV\ (a\ +\ t)))}
\end{leancode}
\leansource{\texttt{shiftedEuler\_\allowbreak{}variable\_\allowbreak{}tendsto} \quad WassersteinMixing/\allowbreak{}SemigroupEulerActual.lean, lines 55–76}
\begin{leancode}
\leanline{0}{theorem\ shiftedEuler\_variable\_tendsto\ \{I\ :\ Type*\}\ \{l\ :\ Filter\ I\}}
\leanline{4}{(P\ :\ Parameters)\ (V\ :\ E\ P.dimension\ →\ ℝ)\ (hV\ :\ PotentialAssumptions\ P\ V)}
\leanline{4}{(a\ T\ :\ ℝ)\ (ha\ :\ 0\ ≤\ a)\ (hT\ :\ 0\ ≤\ T)}
\leanline{4}{(step\ :\ I\ →\ ℝ)\ (count\ :\ I\ →\ ℕ)\ (hstep\ :\ ∀\ i,\ 0\ ≤\ step\ i)}
\leanline{4}{(hmesh\ :\ Tendsto\ step\ l\ ({\leanfallback 𝓝}\ 0))\ (hbound\ :\ ∀\ i,\ (count\ i\ :\ ℝ)\ *\ step\ i\ ≤\ T)}
\leanline{4}{(t\ :\ ℝ)\ (htime\ :\ Tendsto\ (fun\ i\ ↦\ (count\ i\ :\ ℝ)\ *\ step\ i)\ l\ ({\leanfallback 𝓝}\ t))\ :}
\leanline{4}{letI\ :=\ actualProbability\_isProbability\ P\ V\ hV}
\leanline{4}{Tendsto\ (fun\ i\ ↦\ ULA.additiveEulerChain}
\leanline{6}{(ULA.lpDrift\ (actualProbability\ P\ V)\ (gradient\ V)\ hV.gradientLipschitz)}
\leanline{6}{(step\ i)\ (state\ P\ V\ hV\ a)}
\leanline{6}{(ULA.brownianLangevinNoise\ (actualProbability\ P\ V)\ (brownian\ P)}
\leanline{8}{(brownian\_preBrownian\ P\ V\ hV)\ a)\ (count\ i))\ l\ ({\leanfallback 𝓝}\ (state\ P\ V\ hV\ (a\ +\ t)))}
\end{leancode}
\leanlabel{How the Lean proof follows the argument}
\texttt{shifted\_\allowbreak{}state\_\allowbreak{}integral} derives the actual restarted equation by interval-integral subtraction. \texttt{shifted\allowbreak{}Euler\_\allowbreak{}tendsto} uses equal elapsed subdivisions; \texttt{shifted\allowbreak{}Euler\_\allowbreak{}variable\_\allowbreak{}tendsto} applies the general moving-step theorem to the actual Lipschitz drift and shifted Brownian input.
\end{leanbox}


Earlier results used here: \cite[Section~3.13.2]{LangevinLean2026ULASC}; \cite[Section~3.13.5]{LangevinLean2026ULASC}; \cite[Section~3.16.1]{LangevinLean2026ULASC}; \cite[Section~3.16.9]{LangevinLean2026ULASC}; \cite[Section~3.25.1]{LangevinLean2026ULASC}; \cite[Section~3.25.3]{LangevinLean2026ULASC}; \cite[Section~3.44.2]{LangevinLean2026ULASC}.
\section{\texorpdfstring{A finite Euler endpoint as a measurable function of actual inputs}{A finite Euler endpoint as a measurable function of actual inputs}}
\label{proof-module-WassersteinMixing-SemigroupEulerMap-lean}


A finite Euler endpoint is a measurable function of its starting state and the whole future coordinate path. This exact pointwise description links the actual \(L^2\) recursion to whole-path independence and conditional averaging.

\subsection{\texorpdfstring{The actual pointwise recursion and its representative}{The actual pointwise recursion and its representative}}
\label{report-step-WassersteinMixing-SemigroupEulerMap-lean-1}


\textbf{Construction.} Starting from \(x\) and a future path \(w\), recursively add the true gradient Euler drift and consecutive increments of \(\sqrt2w\). The resulting endpoint map is measurable.



\textbf{Claim and proof.} The actual shifted Euler \(L^2\) variable equals this map applied to \((Y_a,W_a)\) almost surely. The actual Brownian increment has the explicit raw future-coordinate representative. Induction on the Euler recursion combines this identity with the almost-everywhere true gradient representative and the preceding endpoint representative. Thus only actual state and future-path evaluations occur.

\begin{leanbox}{Lean: \texttt{pointwiseEulerMap}, \texttt{pointwiseEulerMap\_\allowbreak{}measurable}, \texttt{shiftedNoise\_\allowbreak{}coe} and 1 more}
\leansource{\texttt{pointwiseEulerMap} \quad WassersteinMixing/\allowbreak{}SemigroupEulerMap.lean, lines 20–23}
\begin{leancode}
\leanline{0}{def\ pointwiseEulerMap\ (P\ :\ Parameters)\ (V\ :\ E\ P.dimension\ →\ ℝ)\ (h\ :\ ℝ)\ (n\ :\ ℕ)}
\leanline{4}{(p\ :\ E\ P.dimension\ ×\ (Fin\ P.dimension\ →\ ℝ≥0\ →\ ℝ))\ :\ E\ P.dimension\ :=}
\leanline{2}{ULA.additiveEulerChain\ (gradient\ V)\ h\ p.1}
\leanline{4}{(fun\ s\ ↦\ Real.sqrt\ 2\ •\ WithLp.toLp\ 2\ (fun\ i\ ↦\ p.2\ i\ s.toNNReal))\ n}
\end{leancode}
\leansource{\texttt{pointwiseEulerMap\_\allowbreak{}measurable} \quad WassersteinMixing/\allowbreak{}SemigroupEulerMap.lean, lines 26–43}
\begin{leancode}
\leanline{0}{theorem\ pointwiseEulerMap\_measurable\ (P\ :\ Parameters)\ (V\ :\ E\ P.dimension\ →\ ℝ)}
\leanline{4}{(hV\ :\ PotentialAssumptions\ P\ V)\ (h\ :\ ℝ)\ (n\ :\ ℕ)\ :}
\leanline{4}{Measurable\ (pointwiseEulerMap\ P\ V\ h\ n)}
\end{leancode}
\leansource{\texttt{shiftedNoise\_\allowbreak{}coe} \quad WassersteinMixing/\allowbreak{}SemigroupEulerMap.lean, lines 47–75}
\begin{leancode}
\leanline{0}{theorem\ shiftedNoise\_coe\ (P\ :\ Parameters)\ (V\ :\ E\ P.dimension\ →\ ℝ)}
\leanline{4}{(hV\ :\ PotentialAssumptions\ P\ V)\ (a\ :\ ℝ≥0)\ (s\ :\ ℝ)\ (hs\ :\ 0\ ≤\ s)\ :}
\leanline{4}{ULA.brownianLangevinNoise\ (actualProbability\ P\ V)\ (brownian\ P)}
\leanline{6}{(brownian\_preBrownian\ P\ V\ hV)\ a\ s\ =ᵐ[actualProbability\ P\ V]}
\leanline{6}{fun\ ω\ ↦\ Real.sqrt\ 2\ •\ WithLp.toLp\ 2\ (fun\ i\ ↦\ actualBrownianFuture\ P\ a\ ω\ i\ s.toNNReal)}
\end{leancode}
\leansource{\texttt{shiftedEuler\_\allowbreak{}coe} \quad WassersteinMixing/\allowbreak{}SemigroupEulerMap.lean, lines 79–127}
\begin{leancode}
\leanline{0}{theorem\ shiftedEuler\_coe\ (P\ :\ Parameters)\ (V\ :\ E\ P.dimension\ →\ ℝ)}
\leanline{4}{(hV\ :\ PotentialAssumptions\ P\ V)\ (a\ :\ ℝ≥0)\ (h\ :\ ℝ)\ (hh\ :\ 0\ ≤\ h)\ (n\ :\ ℕ)\ :}
\leanline{4}{letI\ :=\ actualProbability\_isProbability\ P\ V\ hV}
\leanline{4}{ULA.additiveEulerChain}
\leanline{6}{(ULA.lpDrift\ (actualProbability\ P\ V)\ (gradient\ V)\ hV.gradientLipschitz)\ h}
\leanline{6}{(state\ P\ V\ hV\ a)}
\leanline{6}{(ULA.brownianLangevinNoise\ (actualProbability\ P\ V)\ (brownian\ P)}
\leanline{8}{(brownian\_preBrownian\ P\ V\ hV)\ a)\ n\ =ᵐ[actualProbability\ P\ V]}
\leanline{6}{fun\ ω\ ↦\ pointwiseEulerMap\ P\ V\ h\ n\ (state\ P\ V\ hV\ a\ ω,\ actualBrownianFuture\ P\ a\ ω)}
\end{leancode}
\leanlabel{How the Lean proof follows the argument}
\texttt{pointwise\allowbreak{}Euler\allowbreak{}Map\_\allowbreak{}measurable} inducts through the true measurable proposal update. \texttt{shifted\allowbreak{}Noise\_\allowbreak{}coe} identifies the actual additive Brownian increment. \texttt{shifted\allowbreak{}Euler\_\allowbreak{}coe} proves equality of the actual recursion and measurable input map by almost-everywhere induction.
\end{leanbox}


Earlier results used here: \cite[Section~3.13.2]{LangevinLean2026ULASC}; \cite[Section~3.13.3]{LangevinLean2026ULASC}; \cite[Section~3.16.1]{LangevinLean2026ULASC}; \cite[Section~3.24.14]{LangevinLean2026ULASC}; \cite[Section~3.24.15]{LangevinLean2026ULASC}; \cite[Section~3.24.21]{LangevinLean2026ULASC}; \cite[Section~3.25.1]{LangevinLean2026ULASC}.
\section{\texorpdfstring{State-measurable random variables are actual evaluated observables}{State-measurable random variables are actual evaluated observables}}
\label{proof-module-WassersteinMixing-ObservableRange-lean}


Doob–Dynkin factorization identifies every square-integrable function of a state with an observable in the \(L^2\) space of its actual law. On this range, evaluation followed by its adjoint is the identity, identifying state-measurable conditional means with actual adjoint compression.

\subsection{\texorpdfstring{Actual square-integrable Doob–Dynkin factorization}{Actual square-integrable Doob–Dynkin factorization}}
\label{report-step-WassersteinMixing-ObservableRange-lean-1}


\textbf{Claim.} If \(X_*Q=\mu\) and \(Z\in L^2(Q)\) is measurable from \(X\) alone, then \(Z=f(X)\) for some actual \(f\in L^2(\mu)\).



\textbf{Proof.} Choose a state-measurable representative of \(Z\). Measurable factorization writes it as \(g\circ X\). Since this is almost surely the square-integrable original variable, its square norm is integrable. The exact pushforward identity transfers that integrability to \(g\) under \(\mu\). Its \(L^2(\mu)\) equivalence class evaluates back to \(Z\).

\begin{leanbox}{Lean: \texttt{observable\_\allowbreak{}range\_\allowbreak{}of\_\allowbreak{}state\_\allowbreak{}measurable}}
\leansource{\texttt{observable\_\allowbreak{}range\_\allowbreak{}of\_\allowbreak{}state\_\allowbreak{}measurable} \quad WassersteinMixing/\allowbreak{}ObservableRange.lean, lines 27–43}
\begin{leancode}
\leanline{0}{theorem\ observable\_range\_of\_state\_measurable\ (Z\ :\ Lp\ F\ 2\ Q)}
\leanline{4}{(hZ\ :\ AEStronglyMeasurable[mA.comap\ X]\ Z\ Q)\ :}
\leanline{4}{∃\ f\ :\ Lp\ F\ 2\ μ,\ Lp.compMeasurePreservingₗᵢ\ ℝ\ X\ hX\ f\ =\ Z}
\end{leancode}
\leanlabel{How the Lean proof follows the argument}
\texttt{observable\_\allowbreak{}range\_\allowbreak{}of\_\allowbreak{}state\_\allowbreak{}measurable} uses measurable composition factorization on the chosen strongly measurable representative. It transfers \texttt{Mem\allowbreak{}Lp} through the actual measure-preserving state map before constructing the observable class and comparing pullback representatives.
\end{leanbox}
\subsection{\texorpdfstring{Identifying the true conditional mean with compression}{Identifying the true conditional mean with compression}}
\label{report-step-WassersteinMixing-ObservableRange-lean-2}


\textbf{Claim.} For an evaluation isometry \(J\), \(JJ^*Z=Z\) on its range. If all evaluations are past-measurable and \(\mathbb E[Y\mid\mathcal F]\) belongs to that range, then

\[\mathbb E[Y\mid\mathcal F]=J(J^*Y).\]

\textbf{Proof.} The isometry identity \(J^*J=I\) proves the first claim. Pair the conditional expectation against arbitrary evaluated observables. Adaptation lets the conditional projection be removed from that pairing, yielding \(J^*\mathbb E[Y\mid\mathcal F]=J^*Y\). Apply \(J\) and use the first claim for the known range member.

\begin{leanbox}{Lean: \texttt{evaluation\_\allowbreak{}projection\_\allowbreak{}eq\_\allowbreak{}of\_\allowbreak{}mem\_\allowbreak{}range}, \texttt{conditional\_\allowbreak{}eq\_\allowbreak{}evaluation\_\allowbreak{}adjoint}}
\leansource{\texttt{evaluation\_\allowbreak{}projection\_\allowbreak{}eq\_\allowbreak{}of\_\allowbreak{}mem\_\allowbreak{}range} \quad WassersteinMixing/\allowbreak{}ObservableRange.lean, lines 54–59}
\begin{leancode}
\leanline{0}{theorem\ evaluation\_projection\_eq\_of\_mem\_range\ (Z\ :\ Lp\ F\ 2\ Q)\ (hZ\ :\ Z\ ∈\ Set.range\ J)\ :}
\leanline{4}{J\ (J.toContinuousLinearMap.adjoint\ Z)\ =\ Z}
\end{leancode}
\leansource{\texttt{conditional\_\allowbreak{}eq\_\allowbreak{}evaluation\_\allowbreak{}adjoint} \quad WassersteinMixing/\allowbreak{}ObservableRange.lean, lines 63–77}
\begin{leancode}
\leanline{0}{theorem\ conditional\_eq\_evaluation\_adjoint\ \{m\ :\ MeasurableSpace\ Ω\}\ (hm\ :\ m\ ≤\ mΩ)}
\leanline{4}{(Y\ :\ Lp\ F\ 2\ Q)}
\leanline{4}{(hJ\ :\ ∀\ f,\ AEStronglyMeasurable[m]\ (J\ f)\ Q)}
\leanline{4}{(hZ\ :\ (condExpL2\ F\ ℝ\ hm\ Y\ :\ Lp\ F\ 2\ Q)\ ∈\ Set.range\ J)\ :}
\leanline{4}{(condExpL2\ F\ ℝ\ hm\ Y\ :\ Lp\ F\ 2\ Q)\ =\ J\ (J.toContinuousLinearMap.adjoint\ Y)}
\end{leancode}
\leanlabel{How the Lean proof follows the argument}
\texttt{evaluation\_\allowbreak{}projection\_\allowbreak{}eq\_\allowbreak{}of\_\allowbreak{}mem\_\allowbreak{}range} cancels the isometry adjoint on an explicit range witness. \texttt{conditional\_\allowbreak{}eq\_\allowbreak{}evaluation\_\allowbreak{}adjoint} proves equality of extracted adjoint vectors by testing against every observable and using the conditional-expectation pairing identity.
\end{leanbox}
\section{\texorpdfstring{Actual pair laws determine the observable compression}{Actual pair laws determine the observable compression}}
\label{proof-module-WassersteinMixing-SemigroupPair-lean}


The compression of stationary observations depends only on their true two-time joint law. In particular equal pair laws give equal operators, and reversal symmetry of the pair proves self-adjointness.

\subsection{\texorpdfstring{From actual pair laws to operator equality}{From actual pair laws to operator equality}}
\label{report-step-WassersteinMixing-SemigroupPair-lean-1}


\textbf{Claim.} The evaluation pairing is the integral of the actual pointwise observable pairing under the joint state law. Thus equality of pair laws identifies the compressed operators. If that law is unchanged by coordinate swap, the operator is self-adjoint.



\textbf{Proof.} Compare the exact evaluation representatives, then change variables under the measurable pair map. Equal pair laws give equal integrals against every pair of square-integrable observables. The adjoint correlation identity and separation by inner products yield operator equality. Swapping the pair exchanges the two observable arguments and gives the symmetric pairing identity required for self-adjointness.

\begin{leanbox}{Lean: \texttt{evaluation\_\allowbreak{}inner\_\allowbreak{}pair}, \texttt{compression\_\allowbreak{}eq\_\allowbreak{}of\_\allowbreak{}pair\_\allowbreak{}law}, \texttt{compression\_\allowbreak{}isSelfAdjoint\_\allowbreak{}of\_\allowbreak{}pair\_\allowbreak{}swap}}
\leansource{\texttt{evaluation\_\allowbreak{}inner\_\allowbreak{}pair} \quad WassersteinMixing/\allowbreak{}SemigroupPair.lean, lines 19–29}
\begin{leancode}
\leanline{0}{theorem\ evaluation\_inner\_pair\ (X\ Y\ :\ Lp\ E\ 2\ Q)\ (hX\ :\ Q.map\ X\ =\ μ)\ (hY\ :\ Q.map\ Y\ =\ μ)}
\leanline{4}{(f\ g\ :\ Lp\ F\ 2\ μ)\ :}
\leanline{4}{inner\ {\leanfallback 𝕜}\ (evaluation\ ({\leanfallback 𝕜}\ :=\ {\leanfallback 𝕜})\ μ\ Q\ Y\ hY\ f)\ (evaluation\ ({\leanfallback 𝕜}\ :=\ {\leanfallback 𝕜})\ μ\ Q\ X\ hX\ g)\ =}
\leanline{6}{∫\ p,\ inner\ {\leanfallback 𝕜}\ (f\ p.2)\ (g\ p.1)\ ∂Q.map\ (fun\ ω\ ↦\ (X\ ω,\ Y\ ω))}
\end{leancode}
\leansource{\texttt{compression\_\allowbreak{}eq\_\allowbreak{}of\_\allowbreak{}pair\_\allowbreak{}law} \quad WassersteinMixing/\allowbreak{}SemigroupPair.lean, lines 34–47}
\begin{leancode}
\leanline{0}{theorem\ compression\_eq\_of\_pair\_law\ \{Ω'\ :\ Type*\}\ [MeasurableSpace\ Ω']\ (Q'\ :\ Measure\ Ω')}
\leanline{4}{(X\ Y\ :\ Lp\ E\ 2\ Q)\ (X'\ Y'\ :\ Lp\ E\ 2\ Q')}
\leanline{4}{(hX\ :\ Q.map\ X\ =\ μ)\ (hY\ :\ Q.map\ Y\ =\ μ)}
\leanline{4}{(hX'\ :\ Q'.map\ X'\ =\ μ)\ (hY'\ :\ Q'.map\ Y'\ =\ μ)}
\leanline{4}{(hpair\ :\ Q.map\ (fun\ ω\ ↦\ (X\ ω,\ Y\ ω))\ =\ Q'.map\ (fun\ ω\ ↦\ (X'\ ω,\ Y'\ ω)))\ :}
\leanline{4}{ObservableTransition.compression\ (evaluation\ ({\leanfallback 𝕜}\ :=\ {\leanfallback 𝕜})\ (F\ :=\ F)\ μ\ Q\ X\ hX)}
\leanline{6}{(evaluation\ μ\ Q\ Y\ hY)\ =}
\leanline{4}{ObservableTransition.compression\ (evaluation\ μ\ Q'\ X'\ hX')\ (evaluation\ μ\ Q'\ Y'\ hY')}
\end{leancode}
\leansource{\texttt{compression\_\allowbreak{}isSelfAdjoint\_\allowbreak{}of\_\allowbreak{}pair\_\allowbreak{}swap} \quad WassersteinMixing/\allowbreak{}SemigroupPair.lean, lines 52–73}
\begin{leancode}
\leanline{0}{theorem\ compression\_isSelfAdjoint\_of\_pair\_swap\ (X\ Y\ :\ Lp\ E\ 2\ Q)}
\leanline{4}{(hX\ :\ Q.map\ X\ =\ μ)\ (hY\ :\ Q.map\ Y\ =\ μ)}
\leanline{4}{(hswap\ :\ (Q.map\ (fun\ ω\ ↦\ (X\ ω,\ Y\ ω))).map\ Prod.swap\ =}
\leanline{6}{Q.map\ (fun\ ω\ ↦\ (X\ ω,\ Y\ ω)))\ :}
\leanline{4}{IsSelfAdjoint\ (ObservableTransition.compression}
\leanline{6}{(evaluation\ ({\leanfallback 𝕜}\ :=\ {\leanfallback 𝕜})\ (F\ :=\ F)\ μ\ Q\ X\ hX)\ (evaluation\ μ\ Q\ Y\ hY))}
\end{leancode}
\leanlabel{How the Lean proof follows the argument}
\texttt{evaluation\_\allowbreak{}inner\_\allowbreak{}pair} proves the actual pushforward integral identity. \texttt{compression\_\allowbreak{}eq\_\allowbreak{}of\_\allowbreak{}pair\_\allowbreak{}law} uses inner-product extensionality. \texttt{compression\_\allowbreak{}is\allowbreak{}Self\allowbreak{}Adjoint\_\allowbreak{}of\_\allowbreak{}pair\_\allowbreak{}swap} applies the proved pair reversal and the bounded symmetric-operator criterion.
\end{leanbox}
\section{\texorpdfstring{Time homogeneity of the actual diffusion pair law}{Time homogeneity of the actual diffusion pair law}}
\label{proof-module-WassersteinMixing-SemigroupHomogeneous-lean}


The actual state/\allowbreak{}future-path input law is the same at all deterministic cut times. Finite Euler maps preserve that input-law equality; passing their genuine strong limits proves time homogeneity of the actual two-time diffusion law and its observable compression.

\subsection{\texorpdfstring{Equal actual input laws pass to diffusion pair laws}{Equal actual input laws pass to diffusion pair laws}}
\label{report-step-WassersteinMixing-SemigroupHomogeneous-lean-1}


\textbf{Claim.} For \(a,b,t\ge0\),

\[\operatorname{Law}(Y_a,Y_{a+t})=\operatorname{Law}(Y_b,Y_{b+t}),\qquad U_a^*U_{a+t}=P_t.\]

\textbf{Proof.} Each shifted finite Euler endpoint is the same measurable function of its current state and future Brownian path. Those actual paired inputs have equal laws at all cuts, so their anchored Euler output laws agree. The actual Euler endpoints converge strongly to their respective diffusion endpoints. Their pair laws consequently converge weakly, and uniqueness of the weak limit gives equality of the actual diffusion pair laws. Equality of pair laws determines the observable compression; take the second cut to be zero.

\begin{leanbox}{Lean: \texttt{shiftedEuler}, \texttt{shiftedEuler\_\allowbreak{}pair\_\allowbreak{}law}, \texttt{state\_\allowbreak{}pair\_\allowbreak{}law} and 1 more}
\leansource{\texttt{shiftedEuler} \quad WassersteinMixing/\allowbreak{}SemigroupHomogeneous.lean, lines 22–30}
\begin{leancode}
\leanline{0}{def\ shiftedEuler\ (P\ :\ Parameters)\ (V\ :\ E\ P.dimension\ →\ ℝ)}
\leanline{4}{(hV\ :\ PotentialAssumptions\ P\ V)\ (a\ :\ ℝ≥0)\ (h\ :\ ℝ)\ (n\ :\ ℕ)\ :}
\leanline{4}{Lp\ (E\ P.dimension)\ 2\ (actualProbability\ P\ V)\ :=\ by}
\leanline{2}{let\ :=\ actualProbability\_isProbability\ P\ V\ hV}
\leanline{2}{exact\ ULA.additiveEulerChain}
\leanline{4}{(ULA.lpDrift\ (actualProbability\ P\ V)\ (gradient\ V)\ hV.gradientLipschitz)\ h}
\leanline{4}{(state\ P\ V\ hV\ a)}
\leanline{4}{(ULA.brownianLangevinNoise\ (actualProbability\ P\ V)\ (brownian\ P)}
\leanline{6}{(brownian\_preBrownian\ P\ V\ hV)\ a)\ n}
\end{leancode}
\leansource{\texttt{shiftedEuler\_\allowbreak{}pair\_\allowbreak{}law} \quad WassersteinMixing/\allowbreak{}SemigroupHomogeneous.lean, lines 34–53}
\begin{leancode}
\leanline{0}{theorem\ shiftedEuler\_pair\_law\ (P\ :\ Parameters)\ (V\ :\ E\ P.dimension\ →\ ℝ)}
\leanline{4}{(hV\ :\ PotentialAssumptions\ P\ V)\ (a\ b\ :\ ℝ≥0)\ (h\ :\ ℝ)\ (hh\ :\ 0\ ≤\ h)\ (n\ :\ ℕ)\ :}
\leanline{4}{(actualProbability\ P\ V).map\ (fun\ ω\ ↦}
\leanline{6}{(state\ P\ V\ hV\ a\ ω,\ shiftedEuler\ P\ V\ hV\ a\ h\ n\ ω))\ =}
\leanline{4}{(actualProbability\ P\ V).map\ (fun\ ω\ ↦}
\leanline{6}{(state\ P\ V\ hV\ b\ ω,\ shiftedEuler\ P\ V\ hV\ b\ h\ n\ ω))}
\end{leancode}
\leansource{\texttt{state\_\allowbreak{}pair\_\allowbreak{}law} \quad WassersteinMixing/\allowbreak{}SemigroupHomogeneous.lean, lines 57–84}
\begin{leancode}
\leanline{0}{theorem\ state\_pair\_law\ (P\ :\ Parameters)\ (V\ :\ E\ P.dimension\ →\ ℝ)}
\leanline{4}{(hV\ :\ PotentialAssumptions\ P\ V)\ (a\ b\ :\ ℝ≥0)\ (t\ :\ ℝ)\ (ht\ :\ 0\ ≤\ t)\ :}
\leanline{4}{(actualProbability\ P\ V).map\ (fun\ ω\ ↦}
\leanline{6}{(state\ P\ V\ hV\ a\ ω,\ state\ P\ V\ hV\ ((a\ :\ ℝ)\ +\ t)\ ω))\ =}
\leanline{4}{(actualProbability\ P\ V).map\ (fun\ ω\ ↦}
\leanline{6}{(state\ P\ V\ hV\ b\ ω,\ state\ P\ V\ hV\ ((b\ :\ ℝ)\ +\ t)\ ω))}
\end{leancode}
\leansource{\texttt{shifted\_\allowbreak{}compression\_\allowbreak{}eq\_\allowbreak{}op} \quad WassersteinMixing/\allowbreak{}SemigroupHomogeneous.lean, lines 91–100}
\begin{leancode}
\leanline{0}{theorem\ shifted\_compression\_eq\_op\ (P\ :\ Parameters)\ (V\ :\ E\ P.dimension\ →\ ℝ)}
\leanline{4}{(hV\ :\ PotentialAssumptions\ P\ V)\ (a\ :\ ℝ≥0)\ (t\ :\ ℝ)\ (ht\ :\ 0\ ≤\ t)\ :}
\leanline{4}{ObservableTransition.compression\ (evaluation\ ({\leanfallback 𝕜}\ :=\ {\leanfallback 𝕜})\ (F\ :=\ F)\ P\ V\ hV\ a)}
\leanline{6}{(evaluation\ P\ V\ hV\ ((a\ :\ ℝ)\ +\ t))\ =\ op\ P\ V\ hV\ t}
\end{leancode}
\leanlabel{How the Lean proof follows the argument}
\texttt{shifted\allowbreak{}Euler\_\allowbreak{}pair\_\allowbreak{}law} maps the proved equality of state/\allowbreak{}future input laws through the same actual Euler map. \texttt{state\_\allowbreak{}pair\_\allowbreak{}law} passes the actual Euler limits to their pair-law limits. \texttt{shifted\_\allowbreak{}compression\_\allowbreak{}eq\_\allowbreak{}op} invokes equality of compression from actual pair laws.
\end{leanbox}


Earlier results used here: \cite[Section~3.13.2]{LangevinLean2026ULASC}; \cite[Section~3.16.1]{LangevinLean2026ULASC}; \cite[Section~3.25.1]{LangevinLean2026ULASC}; \cite[Section~3.26.3]{LangevinLean2026ULASC}.
\section{\texorpdfstring{The actual conditional Markov identity on the Brownian past}{The actual conditional Markov identity on the Brownian past}}
\label{proof-module-WassersteinMixing-ActualMarkov-lean}


The true diffusion Markov identity is derived from whole-future independence, measurable finite Euler maps, strong Euler convergence, bounded observable convergence and density. It concerns the full initialization/\allowbreak{}Brownian past and every square-integrable target observable.

\subsection{\texorpdfstring{Actual observable representatives and adaptation}{Actual observable representatives and adaptation}}
\label{report-step-WassersteinMixing-ActualMarkov-lean-1}


\textbf{Claim.} Evaluation of a bounded continuous target function is its actual bounded composition with the canonical state. Evaluation of every square-integrable target observable at time \(a\) is measurable from the actual past at \(a\).



\textbf{Proof.} The exact Gibbs state law transports the target observable representative through evaluation. The canonical state is past-measurable; compose its measurable version with the strongly measurable observable and transport back through almost-sure equality.

\begin{leanbox}{Lean: \texttt{actualEvaluation\_\allowbreak{}bcf\_\allowbreak{}eq\_\allowbreak{}boundedComposition}, \texttt{actualEvaluation\_\allowbreak{}adapted\_\allowbreak{}pastAt}}
\leansource{\texttt{actualEvaluation\_\allowbreak{}bcf\_\allowbreak{}eq\_\allowbreak{}boundedComposition} \quad WassersteinMixing/\allowbreak{}ActualMarkov.lean, lines 25–40}
\begin{leancode}
\leanline{0}{private\ theorem\ actualEvaluation\_bcf\_eq\_boundedComposition\ (t\ :\ ℝ)\ (φ\ :\ E\ P.dimension\ →ᵇ\ H)\ :}
\leanline{4}{letI\ :=\ actualTarget\_isProbability\ P\ V\ hV}
\leanline{4}{letI\ :=\ actualProbability\_isProbability\ P\ V\ hV}
\leanline{4}{CanonicalSemigroup.evaluation\ ({\leanfallback 𝕜}\ :=\ ℝ)\ P\ V\ hV\ t}
\leanline{6}{(BoundedContinuousFunction.toLp\ 2\ (ULA.gibbsMeasure\ V)\ ℝ\ φ)\ =}
\leanline{6}{boundedComposition\ (actualProbability\ P\ V)\ (CanonicalSemigroup.state\ P\ V\ hV\ t)\ φ}
\end{leancode}
\leansource{\texttt{actualEvaluation\_\allowbreak{}adapted\_\allowbreak{}pastAt} \quad WassersteinMixing/\allowbreak{}ActualMarkov.lean, lines 45–55}
\begin{leancode}
\leanline{0}{theorem\ actualEvaluation\_adapted\_pastAt\ (a\ :\ ℝ≥0)\ (f\ :\ Lp\ H\ 2\ (ULA.gibbsMeasure\ V))\ :}
\leanline{4}{AEStronglyMeasurable[ULA.lpLangevinPast\ (actualProbability\ P\ V)}
\leanline{6}{(CanonicalSemigroup.initial\ P\ V\ hV)\ (actualBrownianNoise\ P\ V\ hV)\ (a\ :\ ℝ)]}
\leanline{6}{(CanonicalSemigroup.evaluation\ ({\leanfallback 𝕜}\ :=\ ℝ)\ P\ V\ hV\ a\ f)\ (actualProbability\ P\ V)}
\end{leancode}
\leanlabel{How the Lean proof follows the argument}
The bounded-continuous helper compares exact pullback and bounded-composition representatives. \texttt{actual\allowbreak{}Evaluation\_\allowbreak{}adapted\_\allowbreak{}past\allowbreak{}At} combines actual state adaptation with measurable composition and the measure-preserving evaluation representative.
\end{leanbox}


Earlier results used here: \cite[Section~3.26.3]{LangevinLean2026ULASC}; \cite[Section~3.54.4]{LangevinLean2026ULASC}.
\subsection{\texorpdfstring{Proving the conditional mean is a function of the actual current state}{Proving the conditional mean is a function of the actual current state}}
\label{report-step-WassersteinMixing-ActualMarkov-lean-2}


\textbf{Claim.} \(\mathbb E[f(Y_{a+t})\mid\mathcal F_a]\) is measurable from \(Y_a\) alone.



\textbf{Proof.} First take bounded continuous \(f\). The finite shifted Euler endpoint is a measurable function of \((Y_a,W_a)\), and the entire future path is independent of the actual past. Conditional averaging therefore makes each finite Euler conditional mean a function of \(Y_a\). The actual Euler endpoint converges strongly to \(Y_{a+t}\); bounded-composition convergence and closedness of the conditional state-measurability property pass this fact to the diffusion. For arbitrary \(L^2(\pi)\) observables, use density of bounded continuous functions, continuity of evaluation and conditional expectation, and the same closed measurable subspace.

\begin{leanbox}{Lean: \texttt{actual\_\allowbreak{}conditional\_\allowbreak{}state\_\allowbreak{}measurable}}
\leansource{\texttt{actual\_\allowbreak{}conditional\_\allowbreak{}state\_\allowbreak{}measurable} \quad WassersteinMixing/\allowbreak{}ActualMarkov.lean, lines 59–114}
\begin{leancode}
\leanline{0}{theorem\ actual\_conditional\_state\_measurable\ (a\ :\ ℝ≥0)\ (t\ :\ ℝ)\ (ht\ :\ 0\ ≤\ t)}
\leanline{4}{(f\ :\ Lp\ H\ 2\ (ULA.gibbsMeasure\ V))\ :}
\leanline{4}{AEStronglyMeasurable[(inferInstance\ :\ MeasurableSpace\ (E\ P.dimension)).comap}
\leanline{6}{(CanonicalSemigroup.state\ P\ V\ hV\ a)]}
\leanline{6}{(condExpL2\ H\ ℝ\ (ULA.lpLangevinPast\_le\ (actualProbability\ P\ V)}
\leanline{8}{(CanonicalSemigroup.initial\ P\ V\ hV)\ (actualBrownianNoise\ P\ V\ hV)\ (a\ :\ ℝ))}
\leanline{8}{(CanonicalSemigroup.evaluation\ ({\leanfallback 𝕜}\ :=\ ℝ)\ P\ V\ hV\ (a\ +\ t)\ f)\ :\ Lp\ H\ 2\ (actualProbability\ P\ V))}
\leanline{6}{(actualProbability\ P\ V)}
\end{leancode}
\leanlabel{How the Lean proof follows the argument}
\texttt{actual\_\allowbreak{}conditional\_\allowbreak{}state\_\allowbreak{}measurable} applies the independent-input result to \texttt{pointwise\allowbreak{}Euler\allowbreak{}Map}, uses actual shifted Euler convergence, then applies \texttt{cond\allowbreak{}Exp\allowbreak{}L2\_\allowbreak{}state\_\allowbreak{}measurable\_\allowbreak{}of\_\allowbreak{}tendsto}. A closed-preimage argument and bounded-continuous density extend the result to every target observable.
\end{leanbox}


Earlier results used here: \cite[Section~3.13.2]{LangevinLean2026ULASC}; \cite[Section~3.16.1]{LangevinLean2026ULASC}; \cite[Section~3.25.1]{LangevinLean2026ULASC}; \cite[Section~3.26.3]{LangevinLean2026ULASC}; \cite[Section~3.54.4]{LangevinLean2026ULASC}; \cite[Section~3.54.5]{LangevinLean2026ULASC}.
\subsection{\texorpdfstring{The exact actual transition in the conditional identity}{The exact actual transition in the conditional identity}}
\label{report-step-WassersteinMixing-ActualMarkov-lean-3}


\textbf{Claim.}

\[\mathbb E[f(Y_{a+t})\mid\mathcal F_a]=(P_tf)(Y_a),\qquad a,t\ge0.\]

\textbf{Proof.} The preceding state-measurability result and square-integrable Doob–Dynkin factorization put the conditional mean in the range of actual current-state evaluation. Since this evaluation is adapted, its adjoint compression identifies that conditional mean exactly. Actual time homogeneity identifies the compression from \(a\) to \(a+t\) with the canonical time-zero operator \(P_t\). Substitution gives the displayed identity, without an assumed Markov property.

\begin{leanbox}{Lean: \texttt{actualEvaluation\_\allowbreak{}markov\_\allowbreak{}pastAt}}
\leansource{\texttt{actualEvaluation\_\allowbreak{}markov\_\allowbreak{}pastAt} \quad WassersteinMixing/\allowbreak{}ActualMarkov.lean, lines 120–145}
\begin{leancode}
\leanline{0}{theorem\ actualEvaluation\_markov\_pastAt\ (a\ :\ ℝ≥0)\ (t\ :\ ℝ)\ (ht\ :\ 0\ ≤\ t)}
\leanline{4}{(f\ :\ Lp\ H\ 2\ (ULA.gibbsMeasure\ V))\ :}
\leanline{4}{(condExpL2\ H\ ℝ\ (ULA.lpLangevinPast\_le\ (actualProbability\ P\ V)}
\leanline{6}{(CanonicalSemigroup.initial\ P\ V\ hV)\ (actualBrownianNoise\ P\ V\ hV)\ (a\ :\ ℝ))}
\leanline{6}{(CanonicalSemigroup.evaluation\ ({\leanfallback 𝕜}\ :=\ ℝ)\ P\ V\ hV\ (a\ +\ t)\ f)\ :\ Lp\ H\ 2\ (actualProbability\ P\ V))\ =}
\leanline{6}{CanonicalSemigroup.evaluation\ ({\leanfallback 𝕜}\ :=\ ℝ)\ P\ V\ hV\ a}
\leanline{8}{(CanonicalSemigroup.op\ ({\leanfallback 𝕜}\ :=\ ℝ)\ P\ V\ hV\ t\ f)}
\end{leancode}
\leanlabel{How the Lean proof follows the argument}
\texttt{actual\allowbreak{}Evaluation\_\allowbreak{}markov\_\allowbreak{}past\allowbreak{}At} uses \texttt{observable\_\allowbreak{}range\_\allowbreak{}of\_\allowbreak{}state\_\allowbreak{}measurable}, the actual evaluation adaptation theorem, \texttt{conditional\_\allowbreak{}eq\_\allowbreak{}evaluation\_\allowbreak{}adjoint} and \texttt{shifted\_\allowbreak{}compression\_\allowbreak{}eq\_\allowbreak{}op}. Every premise is a previously proved fact about the constructed Brownian solution.
\end{leanbox}


Earlier results used here: \cite[Section~3.26.3]{LangevinLean2026ULASC}; \cite[Section~3.54.5]{LangevinLean2026ULASC}.
\section{\texorpdfstring{Strong limits of bounded increasing positive operators}{Strong limits of bounded increasing positive operators}}
\label{proof-module-WassersteinMixing-MonotoneOperators-lean}


A sequence of positive operators with \(0\le A_n\le A_{n+1}\le cI\) has a bounded linear strong limit. This supplies an actual corrector operator from increasing finite sums, without choosing an eigenbasis.

\subsection{\texorpdfstring{The quadratic estimate behind convergence}{The quadratic estimate behind convergence}}
\label{report-step-WassersteinMixing-MonotoneOperators-lean-1}


\textbf{Claim.} If \(0\le A\le cI\), then \(\|Ax\|^2\le c\langle Ax,x\rangle_{\mathbb R}\).



\textbf{Proof.} Positivity implies self-adjointness. The positive operators \(cI-A\) and \(A\) commute, so their product is positive. Its quadratic form is \(c\langle Ax,x\rangle-\langle A^2x,x\rangle\). Self-adjointness identifies the last term with \(\|Ax\|^2\). Rearrangement yields the bound.

\begin{leanbox}{Lean: \texttt{instance\_\allowbreak{}at\_\allowbreak{}20}, \texttt{positive\_\allowbreak{}operator\_\allowbreak{}square\_\allowbreak{}le}}
\leansource{\texttt{instance\_\allowbreak{}at\_\allowbreak{}20} \quad WassersteinMixing/\allowbreak{}MonotoneOperators.lean, lines 20–20}
\begin{leancode}
\leanline{0}{local\ instance\ :\ InnerProductSpace\ ℝ\ H\ :=\ InnerProductSpace.rclikeToReal\ ℂ\ H}
\end{leancode}
\leansource{\texttt{positive\_\allowbreak{}operator\_\allowbreak{}square\_\allowbreak{}le} \quad WassersteinMixing/\allowbreak{}MonotoneOperators.lean, lines 23–39}
\begin{leancode}
\leanline{0}{theorem\ positive\_operator\_square\_le\ (A\ :\ H\ →L[ℂ]\ H)\ (c\ :\ ℝ)}
\leanline{4}{(hA\ :\ 0\ ≤\ A)\ (hupper\ :\ A\ ≤\ c\ •\ (1\ :\ H\ →L[ℂ]\ H))\ (x\ :\ H)\ :}
\leanline{4}{‖A\ x‖\ \textasciicircum{}\ 2\ ≤\ c\ *\ inner\ ℝ\ (A\ x)\ x}
\end{leancode}
\leanlabel{How the Lean proof follows the argument}
The local scalar restriction gives the real part of the complex inner product. \texttt{positive\_\allowbreak{}operator\_\allowbreak{}square\_\allowbreak{}le} uses positivity of the commuting product and the proved symmetric pairing identity.
\end{leanbox}
\subsection{\texorpdfstring{Constructing the genuine bounded strong limit}{Constructing the genuine bounded strong limit}}
\label{report-step-WassersteinMixing-MonotoneOperators-lean-2}


\textbf{Claim.} There is a bounded linear \(G\) with \(A_nx\to Gx\) for every \(x\), \(\|G\|\le c\), and \(0\le G\le cI\).



\textbf{Proof.} The increasing real numbers \(\langle A_nx,x\rangle\) are bounded by \(c\|x\|^2\), hence are Cauchy. For \(m\ge n\), apply the preceding estimate to \(A_m-A_n\); its scalar quadratic difference tending to zero forces \(A_nx\) to be Cauchy. Completeness gives the pointwise limit. Passing sums and scalar multiples through limits proves linearity, and the uniform norm bound proves continuity. Symmetry and the lower and upper quadratic inequalities pass to the limit.

\begin{leanbox}{Lean: \texttt{monotone\_\allowbreak{}operator\_\allowbreak{}cauchy}, \texttt{exists\_\allowbreak{}monotone\_\allowbreak{}operator\_\allowbreak{}limit}, \texttt{strong\_\allowbreak{}limit\_\allowbreak{}order}}
\leansource{\texttt{monotone\_\allowbreak{}operator\_\allowbreak{}cauchy} \quad WassersteinMixing/\allowbreak{}MonotoneOperators.lean, lines 43–89}
\begin{leancode}
\leanline{0}{theorem\ monotone\_operator\_cauchy\ (A\ :\ ℕ\ →\ H\ →L[ℂ]\ H)\ (c\ :\ ℝ)\ (hc\ :\ 0\ ≤\ c)}
\leanline{4}{(hmono\ :\ Monotone\ A)\ (hpos\ :\ ∀\ n,\ 0\ ≤\ A\ n)}
\leanline{4}{(hupper\ :\ ∀\ n,\ A\ n\ ≤\ c\ •\ (1\ :\ H\ →L[ℂ]\ H))\ (x\ :\ H)\ :}
\leanline{4}{CauchySeq\ (fun\ n\ ↦\ A\ n\ x)}
\end{leancode}
\leansource{\texttt{exists\_\allowbreak{}monotone\_\allowbreak{}operator\_\allowbreak{}limit} \quad WassersteinMixing/\allowbreak{}MonotoneOperators.lean, lines 93–119}
\begin{leancode}
\leanline{0}{theorem\ exists\_monotone\_operator\_limit\ (A\ :\ ℕ\ →\ H\ →L[ℂ]\ H)\ (c\ :\ ℝ)\ (hc\ :\ 0\ ≤\ c)}
\leanline{4}{(hmono\ :\ Monotone\ A)\ (hpos\ :\ ∀\ n,\ 0\ ≤\ A\ n)}
\leanline{4}{(hupper\ :\ ∀\ n,\ A\ n\ ≤\ c\ •\ (1\ :\ H\ →L[ℂ]\ H))\ :}
\leanline{4}{∃\ G\ :\ H\ →L[ℂ]\ H,\ ‖G‖\ ≤\ c\ ∧}
\leanline{6}{∀\ x,\ Tendsto\ (fun\ n\ ↦\ A\ n\ x)\ atTop\ ({\leanfallback 𝓝}\ (G\ x))}
\end{leancode}
\leansource{\texttt{strong\_\allowbreak{}limit\_\allowbreak{}order} \quad WassersteinMixing/\allowbreak{}MonotoneOperators.lean, lines 122–151}
\begin{leancode}
\leanline{0}{theorem\ strong\_limit\_order\ (A\ :\ ℕ\ →\ H\ →L[ℂ]\ H)\ (G\ :\ H\ →L[ℂ]\ H)\ (c\ :\ ℝ)}
\leanline{4}{(hpos\ :\ ∀\ n,\ 0\ ≤\ A\ n)\ (hupper\ :\ ∀\ n,\ A\ n\ ≤\ c\ •\ (1\ :\ H\ →L[ℂ]\ H))}
\leanline{4}{(hlim\ :\ ∀\ x,\ Tendsto\ (fun\ n\ ↦\ A\ n\ x)\ atTop\ ({\leanfallback 𝓝}\ (G\ x)))\ :}
\leanline{4}{0\ ≤\ G\ ∧\ G\ ≤\ c\ •\ (1\ :\ H\ →L[ℂ]\ H)}
\end{leancode}
\leanlabel{How the Lean proof follows the argument}
\texttt{monotone\_\allowbreak{}operator\_\allowbreak{}cauchy} uses monotone convergence of the scalar quadratic forms. \texttt{exists\_\allowbreak{}monotone\_\allowbreak{}operator\_\allowbreak{}limit} constructs the linear map from the actual pointwise limits. \texttt{strong\_\allowbreak{}limit\_\allowbreak{}order} transfers symmetry and both form inequalities through those limits.
\end{leanbox}
\subsection{\texorpdfstring{A commuting operator preserves the sharp energy factor}{A commuting operator preserves the sharp energy factor}}
\label{report-step-WassersteinMixing-MonotoneOperators-lean-3}


\textbf{Claim.} If \(G\) is self-adjoint, \(\|G\|\le h\), \(B\ge0\), and \(GB=BG\), then

\[\langle Gx,BGx\rangle_{\mathbb R}\le h^2\langle x,Bx\rangle_{\mathbb R}.\]

\textbf{Proof.} The operator \(h^2I-G^2\) is positive because its quadratic form is \(h^2\|x\|^2-\|Gx\|^2\). It commutes with \(B\), so their product is positive. Its nonnegative quadratic form, together with symmetry and commutation, is exactly the displayed difference.

\begin{leanbox}{Lean: \texttt{commuting\_\allowbreak{}energy\_\allowbreak{}bound}}
\leansource{\texttt{commuting\_\allowbreak{}energy\_\allowbreak{}bound} \quad WassersteinMixing/\allowbreak{}MonotoneOperators.lean, lines 155–190}
\begin{leancode}
\leanline{0}{theorem\ commuting\_energy\_bound\ (G\ B\ :\ H\ →L[ℂ]\ H)\ (h\ :\ ℝ)\ (hh\ :\ 0\ ≤\ h)}
\leanline{4}{(hG\ :\ IsSelfAdjoint\ G)\ (hnorm\ :\ ‖G‖\ ≤\ h)\ (hB\ :\ 0\ ≤\ B)}
\leanline{4}{(hcomm\ :\ Commute\ G\ B)\ (x\ :\ H)\ :}
\leanline{4}{inner\ ℝ\ (G\ x)\ (B\ (G\ x))\ ≤\ h\ \textasciicircum{}\ 2\ *\ inner\ ℝ\ x\ (B\ x)}
\end{leancode}
\leanlabel{How the Lean proof follows the argument}
\texttt{commuting\_\allowbreak{}energy\_\allowbreak{}bound} constructs the positive squared-norm deficit, proves its commutation with the energy operator and expands the resulting positive product pairing.
\end{leanbox}
\section{\texorpdfstring{A bounded Poisson corrector without a spectral basis}{A bounded Poisson corrector without a spectral basis}}
\label{proof-module-WassersteinMixing-SemigroupPoisson-lean}


Let \((P_t)_{t\ge0}\) be a strongly continuous self-adjoint contraction semigroup on a complex Hilbert space. Define \(D_h=\int_0^h(I-P_t)\,dt\) and finite sums \(G_N=\sum_{k<N}P_{kh}D_h\). Positivity and telescoping construct a genuine bounded strong limit \(G_h\), with the exact discrete Poisson equation and energy domination.

\subsection{\texorpdfstring{Semigroup positivity, order and integrability}{Semigroup positivity, order and integrability}}
\label{report-step-WassersteinMixing-SemigroupPoisson-lean-1}


\textbf{Claim.} The transition operators commute, \(0\preceq P_t\preceq I\), and \(P_t\preceq P_s\) when \(0\le s\le t\). Every orbit is integrable on a compact time interval.



\textbf{Proof.} The semigroup law and commutativity of addition give commutation. Since \(P_t=P_{t/2}^*P_{t/2}\), it is positive. Contractivity and Cauchy–Schwarz bound its quadratic form by \(\|x\|^2\), giving \(P_t\preceq I\). The positive commuting product \(P_s(I-P_{t-s})\) is \(P_s-P_t\). Strong continuity supplies compact-interval Bochner integrability.

\begin{leanbox}{Lean: \texttt{instance\_\allowbreak{}at\_\allowbreak{}26}, \texttt{ReversibleSemigroup}, \texttt{innerLeft} and 7 more}
\leansource{\texttt{instance\_\allowbreak{}at\_\allowbreak{}26} \quad WassersteinMixing/\allowbreak{}SemigroupPoisson.lean, lines 26–26}
\begin{leancode}
\leanline{0}{local\ instance\ :\ InnerProductSpace\ ℝ\ H\ :=\ InnerProductSpace.rclikeToReal\ ℂ\ H}
\end{leancode}
\leansource{\texttt{ReversibleSemigroup} \quad WassersteinMixing/\allowbreak{}SemigroupPoisson.lean, lines 29–38}
\begin{leancode}
\leanline{0}{structure\ ReversibleSemigroup\ (H\ :\ Type*)\ [NormedAddCommGroup\ H]}
\leanline{4}{[InnerProductSpace\ ℂ\ H]\ [CompleteSpace\ H]\ where}
\leanline{2}{op\ :\ ℝ\ →\ H\ →L[ℂ]\ H}
\leanline{2}{zero\ :\ op\ 0\ =\ 1}
\leanline{2}{add\ :\ ∀\ \{s\ t\},\ 0\ ≤\ s\ →\ 0\ ≤\ t\ →\ op\ (s\ +\ t)\ =\ op\ s\ *\ op\ t}
\leanline{2}{selfAdjoint\ :\ ∀\ \{t\},\ 0\ ≤\ t\ →\ IsSelfAdjoint\ (op\ t)}
\leanline{2}{norm\_le\ :\ ∀\ \{t\},\ 0\ ≤\ t\ →\ ∀\ x,\ ‖op\ t\ x‖\ ≤\ ‖x‖}
\leanline{2}{continuous\ :\ ∀\ x,\ ContinuousOn\ (fun\ t\ ↦\ op\ t\ x)\ (Set.Ici\ 0)}
\leanline{0}{\strut}
\leanline{0}{namespace\ ReversibleSemigroup}
\end{leancode}
\leansource{\texttt{innerLeft} \quad WassersteinMixing/\allowbreak{}SemigroupPoisson.lean, lines 42–43}
\begin{leancode}
\leanline{0}{private\ def\ innerLeft\ (x\ :\ H)\ :\ H\ →L[ℝ]\ ℂ\ :=}
\leanline{2}{RCLike.conjCLE.toContinuousLinearMap\ ∘L\ (innerSL\ ℂ\ x).restrictScalars\ ℝ}
\end{leancode}
\leansource{\texttt{innerLeft\_\allowbreak{}apply} \quad WassersteinMixing/\allowbreak{}SemigroupPoisson.lean, lines 45–46}
\begin{leancode}
\leanline{0}{private\ theorem\ innerLeft\_apply\ (x\ y\ :\ H)\ :\ innerLeft\ x\ y\ =\ inner\ ℂ\ y\ x}
\end{leancode}
\leansource{\texttt{linear\_\allowbreak{}intervalIntegrable} \quad WassersteinMixing/\allowbreak{}SemigroupPoisson.lean, lines 48–52}
\begin{leancode}
\leanline{0}{private\ theorem\ linear\_intervalIntegrable\ \{K\ :\ Type*\}\ [NormedAddCommGroup\ K]}
\leanline{4}{[NormedSpace\ ℝ\ K]\ (L\ :\ H\ →L[ℝ]\ K)\ \{f\ :\ ℝ\ →\ H\}\ \{a\ b\ :\ ℝ\}}
\leanline{4}{(hf\ :\ IntervalIntegrable\ f\ volume\ a\ b)\ :}
\leanline{4}{IntervalIntegrable\ (fun\ t\ ↦\ L\ (f\ t))\ volume\ a\ b}
\end{leancode}
\leansource{\texttt{commute} \quad WassersteinMixing/\allowbreak{}SemigroupPoisson.lean, lines 54–57}
\begin{leancode}
\leanline{0}{theorem\ commute\ \{s\ t\ :\ ℝ\}\ (hs\ :\ 0\ ≤\ s)\ (ht\ :\ 0\ ≤\ t)\ :}
\leanline{4}{Commute\ (S.op\ s)\ (S.op\ t)}
\end{leancode}
\leansource{\texttt{op\_\allowbreak{}nonneg} \quad WassersteinMixing/\allowbreak{}SemigroupPoisson.lean, lines 59–66}
\begin{leancode}
\leanline{0}{theorem\ op\_nonneg\ \{t\ :\ ℝ\}\ (ht\ :\ 0\ ≤\ t)\ :\ 0\ ≤\ S.op\ t}
\end{leancode}
\leansource{\texttt{op\_\allowbreak{}le\_\allowbreak{}one} \quad WassersteinMixing/\allowbreak{}SemigroupPoisson.lean, lines 68–82}
\begin{leancode}
\leanline{0}{theorem\ op\_le\_one\ \{t\ :\ ℝ\}\ (ht\ :\ 0\ ≤\ t)\ :\ S.op\ t\ ≤\ 1}
\end{leancode}
\leansource{\texttt{antitone} \quad WassersteinMixing/\allowbreak{}SemigroupPoisson.lean, lines 84–95}
\begin{leancode}
\leanline{0}{theorem\ antitone\ :\ AntitoneOn\ S.op\ (Set.Ici\ 0)}
\end{leancode}
\leansource{\texttt{orbit\_\allowbreak{}integrable} \quad WassersteinMixing/\allowbreak{}SemigroupPoisson.lean, lines 97–101}
\begin{leancode}
\leanline{0}{theorem\ orbit\_integrable\ (x\ :\ H)\ \{h\ :\ ℝ\}\ (hh\ :\ 0\ ≤\ h)\ :}
\leanline{4}{IntervalIntegrable\ (fun\ t\ ↦\ S.op\ t\ x)\ volume\ 0\ h}
\end{leancode}
\leanlabel{How the Lean proof follows the argument}
\texttt{Reversible\allowbreak{}Semigroup} packages operator properties that must be proved for a concrete process. The order lemmas derive positivity and monotonicity from those fields; \texttt{orbit\_\allowbreak{}integrable} restricts the strong continuity theorem to the compact interval. The private inner functional and integrability helper justify the later pairing integrals.
\end{leanbox}
\subsection{\texorpdfstring{The integrated defect and its telescoping order bound}{The integrated defect and its telescoping order bound}}
\label{report-step-WassersteinMixing-SemigroupPoisson-lean-2}


\textbf{Claim.} The actual integral \(D_hx=\int_0^h(x-P_tx)\,dt\) is a bounded linear map, is symmetric, and commutes with every \(P_s\). Moreover

\[0\preceq P_sD_h\preceq h(P_s-P_{s+h}).\]

\textbf{Proof.} The integrand has norm at most \(2\|x\|\), giving boundedness and allowing linearity to pass through the integral. Bounded linear operators and inner-product functionals commute with that integral. Semigroup commutation gives the operator commutation, and self-adjointness gives symmetry. Finally \(P_sD_hx=\int_0^h(P_s-P_{s+t})x\,dt\). Time monotonicity bounds the integrand between zero and \(P_s-P_{s+h}\), proving the order estimate.

\begin{leanbox}{Lean: \texttt{defect}, \texttt{defect\_\allowbreak{}integrable}, \texttt{defect\_\allowbreak{}norm\_\allowbreak{}le} and 6 more}
\leansource{\texttt{defect} \quad WassersteinMixing/\allowbreak{}SemigroupPoisson.lean, lines 104–104}
\begin{leancode}
\leanline{0}{def\ defect\ (h\ :\ ℝ)\ (x\ :\ H)\ :\ H\ :=\ ∫\ t\ in\ 0..h,\ x\ -\ S.op\ t\ x}
\end{leancode}
\leansource{\texttt{defect\_\allowbreak{}integrable} \quad WassersteinMixing/\allowbreak{}SemigroupPoisson.lean, lines 106–108}
\begin{leancode}
\leanline{0}{theorem\ defect\_integrable\ (x\ :\ H)\ \{h\ :\ ℝ\}\ (hh\ :\ 0\ ≤\ h)\ :}
\leanline{4}{IntervalIntegrable\ (fun\ t\ ↦\ x\ -\ S.op\ t\ x)\ volume\ 0\ h}
\end{leancode}
\leansource{\texttt{defect\_\allowbreak{}norm\_\allowbreak{}le} \quad WassersteinMixing/\allowbreak{}SemigroupPoisson.lean, lines 110–119}
\begin{leancode}
\leanline{0}{theorem\ defect\_norm\_le\ (x\ :\ H)\ \{h\ :\ ℝ\}\ (hh\ :\ 0\ ≤\ h)\ :}
\leanline{4}{‖S.defect\ h\ x‖\ ≤\ (2\ *\ h)\ *\ ‖x‖}
\end{leancode}
\leansource{\texttt{defectOp} \quad WassersteinMixing/\allowbreak{}SemigroupPoisson.lean, lines 122–136}
\begin{leancode}
\leanline{0}{def\ defectOp\ (h\ :\ ℝ)\ (hh\ :\ 0\ ≤\ h)\ :\ H\ →L[ℂ]\ H\ :=}
\leanline{2}{LinearMap.mkContinuous}
\leanline{4}{\{\ toFun\ :=\ S.defect\ h}
\leanline{6}{map\_add'\ :=\ by}
\leanline{8}{intro\ x\ y}
\leanline{8}{simp\ only\ [defect,\ map\_add]}
\leanline{8}{simp\_rw\ [show\ ∀\ t,\ x\ +\ y\ -\ (S.op\ t\ x\ +\ S.op\ t\ y)\ =}
\leanline{10}{(x\ -\ S.op\ t\ x)\ +\ (y\ -\ S.op\ t\ y)\ by\ intro\ t;\ abel]}
\leanline{8}{exact\ intervalIntegral.integral\_add\ (S.defect\_integrable\ x\ hh)}
\leanline{10}{(S.defect\_integrable\ y\ hh)}
\leanline{6}{map\_smul'\ :=\ by}
\leanline{8}{intro\ c\ x}
\leanline{8}{simp\ only\ [defect,\ map\_smul,\ RingHom.id\_apply,\ ←\ smul\_sub]}
\leanline{8}{exact\ intervalIntegral.integral\_smul\ c\ \_\ \}}
\leanline{4}{(2\ *\ h)\ (fun\ x\ ↦\ S.defect\_norm\_le\ x\ hh)}
\end{leancode}
\leansource{\texttt{defectOp\_\allowbreak{}apply} \quad WassersteinMixing/\allowbreak{}SemigroupPoisson.lean, lines 138–139}
\begin{leancode}
\leanline{0}{@[simp]\ theorem\ defectOp\_apply\ (h\ :\ ℝ)\ (hh\ :\ 0\ ≤\ h)\ (x\ :\ H)\ :}
\leanline{4}{S.defectOp\ h\ hh\ x\ =\ S.defect\ h\ x}
\end{leancode}
\leansource{\texttt{defectOp\_\allowbreak{}commute} \quad WassersteinMixing/\allowbreak{}SemigroupPoisson.lean, lines 141–153}
\begin{leancode}
\leanline{0}{theorem\ defectOp\_commute\ (h\ :\ ℝ)\ (hh\ :\ 0\ ≤\ h)\ \{s\ :\ ℝ\}\ (hs\ :\ 0\ ≤\ s)\ :}
\leanline{4}{Commute\ (S.op\ s)\ (S.defectOp\ h\ hh)}
\end{leancode}
\leansource{\texttt{defectOp\_\allowbreak{}symmetric} \quad WassersteinMixing/\allowbreak{}SemigroupPoisson.lean, lines 155–175}
\begin{leancode}
\leanline{0}{theorem\ defectOp\_symmetric\ (h\ :\ ℝ)\ (hh\ :\ 0\ ≤\ h)\ :}
\leanline{4}{(S.defectOp\ h\ hh).IsSymmetric}
\end{leancode}
\leansource{\texttt{shiftedDefect\_\allowbreak{}symmetric} \quad WassersteinMixing/\allowbreak{}SemigroupPoisson.lean, lines 178–191}
\begin{leancode}
\leanline{0}{theorem\ shiftedDefect\_symmetric\ (h\ :\ ℝ)\ (hh\ :\ 0\ ≤\ h)\ \{s\ :\ ℝ\}\ (hs\ :\ 0\ ≤\ s)\ :}
\leanline{4}{(S.op\ s\ *\ S.defectOp\ h\ hh).IsSymmetric}
\end{leancode}
\leansource{\texttt{shiftedDefect\_\allowbreak{}bounds} \quad WassersteinMixing/\allowbreak{}SemigroupPoisson.lean, lines 195–273}
\begin{leancode}
\leanline{0}{theorem\ shiftedDefect\_bounds\ (h\ :\ ℝ)\ (hh\ :\ 0\ ≤\ h)\ \{s\ :\ ℝ\}\ (hs\ :\ 0\ ≤\ s)\ :}
\leanline{4}{0\ ≤\ S.op\ s\ *\ S.defectOp\ h\ hh\ ∧}
\leanline{6}{S.op\ s\ *\ S.defectOp\ h\ hh\ ≤\ h\ •\ (S.op\ s\ -\ S.op\ (s\ +\ h))}
\end{leancode}
\leanlabel{How the Lean proof follows the argument}
\texttt{defect} is a genuine vector-valued interval integral; \texttt{defect\allowbreak{}Op} is its constructed bounded linear form. The commutation and symmetry proofs move continuous linear maps through the justified integrals. \texttt{shifted\allowbreak{}Defect\_\allowbreak{}bounds} integrates nonnegative quadratic forms and their uniform upper bound.
\end{leanbox}
\subsection{\texorpdfstring{Constructing the corrector by bounded monotone convergence}{Constructing the corrector by bounded monotone convergence}}
\label{report-step-WassersteinMixing-SemigroupPoisson-lean-3}


\textbf{Claim.} The finite sums are increasing and satisfy

\[0\preceq G_N\preceq h(I-P_{Nh})\preceq hI.\]

They converge strongly to a bounded positive operator \(G_h\), with \(\|G_h\|\le h\).



\textbf{Proof.} Each summand is positive by the previous order bound. Summing that bound telescopes the consecutive semigroup differences. Thus the positive sums are increasing and uniformly bounded by \(hI\). The proved monotone-operator convergence theorem constructs their bounded linear pointwise limit, and positivity and the upper bound pass to it.

\begin{leanbox}{Lean: \texttt{partialCorrector}, \texttt{partialCorrector\_\allowbreak{}succ}, \texttt{partialCorrector\_\allowbreak{}nonneg} and 7 more}
\leansource{\texttt{partialCorrector} \quad WassersteinMixing/\allowbreak{}SemigroupPoisson.lean, lines 276–277}
\begin{leancode}
\leanline{0}{def\ partialCorrector\ (h\ :\ ℝ)\ (hh\ :\ 0\ ≤\ h)\ (N\ :\ ℕ)\ :\ H\ →L[ℂ]\ H\ :=}
\leanline{2}{∑\ n\ ∈\ Finset.range\ N,\ S.op\ (n\ *\ h)\ *\ S.defectOp\ h\ hh}
\end{leancode}
\leansource{\texttt{partialCorrector\_\allowbreak{}succ} \quad WassersteinMixing/\allowbreak{}SemigroupPoisson.lean, lines 279–282}
\begin{leancode}
\leanline{0}{theorem\ partialCorrector\_succ\ (h\ :\ ℝ)\ (hh\ :\ 0\ ≤\ h)\ (N\ :\ ℕ)\ :}
\leanline{4}{S.partialCorrector\ h\ hh\ (N\ +\ 1)\ =\ S.partialCorrector\ h\ hh\ N\ +}
\leanline{6}{S.op\ (N\ *\ h)\ *\ S.defectOp\ h\ hh}
\end{leancode}
\leansource{\texttt{partialCorrector\_\allowbreak{}nonneg} \quad WassersteinMixing/\allowbreak{}SemigroupPoisson.lean, lines 284–288}
\begin{leancode}
\leanline{0}{theorem\ partialCorrector\_nonneg\ (h\ :\ ℝ)\ (hh\ :\ 0\ ≤\ h)\ (N\ :\ ℕ)\ :}
\leanline{4}{0\ ≤\ S.partialCorrector\ h\ hh\ N}
\end{leancode}
\leansource{\texttt{partialCorrector\_\allowbreak{}mono} \quad WassersteinMixing/\allowbreak{}SemigroupPoisson.lean, lines 290–296}
\begin{leancode}
\leanline{0}{theorem\ partialCorrector\_mono\ (h\ :\ ℝ)\ (hh\ :\ 0\ ≤\ h)\ :}
\leanline{4}{Monotone\ (S.partialCorrector\ h\ hh)}
\end{leancode}
\leansource{\texttt{partialCorrector\_\allowbreak{}le\_\allowbreak{}telescoping} \quad WassersteinMixing/\allowbreak{}SemigroupPoisson.lean, lines 299–308}
\begin{leancode}
\leanline{0}{theorem\ partialCorrector\_le\_telescoping\ (h\ :\ ℝ)\ (hh\ :\ 0\ ≤\ h)\ (N\ :\ ℕ)\ :}
\leanline{4}{S.partialCorrector\ h\ hh\ N\ ≤\ h\ •\ (1\ -\ S.op\ (N\ *\ h))}
\end{leancode}
\leansource{\texttt{partialCorrector\_\allowbreak{}le} \quad WassersteinMixing/\allowbreak{}SemigroupPoisson.lean, lines 310–314}
\begin{leancode}
\leanline{0}{theorem\ partialCorrector\_le\ (h\ :\ ℝ)\ (hh\ :\ 0\ ≤\ h)\ (N\ :\ ℕ)\ :}
\leanline{4}{S.partialCorrector\ h\ hh\ N\ ≤\ h\ •\ (1\ :\ H\ →L[ℂ]\ H)}
\end{leancode}
\leansource{\texttt{corrector} \quad WassersteinMixing/\allowbreak{}SemigroupPoisson.lean, lines 317–320}
\begin{leancode}
\leanline{0}{def\ corrector\ (h\ :\ ℝ)\ (hh\ :\ 0\ ≤\ h)\ :\ H\ →L[ℂ]\ H\ :=}
\leanline{2}{Classical.choose\ (exists\_monotone\_operator\_limit\ (S.partialCorrector\ h\ hh)\ h\ hh}
\leanline{4}{(S.partialCorrector\_mono\ h\ hh)\ (S.partialCorrector\_nonneg\ h\ hh)}
\leanline{4}{(S.partialCorrector\_le\ h\ hh))}
\end{leancode}
\leansource{\texttt{corrector\_\allowbreak{}norm\_\allowbreak{}le} \quad WassersteinMixing/\allowbreak{}SemigroupPoisson.lean, lines 322–325}
\begin{leancode}
\leanline{0}{theorem\ corrector\_norm\_le\ (h\ :\ ℝ)\ (hh\ :\ 0\ ≤\ h)\ :\ ‖S.corrector\ h\ hh‖\ ≤\ h}
\end{leancode}
\leansource{\texttt{partialCorrector\_\allowbreak{}tendsto} \quad WassersteinMixing/\allowbreak{}SemigroupPoisson.lean, lines 327–331}
\begin{leancode}
\leanline{0}{theorem\ partialCorrector\_tendsto\ (h\ :\ ℝ)\ (hh\ :\ 0\ ≤\ h)\ (x\ :\ H)\ :}
\leanline{4}{Tendsto\ (fun\ N\ ↦\ S.partialCorrector\ h\ hh\ N\ x)\ atTop\ ({\leanfallback 𝓝}\ (S.corrector\ h\ hh\ x))}
\end{leancode}
\leansource{\texttt{corrector\_\allowbreak{}order} \quad WassersteinMixing/\allowbreak{}SemigroupPoisson.lean, lines 333–337}
\begin{leancode}
\leanline{0}{theorem\ corrector\_order\ (h\ :\ ℝ)\ (hh\ :\ 0\ ≤\ h)\ :}
\leanline{4}{0\ ≤\ S.corrector\ h\ hh\ ∧\ S.corrector\ h\ hh\ ≤\ h\ •\ (1\ :\ H\ →L[ℂ]\ H)}
\end{leancode}
\leanlabel{How the Lean proof follows the argument}
\texttt{partial\allowbreak{}Corrector} contains only finite operator sums. The successive-sum, positivity and telescoping lemmas supply all hypotheses of \texttt{exists\_\allowbreak{}monotone\_\allowbreak{}operator\_\allowbreak{}limit}. \texttt{corrector} chooses that proved operator, while \texttt{partial\allowbreak{}Corrector\_\allowbreak{}tendsto} and \texttt{corrector\_\allowbreak{}order} retain its actual convergence and order properties.
\end{leanbox}
\subsection{\texorpdfstring{The exact Poisson equation and Dirichlet estimate}{The exact Poisson equation and Dirichlet estimate}}
\label{report-step-WassersteinMixing-SemigroupPoisson-lean-4}


\textbf{Claim.} The constructed corrector commutes with the semigroup and satisfies

\[(I-P_h)G_hx=\int_0^h(x-P_tx)\,dt,\qquad \langle G_hx,(I-P_s)G_hx\rangle_{\mathbb R}\le h^2\langle x,(I-P_s)x\rangle_{\mathbb R}.\]

\textbf{Proof.} The finite sums obey \(G_{N+1}=D_h+P_hG_N\). Pass to their strong limit to obtain the Poisson identity. Each finite sum commutes with every transition, and bounded operator continuity passes this identity to the limit. The limit is self-adjoint and bounded by \(h\); apply the proved commuting energy inequality with the positive operator \(I-P_s\). No spectral gap or discrete spectral representation is used.

\begin{leanbox}{Lean: \texttt{partialCorrector\_\allowbreak{}recursion}, \texttt{corrector\_\allowbreak{}poisson}, \texttt{partialCorrector\_\allowbreak{}commute} and 2 more}
\leansource{\texttt{partialCorrector\_\allowbreak{}recursion} \quad WassersteinMixing/\allowbreak{}SemigroupPoisson.lean, lines 339–352}
\begin{leancode}
\leanline{0}{theorem\ partialCorrector\_recursion\ (h\ :\ ℝ)\ (hh\ :\ 0\ ≤\ h)\ (N\ :\ ℕ)\ :}
\leanline{4}{S.partialCorrector\ h\ hh\ (N\ +\ 1)\ =\ S.defectOp\ h\ hh\ +}
\leanline{6}{S.op\ h\ *\ S.partialCorrector\ h\ hh\ N}
\end{leancode}
\leansource{\texttt{corrector\_\allowbreak{}poisson} \quad WassersteinMixing/\allowbreak{}SemigroupPoisson.lean, lines 355–367}
\begin{leancode}
\leanline{0}{theorem\ corrector\_poisson\ (h\ :\ ℝ)\ (hh\ :\ 0\ ≤\ h)\ (x\ :\ H)\ :}
\leanline{4}{S.corrector\ h\ hh\ x\ -\ S.op\ h\ (S.corrector\ h\ hh\ x)\ =\ S.defect\ h\ x}
\end{leancode}
\leansource{\texttt{partialCorrector\_\allowbreak{}commute} \quad WassersteinMixing/\allowbreak{}SemigroupPoisson.lean, lines 369–375}
\begin{leancode}
\leanline{0}{theorem\ partialCorrector\_commute\ (h\ :\ ℝ)\ (hh\ :\ 0\ ≤\ h)\ \{s\ :\ ℝ\}\ (hs\ :\ 0\ ≤\ s)}
\leanline{4}{(N\ :\ ℕ)\ :\ Commute\ (S.op\ s)\ (S.partialCorrector\ h\ hh\ N)}
\end{leancode}
\leansource{\texttt{corrector\_\allowbreak{}commute} \quad WassersteinMixing/\allowbreak{}SemigroupPoisson.lean, lines 378–388}
\begin{leancode}
\leanline{0}{theorem\ corrector\_commute\ (h\ :\ ℝ)\ (hh\ :\ 0\ ≤\ h)\ \{s\ :\ ℝ\}\ (hs\ :\ 0\ ≤\ s)\ :}
\leanline{4}{Commute\ (S.op\ s)\ (S.corrector\ h\ hh)}
\end{leancode}
\leansource{\texttt{corrector\_\allowbreak{}energy\_\allowbreak{}le} \quad WassersteinMixing/\allowbreak{}SemigroupPoisson.lean, lines 392–402}
\begin{leancode}
\leanline{0}{theorem\ corrector\_energy\_le\ (h\ :\ ℝ)\ (hh\ :\ 0\ ≤\ h)\ \{s\ :\ ℝ\}\ (hs\ :\ 0\ ≤\ s)\ (x\ :\ H)\ :}
\leanline{4}{inner\ ℝ\ (S.corrector\ h\ hh\ x)}
\leanline{6}{(S.corrector\ h\ hh\ x\ -\ S.op\ s\ (S.corrector\ h\ hh\ x))\ ≤}
\leanline{6}{h\ \textasciicircum{}\ 2\ *\ inner\ ℝ\ x\ (x\ -\ S.op\ s\ x)}
\end{leancode}
\leanlabel{How the Lean proof follows the argument}
\texttt{partial\allowbreak{}Corrector\_\allowbreak{}recursion} gives the exact finite identity. \texttt{corrector\_\allowbreak{}poisson} passes to the actual strong limit. \texttt{corrector\_\allowbreak{}commute} transfers finite commutation, and \texttt{corrector\_\allowbreak{}energy\_\allowbreak{}le} specializes \texttt{commuting\_\allowbreak{}energy\_\allowbreak{}bound} to the semigroup energy operator.
\end{leanbox}
\section{\texorpdfstring{Transferring the constructed corrector to real observables}{Transferring the constructed corrector to real observables}}
\label{proof-module-WassersteinMixing-ComplexRealCorrector-lean}


An isometric real embedding \(i:R\to H\) into a complex Hilbert space has orthogonal projection \(ii^*\) onto its closed real range. If a semigroup preserves this range, its actual integrated defect, finite corrector sums and strong limit preserve it too. Extraction transfers the same Poisson identity and constants to the real space.

\subsection{\texorpdfstring{The genuine closed real subspace}{The genuine closed real subspace}}
\label{report-step-WassersteinMixing-ComplexRealCorrector-lean-1}


\textbf{Claim.} The real adjoint satisfies \(i^*i=I\), and \(\Pi=ii^*\) fixes exactly \(i(R)\). This image is closed.



\textbf{Proof.} The adjoint identity follows from preservation of the real inner product by an isometry. Hence \(\Pi(ix)=ix\). Conversely a fixed vector \(y\) equals \(i(i^*y)\), so it belongs to the image. The image is therefore the equalizer of the continuous maps \(\Pi\) and the identity, which is closed.

\begin{leanbox}{Lean: \texttt{instance\_\allowbreak{}at\_\allowbreak{}26}, \texttt{projection}, \texttt{adjoint\_\allowbreak{}embed} and 3 more}
\leansource{\texttt{instance\_\allowbreak{}at\_\allowbreak{}26} \quad WassersteinMixing/\allowbreak{}ComplexRealCorrector.lean, lines 26–26}
\begin{leancode}
\leanline{0}{local\ instance\ :\ InnerProductSpace\ ℝ\ H\ :=\ InnerProductSpace.rclikeToReal\ ℂ\ H}
\end{leancode}
\leansource{\texttt{projection} \quad WassersteinMixing/\allowbreak{}ComplexRealCorrector.lean, lines 29–30}
\begin{leancode}
\leanline{0}{def\ projection\ (i\ :\ R\ →ₗᵢ[ℝ]\ H)\ :\ H\ →L[ℝ]\ H\ :=}
\leanline{2}{i.toContinuousLinearMap\ ∘L\ i.toContinuousLinearMap.adjoint}
\end{leancode}
\leansource{\texttt{adjoint\_\allowbreak{}embed} \quad WassersteinMixing/\allowbreak{}ComplexRealCorrector.lean, lines 33–37}
\begin{leancode}
\leanline{0}{theorem\ adjoint\_embed\ (i\ :\ R\ →ₗᵢ[ℝ]\ H)\ (x\ :\ R)\ :}
\leanline{4}{i.toContinuousLinearMap.adjoint\ (i\ x)\ =\ x}
\end{leancode}
\leansource{\texttt{projection\_\allowbreak{}embed} \quad WassersteinMixing/\allowbreak{}ComplexRealCorrector.lean, lines 39–42}
\begin{leancode}
\leanline{0}{@[simp]\ theorem\ projection\_embed\ (i\ :\ R\ →ₗᵢ[ℝ]\ H)\ (x\ :\ R)\ :}
\leanline{4}{projection\ i\ (i\ x)\ =\ i\ x}
\end{leancode}
\leansource{\texttt{projection\_\allowbreak{}fixed\_\allowbreak{}iff} \quad WassersteinMixing/\allowbreak{}ComplexRealCorrector.lean, lines 45–51}
\begin{leancode}
\leanline{0}{theorem\ projection\_fixed\_iff\ (i\ :\ R\ →ₗᵢ[ℝ]\ H)\ (y\ :\ H)\ :}
\leanline{4}{projection\ i\ y\ =\ y\ ↔\ y\ ∈\ Set.range\ i}
\end{leancode}
\leansource{\texttt{isClosed\_\allowbreak{}range} \quad WassersteinMixing/\allowbreak{}ComplexRealCorrector.lean, lines 54–59}
\begin{leancode}
\leanline{0}{theorem\ isClosed\_range\ (i\ :\ R\ →ₗᵢ[ℝ]\ H)\ :\ IsClosed\ (Set.range\ i)}
\end{leancode}
\leanlabel{How the Lean proof follows the argument}
The local real inner product is the real part of the complex one. \texttt{projection} is the actual adjoint composition. \texttt{projection\_\allowbreak{}fixed\_\allowbreak{}iff} characterizes its fixed vectors, and \texttt{is\allowbreak{}Closed\_\allowbreak{}range} uses continuity of the two equalizer maps.
\end{leanbox}
\subsection{\texorpdfstring{Preservation passes through integrals, sums and limits}{Preservation passes through integrals, sums and limits}}
\label{report-step-WassersteinMixing-ComplexRealCorrector-lean-2}


\textbf{Claim.} If \(P_ti(R)\subseteq i(R)\), then \(D_hi(R)\subseteq i(R)\), \(G_Ni(R)\subseteq i(R)\) and \(G_hi(R)\subseteq i(R)\).



\textbf{Proof.} Apply the projection to \(ix-P_tix\); it fixes each integrand. Continuity of the projection lets it pass through the Bochner integral. Induction on the finite-sum recursion gives the same property for every \(G_N\). Taking the strong limit preserves the projection equality, so the constructed corrector also lies in the genuine closed real image.

\begin{leanbox}{Lean: \texttt{PreservesReal}, \texttt{op\_\allowbreak{}fixed}, \texttt{defect\_\allowbreak{}fixed} and 2 more}
\leansource{\texttt{PreservesReal} \quad WassersteinMixing/\allowbreak{}ComplexRealCorrector.lean, lines 63–64}
\begin{leancode}
\leanline{0}{def\ PreservesReal\ (S\ :\ ReversibleSemigroup\ H)\ (i\ :\ R\ →ₗᵢ[ℝ]\ H)\ :\ Prop\ :=}
\leanline{2}{∀\ t,\ 0\ ≤\ t\ →\ ∀\ x,\ projection\ i\ (S.op\ t\ (i\ x))\ =\ S.op\ t\ (i\ x)}
\end{leancode}
\leansource{\texttt{op\_\allowbreak{}fixed} \quad WassersteinMixing/\allowbreak{}ComplexRealCorrector.lean, lines 68–74}
\begin{leancode}
\leanline{0}{theorem\ op\_fixed\ (hreal\ :\ PreservesReal\ S\ i)\ \{t\ :\ ℝ\}\ (ht\ :\ 0\ ≤\ t)}
\leanline{4}{(y\ :\ H)\ (hy\ :\ projection\ i\ y\ =\ y)\ :}
\leanline{4}{projection\ i\ (S.op\ t\ y)\ =\ S.op\ t\ y}
\end{leancode}
\leansource{\texttt{defect\_\allowbreak{}fixed} \quad WassersteinMixing/\allowbreak{}ComplexRealCorrector.lean, lines 77–85}
\begin{leancode}
\leanline{0}{theorem\ defect\_fixed\ (hreal\ :\ PreservesReal\ S\ i)\ (h\ :\ ℝ)\ (hh\ :\ 0\ ≤\ h)\ (x\ :\ R)\ :}
\leanline{4}{projection\ i\ (S.defect\ h\ (i\ x))\ =\ S.defect\ h\ (i\ x)}
\end{leancode}
\leansource{\texttt{partialCorrector\_\allowbreak{}fixed} \quad WassersteinMixing/\allowbreak{}ComplexRealCorrector.lean, lines 88–97}
\begin{leancode}
\leanline{0}{theorem\ partialCorrector\_fixed\ (hreal\ :\ PreservesReal\ S\ i)\ (h\ :\ ℝ)\ (hh\ :\ 0\ ≤\ h)}
\leanline{4}{(N\ :\ ℕ)\ (x\ :\ R)\ :}
\leanline{4}{projection\ i\ (S.partialCorrector\ h\ hh\ N\ (i\ x))\ =\ S.partialCorrector\ h\ hh\ N\ (i\ x)}
\end{leancode}
\leansource{\texttt{corrector\_\allowbreak{}fixed} \quad WassersteinMixing/\allowbreak{}ComplexRealCorrector.lean, lines 100–109}
\begin{leancode}
\leanline{0}{theorem\ corrector\_fixed\ (hreal\ :\ PreservesReal\ S\ i)\ (h\ :\ ℝ)\ (hh\ :\ 0\ ≤\ h)\ (x\ :\ R)\ :}
\leanline{4}{projection\ i\ (S.corrector\ h\ hh\ (i\ x))\ =\ S.corrector\ h\ hh\ (i\ x)}
\end{leancode}
\leanlabel{How the Lean proof follows the argument}
\texttt{Preserves\allowbreak{}Real} expresses only a semigroup range property. \texttt{defect\_\allowbreak{}fixed} commutes the projection through the justified integral. \texttt{partial\allowbreak{}Corrector\_\allowbreak{}fixed} uses the finite recursion, and \texttt{corrector\_\allowbreak{}fixed} uses uniqueness of limits. The actual diffusion proves the range property separately.
\end{leanbox}
\subsection{\texorpdfstring{The exact real corrector and its unchanged constants}{The exact real corrector and its unchanged constants}}
\label{report-step-WassersteinMixing-ComplexRealCorrector-lean-3}


\textbf{Construction.} Define \(P_t^{\mathbb R}=i^*P_ti\) and \(G_h^{\mathbb R}=i^*G_hi\). Range preservation gives \(iP_t^{\mathbb R}=P_ti\) and \(iG_h^{\mathbb R}=G_hi\). Consequently

\[\|G_h^{\mathbb R}\|\le h,\qquad (I-P_h^{\mathbb R})G_h^{\mathbb R}x=\int_0^h(x-P_t^{\mathbb R}x)\,dt.\]

\textbf{Proof.} Apply the isometry to the extracted operators and use the fixed projection identities. Norm preservation transfers the bound without a dimension loss. Adjoint extraction commutes with the integral. Injectivity of the embedding transfers the Poisson equation; preservation of its real inner product transfers the full Dirichlet energy estimate with factor \(h^2\).

\begin{leanbox}{Lean: \texttt{realOp}, \texttt{realCorrector}, \texttt{embed\_\allowbreak{}realOp} and 8 more}
\leansource{\texttt{realOp} \quad WassersteinMixing/\allowbreak{}ComplexRealCorrector.lean, lines 112–113}
\begin{leancode}
\leanline{0}{def\ realOp\ (t\ :\ ℝ)\ :\ R\ →L[ℝ]\ R\ :=}
\leanline{2}{i.toContinuousLinearMap.adjoint\ ∘L\ (S.op\ t).restrictScalars\ ℝ\ ∘L\ i.toContinuousLinearMap}
\end{leancode}
\leansource{\texttt{realCorrector} \quad WassersteinMixing/\allowbreak{}ComplexRealCorrector.lean, lines 116–118}
\begin{leancode}
\leanline{0}{def\ realCorrector\ (h\ :\ ℝ)\ (hh\ :\ 0\ ≤\ h)\ :\ R\ →L[ℝ]\ R\ :=}
\leanline{2}{i.toContinuousLinearMap.adjoint\ ∘L\ (S.corrector\ h\ hh).restrictScalars\ ℝ\ ∘L}
\leanline{4}{i.toContinuousLinearMap}
\end{leancode}
\leansource{\texttt{embed\_\allowbreak{}realOp} \quad WassersteinMixing/\allowbreak{}ComplexRealCorrector.lean, lines 121–122}
\begin{leancode}
\leanline{0}{theorem\ embed\_realOp\ (hreal\ :\ PreservesReal\ S\ i)\ \{t\ :\ ℝ\}\ (ht\ :\ 0\ ≤\ t)\ (x\ :\ R)\ :}
\leanline{4}{i\ (realOp\ S\ i\ t\ x)\ =\ S.op\ t\ (i\ x)}
\end{leancode}
\leansource{\texttt{embed\_\allowbreak{}realCorrector} \quad WassersteinMixing/\allowbreak{}ComplexRealCorrector.lean, lines 125–127}
\begin{leancode}
\leanline{0}{theorem\ embed\_realCorrector\ (hreal\ :\ PreservesReal\ S\ i)\ (h\ :\ ℝ)\ (hh\ :\ 0\ ≤\ h)\ (x\ :\ R)\ :}
\leanline{4}{i\ (realCorrector\ S\ i\ h\ hh\ x)\ =\ S.corrector\ h\ hh\ (i\ x)}
\end{leancode}
\leansource{\texttt{realCorrector\_\allowbreak{}norm\_\allowbreak{}le} \quad WassersteinMixing/\allowbreak{}ComplexRealCorrector.lean, lines 130–136}
\begin{leancode}
\leanline{0}{theorem\ realCorrector\_norm\_le\ (hreal\ :\ PreservesReal\ S\ i)\ (h\ :\ ℝ)\ (hh\ :\ 0\ ≤\ h)\ (x\ :\ R)\ :}
\leanline{4}{‖realCorrector\ S\ i\ h\ hh\ x‖\ ≤\ h\ *\ ‖x‖}
\end{leancode}
\leansource{\texttt{realCorrector\_\allowbreak{}opNorm\_\allowbreak{}le} \quad WassersteinMixing/\allowbreak{}ComplexRealCorrector.lean, lines 139–142}
\begin{leancode}
\leanline{0}{theorem\ realCorrector\_opNorm\_le\ (hreal\ :\ PreservesReal\ S\ i)\ (h\ :\ ℝ)\ (hh\ :\ 0\ ≤\ h)\ :}
\leanline{4}{‖realCorrector\ S\ i\ h\ hh‖\ ≤\ h}
\end{leancode}
\leansource{\texttt{realDefect} \quad WassersteinMixing/\allowbreak{}ComplexRealCorrector.lean, lines 145–146}
\begin{leancode}
\leanline{0}{def\ realDefect\ (h\ :\ ℝ)\ (x\ :\ R)\ :\ R\ :=}
\leanline{2}{i.toContinuousLinearMap.adjoint\ (S.defect\ h\ (i\ x))}
\end{leancode}
\leansource{\texttt{realDefect\_\allowbreak{}eq\_\allowbreak{}integral} \quad WassersteinMixing/\allowbreak{}ComplexRealCorrector.lean, lines 149–157}
\begin{leancode}
\leanline{0}{theorem\ realDefect\_eq\_integral\ (h\ :\ ℝ)\ (hh\ :\ 0\ ≤\ h)\ (x\ :\ R)\ :}
\leanline{4}{realDefect\ S\ i\ h\ x\ =\ ∫\ t\ in\ 0..h,\ x\ -\ realOp\ S\ i\ t\ x}
\end{leancode}
\leansource{\texttt{embed\_\allowbreak{}realDefect} \quad WassersteinMixing/\allowbreak{}ComplexRealCorrector.lean, lines 160–161}
\begin{leancode}
\leanline{0}{theorem\ embed\_realDefect\ (hreal\ :\ PreservesReal\ S\ i)\ (h\ :\ ℝ)\ (hh\ :\ 0\ ≤\ h)\ (x\ :\ R)\ :}
\leanline{4}{i\ (realDefect\ S\ i\ h\ x)\ =\ S.defect\ h\ (i\ x)}
\end{leancode}
\leansource{\texttt{realCorrector\_\allowbreak{}poisson} \quad WassersteinMixing/\allowbreak{}ComplexRealCorrector.lean, lines 164–172}
\begin{leancode}
\leanline{0}{theorem\ realCorrector\_poisson\ (hreal\ :\ PreservesReal\ S\ i)\ (h\ :\ ℝ)\ (hh\ :\ 0\ ≤\ h)\ (x\ :\ R)\ :}
\leanline{4}{realCorrector\ S\ i\ h\ hh\ x\ -\ realOp\ S\ i\ h\ (realCorrector\ S\ i\ h\ hh\ x)\ =}
\leanline{6}{∫\ t\ in\ 0..h,\ x\ -\ realOp\ S\ i\ t\ x}
\end{leancode}
\leansource{\texttt{realCorrector\_\allowbreak{}energy\_\allowbreak{}le} \quad WassersteinMixing/\allowbreak{}ComplexRealCorrector.lean, lines 175–185}
\begin{leancode}
\leanline{0}{theorem\ realCorrector\_energy\_le\ (hreal\ :\ PreservesReal\ S\ i)\ (h\ :\ ℝ)\ (hh\ :\ 0\ ≤\ h)}
\leanline{4}{\{s\ :\ ℝ\}\ (hs\ :\ 0\ ≤\ s)\ (x\ :\ R)\ :}
\leanline{4}{inner\ ℝ\ (realCorrector\ S\ i\ h\ hh\ x)}
\leanline{6}{(realCorrector\ S\ i\ h\ hh\ x\ -\ realOp\ S\ i\ s\ (realCorrector\ S\ i\ h\ hh\ x))\ ≤}
\leanline{6}{h\ \textasciicircum{}\ 2\ *\ inner\ ℝ\ x\ (x\ -\ realOp\ S\ i\ s\ x)}
\end{leancode}
\leanlabel{How the Lean proof follows the argument}
\texttt{real\allowbreak{}Op}, \texttt{real\allowbreak{}Corrector} and \texttt{real\allowbreak{}Defect} are explicit adjoint extractions of constructed objects. The embedding lemmas identify them exactly. \texttt{real\allowbreak{}Corrector\_\allowbreak{}poisson} and \texttt{real\allowbreak{}Corrector\_\allowbreak{}energy\_\allowbreak{}le} transfer the already proved complex identities through the isometry.
\end{leanbox}
\subsection{\texorpdfstring{Actual coordinate embeddings on square-integrable observables}{Actual coordinate embeddings on square-integrable observables}}
\label{report-step-WassersteinMixing-ComplexRealCorrector-lean-4}


\textbf{Construction.} A pointwise real linear isometry induces an isometry on \(L^2\) by composition. For a finite-dimensional real Hilbert space, use an orthonormal basis, take its real coordinates and include them in complex coordinate space.



\textbf{Proof.} The pointwise norm equality gives equality of the integral square norms, so the lifted map is a genuine isometry of equivalence classes. The coordinate representation is an isometry by the orthonormal basis theorem; real-to-complex inclusion preserves the sum of squared coordinate norms. Composing these maps and then lifting to \(L^2\) preserves every observable norm exactly.

\begin{leanbox}{Lean: \texttt{mapLp}, \texttt{mapLp\_\allowbreak{}coeFn}, \texttt{euclideanOfReal} and 3 more}
\leansource{\texttt{mapLp} \quad WassersteinMixing/\allowbreak{}ComplexRealCorrector.lean, lines 197–205}
\begin{leancode}
\leanline{0}{def\ mapLp\ (i\ :\ E\ →ₗᵢ[ℝ]\ F)\ (μ\ :\ Measure\ A)\ :\ Lp\ E\ 2\ μ\ →ₗᵢ[ℝ]\ Lp\ F\ 2\ μ\ where}
\leanline{2}{toLinearMap\ :=\ i.toContinuousLinearMap.compLpₗ\ 2\ μ}
\leanline{2}{norm\_map'\ f\ :=\ by}
\leanline{4}{change\ ‖i.toContinuousLinearMap.compLp\ f‖\ =\ ‖f‖}
\leanline{4}{rw\ [Lp.norm\_def,\ Lp.norm\_def]}
\leanline{4}{congr\ 1}
\leanline{4}{apply\ eLpNorm\_congr\_norm\_ae\ (Lp.aestronglyMeasurable\ \_)\ (Lp.aestronglyMeasurable\ \_)}
\leanline{4}{filter\_upwards\ [i.toContinuousLinearMap.coeFn\_compLp\ f]\ with\ a\ ha}
\leanline{4}{rw\ [ha,\ LinearIsometry.coe\_toContinuousLinearMap,\ i.norm\_map]}
\end{leancode}
\leansource{\texttt{mapLp\_\allowbreak{}coeFn} \quad WassersteinMixing/\allowbreak{}ComplexRealCorrector.lean, lines 208–212}
\begin{leancode}
\leanline{0}{theorem\ mapLp\_coeFn\ (i\ :\ E\ →ₗᵢ[ℝ]\ F)\ (μ\ :\ Measure\ A)\ (f\ :\ Lp\ E\ 2\ μ)\ :}
\leanline{4}{mapLp\ i\ μ\ f\ =ᵐ[μ]\ fun\ a\ ↦\ i\ (f\ a)}
\end{leancode}
\leansource{\texttt{euclideanOfReal} \quad WassersteinMixing/\allowbreak{}ComplexRealCorrector.lean, lines 215–229}
\begin{leancode}
\leanline{0}{def\ euclideanOfReal\ (ι\ :\ Type*)\ [Fintype\ ι]\ :}
\leanline{4}{EuclideanSpace\ ℝ\ ι\ →ₗᵢ[ℝ]\ EuclideanSpace\ ℂ\ ι\ where}
\leanline{2}{toFun\ x\ :=\ WithLp.toLp\ 2\ (fun\ j\ ↦\ (x\ j\ :\ ℂ))}
\leanline{2}{map\_add'\ x\ y\ :=\ by}
\leanline{4}{ext\ j}
\leanline{4}{simp}
\leanline{2}{map\_smul'\ c\ x\ :=\ by}
\leanline{4}{ext\ j}
\leanline{4}{simp}
\leanline{2}{norm\_map'\ x\ :=\ by}
\leanline{4}{rw\ [EuclideanSpace.norm\_eq,\ EuclideanSpace.norm\_eq]}
\leanline{4}{congr\ 1}
\leanline{4}{apply\ Finset.sum\_congr\ rfl}
\leanline{4}{intro\ j\ \_}
\leanline{4}{simp}
\end{leancode}
\leansource{\texttt{coordinate} \quad WassersteinMixing/\allowbreak{}ComplexRealCorrector.lean, lines 233–237}
\begin{leancode}
\leanline{0}{def\ coordinate\ (R\ :\ Type*)\ [NormedAddCommGroup\ R]\ [InnerProductSpace\ ℝ\ R]}
\leanline{4}{[FiniteDimensional\ ℝ\ R]\ :}
\leanline{4}{R\ →ₗᵢ[ℝ]\ EuclideanSpace\ ℂ\ (Fin\ (Module.finrank\ ℝ\ R))\ :=}
\leanline{2}{(euclideanOfReal\ (Fin\ (Module.finrank\ ℝ\ R))).comp}
\leanline{4}{(stdOrthonormalBasis\ ℝ\ R).repr.toLinearIsometry}
\end{leancode}
\leansource{\texttt{coordinateLp} \quad WassersteinMixing/\allowbreak{}ComplexRealCorrector.lean, lines 240–243}
\begin{leancode}
\leanline{0}{def\ coordinateLp\ (R\ :\ Type*)\ [NormedAddCommGroup\ R]\ [InnerProductSpace\ ℝ\ R]}
\leanline{4}{[FiniteDimensional\ ℝ\ R]\ (μ\ :\ Measure\ A)\ :}
\leanline{4}{Lp\ R\ 2\ μ\ →ₗᵢ[ℝ]\ Lp\ (EuclideanSpace\ ℂ\ (Fin\ (Module.finrank\ ℝ\ R)))\ 2\ μ\ :=}
\leanline{2}{mapLp\ (coordinate\ R)\ μ}
\end{leancode}
\leansource{\texttt{coordinateLp\_\allowbreak{}norm} \quad WassersteinMixing/\allowbreak{}ComplexRealCorrector.lean, lines 247–249}
\begin{leancode}
\leanline{0}{theorem\ coordinateLp\_norm\ (R\ :\ Type*)\ [NormedAddCommGroup\ R]\ [InnerProductSpace\ ℝ\ R]}
\leanline{4}{[FiniteDimensional\ ℝ\ R]\ (μ\ :\ Measure\ A)\ (f\ :\ Lp\ R\ 2\ μ)\ :}
\leanline{4}{‖coordinateLp\ R\ μ\ f‖\ =\ ‖f‖}
\end{leancode}
\leanlabel{How the Lean proof follows the argument}
\texttt{map\allowbreak{}Lp} uses actual continuous linear pointwise composition and proves its norm equality by almost-everywhere representative comparison. \texttt{euclidean\allowbreak{}Of\allowbreak{}Real}, \texttt{coordinate} and \texttt{coordinate\allowbreak{}Lp} construct the explicit finite-dimensional embedding; \texttt{coordinate\allowbreak{}Lp\_\allowbreak{}norm} records its exact norm equality.
\end{leanbox}
\section{\texorpdfstring{Real observables and the adjoints of actual evaluations}{Real observables and the adjoints of actual evaluations}}
\label{proof-module-WassersteinMixing-ComplexRealTransition-lean}


For an actual measure-preserving observation, pointwise real embedding and pointwise adjoint extraction commute with evaluation. Taking adjoints therefore proves that the complex transition compression of an embedded real observable is exactly the embedding of its real transition.

\subsection{\texorpdfstring{Scalar restriction and pointwise adjoint extraction}{Scalar restriction and pointwise adjoint extraction}}
\label{report-step-WassersteinMixing-ComplexRealTransition-lean-1}


\textbf{Claim.} Complex adjunction agrees with real adjunction after restricting scalars. On \(L^2\), the real inner product is the integral of the pointwise real inner product and equals the real part of the complex pairing. The adjoint of the lifted embedding is the lift of its pointwise adjoint.



\textbf{Proof.} Test scalar-restricted adjoints against arbitrary vectors and take the real part of the complex adjoint identity. For square-integrable functions, the same identity follows by integrating their pointwise pairings. The real-part functional commutes with the integrable complex pairing; the defining pointwise adjoint equality then identifies the lifted adjoint.

\begin{leanbox}{Lean: \texttt{instance\_\allowbreak{}at\_\allowbreak{}23}, \texttt{instance\_\allowbreak{}at\_\allowbreak{}24}, \texttt{adjoint\_\allowbreak{}restrictScalars} and 4 more}
\leansource{\texttt{instance\_\allowbreak{}at\_\allowbreak{}23} \quad WassersteinMixing/\allowbreak{}ComplexRealTransition.lean, lines 23–23}
\begin{leancode}
\leanline{0}{local\ instance\ :\ InnerProductSpace\ ℝ\ H\ :=\ InnerProductSpace.rclikeToReal\ ℂ\ H}
\end{leancode}
\leansource{\texttt{instance\_\allowbreak{}at\_\allowbreak{}24} \quad WassersteinMixing/\allowbreak{}ComplexRealTransition.lean, lines 24–24}
\begin{leancode}
\leanline{0}{local\ instance\ :\ InnerProductSpace\ ℝ\ K\ :=\ InnerProductSpace.rclikeToReal\ ℂ\ K}
\end{leancode}
\leansource{\texttt{adjoint\_\allowbreak{}restrictScalars} \quad WassersteinMixing/\allowbreak{}ComplexRealTransition.lean, lines 27–34}
\begin{leancode}
\leanline{0}{theorem\ adjoint\_restrictScalars\ (T\ :\ H\ →L[ℂ]\ K)\ :}
\leanline{4}{(T.restrictScalars\ ℝ).adjoint\ =\ T.adjoint.restrictScalars\ ℝ}
\end{leancode}
\leansource{\texttt{instance\_\allowbreak{}at\_\allowbreak{}42} \quad WassersteinMixing/\allowbreak{}ComplexRealTransition.lean, lines 42–42}
\begin{leancode}
\leanline{0}{local\ instance\ :\ InnerProductSpace\ ℝ\ H\ :=\ InnerProductSpace.rclikeToReal\ ℂ\ H}
\end{leancode}
\leansource{\texttt{real\_\allowbreak{}inner\_\allowbreak{}Lp} \quad WassersteinMixing/\allowbreak{}ComplexRealTransition.lean, lines 46–48}
\begin{leancode}
\leanline{0}{theorem\ real\_inner\_Lp\ (μ\ :\ Measure\ A)\ (f\ g\ :\ Lp\ H\ 2\ μ)\ :}
\leanline{4}{inner\ ℝ\ f\ g\ =\ ∫\ a,\ inner\ ℝ\ (f\ a)\ (g\ a)\ ∂μ}
\end{leancode}
\leansource{\texttt{real\_\allowbreak{}inner\_\allowbreak{}Lp\_\allowbreak{}eq\_\allowbreak{}re} \quad WassersteinMixing/\allowbreak{}ComplexRealTransition.lean, lines 52–55}
\begin{leancode}
\leanline{0}{theorem\ real\_inner\_Lp\_eq\_re\ (μ\ :\ Measure\ A)\ (f\ g\ :\ Lp\ H\ 2\ μ)\ :}
\leanline{4}{inner\ ℝ\ f\ g\ =\ (inner\ ℂ\ f\ g).re}
\end{leancode}
\leansource{\texttt{mapLp\_\allowbreak{}adjoint} \quad WassersteinMixing/\allowbreak{}ComplexRealTransition.lean, lines 58–77}
\begin{leancode}
\leanline{0}{theorem\ mapLp\_adjoint\ (i\ :\ R\ →ₗᵢ[ℝ]\ H)\ (μ\ :\ Measure\ A)\ :}
\leanline{4}{(mapLp\ i\ μ).toContinuousLinearMap.adjoint\ =}
\leanline{6}{i.toContinuousLinearMap.adjoint.compLpL\ 2\ μ}
\end{leancode}
\leanlabel{How the Lean proof follows the argument}
The local scalar restrictions specify the underlying real inner products. \texttt{adjoint\_\allowbreak{}restrict\allowbreak{}Scalars} uses uniqueness of adjoints. \texttt{real\_\allowbreak{}inner\_\allowbreak{}Lp\_\allowbreak{}eq\_\allowbreak{}re} uses the integral of the real part. \texttt{map\allowbreak{}Lp\_\allowbreak{}adjoint} tests against arbitrary real observables and compares exact almost-everywhere representatives.
\end{leanbox}
\subsection{\texorpdfstring{Naturality of actual evaluation and its adjoint}{Naturality of actual evaluation and its adjoint}}
\label{report-step-WassersteinMixing-ComplexRealTransition-lean-2}


\textbf{Claim.} If \(\phi_*\nu=\mu\), then pointwise embedding and extraction commute with pullback by \(\phi\). The adjoint pullback satisfies the corresponding embedding identity.



\textbf{Proof.} Both sides have the same almost-everywhere representative, obtained by composing the pointwise map with \(\phi\). A measure-preserving map transports the exceptional null sets correctly. Express extraction as the proved lifted adjoint and take adjoints of its commuting diagram. Scalar-restriction compatibility and the double-adjoint identity give the asserted adjoint-evaluation diagram.

\begin{leanbox}{Lean: \texttt{extraction\_\allowbreak{}compMeasurePreserving}, \texttt{mapLp\_\allowbreak{}compMeasurePreserving}, \texttt{evaluation\_\allowbreak{}adjoint\_\allowbreak{}mapLp}}
\leansource{\texttt{extraction\_\allowbreak{}compMeasurePreserving} \quad WassersteinMixing/\allowbreak{}ComplexRealTransition.lean, lines 80–94}
\begin{leancode}
\leanline{0}{theorem\ extraction\_compMeasurePreserving\ (i\ :\ R\ →ₗᵢ[ℝ]\ H)}
\leanline{4}{(μ\ :\ Measure\ A)\ (ν\ :\ Measure\ B)\ (φ\ :\ B\ →\ A)\ (hφ\ :\ MeasurePreserving\ φ\ ν\ μ)}
\leanline{4}{(f\ :\ Lp\ H\ 2\ μ)\ :}
\leanline{4}{i.toContinuousLinearMap.adjoint.compLpL\ 2\ ν\ (Lp.compMeasurePreserving\ φ\ hφ\ f)\ =}
\leanline{6}{Lp.compMeasurePreserving\ φ\ hφ\ (i.toContinuousLinearMap.adjoint.compLpL\ 2\ μ\ f)}
\end{leancode}
\leansource{\texttt{mapLp\_\allowbreak{}compMeasurePreserving} \quad WassersteinMixing/\allowbreak{}ComplexRealTransition.lean, lines 98–109}
\begin{leancode}
\leanline{0}{theorem\ mapLp\_compMeasurePreserving\ (i\ :\ R\ →ₗᵢ[ℝ]\ H)}
\leanline{4}{(μ\ :\ Measure\ A)\ (ν\ :\ Measure\ B)\ (φ\ :\ B\ →\ A)\ (hφ\ :\ MeasurePreserving\ φ\ ν\ μ)}
\leanline{4}{(f\ :\ Lp\ R\ 2\ μ)\ :}
\leanline{4}{mapLp\ i\ ν\ (Lp.compMeasurePreserving\ φ\ hφ\ f)\ =}
\leanline{6}{Lp.compMeasurePreserving\ φ\ hφ\ (mapLp\ i\ μ\ f)}
\end{leancode}
\leansource{\texttt{evaluation\_\allowbreak{}adjoint\_\allowbreak{}mapLp} \quad WassersteinMixing/\allowbreak{}ComplexRealTransition.lean, lines 112–131}
\begin{leancode}
\leanline{0}{theorem\ evaluation\_adjoint\_mapLp\ (i\ :\ R\ →ₗᵢ[ℝ]\ H)}
\leanline{4}{(μ\ :\ Measure\ A)\ (ν\ :\ Measure\ B)\ (φ\ :\ B\ →\ A)\ (hφ\ :\ MeasurePreserving\ φ\ ν\ μ)}
\leanline{4}{(f\ :\ Lp\ R\ 2\ ν)\ :}
\leanline{4}{(Lp.compMeasurePreservingₗᵢ\ ℂ\ φ\ hφ).toContinuousLinearMap.adjoint\ (mapLp\ i\ ν\ f)\ =}
\leanline{6}{mapLp\ i\ μ\ ((Lp.compMeasurePreservingₗᵢ\ ℝ\ φ\ hφ).toContinuousLinearMap.adjoint\ f)}
\end{leancode}
\leanlabel{How the Lean proof follows the argument}
The two composition theorems use \texttt{Lp.\allowbreak{}ext} and transport representative equalities through the measure-preserving map. \texttt{evaluation\_\allowbreak{}adjoint\_\allowbreak{}map\allowbreak{}Lp} takes adjoints of the exact extraction diagram and applies \texttt{adjoint\_\allowbreak{}restrict\allowbreak{}Scalars}.
\end{leanbox}
\subsection{\texorpdfstring{Actual transition compression preserves the real image}{Actual transition compression preserves the real image}}
\label{report-step-WassersteinMixing-ComplexRealTransition-lean-3}


\textbf{Claim.} Let \(U_0,U_1\) be evaluation isometries of two observations with the same marginal. Their complex compression \(U_0^*U_1\) sends the embedded real observable \(if\) to \(i(U_{0,\mathbb R}^*U_{1,\mathbb R}f)\), hence preserves the actual real image.



\textbf{Proof.} Move embedding through the second evaluation using pointwise naturality, then through the first adjoint using the previous theorem. The resulting vector is explicitly the embedding of the genuine real compression, which supplies its witness in the real range.

\begin{leanbox}{Lean: \texttt{compression\_\allowbreak{}mapLp}, \texttt{compression\_\allowbreak{}mem\_\allowbreak{}range}}
\leansource{\texttt{compression\_\allowbreak{}mapLp} \quad WassersteinMixing/\allowbreak{}ComplexRealTransition.lean, lines 135–146}
\begin{leancode}
\leanline{0}{theorem\ compression\_mapLp\ (i\ :\ R\ →ₗᵢ[ℝ]\ H)}
\leanline{4}{(μ\ :\ Measure\ A)\ (ν\ :\ Measure\ B)\ (φ\ ψ\ :\ B\ →\ A)}
\leanline{4}{(hφ\ :\ MeasurePreserving\ φ\ ν\ μ)\ (hψ\ :\ MeasurePreserving\ ψ\ ν\ μ)}
\leanline{4}{(f\ :\ Lp\ R\ 2\ μ)\ :}
\leanline{4}{ObservableTransition.compression\ (Lp.compMeasurePreservingₗᵢ\ ℂ\ φ\ hφ)}
\leanline{6}{(Lp.compMeasurePreservingₗᵢ\ ℂ\ ψ\ hψ)\ (mapLp\ i\ μ\ f)\ =}
\leanline{6}{mapLp\ i\ μ\ (ObservableTransition.compression\ (Lp.compMeasurePreservingₗᵢ\ ℝ\ φ\ hφ)}
\leanline{8}{(Lp.compMeasurePreservingₗᵢ\ ℝ\ ψ\ hψ)\ f)}
\end{leancode}
\leansource{\texttt{compression\_\allowbreak{}mem\_\allowbreak{}range} \quad WassersteinMixing/\allowbreak{}ComplexRealTransition.lean, lines 149–155}
\begin{leancode}
\leanline{0}{theorem\ compression\_mem\_range\ (i\ :\ R\ →ₗᵢ[ℝ]\ H)}
\leanline{4}{(μ\ :\ Measure\ A)\ (ν\ :\ Measure\ B)\ (φ\ ψ\ :\ B\ →\ A)}
\leanline{4}{(hφ\ :\ MeasurePreserving\ φ\ ν\ μ)\ (hψ\ :\ MeasurePreserving\ ψ\ ν\ μ)}
\leanline{4}{(f\ :\ Lp\ R\ 2\ μ)\ :}
\leanline{4}{ObservableTransition.compression\ (Lp.compMeasurePreservingₗᵢ\ ℂ\ φ\ hφ)}
\leanline{6}{(Lp.compMeasurePreservingₗᵢ\ ℂ\ ψ\ hψ)\ (mapLp\ i\ μ\ f)\ ∈\ Set.range\ (mapLp\ i\ μ)}
\end{leancode}
\leanlabel{How the Lean proof follows the argument}
\texttt{compression\_\allowbreak{}map\allowbreak{}Lp} composes the evaluation and adjoint-evaluation identities. \texttt{compression\_\allowbreak{}mem\_\allowbreak{}range} exhibits the actual real compression as the preimage, so range preservation is proved from the observation maps rather than introduced as an assumption.
\end{leanbox}
\section{\texorpdfstring{Real scalar restriction of the actual complex transition}{Real scalar restriction of the actual complex transition}}
\label{proof-module-WassersteinMixing-SemigroupScalar-lean}


The real and complex transition compressions of the same actual observations agree after restricting scalars. This identifies their genuine operators and transfers the real conditional Markov composition theorem to the complex observable Hilbert space.

\subsection{\texorpdfstring{The real part of the actual compression pairing}{The real part of the actual compression pairing}}
\label{report-step-WassersteinMixing-SemigroupScalar-lean-1}


\textbf{Claim.} The actual complex operator, viewed as a real linear map, equals the actual real operator on the same complex-valued observable space.



\textbf{Proof.} Test the two outputs against any observable using the real inner product. The real \(L^2\) pairing equals the real part of its complex pairing. Apply the actual complex compression identity, then return to the real pairing of the actual evaluations. Their pointwise pullbacks are identical over both scalar fields. The actual real compression identity now gives exactly the second operator pairing. Separation by all real inner-product tests proves equality.

\begin{leanbox}{Lean: \texttt{op\_\allowbreak{}restrictScalars}}
\leansource{\texttt{op\_\allowbreak{}restrictScalars} \quad WassersteinMixing/\allowbreak{}SemigroupScalar.lean, lines 24–45}
\begin{leancode}
\leanline{0}{theorem\ op\_restrictScalars\ (P\ :\ Parameters)\ (V\ :\ E\ P.dimension\ →\ ℝ)}
\leanline{4}{(hV\ :\ PotentialAssumptions\ P\ V)\ (t\ :\ ℝ)\ :}
\leanline{4}{letI\ :\ InnerProductSpace\ ℝ\ F\ :=\ InnerProductSpace.rclikeToReal\ ℂ\ F}
\leanline{4}{(op\ ({\leanfallback 𝕜}\ :=\ ℂ)\ (F\ :=\ F)\ P\ V\ hV\ t).restrictScalars\ ℝ\ =}
\leanline{6}{op\ ({\leanfallback 𝕜}\ :=\ ℝ)\ (F\ :=\ F)\ P\ V\ hV\ t}
\end{leancode}
\leanlabel{How the Lean proof follows the argument}
\texttt{op\_\allowbreak{}restrict\allowbreak{}Scalars} uses \texttt{real\_\allowbreak{}inner\_\allowbreak{}Lp\_\allowbreak{}eq\_\allowbreak{}re} and both scalar versions of \texttt{compression\_\allowbreak{}inner}. The proof compares the actual evaluation maps; it does not postulate equality of the two transitions.
\end{leanbox}
\section{\texorpdfstring{Uniform randomization realizes exact Metropolis weights}{Uniform randomization realizes exact Metropolis weights}}
\label{proof-module-WassersteinMixing-UniformDecision-lean}


An independent uniform variable on \([0,1]\) realizes acceptance probability \(a(x)\) exactly. This connects a concrete random Metropolis update with the weighted accepted and rejected measures used in detailed balance.

\subsection{\texorpdfstring{The actual acceptance decision and its law}{The actual acceptance decision and its law}}
\label{report-step-WassersteinMixing-UniformDecision-lean-1}


\textbf{Claim.} Let \(0\le a\le1\), and let \(U\) be independent and uniform on \([0,1]\). Output \(f(x)\) when \(U\le a(x)\), and \(g(x)\) otherwise. Its law under \(\mu\otimes\operatorname{Unif}[0,1]\) is

\[f_*(a\mu)+g_*((1-a)\mu).\]

\textbf{Proof.} The threshold event is measurable. For a target event, integrate first over the uniform coordinate. Depending on whether each output lies in the event, its section is the entire interval, an initial interval of length \(a(x)\), a terminal interval of length \(1-a(x)\), or empty. These four masses give precisely the displayed sum after integrating over \(x\).

\begin{leanbox}{Lean: \texttt{uniformDecision}, \texttt{uniformDecision\_\allowbreak{}measurable}, \texttt{uniformDecision\_\allowbreak{}law}}
\leansource{\texttt{uniformDecision} \quad WassersteinMixing/\allowbreak{}UniformDecision.lean, lines 23–24}
\begin{leancode}
\leanline{0}{def\ uniformDecision\ (a\ :\ A\ →\ ℝ)\ (f\ g\ :\ A\ →\ B)\ (p\ :\ A\ ×\ I)\ :\ B\ :=}
\leanline{2}{if\ (p.2\ :\ ℝ)\ ≤\ a\ p.1\ then\ f\ p.1\ else\ g\ p.1}
\end{leancode}
\leansource{\texttt{uniformDecision\_\allowbreak{}measurable} \quad WassersteinMixing/\allowbreak{}UniformDecision.lean, lines 26–31}
\begin{leancode}
\leanline{0}{theorem\ uniformDecision\_measurable\ (a\ :\ A\ →\ ℝ)\ (f\ g\ :\ A\ →\ B)}
\leanline{4}{(ha\ :\ Measurable\ a)\ (hf\ :\ Measurable\ f)\ (hg\ :\ Measurable\ g)\ :}
\leanline{4}{Measurable\ (uniformDecision\ a\ f\ g)}
\end{leancode}
\leansource{\texttt{uniformDecision\_\allowbreak{}law} \quad WassersteinMixing/\allowbreak{}UniformDecision.lean, lines 34–79}
\begin{leancode}
\leanline{0}{theorem\ uniformDecision\_law\ (μ\ :\ Measure\ A)\ [SFinite\ μ]}
\leanline{4}{(a\ :\ A\ →\ ℝ)\ (f\ g\ :\ A\ →\ B)}
\leanline{4}{(ha\ :\ Measurable\ a)\ (ha0\ :\ ∀\ x,\ 0\ ≤\ a\ x)\ (ha1\ :\ ∀\ x,\ a\ x\ ≤\ 1)}
\leanline{4}{(hf\ :\ Measurable\ f)\ (hg\ :\ Measurable\ g)\ :}
\leanline{4}{(μ.prod\ (volume\ :\ Measure\ I)).map\ (uniformDecision\ a\ f\ g)\ =}
\leanline{6}{(μ.withDensity\ (fun\ x\ ↦\ ENNReal.ofReal\ (a\ x))).map\ f\ +}
\leanline{8}{(μ.withDensity\ (fun\ x\ ↦\ ENNReal.ofReal\ (1\ -\ a\ x))).map\ g}
\end{leancode}
\leanlabel{How the Lean proof follows the argument}
\texttt{uniform\allowbreak{}Decision} is the actual threshold conditional. \texttt{uniform\allowbreak{}Decision\_\allowbreak{}measurable} checks its measurable branches and decision event. \texttt{uniform\allowbreak{}Decision\_\allowbreak{}law} proves equality on measurable events using the exact unit-interval masses and the weighted-measure formulas.
\end{leanbox}
\section{\texorpdfstring{Actual paired independent innovation paths}{Actual paired independent innovation paths}}
\label{proof-module-WassersteinMixing-ProductNoise-lean}


Two independent paths of independent identically distributed innovations can be zipped coordinatewise into a single path of independent pairs. The exact identity persists when an independent initial state is present.

\subsection{\texorpdfstring{Zipping the two actual infinite product laws}{Zipping the two actual infinite product laws}}
\label{report-step-WassersteinMixing-ProductNoise-lean-1}


\textbf{Claim.} If \((\xi_n)_n\) and \((U_n)_n\) are independent paths with laws \(\mu^{\otimes\mathbb N}\) and \(\nu^{\otimes\mathbb N}\), their zipped path has law \((\mu\otimes\nu)^{\otimes\mathbb N}\).



\textbf{Proof.} Restrict to any finite set of times. The two restricted laws are finite product laws. Reordering their coordinates gives the finite product of pair laws. The infinite product measure is determined by these finite coordinate cylinders, proving the full path identity.

\begin{leanbox}{Lean: \texttt{paired\_\allowbreak{}infinitePi}}
\leansource{\texttt{paired\_\allowbreak{}infinitePi} \quad WassersteinMixing/\allowbreak{}ProductNoise.lean, lines 17–54}
\begin{leancode}
\leanline{0}{theorem\ paired\_infinitePi\ (μ\ :\ Measure\ A)\ (ν\ :\ Measure\ B)}
\leanline{4}{[IsProbabilityMeasure\ μ]\ [IsProbabilityMeasure\ ν]\ :}
\leanline{4}{((Measure.infinitePi\ fun\ \_\ :\ ℕ\ ↦\ μ).prod}
\leanline{6}{(Measure.infinitePi\ fun\ \_\ :\ ℕ\ ↦\ ν)).map}
\leanline{8}{(fun\ p\ n\ ↦\ (p.1\ n,\ p.2\ n))\ =}
\leanline{6}{Measure.infinitePi\ (fun\ \_\ :\ ℕ\ ↦\ μ.prod\ ν)}
\end{leancode}
\leanlabel{How the Lean proof follows the argument}
\texttt{paired\_\allowbreak{}infinite\allowbreak{}Pi} first identifies the finite restriction pushforwards and then applies equality of infinite product measures from their cylinder masses. The proof uses actual measure-preserving coordinate reorderings.
\end{leanbox}
\subsection{\texorpdfstring{Appending a fresh decision path to an actual source}{Appending a fresh decision path to an actual source}}
\label{report-step-WassersteinMixing-ProductNoise-lean-2}


\textbf{Claim.} If an existing source has joint law \(\mu\otimes\eta^{\otimes\mathbb N}\) for its initial state and innovation path, adjoining an independent path with law \(\nu^{\otimes\mathbb N}\) gives initial/\allowbreak{}paired-path law \(\mu\otimes(\eta\otimes\nu)^{\otimes\mathbb N}\).



\textbf{Proof.} Take the actual product extension. The original joint pushforward is preserved on its first coordinate. Reassociate the three factors, apply the zipped-path identity to the last two, and map the result back to the initial/\allowbreak{}paired-path representation.

\begin{leanbox}{Lean: \texttt{augment\_\allowbreak{}initial\_\allowbreak{}noise\_\allowbreak{}law}}
\leansource{\texttt{augment\_\allowbreak{}initial\_\allowbreak{}noise\_\allowbreak{}law} \quad WassersteinMixing/\allowbreak{}ProductNoise.lean, lines 65–102}
\begin{leancode}
\leanline{0}{theorem\ augment\_initial\_noise\_law\ (Q\ :\ Measure\ Ω)\ [IsProbabilityMeasure\ Q]}
\leanline{4}{(μ\ :\ Measure\ X)\ (η\ :\ Measure\ A)\ (ν\ :\ Measure\ B)}
\leanline{4}{[IsProbabilityMeasure\ μ]\ [IsProbabilityMeasure\ η]\ [IsProbabilityMeasure\ ν]}
\leanline{4}{(initial\ :\ Ω\ →\ X)\ (noise\ :\ Ω\ →\ ℕ\ →\ A)}
\leanline{4}{(hi\ :\ AEMeasurable\ initial\ Q)\ (hn\ :\ AEMeasurable\ noise\ Q)}
\leanline{4}{(hjoint\ :\ Q.map\ (fun\ ω\ ↦\ (initial\ ω,\ noise\ ω))\ =}
\leanline{6}{μ.prod\ (Measure.infinitePi\ fun\ \_\ :\ ℕ\ ↦\ η))\ :}
\leanline{4}{(Q.prod\ (Measure.infinitePi\ fun\ \_\ :\ ℕ\ ↦\ ν)).map}
\leanline{6}{(fun\ p\ ↦\ (initial\ p.1,\ fun\ n\ ↦\ (noise\ p.1\ n,\ p.2\ n)))\ =}
\leanline{8}{μ.prod\ (Measure.infinitePi\ fun\ \_\ :\ ℕ\ ↦\ η.prod\ ν)}
\end{leancode}
\leanlabel{How the Lean proof follows the argument}
\texttt{augment\_\allowbreak{}initial\_\allowbreak{}noise\_\allowbreak{}law} combines the actual product pushforward theorem, the measure-preserving product associator and \texttt{paired\_\allowbreak{}infinite\allowbreak{}Pi}. Almost-everywhere measurable input representatives are transported through the product projection before mapping.
\end{leanbox}
\section{\texorpdfstring{Actual Gaussian Metropolis detailed balance}{Actual Gaussian Metropolis detailed balance}}
\label{proof-module-WassersteinMixing-SemigroupMALA-lean}


The genuine Gaussian Langevin proposal and Metropolis acceptance produce a stationary pair law with Gibbs marginals. Detailed balance is proved directly from the symmetric minimum of the forward and backward densities, and its actual \(L^2\) transition is self-adjoint.



Earlier results used here: \cite[Section~3.39]{LangevinLean2026ULASC}.
\subsection{\texorpdfstring{Accepted flow and rejected diagonal mass}{Accepted flow and rejected diagonal mass}}
\label{report-step-WassersteinMixing-SemigroupMALA-lean-1}


\textbf{Construction.} Given proposal-pair law \(J\) and acceptance \(a(x,y)\), retain the accepted weighted pair \(aJ\), and send the rejected pair to \((x,x)\). Its first marginal is the original first marginal, since the weights \(a\) and \(1-a\) sum to one; its second marginal is the actual Metropolis output.



\textbf{Proof of reversal.} Rejected diagonal mass is unchanged by swapping. For a swap-invariant base measure and forward density \(w\), Metropolis weighting gives accepted density \(\min\{w(x,y),w(y,x)\}\), which is symmetric. Thus the entire output pair is reversible.

\begin{leanbox}{Lean: \texttt{metropolisJoint}, \texttt{metropolisJoint\_\allowbreak{}fst}, \texttt{metropolisJoint\_\allowbreak{}snd} and 2 more}
\leansource{\texttt{metropolisJoint} \quad WassersteinMixing/\allowbreak{}SemigroupMALA.lean, lines 28–31}
\begin{leancode}
\leanline{0}{def\ metropolisJoint\ (proposal\ :\ Measure\ (E\ ×\ E))\ (acceptance\ :\ E\ ×\ E\ →\ ℝ)\ :}
\leanline{4}{Measure\ (E\ ×\ E)\ :=}
\leanline{2}{ULA.acceptedFlow\ proposal\ acceptance\ +}
\leanline{4}{(ULA.rejectedFlow\ proposal\ acceptance).map\ (fun\ p\ ↦\ (p.1,\ p.1))}
\end{leancode}
\leansource{\texttt{metropolisJoint\_\allowbreak{}fst} \quad WassersteinMixing/\allowbreak{}SemigroupMALA.lean, lines 34–44}
\begin{leancode}
\leanline{0}{theorem\ metropolisJoint\_fst\ (proposal\ :\ Measure\ (E\ ×\ E))\ (acceptance\ :\ E\ ×\ E\ →\ ℝ)}
\leanline{4}{(ha\ :\ Measurable\ acceptance)\ (ha0\ :\ ∀\ p,\ 0\ ≤\ acceptance\ p)}
\leanline{4}{(ha1\ :\ ∀\ p,\ acceptance\ p\ ≤\ 1)\ :}
\leanline{4}{(metropolisJoint\ proposal\ acceptance).map\ Prod.fst\ =\ proposal.map\ Prod.fst}
\end{leancode}
\leansource{\texttt{metropolisJoint\_\allowbreak{}snd} \quad WassersteinMixing/\allowbreak{}SemigroupMALA.lean, lines 47–53}
\begin{leancode}
\leanline{0}{theorem\ metropolisJoint\_snd\ (proposal\ :\ Measure\ (E\ ×\ E))\ (acceptance\ :\ E\ ×\ E\ →\ ℝ)\ :}
\leanline{4}{(metropolisJoint\ proposal\ acceptance).map\ Prod.snd\ =}
\leanline{6}{ULA.metropolisOutput\ proposal\ acceptance}
\end{leancode}
\leansource{\texttt{metropolisJoint\_\allowbreak{}swap} \quad WassersteinMixing/\allowbreak{}SemigroupMALA.lean, lines 56–64}
\begin{leancode}
\leanline{0}{theorem\ metropolisJoint\_swap\ (proposal\ :\ Measure\ (E\ ×\ E))\ (acceptance\ :\ E\ ×\ E\ →\ ℝ)}
\leanline{4}{(haccepted\ :\ (ULA.acceptedFlow\ proposal\ acceptance).map\ Prod.swap\ =}
\leanline{6}{ULA.acceptedFlow\ proposal\ acceptance)\ :}
\leanline{4}{(metropolisJoint\ proposal\ acceptance).map\ Prod.swap\ =}
\leanline{6}{metropolisJoint\ proposal\ acceptance}
\end{leancode}
\leansource{\texttt{acceptedFlow\_\allowbreak{}swap\_\allowbreak{}of\_\allowbreak{}density} \quad WassersteinMixing/\allowbreak{}SemigroupMALA.lean, lines 67–87}
\begin{leancode}
\leanline{0}{theorem\ acceptedFlow\_swap\_of\_density\ (base\ :\ Measure\ (E\ ×\ E))}
\leanline{4}{(hbase\ :\ base.map\ Prod.swap\ =\ base)\ (w\ :\ E\ ×\ E\ →\ ℝ)}
\leanline{4}{(hw\ :\ Measurable\ w)\ (hw0\ :\ ∀\ p,\ 0\ <\ w\ p)\ :}
\leanline{4}{(ULA.acceptedFlow\ (base.withDensity\ (fun\ p\ ↦\ ENNReal.ofReal\ (w\ p)))}
\leanline{6}{(ULA.metropolisAcceptance\ w)).map\ Prod.swap\ =}
\leanline{4}{ULA.acceptedFlow\ (base.withDensity\ (fun\ p\ ↦\ ENNReal.ofReal\ (w\ p)))}
\leanline{6}{(ULA.metropolisAcceptance\ w)}
\end{leancode}
\leanlabel{How the Lean proof follows the argument}
\texttt{metropolis\allowbreak{}Joint} uses the actual accepted measure and mapped rejected measure. Its marginal lemmas use exact weighted partition and map composition. \texttt{accepted\allowbreak{}Flow\_\allowbreak{}swap\_\allowbreak{}of\_\allowbreak{}density} proves the symmetric minimum-density identity, and \texttt{metropolis\allowbreak{}Joint\_\allowbreak{}swap} combines accepted reversal with the diagonal rejection map.
\end{leanbox}


Earlier results used here: \cite[Section~3.35.1]{LangevinLean2026ULASC}; \cite[Section~3.35.11]{LangevinLean2026ULASC}; \cite[Section~3.35.2]{LangevinLean2026ULASC}; \cite[Section~3.35.3]{LangevinLean2026ULASC}; \cite[Section~3.35.4]{LangevinLean2026ULASC}; \cite[Section~3.35.8]{LangevinLean2026ULASC}; \cite[Section~3.35.9]{LangevinLean2026ULASC}.
\subsection{\texorpdfstring{The actual Gibbs-start Gaussian pair}{The actual Gibbs-start Gaussian pair}}
\label{report-step-WassersteinMixing-SemigroupMALA-lean-2}


\textbf{Claim.} For positive step size, the Gaussian MALA pair has both marginals equal to \(\pi\), is swap invariant, and is a probability law.



\textbf{Proof.} The actual proposal joint law has the Gibbs/\allowbreak{}Gaussian density with respect to the product reference measure. The logarithmic acceptance ratio is the actual ratio of its backward and forward densities, so the preceding minimum-density argument applies. The first marginal is Gibbs by exact acceptance/\allowbreak{}rejection partition. Swap invariance makes the second marginal equal to the first; total mass then equals one. Measurability follows from continuity of the potential and its Lipschitz gradient.

\begin{leanbox}{Lean: \texttt{jointLaw}, \texttt{acceptance\_\allowbreak{}measurable}, \texttt{jointLaw\_\allowbreak{}fst} and 4 more}
\leansource{\texttt{jointLaw} \quad WassersteinMixing/\allowbreak{}SemigroupMALA.lean, lines 96–98}
\begin{leancode}
\leanline{0}{def\ jointLaw\ (V\ :\ E\ →\ ℝ)\ (h\ :\ ℝ)\ :\ Measure\ (E\ ×\ E)\ :=}
\leanline{2}{metropolisJoint\ (ULA.malaProposalJoint\ V\ h\ (ULA.gibbsMeasure\ V))}
\leanline{4}{(fun\ p\ ↦\ ULA.malaAcceptance\ V\ h\ p.1\ p.2)}
\end{leancode}
\leansource{\texttt{acceptance\_\allowbreak{}measurable} \quad WassersteinMixing/\allowbreak{}SemigroupMALA.lean, lines 100–106}
\begin{leancode}
\leanline{0}{private\ theorem\ acceptance\_measurable\ (V\ :\ E\ →\ ℝ)\ (h\ :\ ℝ)\ (L\ :\ ℝ≥0)}
\leanline{4}{(hdiff\ :\ Differentiable\ ℝ\ V)\ (hLip\ :\ LipschitzWith\ L\ (gradient\ V))\ :}
\leanline{4}{Measurable\ (fun\ p\ :\ E\ ×\ E\ ↦\ ULA.malaAcceptance\ V\ h\ p.1\ p.2)}
\end{leancode}
\leansource{\texttt{jointLaw\_\allowbreak{}fst} \quad WassersteinMixing/\allowbreak{}SemigroupMALA.lean, lines 109–119}
\begin{leancode}
\leanline{0}{theorem\ jointLaw\_fst\ (V\ :\ E\ →\ ℝ)\ (h\ :\ ℝ)\ (L\ :\ ℝ≥0)}
\leanline{4}{(hdiff\ :\ Differentiable\ ℝ\ V)\ (hLip\ :\ LipschitzWith\ L\ (gradient\ V))\ :}
\leanline{4}{(jointLaw\ V\ h).map\ Prod.fst\ =\ ULA.gibbsMeasure\ V}
\end{leancode}
\leansource{\texttt{jointLaw\_\allowbreak{}snd\_\allowbreak{}eq\_\allowbreak{}step} \quad WassersteinMixing/\allowbreak{}SemigroupMALA.lean, lines 122–126}
\begin{leancode}
\leanline{0}{theorem\ jointLaw\_snd\_eq\_step\ (V\ :\ E\ →\ ℝ)\ (h\ :\ ℝ)\ (L\ :\ ℝ≥0)}
\leanline{4}{(hdiff\ :\ Differentiable\ ℝ\ V)\ (hLip\ :\ LipschitzWith\ L\ (gradient\ V))\ :}
\leanline{4}{(jointLaw\ V\ h).map\ Prod.snd\ =\ ULA.malaOneStep\ V\ h\ (ULA.gibbsMeasure\ V)}
\end{leancode}
\leansource{\texttt{jointLaw\_\allowbreak{}swap} \quad WassersteinMixing/\allowbreak{}SemigroupMALA.lean, lines 129–151}
\begin{leancode}
\leanline{0}{theorem\ jointLaw\_swap\ (V\ :\ E\ →\ ℝ)\ (h\ :\ ℝ)\ (L\ :\ ℝ≥0)}
\leanline{4}{(hh\ :\ 0\ <\ h)\ (hdiff\ :\ Differentiable\ ℝ\ V)}
\leanline{4}{(hLip\ :\ LipschitzWith\ L\ (gradient\ V))\ :}
\leanline{4}{(jointLaw\ V\ h).map\ Prod.swap\ =\ jointLaw\ V\ h}
\end{leancode}
\leansource{\texttt{jointLaw\_\allowbreak{}snd} \quad WassersteinMixing/\allowbreak{}SemigroupMALA.lean, lines 154–159}
\begin{leancode}
\leanline{0}{theorem\ jointLaw\_snd\ (V\ :\ E\ →\ ℝ)\ (h\ :\ ℝ)\ (L\ :\ ℝ≥0)}
\leanline{4}{(hh\ :\ 0\ <\ h)\ (hdiff\ :\ Differentiable\ ℝ\ V)}
\leanline{4}{(hLip\ :\ LipschitzWith\ L\ (gradient\ V))\ :}
\leanline{4}{(jointLaw\ V\ h).map\ Prod.snd\ =\ ULA.gibbsMeasure\ V}
\end{leancode}
\leansource{\texttt{jointLaw\_\allowbreak{}isProbability} \quad WassersteinMixing/\allowbreak{}SemigroupMALA.lean, lines 163–170}
\begin{leancode}
\leanline{0}{theorem\ jointLaw\_isProbability\ (V\ :\ E\ →\ ℝ)\ (h\ :\ ℝ)\ (L\ :\ ℝ≥0)}
\leanline{4}{[IsProbabilityMeasure\ (ULA.gibbsMeasure\ V)]}
\leanline{4}{(hdiff\ :\ Differentiable\ ℝ\ V)\ (hLip\ :\ LipschitzWith\ L\ (gradient\ V))\ :}
\leanline{4}{IsProbabilityMeasure\ (jointLaw\ V\ h)}
\end{leancode}
\leanlabel{How the Lean proof follows the argument}
\texttt{joint\allowbreak{}Law} specializes the generic Metropolis pair to \texttt{mala\allowbreak{}Proposal\allowbreak{}Joint}. The acceptance helper proves measurability from the true potential and gradient. \texttt{joint\allowbreak{}Law\_\allowbreak{}swap} uses the actual proposal density theorem; its marginal and probability consequences are derived from that measure identity.
\end{leanbox}


Earlier results used here: \cite[Section~3.26.2]{LangevinLean2026ULASC}; \cite[Section~3.26.3]{LangevinLean2026ULASC}; \cite[Section~3.34.2]{LangevinLean2026ULASC}; \cite[Section~3.35.5]{LangevinLean2026ULASC}; \cite[Section~3.35.9]{LangevinLean2026ULASC}; \cite[Section~3.36.1]{LangevinLean2026ULASC}; \cite[Section~3.36.3]{LangevinLean2026ULASC}; \cite[Section~3.36.4]{LangevinLean2026ULASC}; \cite[Section~3.36.5]{LangevinLean2026ULASC}; \cite[Section~3.36.8]{LangevinLean2026ULASC}; \cite[Section~3.37.12]{LangevinLean2026ULASC}; \cite[Section~3.39.1]{LangevinLean2026ULASC}; \cite[Section~3.39.2]{LangevinLean2026ULASC}; \cite[Section~3.39.3]{LangevinLean2026ULASC}; \cite[Section~3.39.4]{LangevinLean2026ULASC}; \cite[Section~3.39.5]{LangevinLean2026ULASC}.
\subsection{\texorpdfstring{The actual self-adjoint contraction and its powers}{The actual self-adjoint contraction and its powers}}
\label{report-step-WassersteinMixing-SemigroupMALA-lean-3}


\textbf{Claim.} The transition compressed from this actual stationary pair is a self-adjoint contraction. Every integer power is self-adjoint and contractive. Actual bounded-continuous two-time correlations are symmetric.



\textbf{Proof.} Apply the stationary-pair construction to the proved marginals and swap identity. Powers of a self-adjoint operator remain self-adjoint; induction applies the contraction bound at each factor. The correlation symmetry follows by mapping the integrable pair function through the actual swap-invariant law.

\begin{leanbox}{Lean: \texttt{operator}, \texttt{operator\_\allowbreak{}norm\_\allowbreak{}le}, \texttt{operator\_\allowbreak{}isSelfAdjoint} and 3 more}
\leansource{\texttt{operator} \quad WassersteinMixing/\allowbreak{}SemigroupMALA.lean, lines 176–180}
\begin{leancode}
\leanline{0}{def\ operator\ (V\ :\ E\ →\ ℝ)\ (h\ :\ ℝ)\ (L\ :\ ℝ≥0)\ (hh\ :\ 0\ <\ h)}
\leanline{4}{(hdiff\ :\ Differentiable\ ℝ\ V)\ (hLip\ :\ LipschitzWith\ L\ (gradient\ V))\ :}
\leanline{4}{Lp\ F\ 2\ (ULA.gibbsMeasure\ V)\ →L[{\leanfallback 𝕜}]\ Lp\ F\ 2\ (ULA.gibbsMeasure\ V)\ :=}
\leanline{2}{StationaryPair.transition\ ({\leanfallback 𝕜}\ :=\ {\leanfallback 𝕜})\ (ULA.gibbsMeasure\ V)\ (jointLaw\ V\ h)}
\leanline{4}{(jointLaw\_fst\ V\ h\ L\ hdiff\ hLip)\ (jointLaw\_snd\ V\ h\ L\ hh\ hdiff\ hLip)}
\end{leancode}
\leansource{\texttt{operator\_\allowbreak{}norm\_\allowbreak{}le} \quad WassersteinMixing/\allowbreak{}SemigroupMALA.lean, lines 183–187}
\begin{leancode}
\leanline{0}{theorem\ operator\_norm\_le\ (V\ :\ E\ →\ ℝ)\ (h\ :\ ℝ)\ (L\ :\ ℝ≥0)\ (hh\ :\ 0\ <\ h)}
\leanline{4}{(hdiff\ :\ Differentiable\ ℝ\ V)\ (hLip\ :\ LipschitzWith\ L\ (gradient\ V))\ :}
\leanline{4}{‖operator\ ({\leanfallback 𝕜}\ :=\ {\leanfallback 𝕜})\ (F\ :=\ F)\ V\ h\ L\ hh\ hdiff\ hLip‖\ ≤\ 1}
\end{leancode}
\leansource{\texttt{operator\_\allowbreak{}isSelfAdjoint} \quad WassersteinMixing/\allowbreak{}SemigroupMALA.lean, lines 191–196}
\begin{leancode}
\leanline{0}{theorem\ operator\_isSelfAdjoint\ (V\ :\ E\ →\ ℝ)\ (h\ :\ ℝ)\ (L\ :\ ℝ≥0)\ (hh\ :\ 0\ <\ h)}
\leanline{4}{(hdiff\ :\ Differentiable\ ℝ\ V)\ (hLip\ :\ LipschitzWith\ L\ (gradient\ V))\ :}
\leanline{4}{IsSelfAdjoint\ (operator\ ({\leanfallback 𝕜}\ :=\ {\leanfallback 𝕜})\ (F\ :=\ F)\ V\ h\ L\ hh\ hdiff\ hLip)}
\end{leancode}
\leansource{\texttt{operator\_\allowbreak{}pow\_\allowbreak{}isSelfAdjoint} \quad WassersteinMixing/\allowbreak{}SemigroupMALA.lean, lines 199–202}
\begin{leancode}
\leanline{0}{theorem\ operator\_pow\_isSelfAdjoint\ (V\ :\ E\ →\ ℝ)\ (h\ :\ ℝ)\ (L\ :\ ℝ≥0)\ (hh\ :\ 0\ <\ h)}
\leanline{4}{(hdiff\ :\ Differentiable\ ℝ\ V)\ (hLip\ :\ LipschitzWith\ L\ (gradient\ V))\ (n\ :\ ℕ)\ :}
\leanline{4}{IsSelfAdjoint\ (operator\ ({\leanfallback 𝕜}\ :=\ {\leanfallback 𝕜})\ (F\ :=\ F)\ V\ h\ L\ hh\ hdiff\ hLip\ \textasciicircum{}\ n)}
\end{leancode}
\leansource{\texttt{operator\_\allowbreak{}pow\_\allowbreak{}norm\_\allowbreak{}le} \quad WassersteinMixing/\allowbreak{}SemigroupMALA.lean, lines 205–217}
\begin{leancode}
\leanline{0}{theorem\ operator\_pow\_norm\_le\ (V\ :\ E\ →\ ℝ)\ (h\ :\ ℝ)\ (L\ :\ ℝ≥0)\ (hh\ :\ 0\ <\ h)}
\leanline{4}{(hdiff\ :\ Differentiable\ ℝ\ V)\ (hLip\ :\ LipschitzWith\ L\ (gradient\ V))\ (n\ :\ ℕ)\ :}
\leanline{4}{‖operator\ ({\leanfallback 𝕜}\ :=\ {\leanfallback 𝕜})\ (F\ :=\ F)\ V\ h\ L\ hh\ hdiff\ hLip\ \textasciicircum{}\ n‖\ ≤\ 1}
\end{leancode}
\leansource{\texttt{jointLaw\_\allowbreak{}correlation\_\allowbreak{}symmetric} \quad WassersteinMixing/\allowbreak{}SemigroupMALA.lean, lines 222–232}
\begin{leancode}
\leanline{0}{theorem\ jointLaw\_correlation\_symmetric\ (V\ :\ E\ →\ ℝ)\ (h\ :\ ℝ)\ (L\ :\ ℝ≥0)\ (hh\ :\ 0\ <\ h)}
\leanline{4}{(hdiff\ :\ Differentiable\ ℝ\ V)\ (hLip\ :\ LipschitzWith\ L\ (gradient\ V))}
\leanline{4}{(f\ g\ :\ E\ →\ F)\ (hf\ :\ Continuous\ f)\ (hg\ :\ Continuous\ g)\ :}
\leanline{4}{(∫\ p,\ inner\ {\leanfallback 𝕜}\ (f\ p.2)\ (g\ p.1)\ ∂jointLaw\ V\ h)\ =}
\leanline{6}{∫\ p,\ inner\ {\leanfallback 𝕜}\ (f\ p.1)\ (g\ p.2)\ ∂jointLaw\ V\ h}
\end{leancode}
\leanlabel{How the Lean proof follows the argument}
\texttt{operator} instantiates the generic stationary-pair adjoint compression over either real or complex targets. Its bounds and self-adjointness use proved law properties. The power bounds are inductive consequences; the final real correlation lemma uses actual pair-law reversal directly.
\end{leanbox}


Earlier results used here: \cite[Section~3.26.3]{LangevinLean2026ULASC}.
\section{\texorpdfstring{Reversal symmetry of actual anchored kernel powers}{Reversal symmetry of actual anchored kernel powers}}
\label{proof-module-WassersteinMixing-SemigroupPowers-lean}


One-step detailed balance implies reversal of every actual finite-step pair \(\mu(dx)K^n(x,dy)\). The proof uses nonnegative observable correlations and actual kernel integration, so the original initial coordinate remains present throughout.

\subsection{\texorpdfstring{Iterating the actual detailed-balance correlation}{Iterating the actual detailed-balance correlation}}
\label{report-step-WassersteinMixing-SemigroupPowers-lean-1}


\textbf{Claim.} With \((Kf)(x)=\int f(y)K(x,dy)\), one-step pair reversal gives

\[\int f\,Kg\,d\mu=\int Kf\,g\,d\mu,\qquad\int f\,K^ng\,d\mu=\int K^nf\,g\,d\mu\]

for measurable nonnegative observables.



\textbf{Proof.} Product-kernel integration expresses the first pairing as the actual two-time integral. Swap invariance exchanges the two coordinates. For the iterated identity, move one averaging operator to the other observable by one-step balance, then apply the induction hypothesis. Kernel integration preserves measurability at every iterate.

\begin{leanbox}{Lean: \texttt{average}, \texttt{average\_\allowbreak{}measurable}, \texttt{average\_\allowbreak{}iterate\_\allowbreak{}measurable} and 2 more}
\leansource{\texttt{average} \quad WassersteinMixing/\allowbreak{}SemigroupPowers.lean, lines 26–26}
\begin{leancode}
\leanline{0}{def\ average\ (κ\ :\ Kernel\ E\ E)\ (f\ :\ E\ →\ ℝ≥0∞)\ (x\ :\ E)\ :\ ℝ≥0∞\ :=\ ∫⁻\ y,\ f\ y\ ∂κ\ x}
\end{leancode}
\leansource{\texttt{average\_\allowbreak{}measurable} \quad WassersteinMixing/\allowbreak{}SemigroupPowers.lean, lines 29–30}
\begin{leancode}
\leanline{0}{theorem\ average\_measurable\ (κ\ :\ Kernel\ E\ E)\ (f\ :\ E\ →\ ℝ≥0∞)\ (hf\ :\ Measurable\ f)\ :}
\leanline{4}{Measurable\ (average\ κ\ f)}
\end{leancode}
\leansource{\texttt{average\_\allowbreak{}iterate\_\allowbreak{}measurable} \quad WassersteinMixing/\allowbreak{}SemigroupPowers.lean, lines 33–39}
\begin{leancode}
\leanline{0}{theorem\ average\_iterate\_measurable\ (κ\ :\ Kernel\ E\ E)\ (f\ :\ E\ →\ ℝ≥0∞)}
\leanline{4}{(hf\ :\ Measurable\ f)\ (n\ :\ ℕ)\ :\ Measurable\ ((average\ κ)\textasciicircum{}[n]\ f)}
\end{leancode}
\leansource{\texttt{average\_\allowbreak{}balance} \quad WassersteinMixing/\allowbreak{}SemigroupPowers.lean, lines 42–63}
\begin{leancode}
\leanline{0}{theorem\ average\_balance\ (μ\ :\ Measure\ E)\ [SFinite\ μ]\ (κ\ :\ Kernel\ E\ E)}
\leanline{4}{[IsSFiniteKernel\ κ]\ (hswap\ :\ (μ\ ⊗ₘ\ κ).map\ Prod.swap\ =\ μ\ ⊗ₘ\ κ)}
\leanline{4}{(f\ g\ :\ E\ →\ ℝ≥0∞)\ (hf\ :\ Measurable\ f)\ (hg\ :\ Measurable\ g)\ :}
\leanline{4}{(∫⁻\ x,\ f\ x\ *\ average\ κ\ g\ x\ ∂μ)\ =\ ∫⁻\ x,\ average\ κ\ f\ x\ *\ g\ x\ ∂μ}
\end{leancode}
\leansource{\texttt{average\_\allowbreak{}iterate\_\allowbreak{}balance} \quad WassersteinMixing/\allowbreak{}SemigroupPowers.lean, lines 67–81}
\begin{leancode}
\leanline{0}{theorem\ average\_iterate\_balance\ (μ\ :\ Measure\ E)\ [SFinite\ μ]\ (κ\ :\ Kernel\ E\ E)}
\leanline{4}{[IsSFiniteKernel\ κ]\ (hswap\ :\ (μ\ ⊗ₘ\ κ).map\ Prod.swap\ =\ μ\ ⊗ₘ\ κ)}
\leanline{4}{(f\ g\ :\ E\ →\ ℝ≥0∞)\ (hf\ :\ Measurable\ f)\ (hg\ :\ Measurable\ g)\ (n\ :\ ℕ)\ :}
\leanline{4}{(∫⁻\ x,\ f\ x\ *\ ((average\ κ)\textasciicircum{}[n]\ g)\ x\ ∂μ)\ =}
\leanline{6}{∫⁻\ x,\ ((average\ κ)\textasciicircum{}[n]\ f)\ x\ *\ g\ x\ ∂μ}
\end{leancode}
\leanlabel{How the Lean proof follows the argument}
\texttt{average} is actual kernel integration. Its measurability lemmas justify each iteration. \texttt{average\_\allowbreak{}balance} changes variables under the swap law and \texttt{average\_\allowbreak{}iterate\_\allowbreak{}balance} proves the iterated pairing by induction.
\end{leanbox}
\subsection{\texorpdfstring{Actual powers, rectangles and their marginals}{Actual powers, rectangles and their marginals}}
\label{report-step-WassersteinMixing-SemigroupPowers-lean-2}


\textbf{Claim.} Define \(K^0\) as the identity kernel and use actual kernel composition for successors. The pair law \(\mu\otimes K^n\) is swap invariant and has both marginals \(\mu\).



\textbf{Proof.} Kernel integration identifies the averaging action of the power with the iterated averaging function. Pair rectangle masses are indicator correlations, so the proved iterated balance gives equality with their reversed rectangles. Equality on rectangles proves full measure reversal. The first marginal is \(\mu\) because the kernel is Markov; swapping gives the second.

\begin{leanbox}{Lean: \texttt{kernelPower}, \texttt{kernelPower\_\allowbreak{}isMarkov}, \texttt{kernelPower\_\allowbreak{}average} and 4 more}
\leansource{\texttt{kernelPower} \quad WassersteinMixing/\allowbreak{}SemigroupPowers.lean, lines 84–86}
\begin{leancode}
\leanline{0}{def\ kernelPower\ (κ\ :\ Kernel\ E\ E)\ :\ ℕ\ →\ Kernel\ E\ E}
\leanline{2}{|\ 0\ =>\ Kernel.id}
\leanline{2}{|\ n\ +\ 1\ =>\ kernelPower\ κ\ n\ ∘ₖ\ κ}
\end{leancode}
\leansource{\texttt{kernelPower\_\allowbreak{}isMarkov} \quad WassersteinMixing/\allowbreak{}SemigroupPowers.lean, lines 88–92}
\begin{leancode}
\leanline{0}{instance\ kernelPower\_isMarkov\ (κ\ :\ Kernel\ E\ E)\ [IsMarkovKernel\ κ]\ (n\ :\ ℕ)\ :}
\leanline{4}{IsMarkovKernel\ (kernelPower\ κ\ n)\ :=\ by}
\leanline{2}{induction\ n\ with}
\leanline{2}{|\ zero\ =>\ unfold\ kernelPower;\ infer\_instance}
\leanline{2}{|\ succ\ n\ ih\ =>\ unfold\ kernelPower;\ infer\_instance}
\end{leancode}
\leansource{\texttt{kernelPower\_\allowbreak{}average} \quad WassersteinMixing/\allowbreak{}SemigroupPowers.lean, lines 95–107}
\begin{leancode}
\leanline{0}{theorem\ kernelPower\_average\ (κ\ :\ Kernel\ E\ E)\ (f\ :\ E\ →\ ℝ≥0∞)\ (hf\ :\ Measurable\ f)}
\leanline{4}{(n\ :\ ℕ)\ (x\ :\ E)\ :\ average\ (kernelPower\ κ\ n)\ f\ x\ =\ ((average\ κ)\textasciicircum{}[n]\ f)\ x}
\end{leancode}
\leansource{\texttt{kernelPower\_\allowbreak{}rectangle} \quad WassersteinMixing/\allowbreak{}SemigroupPowers.lean, lines 114–125}
\begin{leancode}
\leanline{0}{theorem\ kernelPower\_rectangle\ (n\ :\ ℕ)\ (s\ t\ :\ Set\ E)}
\leanline{4}{(hs\ :\ MeasurableSet\ s)\ (ht\ :\ MeasurableSet\ t)\ :}
\leanline{4}{(μ\ ⊗ₘ\ kernelPower\ κ\ n)\ (s\ ×ˢ\ t)\ =}
\leanline{6}{∫⁻\ x,\ s.indicator\ (fun\ \_\ ↦\ (1\ :\ ℝ≥0∞))\ x\ *}
\leanline{8}{((average\ κ)\textasciicircum{}[n]\ (t.indicator\ (fun\ \_\ ↦\ (1\ :\ ℝ≥0∞))))\ x\ ∂μ}
\end{leancode}
\leansource{\texttt{kernelPower\_\allowbreak{}pair\_\allowbreak{}swap} \quad WassersteinMixing/\allowbreak{}SemigroupPowers.lean, lines 129–145}
\begin{leancode}
\leanline{0}{theorem\ kernelPower\_pair\_swap}
\leanline{4}{(hswap\ :\ (μ\ ⊗ₘ\ κ).map\ Prod.swap\ =\ μ\ ⊗ₘ\ κ)\ (n\ :\ ℕ)\ :}
\leanline{4}{(μ\ ⊗ₘ\ kernelPower\ κ\ n).map\ Prod.swap\ =\ μ\ ⊗ₘ\ kernelPower\ κ\ n}
\end{leancode}
\leansource{\texttt{kernelPower\_\allowbreak{}pair\_\allowbreak{}fst} \quad WassersteinMixing/\allowbreak{}SemigroupPowers.lean, lines 148–150}
\begin{leancode}
\leanline{0}{theorem\ kernelPower\_pair\_fst\ (n\ :\ ℕ)\ :}
\leanline{4}{(μ\ ⊗ₘ\ kernelPower\ κ\ n).map\ Prod.fst\ =\ μ}
\end{leancode}
\leansource{\texttt{kernelPower\_\allowbreak{}pair\_\allowbreak{}snd} \quad WassersteinMixing/\allowbreak{}SemigroupPowers.lean, lines 153–157}
\begin{leancode}
\leanline{0}{theorem\ kernelPower\_pair\_snd}
\leanline{4}{(hswap\ :\ (μ\ ⊗ₘ\ κ).map\ Prod.swap\ =\ μ\ ⊗ₘ\ κ)\ (n\ :\ ℕ)\ :}
\leanline{4}{(μ\ ⊗ₘ\ kernelPower\ κ\ n).map\ Prod.snd\ =\ μ}
\end{leancode}
\leanlabel{How the Lean proof follows the argument}
\texttt{kernel\allowbreak{}Power} is the literal recursive kernel composition. \texttt{kernel\allowbreak{}Power\_\allowbreak{}average} and \texttt{kernel\allowbreak{}Power\_\allowbreak{}rectangle} connect it with the actual correlations. \texttt{kernel\allowbreak{}Power\_\allowbreak{}pair\_\allowbreak{}swap} uses product-measure extensionality. Its marginal theorems derive the exact laws.
\end{leanbox}


Earlier results used here: \cite[Section~3.35.5]{LangevinLean2026ULASC}.
\subsection{\texorpdfstring{Applying the power construction to the proved MALA pair}{Applying the power construction to the proved MALA pair}}
\label{report-step-WassersteinMixing-SemigroupPowers-lean-3}


\textbf{Construction.} Disintegrate the actual reversible Gibbs MALA pair to obtain a Markov kernel. Its initial-law product is exactly the original pair. The actual kernel-power pair is therefore a probability measure, is reversible, and has both Gibbs marginals.



\textbf{Proof.} The disintegration theorem identifies the original pair with its first marginal times its conditional kernel. Its first marginal is already proved to be Gibbs. Apply the generic power reversal and marginal theorems to this exact one-step identity.

\begin{leanbox}{Lean: \texttt{kernel}, \texttt{kernel\_\allowbreak{}isMarkov}, \texttt{kernel\_\allowbreak{}pair\_\allowbreak{}eq} and 6 more}
\leansource{\texttt{kernel} \quad WassersteinMixing/\allowbreak{}SemigroupPowers.lean, lines 169–173}
\begin{leancode}
\leanline{0}{def\ kernel\ (V\ :\ E\ →\ ℝ)\ (h\ :\ ℝ)\ (L\ :\ ℝ≥0)}
\leanline{4}{[IsProbabilityMeasure\ (ULA.gibbsMeasure\ V)]}
\leanline{4}{(hdiff\ :\ Differentiable\ ℝ\ V)\ (hLip\ :\ LipschitzWith\ L\ (gradient\ V))\ :\ Kernel\ E\ E\ :=\ by}
\leanline{2}{have\ :=\ jointLaw\_isProbability\ V\ h\ L\ hdiff\ hLip}
\leanline{2}{exact\ (jointLaw\ V\ h).condKernel}
\end{leancode}
\leansource{\texttt{kernel\_\allowbreak{}isMarkov} \quad WassersteinMixing/\allowbreak{}SemigroupPowers.lean, lines 175–181}
\begin{leancode}
\leanline{0}{instance\ kernel\_isMarkov\ (V\ :\ E\ →\ ℝ)\ (h\ :\ ℝ)\ (L\ :\ ℝ≥0)}
\leanline{4}{[IsProbabilityMeasure\ (ULA.gibbsMeasure\ V)]}
\leanline{4}{(hdiff\ :\ Differentiable\ ℝ\ V)\ (hLip\ :\ LipschitzWith\ L\ (gradient\ V))\ :}
\leanline{4}{IsMarkovKernel\ (kernel\ V\ h\ L\ hdiff\ hLip)\ :=\ by}
\leanline{2}{have\ :=\ jointLaw\_isProbability\ V\ h\ L\ hdiff\ hLip}
\leanline{2}{unfold\ kernel}
\leanline{2}{infer\_instance}
\end{leancode}
\leansource{\texttt{kernel\_\allowbreak{}pair\_\allowbreak{}eq} \quad WassersteinMixing/\allowbreak{}SemigroupPowers.lean, lines 184–193}
\begin{leancode}
\leanline{0}{theorem\ kernel\_pair\_eq\ (V\ :\ E\ →\ ℝ)\ (h\ :\ ℝ)\ (L\ :\ ℝ≥0)}
\leanline{4}{[IsProbabilityMeasure\ (ULA.gibbsMeasure\ V)]}
\leanline{4}{(hdiff\ :\ Differentiable\ ℝ\ V)\ (hLip\ :\ LipschitzWith\ L\ (gradient\ V))\ :}
\leanline{4}{ULA.gibbsMeasure\ V\ ⊗ₘ\ kernel\ V\ h\ L\ hdiff\ hLip\ =\ jointLaw\ V\ h}
\end{leancode}
\leansource{\texttt{kernel\_\allowbreak{}pair\_\allowbreak{}swap} \quad WassersteinMixing/\allowbreak{}SemigroupPowers.lean, lines 196–202}
\begin{leancode}
\leanline{0}{theorem\ kernel\_pair\_swap\ (V\ :\ E\ →\ ℝ)\ (h\ :\ ℝ)\ (L\ :\ ℝ≥0)}
\leanline{4}{[IsProbabilityMeasure\ (ULA.gibbsMeasure\ V)]\ (hh\ :\ 0\ <\ h)}
\leanline{4}{(hdiff\ :\ Differentiable\ ℝ\ V)\ (hLip\ :\ LipschitzWith\ L\ (gradient\ V))\ :}
\leanline{4}{(ULA.gibbsMeasure\ V\ ⊗ₘ\ kernel\ V\ h\ L\ hdiff\ hLip).map\ Prod.swap\ =}
\leanline{6}{ULA.gibbsMeasure\ V\ ⊗ₘ\ kernel\ V\ h\ L\ hdiff\ hLip}
\end{leancode}
\leansource{\texttt{powerJointLaw} \quad WassersteinMixing/\allowbreak{}SemigroupPowers.lean, lines 206–210}
\begin{leancode}
\leanline{0}{def\ powerJointLaw\ (V\ :\ E\ →\ ℝ)\ (h\ :\ ℝ)\ (L\ :\ ℝ≥0)}
\leanline{4}{[IsProbabilityMeasure\ (ULA.gibbsMeasure\ V)]}
\leanline{4}{(hdiff\ :\ Differentiable\ ℝ\ V)\ (hLip\ :\ LipschitzWith\ L\ (gradient\ V))\ (n\ :\ ℕ)\ :}
\leanline{4}{Measure\ (E\ ×\ E)\ :=}
\leanline{2}{ULA.gibbsMeasure\ V\ ⊗ₘ\ StationaryKernelPower.kernelPower\ (kernel\ V\ h\ L\ hdiff\ hLip)\ n}
\end{leancode}
\leansource{\texttt{powerJointLaw\_\allowbreak{}isProbability} \quad WassersteinMixing/\allowbreak{}SemigroupPowers.lean, lines 212–217}
\begin{leancode}
\leanline{0}{instance\ powerJointLaw\_isProbability\ (V\ :\ E\ →\ ℝ)\ (h\ :\ ℝ)\ (L\ :\ ℝ≥0)}
\leanline{4}{[IsProbabilityMeasure\ (ULA.gibbsMeasure\ V)]}
\leanline{4}{(hdiff\ :\ Differentiable\ ℝ\ V)\ (hLip\ :\ LipschitzWith\ L\ (gradient\ V))\ (n\ :\ ℕ)\ :}
\leanline{4}{IsProbabilityMeasure\ (powerJointLaw\ V\ h\ L\ hdiff\ hLip\ n)\ :=\ by}
\leanline{2}{unfold\ powerJointLaw}
\leanline{2}{infer\_instance}
\end{leancode}
\leansource{\texttt{powerJointLaw\_\allowbreak{}swap} \quad WassersteinMixing/\allowbreak{}SemigroupPowers.lean, lines 220–225}
\begin{leancode}
\leanline{0}{theorem\ powerJointLaw\_swap\ (V\ :\ E\ →\ ℝ)\ (h\ :\ ℝ)\ (L\ :\ ℝ≥0)}
\leanline{4}{[IsProbabilityMeasure\ (ULA.gibbsMeasure\ V)]\ (hh\ :\ 0\ <\ h)}
\leanline{4}{(hdiff\ :\ Differentiable\ ℝ\ V)\ (hLip\ :\ LipschitzWith\ L\ (gradient\ V))\ (n\ :\ ℕ)\ :}
\leanline{4}{(powerJointLaw\ V\ h\ L\ hdiff\ hLip\ n).map\ Prod.swap\ =\ powerJointLaw\ V\ h\ L\ hdiff\ hLip\ n}
\end{leancode}
\leansource{\texttt{powerJointLaw\_\allowbreak{}fst} \quad WassersteinMixing/\allowbreak{}SemigroupPowers.lean, lines 228–233}
\begin{leancode}
\leanline{0}{theorem\ powerJointLaw\_fst\ (V\ :\ E\ →\ ℝ)\ (h\ :\ ℝ)\ (L\ :\ ℝ≥0)}
\leanline{4}{[IsProbabilityMeasure\ (ULA.gibbsMeasure\ V)]}
\leanline{4}{(hdiff\ :\ Differentiable\ ℝ\ V)\ (hLip\ :\ LipschitzWith\ L\ (gradient\ V))\ (n\ :\ ℕ)\ :}
\leanline{4}{(powerJointLaw\ V\ h\ L\ hdiff\ hLip\ n).map\ Prod.fst\ =\ ULA.gibbsMeasure\ V}
\end{leancode}
\leansource{\texttt{powerJointLaw\_\allowbreak{}snd} \quad WassersteinMixing/\allowbreak{}SemigroupPowers.lean, lines 236–241}
\begin{leancode}
\leanline{0}{theorem\ powerJointLaw\_snd\ (V\ :\ E\ →\ ℝ)\ (h\ :\ ℝ)\ (L\ :\ ℝ≥0)}
\leanline{4}{[IsProbabilityMeasure\ (ULA.gibbsMeasure\ V)]\ (hh\ :\ 0\ <\ h)}
\leanline{4}{(hdiff\ :\ Differentiable\ ℝ\ V)\ (hLip\ :\ LipschitzWith\ L\ (gradient\ V))\ (n\ :\ ℕ)\ :}
\leanline{4}{(powerJointLaw\ V\ h\ L\ hdiff\ hLip\ n).map\ Prod.snd\ =\ ULA.gibbsMeasure\ V}
\end{leancode}
\leanlabel{How the Lean proof follows the argument}
\texttt{kernel} is the actual conditional kernel of \texttt{joint\allowbreak{}Law}, not a postulated transition. \texttt{kernel\_\allowbreak{}pair\_\allowbreak{}eq} applies disintegration. \texttt{power\allowbreak{}Joint\allowbreak{}Law} and its law theorems instantiate the proved generic kernel-power results.
\end{leanbox}


Earlier results used here: \cite[Section~3.26.3]{LangevinLean2026ULASC}.
\section{\texorpdfstring{The explicit actual Gaussian Metropolis kernel}{The explicit actual Gaussian Metropolis kernel}}
\label{proof-module-WassersteinMixing-SemigroupMALAKernel-lean}


The transition kernel samples an actual fresh standard Gaussian, maps accepted mass to the Langevin proposal and rejected mass to the current state. Its action on every initial law equals the existing Gaussian MALA law step exactly.

\subsection{\texorpdfstring{Constructing the actual fresh-noise transition}{Constructing the actual fresh-noise transition}}
\label{report-step-WassersteinMixing-SemigroupMALAKernel-lean-1}


\textbf{Construction.} At state \(x\), take the product \(\delta_x\otimes\gamma\), where \(\gamma\) is standard Gaussian. Weight by actual Metropolis acceptance and rejection probabilities and map to the proposed or unchanged state. The resulting explicit kernel at \(x\) is exactly the MALA law step applied to \(\delta_x\).



\textbf{Proof.} The product kernel is Markov. Measurability of the actual gradient proposal and acceptance permits weighted mapping. Evaluating at a state and substituting the Dirac initial law yields the same accepted/\allowbreak{}rejected formula as the actual MALA step.

\begin{leanbox}{Lean: \texttt{gaussianPairKernel}, \texttt{gaussianPairKernel\_\allowbreak{}isMarkov}, \texttt{explicitKernel} and 1 more}
\leansource{\texttt{gaussianPairKernel} \quad WassersteinMixing/\allowbreak{}SemigroupMALAKernel.lean, lines 24–25}
\begin{leancode}
\leanline{0}{def\ gaussianPairKernel\ :\ Kernel\ E\ (E\ ×\ E)\ :=}
\leanline{2}{Kernel.id\ ×ₖ\ Kernel.const\ E\ (stdGaussian\ E)}
\end{leancode}
\leansource{\texttt{gaussianPairKernel\_\allowbreak{}isMarkov} \quad WassersteinMixing/\allowbreak{}SemigroupMALAKernel.lean, lines 27–29}
\begin{leancode}
\leanline{0}{instance\ gaussianPairKernel\_isMarkov\ :\ IsMarkovKernel\ (gaussianPairKernel\ (E\ :=\ E))\ :=\ by}
\leanline{2}{unfold\ gaussianPairKernel}
\leanline{2}{infer\_instance}
\end{leancode}
\leansource{\texttt{explicitKernel} \quad WassersteinMixing/\allowbreak{}SemigroupMALAKernel.lean, lines 32–37}
\begin{leancode}
\leanline{0}{def\ explicitKernel\ (V\ :\ E\ →\ ℝ)\ (h\ :\ ℝ)\ :\ Kernel\ E\ E\ :=}
\leanline{2}{((gaussianPairKernel\ (E\ :=\ E)).withDensity}
\leanline{6}{(fun\ \_\ p\ ↦\ ENNReal.ofReal\ (ULA.malaProposalAcceptance\ V\ h\ p))).map}
\leanline{8}{(fun\ p\ ↦\ ULA.malaProposal\ V\ h\ p.1\ p.2)\ +}
\leanline{2}{((gaussianPairKernel\ (E\ :=\ E)).withDensity}
\leanline{6}{(fun\ \_\ p\ ↦\ ENNReal.ofReal\ (1\ -\ ULA.malaProposalAcceptance\ V\ h\ p))).map\ Prod.fst}
\end{leancode}
\leansource{\texttt{explicitKernel\_\allowbreak{}apply} \quad WassersteinMixing/\allowbreak{}SemigroupMALAKernel.lean, lines 40–54}
\begin{leancode}
\leanline{0}{theorem\ explicitKernel\_apply\ (V\ :\ E\ →\ ℝ)\ (h\ :\ ℝ)\ (L\ :\ ℝ≥0)}
\leanline{4}{(hdiff\ :\ Differentiable\ ℝ\ V)\ (hLip\ :\ LipschitzWith\ L\ (gradient\ V))\ (x\ :\ E)\ :}
\leanline{4}{explicitKernel\ V\ h\ x\ =\ ULA.malaOneStep\ V\ h\ (Measure.dirac\ x)}
\end{leancode}
\leanlabel{How the Lean proof follows the argument}
\texttt{gaussian\allowbreak{}Pair\allowbreak{}Kernel} is the product of the identity and constant-Gaussian kernels. \texttt{explicit\allowbreak{}Kernel} is the sum of its actual weighted mapped kernels. \texttt{explicit\allowbreak{}Kernel\_\allowbreak{}apply} compares their event masses with the Dirac-start MALA step.
\end{leanbox}


Earlier results used here: \cite[Section~3.34.1]{LangevinLean2026ULASC}; \cite[Section~3.36.2]{LangevinLean2026ULASC}; \cite[Section~3.36.5]{LangevinLean2026ULASC}; \cite[Section~3.36.6]{LangevinLean2026ULASC}; \cite[Section~3.36.8]{LangevinLean2026ULASC}.
\subsection{\texorpdfstring{The exact law action and mass normalization}{The exact law action and mass normalization}}
\label{report-step-WassersteinMixing-SemigroupMALAKernel-lean-2}


\textbf{Claim.} For every sigma-finite initial law \(\mu\), the explicit kernel action is the actual MALA output law. It has total mass one at every state.



\textbf{Proof.} Composing the fresh-Gaussian pair kernel with \(\mu\) gives \(\mu\otimes\gamma\). Mapping and weighting commute with actual measure/\allowbreak{}kernel composition, so the explicit output is exactly the accepted/\allowbreak{}rejected law formula for \(\mu\). For a Dirac input, acceptance plus rejection partition the entire Gaussian mass, proving the Markov normalization.

\begin{leanbox}{Lean: \texttt{gaussianPairKernel\_\allowbreak{}comp}, \texttt{explicitKernel\_\allowbreak{}comp}, \texttt{explicitKernel\_\allowbreak{}isMarkov}}
\leansource{\texttt{gaussianPairKernel\_\allowbreak{}comp} \quad WassersteinMixing/\allowbreak{}SemigroupMALAKernel.lean, lines 58–60}
\begin{leancode}
\leanline{0}{theorem\ gaussianPairKernel\_comp\ (μ\ :\ Measure\ E)\ [SFinite\ μ]\ :}
\leanline{4}{gaussianPairKernel\ (E\ :=\ E)\ ∘ₘ\ μ\ =\ μ.prod\ (stdGaussian\ E)}
\end{leancode}
\leansource{\texttt{explicitKernel\_\allowbreak{}comp} \quad WassersteinMixing/\allowbreak{}SemigroupMALAKernel.lean, lines 64–77}
\begin{leancode}
\leanline{0}{theorem\ explicitKernel\_comp\ (V\ :\ E\ →\ ℝ)\ (h\ :\ ℝ)\ (L\ :\ ℝ≥0)}
\leanline{4}{(hdiff\ :\ Differentiable\ ℝ\ V)\ (hLip\ :\ LipschitzWith\ L\ (gradient\ V))}
\leanline{4}{(μ\ :\ Measure\ E)\ [SFinite\ μ]\ :}
\leanline{4}{explicitKernel\ V\ h\ ∘ₘ\ μ\ =\ ULA.malaOneStep\ V\ h\ μ}
\end{leancode}
\leansource{\texttt{explicitKernel\_\allowbreak{}isMarkov} \quad WassersteinMixing/\allowbreak{}SemigroupMALAKernel.lean, lines 80–98}
\begin{leancode}
\leanline{0}{theorem\ explicitKernel\_isMarkov\ (V\ :\ E\ →\ ℝ)\ (h\ :\ ℝ)\ (L\ :\ ℝ≥0)}
\leanline{4}{(hdiff\ :\ Differentiable\ ℝ\ V)\ (hLip\ :\ LipschitzWith\ L\ (gradient\ V))\ :}
\leanline{4}{IsMarkovKernel\ (explicitKernel\ V\ h)}
\end{leancode}
\leanlabel{How the Lean proof follows the argument}
\texttt{gaussian\allowbreak{}Pair\allowbreak{}Kernel\_\allowbreak{}comp} identifies the actual independent-product sampling law. \texttt{explicit\allowbreak{}Kernel\_\allowbreak{}comp} uses map, density and sum composition identities. \texttt{explicit\allowbreak{}Kernel\_\allowbreak{}is\allowbreak{}Markov} proves unit mass from the true accepted/\allowbreak{}rejected partition.
\end{leanbox}


Earlier results used here: \cite[Section~3.34.1]{LangevinLean2026ULASC}; \cite[Section~3.35.1]{LangevinLean2026ULASC}; \cite[Section~3.35.2]{LangevinLean2026ULASC}; \cite[Section~3.35.4]{LangevinLean2026ULASC}; \cite[Section~3.36.2]{LangevinLean2026ULASC}; \cite[Section~3.36.3]{LangevinLean2026ULASC}; \cite[Section~3.36.4]{LangevinLean2026ULASC}; \cite[Section~3.36.5]{LangevinLean2026ULASC}; \cite[Section~3.36.6]{LangevinLean2026ULASC}; \cite[Section~3.36.8]{LangevinLean2026ULASC}.
\section{\texorpdfstring{An actual common-randomness Euler and Metropolis update}{An actual common-randomness Euler and Metropolis update}}
\label{proof-module-WassersteinMixing-SemigroupRandomMALA-lean}


Independent Gaussian proposals and independent uniform decisions are combined into a single innovation law. A deterministic paired update advances Euler and MALA states with common Gaussian input, and its actual law is the established coupled-step law.



Earlier results used here: \cite[Section~3.33]{LangevinLean2026ULASC}.
\subsection{\texorpdfstring{The paired actual sampling map}{The paired actual sampling map}}
\label{report-step-WassersteinMixing-SemigroupRandomMALA-lean-1}


\textbf{Construction.} At paired states \((x,y)\), update the Euler coordinate by \(x-h\nabla V(x)+\sqrt{2h}\,\xi\). In the second coordinate use the actual Gaussian proposal from \(y\), accepted by an independent uniform threshold and otherwise leave \(y\) unchanged. The map is measurable and its first coordinate is exactly the Euler update, including along the full recursion.



\textbf{Proof.} The innovation measure is the product of standard Gaussian and unit-interval probability. Continuity of the actual proposal and measurability of its acceptance event give measurable branches; projection to the first coordinate removes the acceptance decision.

\begin{leanbox}{Lean: \texttt{innovationLaw}, \texttt{innovationLaw\_\allowbreak{}isProbability}, \texttt{update} and 3 more}
\leansource{\texttt{innovationLaw} \quad WassersteinMixing/\allowbreak{}SemigroupRandomMALA.lean, lines 25–25}
\begin{leancode}
\leanline{0}{def\ innovationLaw\ :\ Measure\ (E\ ×\ I)\ :=\ (stdGaussian\ E).prod\ volume}
\end{leancode}
\leansource{\texttt{innovationLaw\_\allowbreak{}isProbability} \quad WassersteinMixing/\allowbreak{}SemigroupRandomMALA.lean, lines 27–29}
\begin{leancode}
\leanline{0}{instance\ innovationLaw\_isProbability\ :\ IsProbabilityMeasure\ (innovationLaw\ (E\ :=\ E))\ :=\ by}
\leanline{2}{unfold\ innovationLaw}
\leanline{2}{infer\_instance}
\end{leancode}
\leansource{\texttt{update} \quad WassersteinMixing/\allowbreak{}SemigroupRandomMALA.lean, lines 32–36}
\begin{leancode}
\leanline{0}{def\ update\ (V\ :\ E\ →\ ℝ)\ (h\ :\ ℝ)\ (x\ :\ E\ ×\ E)\ (z\ :\ E\ ×\ I)\ :\ E\ ×\ E\ :=}
\leanline{2}{uniformDecision}
\leanline{4}{(fun\ p\ :\ (E\ ×\ E)\ ×\ E\ ↦\ ULA.malaProposalAcceptance\ V\ h\ (p.1.2,\ p.2))}
\leanline{4}{(fun\ p\ ↦\ (ULA.malaProposal\ V\ h\ p.1.1\ p.2,\ ULA.malaProposal\ V\ h\ p.1.2\ p.2))}
\leanline{4}{(fun\ p\ ↦\ (ULA.malaProposal\ V\ h\ p.1.1\ p.2,\ p.1.2))\ ((x,\ z.1),\ z.2)}
\end{leancode}
\leansource{\texttt{update\_\allowbreak{}fst} \quad WassersteinMixing/\allowbreak{}SemigroupRandomMALA.lean, lines 38–41}
\begin{leancode}
\leanline{0}{@[simp]\ theorem\ update\_fst\ (V\ :\ E\ →\ ℝ)\ (h\ :\ ℝ)\ (x\ :\ E\ ×\ E)\ (z\ :\ E\ ×\ I)\ :}
\leanline{4}{(update\ V\ h\ x\ z).1\ =\ ULA.malaProposal\ V\ h\ x.1\ z.1}
\end{leancode}
\leansource{\texttt{chain\_\allowbreak{}fst} \quad WassersteinMixing/\allowbreak{}SemigroupRandomMALA.lean, lines 44–50}
\begin{leancode}
\leanline{0}{theorem\ chain\_fst\ (V\ :\ E\ →\ ℝ)\ (h\ :\ ℝ)\ (x\ :\ E\ ×\ E)\ (w\ :\ ℕ\ →\ E\ ×\ I)\ (n\ :\ ℕ)\ :}
\leanline{4}{(ULA.noiseDrivenChain\ (update\ V\ h)\ x\ w\ n).1\ =}
\leanline{6}{ULA.noiseDrivenChain\ (fun\ x\ z\ ↦\ ULA.malaProposal\ V\ h\ x\ z)\ x.1}
\leanline{8}{(fun\ k\ ↦\ (w\ k).1)\ n}
\end{leancode}
\leansource{\texttt{update\_\allowbreak{}measurable} \quad WassersteinMixing/\allowbreak{}SemigroupRandomMALA.lean, lines 52–62}
\begin{leancode}
\leanline{0}{theorem\ update\_measurable\ (V\ :\ E\ →\ ℝ)\ (h\ :\ ℝ)\ (L\ :\ ℝ≥0)}
\leanline{4}{(hdiff\ :\ Differentiable\ ℝ\ V)\ (hLip\ :\ LipschitzWith\ L\ (gradient\ V))\ :}
\leanline{4}{Measurable\ (Function.uncurry\ (update\ V\ h))}
\end{leancode}
\leanlabel{How the Lean proof follows the argument}
\texttt{innovation\allowbreak{}Law} is the actual Gaussian/\allowbreak{}uniform product. \texttt{update} is the paired threshold update. \texttt{update\_\allowbreak{}fst} and \texttt{chain\_\allowbreak{}fst} identify the actual Euler coordinate; \texttt{update\_\allowbreak{}measurable} combines proposal and threshold measurability.
\end{leanbox}


Earlier results used here: \cite[Section~3.20.1]{LangevinLean2026ULASC}; \cite[Section~3.20.3]{LangevinLean2026ULASC}; \cite[Section~3.34.1]{LangevinLean2026ULASC}; \cite[Section~3.36.2]{LangevinLean2026ULASC}; \cite[Section~3.36.5]{LangevinLean2026ULASC}; \cite[Section~3.36.6]{LangevinLean2026ULASC}.
\subsection{\texorpdfstring{Identifying every actual paired Markov law}{Identifying every actual paired Markov law}}
\label{report-step-WassersteinMixing-SemigroupRandomMALA-lean-2}


\textbf{Claim.} The sampled paired update has exactly the coupled Euler/\allowbreak{}MALA law, and its iterated laws agree at every finite time.



\textbf{Proof.} Integrate the uniform decision using the proved acceptance/\allowbreak{}rejection identity. The resulting weighted Gaussian maps are precisely the coupled-step construction. Starting from the same paired initial law, substitute that exact one-step equality at each recursion step to obtain equality of all iterated laws.

\begin{leanbox}{Lean: \texttt{markovStep\_\allowbreak{}eq\_\allowbreak{}coupledStep}, \texttt{chainLaw\_\allowbreak{}eq\_\allowbreak{}coupledChain}}
\leansource{\texttt{markovStep\_\allowbreak{}eq\_\allowbreak{}coupledStep} \quad WassersteinMixing/\allowbreak{}SemigroupRandomMALA.lean, lines 66–81}
\begin{leancode}
\leanline{0}{theorem\ markovStep\_eq\_coupledStep\ (V\ :\ E\ →\ ℝ)\ (h\ :\ ℝ)\ (L\ :\ ℝ≥0)}
\leanline{4}{(hdiff\ :\ Differentiable\ ℝ\ V)\ (hLip\ :\ LipschitzWith\ L\ (gradient\ V))}
\leanline{4}{(ρ\ :\ Measure\ (E\ ×\ E))\ [IsProbabilityMeasure\ ρ]\ :}
\leanline{4}{ULA.markovStep\ (update\ V\ h)\ innovationLaw\ ρ\ =\ ULA.malaCoupledStep\ V\ h\ ρ}
\end{leancode}
\leansource{\texttt{chainLaw\_\allowbreak{}eq\_\allowbreak{}coupledChain} \quad WassersteinMixing/\allowbreak{}SemigroupRandomMALA.lean, lines 84–100}
\begin{leancode}
\leanline{0}{theorem\ chainLaw\_eq\_coupledChain\ (V\ :\ E\ →\ ℝ)\ (h\ :\ ℝ)\ (L\ :\ ℝ≥0)}
\leanline{4}{(hdiff\ :\ Differentiable\ ℝ\ V)\ (hLip\ :\ LipschitzWith\ L\ (gradient\ V))}
\leanline{4}{(μ\ :\ Measure\ E)\ [IsProbabilityMeasure\ μ]\ (n\ :\ ℕ)\ :}
\leanline{4}{ULA.chainLaw\ (update\ V\ h)\ (μ.map\ (fun\ x\ ↦\ (x,\ x)))}
\leanline{6}{(Measure.infinitePi\ (fun\ \_\ :\ ℕ\ ↦\ innovationLaw))\ n\ =}
\leanline{8}{ULA.malaCoupledChain\ V\ h\ μ\ n}
\end{leancode}
\leanlabel{How the Lean proof follows the argument}
\texttt{markov\allowbreak{}Step\_\allowbreak{}eq\_\allowbreak{}coupled\allowbreak{}Step} invokes \texttt{uniform\allowbreak{}Decision\_\allowbreak{}law} for the actual acceptance and proposal functions. \texttt{chain\allowbreak{}Law\_\allowbreak{}eq\_\allowbreak{}coupled\allowbreak{}Chain} inducts on the actual Markov-law recursion.
\end{leanbox}


Earlier results used here: \cite[Section~3.20.5]{LangevinLean2026ULASC}; \cite[Section~3.33.5]{LangevinLean2026ULASC}; \cite[Section~3.33.6]{LangevinLean2026ULASC}; \cite[Section~3.33.9]{LangevinLean2026ULASC}; \cite[Section~3.36.14]{LangevinLean2026ULASC}; \cite[Section~3.36.15]{LangevinLean2026ULASC}; \cite[Section~3.36.3]{LangevinLean2026ULASC}; \cite[Section~3.36.4]{LangevinLean2026ULASC}; \cite[Section~3.36.5]{LangevinLean2026ULASC}; \cite[Section~3.36.6]{LangevinLean2026ULASC}; \cite[Section~3.36.9]{LangevinLean2026ULASC}.
\section{\texorpdfstring{The actual Brownian and independent-uniform MALA source}{The actual Brownian and independent-uniform MALA source}}
\label{proof-module-WassersteinMixing-SemigroupRandomMALACanonical-lean}


Adjoin an independent infinite uniform path to the actual Gibbs/\allowbreak{}Brownian probability space. Its normalized Brownian innovations and uniform decisions have the required product path law, so the deterministic common-randomness recursion realizes the actual coupled Euler/\allowbreak{}MALA laws.

\subsection{\texorpdfstring{The genuine augmented source and Gaussian path law}{The genuine augmented source and Gaussian path law}}
\label{report-step-WassersteinMixing-SemigroupRandomMALACanonical-lean-1}


\textbf{Construction.} Take the product of the actual Gibbs/\allowbreak{}Brownian source with independent uniforms \((U_n)_n\). The Brownian innovations at step \(r>0\) have the standard Gaussian infinite-product law and are independent of the actual initial state.



\textbf{Proof.} Coordinate Brownian increments are independent centered Gaussians with the exact interval variance. Their prescribed normalization gives standard Gaussian vectors. The initial Gibbs factor was independent of the Brownian source. The actual product extension retains both these laws and adds only independent uniform decisions.

\begin{leanbox}{Lean: \texttt{Space}, \texttt{uniformPath}, \texttt{uniformPath\_\allowbreak{}isProbability} and 5 more}
\leansource{\texttt{Space} \quad WassersteinMixing/\allowbreak{}SemigroupRandomMALACanonical.lean, lines 20–20}
\begin{leancode}
\leanline{0}{abbrev\ Space\ (P\ :\ Parameters)\ :=\ ActualSpace\ P\ ×\ (ℕ\ →\ I)}
\end{leancode}
\leansource{\texttt{uniformPath} \quad WassersteinMixing/\allowbreak{}SemigroupRandomMALACanonical.lean, lines 22–22}
\begin{leancode}
\leanline{0}{def\ uniformPath\ :\ Measure\ (ℕ\ →\ I)\ :=\ Measure.infinitePi\ (fun\ \_\ :\ ℕ\ ↦\ (volume\ :\ Measure\ I))}
\end{leancode}
\leansource{\texttt{uniformPath\_\allowbreak{}isProbability} \quad WassersteinMixing/\allowbreak{}SemigroupRandomMALACanonical.lean, lines 24–26}
\begin{leancode}
\leanline{0}{instance\ uniformPath\_isProbability\ :\ IsProbabilityMeasure\ uniformPath\ :=\ by}
\leanline{2}{unfold\ uniformPath}
\leanline{2}{infer\_instance}
\end{leancode}
\leansource{\texttt{probability} \quad WassersteinMixing/\allowbreak{}SemigroupRandomMALACanonical.lean, lines 28–29}
\begin{leancode}
\leanline{0}{def\ probability\ (P\ :\ Parameters)\ (V\ :\ E\ P.dimension\ →\ ℝ)\ :\ Measure\ (Space\ P)\ :=}
\leanline{2}{(actualProbability\ P\ V).prod\ uniformPath}
\end{leancode}
\leansource{\texttt{probability\_\allowbreak{}isProbability} \quad WassersteinMixing/\allowbreak{}SemigroupRandomMALACanonical.lean, lines 31–35}
\begin{leancode}
\leanline{0}{theorem\ probability\_isProbability\ (P\ :\ Parameters)\ (V\ :\ E\ P.dimension\ →\ ℝ)}
\leanline{4}{(hV\ :\ PotentialAssumptions\ P\ V)\ :\ IsProbabilityMeasure\ (probability\ P\ V)}
\end{leancode}
\leansource{\texttt{gaussian} \quad WassersteinMixing/\allowbreak{}SemigroupRandomMALACanonical.lean, lines 38–45}
\begin{leancode}
\leanline{0}{def\ gaussian\ (P\ :\ Parameters)\ (V\ :\ E\ P.dimension\ →\ ℝ)\ (hV\ :\ PotentialAssumptions\ P\ V)}
\leanline{4}{(r\ :\ ℝ)\ :\ ℕ\ →\ Lp\ (E\ P.dimension)\ 2\ (actualProbability\ P\ V)\ :=\ by}
\leanline{2}{let\ :=\ actualTarget\_isProbability\ P\ V\ hV}
\leanline{2}{exact\ ULA.Hilbert.productNoise}
\leanline{4}{(ULA.CanonicalBrownian.vectorProbability\ (Fin\ P.dimension))\ (diagonalGibbsCoupling\ P\ V)}
\leanline{4}{(ULA.Hilbert.canonicalEulerNoise\ r)}
\leanline{4}{(ULA.Hilbert.sourceEulerNoise\_memLp\ \_\ ULA.CanonicalBrownian.coordinate}
\leanline{6}{ULA.CanonicalBrownian.coordinate\_isPreBrownian\ r)}
\end{leancode}
\leansource{\texttt{gaussian\_\allowbreak{}path\_\allowbreak{}law} \quad WassersteinMixing/\allowbreak{}SemigroupRandomMALACanonical.lean, lines 47–61}
\begin{leancode}
\leanline{0}{theorem\ gaussian\_path\_law\ (P\ :\ Parameters)\ (V\ :\ E\ P.dimension\ →\ ℝ)}
\leanline{4}{(hV\ :\ PotentialAssumptions\ P\ V)\ (r\ :\ ℝ)\ (hr\ :\ 0\ <\ r)\ :}
\leanline{4}{(actualProbability\ P\ V).map\ (fun\ ω\ n\ ↦\ gaussian\ P\ V\ hV\ r\ n\ ω)\ =}
\leanline{6}{Measure.infinitePi\ (fun\ \_\ :\ ℕ\ ↦\ stdGaussian\ (E\ P.dimension))}
\end{leancode}
\leansource{\texttt{initial\_\allowbreak{}gaussian\_\allowbreak{}law} \quad WassersteinMixing/\allowbreak{}SemigroupRandomMALACanonical.lean, lines 63–84}
\begin{leancode}
\leanline{0}{theorem\ initial\_gaussian\_law\ (P\ :\ Parameters)\ (V\ :\ E\ P.dimension\ →\ ℝ)}
\leanline{4}{(hV\ :\ PotentialAssumptions\ P\ V)\ (r\ :\ ℝ)\ (hr\ :\ 0\ <\ r)\ :}
\leanline{4}{(actualProbability\ P\ V).map}
\leanline{6}{(fun\ ω\ ↦\ (actualInitial\ P\ V\ hV\ ω,\ fun\ n\ ↦\ gaussian\ P\ V\ hV\ r\ n\ ω))\ =}
\leanline{6}{(ULA.gibbsMeasure\ V).prod}
\leanline{8}{(Measure.infinitePi\ (fun\ \_\ :\ ℕ\ ↦\ stdGaussian\ (E\ P.dimension)))}
\end{leancode}
\leanlabel{How the Lean proof follows the argument}
\texttt{Space}, \texttt{uniform\allowbreak{}Path} and \texttt{probability} define the actual augmented space and probability. \texttt{gaussian\_\allowbreak{}path\_\allowbreak{}law} and \texttt{initial\_\allowbreak{}gaussian\_\allowbreak{}law} specialize the actual Brownian innovation law and independence theorems.
\end{leanbox}


Earlier results used here: \cite[Section~3.17.2]{LangevinLean2026ULASC}; \cite[Section~3.22.21]{LangevinLean2026ULASC}; \cite[Section~3.23.5]{LangevinLean2026ULASC}; \cite[Section~3.26.3]{LangevinLean2026ULASC}; \cite[Section~3.32.14]{LangevinLean2026ULASC}; \cite[Section~3.32.15]{LangevinLean2026ULASC}; \cite[Section~3.32.16]{LangevinLean2026ULASC}; \cite[Section~3.32.5]{LangevinLean2026ULASC}; \cite[Section~3.41.1]{LangevinLean2026ULASC}; \cite[Section~3.45.3]{LangevinLean2026ULASC}; \cite[Section~3.8.17]{LangevinLean2026ULASC}; \cite[Section~3.8.19]{LangevinLean2026ULASC}; \cite[Section~3.8.20]{LangevinLean2026ULASC}; \cite[Section~3.8.22]{LangevinLean2026ULASC}.
\subsection{\texorpdfstring{The exact canonical sampled recursion law}{The exact canonical sampled recursion law}}
\label{report-step-WassersteinMixing-SemigroupRandomMALACanonical-lean-2}


\textbf{Claim.} The augmented initial/\allowbreak{}innovation input has the product law of the diagonal Gibbs initialization and independent Gaussian/\allowbreak{}uniform pairs. The canonical recursion therefore has exactly the actual coupled Euler/\allowbreak{}MALA finite-time law.



\textbf{Proof.} Zip the actual Gaussian path with the independent uniform path using the proved infinite-product identity. Apply the deterministic sampled-chain pushforward theorem to this exact input law. The one-step sampled law has already been identified with the actual coupled Gaussian Metropolis step.

\begin{leanbox}{Lean: \texttt{input}, \texttt{input\_\allowbreak{}law}, \texttt{state} and 1 more}
\leansource{\texttt{input} \quad WassersteinMixing/\allowbreak{}SemigroupRandomMALACanonical.lean, lines 86–88}
\begin{leancode}
\leanline{0}{def\ input\ (P\ :\ Parameters)\ (V\ :\ E\ P.dimension\ →\ ℝ)\ (hV\ :\ PotentialAssumptions\ P\ V)}
\leanline{4}{(r\ :\ ℝ)\ (ω\ :\ Space\ P)\ :\ (E\ P.dimension)\ ×\ (ℕ\ →\ (E\ P.dimension)\ ×\ I)\ :=}
\leanline{2}{(actualInitial\ P\ V\ hV\ ω.1,\ fun\ n\ ↦\ (gaussian\ P\ V\ hV\ r\ n\ ω.1,\ ω.2\ n))}
\end{leancode}
\leansource{\texttt{input\_\allowbreak{}law} \quad WassersteinMixing/\allowbreak{}SemigroupRandomMALACanonical.lean, lines 90–102}
\begin{leancode}
\leanline{0}{theorem\ input\_law\ (P\ :\ Parameters)\ (V\ :\ E\ P.dimension\ →\ ℝ)}
\leanline{4}{(hV\ :\ PotentialAssumptions\ P\ V)\ (r\ :\ ℝ)\ (hr\ :\ 0\ <\ r)\ :}
\leanline{4}{(probability\ P\ V).map\ (input\ P\ V\ hV\ r)\ =}
\leanline{6}{(ULA.gibbsMeasure\ V).prod}
\leanline{8}{(Measure.infinitePi\ (fun\ \_\ :\ ℕ\ ↦\ innovationLaw))}
\end{leancode}
\leansource{\texttt{state} \quad WassersteinMixing/\allowbreak{}SemigroupRandomMALACanonical.lean, lines 105–108}
\begin{leancode}
\leanline{0}{def\ state\ (P\ :\ Parameters)\ (V\ :\ E\ P.dimension\ →\ ℝ)\ (hV\ :\ PotentialAssumptions\ P\ V)}
\leanline{4}{(r\ :\ ℝ)\ (n\ :\ ℕ)\ (ω\ :\ Space\ P)\ :\ E\ P.dimension\ ×\ E\ P.dimension\ :=}
\leanline{2}{ULA.noiseDrivenChain\ (RandomMALA.update\ V\ r)}
\leanline{4}{((input\ P\ V\ hV\ r\ ω).1,\ (input\ P\ V\ hV\ r\ ω).1)\ (input\ P\ V\ hV\ r\ ω).2\ n}
\end{leancode}
\leansource{\texttt{state\_\allowbreak{}law} \quad WassersteinMixing/\allowbreak{}SemigroupRandomMALACanonical.lean, lines 110–140}
\begin{leancode}
\leanline{0}{theorem\ state\_law\ (P\ :\ Parameters)\ (V\ :\ E\ P.dimension\ →\ ℝ)}
\leanline{4}{(hV\ :\ PotentialAssumptions\ P\ V)\ (r\ :\ ℝ)\ (hr\ :\ 0\ <\ r)\ (n\ :\ ℕ)\ :}
\leanline{4}{(probability\ P\ V).map\ (state\ P\ V\ hV\ r\ n)\ =}
\leanline{6}{ULA.malaCoupledChain\ V\ r\ (ULA.gibbsMeasure\ V)\ n}
\end{leancode}
\leanlabel{How the Lean proof follows the argument}
\texttt{input\_\allowbreak{}law} uses \texttt{augment\_\allowbreak{}initial\_\allowbreak{}noise\_\allowbreak{}law}. \texttt{state} applies the genuine paired recursion to this input; \texttt{state\_\allowbreak{}law} uses its exact deterministic-to-Markov law theorem and \texttt{chain\allowbreak{}Law\_\allowbreak{}eq\_\allowbreak{}coupled\allowbreak{}Chain}.
\end{leanbox}


Earlier results used here: \cite[Section~3.20.1]{LangevinLean2026ULASC}; \cite[Section~3.20.4]{LangevinLean2026ULASC}; \cite[Section~3.20.5]{LangevinLean2026ULASC}; \cite[Section~3.26.3]{LangevinLean2026ULASC}; \cite[Section~3.36.14]{LangevinLean2026ULASC}.
\section{\texorpdfstring{Moments and coupling error of the actual random MALA path}{Moments and coupling error of the actual random MALA path}}
\label{proof-module-WassersteinMixing-SemigroupRandomMALABounds-lean}


The canonical sampled MALA recursion is converted to a genuine \(L^2\) random variable. Its marginal remains exactly Gibbs, and its mean distance from the synchronous Euler coordinate satisfies the proved cumulative rejection estimate.

\subsection{\texorpdfstring{Actual representatives, Gibbs laws and square integrability}{Actual representatives, Gibbs laws and square integrability}}
\label{report-step-WassersteinMixing-SemigroupRandomMALABounds-lean-1}


\textbf{Claim.} The sampled input and paired states are measurable. The first coordinate is the actual Euler recursion; the second has Gibbs law and finite second moment, hence defines an actual \(L^2\) variable with that same law.



\textbf{Proof.} Compose the measurable sampled update through finite recursion. Projecting its exact law gives the established MALA marginal, which is Gibbs by invariance at each finite step. Pull the Gibbs finite second moment through this exact pushforward to obtain square integrability; passing to the \(L^2\) class preserves its almost-everywhere representative and law.

\begin{leanbox}{Lean: \texttt{input\_\allowbreak{}aemeasurable}, \texttt{state\_\allowbreak{}aemeasurable}, \texttt{euler} and 6 more}
\leansource{\texttt{input\_\allowbreak{}aemeasurable} \quad WassersteinMixing/\allowbreak{}SemigroupRandomMALABounds.lean, lines 21–25}
\begin{leancode}
\leanline{0}{theorem\ input\_aemeasurable\ (r\ :\ ℝ)\ :}
\leanline{4}{AEMeasurable\ (input\ P\ V\ hV\ r)\ (probability\ P\ V)}
\end{leancode}
\leansource{\texttt{state\_\allowbreak{}aemeasurable} \quad WassersteinMixing/\allowbreak{}SemigroupRandomMALABounds.lean, lines 27–33}
\begin{leancode}
\leanline{0}{theorem\ state\_aemeasurable\ (r\ :\ ℝ)\ (n\ :\ ℕ)\ :}
\leanline{4}{AEMeasurable\ (state\ P\ V\ hV\ r\ n)\ (probability\ P\ V)}
\end{leancode}
\leansource{\texttt{euler} \quad WassersteinMixing/\allowbreak{}SemigroupRandomMALABounds.lean, lines 36–39}
\begin{leancode}
\leanline{0}{def\ euler\ (r\ :\ ℝ)\ (n\ :\ ℕ)\ :\ Lp\ (E\ P.dimension)\ 2\ (actualProbability\ P\ V)\ :=\ by}
\leanline{2}{let\ :=\ actualProbability\_isProbability\ P\ V\ hV}
\leanline{2}{exact\ ULA.Hilbert.lpChain\ (actualProbability\ P\ V)\ V\ r\ hV.gradientLipschitz}
\leanline{4}{(actualInitial\ P\ V\ hV)\ (gaussian\ P\ V\ hV\ r)\ n}
\end{leancode}
\leansource{\texttt{state\_\allowbreak{}fst\_\allowbreak{}eq\_\allowbreak{}euler} \quad WassersteinMixing/\allowbreak{}SemigroupRandomMALABounds.lean, lines 41–55}
\begin{leancode}
\leanline{0}{theorem\ state\_fst\_eq\_euler\ (r\ :\ ℝ)\ (n\ :\ ℕ)\ :}
\leanline{4}{(fun\ ω\ ↦\ (state\ P\ V\ hV\ r\ n\ ω).1)\ =ᵐ[probability\ P\ V]}
\leanline{6}{(fun\ ω\ ↦\ euler\ P\ V\ hV\ r\ n\ ω.1)}
\end{leancode}
\leansource{\texttt{state\_\allowbreak{}snd\_\allowbreak{}law} \quad WassersteinMixing/\allowbreak{}SemigroupRandomMALABounds.lean, lines 58–67}
\begin{leancode}
\leanline{0}{theorem\ state\_snd\_law\ (r\ :\ ℝ)\ (hr\ :\ 0\ <\ r)\ (n\ :\ ℕ)\ :}
\leanline{4}{(probability\ P\ V).map\ (fun\ ω\ ↦\ (state\ P\ V\ hV\ r\ n\ ω).2)\ =\ ULA.gibbsMeasure\ V}
\end{leancode}
\leansource{\texttt{state\_\allowbreak{}snd\_\allowbreak{}memLp} \quad WassersteinMixing/\allowbreak{}SemigroupRandomMALABounds.lean, lines 69–75}
\begin{leancode}
\leanline{0}{theorem\ state\_snd\_memLp\ (r\ :\ ℝ)\ (hr\ :\ 0\ <\ r)\ (n\ :\ ℕ)\ :}
\leanline{4}{MemLp\ (fun\ ω\ ↦\ (state\ P\ V\ hV\ r\ n\ ω).2)\ 2\ (probability\ P\ V)}
\end{leancode}
\leansource{\texttt{mala} \quad WassersteinMixing/\allowbreak{}SemigroupRandomMALABounds.lean, lines 78–79}
\begin{leancode}
\leanline{0}{def\ mala\ (r\ :\ ℝ)\ (hr\ :\ 0\ <\ r)\ (n\ :\ ℕ)\ :\ Lp\ (E\ P.dimension)\ 2\ (probability\ P\ V)\ :=}
\leanline{2}{(state\_snd\_memLp\ P\ V\ hV\ r\ hr\ n).toLp\ (fun\ ω\ ↦\ (state\ P\ V\ hV\ r\ n\ ω).2)}
\end{leancode}
\leansource{\texttt{mala\_\allowbreak{}coe} \quad WassersteinMixing/\allowbreak{}SemigroupRandomMALABounds.lean, lines 81–83}
\begin{leancode}
\leanline{0}{theorem\ mala\_coe\ (r\ :\ ℝ)\ (hr\ :\ 0\ <\ r)\ (n\ :\ ℕ)\ :}
\leanline{4}{mala\ P\ V\ hV\ r\ hr\ n\ =ᵐ[probability\ P\ V]\ fun\ ω\ ↦\ (state\ P\ V\ hV\ r\ n\ ω).2}
\end{leancode}
\leansource{\texttt{mala\_\allowbreak{}law} \quad WassersteinMixing/\allowbreak{}SemigroupRandomMALABounds.lean, lines 85–87}
\begin{leancode}
\leanline{0}{theorem\ mala\_law\ (r\ :\ ℝ)\ (hr\ :\ 0\ <\ r)\ (n\ :\ ℕ)\ :}
\leanline{4}{(probability\ P\ V).map\ (mala\ P\ V\ hV\ r\ hr\ n)\ =\ ULA.gibbsMeasure\ V}
\end{leancode}
\leanlabel{How the Lean proof follows the argument}
The initial measurability lemmas justify the sampled pushforwards. \texttt{state\_\allowbreak{}snd\_\allowbreak{}law} and \texttt{state\_\allowbreak{}snd\_\allowbreak{}mem\allowbreak{}Lp} establish the true marginal and moment before defining \texttt{mala}. \texttt{mala\_\allowbreak{}coe} and \texttt{mala\_\allowbreak{}law} retain the representative and exact law.
\end{leanbox}


Earlier results used here: \cite[Section~3.19.1]{LangevinLean2026ULASC}; \cite[Section~3.20.1]{LangevinLean2026ULASC}; \cite[Section~3.20.4]{LangevinLean2026ULASC}; \cite[Section~3.23.1]{LangevinLean2026ULASC}; \cite[Section~3.26.3]{LangevinLean2026ULASC}; \cite[Section~3.36.17]{LangevinLean2026ULASC}; \cite[Section~3.39.6]{LangevinLean2026ULASC}; \cite[Section~3.42.1]{LangevinLean2026ULASC}.
\subsection{\texorpdfstring{The actual accumulated rejection distance}{The actual accumulated rejection distance}}
\label{report-step-WassersteinMixing-SemigroupRandomMALABounds-lean-2}


\textbf{Claim.} Under the stated small-step conditions, the expected distance between the actual MALA state and the lifted Euler state is at most \(nCr\sqrt r\), where \(C\) is the proved Gaussian rejection constant for the potential.



\textbf{Proof.} Both coordinates are square integrable, so their difference norm is integrable. The actual joint law equals the established coupled-chain law. Transfer its cumulative rejection-cost estimate through that exact law and through the almost-everywhere representatives of the \(L^2\) states.

\begin{leanbox}{Lean: \texttt{distance\_\allowbreak{}integrable}, \texttt{distance\_\allowbreak{}integral\_\allowbreak{}le}}
\leansource{\texttt{distance\_\allowbreak{}integrable} \quad WassersteinMixing/\allowbreak{}SemigroupRandomMALABounds.lean, lines 89–99}
\begin{leancode}
\leanline{0}{theorem\ distance\_integrable\ (r\ :\ ℝ)\ (hr\ :\ 0\ <\ r)\ (hr1\ :\ r\ ≤\ 1)}
\leanline{4}{(hsmall\ :\ r\ ≤\ 2\ /\ (P.alpha\ +\ P.beta))\ (n\ :\ ℕ)\ :}
\leanline{4}{Integrable\ (fun\ ω\ ↦\ ‖(state\ P\ V\ hV\ r\ n\ ω).1\ -\ (state\ P\ V\ hV\ r\ n\ ω).2‖)}
\leanline{6}{(probability\ P\ V)}
\end{leancode}
\leansource{\texttt{distance\_\allowbreak{}integral\_\allowbreak{}le} \quad WassersteinMixing/\allowbreak{}SemigroupRandomMALABounds.lean, lines 103–114}
\begin{leancode}
\leanline{0}{theorem\ distance\_integral\_le\ (r\ :\ ℝ)\ (hr\ :\ 0\ <\ r)\ (hr1\ :\ r\ ≤\ 1)}
\leanline{4}{(hsmall\ :\ r\ ≤\ 2\ /\ (P.alpha\ +\ P.beta))\ (n\ :\ ℕ)\ :}
\leanline{4}{(∫\ ω,\ ‖(state\ P\ V\ hV\ r\ n\ ω).1\ -\ (state\ P\ V\ hV\ r\ n\ ω).2‖\ ∂probability\ P\ V)\ ≤}
\leanline{6}{(n\ :\ ℝ)\ *\ ULA.malaRejectionConstant\ V\ ⟨P.beta,\ P.beta\_pos.le⟩\ *\ r\ *\ Real.sqrt\ r}
\end{leancode}
\leanlabel{How the Lean proof follows the argument}
\texttt{distance\_\allowbreak{}integrable} uses the two true \(L^2\) coordinates. \texttt{distance\_\allowbreak{}integral\_\allowbreak{}le} rewrites the actual expectation under \texttt{state\_\allowbreak{}law} and invokes the proved coupled MALA/\allowbreak{}Euler rejection estimate.
\end{leanbox}


Earlier results used here: \cite[Section~3.34.22]{LangevinLean2026ULASC}; \cite[Section~3.36.25]{LangevinLean2026ULASC}; \cite[Section~3.36.26]{LangevinLean2026ULASC}; \cite[Section~3.39.6]{LangevinLean2026ULASC}.
\section{\texorpdfstring{Preserving actual dynamics under an independent source extension}{Preserving actual dynamics under an independent source extension}}
\label{proof-module-WassersteinMixing-SemigroupLift-lean}


Appending an independent coordinate to the probability space changes neither the Langevin dynamics nor their laws. Measure-preserving pullback commutes with the true drift, Bochner time integral, Euler recursion and stationary observable evaluation.



Earlier results used here: \cite[Section~3.14]{LangevinLean2026Kinetic}; \cite[Section~3.16]{LangevinLean2026ULASC}.
\subsection{\texorpdfstring{Lifting genuine solutions and observables}{Lifting genuine solutions and observables}}
\label{report-step-WassersteinMixing-SemigroupLift-lean-1}


\textbf{Claim.} Pulling a Langevin solution and its additive noise through the first projection of an independent product gives the same integral equation on the extended space. Euler endpoints and stationary evaluations lift in the same way, and their laws are unchanged.



\textbf{Proof.} The first projection is measure-preserving because the added measure is a probability. Its \(L^2\) pullback is a linear isometry. The actual Lipschitz drift commutes with this pullback by comparing representatives. Bounded linear pullback commutes with the Bochner integral; induction gives the Euler identity. Exact map composition and measure preservation give the law and observable-evaluation identities.

\begin{leanbox}{Lean: \texttt{langevin\_\allowbreak{}integral\_\allowbreak{}sourceLift}, \texttt{additiveEulerChain\_\allowbreak{}sourceLift}, \texttt{sourceLift\_\allowbreak{}law} and 1 more}
\leansource{\texttt{langevin\_\allowbreak{}integral\_\allowbreak{}sourceLift} \quad WassersteinMixing/\allowbreak{}SemigroupLift.lean, lines 28–44}
\begin{leancode}
\leanline{0}{theorem\ langevin\_integral\_sourceLift\ (g\ :\ E\ →\ E)\ (L\ :\ ℝ≥0)\ (hg\ :\ LipschitzWith\ L\ g)}
\leanline{4}{(initial\ :\ Lp\ E\ 2\ P)\ (B\ Y\ :\ ℝ\ →\ Lp\ E\ 2\ P)\ (hY\ :\ Continuous\ Y)}
\leanline{4}{(T\ :\ ℝ)\ (heq\ :\ ∀\ t\ ∈\ Icc\ 0\ T,\ Y\ t\ -\ initial\ =}
\leanline{6}{-(∫\ s\ in\ 0..t,\ ULA.lpDrift\ P\ g\ hg\ (Y\ s))\ +\ B\ t)\ :}
\leanline{4}{∀\ t\ ∈\ Icc\ 0\ T,}
\leanline{6}{KineticLangevin.sourcePullback\ φ\ hφ\ (Y\ t)\ -}
\leanline{10}{KineticLangevin.sourcePullback\ φ\ hφ\ initial\ =}
\leanline{8}{-(∫\ s\ in\ 0..t,\ ULA.lpDrift\ Q\ g\ hg\ (KineticLangevin.sourcePullback\ φ\ hφ\ (Y\ s)))\ +}
\leanline{10}{KineticLangevin.sourcePullback\ φ\ hφ\ (B\ t)}
\end{leancode}
\leansource{\texttt{additiveEulerChain\_\allowbreak{}sourceLift} \quad WassersteinMixing/\allowbreak{}SemigroupLift.lean, lines 47–58}
\begin{leancode}
\leanline{0}{theorem\ additiveEulerChain\_sourceLift\ (g\ :\ E\ →\ E)\ (L\ :\ ℝ≥0)\ (hg\ :\ LipschitzWith\ L\ g)}
\leanline{4}{(step\ :\ ℝ)\ (initial\ :\ Lp\ E\ 2\ P)\ (B\ :\ ℝ\ →\ Lp\ E\ 2\ P)\ (n\ :\ ℕ)\ :}
\leanline{4}{KineticLangevin.sourcePullback\ φ\ hφ}
\leanline{6}{(ULA.additiveEulerChain\ (ULA.lpDrift\ P\ g\ hg)\ step\ initial\ B\ n)\ =}
\leanline{4}{ULA.additiveEulerChain\ (ULA.lpDrift\ Q\ g\ hg)\ step}
\leanline{6}{(KineticLangevin.sourcePullback\ φ\ hφ\ initial)}
\leanline{6}{(fun\ t\ ↦\ KineticLangevin.sourcePullback\ φ\ hφ\ (B\ t))\ n}
\end{leancode}
\leansource{\texttt{sourceLift\_\allowbreak{}law} \quad WassersteinMixing/\allowbreak{}SemigroupLift.lean, lines 65–67}
\begin{leancode}
\leanline{0}{theorem\ sourceLift\_law\ (X\ :\ Lp\ E\ 2\ P)\ :}
\leanline{4}{Q.map\ (KineticLangevin.sourcePullback\ φ\ hφ\ X)\ =\ P.map\ X}
\end{leancode}
\leansource{\texttt{stationaryEvaluation\_\allowbreak{}sourceLift} \quad WassersteinMixing/\allowbreak{}SemigroupLift.lean, lines 72–89}
\begin{leancode}
\leanline{0}{theorem\ stationaryEvaluation\_sourceLift\ (μ\ :\ Measure\ E)\ (X\ :\ Lp\ E\ 2\ P)}
\leanline{4}{(hX\ :\ P.map\ X\ =\ μ)\ (f\ :\ Lp\ F\ 2\ μ)\ :}
\leanline{4}{StationaryObservables.evaluation\ ({\leanfallback 𝕜}\ :=\ {\leanfallback 𝕜})\ μ\ Q}
\leanline{6}{(KineticLangevin.sourcePullback\ φ\ hφ\ X)\ ((sourceLift\_law\ φ\ hφ\ X).trans\ hX)\ f\ =}
\leanline{6}{Lp.compMeasurePreservingₗᵢ\ {\leanfallback 𝕜}\ φ\ hφ}
\leanline{8}{(StationaryObservables.evaluation\ ({\leanfallback 𝕜}\ :=\ {\leanfallback 𝕜})\ μ\ P\ X\ hX\ f)}
\end{leancode}
\leanlabel{How the Lean proof follows the argument}
\texttt{langevin\_\allowbreak{}integral\_\allowbreak{}source\allowbreak{}Lift} transports the actual integral equation through the pullback isometry. \texttt{additive\allowbreak{}Euler\allowbreak{}Chain\_\allowbreak{}source\allowbreak{}Lift} inducts on its recursion. \texttt{source\allowbreak{}Lift\_\allowbreak{}law} and \texttt{stationary\allowbreak{}Evaluation\_\allowbreak{}source\allowbreak{}Lift} use exact measure-preserving composition and representative equality.
\end{leanbox}


Earlier results used here: \cite[Section~3.13.2]{LangevinLean2026ULASC}; \cite[Section~3.13.5]{LangevinLean2026ULASC}; \cite[Section~3.16.1]{LangevinLean2026ULASC}; \cite[Section~3.16.3]{LangevinLean2026ULASC}; \cite[Section~3.14.1]{LangevinLean2026Kinetic}; \cite[Section~3.14.2]{LangevinLean2026Kinetic}.
\section{\texorpdfstring{Strong-limit identities for actual observation compressions}{Strong-limit identities for actual observation compressions}}
\label{proof-module-WassersteinMixing-SemigroupLimits-lean}


Common isometric source lifts do not change compression. Self-adjointness and composition identities pass through strong operator limits when the relevant approximants have uniform contraction bounds.

\subsection{\texorpdfstring{Isometric cancellation and strong-limit algebra}{Isometric cancellation and strong-limit algebra}}
\label{report-step-WassersteinMixing-SemigroupLimits-lean-1}


\textbf{Claim.} Replacing both observation maps by their common isometric lift leaves their adjoint compression unchanged. A strong limit of self-adjoint operators is self-adjoint. If two strongly convergent operator families are contractions, their compositions converge strongly to the composition of their limits.



\textbf{Proof.} The lift cancels because its adjoint times itself is the identity. For self-adjointness, pass to the limit in the two inner-product pairings against fixed vectors. For composition, split the error into the first contraction applied to the second orbit error and the first-family error evaluated at the fixed limiting vector. Both terms tend to zero. Uniqueness of limits then identifies any separately known composition limit.

\begin{leanbox}{Lean: \texttt{compression\_\allowbreak{}isometricLift}, \texttt{strongLimit\_\allowbreak{}isSelfAdjoint}, \texttt{strongLimit\_\allowbreak{}comp} and 1 more}
\leansource{\texttt{compression\_\allowbreak{}isometricLift} \quad WassersteinMixing/\allowbreak{}SemigroupLimits.lean, lines 19–30}
\begin{leancode}
\leanline{0}{theorem\ compression\_isometricLift\ \{{\leanfallback 𝕜}\ A\ B\ C\ :\ Type*\}\ [RCLike\ {\leanfallback 𝕜}]}
\leanline{4}{[NormedAddCommGroup\ A]\ [InnerProductSpace\ {\leanfallback 𝕜}\ A]\ [CompleteSpace\ A]}
\leanline{4}{[NormedAddCommGroup\ B]\ [InnerProductSpace\ {\leanfallback 𝕜}\ B]\ [CompleteSpace\ B]}
\leanline{4}{[NormedAddCommGroup\ C]\ [InnerProductSpace\ {\leanfallback 𝕜}\ C]\ [CompleteSpace\ C]}
\leanline{4}{(initial\ current\ :\ A\ →ₗᵢ[{\leanfallback 𝕜}]\ B)\ (lift\ :\ B\ →ₗᵢ[{\leanfallback 𝕜}]\ C)\ :}
\leanline{4}{ObservableTransition.compression\ (lift.comp\ initial)\ (lift.comp\ current)\ =}
\leanline{6}{ObservableTransition.compression\ initial\ current}
\end{leancode}
\leansource{\texttt{strongLimit\_\allowbreak{}isSelfAdjoint} \quad WassersteinMixing/\allowbreak{}SemigroupLimits.lean, lines 36–49}
\begin{leancode}
\leanline{0}{theorem\ strongLimit\_isSelfAdjoint\ [l.NeBot]}
\leanline{4}{(A\ :\ I\ →\ H\ →L[{\leanfallback 𝕜}]\ H)\ (T\ :\ H\ →L[{\leanfallback 𝕜}]\ H)}
\leanline{4}{(hA\ :\ ∀\ i,\ IsSelfAdjoint\ (A\ i))}
\leanline{4}{(hlim\ :\ ∀\ f,\ Tendsto\ (fun\ i\ ↦\ A\ i\ f)\ l\ ({\leanfallback 𝓝}\ (T\ f)))\ :\ IsSelfAdjoint\ T}
\end{leancode}
\leansource{\texttt{strongLimit\_\allowbreak{}comp} \quad WassersteinMixing/\allowbreak{}SemigroupLimits.lean, lines 54–69}
\begin{leancode}
\leanline{0}{theorem\ strongLimit\_comp}
\leanline{4}{(A\ B\ :\ I\ →\ H\ →L[{\leanfallback 𝕜}]\ H)\ (T\ S\ :\ H\ →L[{\leanfallback 𝕜}]\ H)}
\leanline{4}{(hcontract\ :\ ∀\ i\ f,\ ‖A\ i\ f‖\ ≤\ ‖f‖)}
\leanline{4}{(hA\ :\ ∀\ f,\ Tendsto\ (fun\ i\ ↦\ A\ i\ f)\ l\ ({\leanfallback 𝓝}\ (T\ f)))}
\leanline{4}{(hB\ :\ ∀\ f,\ Tendsto\ (fun\ i\ ↦\ B\ i\ f)\ l\ ({\leanfallback 𝓝}\ (S\ f)))\ (f\ :\ H)\ :}
\leanline{4}{Tendsto\ (fun\ i\ ↦\ A\ i\ (B\ i\ f))\ l\ ({\leanfallback 𝓝}\ (T\ (S\ f)))}
\end{leancode}
\leansource{\texttt{strongLimit\_\allowbreak{}comp\_\allowbreak{}eq} \quad WassersteinMixing/\allowbreak{}SemigroupLimits.lean, lines 74–85}
\begin{leancode}
\leanline{0}{theorem\ strongLimit\_comp\_eq\ [l.NeBot]}
\leanline{4}{(A\ B\ C\ :\ I\ →\ H\ →L[{\leanfallback 𝕜}]\ H)\ (T\ S\ R\ :\ H\ →L[{\leanfallback 𝕜}]\ H)}
\leanline{4}{(hcontract\ :\ ∀\ i\ f,\ ‖A\ i\ f‖\ ≤\ ‖f‖)}
\leanline{4}{(hcomp\ :\ ∀\ i,\ C\ i\ =\ A\ i\ ∘L\ B\ i)}
\leanline{4}{(hA\ :\ ∀\ f,\ Tendsto\ (fun\ i\ ↦\ A\ i\ f)\ l\ ({\leanfallback 𝓝}\ (T\ f)))}
\leanline{4}{(hB\ :\ ∀\ f,\ Tendsto\ (fun\ i\ ↦\ B\ i\ f)\ l\ ({\leanfallback 𝓝}\ (S\ f)))}
\leanline{4}{(hC\ :\ ∀\ f,\ Tendsto\ (fun\ i\ ↦\ C\ i\ f)\ l\ ({\leanfallback 𝓝}\ (R\ f)))\ :\ R\ =\ T\ ∘L\ S}
\end{leancode}
\leanlabel{How the Lean proof follows the argument}
\texttt{compression\_\allowbreak{}isometric\allowbreak{}Lift} uses the exact adjoint composition identities. \texttt{strong\allowbreak{}Limit\_\allowbreak{}is\allowbreak{}Self\allowbreak{}Adjoint} passes symmetric pairings through limits. \texttt{strong\allowbreak{}Limit\_\allowbreak{}comp} uses the contraction norm bound and a two-term difference; \texttt{strong\allowbreak{}Limit\_\allowbreak{}comp\_\allowbreak{}eq} applies uniqueness of the resulting vector limits.
\end{leanbox}
\section{\texorpdfstring{Actual MALA approximations converge to the Langevin operators}{Actual MALA approximations converge to the Langevin operators}}
\label{proof-module-WassersteinMixing-SemigroupApproximation-lean}


On the genuine Brownian/\allowbreak{}uniform extension, Euler converges strongly to the lifted actual diffusion, while the mean MALA/\allowbreak{}Euler discrepancy is bounded by accumulated rejections. Stationary observable continuity then upgrades actual MALA convergence in measure to strong convergence of all \(L^2\) transition orbits.

\subsection{\texorpdfstring{The actual lifted diffusion and Euler limit}{The actual lifted diffusion and Euler limit}}
\label{report-step-WassersteinMixing-SemigroupApproximation-lean-1}


\textbf{Claim.} The actual diffusion lifted to the uniform extension retains its Gibbs law and continuity. The Euler path from the same Brownian innovations converges strongly to it along any bounded moving mesh with \(r_i\to0\) and \(n_ir_i\to t\).



\textbf{Proof.} The source-lifting isometry preserves laws and continuity. Identify the actual Gaussian innovations with the normalized Brownian increments, then identify the Euler recursion with the actual additive-noise Euler construction. Apply its moving-step convergence theorem and lift the result through the isometry. The auxiliary mean-norm bound follows from Cauchy–Schwarz on a probability space.

\begin{leanbox}{Lean: \texttt{stateLift}, \texttt{diffusion}, \texttt{integral\_\allowbreak{}norm\_\allowbreak{}le\_\allowbreak{}lp\_\allowbreak{}norm} and 6 more}
\leansource{\texttt{stateLift} \quad WassersteinMixing/\allowbreak{}SemigroupApproximation.lean, lines 25–28}
\begin{leancode}
\leanline{0}{def\ stateLift\ (P\ :\ Parameters)\ (V\ :\ E\ P.dimension\ →\ ℝ)\ :}
\leanline{4}{Lp\ (E\ P.dimension)\ 2\ (actualProbability\ P\ V)\ →ₗᵢ[ℝ]}
\leanline{6}{Lp\ (E\ P.dimension)\ 2\ (probability\ P\ V)\ :=}
\leanline{2}{Lp.compMeasurePreservingₗᵢ\ ℝ\ Prod.fst\ measurePreserving\_fst}
\end{leancode}
\leansource{\texttt{diffusion} \quad WassersteinMixing/\allowbreak{}SemigroupApproximation.lean, lines 31–33}
\begin{leancode}
\leanline{0}{def\ diffusion\ (P\ :\ Parameters)\ (V\ :\ E\ P.dimension\ →\ ℝ)\ (hV\ :\ PotentialAssumptions\ P\ V)}
\leanline{4}{(t\ :\ ℝ)\ :\ Lp\ (E\ P.dimension)\ 2\ (probability\ P\ V)\ :=}
\leanline{2}{stateLift\ P\ V\ (CanonicalSemigroup.state\ P\ V\ hV\ t)}
\end{leancode}
\leansource{\texttt{integral\_\allowbreak{}norm\_\allowbreak{}le\_\allowbreak{}lp\_\allowbreak{}norm} \quad WassersteinMixing/\allowbreak{}SemigroupApproximation.lean, lines 35–46}
\begin{leancode}
\leanline{0}{private\ theorem\ integral\_norm\_le\_lp\_norm\ \{Ω\ H\ :\ Type*\}\ [MeasurableSpace\ Ω]}
\leanline{4}{[NormedAddCommGroup\ H]\ (Q\ :\ Measure\ Ω)\ [IsProbabilityMeasure\ Q]\ (Z\ :\ Lp\ H\ 2\ Q)\ :}
\leanline{4}{(∫\ ω,\ ‖Z\ ω‖\ ∂Q)\ ≤\ ‖Z‖}
\end{leancode}
\leansource{\texttt{diffusion\_\allowbreak{}law} \quad WassersteinMixing/\allowbreak{}SemigroupApproximation.lean, lines 48–55}
\begin{leancode}
\leanline{0}{theorem\ diffusion\_law\ (P\ :\ Parameters)\ (V\ :\ E\ P.dimension\ →\ ℝ)}
\leanline{4}{(hV\ :\ PotentialAssumptions\ P\ V)\ (t\ :\ ℝ)\ :}
\leanline{4}{(probability\ P\ V).map\ (diffusion\ P\ V\ hV\ t)\ =\ ULA.gibbsMeasure\ V}
\end{leancode}
\leansource{\texttt{continuous\_\allowbreak{}diffusion} \quad WassersteinMixing/\allowbreak{}SemigroupApproximation.lean, lines 57–59}
\begin{leancode}
\leanline{0}{theorem\ continuous\_diffusion\ (P\ :\ Parameters)\ (V\ :\ E\ P.dimension\ →\ ℝ)}
\leanline{4}{(hV\ :\ PotentialAssumptions\ P\ V)\ :\ Continuous\ (diffusion\ P\ V\ hV)}
\end{leancode}
\leansource{\texttt{gaussian\_\allowbreak{}eq\_\allowbreak{}brownian} \quad WassersteinMixing/\allowbreak{}SemigroupApproximation.lean, lines 62–71}
\begin{leancode}
\leanline{0}{theorem\ gaussian\_eq\_brownian\ (P\ :\ Parameters)\ (V\ :\ E\ P.dimension\ →\ ℝ)}
\leanline{4}{(hV\ :\ PotentialAssumptions\ P\ V)\ (r\ :\ ℝ)\ :}
\leanline{4}{gaussian\ P\ V\ hV\ r\ =\ ULA.Hilbert.brownianEulerNoise\ (actualProbability\ P\ V)}
\leanline{6}{(CanonicalSemigroup.brownian\ P)\ (CanonicalSemigroup.brownian\_preBrownian\ P\ V\ hV)\ r}
\end{leancode}
\leansource{\texttt{euler\_\allowbreak{}eq\_\allowbreak{}additive} \quad WassersteinMixing/\allowbreak{}SemigroupApproximation.lean, lines 75–88}
\begin{leancode}
\leanline{0}{theorem\ euler\_eq\_additive\ (P\ :\ Parameters)\ (V\ :\ E\ P.dimension\ →\ ℝ)}
\leanline{4}{(hV\ :\ PotentialAssumptions\ P\ V)\ (r\ :\ ℝ)\ (hr\ :\ 0\ <\ r)\ (n\ :\ ℕ)\ :}
\leanline{4}{letI\ :=\ actualProbability\_isProbability\ P\ V\ hV}
\leanline{4}{euler\ P\ V\ hV\ r\ n\ =\ ULA.additiveEulerChain}
\leanline{6}{(ULA.lpDrift\ (actualProbability\ P\ V)\ (gradient\ V)\ hV.gradientLipschitz)\ r}
\leanline{6}{(CanonicalSemigroup.initial\ P\ V\ hV)}
\leanline{6}{(ULA.brownianLangevinNoise\ (actualProbability\ P\ V)\ (CanonicalSemigroup.brownian\ P)}
\leanline{8}{(CanonicalSemigroup.brownian\_preBrownian\ P\ V\ hV)\ 0)\ n}
\end{leancode}
\leansource{\texttt{euler\_\allowbreak{}variable\_\allowbreak{}tendsto} \quad WassersteinMixing/\allowbreak{}SemigroupApproximation.lean, lines 92–104}
\begin{leancode}
\leanline{0}{theorem\ euler\_variable\_tendsto\ \{I\ :\ Type*\}\ \{l\ :\ Filter\ I\}}
\leanline{4}{(P\ :\ Parameters)\ (V\ :\ E\ P.dimension\ →\ ℝ)\ (hV\ :\ PotentialAssumptions\ P\ V)}
\leanline{4}{(step\ :\ I\ →\ ℝ)\ (count\ :\ I\ →\ ℕ)\ (hstep\ :\ ∀\ i,\ 0\ <\ step\ i)}
\leanline{4}{(hmesh\ :\ Tendsto\ step\ l\ ({\leanfallback 𝓝}\ 0))\ (T\ :\ ℝ)\ (hT\ :\ 0\ ≤\ T)}
\leanline{4}{(hbound\ :\ ∀\ i,\ (count\ i\ :\ ℝ)\ *\ step\ i\ ≤\ T)}
\leanline{4}{(t\ :\ ℝ)\ (htime\ :\ Tendsto\ (fun\ i\ ↦\ (count\ i\ :\ ℝ)\ *\ step\ i)\ l\ ({\leanfallback 𝓝}\ t))\ :}
\leanline{4}{Tendsto\ (fun\ i\ ↦\ euler\ P\ V\ hV\ (step\ i)\ (count\ i))\ l}
\leanline{6}{({\leanfallback 𝓝}\ (CanonicalSemigroup.state\ P\ V\ hV\ t))}
\end{leancode}
\leansource{\texttt{liftedEuler\_\allowbreak{}variable\_\allowbreak{}tendsto} \quad WassersteinMixing/\allowbreak{}SemigroupApproximation.lean, lines 107–116}
\begin{leancode}
\leanline{0}{theorem\ liftedEuler\_variable\_tendsto\ \{I\ :\ Type*\}\ \{l\ :\ Filter\ I\}}
\leanline{4}{(P\ :\ Parameters)\ (V\ :\ E\ P.dimension\ →\ ℝ)\ (hV\ :\ PotentialAssumptions\ P\ V)}
\leanline{4}{(step\ :\ I\ →\ ℝ)\ (count\ :\ I\ →\ ℕ)\ (hstep\ :\ ∀\ i,\ 0\ <\ step\ i)}
\leanline{4}{(hmesh\ :\ Tendsto\ step\ l\ ({\leanfallback 𝓝}\ 0))\ (T\ :\ ℝ)\ (hT\ :\ 0\ ≤\ T)}
\leanline{4}{(hbound\ :\ ∀\ i,\ (count\ i\ :\ ℝ)\ *\ step\ i\ ≤\ T)}
\leanline{4}{(t\ :\ ℝ)\ (htime\ :\ Tendsto\ (fun\ i\ ↦\ (count\ i\ :\ ℝ)\ *\ step\ i)\ l\ ({\leanfallback 𝓝}\ t))\ :}
\leanline{4}{Tendsto\ (fun\ i\ ↦\ stateLift\ P\ V\ (euler\ P\ V\ hV\ (step\ i)\ (count\ i)))\ l}
\leanline{6}{({\leanfallback 𝓝}\ (diffusion\ P\ V\ hV\ t))}
\end{leancode}
\leanlabel{How the Lean proof follows the argument}
\texttt{state\allowbreak{}Lift} is actual product-projection pullback and \texttt{diffusion} lifts the canonical global state. The Gaussian/\allowbreak{}Brownian and Euler/\allowbreak{}additive identities identify the concrete recursions before invoking moving-step Euler convergence.
\end{leanbox}


Earlier results used here: \cite[Section~3.13.2]{LangevinLean2026ULASC}; \cite[Section~3.16.1]{LangevinLean2026ULASC}; \cite[Section~3.25.1]{LangevinLean2026ULASC}; \cite[Section~3.26.3]{LangevinLean2026ULASC}; \cite[Section~3.29.7]{LangevinLean2026ULASC}; \cite[Section~3.32.13]{LangevinLean2026ULASC}; \cite[Section~3.42.5]{LangevinLean2026ULASC}; \cite[Section~3.8.17]{LangevinLean2026ULASC}; \cite[Section~3.8.19]{LangevinLean2026ULASC}; \cite[Section~3.8.20]{LangevinLean2026ULASC}; \cite[Section~3.14.1]{LangevinLean2026Kinetic}.
\subsection{\texorpdfstring{Rejections vanish on a bounded elapsed horizon}{Rejections vanish on a bounded elapsed horizon}}
\label{report-step-WassersteinMixing-SemigroupApproximation-lean-2}


\textbf{Claim.} The expected MALA/\allowbreak{}Euler distance is at most \(nCr\sqrt r\). For bounded elapsed time \(nr\le T\), it tends to zero as \(r\to0\), so the actual MALA endpoint converges in measure to the actual diffusion endpoint.



\textbf{Proof.} Transfer the proved cumulative rejection bound to the genuine canonical random coupling. Its upper bound is at most \(TC\sqrt r\). The actual Euler endpoint converges strongly, hence in mean and in measure, to the diffusion. Add the two mean-distance estimates and use the integral criterion for convergence in measure. Required small-step restrictions hold eventually as the mesh tends to zero.

\begin{leanbox}{Lean: \texttt{mala\_\allowbreak{}euler\_\allowbreak{}integral\_\allowbreak{}le}, \texttt{mala\_\allowbreak{}variable\_\allowbreak{}tendstoInMeasure}}
\leansource{\texttt{mala\_\allowbreak{}euler\_\allowbreak{}integral\_\allowbreak{}le} \quad WassersteinMixing/\allowbreak{}SemigroupApproximation.lean, lines 119–136}
\begin{leancode}
\leanline{0}{theorem\ mala\_euler\_integral\_le\ (P\ :\ Parameters)\ (V\ :\ E\ P.dimension\ →\ ℝ)}
\leanline{4}{(hV\ :\ PotentialAssumptions\ P\ V)\ (r\ :\ ℝ)\ (hr\ :\ 0\ <\ r)\ (hr1\ :\ r\ ≤\ 1)}
\leanline{4}{(hsmall\ :\ r\ ≤\ 2\ /\ (P.alpha\ +\ P.beta))\ (n\ :\ ℕ)\ :}
\leanline{4}{(∫\ ω,\ ‖mala\ P\ V\ hV\ r\ hr\ n\ ω\ -\ stateLift\ P\ V\ (euler\ P\ V\ hV\ r\ n)\ ω‖}
\leanline{6}{∂probability\ P\ V)\ ≤}
\leanline{6}{(n\ :\ ℝ)\ *\ ULA.malaRejectionConstant\ V\ ⟨P.beta,\ P.beta\_pos.le⟩\ *\ r\ *\ Real.sqrt\ r}
\end{leancode}
\leansource{\texttt{mala\_\allowbreak{}variable\_\allowbreak{}tendstoInMeasure} \quad WassersteinMixing/\allowbreak{}SemigroupApproximation.lean, lines 141–196}
\begin{leancode}
\leanline{0}{theorem\ mala\_variable\_tendstoInMeasure\ \{I\ :\ Type*\}\ \{l\ :\ Filter\ I\}}
\leanline{4}{(P\ :\ Parameters)\ (V\ :\ E\ P.dimension\ →\ ℝ)\ (hV\ :\ PotentialAssumptions\ P\ V)}
\leanline{4}{(step\ :\ I\ →\ ℝ)\ (count\ :\ I\ →\ ℕ)\ (hstep\ :\ ∀\ i,\ 0\ <\ step\ i)}
\leanline{4}{(hmesh\ :\ Tendsto\ step\ l\ ({\leanfallback 𝓝}\ 0))\ (T\ :\ ℝ)\ (hT\ :\ 0\ ≤\ T)}
\leanline{4}{(hbound\ :\ ∀\ i,\ (count\ i\ :\ ℝ)\ *\ step\ i\ ≤\ T)}
\leanline{4}{(t\ :\ ℝ)\ (htime\ :\ Tendsto\ (fun\ i\ ↦\ (count\ i\ :\ ℝ)\ *\ step\ i)\ l\ ({\leanfallback 𝓝}\ t))\ :}
\leanline{4}{TendstoInMeasure\ (probability\ P\ V)}
\leanline{6}{(fun\ i\ ω\ ↦\ mala\ P\ V\ hV\ (step\ i)\ (hstep\ i)\ (count\ i)\ ω)\ l}
\leanline{6}{(diffusion\ P\ V\ hV\ t)}
\end{leancode}
\leanlabel{How the Lean proof follows the argument}
\texttt{mala\_\allowbreak{}euler\_\allowbreak{}integral\_\allowbreak{}le} specializes the actual random-coupling estimate. \texttt{mala\_\allowbreak{}variable\_\allowbreak{}tendsto\allowbreak{}In\allowbreak{}Measure} proves eventual step restrictions, controls cumulative rejections by the elapsed horizon, combines the Euler limit and applies the mean-norm convergence criterion.
\end{leanbox}


Earlier results used here: \cite[Section~3.34.22]{LangevinLean2026ULASC}; \cite[Section~3.34.23]{LangevinLean2026ULASC}.
\subsection{\texorpdfstring{Strong convergence of the actual observable operators}{Strong convergence of the actual observable operators}}
\label{report-step-WassersteinMixing-SemigroupApproximation-lean-3}


\textbf{Claim.} For every square-integrable Gibbs observable \(f\), actual MALA evaluations converge in \(L^2\) to diffusion evaluation; their compressed operators converge strongly to \(P_tf\).



\textbf{Proof.} Every MALA endpoint has the exact Gibbs law. The actual diffusion limit has that same law, so stationary-observable convergence upgrades convergence in measure of states to strong convergence of every evaluated observable. Diffusion evaluation on the extension is exactly the source-lift of original evaluation. The common lift is an isometry, so compression removes it. Boundedness of the fixed initial adjoint then passes evaluation convergence to the actual compressed operators.

\begin{leanbox}{Lean: \texttt{malaEvaluation}, \texttt{diffusionEvaluation}, \texttt{malaEvaluation\_\allowbreak{}variable\_\allowbreak{}tendsto} and 4 more}
\leansource{\texttt{malaEvaluation} \quad WassersteinMixing/\allowbreak{}SemigroupApproximation.lean, lines 202–205}
\begin{leancode}
\leanline{0}{def\ malaEvaluation\ (P\ :\ Parameters)\ (V\ :\ E\ P.dimension\ →\ ℝ)}
\leanline{4}{(hV\ :\ PotentialAssumptions\ P\ V)\ (r\ :\ ℝ)\ (hr\ :\ 0\ <\ r)\ (n\ :\ ℕ)\ :}
\leanline{4}{Lp\ F\ 2\ (ULA.gibbsMeasure\ V)\ →ₗᵢ[{\leanfallback 𝕜}]\ Lp\ F\ 2\ (probability\ P\ V)\ :=}
\leanline{2}{StationaryObservables.evaluation\ \_\ \_\ (mala\ P\ V\ hV\ r\ hr\ n)\ (mala\_law\ P\ V\ hV\ r\ hr\ n)}
\end{leancode}
\leansource{\texttt{diffusionEvaluation} \quad WassersteinMixing/\allowbreak{}SemigroupApproximation.lean, lines 208–211}
\begin{leancode}
\leanline{0}{def\ diffusionEvaluation\ (P\ :\ Parameters)\ (V\ :\ E\ P.dimension\ →\ ℝ)}
\leanline{4}{(hV\ :\ PotentialAssumptions\ P\ V)\ (t\ :\ ℝ)\ :}
\leanline{4}{Lp\ F\ 2\ (ULA.gibbsMeasure\ V)\ →ₗᵢ[{\leanfallback 𝕜}]\ Lp\ F\ 2\ (probability\ P\ V)\ :=}
\leanline{2}{StationaryObservables.evaluation\ \_\ \_\ (diffusion\ P\ V\ hV\ t)\ (diffusion\_law\ P\ V\ hV\ t)}
\end{leancode}
\leansource{\texttt{malaEvaluation\_\allowbreak{}variable\_\allowbreak{}tendsto} \quad WassersteinMixing/\allowbreak{}SemigroupApproximation.lean, lines 216–231}
\begin{leancode}
\leanline{0}{theorem\ malaEvaluation\_variable\_tendsto\ [NormedSpace\ ℝ\ F]}
\leanline{4}{\{I\ :\ Type*\}\ \{l\ :\ Filter\ I\}\ [l.IsCountablyGenerated]}
\leanline{4}{(P\ :\ Parameters)\ (V\ :\ E\ P.dimension\ →\ ℝ)\ (hV\ :\ PotentialAssumptions\ P\ V)}
\leanline{4}{(step\ :\ I\ →\ ℝ)\ (count\ :\ I\ →\ ℕ)\ (hstep\ :\ ∀\ i,\ 0\ <\ step\ i)}
\leanline{4}{(hmesh\ :\ Tendsto\ step\ l\ ({\leanfallback 𝓝}\ 0))\ (T\ :\ ℝ)\ (hT\ :\ 0\ ≤\ T)}
\leanline{4}{(hbound\ :\ ∀\ i,\ (count\ i\ :\ ℝ)\ *\ step\ i\ ≤\ T)}
\leanline{4}{(t\ :\ ℝ)\ (htime\ :\ Tendsto\ (fun\ i\ ↦\ (count\ i\ :\ ℝ)\ *\ step\ i)\ l\ ({\leanfallback 𝓝}\ t))}
\leanline{4}{(f\ :\ Lp\ F\ 2\ (ULA.gibbsMeasure\ V))\ :}
\leanline{4}{Tendsto\ (fun\ i\ ↦\ malaEvaluation\ ({\leanfallback 𝕜}\ :=\ {\leanfallback 𝕜})\ P\ V\ hV\ (step\ i)\ (hstep\ i)\ (count\ i)\ f)\ l}
\leanline{6}{({\leanfallback 𝓝}\ (diffusionEvaluation\ ({\leanfallback 𝕜}\ :=\ {\leanfallback 𝕜})\ P\ V\ hV\ t\ f))}
\end{leancode}
\leansource{\texttt{diffusionEvaluation\_\allowbreak{}eq\_\allowbreak{}lift} \quad WassersteinMixing/\allowbreak{}SemigroupApproximation.lean, lines 235–247}
\begin{leancode}
\leanline{0}{theorem\ diffusionEvaluation\_eq\_lift\ (P\ :\ Parameters)\ (V\ :\ E\ P.dimension\ →\ ℝ)}
\leanline{4}{(hV\ :\ PotentialAssumptions\ P\ V)\ (t\ :\ ℝ)\ :}
\leanline{4}{diffusionEvaluation\ ({\leanfallback 𝕜}\ :=\ {\leanfallback 𝕜})\ (F\ :=\ F)\ P\ V\ hV\ t\ =}
\leanline{6}{(Lp.compMeasurePreservingₗᵢ\ {\leanfallback 𝕜}\ Prod.fst}
\leanline{8}{(measurePreserving\_fst\ (μ\ :=\ actualProbability\ P\ V)\ (ν\ :=\ uniformPath))).comp}
\leanline{8}{(CanonicalSemigroup.evaluation\ P\ V\ hV\ t)}
\end{leancode}
\leansource{\texttt{diffusion\_\allowbreak{}compression} \quad WassersteinMixing/\allowbreak{}SemigroupApproximation.lean, lines 251–256}
\begin{leancode}
\leanline{0}{theorem\ diffusion\_compression\ (P\ :\ Parameters)\ (V\ :\ E\ P.dimension\ →\ ℝ)}
\leanline{4}{(hV\ :\ PotentialAssumptions\ P\ V)\ (t\ :\ ℝ)\ :}
\leanline{4}{ObservableTransition.compression\ (diffusionEvaluation\ ({\leanfallback 𝕜}\ :=\ {\leanfallback 𝕜})\ (F\ :=\ F)\ P\ V\ hV\ 0)}
\leanline{6}{(diffusionEvaluation\ P\ V\ hV\ t)\ =\ CanonicalSemigroup.op\ P\ V\ hV\ t}
\end{leancode}
\leansource{\texttt{malaOp} \quad WassersteinMixing/\allowbreak{}SemigroupApproximation.lean, lines 259–263}
\begin{leancode}
\leanline{0}{def\ malaOp\ (P\ :\ Parameters)\ (V\ :\ E\ P.dimension\ →\ ℝ)}
\leanline{4}{(hV\ :\ PotentialAssumptions\ P\ V)\ (r\ :\ ℝ)\ (hr\ :\ 0\ <\ r)\ (n\ :\ ℕ)\ :}
\leanline{4}{Lp\ F\ 2\ (ULA.gibbsMeasure\ V)\ →L[{\leanfallback 𝕜}]\ Lp\ F\ 2\ (ULA.gibbsMeasure\ V)\ :=}
\leanline{2}{ObservableTransition.compression\ (diffusionEvaluation\ P\ V\ hV\ 0)}
\leanline{4}{(malaEvaluation\ P\ V\ hV\ r\ hr\ n)}
\end{leancode}
\leansource{\texttt{malaOp\_\allowbreak{}variable\_\allowbreak{}tendsto} \quad WassersteinMixing/\allowbreak{}SemigroupApproximation.lean, lines 266–279}
\begin{leancode}
\leanline{0}{theorem\ malaOp\_variable\_tendsto\ [NormedSpace\ ℝ\ F]}
\leanline{4}{\{I\ :\ Type*\}\ \{l\ :\ Filter\ I\}\ [l.IsCountablyGenerated]}
\leanline{4}{(P\ :\ Parameters)\ (V\ :\ E\ P.dimension\ →\ ℝ)\ (hV\ :\ PotentialAssumptions\ P\ V)}
\leanline{4}{(step\ :\ I\ →\ ℝ)\ (count\ :\ I\ →\ ℕ)\ (hstep\ :\ ∀\ i,\ 0\ <\ step\ i)}
\leanline{4}{(hmesh\ :\ Tendsto\ step\ l\ ({\leanfallback 𝓝}\ 0))\ (T\ :\ ℝ)\ (hT\ :\ 0\ ≤\ T)}
\leanline{4}{(hbound\ :\ ∀\ i,\ (count\ i\ :\ ℝ)\ *\ step\ i\ ≤\ T)}
\leanline{4}{(t\ :\ ℝ)\ (htime\ :\ Tendsto\ (fun\ i\ ↦\ (count\ i\ :\ ℝ)\ *\ step\ i)\ l\ ({\leanfallback 𝓝}\ t))}
\leanline{4}{(f\ :\ Lp\ F\ 2\ (ULA.gibbsMeasure\ V))\ :}
\leanline{4}{Tendsto\ (fun\ i\ ↦\ malaOp\ ({\leanfallback 𝕜}\ :=\ {\leanfallback 𝕜})\ P\ V\ hV\ (step\ i)\ (hstep\ i)\ (count\ i)\ f)\ l}
\leanline{6}{({\leanfallback 𝓝}\ (CanonicalSemigroup.op\ ({\leanfallback 𝕜}\ :=\ {\leanfallback 𝕜})\ P\ V\ hV\ t\ f))}
\end{leancode}
\leanlabel{How the Lean proof follows the argument}
\texttt{mala\allowbreak{}Evaluation} and \texttt{diffusion\allowbreak{}Evaluation} are actual stationary pullbacks. The evaluation-limit theorem applies the proved stationary convergence result. \texttt{diffusion\allowbreak{}Evaluation\_\allowbreak{}eq\_\allowbreak{}lift} and \texttt{diffusion\_\allowbreak{}compression} remove the independent extension. \texttt{mala\allowbreak{}Op\_\allowbreak{}variable\_\allowbreak{}tendsto} passes the limit through the fixed adjoint compression.
\end{leanbox}


Earlier results used here: \cite[Section~3.26.3]{LangevinLean2026ULASC}.
\section{\texorpdfstring{Retaining the actual initial coordinate through fresh updates}{Retaining the actual initial coordinate through fresh updates}}
\label{proof-module-WassersteinMixing-SemigroupRandomMALAAnchors-lean}


A sampled MALA update is identified with its explicit transition kernel. Extending the state by an unchanged initial coordinate yields an actual anchored recursion, whose law advances by the identity kernel on the anchor and the actual MALA kernel on the endpoint.

\subsection{\texorpdfstring{Actual sampled updates are kernel composition}{Actual sampled updates are kernel composition}}
\label{report-step-WassersteinMixing-SemigroupRandomMALAAnchors-lean-1}


\textbf{Claim.} Sampling an independent innovation and applying a measurable update is exactly composition with its sampled transition kernel. For the scalar MALA update, this kernel is the explicit Gaussian accepted/\allowbreak{}rejected kernel.



\textbf{Proof.} Product integration identifies the sampled Markov law with integration of the output kernel over the current state law. For MALA, integrate the independent uniform threshold exactly; its weighted Gaussian output agrees with the explicit transition formula.

\begin{leanbox}{Lean: \texttt{sampled\_\allowbreak{}markovStep\_\allowbreak{}eq\_\allowbreak{}comp}, \texttt{malaUpdate}, \texttt{malaUpdate\_\allowbreak{}measurable} and 1 more}
\leansource{\texttt{sampled\_\allowbreak{}markovStep\_\allowbreak{}eq\_\allowbreak{}comp} \quad WassersteinMixing/\allowbreak{}SemigroupRandomMALAAnchors.lean, lines 23–35}
\begin{leancode}
\leanline{0}{theorem\ sampled\_markovStep\_eq\_comp\ (F\ :\ A\ →\ Z\ →\ A)}
\leanline{4}{(hF\ :\ Measurable\ (Function.uncurry\ F))\ (η\ :\ Measure\ Z)\ [IsProbabilityMeasure\ η]}
\leanline{4}{(κ\ :\ Kernel\ A\ A)\ (hκ\ :\ ∀\ x,\ κ\ x\ =\ η.map\ (F\ x))}
\leanline{4}{(μ\ :\ Measure\ A)\ [IsProbabilityMeasure\ μ]\ :}
\leanline{4}{ULA.markovStep\ F\ η\ μ\ =\ κ\ ∘ₘ\ μ}
\end{leancode}
\leansource{\texttt{malaUpdate} \quad WassersteinMixing/\allowbreak{}SemigroupRandomMALAAnchors.lean, lines 43–44}
\begin{leancode}
\leanline{0}{def\ malaUpdate\ (V\ :\ E\ →\ ℝ)\ (h\ :\ ℝ)\ (x\ :\ E)\ (z\ :\ E\ ×\ I)\ :\ E\ :=}
\leanline{2}{(update\ V\ h\ (x,\ x)\ z).2}
\end{leancode}
\leansource{\texttt{malaUpdate\_\allowbreak{}measurable} \quad WassersteinMixing/\allowbreak{}SemigroupRandomMALAAnchors.lean, lines 46–50}
\begin{leancode}
\leanline{0}{theorem\ malaUpdate\_measurable\ (V\ :\ E\ →\ ℝ)\ (h\ :\ ℝ)\ (L\ :\ ℝ≥0)}
\leanline{4}{(hdiff\ :\ Differentiable\ ℝ\ V)\ (hLip\ :\ LipschitzWith\ L\ (gradient\ V))\ :}
\leanline{4}{Measurable\ (Function.uncurry\ (malaUpdate\ V\ h))}
\end{leancode}
\leansource{\texttt{malaUpdate\_\allowbreak{}law} \quad WassersteinMixing/\allowbreak{}SemigroupRandomMALAAnchors.lean, lines 53–69}
\begin{leancode}
\leanline{0}{theorem\ malaUpdate\_law\ (V\ :\ E\ →\ ℝ)\ (h\ :\ ℝ)\ (L\ :\ ℝ≥0)}
\leanline{4}{(hdiff\ :\ Differentiable\ ℝ\ V)\ (hLip\ :\ LipschitzWith\ L\ (gradient\ V))\ (x\ :\ E)\ :}
\leanline{4}{(innovationLaw\ (E\ :=\ E)).map\ (malaUpdate\ V\ h\ x)\ =}
\leanline{6}{MALAStationaryPair.explicitKernel\ V\ h\ x}
\end{leancode}
\leanlabel{How the Lean proof follows the argument}
\texttt{sampled\_\allowbreak{}markov\allowbreak{}Step\_\allowbreak{}eq\_\allowbreak{}comp} proves the generic equality by integration on measurable events. \texttt{mala\allowbreak{}Update} is the actual scalar uniform-threshold update; \texttt{mala\allowbreak{}Update\_\allowbreak{}law} applies the exact uniform-decision law to identify its kernel.
\end{leanbox}


Earlier results used here: \cite[Section~3.33.5]{LangevinLean2026ULASC}; \cite[Section~3.36.12]{LangevinLean2026ULASC}; \cite[Section~3.36.9]{LangevinLean2026ULASC}.
\subsection{\texorpdfstring{The actual anchored update and recursion}{The actual anchored update and recursion}}
\label{report-step-WassersteinMixing-SemigroupRandomMALAAnchors-lean-2}


\textbf{Claim.} Updating \((x_0,x)\) to \((x_0,\operatorname{MALA}(x,\xi,U))\) preserves the anchor and has kernel \(I\otimes K\). The resulting deterministic recursion equals the pair of its original initial value and the scalar sampled MALA recursion.



\textbf{Proof.} The first output coordinate is deterministic and the second has the proved explicit MALA law, so the sampled pair law is the Dirac anchor times that transition. Integrating over the current pair gives the parallel-kernel action. Inducting the deterministic recursion proves the exact anchor/\allowbreak{}endpoint identity.

\begin{leanbox}{Lean: \texttt{anchoredUpdate}, \texttt{anchoredUpdate\_\allowbreak{}measurable}, \texttt{anchoredUpdate\_\allowbreak{}law} and 2 more}
\leansource{\texttt{anchoredUpdate} \quad WassersteinMixing/\allowbreak{}SemigroupRandomMALAAnchors.lean, lines 72–73}
\begin{leancode}
\leanline{0}{def\ anchoredUpdate\ (V\ :\ E\ →\ ℝ)\ (h\ :\ ℝ)\ (p\ :\ E\ ×\ E)\ (z\ :\ E\ ×\ I)\ :\ E\ ×\ E\ :=}
\leanline{2}{(p.1,\ malaUpdate\ V\ h\ p.2\ z)}
\end{leancode}
\leansource{\texttt{anchoredUpdate\_\allowbreak{}measurable} \quad WassersteinMixing/\allowbreak{}SemigroupRandomMALAAnchors.lean, lines 75–79}
\begin{leancode}
\leanline{0}{theorem\ anchoredUpdate\_measurable\ (V\ :\ E\ →\ ℝ)\ (h\ :\ ℝ)\ (L\ :\ ℝ≥0)}
\leanline{4}{(hdiff\ :\ Differentiable\ ℝ\ V)\ (hLip\ :\ LipschitzWith\ L\ (gradient\ V))\ :}
\leanline{4}{Measurable\ (Function.uncurry\ (anchoredUpdate\ V\ h))}
\end{leancode}
\leansource{\texttt{anchoredUpdate\_\allowbreak{}law} \quad WassersteinMixing/\allowbreak{}SemigroupRandomMALAAnchors.lean, lines 81–91}
\begin{leancode}
\leanline{0}{theorem\ anchoredUpdate\_law\ (V\ :\ E\ →\ ℝ)\ (h\ :\ ℝ)\ (L\ :\ ℝ≥0)}
\leanline{4}{(hdiff\ :\ Differentiable\ ℝ\ V)\ (hLip\ :\ LipschitzWith\ L\ (gradient\ V))\ (p\ :\ E\ ×\ E)\ :}
\leanline{4}{(innovationLaw\ (E\ :=\ E)).map\ (anchoredUpdate\ V\ h\ p)\ =}
\leanline{6}{(Kernel.id\ {\leanfallback ∥}ₖ\ MALAStationaryPair.explicitKernel\ V\ h)\ p}
\end{leancode}
\leansource{\texttt{anchored\_\allowbreak{}markovStep} \quad WassersteinMixing/\allowbreak{}SemigroupRandomMALAAnchors.lean, lines 93–99}
\begin{leancode}
\leanline{0}{theorem\ anchored\_markovStep\ (V\ :\ E\ →\ ℝ)\ (h\ :\ ℝ)\ (L\ :\ ℝ≥0)}
\leanline{4}{(hdiff\ :\ Differentiable\ ℝ\ V)\ (hLip\ :\ LipschitzWith\ L\ (gradient\ V))}
\leanline{4}{(ρ\ :\ Measure\ (E\ ×\ E))\ [IsProbabilityMeasure\ ρ]\ :}
\leanline{4}{ULA.markovStep\ (anchoredUpdate\ V\ h)\ innovationLaw\ ρ\ =}
\leanline{6}{(Kernel.id\ {\leanfallback ∥}ₖ\ MALAStationaryPair.explicitKernel\ V\ h)\ ∘ₘ\ ρ}
\end{leancode}
\leansource{\texttt{anchored\_\allowbreak{}chain\_\allowbreak{}eq} \quad WassersteinMixing/\allowbreak{}SemigroupRandomMALAAnchors.lean, lines 103–112}
\begin{leancode}
\leanline{0}{theorem\ anchored\_chain\_eq\ (V\ :\ E\ →\ ℝ)\ (h\ :\ ℝ)\ (x\ :\ E)\ (w\ :\ ℕ\ →\ E\ ×\ I)\ (n\ :\ ℕ)\ :}
\leanline{4}{ULA.noiseDrivenChain\ (anchoredUpdate\ V\ h)\ (x,\ x)\ w\ n\ =}
\leanline{6}{(x,\ (ULA.noiseDrivenChain\ (update\ V\ h)\ (x,\ x)\ w\ n).2)}
\end{leancode}
\leanlabel{How the Lean proof follows the argument}
\texttt{anchored\allowbreak{}Update\_\allowbreak{}law} maps the exact scalar update law while keeping the first coordinate. \texttt{anchored\_\allowbreak{}markov\allowbreak{}Step} applies the generic sampled composition result. \texttt{anchored\_\allowbreak{}chain\_\allowbreak{}eq} is induction on the actual update recursion.
\end{leanbox}


Earlier results used here: \cite[Section~3.20.1]{LangevinLean2026ULASC}; \cite[Section~3.20.3]{LangevinLean2026ULASC}; \cite[Section~3.33.5]{LangevinLean2026ULASC}.
\section{\texorpdfstring{The fresh-noise anchored MALA pair and its actual powers}{The fresh-noise anchored MALA pair and its actual powers}}
\label{proof-module-WassersteinMixing-SemigroupMALAPair-lean}


The fresh-noise pair retains the current state next to the actual accepted or rejected output. Its exact kernel identity links the concrete random update to the already proved reversible Gaussian pair, and hence to reversible anchored finite-step laws.

\subsection{\texorpdfstring{Identifying the actual sampled one-step pair}{Identifying the actual sampled one-step pair}}
\label{report-step-WassersteinMixing-SemigroupMALAPair-lean-1}


\textbf{Claim.} The fresh pair kernel is \(x\mapsto\delta_x\otimes K(x,\cdot)\), where \(K\) is the actual explicit MALA kernel. For any initial law, its integrated pair is the accepted Metropolis flow plus rejected diagonal mass. In particular \(\pi\otimes K\) is exactly the proved reversible Gibbs pair.



\textbf{Proof.} Direct accepted/\allowbreak{}rejected partition gives first marginal \(\delta_x\); mapping the second coordinate gives the actual explicit kernel. A pair whose first marginal is a Dirac measure has deterministic first coordinate almost surely, so it equals the corresponding product with its second marginal. Integrate this identity. Moving the actual proposal map through the weighted measures identifies the result with the Gaussian Metropolis pair construction.

\begin{leanbox}{Lean: \texttt{freshJointKernel}, \texttt{freshJointMeasure}, \texttt{freshJointKernel\_\allowbreak{}comp} and 8 more}
\leansource{\texttt{freshJointKernel} \quad WassersteinMixing/\allowbreak{}SemigroupMALAPair.lean, lines 26–32}
\begin{leancode}
\leanline{0}{def\ freshJointKernel\ (V\ :\ E\ →\ ℝ)\ (h\ :\ ℝ)\ :\ Kernel\ E\ (E\ ×\ E)\ :=}
\leanline{2}{((gaussianPairKernel\ (E\ :=\ E)).withDensity}
\leanline{6}{(fun\ \_\ p\ ↦\ ENNReal.ofReal\ (ULA.malaProposalAcceptance\ V\ h\ p))).map}
\leanline{8}{(fun\ p\ ↦\ (p.1,\ ULA.malaProposal\ V\ h\ p.1\ p.2))\ +}
\leanline{2}{((gaussianPairKernel\ (E\ :=\ E)).withDensity}
\leanline{6}{(fun\ \_\ p\ ↦\ ENNReal.ofReal\ (1\ -\ ULA.malaProposalAcceptance\ V\ h\ p))).map}
\leanline{8}{(fun\ p\ ↦\ (p.1,\ p.1))}
\end{leancode}
\leansource{\texttt{freshJointMeasure} \quad WassersteinMixing/\allowbreak{}SemigroupMALAPair.lean, lines 35–41}
\begin{leancode}
\leanline{0}{def\ freshJointMeasure\ (V\ :\ E\ →\ ℝ)\ (h\ :\ ℝ)\ (μ\ :\ Measure\ E)\ :\ Measure\ (E\ ×\ E)\ :=}
\leanline{2}{((μ.prod\ (stdGaussian\ E)).withDensity}
\leanline{6}{(fun\ p\ ↦\ ENNReal.ofReal\ (ULA.malaProposalAcceptance\ V\ h\ p))).map}
\leanline{8}{(fun\ p\ ↦\ (p.1,\ ULA.malaProposal\ V\ h\ p.1\ p.2))\ +}
\leanline{2}{((μ.prod\ (stdGaussian\ E)).withDensity}
\leanline{6}{(fun\ p\ ↦\ ENNReal.ofReal\ (1\ -\ ULA.malaProposalAcceptance\ V\ h\ p))).map}
\leanline{8}{(fun\ p\ ↦\ (p.1,\ p.1))}
\end{leancode}
\leansource{\texttt{freshJointKernel\_\allowbreak{}comp} \quad WassersteinMixing/\allowbreak{}SemigroupMALAPair.lean, lines 43–55}
\begin{leancode}
\leanline{0}{theorem\ freshJointKernel\_comp\ (V\ :\ E\ →\ ℝ)\ (h\ :\ ℝ)\ (L\ :\ ℝ≥0)}
\leanline{4}{(hdiff\ :\ Differentiable\ ℝ\ V)\ (hLip\ :\ LipschitzWith\ L\ (gradient\ V))}
\leanline{4}{(μ\ :\ Measure\ E)\ [SFinite\ μ]\ :}
\leanline{4}{freshJointKernel\ V\ h\ ∘ₘ\ μ\ =\ freshJointMeasure\ V\ h\ μ}
\end{leancode}
\leansource{\texttt{freshJointMeasure\_\allowbreak{}fst} \quad WassersteinMixing/\allowbreak{}SemigroupMALAPair.lean, lines 58–79}
\begin{leancode}
\leanline{0}{theorem\ freshJointMeasure\_fst\ (V\ :\ E\ →\ ℝ)\ (h\ :\ ℝ)\ (L\ :\ ℝ≥0)}
\leanline{4}{(hdiff\ :\ Differentiable\ ℝ\ V)\ (hLip\ :\ LipschitzWith\ L\ (gradient\ V))\ (μ\ :\ Measure\ E)\ :}
\leanline{4}{(freshJointMeasure\ V\ h\ μ).map\ Prod.fst\ =\ μ}
\end{leancode}
\leansource{\texttt{freshJointKernel\_\allowbreak{}fst} \quad WassersteinMixing/\allowbreak{}SemigroupMALAPair.lean, lines 82–90}
\begin{leancode}
\leanline{0}{theorem\ freshJointKernel\_fst\ (V\ :\ E\ →\ ℝ)\ (h\ :\ ℝ)\ (L\ :\ ℝ≥0)}
\leanline{4}{(hdiff\ :\ Differentiable\ ℝ\ V)\ (hLip\ :\ LipschitzWith\ L\ (gradient\ V))\ (x\ :\ E)\ :}
\leanline{4}{(freshJointKernel\ V\ h\ x).map\ Prod.fst\ =\ Measure.dirac\ x}
\end{leancode}
\leansource{\texttt{freshJointKernel\_\allowbreak{}snd} \quad WassersteinMixing/\allowbreak{}SemigroupMALAPair.lean, lines 93–104}
\begin{leancode}
\leanline{0}{theorem\ freshJointKernel\_snd\ (V\ :\ E\ →\ ℝ)\ (h\ :\ ℝ)\ (L\ :\ ℝ≥0)}
\leanline{4}{(hLip\ :\ LipschitzWith\ L\ (gradient\ V))\ (x\ :\ E)\ :}
\leanline{4}{(freshJointKernel\ V\ h\ x).map\ Prod.snd\ =\ explicitKernel\ V\ h\ x}
\end{leancode}
\leansource{\texttt{measure\_\allowbreak{}eq\_\allowbreak{}dirac\_\allowbreak{}prod} \quad WassersteinMixing/\allowbreak{}SemigroupMALAPair.lean, lines 107–128}
\begin{leancode}
\leanline{0}{theorem\ measure\_eq\_dirac\_prod\ (J\ :\ Measure\ (E\ ×\ E))\ (x\ :\ E)}
\leanline{4}{(hfst\ :\ J.map\ Prod.fst\ =\ Measure.dirac\ x)\ :}
\leanline{4}{J\ =\ (Measure.dirac\ x).prod\ (J.map\ Prod.snd)}
\end{leancode}
\leansource{\texttt{freshJointKernel\_\allowbreak{}eq\_\allowbreak{}id\_\allowbreak{}prod} \quad WassersteinMixing/\allowbreak{}SemigroupMALAPair.lean, lines 131–138}
\begin{leancode}
\leanline{0}{theorem\ freshJointKernel\_eq\_id\_prod\ (V\ :\ E\ →\ ℝ)\ (h\ :\ ℝ)\ (L\ :\ ℝ≥0)}
\leanline{4}{(hdiff\ :\ Differentiable\ ℝ\ V)\ (hLip\ :\ LipschitzWith\ L\ (gradient\ V))\ :}
\leanline{4}{freshJointKernel\ V\ h\ =\ Kernel.id\ ×ₖ\ explicitKernel\ V\ h}
\end{leancode}
\leansource{\texttt{freshJointMeasure\_\allowbreak{}eq\_\allowbreak{}metropolis} \quad WassersteinMixing/\allowbreak{}SemigroupMALAPair.lean, lines 141–164}
\begin{leancode}
\leanline{0}{theorem\ freshJointMeasure\_eq\_metropolis\ (V\ :\ E\ →\ ℝ)\ (h\ :\ ℝ)\ (L\ :\ ℝ≥0)}
\leanline{4}{(hdiff\ :\ Differentiable\ ℝ\ V)\ (hLip\ :\ LipschitzWith\ L\ (gradient\ V))\ (μ\ :\ Measure\ E)\ :}
\leanline{4}{freshJointMeasure\ V\ h\ μ\ =}
\leanline{6}{metropolisJoint\ (ULA.malaProposalJoint\ V\ h\ μ)}
\leanline{8}{(fun\ p\ ↦\ ULA.malaAcceptance\ V\ h\ p.1\ p.2)}
\end{leancode}
\leansource{\texttt{explicitKernel\_\allowbreak{}pair\_\allowbreak{}eq} \quad WassersteinMixing/\allowbreak{}SemigroupMALAPair.lean, lines 167–177}
\begin{leancode}
\leanline{0}{theorem\ explicitKernel\_pair\_eq\ (V\ :\ E\ →\ ℝ)\ (h\ :\ ℝ)\ (L\ :\ ℝ≥0)}
\leanline{4}{(hdiff\ :\ Differentiable\ ℝ\ V)\ (hLip\ :\ LipschitzWith\ L\ (gradient\ V))}
\leanline{4}{(μ\ :\ Measure\ E)\ [SFinite\ μ]\ :}
\leanline{4}{μ\ ⊗ₘ\ explicitKernel\ V\ h\ =}
\leanline{6}{metropolisJoint\ (ULA.malaProposalJoint\ V\ h\ μ)}
\leanline{8}{(fun\ p\ ↦\ ULA.malaAcceptance\ V\ h\ p.1\ p.2)}
\end{leancode}
\leansource{\texttt{explicitKernel\_\allowbreak{}gibbs\_\allowbreak{}pair\_\allowbreak{}eq} \quad WassersteinMixing/\allowbreak{}SemigroupMALAPair.lean, lines 180–184}
\begin{leancode}
\leanline{0}{theorem\ explicitKernel\_gibbs\_pair\_eq\ (V\ :\ E\ →\ ℝ)\ (h\ :\ ℝ)\ (L\ :\ ℝ≥0)}
\leanline{4}{[SFinite\ (ULA.gibbsMeasure\ V)]}
\leanline{4}{(hdiff\ :\ Differentiable\ ℝ\ V)\ (hLip\ :\ LipschitzWith\ L\ (gradient\ V))\ :}
\leanline{4}{ULA.gibbsMeasure\ V\ ⊗ₘ\ explicitKernel\ V\ h\ =\ jointLaw\ V\ h}
\end{leancode}
\leanlabel{How the Lean proof follows the argument}
\texttt{fresh\allowbreak{}Joint\allowbreak{}Kernel} and \texttt{fresh\allowbreak{}Joint\allowbreak{}Measure} keep the input coordinate in the actual weighted output. The deterministic-first-coordinate lemma proves their product form. \texttt{explicit\allowbreak{}Kernel\_\allowbreak{}pair\_\allowbreak{}eq} is valid for every sigma-finite initial law; \texttt{explicit\allowbreak{}Kernel\_\allowbreak{}gibbs\_\allowbreak{}pair\_\allowbreak{}eq} gives the exact stationary identification.
\end{leanbox}


Earlier results used here: \cite[Section~3.26.3]{LangevinLean2026ULASC}; \cite[Section~3.34.1]{LangevinLean2026ULASC}; \cite[Section~3.34.2]{LangevinLean2026ULASC}; \cite[Section~3.35.1]{LangevinLean2026ULASC}; \cite[Section~3.35.2]{LangevinLean2026ULASC}; \cite[Section~3.35.4]{LangevinLean2026ULASC}; \cite[Section~3.35.7]{LangevinLean2026ULASC}; \cite[Section~3.36.1]{LangevinLean2026ULASC}; \cite[Section~3.36.2]{LangevinLean2026ULASC}; \cite[Section~3.36.3]{LangevinLean2026ULASC}; \cite[Section~3.36.4]{LangevinLean2026ULASC}; \cite[Section~3.36.5]{LangevinLean2026ULASC}; \cite[Section~3.36.6]{LangevinLean2026ULASC}; \cite[Section~3.39.1]{LangevinLean2026ULASC}.
\subsection{\texorpdfstring{The actual anchored endpoint-update recurrence}{The actual anchored endpoint-update recurrence}}
\label{report-step-WassersteinMixing-SemigroupMALAPair-lean-2}


\textbf{Claim.}

\[\mu\otimes K^{n+1}=(I\otimes K)_*(\mu\otimes K^n),\qquad\mu\otimes K^0=(x\mapsto(x,x))_*\mu.\]

\textbf{Proof.} Associativity proves that the recursive kernel power can advance by its final transition: \(K^{n+1}=K\circ K^n\). Composing the product pair with the kernel that keeps its first coordinate and updates its second gives precisely that power. The zero-step identity is the diagonal law of the identity kernel.

\begin{leanbox}{Lean: \texttt{kernelPower\_\allowbreak{}succ\_\allowbreak{}last}, \texttt{anchored\_\allowbreak{}pair\_\allowbreak{}succ}, \texttt{anchored\_\allowbreak{}pair\_\allowbreak{}zero}}
\leansource{\texttt{kernelPower\_\allowbreak{}succ\_\allowbreak{}last} \quad WassersteinMixing/\allowbreak{}SemigroupMALAPair.lean, lines 192–198}
\begin{leancode}
\leanline{0}{theorem\ kernelPower\_succ\_last\ (κ\ :\ Kernel\ E\ E)\ (n\ :\ ℕ)\ :}
\leanline{4}{kernelPower\ κ\ (n\ +\ 1)\ =\ κ\ ∘ₖ\ kernelPower\ κ\ n}
\end{leancode}
\leansource{\texttt{anchored\_\allowbreak{}pair\_\allowbreak{}succ} \quad WassersteinMixing/\allowbreak{}SemigroupMALAPair.lean, lines 201–205}
\begin{leancode}
\leanline{0}{theorem\ anchored\_pair\_succ\ (μ\ :\ Measure\ E)\ [SFinite\ μ]}
\leanline{4}{(κ\ :\ Kernel\ E\ E)\ [IsMarkovKernel\ κ]\ (n\ :\ ℕ)\ :}
\leanline{4}{μ\ ⊗ₘ\ kernelPower\ κ\ (n\ +\ 1)\ =}
\leanline{6}{(Kernel.id\ {\leanfallback ∥}ₖ\ κ)\ ∘ₘ\ (μ\ ⊗ₘ\ kernelPower\ κ\ n)}
\end{leancode}
\leansource{\texttt{anchored\_\allowbreak{}pair\_\allowbreak{}zero} \quad WassersteinMixing/\allowbreak{}SemigroupMALAPair.lean, lines 208–211}
\begin{leancode}
\leanline{0}{theorem\ anchored\_pair\_zero\ (μ\ :\ Measure\ E)\ [SFinite\ μ]\ (κ\ :\ Kernel\ E\ E)\ :}
\leanline{4}{μ\ ⊗ₘ\ kernelPower\ κ\ 0\ =\ μ.map\ (fun\ x\ ↦\ (x,\ x))}
\end{leancode}
\leanlabel{How the Lean proof follows the argument}
\texttt{kernel\allowbreak{}Power\_\allowbreak{}succ\_\allowbreak{}last} proves the alternate composition order by induction. \texttt{anchored\_\allowbreak{}pair\_\allowbreak{}succ} uses the exact parallel-kernel composition identity, and \texttt{anchored\_\allowbreak{}pair\_\allowbreak{}zero} uses the identity-kernel product law.
\end{leanbox}
\subsection{\texorpdfstring{Actual reversible finite-step endpoint pairs}{Actual reversible finite-step endpoint pairs}}
\label{report-step-WassersteinMixing-SemigroupMALAPair-lean-3}


\textbf{Claim.} The actual explicit anchored pair \(J_n=\pi\otimes K^n\) is a probability law, is invariant under swapping, has both marginals \(\pi\), and follows the endpoint-update recurrence above.



\textbf{Proof.} The explicit kernel is Markov. Its one-step pair is exactly the proved reversible Gaussian pair, so generic kernel-power reversal applies. First-marginal preservation and swap symmetry give both Gibbs marginals. The generic anchored successor theorem supplies the actual update recursion.

\begin{leanbox}{Lean: \texttt{explicitPowerJointLaw}, \texttt{explicitPowerJointLaw\_\allowbreak{}isProbability}, \texttt{explicitPowerJointLaw\_\allowbreak{}swap} and 3 more}
\leansource{\texttt{explicitPowerJointLaw} \quad WassersteinMixing/\allowbreak{}SemigroupMALAPair.lean, lines 220–221}
\begin{leancode}
\leanline{0}{def\ explicitPowerJointLaw\ (V\ :\ E\ →\ ℝ)\ (h\ :\ ℝ)\ (n\ :\ ℕ)\ :\ Measure\ (E\ ×\ E)\ :=}
\leanline{2}{ULA.gibbsMeasure\ V\ ⊗ₘ\ StationaryKernelPower.kernelPower\ (explicitKernel\ V\ h)\ n}
\end{leancode}
\leansource{\texttt{explicitPowerJointLaw\_\allowbreak{}isProbability} \quad WassersteinMixing/\allowbreak{}SemigroupMALAPair.lean, lines 223–229}
\begin{leancode}
\leanline{0}{theorem\ explicitPowerJointLaw\_isProbability\ (V\ :\ E\ →\ ℝ)\ (h\ :\ ℝ)\ (L\ :\ ℝ≥0)}
\leanline{4}{[IsProbabilityMeasure\ (ULA.gibbsMeasure\ V)]}
\leanline{4}{(hdiff\ :\ Differentiable\ ℝ\ V)\ (hLip\ :\ LipschitzWith\ L\ (gradient\ V))\ (n\ :\ ℕ)\ :}
\leanline{4}{IsProbabilityMeasure\ (explicitPowerJointLaw\ V\ h\ n)}
\end{leancode}
\leansource{\texttt{explicitPowerJointLaw\_\allowbreak{}swap} \quad WassersteinMixing/\allowbreak{}SemigroupMALAPair.lean, lines 232–239}
\begin{leancode}
\leanline{0}{theorem\ explicitPowerJointLaw\_swap\ (V\ :\ E\ →\ ℝ)\ (h\ :\ ℝ)\ (L\ :\ ℝ≥0)}
\leanline{4}{[IsProbabilityMeasure\ (ULA.gibbsMeasure\ V)]\ (hh\ :\ 0\ <\ h)}
\leanline{4}{(hdiff\ :\ Differentiable\ ℝ\ V)\ (hLip\ :\ LipschitzWith\ L\ (gradient\ V))\ (n\ :\ ℕ)\ :}
\leanline{4}{(explicitPowerJointLaw\ V\ h\ n).map\ Prod.swap\ =\ explicitPowerJointLaw\ V\ h\ n}
\end{leancode}
\leansource{\texttt{explicitPowerJointLaw\_\allowbreak{}fst} \quad WassersteinMixing/\allowbreak{}SemigroupMALAPair.lean, lines 241–246}
\begin{leancode}
\leanline{0}{theorem\ explicitPowerJointLaw\_fst\ (V\ :\ E\ →\ ℝ)\ (h\ :\ ℝ)\ (L\ :\ ℝ≥0)}
\leanline{4}{[IsProbabilityMeasure\ (ULA.gibbsMeasure\ V)]}
\leanline{4}{(hdiff\ :\ Differentiable\ ℝ\ V)\ (hLip\ :\ LipschitzWith\ L\ (gradient\ V))\ (n\ :\ ℕ)\ :}
\leanline{4}{(explicitPowerJointLaw\ V\ h\ n).map\ Prod.fst\ =\ ULA.gibbsMeasure\ V}
\end{leancode}
\leansource{\texttt{explicitPowerJointLaw\_\allowbreak{}snd} \quad WassersteinMixing/\allowbreak{}SemigroupMALAPair.lean, lines 248–254}
\begin{leancode}
\leanline{0}{theorem\ explicitPowerJointLaw\_snd\ (V\ :\ E\ →\ ℝ)\ (h\ :\ ℝ)\ (L\ :\ ℝ≥0)}
\leanline{4}{[IsProbabilityMeasure\ (ULA.gibbsMeasure\ V)]\ (hh\ :\ 0\ <\ h)}
\leanline{4}{(hdiff\ :\ Differentiable\ ℝ\ V)\ (hLip\ :\ LipschitzWith\ L\ (gradient\ V))\ (n\ :\ ℕ)\ :}
\leanline{4}{(explicitPowerJointLaw\ V\ h\ n).map\ Prod.snd\ =\ ULA.gibbsMeasure\ V}
\end{leancode}
\leansource{\texttt{explicitPowerJointLaw\_\allowbreak{}succ} \quad WassersteinMixing/\allowbreak{}SemigroupMALAPair.lean, lines 257–263}
\begin{leancode}
\leanline{0}{theorem\ explicitPowerJointLaw\_succ\ (V\ :\ E\ →\ ℝ)\ (h\ :\ ℝ)\ (L\ :\ ℝ≥0)}
\leanline{4}{[IsProbabilityMeasure\ (ULA.gibbsMeasure\ V)]}
\leanline{4}{(hdiff\ :\ Differentiable\ ℝ\ V)\ (hLip\ :\ LipschitzWith\ L\ (gradient\ V))\ (n\ :\ ℕ)\ :}
\leanline{4}{explicitPowerJointLaw\ V\ h\ (n\ +\ 1)\ =}
\leanline{6}{(Kernel.id\ {\leanfallback ∥}ₖ\ explicitKernel\ V\ h)\ ∘ₘ\ explicitPowerJointLaw\ V\ h\ n}
\end{leancode}
\leanlabel{How the Lean proof follows the argument}
\texttt{explicit\allowbreak{}Power\allowbreak{}Joint\allowbreak{}Law} uses the explicit fresh-Gaussian kernel directly, avoiding a conditional-kernel identification. Its probability, swap, marginal and successor theorems specialize the proved actual kernel identities.
\end{leanbox}


Earlier results used here: \cite[Section~3.26.3]{LangevinLean2026ULASC}; \cite[Section~3.35.5]{LangevinLean2026ULASC}.
\section{\texorpdfstring{The exact anchored law of the actual random MALA path}{The exact anchored law of the actual random MALA path}}
\label{proof-module-WassersteinMixing-SemigroupAnchoredMALALaw-lean}


The finite-time pair of the actual initial state and sampled MALA endpoint has law \(\pi\otimes K^n\), with the explicit fresh-noise kernel. This identifies actual random endpoint pairs with the proved reversible kernel powers.

\subsection{\texorpdfstring{Induction on the actual anchored law}{Induction on the actual anchored law}}
\label{report-step-WassersteinMixing-SemigroupAnchoredMALALaw-lean-1}


\textbf{Claim.}

\[\operatorname{Law}(X_0,X_n^{\mathrm{MALA}})=\pi\otimes K^n,\qquad \operatorname{Law}(X_0,X_n^{\mathrm{MALA}})=\operatorname{Law}(X_n^{\mathrm{MALA}},X_0).\]

\textbf{Proof.} At time zero the anchor/\allowbreak{}endpoint law is the diagonal Gibbs law, exactly \(\pi\otimes K^0\). Both the sampled anchored recursion and the explicit kernel-power pair advance by the same parallel-kernel update. Induction gives equality of the full pair laws. The actual product initial/\allowbreak{}innovation law transfers this equality to the canonical sampled path. Reversal follows from the proved explicit-power detailed balance.

\begin{leanbox}{Lean: \texttt{anchored\_\allowbreak{}chainLaw}, \texttt{recursion\_\allowbreak{}law}, \texttt{anchored\_\allowbreak{}state\_\allowbreak{}law} and 1 more}
\leansource{\texttt{anchored\_\allowbreak{}chainLaw} \quad WassersteinMixing/\allowbreak{}SemigroupAnchoredMALALaw.lean, lines 23–39}
\begin{leancode}
\leanline{0}{theorem\ anchored\_chainLaw\ (V\ :\ E\ →\ ℝ)\ (h\ :\ ℝ)\ (L\ :\ ℝ≥0)}
\leanline{4}{(hdiff\ :\ Differentiable\ ℝ\ V)\ (hLip\ :\ LipschitzWith\ L\ (gradient\ V))}
\leanline{4}{(μ\ :\ Measure\ E)\ [IsProbabilityMeasure\ μ]\ (n\ :\ ℕ)\ :}
\leanline{4}{ULA.chainLaw\ (anchoredUpdate\ V\ h)\ (μ.map\ (fun\ x\ ↦\ (x,\ x)))}
\leanline{6}{(Measure.infinitePi\ (fun\ \_\ :\ ℕ\ ↦\ innovationLaw))\ n\ =}
\leanline{8}{μ\ ⊗ₘ\ StationaryKernelPower.kernelPower\ (MALAStationaryPair.explicitKernel\ V\ h)\ n}
\end{leancode}
\leansource{\texttt{recursion\_\allowbreak{}law} \quad WassersteinMixing/\allowbreak{}SemigroupAnchoredMALALaw.lean, lines 48–70}
\begin{leancode}
\leanline{0}{theorem\ recursion\_law\ \{H\ :\ Type*\}\ [MeasurableSpace\ H]}
\leanline{4}{(F\ :\ H\ →\ ((E\ P.dimension)\ ×\ I)\ →\ H)\ (hF\ :\ Measurable\ (Function.uncurry\ F))}
\leanline{4}{(j\ :\ E\ P.dimension\ →\ H)\ (hj\ :\ Measurable\ j)}
\leanline{4}{(r\ :\ ℝ)\ (hr\ :\ 0\ <\ r)\ (n\ :\ ℕ)\ :}
\leanline{4}{(probability\ P\ V).map\ (fun\ ω\ ↦\ ULA.noiseDrivenChain\ F}
\leanline{6}{(j\ (input\ P\ V\ hV\ r\ ω).1)\ (input\ P\ V\ hV\ r\ ω).2\ n)\ =}
\leanline{6}{ULA.chainLaw\ F\ ((ULA.gibbsMeasure\ V).map\ j)}
\leanline{8}{(Measure.infinitePi\ (fun\ \_\ :\ ℕ\ ↦\ innovationLaw))\ n}
\end{leancode}
\leansource{\texttt{anchored\_\allowbreak{}state\_\allowbreak{}law} \quad WassersteinMixing/\allowbreak{}SemigroupAnchoredMALALaw.lean, lines 73–86}
\begin{leancode}
\leanline{0}{theorem\ anchored\_state\_law\ (r\ :\ ℝ)\ (hr\ :\ 0\ <\ r)\ (n\ :\ ℕ)\ :}
\leanline{4}{(probability\ P\ V).map}
\leanline{6}{(fun\ ω\ ↦\ (actualInitial\ P\ V\ hV\ ω.1,\ (state\ P\ V\ hV\ r\ n\ ω).2))\ =}
\leanline{8}{MALAStationaryPair.explicitPowerJointLaw\ V\ r\ n}
\end{leancode}
\leansource{\texttt{anchored\_\allowbreak{}state\_\allowbreak{}swap} \quad WassersteinMixing/\allowbreak{}SemigroupAnchoredMALALaw.lean, lines 89–97}
\begin{leancode}
\leanline{0}{theorem\ anchored\_state\_swap\ (r\ :\ ℝ)\ (hr\ :\ 0\ <\ r)\ (n\ :\ ℕ)\ :}
\leanline{4}{((probability\ P\ V).map}
\leanline{6}{(fun\ ω\ ↦\ (actualInitial\ P\ V\ hV\ ω.1,\ (state\ P\ V\ hV\ r\ n\ ω).2))).map\ Prod.swap\ =}
\leanline{6}{(probability\ P\ V).map}
\leanline{8}{(fun\ ω\ ↦\ (actualInitial\ P\ V\ hV\ ω.1,\ (state\ P\ V\ hV\ r\ n\ ω).2))}
\end{leancode}
\leanlabel{How the Lean proof follows the argument}
\texttt{anchored\_\allowbreak{}chain\allowbreak{}Law} compares the zero-step and successor laws. \texttt{recursion\_\allowbreak{}law} applies the actual sampled-input law. \texttt{anchored\_\allowbreak{}state\_\allowbreak{}law} identifies the true canonical pair and \texttt{anchored\_\allowbreak{}state\_\allowbreak{}swap} substitutes the proved explicit-power swap identity.
\end{leanbox}


Earlier results used here: \cite[Section~3.20.1]{LangevinLean2026ULASC}; \cite[Section~3.20.4]{LangevinLean2026ULASC}; \cite[Section~3.20.5]{LangevinLean2026ULASC}; \cite[Section~3.26.3]{LangevinLean2026ULASC}; \cite[Section~3.33.6]{LangevinLean2026ULASC}; \cite[Section~3.33.9]{LangevinLean2026ULASC}.
\section{\texorpdfstring{Actual floor meshes and elapsed-time limits}{Actual floor meshes and elapsed-time limits}}
\label{proof-module-WassersteinMixing-SemigroupMesh-lean}


The integer mesh count \(N_r(t)=\lfloor t/r\rfloor\) gives actual finite-step approximants at elapsed times tending to \(t\). Its sum at two target times has the analogous bound and converges to their sum.

\subsection{\texorpdfstring{The true integer count and vanishing remainder}{The true integer count and vanishing remainder}}
\label{report-step-WassersteinMixing-SemigroupMesh-lean-1}


\textbf{Claim.} For \(r>0\) and \(t\ge0\),

\[N_r(t)r\le t,\qquad 0\le t-N_r(t)r<r.\]

Hence \(N_r(t)r\to t\) as \(r\downarrow0\). For two nonnegative times, \((N_r(s)+N_r(t))r\le s+t\) and these elapsed sums tend to \(s+t\).



\textbf{Proof.} The defining floor inequalities bound the ratio \(t/r\) between its natural integer floor and the next integer. Multiply by the positive step to get the elapsed-time and remainder estimates. Squeeze the remainder between zero and the step, then add the two convergent elapsed times.

\begin{leanbox}{Lean: \texttt{meshCount}, \texttt{meshCount\_\allowbreak{}elapsed\_\allowbreak{}le}, \texttt{meshCount\_\allowbreak{}remainder\_\allowbreak{}lt} and 3 more}
\leansource{\texttt{meshCount} \quad WassersteinMixing/\allowbreak{}SemigroupMesh.lean, lines 19–19}
\begin{leancode}
\leanline{0}{def\ meshCount\ (step\ time\ :\ ℝ)\ :\ ℕ\ :=\ ⌊time\ /\ step⌋₊}
\end{leancode}
\leansource{\texttt{meshCount\_\allowbreak{}elapsed\_\allowbreak{}le} \quad WassersteinMixing/\allowbreak{}SemigroupMesh.lean, lines 21–24}
\begin{leancode}
\leanline{0}{theorem\ meshCount\_elapsed\_le\ (step\ time\ :\ ℝ)\ (hs\ :\ 0\ <\ step)\ (ht\ :\ 0\ ≤\ time)\ :}
\leanline{4}{(meshCount\ step\ time\ :\ ℝ)\ *\ step\ ≤\ time}
\end{leancode}
\leansource{\texttt{meshCount\_\allowbreak{}remainder\_\allowbreak{}lt} \quad WassersteinMixing/\allowbreak{}SemigroupMesh.lean, lines 26–31}
\begin{leancode}
\leanline{0}{theorem\ meshCount\_remainder\_lt\ (step\ time\ :\ ℝ)\ (hs\ :\ 0\ <\ step)\ :}
\leanline{4}{time\ -\ (meshCount\ step\ time\ :\ ℝ)\ *\ step\ <\ step}
\end{leancode}
\leansource{\texttt{meshCount\_\allowbreak{}elapsed\_\allowbreak{}tendsto} \quad WassersteinMixing/\allowbreak{}SemigroupMesh.lean, lines 34–42}
\begin{leancode}
\leanline{0}{theorem\ meshCount\_elapsed\_tendsto\ \{I\ :\ Type*\}\ \{l\ :\ Filter\ I\}}
\leanline{4}{(step\ :\ I\ →\ ℝ)\ (hs\ :\ ∀\ i,\ 0\ <\ step\ i)\ (hlim\ :\ Tendsto\ step\ l\ ({\leanfallback 𝓝}\ 0))}
\leanline{4}{(time\ :\ ℝ)\ (ht\ :\ 0\ ≤\ time)\ :}
\leanline{4}{Tendsto\ (fun\ i\ ↦\ (meshCount\ (step\ i)\ time\ :\ ℝ)\ *\ step\ i)\ l\ ({\leanfallback 𝓝}\ time)}
\end{leancode}
\leansource{\texttt{meshCount\_\allowbreak{}sum\_\allowbreak{}elapsed\_\allowbreak{}le} \quad WassersteinMixing/\allowbreak{}SemigroupMesh.lean, lines 44–48}
\begin{leancode}
\leanline{0}{theorem\ meshCount\_sum\_elapsed\_le\ (step\ s\ t\ :\ ℝ)\ (hs\ :\ 0\ <\ step)}
\leanline{4}{(hs0\ :\ 0\ ≤\ s)\ (ht0\ :\ 0\ ≤\ t)\ :}
\leanline{4}{((meshCount\ step\ s\ +\ meshCount\ step\ t\ :\ ℕ)\ :\ ℝ)\ *\ step\ ≤\ s\ +\ t}
\end{leancode}
\leansource{\texttt{meshCount\_\allowbreak{}sum\_\allowbreak{}elapsed\_\allowbreak{}tendsto} \quad WassersteinMixing/\allowbreak{}SemigroupMesh.lean, lines 52–62}
\begin{leancode}
\leanline{0}{theorem\ meshCount\_sum\_elapsed\_tendsto\ \{I\ :\ Type*\}\ \{l\ :\ Filter\ I\}}
\leanline{4}{(step\ :\ I\ →\ ℝ)\ (hs\ :\ ∀\ i,\ 0\ <\ step\ i)\ (hlim\ :\ Tendsto\ step\ l\ ({\leanfallback 𝓝}\ 0))}
\leanline{4}{(s\ t\ :\ ℝ)\ (hs0\ :\ 0\ ≤\ s)\ (ht0\ :\ 0\ ≤\ t)\ :}
\leanline{4}{Tendsto\ (fun\ i\ ↦\ ((meshCount\ (step\ i)\ s\ +\ meshCount\ (step\ i)\ t\ :\ ℕ)\ :\ ℝ)\ *\ step\ i)}
\leanline{6}{l\ ({\leanfallback 𝓝}\ (s\ +\ t))}
\end{leancode}
\leanlabel{How the Lean proof follows the argument}
\texttt{mesh\allowbreak{}Count} is the actual natural floor of the time/\allowbreak{}step ratio. The floor lower and strict upper inequalities prove the bounds. The elapsed-time limit is a squeeze argument; the sum lemmas combine the exact count arithmetic and the two limits.
\end{leanbox}
\section{\texorpdfstring{Self-adjointness of the actual Langevin transition}{Self-adjointness of the actual Langevin transition}}
\label{proof-module-WassersteinMixing-SemigroupReversible-lean}


Actual anchored random MALA pairs are reversible. Their corresponding Gibbs observable compressions converge strongly to the actual Langevin operator. Passing the symmetric pairing identity to that genuine limit proves actual diffusion reversibility.

\subsection{\texorpdfstring{The true reversible approximants and their diffusion limit}{The true reversible approximants and their diffusion limit}}
\label{report-step-WassersteinMixing-SemigroupReversible-lean-1}


\textbf{Claim.} The actual canonical Langevin operator \(P_t\) is self-adjoint for \(t\ge0\).



\textbf{Proof.} The zero-time diffusion variable is the actual initial coordinate. The exact anchored sampled MALA pair law therefore identifies its initial/\allowbreak{}endpoint pair with the reversible explicit kernel-power pair. The pair-law theorem makes each actual MALA compression self-adjoint. Choose the floor-mesh counts approaching \(t\). The actual MALA operator approximation theorem gives strong convergence to \(P_t\); self-adjointness passes to that strong limit by testing all pairings against fixed vectors.

\begin{leanbox}{Lean: \texttt{diffusion\_\allowbreak{}zero\_\allowbreak{}coe}, \texttt{mala\_\allowbreak{}pair\_\allowbreak{}swap}, \texttt{malaOp\_\allowbreak{}isSelfAdjoint} and 1 more}
\leansource{\texttt{diffusion\_\allowbreak{}zero\_\allowbreak{}coe} \quad WassersteinMixing/\allowbreak{}SemigroupReversible.lean, lines 21–28}
\begin{leancode}
\leanline{0}{theorem\ diffusion\_zero\_coe\ (P\ :\ Parameters)\ (V\ :\ E\ P.dimension\ →\ ℝ)}
\leanline{4}{(hV\ :\ PotentialAssumptions\ P\ V)\ :}
\leanline{4}{diffusion\ P\ V\ hV\ 0\ =ᵐ[probability\ P\ V]\ fun\ ω\ ↦\ actualInitial\ P\ V\ hV\ ω.1}
\end{leancode}
\leansource{\texttt{mala\_\allowbreak{}pair\_\allowbreak{}swap} \quad WassersteinMixing/\allowbreak{}SemigroupReversible.lean, lines 32–46}
\begin{leancode}
\leanline{0}{theorem\ mala\_pair\_swap\ (P\ :\ Parameters)\ (V\ :\ E\ P.dimension\ →\ ℝ)}
\leanline{4}{(hV\ :\ PotentialAssumptions\ P\ V)\ (r\ :\ ℝ)\ (hr\ :\ 0\ <\ r)\ (n\ :\ ℕ)\ :}
\leanline{4}{((probability\ P\ V).map\ (fun\ ω\ ↦}
\leanline{6}{(diffusion\ P\ V\ hV\ 0\ ω,\ mala\ P\ V\ hV\ r\ hr\ n\ ω))).map\ Prod.swap\ =}
\leanline{6}{(probability\ P\ V).map\ (fun\ ω\ ↦}
\leanline{8}{(diffusion\ P\ V\ hV\ 0\ ω,\ mala\ P\ V\ hV\ r\ hr\ n\ ω))}
\end{leancode}
\leansource{\texttt{malaOp\_\allowbreak{}isSelfAdjoint} \quad WassersteinMixing/\allowbreak{}SemigroupReversible.lean, lines 53–59}
\begin{leancode}
\leanline{0}{theorem\ malaOp\_isSelfAdjoint\ (P\ :\ Parameters)\ (V\ :\ E\ P.dimension\ →\ ℝ)}
\leanline{4}{(hV\ :\ PotentialAssumptions\ P\ V)\ (r\ :\ ℝ)\ (hr\ :\ 0\ <\ r)\ (n\ :\ ℕ)\ :}
\leanline{4}{IsSelfAdjoint\ (malaOp\ ({\leanfallback 𝕜}\ :=\ {\leanfallback 𝕜})\ (F\ :=\ F)\ P\ V\ hV\ r\ hr\ n)}
\end{leancode}
\leansource{\texttt{op\_\allowbreak{}isSelfAdjoint} \quad WassersteinMixing/\allowbreak{}SemigroupReversible.lean, lines 70–87}
\begin{leancode}
\leanline{0}{theorem\ op\_isSelfAdjoint\ (P\ :\ Parameters)\ (V\ :\ E\ P.dimension\ →\ ℝ)}
\leanline{4}{(hV\ :\ PotentialAssumptions\ P\ V)\ (t\ :\ ℝ)\ (ht\ :\ 0\ ≤\ t)\ :}
\leanline{4}{IsSelfAdjoint\ (op\ ({\leanfallback 𝕜}\ :=\ {\leanfallback 𝕜})\ (F\ :=\ F)\ P\ V\ hV\ t)}
\end{leancode}
\leanlabel{How the Lean proof follows the argument}
\texttt{mala\_\allowbreak{}pair\_\allowbreak{}swap} uses the exact canonical anchored law and actual initial representative. \texttt{mala\allowbreak{}Op\_\allowbreak{}is\allowbreak{}Self\allowbreak{}Adjoint} applies pair-law reversal. \texttt{op\_\allowbreak{}is\allowbreak{}Self\allowbreak{}Adjoint} combines the floor-mesh convergence, actual operator approximation and the proved strong-limit self-adjointness theorem.
\end{leanbox}


Earlier results used here: \cite[Section~3.26.3]{LangevinLean2026ULASC}.
\section{\texorpdfstring{The actual reversible strongly continuous contraction semigroup}{The actual reversible strongly continuous contraction semigroup}}
\label{proof-module-WassersteinMixing-SemigroupInstance-lean}


All operator properties required by the corrector construction are assembled for the actual Gibbs Langevin diffusion. The semigroup law, self-adjointness, contraction, identity at zero and strong continuity are proved from the actual process.

\subsection{\texorpdfstring{Assembling the actual operator properties}{Assembling the actual operator properties}}
\label{report-step-WassersteinMixing-SemigroupInstance-lean-1}


\textbf{Claim.} The canonical transition family is a reversible strongly continuous contraction semigroup on complex \(L^2(\pi)\).



\textbf{Proof.} The actual conditional Markov theorem and compression calculation prove the real composition law. Scalar restriction identifies the real and complex operators, so that law transfers to complex observables. Reversibility follows from actual reversible MALA approximants and their strong limit; contraction, the time-zero identity and strong continuity follow from actual stationary evaluation. Each defining property is therefore a theorem about the same canonical Brownian operator family.

\begin{leanbox}{Lean: \texttt{op\_\allowbreak{}add\_\allowbreak{}real}, \texttt{instSemigroupInstanceReal}, \texttt{op\_\allowbreak{}add} and 2 more}
\leansource{\texttt{op\_\allowbreak{}add\_\allowbreak{}real} \quad WassersteinMixing/\allowbreak{}SemigroupInstance.lean, lines 26–31}
\begin{leancode}
\leanline{0}{theorem\ op\_add\_real\ (P\ :\ Parameters)\ (V\ :\ E\ P.dimension\ →\ ℝ)}
\leanline{4}{(hV\ :\ PotentialAssumptions\ P\ V)\ (s\ t\ :\ ℝ)\ (hs\ :\ 0\ ≤\ s)\ (ht\ :\ 0\ ≤\ t)\ :}
\leanline{4}{op\ ({\leanfallback 𝕜}\ :=\ ℝ)\ (F\ :=\ F)\ P\ V\ hV\ (s\ +\ t)\ =\ op\ P\ V\ hV\ s\ ∘L\ op\ P\ V\ hV\ t}
\end{leancode}
\leansource{\texttt{instSemigroupInstanceReal} \quad WassersteinMixing/\allowbreak{}SemigroupInstance.lean, lines 37–38}
\begin{leancode}
\leanline{0}{local\ instance\ instSemigroupInstanceReal\ :\ InnerProductSpace\ ℝ\ F\ :=}
\leanline{2}{InnerProductSpace.rclikeToReal\ ℂ\ F}
\end{leancode}
\leansource{\texttt{op\_\allowbreak{}add} \quad WassersteinMixing/\allowbreak{}SemigroupInstance.lean, lines 42–51}
\begin{leancode}
\leanline{0}{theorem\ op\_add\ (P\ :\ Parameters)\ (V\ :\ E\ P.dimension\ →\ ℝ)}
\leanline{4}{(hV\ :\ PotentialAssumptions\ P\ V)\ (s\ t\ :\ ℝ)\ (hs\ :\ 0\ ≤\ s)\ (ht\ :\ 0\ ≤\ t)\ :}
\leanline{4}{op\ ({\leanfallback 𝕜}\ :=\ ℂ)\ (F\ :=\ F)\ P\ V\ hV\ (s\ +\ t)\ =\ op\ P\ V\ hV\ s\ *\ op\ P\ V\ hV\ t}
\end{leancode}
\leansource{\texttt{semigroup} \quad WassersteinMixing/\allowbreak{}SemigroupInstance.lean, lines 55–62}
\begin{leancode}
\leanline{0}{def\ semigroup\ (P\ :\ Parameters)\ (V\ :\ E\ P.dimension\ →\ ℝ)\ (hV\ :\ PotentialAssumptions\ P\ V)\ :}
\leanline{4}{ReversibleSemigroup\ (Lp\ F\ 2\ (ULA.gibbsMeasure\ V))\ where}
\leanline{2}{op\ :=\ op\ P\ V\ hV}
\leanline{2}{zero\ :=\ op\_zero\ P\ V\ hV}
\leanline{2}{add\ :=\ fun\ \{s\ t\}\ hs\ ht\ ↦\ op\_add\ P\ V\ hV\ s\ t\ hs\ ht}
\leanline{2}{selfAdjoint\ :=\ fun\ \{t\}\ ht\ ↦\ op\_isSelfAdjoint\ P\ V\ hV\ t\ ht}
\leanline{2}{norm\_le\ :=\ fun\ \{t\}\ \_\ f\ ↦\ op\_norm\_le\ P\ V\ hV\ t\ f}
\leanline{2}{continuous\ :=\ fun\ f\ ↦\ continuousOn\_op\ P\ V\ hV\ f}
\end{leancode}
\leansource{\texttt{semigroup\_\allowbreak{}op} \quad WassersteinMixing/\allowbreak{}SemigroupInstance.lean, lines 64–66}
\begin{leancode}
\leanline{0}{@[simp]\ theorem\ semigroup\_op\ (P\ :\ Parameters)\ (V\ :\ E\ P.dimension\ →\ ℝ)}
\leanline{4}{(hV\ :\ PotentialAssumptions\ P\ V)\ (t\ :\ ℝ)\ :}
\leanline{4}{(semigroup\ (F\ :=\ F)\ P\ V\ hV).op\ t\ =\ op\ ({\leanfallback 𝕜}\ :=\ ℂ)\ P\ V\ hV\ t}
\end{leancode}
\leanlabel{How the Lean proof follows the argument}
\texttt{op\_\allowbreak{}add\_\allowbreak{}real} instantiates the actual Markov composition theorem. The local real structure and \texttt{op\_\allowbreak{}restrict\allowbreak{}Scalars} transfer it in \texttt{op\_\allowbreak{}add}. \texttt{semigroup} fills every field of \texttt{Reversible\allowbreak{}Semigroup} with the proved canonical operator results, and \texttt{semigroup\_\allowbreak{}op} records its definitional identification.
\end{leanbox}


Earlier results used here: \cite[Section~3.26.3]{LangevinLean2026ULASC}.
\section{\texorpdfstring{Real-range preservation by the actual Langevin operator}{Real-range preservation by the actual Langevin operator}}
\label{proof-module-WassersteinMixing-ComplexRealActual-lean}


The canonical Brownian diffusion evaluations are actual measure-preserving observations. Their naturality identifies the complex transition on real observables with the actual real transition at every time. The real-range condition needed by the corrector is consequently derived.

\subsection{\texorpdfstring{Exact intertwining of the actual real and complex transitions}{Exact intertwining of the actual real and complex transitions}}
\label{report-step-WassersteinMixing-ComplexRealActual-lean-1}


\textbf{Claim.} For every pointwise real isometry \(i\), every actual time \(t\) and real observable \(f\),

\[P_t^{\mathbb C}(if)=i(P_t^{\mathbb R}f).\]

\textbf{Proof.} Gibbs stationarity makes evaluation at zero and at \(t\) isometric pullbacks. Apply the exact compression embedding theorem to these two actual observations. This proves range preservation and applies in particular to the concrete orthonormal-coordinate inclusion, with no norm or dimension loss.

\begin{leanbox}{Lean: \texttt{instance\_\allowbreak{}at\_\allowbreak{}22}, \texttt{op\_\allowbreak{}mapLp}, \texttt{op\_\allowbreak{}mem\_\allowbreak{}real\_\allowbreak{}range} and 1 more}
\leansource{\texttt{instance\_\allowbreak{}at\_\allowbreak{}22} \quad WassersteinMixing/\allowbreak{}ComplexRealActual.lean, lines 22–22}
\begin{leancode}
\leanline{0}{local\ instance\ :\ InnerProductSpace\ ℝ\ H\ :=\ InnerProductSpace.rclikeToReal\ ℂ\ H}
\end{leancode}
\leansource{\texttt{op\_\allowbreak{}mapLp} \quad WassersteinMixing/\allowbreak{}ComplexRealActual.lean, lines 26–34}
\begin{leancode}
\leanline{0}{theorem\ op\_mapLp\ (i\ :\ R\ →ₗᵢ[ℝ]\ H)\ (P\ :\ Parameters)\ (V\ :\ E\ P.dimension\ →\ ℝ)}
\leanline{4}{(hV\ :\ PotentialAssumptions\ P\ V)\ (t\ :\ ℝ)\ (f\ :\ Lp\ R\ 2\ (ULA.gibbsMeasure\ V))\ :}
\leanline{4}{op\ ({\leanfallback 𝕜}\ :=\ ℂ)\ P\ V\ hV\ t\ (RealComplexEmbedding.mapLp\ i\ (ULA.gibbsMeasure\ V)\ f)\ =}
\leanline{6}{RealComplexEmbedding.mapLp\ i\ (ULA.gibbsMeasure\ V)}
\leanline{8}{(op\ ({\leanfallback 𝕜}\ :=\ ℝ)\ P\ V\ hV\ t\ f)}
\end{leancode}
\leansource{\texttt{op\_\allowbreak{}mem\_\allowbreak{}real\_\allowbreak{}range} \quad WassersteinMixing/\allowbreak{}ComplexRealActual.lean, lines 37–41}
\begin{leancode}
\leanline{0}{theorem\ op\_mem\_real\_range\ (i\ :\ R\ →ₗᵢ[ℝ]\ H)\ (P\ :\ Parameters)\ (V\ :\ E\ P.dimension\ →\ ℝ)}
\leanline{4}{(hV\ :\ PotentialAssumptions\ P\ V)\ (t\ :\ ℝ)\ (f\ :\ Lp\ R\ 2\ (ULA.gibbsMeasure\ V))\ :}
\leanline{4}{op\ ({\leanfallback 𝕜}\ :=\ ℂ)\ P\ V\ hV\ t\ (RealComplexEmbedding.mapLp\ i\ (ULA.gibbsMeasure\ V)\ f)\ ∈}
\leanline{6}{Set.range\ (RealComplexEmbedding.mapLp\ i\ (ULA.gibbsMeasure\ V))}
\end{leancode}
\leansource{\texttt{op\_\allowbreak{}coordinateLp} \quad WassersteinMixing/\allowbreak{}ComplexRealActual.lean, lines 45–52}
\begin{leancode}
\leanline{0}{theorem\ op\_coordinateLp\ (P\ :\ Parameters)\ (V\ :\ E\ P.dimension\ →\ ℝ)}
\leanline{4}{(hV\ :\ PotentialAssumptions\ P\ V)\ (t\ :\ ℝ)}
\leanline{4}{(f\ :\ Lp\ (E\ P.dimension)\ 2\ (ULA.gibbsMeasure\ V))\ :}
\leanline{4}{op\ ({\leanfallback 𝕜}\ :=\ ℂ)\ P\ V\ hV\ t}
\leanline{8}{(RealComplexEmbedding.coordinateLp\ (E\ P.dimension)\ (ULA.gibbsMeasure\ V)\ f)\ =}
\leanline{6}{RealComplexEmbedding.coordinateLp\ (E\ P.dimension)\ (ULA.gibbsMeasure\ V)}
\leanline{8}{(op\ ({\leanfallback 𝕜}\ :=\ ℝ)\ P\ V\ hV\ t\ f)}
\end{leancode}
\leanlabel{How the Lean proof follows the argument}
\texttt{op\_\allowbreak{}map\allowbreak{}Lp} specializes \texttt{compression\_\allowbreak{}map\allowbreak{}Lp} to the actual canonical state evaluations and their proved Gibbs laws. \texttt{op\_\allowbreak{}mem\_\allowbreak{}real\_\allowbreak{}range} gives the explicit range witness. \texttt{op\_\allowbreak{}coordinate\allowbreak{}Lp} uses the concrete coordinate embedding of the physical state space.
\end{leanbox}


Earlier results used here: \cite[Section~3.26.3]{LangevinLean2026ULASC}.
\subsection{\texorpdfstring{The corrector extraction is the actual real transition}{The corrector extraction is the actual real transition}}
\label{report-step-WassersteinMixing-ComplexRealActual-lean-2}


\textbf{Claim.} A semigroup assembled from these actual canonical operators preserves the real coordinate image. Its extracted real transition \(i^*P_ti\) equals the actual real canonical transition itself.



\textbf{Proof.} The preceding intertwining identity puts every output in the embedded real range, which is exactly the projection-fixed condition. Apply the embedding adjoint to the same identity and use \(i^*i=I\) to identify the extracted transition.

\begin{leanbox}{Lean: \texttt{preservesReal\_\allowbreak{}of\_\allowbreak{}op\_\allowbreak{}eq}, \texttt{realOp\_\allowbreak{}eq\_\allowbreak{}op}}
\leansource{\texttt{preservesReal\_\allowbreak{}of\_\allowbreak{}op\_\allowbreak{}eq} \quad WassersteinMixing/\allowbreak{}ComplexRealActual.lean, lines 56–65}
\begin{leancode}
\leanline{0}{theorem\ preservesReal\_of\_op\_eq\ (i\ :\ R\ →ₗᵢ[ℝ]\ H)}
\leanline{4}{(P\ :\ Parameters)\ (V\ :\ E\ P.dimension\ →\ ℝ)\ (hV\ :\ PotentialAssumptions\ P\ V)}
\leanline{4}{(S\ :\ ReversibleSemigroup\ (Lp\ H\ 2\ (ULA.gibbsMeasure\ V)))}
\leanline{4}{(hS\ :\ ∀\ t,\ 0\ ≤\ t\ →\ S.op\ t\ =\ op\ ({\leanfallback 𝕜}\ :=\ ℂ)\ P\ V\ hV\ t)\ :}
\leanline{4}{ComplexRealCorrector.PreservesReal\ S}
\leanline{6}{(RealComplexEmbedding.mapLp\ i\ (ULA.gibbsMeasure\ V))}
\end{leancode}
\leansource{\texttt{realOp\_\allowbreak{}eq\_\allowbreak{}op} \quad WassersteinMixing/\allowbreak{}ComplexRealActual.lean, lines 69–79}
\begin{leancode}
\leanline{0}{theorem\ realOp\_eq\_op\ (i\ :\ R\ →ₗᵢ[ℝ]\ H)}
\leanline{4}{(P\ :\ Parameters)\ (V\ :\ E\ P.dimension\ →\ ℝ)\ (hV\ :\ PotentialAssumptions\ P\ V)}
\leanline{4}{(S\ :\ ReversibleSemigroup\ (Lp\ H\ 2\ (ULA.gibbsMeasure\ V)))}
\leanline{4}{(t\ :\ ℝ)\ (hS\ :\ S.op\ t\ =\ op\ ({\leanfallback 𝕜}\ :=\ ℂ)\ P\ V\ hV\ t)\ :}
\leanline{4}{ComplexRealCorrector.realOp\ S\ (RealComplexEmbedding.mapLp\ i\ (ULA.gibbsMeasure\ V))\ t\ =}
\leanline{6}{op\ ({\leanfallback 𝕜}\ :=\ ℝ)\ P\ V\ hV\ t}
\end{leancode}
\leanlabel{How the Lean proof follows the argument}
\texttt{preserves\allowbreak{}Real\_\allowbreak{}of\_\allowbreak{}op\_\allowbreak{}eq} rewrites the supplied operator identification and applies the proved actual range theorem. \texttt{real\allowbreak{}Op\_\allowbreak{}eq\_\allowbreak{}op} rewrites through \texttt{op\_\allowbreak{}map\allowbreak{}Lp} and cancels the isometric embedding with \texttt{adjoint\_\allowbreak{}embed}.
\end{leanbox}


Earlier results used here: \cite[Section~3.26.3]{LangevinLean2026ULASC}.
\section{\texorpdfstring{Stationary evaluation of the constructed Poisson corrector}{Stationary evaluation of the constructed Poisson corrector}}
\label{proof-module-WassersteinMixing-PoissonProcess-lean}


A real observable transition and its genuine bounded Poisson corrector can be evaluated along a stationary Markov process. The resulting corrector increments, exact coboundary decomposition and centered martingale residual inherit the same constants, including the residual coefficient \(18\).

\subsection{\texorpdfstring{The stationary Markov evaluation identity}{The stationary Markov evaluation identity}}
\label{report-step-WassersteinMixing-PoissonProcess-lean-1}


\textbf{Claim.} For stationary isometric evaluations \(U_n\) and transition \(P\), with \(\mathbb E[U_{n+1}f\mid\mathcal F_n]=U_nPf\),

\[\|U_{n+1}f-U_nf\|_2^2=2\langle f,(I-P)f\rangle.\]

\textbf{Proof.} Adaptation makes current evaluation fixed by conditional expectation. Its orthogonal-projection pairing identity replaces the two-time cross term by \(\langle U_nf,U_nPf\rangle\). The evaluation isometry transfers this pairing and both square norms to the observable space. Expanding the square gives the formula.

\begin{leanbox}{Lean: \texttt{PoissonProcess}, \texttt{condExp\_\allowbreak{}eval\_\allowbreak{}self}, \texttt{eval\_\allowbreak{}increment\_\allowbreak{}norm\_\allowbreak{}sq}}
\leansource{\texttt{PoissonProcess} \quad WassersteinMixing/\allowbreak{}PoissonProcess.lean, lines 29–43}
\begin{leancode}
\leanline{0}{structure\ PoissonProcess\ (Q\ :\ Measure\ Ω)\ (h\ :\ ℝ)\ (hh\ :\ 0\ <\ h)\ where}
\leanline{2}{transition\ :\ H\ →L[ℝ]\ H}
\leanline{2}{corrector\ :\ H\ →L[ℝ]\ H}
\leanline{2}{poissonSource\ :\ H\ →L[ℝ]\ H}
\leanline{2}{corrector\_norm\ :\ ∀\ f,\ ‖corrector\ f‖\ ≤\ h\ *\ ‖f‖}
\leanline{2}{corrector\_energy\ :\ ∀\ f,}
\leanline{4}{inner\ ℝ\ (corrector\ f)\ (corrector\ f\ -\ transition\ (corrector\ f))\ ≤}
\leanline{6}{h\ \textasciicircum{}\ 2\ *\ inner\ ℝ\ f\ (f\ -\ transition\ f)}
\leanline{2}{poisson\_identity\ :\ ∀\ f,\ corrector\ f\ -\ transition\ (corrector\ f)\ =\ poissonSource\ f}
\leanline{2}{filtration\ :\ Filtration\ ℕ\ ‹MeasurableSpace\ Ω›}
\leanline{2}{eval\ :\ ℕ\ →\ H\ →ₗᵢ[ℝ]\ Lp\ E\ 2\ Q}
\leanline{2}{adapted\ :\ ∀\ n\ f,\ AEStronglyMeasurable[filtration\ n]\ (eval\ n\ f)\ Q}
\leanline{2}{markov\ :\ ∀\ n\ f,}
\leanline{4}{(condExpL2\ E\ ℝ\ (filtration.le\ n)\ (eval\ (n\ +\ 1)\ f)\ :\ Lp\ E\ 2\ Q)\ =}
\leanline{6}{eval\ n\ (transition\ f)}
\end{leancode}
\leansource{\texttt{condExp\_\allowbreak{}eval\_\allowbreak{}self} \quad WassersteinMixing/\allowbreak{}PoissonProcess.lean, lines 52–57}
\begin{leancode}
\leanline{0}{theorem\ condExp\_eval\_self\ (f\ :\ H)\ (n\ :\ ℕ)\ :}
\leanline{4}{(condExpL2\ E\ ℝ\ (S.filtration.le\ n)\ (S.eval\ n\ f)\ :\ Lp\ E\ 2\ Q)\ =\ S.eval\ n\ f}
\end{leancode}
\leansource{\texttt{eval\_\allowbreak{}increment\_\allowbreak{}norm\_\allowbreak{}sq} \quad WassersteinMixing/\allowbreak{}PoissonProcess.lean, lines 61–70}
\begin{leancode}
\leanline{0}{theorem\ eval\_increment\_norm\_sq\ (f\ :\ H)\ (n\ :\ ℕ)\ :}
\leanline{4}{‖S.eval\ (n\ +\ 1)\ f\ -\ S.eval\ n\ f‖\ \textasciicircum{}\ 2\ =}
\leanline{6}{2\ *\ inner\ ℝ\ f\ (f\ -\ S.transition\ f)}
\end{leancode}
\leanlabel{How the Lean proof follows the argument}
\texttt{Poisson\allowbreak{}Process} packages a transition, the constructed-corrector properties and actual evaluation/\allowbreak{}Markov data for reuse. \texttt{cond\allowbreak{}Exp\_\allowbreak{}eval\_\allowbreak{}self} uses measurability of current evaluations. \texttt{eval\_\allowbreak{}increment\_\allowbreak{}norm\_\allowbreak{}sq} applies the actual conditional-expectation inner-product identity.
\end{leanbox}
\subsection{\texorpdfstring{The evaluated corrector and its increment energy}{The evaluated corrector and its increment energy}}
\label{report-step-WassersteinMixing-PoissonProcess-lean-2}


\textbf{Claim.} For \(g_n=U_nG_hf\),

\[\|g_n\|_2\le h\|f\|,\qquad \|g_{n+1}-g_n\|_2^2\le h^2\|U_{n+1}f-U_nf\|_2^2.\]

\textbf{Proof.} The norm bound is the corrector operator bound transported through the isometry. For the increment, use the preceding stationary increment formula for \(G_hf\), then apply the proved corrector Dirichlet estimate and the same formula for \(f\). Thus an actual score-increment bound \(6h\beta^2d\) implies the corrector-increment bound \(6h^3\beta^2d\). Adaptation follows from the actual evaluation property.

\begin{leanbox}{Lean: \texttt{correctorAt}, \texttt{correctorAt\_\allowbreak{}adapted}, \texttt{correctorAt\_\allowbreak{}norm\_\allowbreak{}le} and 3 more}
\leansource{\texttt{correctorAt} \quad WassersteinMixing/\allowbreak{}PoissonProcess.lean, lines 73–74}
\begin{leancode}
\leanline{0}{def\ correctorAt\ (f\ :\ H)\ (n\ :\ ℕ)\ :\ Lp\ E\ 2\ Q\ :=}
\leanline{2}{S.eval\ n\ (S.corrector\ f)}
\end{leancode}
\leansource{\texttt{correctorAt\_\allowbreak{}adapted} \quad WassersteinMixing/\allowbreak{}PoissonProcess.lean, lines 78–80}
\begin{leancode}
\leanline{0}{theorem\ correctorAt\_adapted\ (f\ :\ H)\ (n\ :\ ℕ)\ :}
\leanline{4}{AEStronglyMeasurable[S.filtration\ n]\ (S.correctorAt\ f\ n)\ Q}
\end{leancode}
\leansource{\texttt{correctorAt\_\allowbreak{}norm\_\allowbreak{}le} \quad WassersteinMixing/\allowbreak{}PoissonProcess.lean, lines 83–86}
\begin{leancode}
\leanline{0}{theorem\ correctorAt\_norm\_le\ (f\ :\ H)\ (n\ :\ ℕ)\ :}
\leanline{4}{‖S.correctorAt\ f\ n‖\ ≤\ h\ *\ ‖f‖}
\end{leancode}
\leansource{\texttt{correctorAt\_\allowbreak{}norm\_\allowbreak{}sq\_\allowbreak{}le} \quad WassersteinMixing/\allowbreak{}PoissonProcess.lean, lines 89–94}
\begin{leancode}
\leanline{0}{theorem\ correctorAt\_norm\_sq\_le\ (f\ :\ H)\ (B\ :\ ℝ)\ (hscore\ :\ ‖f‖\ \textasciicircum{}\ 2\ ≤\ B)\ (n\ :\ ℕ)\ :}
\leanline{4}{‖S.correctorAt\ f\ n‖\ \textasciicircum{}\ 2\ ≤\ h\ \textasciicircum{}\ 2\ *\ B}
\end{leancode}
\leansource{\texttt{correctorAt\_\allowbreak{}increment\_\allowbreak{}norm\_\allowbreak{}sq\_\allowbreak{}le} \quad WassersteinMixing/\allowbreak{}PoissonProcess.lean, lines 98–105}
\begin{leancode}
\leanline{0}{theorem\ correctorAt\_increment\_norm\_sq\_le\ (f\ :\ H)\ (n\ :\ ℕ)\ :}
\leanline{4}{‖S.correctorAt\ f\ n\ -\ S.correctorAt\ f\ (n\ +\ 1)‖\ \textasciicircum{}\ 2\ ≤}
\leanline{6}{h\ \textasciicircum{}\ 2\ *\ ‖S.eval\ (n\ +\ 1)\ f\ -\ S.eval\ n\ f‖\ \textasciicircum{}\ 2}
\end{leancode}
\leansource{\texttt{correctorAt\_\allowbreak{}increment\_\allowbreak{}energy\_\allowbreak{}le} \quad WassersteinMixing/\allowbreak{}PoissonProcess.lean, lines 108–112}
\begin{leancode}
\leanline{0}{theorem\ correctorAt\_increment\_energy\_le\ (f\ :\ H)\ (β\ d\ :\ ℝ)\ (n\ :\ ℕ)}
\leanline{4}{(hincrement\ :\ ‖S.eval\ (n\ +\ 1)\ f\ -\ S.eval\ n\ f‖\ \textasciicircum{}\ 2\ ≤\ 6\ *\ h\ *\ β\ \textasciicircum{}\ 2\ *\ d)\ :}
\leanline{4}{‖S.correctorAt\ f\ n\ -\ S.correctorAt\ f\ (n\ +\ 1)‖\ \textasciicircum{}\ 2\ ≤\ 6\ *\ h\ \textasciicircum{}\ 3\ *\ β\ \textasciicircum{}\ 2\ *\ d}
\end{leancode}
\leanlabel{How the Lean proof follows the argument}
\texttt{corrector\allowbreak{}At} evaluates the existing corrector. Its norm and adaptation lemmas use the corresponding constructed properties. \texttt{corrector\allowbreak{}At\_\allowbreak{}increment\_\allowbreak{}norm\_\allowbreak{}sq\_\allowbreak{}le} rewrites both increments as energy pairings; the final energy lemma inserts the primitive score increment estimate.
\end{leanbox}
\subsection{\texorpdfstring{The true centered residual and coefficient eighteen}{The true centered residual and coefficient eighteen}}
\label{report-step-WassersteinMixing-PoissonProcess-lean-3}


\textbf{Claim.} For a local defect \(D_n\), define \(M_n=D_n-(g_n-g_{n+1})\). Then \(D_n=g_n-g_{n+1}+M_n\). If its conditional mean is the evaluated Poisson source, \(\mathbb E[M_n\mid\mathcal F_n]=0\), and \(M_n\) is orthogonal to every current adapted square-integrable vector. If \(\|D_n\|_2^2\le3h^3\beta^2d\), then \(\|M_n\|_2^2\le18h^3\beta^2d\).



\textbf{Proof.} The decomposition is algebraic. Conditional expectation of the corrector difference is the Poisson source by the Markov and Poisson identities, giving cancellation. Orthogonality follows from the conditional-expectation projection. Finally \(\|a-b\|^2\le2\|a\|^2+2\|b\|^2\) combines the independent constants \(3\) and \(6\).

\begin{leanbox}{Lean: \texttt{martingaleAt}, \texttt{defect\_\allowbreak{}decomposition}, \texttt{martingaleAt\_\allowbreak{}condExp\_\allowbreak{}zero} and 2 more}
\leansource{\texttt{martingaleAt} \quad WassersteinMixing/\allowbreak{}PoissonProcess.lean, lines 115–116}
\begin{leancode}
\leanline{0}{def\ martingaleAt\ (f\ :\ H)\ (D\ :\ ℕ\ →\ Lp\ E\ 2\ Q)\ (n\ :\ ℕ)\ :\ Lp\ E\ 2\ Q\ :=}
\leanline{2}{D\ n\ -\ (S.correctorAt\ f\ n\ -\ S.correctorAt\ f\ (n\ +\ 1))}
\end{leancode}
\leansource{\texttt{defect\_\allowbreak{}decomposition} \quad WassersteinMixing/\allowbreak{}PoissonProcess.lean, lines 119–122}
\begin{leancode}
\leanline{0}{theorem\ defect\_decomposition\ (f\ :\ H)\ (D\ :\ ℕ\ →\ Lp\ E\ 2\ Q)\ (n\ :\ ℕ)\ :}
\leanline{4}{D\ n\ =\ S.correctorAt\ f\ n\ -\ S.correctorAt\ f\ (n\ +\ 1)\ +\ S.martingaleAt\ f\ D\ n}
\end{leancode}
\leansource{\texttt{martingaleAt\_\allowbreak{}condExp\_\allowbreak{}zero} \quad WassersteinMixing/\allowbreak{}PoissonProcess.lean, lines 126–132}
\begin{leancode}
\leanline{0}{theorem\ martingaleAt\_condExp\_zero\ (f\ :\ H)\ (D\ :\ ℕ\ →\ Lp\ E\ 2\ Q)\ (n\ :\ ℕ)}
\leanline{4}{(hmean\ :\ (condExpL2\ E\ ℝ\ (S.filtration.le\ n)\ (D\ n)\ :\ Lp\ E\ 2\ Q)\ =}
\leanline{6}{S.eval\ n\ (S.poissonSource\ f))\ :}
\leanline{4}{(condExpL2\ E\ ℝ\ (S.filtration.le\ n)\ (S.martingaleAt\ f\ D\ n)\ :\ Lp\ E\ 2\ Q)\ =\ 0}
\end{leancode}
\leansource{\texttt{martingaleAt\_\allowbreak{}inner\_\allowbreak{}eq\_\allowbreak{}zero} \quad WassersteinMixing/\allowbreak{}PoissonProcess.lean, lines 135–144}
\begin{leancode}
\leanline{0}{theorem\ martingaleAt\_inner\_eq\_zero\ (f\ :\ H)\ (D\ :\ ℕ\ →\ Lp\ E\ 2\ Q)\ (n\ :\ ℕ)}
\leanline{4}{(hmean\ :\ (condExpL2\ E\ ℝ\ (S.filtration.le\ n)\ (D\ n)\ :\ Lp\ E\ 2\ Q)\ =}
\leanline{6}{S.eval\ n\ (S.poissonSource\ f))}
\leanline{4}{(predictable\ :\ Lp\ E\ 2\ Q)}
\leanline{4}{(hadapted\ :\ AEStronglyMeasurable[S.filtration\ n]\ predictable\ Q)\ :}
\leanline{4}{inner\ ℝ\ predictable\ (S.martingaleAt\ f\ D\ n)\ =\ 0}
\end{leancode}
\leansource{\texttt{martingaleAt\_\allowbreak{}energy\_\allowbreak{}le} \quad WassersteinMixing/\allowbreak{}PoissonProcess.lean, lines 148–156}
\begin{leancode}
\leanline{0}{theorem\ martingaleAt\_energy\_le\ (f\ :\ H)\ (D\ :\ ℕ\ →\ Lp\ E\ 2\ Q)\ (β\ d\ :\ ℝ)\ (n\ :\ ℕ)}
\leanline{4}{(hdefect\ :\ ‖D\ n‖\ \textasciicircum{}\ 2\ ≤\ 3\ *\ h\ \textasciicircum{}\ 3\ *\ β\ \textasciicircum{}\ 2\ *\ d)}
\leanline{4}{(hincrement\ :\ ‖S.eval\ (n\ +\ 1)\ f\ -\ S.eval\ n\ f‖\ \textasciicircum{}\ 2\ ≤\ 6\ *\ h\ *\ β\ \textasciicircum{}\ 2\ *\ d)\ :}
\leanline{4}{‖S.martingaleAt\ f\ D\ n‖\ \textasciicircum{}\ 2\ ≤\ 18\ *\ h\ \textasciicircum{}\ 3\ *\ β\ \textasciicircum{}\ 2\ *\ d}
\end{leancode}
\leanlabel{How the Lean proof follows the argument}
\texttt{martingale\allowbreak{}At} defines the actual residual. The centering proof uses linearity of \texttt{cond\allowbreak{}Exp\allowbreak{}L2}, adaptation, the Markov identity and the proved Poisson equation. The energy theorem combines the previously proved defect and corrector-increment bounds; no residual-energy conclusion is a primitive input.
\end{leanbox}
\section{\texorpdfstring{From the constructed semigroup corrector to the comparison process}{From the constructed semigroup corrector to the comparison process}}
\label{proof-module-WassersteinMixing-SemigroupProcess-lean}


The bounded corrector of a reversible semigroup, its derived real-range preservation and actual stationary Markov evaluations supply a real Poisson comparison process. Corrector norms, Dirichlet domination and the exact integrated source are consequences of the construction.

\subsection{\texorpdfstring{The genuine source integral and evaluated process}{The genuine source integral and evaluated process}}
\label{report-step-WassersteinMixing-SemigroupProcess-lean-1}


\textbf{Construction.} Extract the real corrector \(G_h^{\mathbb R}\) and real transition \(P_h^{\mathbb R}\), and define the source operator by \((I-P_h^{\mathbb R})G_h^{\mathbb R}\). The already proved Poisson identity identifies this operator with \(f\mapsto\int_0^h(f-P_t^{\mathbb R}f)\,dt\).



\textbf{Proof.} The real-extraction theorems provide the exact integral identity, the norm bound \(h\) and energy factor \(h^2\). Combine them with the supplied actual evaluation isometries, their adaptation and their conditional Markov identity. The resulting process uses the constructed corrector itself and retains all these proved estimates.

\begin{leanbox}{Lean: \texttt{semigroupPoissonSource}, \texttt{semigroupPoissonSource\_\allowbreak{}eq\_\allowbreak{}integral}, \texttt{semigroupPoissonProcess}}
\leansource{\texttt{semigroupPoissonSource} \quad WassersteinMixing/\allowbreak{}SemigroupProcess.lean, lines 26–28}
\begin{leancode}
\leanline{0}{def\ semigroupPoissonSource\ (S\ :\ ReversibleSemigroup\ H)\ (i\ :\ R\ →ₗᵢ[ℝ]\ H)}
\leanline{4}{(h\ :\ ℝ)\ (hh\ :\ 0\ ≤\ h)\ :\ R\ →L[ℝ]\ R\ :=}
\leanline{2}{realCorrector\ S\ i\ h\ hh\ -\ realOp\ S\ i\ h\ ∘L\ realCorrector\ S\ i\ h\ hh}
\end{leancode}
\leansource{\texttt{semigroupPoissonSource\_\allowbreak{}eq\_\allowbreak{}integral} \quad WassersteinMixing/\allowbreak{}SemigroupProcess.lean, lines 30–34}
\begin{leancode}
\leanline{0}{theorem\ semigroupPoissonSource\_eq\_integral\ (S\ :\ ReversibleSemigroup\ H)}
\leanline{4}{(i\ :\ R\ →ₗᵢ[ℝ]\ H)\ (hreal\ :\ PreservesReal\ S\ i)\ (h\ :\ ℝ)\ (hh\ :\ 0\ ≤\ h)\ (f\ :\ R)\ :}
\leanline{4}{semigroupPoissonSource\ S\ i\ h\ hh\ f\ =\ ∫\ t\ in\ 0..h,\ f\ -\ realOp\ S\ i\ t\ f}
\end{leancode}
\leansource{\texttt{semigroupPoissonProcess} \quad WassersteinMixing/\allowbreak{}SemigroupProcess.lean, lines 38–54}
\begin{leancode}
\leanline{0}{def\ semigroupPoissonProcess\ (Q\ :\ Measure\ Ω)\ (S\ :\ ReversibleSemigroup\ H)}
\leanline{4}{(i\ :\ R\ →ₗᵢ[ℝ]\ H)\ (hreal\ :\ PreservesReal\ S\ i)\ (h\ :\ ℝ)\ (hh\ :\ 0\ <\ h)}
\leanline{4}{(filtration\ :\ Filtration\ ℕ\ mΩ)\ (eval\ :\ ℕ\ →\ R\ →ₗᵢ[ℝ]\ Lp\ E\ 2\ Q)}
\leanline{4}{(hadapted\ :\ ∀\ n\ f,\ AEStronglyMeasurable[filtration\ n]\ (eval\ n\ f)\ Q)}
\leanline{4}{(hmarkov\ :\ ∀\ n\ f,}
\leanline{6}{(condExpL2\ E\ ℝ\ (filtration.le\ n)\ (eval\ (n\ +\ 1)\ f)\ :\ Lp\ E\ 2\ Q)\ =}
\leanline{8}{eval\ n\ (realOp\ S\ i\ h\ f))\ :\ PoissonProcess\ R\ E\ Q\ h\ hh\ where}
\leanline{2}{transition\ :=\ realOp\ S\ i\ h}
\leanline{2}{corrector\ :=\ realCorrector\ S\ i\ h\ hh.le}
\leanline{2}{poissonSource\ :=\ semigroupPoissonSource\ S\ i\ h\ hh.le}
\leanline{2}{corrector\_norm\ :=\ realCorrector\_norm\_le\ S\ i\ hreal\ h\ hh.le}
\leanline{2}{corrector\_energy\ :=\ realCorrector\_energy\_le\ S\ i\ hreal\ h\ hh.le\ hh.le}
\leanline{2}{poisson\_identity\ :=\ fun\ \_\ ↦\ rfl}
\leanline{2}{filtration\ :=\ filtration}
\leanline{2}{eval\ :=\ eval}
\leanline{2}{adapted\ :=\ hadapted}
\leanline{2}{markov\ :=\ hmarkov}
\end{leancode}
\leanlabel{How the Lean proof follows the argument}
\texttt{semigroup\allowbreak{}Poisson\allowbreak{}Source} is the bounded operator difference. \texttt{semigroup\allowbreak{}Poisson\allowbreak{}Source\_\allowbreak{}eq\_\allowbreak{}integral} invokes \texttt{real\allowbreak{}Corrector\_\allowbreak{}poisson}. \texttt{semigroup\allowbreak{}Poisson\allowbreak{}Process} fills every corrector field from the proved real-extraction theorems, leaving only actual evaluation and Markov facts to the concrete process construction.
\end{leanbox}
\section{\texorpdfstring{The genuine Gibbs score observable}{The genuine Gibbs score observable}}
\label{proof-module-WassersteinMixing-ActualScoreObservable-lean}


The score \(f=\nabla V\) is constructed as a member of \(L^2(\pi;\mathbb R^d)\). Actual stationary evaluation recovers the previously constructed score process, and the Hilbert-space norm is bounded by \(\beta d\).

\subsection{\texorpdfstring{Constructing the square-integrable gradient}{Constructing the square-integrable gradient}}
\label{report-step-WassersteinMixing-ActualScoreObservable-lean-1}


\textbf{Claim.} \(\nabla V\in L^2(\pi)\), where \(\pi\) is the normalized Gibbs law of the actual potential.



\textbf{Proof.} Strong convexity makes the Gibbs partition function finite. The shared Gibbs gradient-moment theorem then proves square integrability of the Lipschitz gradient. Passing to its equivalence class preserves the actual gradient as an almost-everywhere representative.

\begin{leanbox}{Lean: \texttt{actual\_\allowbreak{}gradient\_\allowbreak{}memLp}, \texttt{actualObservableScore}, \texttt{actualObservableScore\_\allowbreak{}coe}}
\leansource{\texttt{actual\_\allowbreak{}gradient\_\allowbreak{}memLp} \quad WassersteinMixing/\allowbreak{}ActualScoreObservable.lean, lines 21–25}
\begin{leancode}
\leanline{0}{theorem\ actual\_gradient\_memLp\ :\ MemLp\ (gradient\ V)\ 2\ (ULA.gibbsMeasure\ V)}
\end{leancode}
\leansource{\texttt{actualObservableScore} \quad WassersteinMixing/\allowbreak{}ActualScoreObservable.lean, lines 28–29}
\begin{leancode}
\leanline{0}{def\ actualObservableScore\ :\ Lp\ (E\ P.dimension)\ 2\ (ULA.gibbsMeasure\ V)\ :=}
\leanline{2}{(actual\_gradient\_memLp\ P\ V\ hV).toLp\ (gradient\ V)}
\end{leancode}
\leansource{\texttt{actualObservableScore\_\allowbreak{}coe} \quad WassersteinMixing/\allowbreak{}ActualScoreObservable.lean, lines 31–33}
\begin{leancode}
\leanline{0}{theorem\ actualObservableScore\_coe\ :}
\leanline{4}{actualObservableScore\ P\ V\ hV\ =ᵐ[ULA.gibbsMeasure\ V]\ gradient\ V}
\end{leancode}
\leanlabel{How the Lean proof follows the argument}
\texttt{actual\_\allowbreak{}gradient\_\allowbreak{}mem\allowbreak{}Lp} specializes the Gibbs moment theorem with the actual potential constants. \texttt{actual\allowbreak{}Observable\allowbreak{}Score} uses the resulting \texttt{to\allowbreak{}Lp} constructor and \texttt{actual\allowbreak{}Observable\allowbreak{}Score\_\allowbreak{}coe} records its precise representative.
\end{leanbox}


Earlier results used here: \cite[Section~3.26.10]{LangevinLean2026ULASC}; \cite[Section~3.26.3]{LangevinLean2026ULASC}; \cite[Section~3.27.7]{LangevinLean2026ULASC}.
\subsection{\texorpdfstring{Evaluating the actual score without a new identification hypothesis}{Evaluating the actual score without a new identification hypothesis}}
\label{report-step-WassersteinMixing-ActualScoreObservable-lean-2}


\textbf{Claim.} The stationary evaluation map \(U_n f=f(Y_{nh})\) is an isometry, agrees with global evaluation at \(nh\), and satisfies \(U_nf=\nabla V(Y_{nh})\). Consequently \(\|f\|_{L^2(\pi)}^2\le\beta d\).



\textbf{Proof.} The actual grid state has law \(\pi\), giving the isometry. The proved global/\allowbreak{}grid path identity identifies the two evaluation maps. Pull the gradient representative through this measure-preserving state map to identify its evaluated class with the actual drift class. Apply the actual score moment at one time and use the isometry.

\begin{leanbox}{Lean: \texttt{actualGridEvaluation}, \texttt{actualGridEvaluation\_\allowbreak{}eq\_\allowbreak{}global}, \texttt{actualGridEvaluation\_\allowbreak{}score} and 1 more}
\leansource{\texttt{actualGridEvaluation} \quad WassersteinMixing/\allowbreak{}ActualScoreObservable.lean, lines 36–40}
\begin{leancode}
\leanline{0}{def\ actualGridEvaluation\ (n\ :\ ℕ)\ :}
\leanline{4}{Lp\ (E\ P.dimension)\ 2\ (ULA.gibbsMeasure\ V)\ →ₗᵢ[ℝ]}
\leanline{6}{Lp\ (E\ P.dimension)\ 2\ (actualProbability\ P\ V)\ :=}
\leanline{2}{StationaryObservables.evaluation\ (ULA.gibbsMeasure\ V)\ (actualProbability\ P\ V)}
\leanline{4}{(actualReference\ P\ V\ hV\ n)\ (actualReference\_stationary\ P\ V\ hV\ n)}
\end{leancode}
\leansource{\texttt{actualGridEvaluation\_\allowbreak{}eq\_\allowbreak{}global} \quad WassersteinMixing/\allowbreak{}ActualScoreObservable.lean, lines 42–46}
\begin{leancode}
\leanline{0}{theorem\ actualGridEvaluation\_eq\_global\ (n\ :\ ℕ)\ :}
\leanline{4}{actualGridEvaluation\ P\ V\ hV\ n\ =}
\leanline{6}{CanonicalSemigroup.evaluation\ ({\leanfallback 𝕜}\ :=\ ℝ)\ (F\ :=\ E\ P.dimension)\ P\ V\ hV\ (n\ *\ P.step)}
\end{leancode}
\leansource{\texttt{actualGridEvaluation\_\allowbreak{}score} \quad WassersteinMixing/\allowbreak{}ActualScoreObservable.lean, lines 49–61}
\begin{leancode}
\leanline{0}{theorem\ actualGridEvaluation\_score\ (n\ :\ ℕ)\ :}
\leanline{4}{actualGridEvaluation\ P\ V\ hV\ n\ (actualObservableScore\ P\ V\ hV)\ =\ actualScore\ P\ V\ hV\ n}
\end{leancode}
\leansource{\texttt{actualObservableScore\_\allowbreak{}energy} \quad WassersteinMixing/\allowbreak{}ActualScoreObservable.lean, lines 64–69}
\begin{leancode}
\leanline{0}{theorem\ actualObservableScore\_energy\ :}
\leanline{4}{‖actualObservableScore\ P\ V\ hV‖\ \textasciicircum{}\ 2\ ≤\ P.beta\ *\ P.dimension}
\end{leancode}
\leanlabel{How the Lean proof follows the argument}
\texttt{actual\allowbreak{}Grid\allowbreak{}Evaluation\_\allowbreak{}score} transports the almost-everywhere target representative through the actual stationary map and compares it with \texttt{lp\allowbreak{}Drift\_\allowbreak{}coe}. \texttt{actual\allowbreak{}Observable\allowbreak{}Score\_\allowbreak{}energy} transfers the actual probability-space moment through the evaluation isometry.
\end{leanbox}


Earlier results used here: \cite[Section~3.13.3]{LangevinLean2026ULASC}; \cite[Section~3.26.3]{LangevinLean2026ULASC}.
\section{\texorpdfstring{The actual Langevin Poisson comparison process}{The actual Langevin Poisson comparison process}}
\label{proof-module-WassersteinMixing-ActualPoisson-lean}


The actual reversible canonical semigroup, its proved real-coordinate preservation, its bounded constructed corrector and the actual Brownian filtration combine into a genuine Poisson process for the physical real score observables. All analytic components are derived from the potential.

\subsection{\texorpdfstring{The concrete complex observable space and actual real extraction}{The concrete complex observable space and actual real extraction}}
\label{report-step-WassersteinMixing-ActualPoisson-lean-1}


\textbf{Construction.} Embed the physical real observable space isometrically in complex coordinate \(L^2(\pi)\), and use the actual canonical reversible semigroup there. The actual transition preserves the real image, and its adjoint extraction equals the true real canonical transition.



\textbf{Proof.} The complex coordinate inclusion is an exact isometry. Actual evaluation compression intertwines it with the real operator, proving preservation. The extraction identity follows by canceling the isometric embedding with its adjoint. Thus the corrector constructed on the complex space acts on genuine physical real observables.

\begin{leanbox}{Lean: \texttt{ComplexObservableValue}, \texttt{actualObservableSemigroup}, \texttt{actualRealEmbedding} and 2 more}
\leansource{\texttt{ComplexObservableValue} \quad WassersteinMixing/\allowbreak{}ActualPoisson.lean, lines 22–23}
\begin{leancode}
\leanline{0}{abbrev\ ComplexObservableValue\ (P\ :\ Parameters)\ :=}
\leanline{2}{EuclideanSpace\ ℂ\ (Fin\ (Module.finrank\ ℝ\ (E\ P.dimension)))}
\end{leancode}
\leansource{\texttt{actualObservableSemigroup} \quad WassersteinMixing/\allowbreak{}ActualPoisson.lean, lines 25–28}
\begin{leancode}
\leanline{0}{def\ actualObservableSemigroup\ (P\ :\ Parameters)\ (V\ :\ E\ P.dimension\ →\ ℝ)}
\leanline{4}{(hV\ :\ PotentialAssumptions\ P\ V)\ :}
\leanline{4}{ReversibleSemigroup\ (Lp\ (ComplexObservableValue\ P)\ 2\ (ULA.gibbsMeasure\ V))\ :=}
\leanline{2}{CanonicalSemigroup.semigroup\ P\ V\ hV}
\end{leancode}
\leansource{\texttt{actualRealEmbedding} \quad WassersteinMixing/\allowbreak{}ActualPoisson.lean, lines 30–33}
\begin{leancode}
\leanline{0}{def\ actualRealEmbedding\ (P\ :\ Parameters)\ (V\ :\ E\ P.dimension\ →\ ℝ)\ :}
\leanline{4}{Lp\ (E\ P.dimension)\ 2\ (ULA.gibbsMeasure\ V)\ →ₗᵢ[ℝ]}
\leanline{6}{Lp\ (ComplexObservableValue\ P)\ 2\ (ULA.gibbsMeasure\ V)\ :=}
\leanline{2}{RealComplexEmbedding.coordinateLp\ (E\ P.dimension)\ (ULA.gibbsMeasure\ V)}
\end{leancode}
\leansource{\texttt{actualObservableSemigroup\_\allowbreak{}preservesReal} \quad WassersteinMixing/\allowbreak{}ActualPoisson.lean, lines 35–40}
\begin{leancode}
\leanline{0}{theorem\ actualObservableSemigroup\_preservesReal\ (P\ :\ Parameters)\ (V\ :\ E\ P.dimension\ →\ ℝ)}
\leanline{4}{(hV\ :\ PotentialAssumptions\ P\ V)\ :}
\leanline{4}{ComplexRealCorrector.PreservesReal\ (actualObservableSemigroup\ P\ V\ hV)}
\leanline{6}{(actualRealEmbedding\ P\ V)}
\end{leancode}
\leansource{\texttt{actualRealOp\_\allowbreak{}eq\_\allowbreak{}op} \quad WassersteinMixing/\allowbreak{}ActualPoisson.lean, lines 42–47}
\begin{leancode}
\leanline{0}{theorem\ actualRealOp\_eq\_op\ (P\ :\ Parameters)\ (V\ :\ E\ P.dimension\ →\ ℝ)}
\leanline{4}{(hV\ :\ PotentialAssumptions\ P\ V)\ (t\ :\ ℝ)\ :}
\leanline{4}{ComplexRealCorrector.realOp\ (actualObservableSemigroup\ P\ V\ hV)\ (actualRealEmbedding\ P\ V)\ t\ =}
\leanline{6}{CanonicalSemigroup.op\ ({\leanfallback 𝕜}\ :=\ ℝ)\ (F\ :=\ E\ P.dimension)\ P\ V\ hV\ t}
\end{leancode}
\leanlabel{How the Lean proof follows the argument}
\texttt{actual\allowbreak{}Observable\allowbreak{}Semigroup} is the proved canonical semigroup instance, and \texttt{actual\allowbreak{}Real\allowbreak{}Embedding} is the actual coordinate isometry. The preservation and extraction theorems instantiate the already proved actual real-range bridge.
\end{leanbox}


Earlier results used here: \cite[Section~3.26.3]{LangevinLean2026ULASC}.
\subsection{\texorpdfstring{The actual filtration, evaluations and Markov law}{The actual filtration, evaluations and Markov law}}
\label{report-step-WassersteinMixing-ActualPoisson-lean-2}


\textbf{Claim.} Actual grid evaluations are adapted to the actual Brownian filtration and satisfy the exact one-step conditional Markov identity with the extracted transition. They therefore form a Poisson process with the constructed corrector, source identity, norm bound and energy domination.



\textbf{Proof.} Specialize the actual arbitrary-time adaptation and Markov theorems to \(nh\) and step \(h\). The true grid/\allowbreak{}global evaluation identity identifies both endpoints, and the extracted transition is already equal to the actual real operator. Supply these derived facts to the semigroup-corrector process construction.

\begin{leanbox}{Lean: \texttt{actualGridEvaluation\_\allowbreak{}adapted}, \texttt{actualGridEvaluation\_\allowbreak{}markov}, \texttt{actualPoissonProcess}}
\leansource{\texttt{actualGridEvaluation\_\allowbreak{}adapted} \quad WassersteinMixing/\allowbreak{}ActualPoisson.lean, lines 49–60}
\begin{leancode}
\leanline{0}{theorem\ actualGridEvaluation\_adapted\ (P\ :\ Parameters)\ (V\ :\ E\ P.dimension\ →\ ℝ)}
\leanline{4}{(hV\ :\ PotentialAssumptions\ P\ V)\ (n\ :\ ℕ)\ (f\ :\ Lp\ (E\ P.dimension)\ 2\ (ULA.gibbsMeasure\ V))\ :}
\leanline{4}{AEStronglyMeasurable[actualFiltration\ P\ V\ hV\ n]}
\leanline{6}{(actualGridEvaluation\ P\ V\ hV\ n\ f)\ (actualProbability\ P\ V)}
\end{leancode}
\leansource{\texttt{actualGridEvaluation\_\allowbreak{}markov} \quad WassersteinMixing/\allowbreak{}ActualPoisson.lean, lines 62–77}
\begin{leancode}
\leanline{0}{theorem\ actualGridEvaluation\_markov\ (P\ :\ Parameters)\ (V\ :\ E\ P.dimension\ →\ ℝ)}
\leanline{4}{(hV\ :\ PotentialAssumptions\ P\ V)\ (n\ :\ ℕ)\ (f\ :\ Lp\ (E\ P.dimension)\ 2\ (ULA.gibbsMeasure\ V))\ :}
\leanline{4}{(condExpL2\ (E\ P.dimension)\ ℝ\ ((actualFiltration\ P\ V\ hV).le\ n)}
\leanline{6}{(actualGridEvaluation\ P\ V\ hV\ (n\ +\ 1)\ f)\ :\ Lp\ (E\ P.dimension)\ 2\ (actualProbability\ P\ V))\ =}
\leanline{8}{actualGridEvaluation\ P\ V\ hV\ n}
\leanline{10}{(ComplexRealCorrector.realOp\ (actualObservableSemigroup\ P\ V\ hV)}
\leanline{12}{(actualRealEmbedding\ P\ V)\ P.step\ f)}
\end{leancode}
\leansource{\texttt{actualPoissonProcess} \quad WassersteinMixing/\allowbreak{}ActualPoisson.lean, lines 81–88}
\begin{leancode}
\leanline{0}{def\ actualPoissonProcess\ (P\ :\ Parameters)\ (V\ :\ E\ P.dimension\ →\ ℝ)}
\leanline{4}{(hV\ :\ PotentialAssumptions\ P\ V)\ :}
\leanline{4}{PoissonProcess\ (Lp\ (E\ P.dimension)\ 2\ (ULA.gibbsMeasure\ V))\ (E\ P.dimension)}
\leanline{6}{(actualProbability\ P\ V)\ P.step\ P.step\_pos\ :=}
\leanline{2}{semigroupPoissonProcess\ (actualProbability\ P\ V)\ (actualObservableSemigroup\ P\ V\ hV)}
\leanline{4}{(actualRealEmbedding\ P\ V)\ (actualObservableSemigroup\_preservesReal\ P\ V\ hV)}
\leanline{4}{P.step\ P.step\_pos\ (actualFiltration\ P\ V\ hV)\ (actualGridEvaluation\ P\ V\ hV)}
\leanline{4}{(actualGridEvaluation\_adapted\ P\ V\ hV)\ (actualGridEvaluation\_markov\ P\ V\ hV)}
\end{leancode}
\leanlabel{How the Lean proof follows the argument}
\texttt{actual\allowbreak{}Grid\allowbreak{}Evaluation\_\allowbreak{}adapted} and \texttt{actual\allowbreak{}Grid\allowbreak{}Evaluation\_\allowbreak{}markov} specialize the checked actual arbitrary-time theorems. \texttt{actual\allowbreak{}Poisson\allowbreak{}Process} invokes \texttt{semigroup\allowbreak{}Poisson\allowbreak{}Process} with only constructed objects and derived analytic properties.
\end{leanbox}


Earlier results used here: \cite[Section~3.26.3]{LangevinLean2026ULASC}.
\section{\texorpdfstring{Conditional expectation of a genuine integrated defect}{Conditional expectation of a genuine integrated defect}}
\label{proof-module-WassersteinMixing-ConditionalDefect-lean}


The actual Hilbert conditional-expectation projection commutes with Bochner time integration. If each increment has conditional mean equal to an evaluated source, the time-integrated defect has conditional mean equal to the evaluated integrated source.

\subsection{\texorpdfstring{Commuting the actual bounded projections and time integral}{Commuting the actual bounded projections and time integral}}
\label{report-step-WassersteinMixing-ConditionalDefect-lean-1}


\textbf{Claim.} If \(\mathbb E[I(t)\mid\mathcal F]=U s(t)\) for \(0\le t\le h\), where \(U\) is an evaluation isometry and the increment curve is integrable, then

\[\mathbb E\left[\int_0^h I(t)\,dt\,\middle|\,\mathcal F\right]=U\left(\int_0^h s(t)\,dt\right).\]

\textbf{Proof.} Express conditional expectation as the orthogonal projection onto the closed past-measurable \(L^2\) subspace. It is bounded linear, so it commutes with the justified integral. Substitute the pointwise conditional-mean identities inside the integral. Evaluation is also bounded linear and commutes with the same integral, yielding the displayed result.

\begin{leanbox}{Lean: \texttt{conditionalDefect\_\allowbreak{}eq\_\allowbreak{}evaluation\_\allowbreak{}integral}}
\leansource{\texttt{conditionalDefect\_\allowbreak{}eq\_\allowbreak{}evaluation\_\allowbreak{}integral} \quad WassersteinMixing/\allowbreak{}ConditionalDefect.lean, lines 24–39}
\begin{leancode}
\leanline{0}{theorem\ conditionalDefect\_eq\_evaluation\_integral}
\leanline{4}{(eval\ :\ R\ →ₗᵢ[ℝ]\ Lp\ E\ 2\ Q)\ (increment\ :\ ℝ\ →\ Lp\ E\ 2\ Q)\ (source\ :\ ℝ\ →\ R)}
\leanline{4}{(h\ :\ ℝ)\ (hh\ :\ 0\ ≤\ h)\ (hi\ :\ IntervalIntegrable\ increment\ volume\ 0\ h)}
\leanline{4}{(hmean\ :\ ∀\ t\ ∈\ Set.Icc\ 0\ h,}
\leanline{6}{(condExpL2\ E\ ℝ\ hm\ (increment\ t)\ :\ Lp\ E\ 2\ Q)\ =\ eval\ (source\ t))\ :}
\leanline{4}{(condExpL2\ E\ ℝ\ hm\ (stationaryLocalDefect\ increment\ h)\ :\ Lp\ E\ 2\ Q)\ =}
\leanline{6}{eval\ (∫\ t\ in\ 0..h,\ source\ t)}
\end{leancode}
\leanlabel{How the Lean proof follows the argument}
\texttt{conditional\allowbreak{}Defect\_\allowbreak{}eq\_\allowbreak{}evaluation\_\allowbreak{}integral} uses the actual \texttt{lp\allowbreak{}Meas} orthogonal projection and its interval-integral commutation theorem, then the evaluation isometry integral identity. Pointwise substitution is justified on the full closed time interval.
\end{leanbox}
\section{\texorpdfstring{The exact conditional mean of the actual Euler defect}{The exact conditional mean of the actual Euler defect}}
\label{proof-module-WassersteinMixing-ActualDefectMean-lean}


The actual score is the true target gradient evaluated on the actual diffusion. The proved conditional Markov identity identifies the mean of every genuine score increment, and integration gives the precise semigroup source for the actual local Euler defect.

\subsection{\texorpdfstring{Identifying the true score evaluations}{Identifying the true score evaluations}}
\label{report-step-WassersteinMixing-ActualDefectMean-lean-1}


\textbf{Claim.} For the actual target observable \(f=\nabla V\), \(U_tf=\nabla V(Y_t)\), and the actual local increment is \(I_n(s)=U_{nh}f-U_{nh+s}f\).



\textbf{Proof.} Transport the actual gradient representative through the measure-preserving Gibbs state map, and compare it with the true \(L^2\) drift representative. The proved global/\allowbreak{}segment path identity then identifies the two actual score evaluations defining the local increment.

\begin{leanbox}{Lean: \texttt{actualEvaluation\_\allowbreak{}score}, \texttt{actualScoreIncrement\_\allowbreak{}eq\_\allowbreak{}evaluation}}
\leansource{\texttt{actualEvaluation\_\allowbreak{}score} \quad WassersteinMixing/\allowbreak{}ActualDefectMean.lean, lines 23–38}
\begin{leancode}
\leanline{0}{theorem\ actualEvaluation\_score\ (t\ :\ ℝ)\ :}
\leanline{4}{letI\ :=\ actualProbability\_isProbability\ P\ V\ hV}
\leanline{4}{CanonicalSemigroup.evaluation\ ({\leanfallback 𝕜}\ :=\ ℝ)\ P\ V\ hV\ t\ (actualObservableScore\ P\ V\ hV)\ =}
\leanline{6}{ULA.lpDrift\ (actualProbability\ P\ V)\ (gradient\ V)\ hV.gradientLipschitz}
\leanline{8}{(CanonicalSemigroup.state\ P\ V\ hV\ t)}
\end{leancode}
\leansource{\texttt{actualScoreIncrement\_\allowbreak{}eq\_\allowbreak{}evaluation} \quad WassersteinMixing/\allowbreak{}ActualDefectMean.lean, lines 41–49}
\begin{leancode}
\leanline{0}{theorem\ actualScoreIncrement\_eq\_evaluation\ (n\ :\ ℕ)\ (s\ :\ ℝ)\ (hs\ :\ s\ ∈\ Icc\ 0\ P.step)\ :}
\leanline{4}{actualScoreIncrement\ P\ V\ hV\ n\ s\ =}
\leanline{6}{actualGridEvaluation\ P\ V\ hV\ n\ (actualObservableScore\ P\ V\ hV)\ -}
\leanline{8}{CanonicalSemigroup.evaluation\ ({\leanfallback 𝕜}\ :=\ ℝ)\ P\ V\ hV\ ((n\ :\ ℝ)\ *\ P.step\ +\ s)}
\leanline{10}{(actualObservableScore\ P\ V\ hV)}
\end{leancode}
\leanlabel{How the Lean proof follows the argument}
\texttt{actual\allowbreak{}Evaluation\_\allowbreak{}score} compares the exact target-gradient and drift representatives. \texttt{actual\allowbreak{}Score\allowbreak{}Increment\_\allowbreak{}eq\_\allowbreak{}evaluation} uses the true grid evaluation and the proved canonical/\allowbreak{}segment path equality.
\end{leanbox}


Earlier results used here: \cite[Section~3.13.2]{LangevinLean2026ULASC}; \cite[Section~3.13.3]{LangevinLean2026ULASC}; \cite[Section~3.26.3]{LangevinLean2026ULASC}.
\subsection{\texorpdfstring{The actual conditional source integral}{The actual conditional source integral}}
\label{report-step-WassersteinMixing-ActualDefectMean-lean-2}


\textbf{Claim.}

\[\mathbb E[I_n(s)\mid\mathcal F_n]=U_{nh}(f-P_sf),\qquad \mathbb E[D_n\mid\mathcal F_n]=U_{nh}\int_0^h(f-P_sf)\,ds.\]

\textbf{Proof.} The current evaluation is fixed by conditioning. Apply the actual conditional Markov theorem to the later score evaluation and subtract, obtaining the increment mean. The actual increment curve is continuous and hence integrable on the compact step interval. Commute the conditional projection and evaluation through its true integral to obtain the defect mean.

\begin{leanbox}{Lean: \texttt{actualScoreIncrement\_\allowbreak{}conditionalMean}, \texttt{actualDefect\_\allowbreak{}conditionalMean}}
\leansource{\texttt{actualScoreIncrement\_\allowbreak{}conditionalMean} \quad WassersteinMixing/\allowbreak{}ActualDefectMean.lean, lines 53–76}
\begin{leancode}
\leanline{0}{theorem\ actualScoreIncrement\_conditionalMean\ (n\ :\ ℕ)\ (s\ :\ ℝ)\ (hs\ :\ s\ ∈\ Icc\ 0\ P.step)\ :}
\leanline{4}{(condExpL2\ (E\ P.dimension)\ ℝ\ (actualPast\_le\ P\ V\ hV\ n)}
\leanline{6}{(actualScoreIncrement\ P\ V\ hV\ n\ s)\ :\ Lp\ (E\ P.dimension)\ 2\ (actualProbability\ P\ V))\ =}
\leanline{6}{actualGridEvaluation\ P\ V\ hV\ n}
\leanline{8}{(actualObservableScore\ P\ V\ hV\ -}
\leanline{10}{CanonicalSemigroup.op\ ({\leanfallback 𝕜}\ :=\ ℝ)\ P\ V\ hV\ s\ (actualObservableScore\ P\ V\ hV))}
\end{leancode}
\leansource{\texttt{actualDefect\_\allowbreak{}conditionalMean} \quad WassersteinMixing/\allowbreak{}ActualDefectMean.lean, lines 81–92}
\begin{leancode}
\leanline{0}{theorem\ actualDefect\_conditionalMean\ (n\ :\ ℕ)\ :}
\leanline{4}{(condExpL2\ (E\ P.dimension)\ ℝ\ (actualPast\_le\ P\ V\ hV\ n)}
\leanline{6}{(actualDefect\ P\ V\ hV\ n)\ :\ Lp\ (E\ P.dimension)\ 2\ (actualProbability\ P\ V))\ =}
\leanline{6}{actualGridEvaluation\ P\ V\ hV\ n}
\leanline{8}{(∫\ s\ in\ 0..P.step,\ actualObservableScore\ P\ V\ hV\ -}
\leanline{10}{CanonicalSemigroup.op\ ({\leanfallback 𝕜}\ :=\ ℝ)\ P\ V\ hV\ s\ (actualObservableScore\ P\ V\ hV))}
\end{leancode}
\leanlabel{How the Lean proof follows the argument}
\texttt{actual\allowbreak{}Score\allowbreak{}Increment\_\allowbreak{}conditional\allowbreak{}Mean} applies the actual Markov theorem at the true grid time, including its zero-time identity. \texttt{actual\allowbreak{}Defect\_\allowbreak{}conditional\allowbreak{}Mean} invokes \texttt{conditional\allowbreak{}Defect\_\allowbreak{}eq\_\allowbreak{}evaluation\_\allowbreak{}integral} with the actual continuous score-increment curve.
\end{leanbox}
\section{\texorpdfstring{The shifted energy floor and the coefficient six}{The shifted energy floor and the coefficient six}}
\label{proof-module-WassersteinMixing-Recurrence-lean}


Let \(q=1-\alpha h\), \(C=h^2\beta d\), \(F=19h^3\beta^2d\) and \(E_*=19h^2\kappa\beta d\), with \(\kappa=\beta/\alpha\). The exact fixed-point identity \(qE_*+F=E_*\) turns a supplied scalar energy recurrence into uniform and finite-time bounds. The final norm comparison uses only triangle inequalities and these two squared budgets.



Earlier results used here: \cite[Section~3.15]{LangevinLean2026ULASC}.
\subsection{\texorpdfstring{The budgets and their exact fixed point}{The budgets and their exact fixed point}}
\label{report-step-WassersteinMixing-Recurrence-lean-1}


\textbf{Claim.} Under the numerical parameter restrictions, \(C\ge0\), \(E_*\ge0\), \(C\le E_*\), and \(qE_*+F=E_*\).




\textbf{Proof.} The budgets are products of nonnegative factors. Since \(\kappa\ge1\), \(E_*=19\kappa C\ge C\). For the fixed point, \((1-q)E_*=\alpha h\cdot19h^2(\beta/\alpha)\beta d=19h^3\beta^2d=F\). Division by \(\alpha\) is valid because \(\alpha>0\).

\begin{leanbox}{Lean: \texttt{Parameters.\allowbreak{}correctorBudget}, \texttt{Parameters.\allowbreak{}energyForcing}, \texttt{Parameters.\allowbreak{}energyFloor} and 4 more}
\leansource{\texttt{Parameters.\allowbreak{}correctorBudget} \quad WassersteinMixing/\allowbreak{}Recurrence.lean, lines 17–18}
\begin{leancode}
\leanline{0}{def\ Parameters.correctorBudget\ (P\ :\ Parameters)\ :\ ℝ\ :=}
\leanline{2}{P.step\ \textasciicircum{}\ 2\ *\ P.beta\ *\ P.dimension}
\end{leancode}
\leansource{\texttt{Parameters.\allowbreak{}energyForcing} \quad WassersteinMixing/\allowbreak{}Recurrence.lean, lines 21–22}
\begin{leancode}
\leanline{0}{def\ Parameters.energyForcing\ (P\ :\ Parameters)\ :\ ℝ\ :=}
\leanline{2}{19\ *\ P.step\ \textasciicircum{}\ 3\ *\ P.beta\ \textasciicircum{}\ 2\ *\ P.dimension}
\end{leancode}
\leansource{\texttt{Parameters.\allowbreak{}energyFloor} \quad WassersteinMixing/\allowbreak{}Recurrence.lean, lines 25–26}
\begin{leancode}
\leanline{0}{def\ Parameters.energyFloor\ (P\ :\ Parameters)\ :\ ℝ\ :=}
\leanline{2}{19\ *\ P.step\ \textasciicircum{}\ 2\ *\ P.condition\ *\ P.beta\ *\ P.dimension}
\end{leancode}
\leansource{\texttt{Parameters.\allowbreak{}correctorBudget\_\allowbreak{}nonneg} \quad WassersteinMixing/\allowbreak{}Recurrence.lean, lines 28–30}
\begin{leancode}
\leanline{0}{theorem\ Parameters.correctorBudget\_nonneg\ (P\ :\ Parameters)\ :\ 0\ ≤\ P.correctorBudget}
\end{leancode}
\leansource{\texttt{Parameters.\allowbreak{}energyFloor\_\allowbreak{}nonneg} \quad WassersteinMixing/\allowbreak{}Recurrence.lean, lines 32–36}
\begin{leancode}
\leanline{0}{theorem\ Parameters.energyFloor\_nonneg\ (P\ :\ Parameters)\ :\ 0\ ≤\ P.energyFloor}
\end{leancode}
\leansource{\texttt{Parameters.\allowbreak{}energy\_\allowbreak{}fixed\_\allowbreak{}point} \quad WassersteinMixing/\allowbreak{}Recurrence.lean, lines 39–44}
\begin{leancode}
\leanline{0}{theorem\ Parameters.energy\_fixed\_point\ (P\ :\ Parameters)\ :}
\leanline{4}{P.contraction\ *\ P.energyFloor\ +\ P.energyForcing\ =\ P.energyFloor}
\end{leancode}
\leansource{\texttt{Parameters.\allowbreak{}correctorBudget\_\allowbreak{}le\_\allowbreak{}floor} \quad WassersteinMixing/\allowbreak{}Recurrence.lean, lines 46–56}
\begin{leancode}
\leanline{0}{theorem\ Parameters.correctorBudget\_le\_floor\ (P\ :\ Parameters)\ :}
\leanline{4}{P.correctorBudget\ ≤\ P.energyFloor}
\end{leancode}
\leanlabel{How the Lean proof follows the argument}
The three parameter definitions retain the exact constants. The sign lemmas use positive curvature and nonnegative dimension. \texttt{energy\_\allowbreak{}fixed\_\allowbreak{}point} clears the nonzero denominator and proves the polynomial identity. \texttt{corrector\allowbreak{}Budget\_\allowbreak{}le\_\allowbreak{}floor} rewrites the floor as \(19\kappa\) times the budget.
\end{leanbox}
\subsection{\texorpdfstring{Uniform and finite-time energy envelopes}{Uniform and finite-time energy envelopes}}
\label{report-step-WassersteinMixing-Recurrence-lean-2}


\textbf{Claim.} Suppose \(E_{n+1}\le qE_n+F\). If \(E_0\le C\), then \(E_n\le E_*\) for every \(n\). Without the initial restriction,

\[E_n\le q^nE_0+(1-q^n)E_*.\]


\textbf{Proof.} The uniform bound follows by induction. Initially \(E_0\le C\le E_*\). Multiplication of the induction hypothesis by \(q\ge0\), followed by the recurrence and fixed-point identity, gives \(E_{n+1}\le qE_*+F=E_*\). For the finite-time bound, iterating the affine recurrence yields its geometric sum. Substitute \(F=(1-q)E_*\); since \(q<1\), cancellation of \(1-q\) gives the displayed expression.

\begin{leanbox}{Lean: \texttt{shifted\_\allowbreak{}energy\_\allowbreak{}uniform}, \texttt{shifted\_\allowbreak{}energy\_\allowbreak{}finite}}
\leansource{\texttt{shifted\_\allowbreak{}energy\_\allowbreak{}uniform} \quad WassersteinMixing/\allowbreak{}Recurrence.lean, lines 60–69}
\begin{leancode}
\leanline{0}{theorem\ shifted\_energy\_uniform\ (P\ :\ Parameters)\ (energy\ :\ ℕ\ →\ ℝ)}
\leanline{4}{(hzero\ :\ energy\ 0\ ≤\ P.correctorBudget)}
\leanline{4}{(hstep\ :\ ∀\ n,\ energy\ (n\ +\ 1)\ ≤\ P.contraction\ *\ energy\ n\ +\ P.energyForcing)}
\leanline{4}{(n\ :\ ℕ)\ :\ energy\ n\ ≤\ P.energyFloor}
\end{leancode}
\leansource{\texttt{shifted\_\allowbreak{}energy\_\allowbreak{}finite} \quad WassersteinMixing/\allowbreak{}Recurrence.lean, lines 72–84}
\begin{leancode}
\leanline{0}{theorem\ shifted\_energy\_finite\ (P\ :\ Parameters)\ (energy\ :\ ℕ\ →\ ℝ)}
\leanline{4}{(hstep\ :\ ∀\ n,\ energy\ (n\ +\ 1)\ ≤\ P.contraction\ *\ energy\ n\ +\ P.energyForcing)}
\leanline{4}{(n\ :\ ℕ)\ :}
\leanline{4}{energy\ n\ ≤\ P.contraction\ \textasciicircum{}\ n\ *\ energy\ 0\ +}
\leanline{6}{(1\ -\ P.contraction\ \textasciicircum{}\ n)\ *\ P.energyFloor}
\end{leancode}
\leanlabel{How the Lean proof follows the argument}
\texttt{shifted\_\allowbreak{}energy\_\allowbreak{}uniform} is an explicit natural-number induction. \texttt{shifted\_\allowbreak{}energy\_\allowbreak{}finite} invokes the shared exact affine-recurrence formula and substitutes the proved fixed point. These are scalar implications. The stationary comparison derives their recurrence from the constructed corrector, actual martingale centering and genuine curvature bounds; the realized comparison supplies all primitive inputs from the actual potential dynamics.
\end{leanbox}


Earlier results used here: \cite[Section~3.15.2]{LangevinLean2026ULASC}.
\subsection{\texorpdfstring{Converting the two squared budgets to the literal bias constant}{Converting the two squared budgets to the literal bias constant}}
\label{report-step-WassersteinMixing-Recurrence-lean-3}


\textbf{Claim.} If \(\|z+g\|^2\le E_*\) and \(\|g\|^2\le C\), then

\[\|z\|\le6h\sqrt{\kappa\beta d}.\]


\textbf{Proof.} Put \(R=h\sqrt{\kappa\beta d}\ge0\). Then \(E_*=19R^2\), while \(C\le R^2\) because \(\kappa\ge1\). Hence \(\|z+g\|\le5R\), using \(19\le25\), and \(\|g\|\le R\). The triangle inequality for \(z=(z+g)-g\) gives \(\|z\|\le6R\). This proves the source's coefficient \(6\) directly, including a zero-dimensional budget.

\begin{leanbox}{Lean: \texttt{shifted\_\allowbreak{}norm\_\allowbreak{}le\_\allowbreak{}bias}}
\leansource{\texttt{shifted\_\allowbreak{}norm\_\allowbreak{}le\_\allowbreak{}bias} \quad WassersteinMixing/\allowbreak{}Recurrence.lean, lines 88–121}
\begin{leancode}
\leanline{0}{theorem\ shifted\_norm\_le\_bias\ \{H\ :\ Type*\}\ [NormedAddCommGroup\ H]}
\leanline{4}{(P\ :\ Parameters)\ (z\ g\ :\ H)}
\leanline{4}{(hu\ :\ ‖z\ +\ g‖\ \textasciicircum{}\ 2\ ≤\ P.energyFloor)}
\leanline{4}{(hg\ :\ ‖g‖\ \textasciicircum{}\ 2\ ≤\ P.correctorBudget)\ :\ ‖z‖\ ≤\ P.bias}
\end{leancode}
\leanlabel{How the Lean proof follows the argument}
\texttt{shifted\_\allowbreak{}norm\_\allowbreak{}le\_\allowbreak{}bias} proves the squared identities for the nonnegative radius, compares each squared norm using \texttt{sq\_\allowbreak{}le\_\allowbreak{}sq\ensuremath{{}_{0}}}, and applies \texttt{norm\_\allowbreak{}sub\_\allowbreak{}le}. The proof does not use an unproved manipulation of a limit superior.
\end{leanbox}
\section{\texorpdfstring{The stationary comparison and the sharp shifted-energy constants}{The stationary comparison and the sharp shifted-energy constants}}
\label{proof-module-WassersteinMixing-PoissonComparison-lean}


Combine a stationary evaluated Poisson corrector with a synchronous error update, an actual curvature operator and primitive score increments. The defect, martingale, shifted energy and transport bounds are all derived. The constants \(3\), \(18\), \(19\) and \(6\) appear as conclusions.



Earlier results used here: \cite[Section~3.17]{LangevinLean2026ULASC}.
\subsection{\texorpdfstring{The primitive comparison and shifted drift}{The primitive comparison and shifted drift}}
\label{report-step-WassersteinMixing-PoissonComparison-lean-1}


\textbf{Set-up.} Let \(e_n=Y_n-X_n\), \(D_n=\int_0^h I_n(s)\,ds\), and let \(g_n\) be the evaluated corrector. Assume the exact update \(e_{n+1}=e_n-hA_ne_n+D_n\), stationary reference laws and equal initial states. The primitive score bounds are \(\|f\|^2\le\beta d\), \(\|I_n(s)\|^2\le6s\beta^2d\); the defect conditional mean is its evaluated Poisson source. Define

\[Z_n=(e_n+g_n)-hA_n(e_n+g_n)+hA_ng_n.\]

The symmetric curvature bounds and adaptation refer to the genuine \(L^2\) operator and current filtration.

\begin{leanbox}{Lean: \texttt{PoissonComparison}, \texttt{error}, \texttt{defect} and 1 more}
\leansource{\texttt{PoissonComparison} \quad WassersteinMixing/\allowbreak{}PoissonComparison.lean, lines 33–67}
\begin{leancode}
\leanline{0}{structure\ PoissonComparison\ (S\ :\ PoissonProcess\ H\ E\ Q\ P.step\ P.step\_pos)\ (score\ :\ H)\ where}
\leanline{2}{reference\ :\ ℕ\ →\ Lp\ E\ 2\ Q}
\leanline{2}{chain\ :\ ℕ\ →\ Lp\ E\ 2\ Q}
\leanline{2}{increment\ :\ ℕ\ →\ ℝ\ →\ Lp\ E\ 2\ Q}
\leanline{2}{curvature\ :\ ℕ\ →\ Lp\ E\ 2\ Q\ →L[ℝ]\ Lp\ E\ 2\ Q}
\leanline{2}{score\_bound\ :\ ‖score‖\ \textasciicircum{}\ 2\ ≤\ P.beta\ *\ P.dimension}
\leanline{2}{reference\_stationary\ :\ ∀\ n,\ Q.map\ (reference\ n)\ =\ Q.map\ (reference\ 0)}
\leanline{2}{initial\_equal\ :\ reference\ 0\ =\ chain\ 0}
\leanline{2}{increment\_memLp\ :\ ∀\ n,}
\leanline{4}{MemLp\ (increment\ n)\ 2\ (volume.restrict\ (Set.Ioc\ 0\ P.step))}
\leanline{2}{increment\_bound\ :\ ∀\ n,\ ∀\ s\ ∈\ Set.Ioc\ 0\ P.step,}
\leanline{4}{‖increment\ n\ s‖\ \textasciicircum{}\ 2\ ≤\ 6\ *\ s\ *\ P.beta\ \textasciicircum{}\ 2\ *\ P.dimension}
\leanline{2}{increment\_endpoint\ :\ ∀\ n,}
\leanline{4}{increment\ n\ P.step\ =\ S.eval\ n\ score\ -\ S.eval\ (n\ +\ 1)\ score}
\leanline{2}{defect\_mean\ :\ ∀\ n,}
\leanline{4}{(condExpL2\ E\ ℝ\ (S.filtration.le\ n)}
\leanline{6}{(stationaryLocalDefect\ (increment\ n)\ P.step)\ :\ Lp\ E\ 2\ Q)\ =}
\leanline{6}{S.eval\ n\ (S.poissonSource\ score)}
\leanline{2}{error\_step\ :\ ∀\ n,}
\leanline{4}{reference\ (n\ +\ 1)\ -\ chain\ (n\ +\ 1)\ =}
\leanline{6}{(reference\ n\ -\ chain\ n)\ -\ P.step\ •\ curvature\ n\ (reference\ n\ -\ chain\ n)\ +}
\leanline{8}{stationaryLocalDefect\ (increment\ n)\ P.step}
\leanline{2}{error\_adapted\ :\ ∀\ n,}
\leanline{4}{AEStronglyMeasurable[S.filtration\ n]\ (reference\ n\ -\ chain\ n)\ Q}
\leanline{2}{curvature\_adapted\ :\ ∀\ n\ (U\ :\ Lp\ E\ 2\ Q),}
\leanline{4}{AEStronglyMeasurable[S.filtration\ n]\ U\ Q\ →}
\leanline{6}{AEStronglyMeasurable[S.filtration\ n]\ (curvature\ n\ U)\ Q}
\leanline{2}{curvature\_symmetric\ :\ ∀\ n\ u\ v,}
\leanline{4}{inner\ ℝ\ (curvature\ n\ u)\ v\ =\ inner\ ℝ\ u\ (curvature\ n\ v)}
\leanline{2}{curvature\_lower\ :\ ∀\ n\ u,}
\leanline{4}{P.alpha\ *\ ‖u‖\ \textasciicircum{}\ 2\ ≤\ inner\ ℝ\ u\ (curvature\ n\ u)}
\leanline{2}{curvature\_upper\ :\ ∀\ n\ u,}
\leanline{4}{inner\ ℝ\ u\ (curvature\ n\ u)\ ≤\ P.beta\ *\ ‖u‖\ \textasciicircum{}\ 2}
\leanline{0}{\strut}
\leanline{0}{namespace\ PoissonComparison}
\end{leancode}
\leansource{\texttt{error} \quad WassersteinMixing/\allowbreak{}PoissonComparison.lean, lines 72–72}
\begin{leancode}
\leanline{0}{def\ error\ (n\ :\ ℕ)\ :\ Lp\ E\ 2\ Q\ :=\ C.reference\ n\ -\ C.chain\ n}
\end{leancode}
\leansource{\texttt{defect} \quad WassersteinMixing/\allowbreak{}PoissonComparison.lean, lines 75–75}
\begin{leancode}
\leanline{0}{def\ defect\ (n\ :\ ℕ)\ :\ Lp\ E\ 2\ Q\ :=\ stationaryLocalDefect\ (C.increment\ n)\ P.step}
\end{leancode}
\leansource{\texttt{predictable} \quad WassersteinMixing/\allowbreak{}PoissonComparison.lean, lines 78–83}
\begin{leancode}
\leanline{0}{def\ predictable\ (n\ :\ ℕ)\ :\ Lp\ E\ 2\ Q\ :=}
\leanline{2}{(C.error\ n\ +\ S.correctorAt\ score\ n)\ -}
\leanline{4}{P.step\ •\ C.curvature\ n\ (C.error\ n\ +\ S.correctorAt\ score\ n)\ +}
\leanline{6}{P.step\ •\ C.curvature\ n\ (S.correctorAt\ score\ n)}
\leanline{0}{\strut}
\leanline{0}{include\ C\ in}
\end{leancode}
\leanlabel{How the Lean proof follows the argument}
\texttt{Poisson\allowbreak{}Comparison} contains the primitive random variables, integral increments, exact equation, moment/\allowbreak{}conditional-mean data, and curvature/\allowbreak{}adaptation facts. \texttt{error}, \texttt{defect} and \texttt{predictable} are their literal expressions. No corrector, martingale or final transport bound is a field.
\end{leanbox}
\subsection{\texorpdfstring{Deriving the component bounds and adaptation}{Deriving the component bounds and adaptation}}
\label{report-step-WassersteinMixing-PoissonComparison-lean-2}


\textbf{Claim.}

\[\|g_n\|^2\le h^2\beta d,\quad\|D_n\|^2\le3h^3\beta^2d,\quad\|M_n\|^2\le18h^3\beta^2d.\]

\textbf{Proof.} Use the corrector norm bound and score moment for the first estimate. Time Cauchy–Schwarz and integration of the primitive increment bound give the second. The endpoint increment identifies the actual score difference, so the process corrector-increment bound and squared triangle inequality give the third. The shifted drift is adapted because the error and corrector are adapted and the actual curvature preserves the current measurable subspace.

\begin{leanbox}{Lean: \texttt{corrector\_\allowbreak{}energy}, \texttt{defect\_\allowbreak{}energy}, \texttt{score\_\allowbreak{}increment\_\allowbreak{}energy} and 2 more}
\leansource{\texttt{corrector\_\allowbreak{}energy} \quad WassersteinMixing/\allowbreak{}PoissonComparison.lean, lines 84–87}
\begin{leancode}
\leanline{0}{theorem\ corrector\_energy\ (n\ :\ ℕ)\ :}
\leanline{4}{‖S.correctorAt\ score\ n‖\ \textasciicircum{}\ 2\ ≤\ P.correctorBudget}
\end{leancode}
\leansource{\texttt{defect\_\allowbreak{}energy} \quad WassersteinMixing/\allowbreak{}PoissonComparison.lean, lines 89–93}
\begin{leancode}
\leanline{0}{theorem\ defect\_energy\ (n\ :\ ℕ)\ :\ ‖C.defect\ n‖\ \textasciicircum{}\ 2\ ≤\ 3\ *\ P.step\ \textasciicircum{}\ 3\ *\ P.beta\ \textasciicircum{}\ 2\ *\ P.dimension}
\end{leancode}
\leansource{\texttt{score\_\allowbreak{}increment\_\allowbreak{}energy} \quad WassersteinMixing/\allowbreak{}PoissonComparison.lean, lines 94–99}
\begin{leancode}
\leanline{0}{theorem\ score\_increment\_energy\ (n\ :\ ℕ)\ :}
\leanline{4}{‖S.eval\ (n\ +\ 1)\ score\ -\ S.eval\ n\ score‖\ \textasciicircum{}\ 2\ ≤}
\leanline{6}{6\ *\ P.step\ *\ P.beta\ \textasciicircum{}\ 2\ *\ P.dimension}
\end{leancode}
\leansource{\texttt{martingale\_\allowbreak{}energy} \quad WassersteinMixing/\allowbreak{}PoissonComparison.lean, lines 101–104}
\begin{leancode}
\leanline{0}{theorem\ martingale\_energy\ (n\ :\ ℕ)\ :}
\leanline{4}{‖S.martingaleAt\ score\ C.defect\ n‖\ \textasciicircum{}\ 2\ ≤\ 18\ *\ P.step\ \textasciicircum{}\ 3\ *\ P.beta\ \textasciicircum{}\ 2\ *\ P.dimension}
\end{leancode}
\leansource{\texttt{predictable\_\allowbreak{}adapted} \quad WassersteinMixing/\allowbreak{}PoissonComparison.lean, lines 106–122}
\begin{leancode}
\leanline{0}{theorem\ predictable\_adapted\ (n\ :\ ℕ)\ :}
\leanline{4}{AEStronglyMeasurable[S.filtration\ n]\ (C.predictable\ n)\ Q}
\end{leancode}
\leanlabel{How the Lean proof follows the argument}
The component energy theorems invoke the constructed process and local integral estimates. \texttt{score\_\allowbreak{}increment\_\allowbreak{}energy} uses the exact endpoint identity. \texttt{predictable\_\allowbreak{}adapted} proves closure inside the measurable submodule under the actual curvature action, sums and scalar multiplication.
\end{leanbox}
\subsection{\texorpdfstring{The coefficient nineteen and the final constant six}{The coefficient nineteen and the final constant six}}
\label{report-step-WassersteinMixing-PoissonComparison-lean-3}


\textbf{Claim.} With \(u_n=\|e_n+g_n\|^2\),

\[u_{n+1}\le(1-\alpha h)u_n+19h^3\beta^2d,\qquad \|e_n\|\le6h\sqrt{(\beta/\alpha)\beta d}.\]

\textbf{Proof.} The exact coboundary decomposition gives \(e_{n+1}+g_{n+1}=Z_n+M_n\). Conditional centering makes their cross term zero. Weighted curvature interpolation bounds \(\|Z_n\|^2\) by \((1-\alpha h)u_n+h\beta\|g_n\|^2\). Adding the martingale bound gives the coefficient \(1+18=19\). Equal initial states bound \(u_0\) by the corrector budget. The scalar affine recurrence gives a uniform energy floor. Removing the corrector by triangle inequality and the proved numerical comparison gives the final constant \(6\).

\begin{leanbox}{Lean: \texttt{shifted\_\allowbreak{}equation}, \texttt{shifted\_\allowbreak{}energy\_\allowbreak{}step}, \texttt{shifted\_\allowbreak{}energy\_\allowbreak{}bound} and 1 more}
\leansource{\texttt{shifted\_\allowbreak{}equation} \quad WassersteinMixing/\allowbreak{}PoissonComparison.lean, lines 125–131}
\begin{leancode}
\leanline{0}{theorem\ shifted\_equation\ (n\ :\ ℕ)\ :}
\leanline{4}{C.error\ (n\ +\ 1)\ +\ S.correctorAt\ score\ (n\ +\ 1)\ =}
\leanline{6}{C.predictable\ n\ +\ S.martingaleAt\ score\ C.defect\ n}
\end{leancode}
\leansource{\texttt{shifted\_\allowbreak{}energy\_\allowbreak{}step} \quad WassersteinMixing/\allowbreak{}PoissonComparison.lean, lines 134–148}
\begin{leancode}
\leanline{0}{theorem\ shifted\_energy\_step\ (n\ :\ ℕ)\ :}
\leanline{4}{‖C.error\ (n\ +\ 1)\ +\ S.correctorAt\ score\ (n\ +\ 1)‖\ \textasciicircum{}\ 2\ ≤}
\leanline{6}{P.contraction\ *\ ‖C.error\ n\ +\ S.correctorAt\ score\ n‖\ \textasciicircum{}\ 2\ +\ P.energyForcing}
\end{leancode}
\leansource{\texttt{shifted\_\allowbreak{}energy\_\allowbreak{}bound} \quad WassersteinMixing/\allowbreak{}PoissonComparison.lean, lines 150–154}
\begin{leancode}
\leanline{0}{theorem\ shifted\_energy\_bound\ (n\ :\ ℕ)\ :}
\leanline{4}{‖C.error\ n\ +\ S.correctorAt\ score\ n‖\ \textasciicircum{}\ 2\ ≤\ P.energyFloor}
\end{leancode}
\leansource{\texttt{error\_\allowbreak{}norm\_\allowbreak{}le\_\allowbreak{}bias} \quad WassersteinMixing/\allowbreak{}PoissonComparison.lean, lines 157–159}
\begin{leancode}
\leanline{0}{theorem\ error\_norm\_le\_bias\ (n\ :\ ℕ)\ :\ ‖C.error\ n‖\ ≤\ P.bias}
\end{leancode}
\leanlabel{How the Lean proof follows the argument}
\texttt{shifted\_\allowbreak{}equation} applies the exact algebraic shifted identity. \texttt{shifted\_\allowbreak{}energy\_\allowbreak{}step} combines weighted interpolation and actual martingale orthogonality. The recurrence and final norm theorems use the independently proved scalar recurrence results.
\end{leanbox}
\subsection{\texorpdfstring{Actual pushforward transport from the synchronous coupling}{Actual pushforward transport from the synchronous coupling}}
\label{report-step-WassersteinMixing-PoissonComparison-lean-4}


\textbf{Claim.} The reference initial law and the actual chain law satisfy \(W_2(\operatorname{Law}(Y_0),\operatorname{Law}(X_n))\le6h\sqrt{(\beta/\alpha)\beta d}\), and this transport distance is finite.



\textbf{Proof.} Stationarity replaces the initial reference law by the current reference law. The actual joint random variables \((Y_n,X_n)\) are an admissible coupling. Their \(L^2\) difference bounds its transport cost by the preceding error estimate. The same actual coupling proves finiteness before converting the extended-valued distance to a real number.

\begin{leanbox}{Lean: \texttt{stationary\_\allowbreak{}transport\_\allowbreak{}bound}, \texttt{stationary\_\allowbreak{}transport\_\allowbreak{}finite}}
\leansource{\texttt{stationary\_\allowbreak{}transport\_\allowbreak{}bound} \quad WassersteinMixing/\allowbreak{}PoissonComparison.lean, lines 163–167}
\begin{leancode}
\leanline{0}{theorem\ stationary\_transport\_bound\ [BorelSpace\ E]\ [SecondCountableTopology\ E]\ (n\ :\ ℕ)\ :}
\leanline{4}{ULA.wasserstein2Real\ (Q.map\ (C.reference\ 0))\ (Q.map\ (C.chain\ n))\ ≤\ P.bias}
\end{leancode}
\leansource{\texttt{stationary\_\allowbreak{}transport\_\allowbreak{}finite} \quad WassersteinMixing/\allowbreak{}PoissonComparison.lean, lines 170–174}
\begin{leancode}
\leanline{0}{theorem\ stationary\_transport\_finite\ [BorelSpace\ E]\ [SecondCountableTopology\ E]\ (n\ :\ ℕ)\ :}
\leanline{4}{ULA.wasserstein2\ (Q.map\ (C.reference\ 0))\ (Q.map\ (C.chain\ n))\ ≠\ ⊤}
\end{leancode}
\leanlabel{How the Lean proof follows the argument}
\texttt{stationary\_\allowbreak{}transport\_\allowbreak{}bound} uses the actual pushforward coupling cost theorem. \texttt{stationary\_\allowbreak{}transport\_\allowbreak{}finite} obtains a finite extended-valued bound from the same \(L^2\) coupling, preventing loss of information during real conversion.
\end{leanbox}


Earlier results used here: \cite[Section~3.17.12]{LangevinLean2026ULASC}; \cite[Section~3.17.13]{LangevinLean2026ULASC}; \cite[Section~3.17.3]{LangevinLean2026ULASC}; \cite[Section~3.17.4]{LangevinLean2026ULASC}.
\section{\texorpdfstring{Actual Gaussian ULA contraction and transport to an invariant law}{Actual Gaussian ULA contraction and transport to an invariant law}}
\label{proof-module-WassersteinMixing-Transport-lean}


Let \(V\) satisfy the stated differentiability, \(\alpha\)-strong convexity and \(\beta\)-Lipschitz gradient conditions, with \(0<h\le1/\beta\). The actual ULA update is \(x\mapsto x-h\nabla V(x)+\sqrt{2h}\,\xi\), with independent Gaussian innovations. This section proves its transport contraction from the potential conditions. The separate Gaussian-past fixed-point construction proves existence, finite second moment and exact invariance of the actual invariant law, discharging the invariant-law hypothesis of the general contraction result.



Earlier results used here: \cite[Section~3.9]{LangevinLean2026Kinetic}; \cite[Section~3.36]{LangevinLean2026ULASC}; \cite[Section~3.33]{LangevinLean2026ULASC}.
\subsection{\texorpdfstring{Common noise gives the exact pathwise factor}{Common noise gives the exact pathwise factor}}
\label{report-step-WassersteinMixing-Transport-lean-1}


\textbf{Claim.} For all \(x,y,z\), the two updates using the same innovation satisfy

\[\|(x-h\nabla V(x)+\sqrt{2h}z)-(y-h\nabla V(y)+\sqrt{2h}z)\|
\le(1-\alpha h)\|x-y\|.\]


\textbf{Proof.} The common noise cancels. The strong-convexity and smoothness gradient-step estimate applies for \(h\le2/(\alpha+\beta)\). This restriction follows from \(h\beta\le1\) and \(\alpha\le\beta\), since \(h(\alpha+\beta)\le2h\beta\le2\). Its exact contraction factor is \(1-\alpha h\ge0\), giving the claim.

\begin{leanbox}{Lean: \texttt{ula\_\allowbreak{}shared\_\allowbreak{}noise\_\allowbreak{}contraction}}
\leansource{\texttt{ula\_\allowbreak{}shared\_\allowbreak{}noise\_\allowbreak{}contraction} \quad WassersteinMixing/\allowbreak{}Transport.lean, lines 27–39}
\begin{leancode}
\leanline{0}{theorem\ ula\_shared\_noise\_contraction\ (P\ :\ Parameters)\ (V\ :\ H\ →\ ℝ)}
\leanline{4}{(hV\ :\ PotentialAssumptions\ P\ V)\ (x\ y\ z\ :\ H)\ :}
\leanline{4}{‖ULA.malaProposal\ V\ P.step\ x\ z\ -\ ULA.malaProposal\ V\ P.step\ y\ z‖\ ≤}
\leanline{6}{P.contraction\ *\ ‖x\ -\ y‖}
\end{leancode}
\leanlabel{How the Lean proof follows the argument}
\texttt{ula\_\allowbreak{}shared\_\allowbreak{}noise\_\allowbreak{}contraction} verifies the admissible shared ULA step range and invokes \texttt{ULA.\allowbreak{}strong\_\allowbreak{}smooth\_\allowbreak{}gradient\_\allowbreak{}step\_\allowbreak{}contract}. Unfolding \texttt{ULA.\allowbreak{}mala\allowbreak{}Proposal} cancels the identical noise terms. The contraction is derived from the potential contract rather than assumed as a coupling field.
\end{leanbox}


Earlier results used here: \cite[Section~3.10.7]{LangevinLean2026ULASC}; \cite[Section~3.34.1]{LangevinLean2026ULASC}.
\subsection{\texorpdfstring{Transport contraction of the actual Markov laws}{Transport contraction of the actual Markov laws}}
\label{report-step-WassersteinMixing-Transport-lean-2}


\textbf{Claim.} Let \(T\) be the independent-innovation Markov transition and let \(\mu,\nu\) be probability laws. Then, in extended nonnegative transport distance,

\[W_2(T\mu,T\nu)\le qW_2(\mu,\nu),\qquad
W_2(T^k\mu,T^k\nu)\le q^kW_2(\mu,\nu),\quad q=1-\alpha h.\]


\textbf{Proof.} Extend any initial coupling by an independent innovation and use that same innovation for both updates. The proved pathwise bound multiplies the coupling's root-mean-square cost by at most \(q\). Passing to the infimum over couplings proves the one-step estimate. Induction gives the iterated estimate because \(q\ge0\). For standard Gaussian innovations, the independent-product transition is exactly the existing Gaussian ULA law transition; induction identifies their iterated laws.

\begin{leanbox}{Lean: \texttt{ula\_\allowbreak{}markovStep\_\allowbreak{}wasserstein2\_\allowbreak{}le}, \texttt{ula\_\allowbreak{}markovLaws\_\allowbreak{}wasserstein2\_\allowbreak{}le}, \texttt{ula\_\allowbreak{}iterate\_\allowbreak{}eq\_\allowbreak{}markovLaws} and 1 more}
\leansource{\texttt{ula\_\allowbreak{}markovStep\_\allowbreak{}wasserstein2\_\allowbreak{}le} \quad WassersteinMixing/\allowbreak{}Transport.lean, lines 43–57}
\begin{leancode}
\leanline{0}{theorem\ ula\_markovStep\_wasserstein2\_le\ (P\ :\ Parameters)\ (V\ :\ H\ →\ ℝ)}
\leanline{4}{(hV\ :\ PotentialAssumptions\ P\ V)\ (η\ :\ Measure\ H)\ [IsProbabilityMeasure\ η]}
\leanline{4}{(μ\ ν\ :\ Measure\ H)\ [IsProbabilityMeasure\ μ]\ [IsProbabilityMeasure\ ν]\ :}
\leanline{4}{ULA.wasserstein2\ (ULA.markovStep\ (ULA.malaProposal\ V\ P.step)\ η\ μ)}
\leanline{6}{(ULA.markovStep\ (ULA.malaProposal\ V\ P.step)\ η\ ν)\ ≤}
\leanline{6}{ENNReal.ofReal\ P.contraction\ *\ ULA.wasserstein2\ μ\ ν}
\end{leancode}
\leansource{\texttt{ula\_\allowbreak{}markovLaws\_\allowbreak{}wasserstein2\_\allowbreak{}le} \quad WassersteinMixing/\allowbreak{}Transport.lean, lines 60–82}
\begin{leancode}
\leanline{0}{theorem\ ula\_markovLaws\_wasserstein2\_le\ (P\ :\ Parameters)\ (V\ :\ H\ →\ ℝ)}
\leanline{4}{(hV\ :\ PotentialAssumptions\ P\ V)\ (η\ :\ Measure\ H)\ [IsProbabilityMeasure\ η]}
\leanline{4}{(μ\ ν\ :\ Measure\ H)\ [IsProbabilityMeasure\ μ]\ [IsProbabilityMeasure\ ν]\ (k\ :\ ℕ)\ :}
\leanline{4}{ULA.wasserstein2\ (ULA.markovLaws\ (ULA.malaProposal\ V\ P.step)\ η\ μ\ k)}
\leanline{6}{(ULA.markovLaws\ (ULA.malaProposal\ V\ P.step)\ η\ ν\ k)\ ≤}
\leanline{6}{(ENNReal.ofReal\ P.contraction)\ \textasciicircum{}\ k\ *\ ULA.wasserstein2\ μ\ ν}
\end{leancode}
\leansource{\texttt{ula\_\allowbreak{}iterate\_\allowbreak{}eq\_\allowbreak{}markovLaws} \quad WassersteinMixing/\allowbreak{}Transport.lean, lines 87–94}
\begin{leancode}
\leanline{0}{theorem\ ula\_iterate\_eq\_markovLaws\ (V\ :\ H\ →\ ℝ)\ (h\ :\ ℝ)\ (μ\ :\ Measure\ H)\ (k\ :\ ℕ)\ :}
\leanline{4}{(ULA.ulaLawStep\ V\ h)\textasciicircum{}[k]\ μ\ =}
\leanline{6}{ULA.markovLaws\ (ULA.malaProposal\ V\ h)\ (stdGaussian\ H)\ μ\ k}
\end{leancode}
\leansource{\texttt{ula\_\allowbreak{}iterate\_\allowbreak{}wasserstein2\_\allowbreak{}le} \quad WassersteinMixing/\allowbreak{}Transport.lean, lines 98–105}
\begin{leancode}
\leanline{0}{theorem\ ula\_iterate\_wasserstein2\_le\ (P\ :\ Parameters)\ (V\ :\ H\ →\ ℝ)}
\leanline{4}{(hV\ :\ PotentialAssumptions\ P\ V)\ (μ\ ν\ :\ Measure\ H)}
\leanline{4}{[IsProbabilityMeasure\ μ]\ [IsProbabilityMeasure\ ν]\ (k\ :\ ℕ)\ :}
\leanline{4}{ULA.wasserstein2\ ((ULA.ulaLawStep\ V\ P.step)\textasciicircum{}[k]\ μ)}
\leanline{6}{((ULA.ulaLawStep\ V\ P.step)\textasciicircum{}[k]\ ν)\ ≤}
\leanline{6}{(ENNReal.ofReal\ P.contraction)\ \textasciicircum{}\ k\ *\ ULA.wasserstein2\ μ\ ν}
\end{leancode}
\leanlabel{How the Lean proof follows the argument}
\texttt{ula\_\allowbreak{}markov\allowbreak{}Step\_\allowbreak{}wasserstein2\_\allowbreak{}le} applies the shared random-map transport theorem to the actual measurable update. \texttt{ula\_\allowbreak{}markov\allowbreak{}Laws\_\allowbreak{}wasserstein2\_\allowbreak{}le} iterates it with probability preservation checked. \texttt{ula\_\allowbreak{}iterate\_\allowbreak{}eq\_\allowbreak{}markov\allowbreak{}Laws} proves equality with the Gaussian ULA implementation, and \texttt{ula\_\allowbreak{}iterate\_\allowbreak{}wasserstein2\_\allowbreak{}le} substitutes that identity into the extended-distance estimate.
\end{leanbox}


Earlier results used here: \cite[Section~3.17.3]{LangevinLean2026ULASC}; \cite[Section~3.33.11]{LangevinLean2026ULASC}; \cite[Section~3.33.5]{LangevinLean2026ULASC}; \cite[Section~3.33.6]{LangevinLean2026ULASC}; \cite[Section~3.34.1]{LangevinLean2026ULASC}; \cite[Section~3.36.5]{LangevinLean2026ULASC}; \cite[Section~3.36.7]{LangevinLean2026ULASC}; \cite[Section~3.9.2]{LangevinLean2026Kinetic}.
\subsection{\texorpdfstring{Finite transport and convergence to the constructed invariant law}{Finite transport and convergence to the constructed invariant law}}
\label{report-step-WassersteinMixing-Transport-lean-3}


\textbf{Claim.} If \(W_2(\mu,\nu)<\infty\), then the iterated transport distance is finite and its ordinary real value obeys the same geometric estimate. If additionally \(T\nu=\nu\), then

\[W_2(T^k\mu,\nu)\le q^kW_2(\mu,\nu).\]


\textbf{Proof.} The extended upper bound is a finite product, so the distance it bounds is finite. Converting a finite extended inequality to real values preserves it. The stationarity equality implies \(T^k\nu=\nu\) by induction; substitute this into the two-law estimate. The independent Gaussian-past fixed-point argument constructs this invariant probability law and proves its finite second moment and exact stationarity, so the final bounds apply to that actual law.

\begin{leanbox}{Lean: \texttt{ula\_\allowbreak{}iterate\_\allowbreak{}wasserstein2Real\_\allowbreak{}le}, \texttt{ula\_\allowbreak{}iterate\_\allowbreak{}transport\_\allowbreak{}finite}, \texttt{ula\_\allowbreak{}to\_\allowbreak{}invariant\_\allowbreak{}wasserstein2Real\_\allowbreak{}le}}
\leansource{\texttt{ula\_\allowbreak{}iterate\_\allowbreak{}wasserstein2Real\_\allowbreak{}le} \quad WassersteinMixing/\allowbreak{}Transport.lean, lines 109–121}
\begin{leancode}
\leanline{0}{theorem\ ula\_iterate\_wasserstein2Real\_le\ (P\ :\ Parameters)\ (V\ :\ H\ →\ ℝ)}
\leanline{4}{(hV\ :\ PotentialAssumptions\ P\ V)\ (μ\ ν\ :\ Measure\ H)}
\leanline{4}{[IsProbabilityMeasure\ μ]\ [IsProbabilityMeasure\ ν]}
\leanline{4}{(hfinite\ :\ ULA.wasserstein2\ μ\ ν\ ≠\ ⊤)\ (k\ :\ ℕ)\ :}
\leanline{4}{ULA.wasserstein2Real\ ((ULA.ulaLawStep\ V\ P.step)\textasciicircum{}[k]\ μ)}
\leanline{6}{((ULA.ulaLawStep\ V\ P.step)\textasciicircum{}[k]\ ν)\ ≤}
\leanline{6}{P.contraction\ \textasciicircum{}\ k\ *\ ULA.wasserstein2Real\ μ\ ν}
\end{leancode}
\leansource{\texttt{ula\_\allowbreak{}iterate\_\allowbreak{}transport\_\allowbreak{}finite} \quad WassersteinMixing/\allowbreak{}Transport.lean, lines 124–132}
\begin{leancode}
\leanline{0}{theorem\ ula\_iterate\_transport\_finite\ (P\ :\ Parameters)\ (V\ :\ H\ →\ ℝ)}
\leanline{4}{(hV\ :\ PotentialAssumptions\ P\ V)\ (μ\ ν\ :\ Measure\ H)}
\leanline{4}{[IsProbabilityMeasure\ μ]\ [IsProbabilityMeasure\ ν]}
\leanline{4}{(hfinite\ :\ ULA.wasserstein2\ μ\ ν\ ≠\ ⊤)\ (k\ :\ ℕ)\ :}
\leanline{4}{ULA.wasserstein2\ ((ULA.ulaLawStep\ V\ P.step)\textasciicircum{}[k]\ μ)}
\leanline{6}{((ULA.ulaLawStep\ V\ P.step)\textasciicircum{}[k]\ ν)\ ≠\ ⊤}
\end{leancode}
\leansource{\texttt{ula\_\allowbreak{}to\_\allowbreak{}invariant\_\allowbreak{}wasserstein2Real\_\allowbreak{}le} \quad WassersteinMixing/\allowbreak{}Transport.lean, lines 136–145}
\begin{leancode}
\leanline{0}{theorem\ ula\_to\_invariant\_wasserstein2Real\_le\ (P\ :\ Parameters)\ (V\ :\ H\ →\ ℝ)}
\leanline{4}{(hV\ :\ PotentialAssumptions\ P\ V)\ (μ\ ν\ :\ Measure\ H)}
\leanline{4}{[IsProbabilityMeasure\ μ]\ [IsProbabilityMeasure\ ν]}
\leanline{4}{(hν\ :\ ULA.ulaLawStep\ V\ P.step\ ν\ =\ ν)}
\leanline{4}{(hfinite\ :\ ULA.wasserstein2\ μ\ ν\ ≠\ ⊤)\ (k\ :\ ℕ)\ :}
\leanline{4}{ULA.wasserstein2Real\ ((ULA.ulaLawStep\ V\ P.step)\textasciicircum{}[k]\ μ)\ ν\ ≤}
\leanline{6}{P.contraction\ \textasciicircum{}\ k\ *\ ULA.wasserstein2Real\ μ\ ν}
\end{leancode}
\leanlabel{How the Lean proof follows the argument}
\texttt{ula\_\allowbreak{}iterate\_\allowbreak{}transport\_\allowbreak{}finite} bounds the output by a finite extended product. \texttt{ula\_\allowbreak{}iterate\_\allowbreak{}wasserstein2\allowbreak{}Real\_\allowbreak{}le} applies \texttt{ENNReal.\allowbreak{}to\allowbreak{}Real\_\allowbreak{}mono} only after finiteness is proved. \texttt{ula\_\allowbreak{}to\_\allowbreak{}invariant\_\allowbreak{}wasserstein2\allowbreak{}Real\_\allowbreak{}le} uses the supplied invariance equality through \texttt{Function.\allowbreak{}iterate\_\allowbreak{}fixed}.



The separate invariant-law construction proves the invariance equality and finite second moment required here, and the final comparison theorems apply this generic result to that constructed law.
\end{leanbox}


Earlier results used here: \cite[Section~3.17.3]{LangevinLean2026ULASC}; \cite[Section~3.17.4]{LangevinLean2026ULASC}; \cite[Section~3.36.7]{LangevinLean2026ULASC}.
\section{\texorpdfstring{Constructing the actual ULA invariant probability law on an independent Gaussian past}{Constructing the actual ULA invariant probability law on an independent Gaussian past}}
\label{proof-module-WassersteinMixing-Invariant-lean}


A contracting additive-noise update has an invariant probability law constructed from a fixed random vector on an independent infinite noise sequence. The operator acts on the complete space \(L^2\); a Banach fixed point supplies existence and a finite second moment. Independence of the most recent noise from the older past proves invariance for the original transition. Applied to the actual Gaussian ULA update, this gives an invariant law and convergence from every finite-second-moment initialization without an invariant-law premise.



Earlier results used here: \cite[Section~3.13]{LangevinLean2026ULASC}; \cite[Section~3.36]{LangevinLean2026ULASC}; \cite[Section~3.21]{LangevinLean2026ULASC}.
\subsection{\texorpdfstring{The independent past, its tail and its newest innovation}{The independent past, its tail and its newest innovation}}
\label{report-step-WassersteinMixing-Invariant-lean-1}


\textbf{Construction.} For an innovation probability law \(\eta\), put \(Q=\eta^{\otimes\mathbb N}\) on sequences \(\omega=(\omega_0,\omega_1,\ldots)\). The tail map is \(\tau\omega=(\omega_1,\omega_2,\ldots)\).




\textbf{Claim.} The map \(\tau\) preserves \(Q\), and the tail is independent of \(\omega_0\).




\textbf{Proof.} Every finite set of coordinates has its product law. Reindexing a finite cylinder by the injective successor map therefore leaves its probability unchanged; the product-measure identity gives measure preservation. Independence follows because every finite set of tail coordinates has indices disjoint from \(\{0\}\). Independence of the full tail process follows from these finite-coordinate independence statements.

\begin{leanbox}{Lean: \texttt{probability}, \texttt{probability\_\allowbreak{}isProbability}, \texttt{tail\_\allowbreak{}measurePreserving} and 1 more}
\leansource{\texttt{probability} \quad WassersteinMixing/\allowbreak{}Invariant.lean, lines 31–31}
\begin{leancode}
\leanline{0}{def\ probability\ (η\ :\ Measure\ E)\ :\ Measure\ (ℕ\ →\ E)\ :=\ Measure.infinitePi\ (fun\ \_\ :\ ℕ\ ↦\ η)}
\end{leancode}
\leansource{\texttt{probability\_\allowbreak{}isProbability} \quad WassersteinMixing/\allowbreak{}Invariant.lean, lines 33–34}
\begin{leancode}
\leanline{0}{instance\ probability\_isProbability\ (η\ :\ Measure\ E)\ [IsProbabilityMeasure\ η]\ :}
\leanline{4}{IsProbabilityMeasure\ (probability\ η)\ :=\ by\ unfold\ probability;\ infer\_instance}
\end{leancode}
\leansource{\texttt{tail\_\allowbreak{}measurePreserving} \quad WassersteinMixing/\allowbreak{}Invariant.lean, lines 37–40}
\begin{leancode}
\leanline{0}{theorem\ tail\_measurePreserving\ (η\ :\ Measure\ E)\ [IsProbabilityMeasure\ η]\ :}
\leanline{4}{MeasurePreserving\ (fun\ ω\ :\ ℕ\ →\ E\ ↦\ fun\ n\ ↦\ ω\ (n\ +\ 1))\ (probability\ η)\ (probability\ η)}
\end{leancode}
\leansource{\texttt{tail\_\allowbreak{}independent\_\allowbreak{}head} \quad WassersteinMixing/\allowbreak{}Invariant.lean, lines 43–58}
\begin{leancode}
\leanline{0}{theorem\ tail\_independent\_head\ (η\ :\ Measure\ E)\ [IsProbabilityMeasure\ η]\ :}
\leanline{4}{IndepFun\ (fun\ ω\ :\ ℕ\ →\ E\ ↦\ fun\ n\ ↦\ ω\ (n\ +\ 1))\ (fun\ ω\ ↦\ ω\ 0)\ (probability\ η)}
\end{leancode}
\leanlabel{How the Lean proof follows the argument}
\texttt{probability} is the actual infinite product measure and has a probability instance. \texttt{tail\_\allowbreak{}measure\allowbreak{}Preserving} applies the injective reindexing identity for \texttt{Measure.\allowbreak{}infinite\allowbreak{}Pi}. \texttt{tail\_\allowbreak{}independent\_\allowbreak{}head} starts from the independent coordinate family, proves the finite index sets disjoint and extends finite-coordinate independence to the whole process.
\end{leanbox}
\subsection{\texorpdfstring{A contraction on the complete space of square-integrable past functions}{A contraction on the complete space of square-integrable past functions}}
\label{report-step-WassersteinMixing-Invariant-lean-2}


\textbf{Construction and claim.} Suppose \(F:E\to E\) is \(q\)-Lipschitz with \(0\le q<1\), the innovation has finite second moment, and \(a\in\mathbb R\). On \(L^2(Q;E)\) define

\[\Phi(X)=F(X\circ\tau)+a\omega_0.\]


The tail pullback is a linear isometry because \(\tau\) preserves \(Q\); the head coordinate is square-integrable because its law is \(\eta\). A Lipschitz map has at most linear growth, so composition with \(F\) preserves square integrability on the probability space. The common additive noise cancels in a difference, giving

\[\|\Phi(X)-\Phi(Y)\|_{L^2}\le q\|X\circ\tau-Y\circ\tau\|_{L^2}
=q\|X-Y\|_{L^2}.\]


Thus \(\Phi\) is a strict contraction on a complete space.

\begin{leanbox}{Lean: \texttt{tailPullback}, \texttt{headLp}, \texttt{pastOperator} and 1 more}
\leansource{\texttt{tailPullback} \quad WassersteinMixing/\allowbreak{}Invariant.lean, lines 64–66}
\begin{leancode}
\leanline{0}{def\ tailPullback\ (η\ :\ Measure\ E)\ [IsProbabilityMeasure\ η]\ :}
\leanline{4}{Lp\ E\ 2\ (probability\ η)\ →ₗᵢ[ℝ]\ Lp\ E\ 2\ (probability\ η)\ :=}
\leanline{2}{Lp.compMeasurePreservingₗᵢ\ ℝ\ \_\ (tail\_measurePreserving\ η)}
\end{leancode}
\leansource{\texttt{headLp} \quad WassersteinMixing/\allowbreak{}Invariant.lean, lines 69–72}
\begin{leancode}
\leanline{0}{def\ headLp\ (η\ :\ Measure\ E)\ [IsProbabilityMeasure\ η]}
\leanline{4}{(hη\ :\ MemLp\ (id\ :\ E\ →\ E)\ 2\ η)\ :\ Lp\ E\ 2\ (probability\ η)\ :=}
\leanline{2}{(hη.comp\_measurePreserving\ (measurePreserving\_eval\_infinitePi\ (fun\ \_\ :\ ℕ\ ↦\ η)\ 0)).toLp}
\leanline{4}{(fun\ ω\ :\ ℕ\ →\ E\ ↦\ ω\ 0)}
\end{leancode}
\leansource{\texttt{pastOperator} \quad WassersteinMixing/\allowbreak{}Invariant.lean, lines 75–79}
\begin{leancode}
\leanline{0}{def\ pastOperator\ (η\ :\ Measure\ E)\ [IsProbabilityMeasure\ η]}
\leanline{4}{(hη\ :\ MemLp\ (id\ :\ E\ →\ E)\ 2\ η)\ (F\ :\ E\ →\ E)\ (q\ :\ ℝ≥0)}
\leanline{4}{(hF\ :\ LipschitzWith\ q\ F)\ (a\ :\ ℝ)\ (X\ :\ Lp\ E\ 2\ (probability\ η))\ :}
\leanline{4}{Lp\ E\ 2\ (probability\ η)\ :=}
\leanline{2}{ULA.lpDrift\ (probability\ η)\ F\ hF\ (tailPullback\ η\ X)\ +\ a\ •\ headLp\ η\ hη}
\end{leancode}
\leansource{\texttt{pastOperator\_\allowbreak{}lipschitz} \quad WassersteinMixing/\allowbreak{}Invariant.lean, lines 83–92}
\begin{leancode}
\leanline{0}{theorem\ pastOperator\_lipschitz\ (η\ :\ Measure\ E)\ [IsProbabilityMeasure\ η]}
\leanline{4}{(hη\ :\ MemLp\ (id\ :\ E\ →\ E)\ 2\ η)\ (F\ :\ E\ →\ E)\ (q\ :\ ℝ≥0)}
\leanline{4}{(hF\ :\ LipschitzWith\ q\ F)\ (a\ :\ ℝ)\ :}
\leanline{4}{LipschitzWith\ q\ (pastOperator\ η\ hη\ F\ q\ hF\ a)}
\end{leancode}
\leanlabel{How the Lean proof follows the argument}
\texttt{tail\allowbreak{}Pullback} is the actual measure-preserving pullback isometry and \texttt{head\allowbreak{}Lp} is the newest coordinate as an \texttt{Lp} vector. \texttt{past\allowbreak{}Operator} uses the shared square-integrable drift composition. \texttt{past\allowbreak{}Operator\_\allowbreak{}lipschitz} cancels the common noise, applies the drift Lipschitz bound and uses the tail isometry.
\end{leanbox}


Earlier results used here: \cite[Section~3.13.2]{LangevinLean2026ULASC}; \cite[Section~3.13.5]{LangevinLean2026ULASC}.
\subsection{\texorpdfstring{The fixed random variable and its invariant law}{The fixed random variable and its invariant law}}
\label{report-step-WassersteinMixing-Invariant-lean-3}


\textbf{Claim.} There is a square-integrable random vector \(X_*\) with

\[X_*(\omega)=F(X_*(\tau\omega))+a\omega_0\qquad Q\text{-almost everywhere}.\]


Its law \(\nu\) is a probability law with finite second moment and is invariant under the independent-noise transition \(x\mapsto F(x)+a\xi\), \(\xi\sim\eta\).




\textbf{Proof.} Banach's fixed-point theorem applied to \(\Phi\) constructs \(X_*\in L^2\). Equality in \(L^2\) yields the displayed almost-everywhere recursion after choosing representatives. The pushforward law is a probability measure, and square integrability transfers to it. Since the tail preserves \(Q\), \(X_*\circ\tau\) has law \(\nu\). Since the tail is independent of \(\omega_0\), the pair \((X_*\circ\tau,\omega_0)\) has product law \(\nu\otimes\eta\). Taking laws in the recursion therefore gives \(T\nu=\nu\) for the original independent-noise transition.

\begin{leanbox}{Lean: \texttt{state}, \texttt{state\_\allowbreak{}recursion}, \texttt{law} and 3 more}
\leansource{\texttt{state} \quad WassersteinMixing/\allowbreak{}Invariant.lean, lines 95–99}
\begin{leancode}
\leanline{0}{def\ state\ (η\ :\ Measure\ E)\ [IsProbabilityMeasure\ η]}
\leanline{4}{(hη\ :\ MemLp\ (id\ :\ E\ →\ E)\ 2\ η)\ (F\ :\ E\ →\ E)\ (q\ :\ ℝ≥0)}
\leanline{4}{(hF\ :\ LipschitzWith\ q\ F)\ (hq\ :\ q\ <\ 1)\ (a\ :\ ℝ)\ :\ Lp\ E\ 2\ (probability\ η)\ :=}
\leanline{2}{ContractingWith.fixedPoint\ (pastOperator\ η\ hη\ F\ q\ hF\ a)}
\leanline{4}{⟨hq,\ pastOperator\_lipschitz\ η\ hη\ F\ q\ hF\ a⟩}
\end{leancode}
\leansource{\texttt{state\_\allowbreak{}recursion} \quad WassersteinMixing/\allowbreak{}Invariant.lean, lines 103–134}
\begin{leancode}
\leanline{0}{theorem\ state\_recursion\ (η\ :\ Measure\ E)\ [IsProbabilityMeasure\ η]}
\leanline{4}{(hη\ :\ MemLp\ (id\ :\ E\ →\ E)\ 2\ η)\ (F\ :\ E\ →\ E)\ (q\ :\ ℝ≥0)}
\leanline{4}{(hF\ :\ LipschitzWith\ q\ F)\ (hq\ :\ q\ <\ 1)\ (a\ :\ ℝ)\ :}
\leanline{4}{state\ η\ hη\ F\ q\ hF\ hq\ a\ =ᵐ[probability\ η]\ fun\ ω\ ↦}
\leanline{6}{F\ (state\ η\ hη\ F\ q\ hF\ hq\ a\ (fun\ n\ ↦\ ω\ (n\ +\ 1)))\ +\ a\ •\ ω\ 0}
\end{leancode}
\leansource{\texttt{law} \quad WassersteinMixing/\allowbreak{}Invariant.lean, lines 137–140}
\begin{leancode}
\leanline{0}{def\ law\ (η\ :\ Measure\ E)\ [IsProbabilityMeasure\ η]}
\leanline{4}{(hη\ :\ MemLp\ (id\ :\ E\ →\ E)\ 2\ η)\ (F\ :\ E\ →\ E)\ (q\ :\ ℝ≥0)}
\leanline{4}{(hF\ :\ LipschitzWith\ q\ F)\ (hq\ :\ q\ <\ 1)\ (a\ :\ ℝ)\ :\ Measure\ E\ :=}
\leanline{2}{(probability\ η).map\ (state\ η\ hη\ F\ q\ hF\ hq\ a)}
\end{leancode}
\leansource{\texttt{law\_\allowbreak{}isProbability} \quad WassersteinMixing/\allowbreak{}Invariant.lean, lines 142–149}
\begin{leancode}
\leanline{0}{instance\ law\_isProbability\ (η\ :\ Measure\ E)\ [IsProbabilityMeasure\ η]}
\leanline{4}{(hη\ :\ MemLp\ (id\ :\ E\ →\ E)\ 2\ η)\ (F\ :\ E\ →\ E)\ (q\ :\ ℝ≥0)}
\leanline{4}{(hF\ :\ LipschitzWith\ q\ F)\ (hq\ :\ q\ <\ 1)\ (a\ :\ ℝ)\ :}
\leanline{4}{IsProbabilityMeasure\ (law\ η\ hη\ F\ q\ hF\ hq\ a)\ :=\ by}
\leanline{2}{unfold\ law}
\leanline{2}{constructor}
\leanline{2}{rw\ [Measure.map\_apply\_of\_aemeasurable\ (Lp.aestronglyMeasurable\ \_).aemeasurable}
\leanline{4}{MeasurableSet.univ,\ Set.preimage\_univ,\ measure\_univ]}
\end{leancode}
\leansource{\texttt{law\_\allowbreak{}memLp} \quad WassersteinMixing/\allowbreak{}Invariant.lean, lines 152–157}
\begin{leancode}
\leanline{0}{theorem\ law\_memLp\ (η\ :\ Measure\ E)\ [IsProbabilityMeasure\ η]}
\leanline{4}{(hη\ :\ MemLp\ (id\ :\ E\ →\ E)\ 2\ η)\ (F\ :\ E\ →\ E)\ (q\ :\ ℝ≥0)}
\leanline{4}{(hF\ :\ LipschitzWith\ q\ F)\ (hq\ :\ q\ <\ 1)\ (a\ :\ ℝ)\ :}
\leanline{4}{MemLp\ (id\ :\ E\ →\ E)\ 2\ (law\ η\ hη\ F\ q\ hF\ hq\ a)}
\end{leancode}
\leansource{\texttt{law\_\allowbreak{}invariant} \quad WassersteinMixing/\allowbreak{}Invariant.lean, lines 161–188}
\begin{leancode}
\leanline{0}{theorem\ law\_invariant\ (η\ :\ Measure\ E)\ [IsProbabilityMeasure\ η]}
\leanline{4}{(hη\ :\ MemLp\ (id\ :\ E\ →\ E)\ 2\ η)\ (F\ :\ E\ →\ E)\ (q\ :\ ℝ≥0)}
\leanline{4}{(hF\ :\ LipschitzWith\ q\ F)\ (hq\ :\ q\ <\ 1)\ (a\ :\ ℝ)\ :}
\leanline{4}{((law\ η\ hη\ F\ q\ hF\ hq\ a).prod\ η).map\ (fun\ p\ ↦\ F\ p.1\ +\ a\ •\ p.2)\ =}
\leanline{6}{law\ η\ hη\ F\ q\ hF\ hq\ a}
\end{leancode}
\leanlabel{How the Lean proof follows the argument}
\texttt{state} invokes \texttt{Contracting\allowbreak{}With.\allowbreak{}fixed\allowbreak{}Point} with the proved contraction. \texttt{state\_\allowbreak{}recursion} proves the literal almost-everywhere recursion using the pullback and drift representative identities. \texttt{law} is its actual pushforward, with probability and finite-second-moment proofs. \texttt{law\_\allowbreak{}invariant} identifies the head/\allowbreak{}tail pair law as a product, then pushes that product through the actual update.
\end{leanbox}


Earlier results used here: \cite[Section~3.13.2]{LangevinLean2026ULASC}; \cite[Section~3.13.3]{LangevinLean2026ULASC}.
\subsection{\texorpdfstring{Specializing the construction to the actual Gaussian ULA transition}{Specializing the construction to the actual Gaussian ULA transition}}
\label{report-step-WassersteinMixing-Invariant-lean-4}


\textbf{Claim.} For the potential contract and numerical parameters, take \(F(x)=x-h\nabla V(x)\), \(q=1-\alpha h\), \(\eta=\mathcal N(0,I)\) and \(a=\sqrt{2h}\). The resulting law \(\widehat\pi\) is a probability measure with finite second moment and satisfies the actual Gaussian ULA invariance equality \(T\widehat\pi=\widehat\pi\). For every probability initialization \(\mu\) with finite second moment,

\[W_2(T^k\mu,\widehat\pi)\le q^kW_2(\mu,\widehat\pi),
\qquad W_2(T^k\mu,\widehat\pi)\longrightarrow0.\]


\textbf{Proof.} Common-noise gradient-step contraction proves the required Lipschitz bound for \(F\), and \(0\le q<1\) is already established. Standard Gaussian innovations have finite second moment, so the general construction applies. Its invariance equality is exactly the existing Gaussian law transition. Both initial laws have finite second moments, so their transport distance is finite. The actual ULA contraction gives the geometric estimate. Finally \(q^k\to0\) and the distance is nonnegative, so squeezing proves attraction.

\begin{leanbox}{Lean: \texttt{ulaGradientStep}, \texttt{ulaGradientStep\_\allowbreak{}lipschitz}, \texttt{invariantLaw} and 5 more}
\leansource{\texttt{ulaGradientStep} \quad WassersteinMixing/\allowbreak{}Invariant.lean, lines 198–199}
\begin{leancode}
\leanline{0}{def\ ulaGradientStep\ (P\ :\ Parameters)\ (V\ :\ E\ →\ ℝ)\ (x\ :\ E)\ :\ E\ :=}
\leanline{2}{x\ -\ P.step\ •\ gradient\ V\ x}
\end{leancode}
\leansource{\texttt{ulaGradientStep\_\allowbreak{}lipschitz} \quad WassersteinMixing/\allowbreak{}Invariant.lean, lines 203–212}
\begin{leancode}
\leanline{0}{theorem\ ulaGradientStep\_lipschitz\ (P\ :\ Parameters)\ (V\ :\ E\ →\ ℝ)}
\leanline{4}{(hV\ :\ PotentialAssumptions\ P\ V)\ :}
\leanline{4}{LipschitzWith\ ⟨P.contraction,\ P.contraction\_nonneg⟩\ (ulaGradientStep\ P\ V)}
\end{leancode}
\leansource{\texttt{invariantLaw} \quad WassersteinMixing/\allowbreak{}Invariant.lean, lines 216–219}
\begin{leancode}
\leanline{0}{def\ invariantLaw\ (P\ :\ Parameters)\ (V\ :\ E\ →\ ℝ)\ (hV\ :\ PotentialAssumptions\ P\ V)\ :\ Measure\ E\ :=}
\leanline{2}{AdditiveInvariant.law\ (stdGaussian\ E)\ (IsGaussian.memLp\_id\ (stdGaussian\ E)\ 2\ (by\ norm\_num))}
\leanline{4}{(ulaGradientStep\ P\ V)\ ⟨P.contraction,\ P.contraction\_nonneg⟩}
\leanline{4}{(ulaGradientStep\_lipschitz\ P\ V\ hV)\ P.contraction\_lt\_one\ (Real.sqrt\ (2\ *\ P.step))}
\end{leancode}
\leansource{\texttt{invariantLaw\_\allowbreak{}isProbability} \quad WassersteinMixing/\allowbreak{}Invariant.lean, lines 221–227}
\begin{leancode}
\leanline{0}{instance\ invariantLaw\_isProbability\ (P\ :\ Parameters)\ (V\ :\ E\ →\ ℝ)}
\leanline{4}{(hV\ :\ PotentialAssumptions\ P\ V)\ :\ IsProbabilityMeasure\ (invariantLaw\ P\ V\ hV)\ :=\ by}
\leanline{2}{unfold\ invariantLaw}
\leanline{2}{exact\ AdditiveInvariant.law\_isProbability\ (stdGaussian\ E)}
\leanline{4}{(IsGaussian.memLp\_id\ (stdGaussian\ E)\ 2\ (by\ norm\_num))}
\leanline{4}{(ulaGradientStep\ P\ V)\ ⟨P.contraction,\ P.contraction\_nonneg⟩}
\leanline{4}{(ulaGradientStep\_lipschitz\ P\ V\ hV)\ P.contraction\_lt\_one\ (Real.sqrt\ (2\ *\ P.step))}
\end{leancode}
\leansource{\texttt{invariantLaw\_\allowbreak{}memLp} \quad WassersteinMixing/\allowbreak{}Invariant.lean, lines 230–232}
\begin{leancode}
\leanline{0}{theorem\ invariantLaw\_memLp\ (P\ :\ Parameters)\ (V\ :\ E\ →\ ℝ)}
\leanline{4}{(hV\ :\ PotentialAssumptions\ P\ V)\ :\ MemLp\ (id\ :\ E\ →\ E)\ 2\ (invariantLaw\ P\ V\ hV)}
\end{leancode}
\leansource{\texttt{invariantLaw\_\allowbreak{}invariant} \quad WassersteinMixing/\allowbreak{}Invariant.lean, lines 236–242}
\begin{leancode}
\leanline{0}{theorem\ invariantLaw\_invariant\ (P\ :\ Parameters)\ (V\ :\ E\ →\ ℝ)}
\leanline{4}{(hV\ :\ PotentialAssumptions\ P\ V)\ :}
\leanline{4}{ULA.ulaLawStep\ V\ P.step\ (invariantLaw\ P\ V\ hV)\ =\ invariantLaw\ P\ V\ hV}
\end{leancode}
\leansource{\texttt{invariantLaw\_\allowbreak{}mixing} \quad WassersteinMixing/\allowbreak{}Invariant.lean, lines 246–253}
\begin{leancode}
\leanline{0}{theorem\ invariantLaw\_mixing\ (P\ :\ Parameters)\ (V\ :\ E\ →\ ℝ)}
\leanline{4}{(hV\ :\ PotentialAssumptions\ P\ V)\ (μ\ :\ Measure\ E)\ [IsProbabilityMeasure\ μ]}
\leanline{4}{(hμ\ :\ MemLp\ (id\ :\ E\ →\ E)\ 2\ μ)\ (k\ :\ ℕ)\ :}
\leanline{4}{ULA.wasserstein2Real\ ((ULA.ulaLawStep\ V\ P.step)\textasciicircum{}[k]\ μ)\ (invariantLaw\ P\ V\ hV)\ ≤}
\leanline{6}{P.contraction\ \textasciicircum{}\ k\ *\ ULA.wasserstein2Real\ μ\ (invariantLaw\ P\ V\ hV)}
\end{leancode}
\leansource{\texttt{invariantLaw\_\allowbreak{}attraction} \quad WassersteinMixing/\allowbreak{}Invariant.lean, lines 256–265}
\begin{leancode}
\leanline{0}{theorem\ invariantLaw\_attraction\ (P\ :\ Parameters)\ (V\ :\ E\ →\ ℝ)}
\leanline{4}{(hV\ :\ PotentialAssumptions\ P\ V)\ (μ\ :\ Measure\ E)\ [IsProbabilityMeasure\ μ]}
\leanline{4}{(hμ\ :\ MemLp\ (id\ :\ E\ →\ E)\ 2\ μ)\ :}
\leanline{4}{Filter.Tendsto\ (fun\ k\ :\ ℕ\ ↦}
\leanline{6}{ULA.wasserstein2Real\ ((ULA.ulaLawStep\ V\ P.step)\textasciicircum{}[k]\ μ)\ (invariantLaw\ P\ V\ hV))}
\leanline{6}{Filter.atTop\ ({\leanfallback 𝓝}\ 0)}
\end{leancode}
\leanlabel{How the Lean proof follows the argument}
\texttt{ula\allowbreak{}Gradient\allowbreak{}Step\_\allowbreak{}lipschitz} specializes the proved common-noise inequality. \texttt{invariant\allowbreak{}Law} is the constructed Gaussian fixed-point law, not a supplied candidate. Its probability, finite moment and actual \texttt{ULA.\allowbreak{}ula\allowbreak{}Law\allowbreak{}Step} invariance are exported separately. \texttt{invariant\allowbreak{}Law\_\allowbreak{}mixing} invokes the proved two-law contraction with finite transport; \texttt{invariant\allowbreak{}Law\_\allowbreak{}attraction} squeezes between zero and the geometric upper bound.
\end{leanbox}


Earlier results used here: \cite[Section~3.17.4]{LangevinLean2026ULASC}; \cite[Section~3.21.2]{LangevinLean2026ULASC}; \cite[Section~3.34.1]{LangevinLean2026ULASC}; \cite[Section~3.36.7]{LangevinLean2026ULASC}.
\section{\texorpdfstring{Transport symmetry and the triangle inequality from actual coupling disintegration}{Transport symmetry and the triangle inequality from actual coupling disintegration}}
\label{proof-module-WassersteinMixing-TransportGeometry-lean}


The quadratic Wasserstein distance is the infimum of root-mean-square displacement over actual couplings. Symmetry follows by swapping coordinates. The triangle inequality is proved by disintegrating one coupling over its shared marginal, constructing a triple law, and applying the \(L^2\) triangle inequality. No optimal coupling or assumed metric triangle inequality is required.



Earlier results used here: \cite[Section~3.22]{LangevinLean2026ULASC}.
\subsection{\texorpdfstring{Swapping the endpoints of a coupling}{Swapping the endpoints of a coupling}}
\label{report-step-WassersteinMixing-TransportGeometry-lean-1}


\textbf{Claim.} Swapping the two coordinates of a coupling of \(\mu,\nu\) produces a coupling of \(\nu,\mu\) with the same cost. Consequently \(W_2(\mu,\nu)=W_2(\nu,\mu)\), both as an extended distance and under the ordinary real convention.




\textbf{Proof.} The coordinate-swap map interchanges the two marginal projections. The displacement changes from \(x-y\) to \(y-x\), preserving its norm and its quadratic cost. Every coupling therefore supplies an equal-cost coupling in the reverse direction. Taking the infimum in each direction gives the two inequalities required for equality. Applying the real-value operation preserves that equality.

\begin{leanbox}{Lean: \texttt{couplingSwap}, \texttt{couplingCost\_\allowbreak{}swap}, \texttt{wasserstein2\_\allowbreak{}symm} and 1 more}
\leansource{\texttt{couplingSwap} \quad WassersteinMixing/\allowbreak{}TransportGeometry.lean, lines 23–26}
\begin{leancode}
\leanline{0}{def\ couplingSwap\ \{μ\ ν\ :\ Measure\ H\}\ (c\ :\ ULA.Coupling\ μ\ ν)\ :\ ULA.Coupling\ ν\ μ\ where}
\leanline{2}{joint\ :=\ c.joint.map\ Prod.swap}
\leanline{2}{fst\ :=\ by\ rw\ [Measure.map\_map\ measurable\_fst\ measurable\_swap];\ exact\ c.snd}
\leanline{2}{snd\ :=\ by\ rw\ [Measure.map\_map\ measurable\_snd\ measurable\_swap];\ exact\ c.fst}
\end{leancode}
\leansource{\texttt{couplingCost\_\allowbreak{}swap} \quad WassersteinMixing/\allowbreak{}TransportGeometry.lean, lines 29–35}
\begin{leancode}
\leanline{0}{theorem\ couplingCost\_swap\ \{μ\ ν\ :\ Measure\ H\}\ (c\ :\ ULA.Coupling\ μ\ ν)\ :}
\leanline{4}{ULA.couplingCost\ (couplingSwap\ c).joint\ =\ ULA.couplingCost\ c.joint}
\end{leancode}
\leansource{\texttt{wasserstein2\_\allowbreak{}symm} \quad WassersteinMixing/\allowbreak{}TransportGeometry.lean, lines 38–45}
\begin{leancode}
\leanline{0}{theorem\ wasserstein2\_symm\ (μ\ ν\ :\ Measure\ H)\ :\ ULA.wasserstein2\ μ\ ν\ =\ ULA.wasserstein2\ ν\ μ}
\end{leancode}
\leansource{\texttt{wasserstein2Real\_\allowbreak{}symm} \quad WassersteinMixing/\allowbreak{}TransportGeometry.lean, lines 48–53}
\begin{leancode}
\leanline{0}{theorem\ wasserstein2Real\_symm\ (μ\ ν\ :\ Measure\ H)\ :}
\leanline{4}{ULA.wasserstein2Real\ μ\ ν\ =\ ULA.wasserstein2Real\ ν\ μ}
\end{leancode}
\leanlabel{How the Lean proof follows the argument}
\texttt{coupling\allowbreak{}Swap} pushes the actual joint measure through \texttt{Prod.\allowbreak{}swap} and proves its marginals. \texttt{coupling\allowbreak{}Cost\_\allowbreak{}swap} uses invariance of the norm of a reversed difference. \texttt{wasserstein2\_\allowbreak{}symm} passes these coupling-cost bounds to the two infima, and \texttt{wasserstein2\allowbreak{}Real\_\allowbreak{}symm} rewrites the equal extended values.
\end{leanbox}


Earlier results used here: \cite[Section~3.17.1]{LangevinLean2026ULASC}; \cite[Section~3.17.2]{LangevinLean2026ULASC}; \cite[Section~3.17.3]{LangevinLean2026ULASC}; \cite[Section~3.17.4]{LangevinLean2026ULASC}; \cite[Section~3.17.5]{LangevinLean2026ULASC}.
\subsection{\texorpdfstring{Constructing the glued triple and checking both pair marginals}{Constructing the glued triple and checking both pair marginals}}
\label{report-step-WassersteinMixing-TransportGeometry-lean-2}


\textbf{Construction.} In a standard Borel space, let \(\gamma_{XY}\) couple \(\mu,\nu\), and let \(\gamma_{YZ}\) couple \(\nu,\xi\). Disintegrate the second coupling as \(\gamma_{YZ}(dy,dz)=\nu(dy)K(y,dz)\). Define

\[R(dx,dy,dz)=\gamma_{XY}(dx,dy)K(y,dz).\]


\textbf{Claim.} This is a probability law with \((X,Y)\)-marginal \(\gamma_{XY}\) and \((Y,Z)\)-marginal \(\gamma_{YZ}\).




\textbf{Proof.} Integrating out \(z\) leaves \(\gamma_{XY}\), since the conditional kernel has mass one. Integrating a measurable function of \((y,z)\) first in \(K(y,dz)\) and then in \(\gamma_{XY}\) is the same as integrating it against \(\nu(dy)K(y,dz)\), because the second marginal of \(\gamma_{XY}\) is \(\nu\). The disintegration identity makes the latter measure \(\gamma_{YZ}\). Applying the first marginal equality to the entire space proves total mass one.

\begin{leanbox}{Lean: \texttt{glueCouplings}, \texttt{glueCouplings\_\allowbreak{}left}, \texttt{glueCouplings\_\allowbreak{}right} and 1 more}
\leansource{\texttt{glueCouplings} \quad WassersteinMixing/\allowbreak{}TransportGeometry.lean, lines 60–62}
\begin{leancode}
\leanline{0}{def\ glueCouplings\ (cXY\ :\ ULA.Coupling\ μ\ ν)\ (cYZ\ :\ ULA.Coupling\ ν\ ξ)\ :\ Measure\ ((H\ ×\ H)\ ×\ H)\ :=\ by}
\leanline{2}{have\ :=\ cYZ.isProbabilityMeasure}
\leanline{2}{exact\ cXY.joint\ ⊗ₘ\ (cYZ.joint.condKernel.comap\ Prod.snd\ measurable\_snd)}
\end{leancode}
\leansource{\texttt{glueCouplings\_\allowbreak{}left} \quad WassersteinMixing/\allowbreak{}TransportGeometry.lean, lines 66–70}
\begin{leancode}
\leanline{0}{theorem\ glueCouplings\_left\ (cXY\ :\ ULA.Coupling\ μ\ ν)\ (cYZ\ :\ ULA.Coupling\ ν\ ξ)\ :}
\leanline{4}{(glueCouplings\ cXY\ cYZ).map\ Prod.fst\ =\ cXY.joint}
\end{leancode}
\leansource{\texttt{glueCouplings\_\allowbreak{}right} \quad WassersteinMixing/\allowbreak{}TransportGeometry.lean, lines 75–95}
\begin{leancode}
\leanline{0}{theorem\ glueCouplings\_right\ (cXY\ :\ ULA.Coupling\ μ\ ν)\ (cYZ\ :\ ULA.Coupling\ ν\ ξ)\ :}
\leanline{4}{(glueCouplings\ cXY\ cYZ).map\ (fun\ p\ ↦\ (p.1.2,\ p.2))\ =\ cYZ.joint}
\end{leancode}
\leansource{\texttt{glueCouplings\_\allowbreak{}isProbability} \quad WassersteinMixing/\allowbreak{}TransportGeometry.lean, lines 99–104}
\begin{leancode}
\leanline{0}{theorem\ glueCouplings\_isProbability\ (cXY\ :\ ULA.Coupling\ μ\ ν)\ (cYZ\ :\ ULA.Coupling\ ν\ ξ)\ :}
\leanline{4}{IsProbabilityMeasure\ (glueCouplings\ cXY\ cYZ)}
\end{leancode}
\leanlabel{How the Lean proof follows the argument}
\texttt{glue\allowbreak{}Couplings} uses the actual conditional kernel of the second joint measure and pulls it back along the shared coordinate. \texttt{glue\allowbreak{}Couplings\_\allowbreak{}left} applies the composition-product marginal identity. \texttt{glue\allowbreak{}Couplings\_\allowbreak{}right} proves the projection by disintegration and the shared marginal, using equality of the relevant nonnegative integrals. \texttt{glue\allowbreak{}Couplings\_\allowbreak{}is\allowbreak{}Probability} checks the total mass.
\end{leanbox}


Earlier results used here: \cite[Section~3.17.1]{LangevinLean2026ULASC}.
\subsection{\texorpdfstring{Minkowski and the two coupling infima}{Minkowski and the two coupling infima}}
\label{report-step-WassersteinMixing-TransportGeometry-lean-3}


\textbf{Claim.} For probability laws on the stated measurable normed space,

\[W_2(\mu,\xi)\le W_2(\mu,\nu)+W_2(\nu,\xi),\]


including infinite values.




\textbf{Proof.} On the constructed triple law, \(X-Z=(X-Y)+(Y-Z)\). The extended \(L^2\) triangle inequality bounds its cost by the sum of the two pair costs, which equal the costs of the supplied couplings by the marginal identities. The \((X,Z)\)-projection is an actual coupling of \(\mu,\xi\), hence

\[W_2(\mu,\xi)\le\operatorname{cost}(\gamma_{XY})+\operatorname{cost}(\gamma_{YZ}).\]


This holds for every pair of couplings. Taking the infimum over each one gives the desired inequality. The argument neither selects an optimal coupling nor replaces an infinite cost by a real number.

\begin{leanbox}{Lean: \texttt{wasserstein2\_\allowbreak{}le\_\allowbreak{}couplingCost\_\allowbreak{}add}, \texttt{wasserstein2\_\allowbreak{}triangle}}
\leansource{\texttt{wasserstein2\_\allowbreak{}le\_\allowbreak{}couplingCost\_\allowbreak{}add} \quad WassersteinMixing/\allowbreak{}TransportGeometry.lean, lines 109–148}
\begin{leancode}
\leanline{0}{theorem\ wasserstein2\_le\_couplingCost\_add}
\leanline{4}{(cXY\ :\ ULA.Coupling\ μ\ ν)\ (cYZ\ :\ ULA.Coupling\ ν\ ξ)\ :}
\leanline{4}{ULA.wasserstein2\ μ\ ξ\ ≤\ ULA.couplingCost\ cXY.joint\ +\ ULA.couplingCost\ cYZ.joint}
\end{leancode}
\leansource{\texttt{wasserstein2\_\allowbreak{}triangle} \quad WassersteinMixing/\allowbreak{}TransportGeometry.lean, lines 152–162}
\begin{leancode}
\leanline{0}{theorem\ wasserstein2\_triangle\ (μ\ ν\ ξ\ :\ Measure\ H)}
\leanline{4}{[IsProbabilityMeasure\ μ]\ [IsProbabilityMeasure\ ν]\ [IsProbabilityMeasure\ ξ]\ :}
\leanline{4}{ULA.wasserstein2\ μ\ ξ\ ≤\ ULA.wasserstein2\ μ\ ν\ +\ ULA.wasserstein2\ ν\ ξ}
\end{leancode}
\leanlabel{How the Lean proof follows the argument}
\texttt{wasserstein2\_\allowbreak{}le\_\allowbreak{}coupling\allowbreak{}Cost\_\allowbreak{}add} verifies both endpoint laws and both pair costs of the glued triple, then applies \texttt{e\allowbreak{}Lp\allowbreak{}Norm\_\allowbreak{}add\_\allowbreak{}le}. \texttt{wasserstein2\_\allowbreak{}triangle} distributes the extended infima over addition and passes the arbitrary-coupling inequality through both infima.
\end{leanbox}


Earlier results used here: \cite[Section~3.17.1]{LangevinLean2026ULASC}; \cite[Section~3.17.2]{LangevinLean2026ULASC}; \cite[Section~3.17.3]{LangevinLean2026ULASC}; \cite[Section~3.17.9]{LangevinLean2026ULASC}.
\subsection{\texorpdfstring{Finiteness and the ordinary real triangle inequality}{Finiteness and the ordinary real triangle inequality}}
\label{report-step-WassersteinMixing-TransportGeometry-lean-4}


\textbf{Claim.} If \(W_2(\mu,\nu)\) and \(W_2(\nu,\xi)\) are finite, then \(W_2(\mu,\xi)\) is finite and the same triangle inequality holds for their ordinary real values.




\textbf{Proof.} The two finite distances have finite sum. The proved extended triangle inequality bounds the endpoint distance by that sum, proving its finiteness. Conversion of an extended inequality with finite upper side preserves order; the real value of the finite sum is the sum of the real values. These facts give the ordinary real triangle inequality with every required finiteness premise explicit.

\begin{leanbox}{Lean: \texttt{wasserstein2\_\allowbreak{}triangle\_\allowbreak{}finite}, \texttt{wasserstein2Real\_\allowbreak{}triangle}}
\leansource{\texttt{wasserstein2\_\allowbreak{}triangle\_\allowbreak{}finite} \quad WassersteinMixing/\allowbreak{}TransportGeometry.lean, lines 166–171}
\begin{leancode}
\leanline{0}{theorem\ wasserstein2\_triangle\_finite\ (μ\ ν\ ξ\ :\ Measure\ H)}
\leanline{4}{[IsProbabilityMeasure\ μ]\ [IsProbabilityMeasure\ ν]\ [IsProbabilityMeasure\ ξ]}
\leanline{4}{(hμν\ :\ ULA.wasserstein2\ μ\ ν\ ≠\ ⊤)\ (hνξ\ :\ ULA.wasserstein2\ ν\ ξ\ ≠\ ⊤)\ :}
\leanline{4}{ULA.wasserstein2\ μ\ ξ\ ≠\ ⊤}
\end{leancode}
\leansource{\texttt{wasserstein2Real\_\allowbreak{}triangle} \quad WassersteinMixing/\allowbreak{}TransportGeometry.lean, lines 175–181}
\begin{leancode}
\leanline{0}{theorem\ wasserstein2Real\_triangle\ (μ\ ν\ ξ\ :\ Measure\ H)}
\leanline{4}{[IsProbabilityMeasure\ μ]\ [IsProbabilityMeasure\ ν]\ [IsProbabilityMeasure\ ξ]}
\leanline{4}{(hμν\ :\ ULA.wasserstein2\ μ\ ν\ ≠\ ⊤)\ (hνξ\ :\ ULA.wasserstein2\ ν\ ξ\ ≠\ ⊤)\ :}
\leanline{4}{ULA.wasserstein2Real\ μ\ ξ\ ≤\ ULA.wasserstein2Real\ μ\ ν\ +\ ULA.wasserstein2Real\ ν\ ξ}
\end{leancode}
\leanlabel{How the Lean proof follows the argument}
\texttt{wasserstein2\_\allowbreak{}triangle\_\allowbreak{}finite} combines the extended bound with finiteness of the sum. \texttt{wasserstein2\allowbreak{}Real\_\allowbreak{}triangle} applies \texttt{ENNReal.\allowbreak{}to\allowbreak{}Real\_\allowbreak{}mono} to that finite upper bound and uses \texttt{ENNReal.\allowbreak{}to\allowbreak{}Real\_\allowbreak{}add} with both finite premises.
\end{leanbox}


Earlier results used here: \cite[Section~3.17.3]{LangevinLean2026ULASC}; \cite[Section~3.17.4]{LangevinLean2026ULASC}.
\section{\texorpdfstring{The exact scale-invariant step and logarithmic iteration count}{The exact scale-invariant step and logarithmic iteration count}}
\label{proof-module-WassersteinMixing-Complexity-lean}


This section proves scalar accuracy arithmetic. For \(\alpha,\beta,d,\varepsilon>0\), define

\[h=\frac1\beta\min\left(1,\frac{\varepsilon}{12\sqrt d}\right),
\qquad K=\frac\beta\alpha\max\left(1,\frac{12\sqrt d}{\varepsilon}\right).\]


The constants \(6\) and \(12\), the maximum, minimum and logarithm are the literal quantities of Corollary 1.1. A convergence estimate for actual sampling laws is not established by this scalar argument.

\subsection{\texorpdfstring{The chosen step and its reciprocal}{The chosen step and its reciprocal}}
\label{report-step-WassersteinMixing-Complexity-lean-1}


\textbf{Claim.} Put \(m=\max(1,12\sqrt d/\varepsilon)\). Then \(h=1/(\beta m)>0\), \(h\le1/\beta\), and \(\alpha hK=1\).




\textbf{Proof.} If \(\varepsilon\le12\sqrt d\), the minimum in the step is its second term and the maximum is \(12\sqrt d/\varepsilon\). Otherwise they both select their first terms. In either case multiplication verifies \(h=1/(\beta m)\). Positivity follows from \(\beta>0\), \(m\ge1\), and the upper bound follows by replacing the minimum by \(1\). Finally substitute \(K=\beta m/\alpha\) and cancel the positive factors.

\begin{leanbox}{Lean: \texttt{paperStep}, \texttt{iterationPrefactor}, \texttt{paperStep\_\allowbreak{}eq\_\allowbreak{}inv} and 3 more}
\leansource{\texttt{paperStep} \quad WassersteinMixing/\allowbreak{}Complexity.lean, lines 19–20}
\begin{leancode}
\leanline{0}{def\ paperStep\ (β\ d\ ε\ :\ ℝ)\ :\ ℝ\ :=}
\leanline{2}{(1\ /\ β)\ *\ min\ 1\ (ε\ /\ (12\ *\ Real.sqrt\ d))}
\end{leancode}
\leansource{\texttt{iterationPrefactor} \quad WassersteinMixing/\allowbreak{}Complexity.lean, lines 23–24}
\begin{leancode}
\leanline{0}{def\ iterationPrefactor\ (α\ β\ d\ ε\ :\ ℝ)\ :\ ℝ\ :=}
\leanline{2}{(β\ /\ α)\ *\ max\ 1\ (12\ *\ Real.sqrt\ d\ /\ ε)}
\end{leancode}
\leansource{\texttt{paperStep\_\allowbreak{}eq\_\allowbreak{}inv} \quad WassersteinMixing/\allowbreak{}Complexity.lean, lines 27–40}
\begin{leancode}
\leanline{0}{theorem\ paperStep\_eq\_inv\ (β\ d\ ε\ :\ ℝ)\ (hβ\ :\ 0\ <\ β)\ (hd\ :\ 0\ <\ d)\ (hε\ :\ 0\ <\ ε)\ :}
\leanline{4}{paperStep\ β\ d\ ε\ =\ 1\ /\ (β\ *\ max\ 1\ (12\ *\ Real.sqrt\ d\ /\ ε))}
\end{leancode}
\leansource{\texttt{paperStep\_\allowbreak{}pos} \quad WassersteinMixing/\allowbreak{}Complexity.lean, lines 43–46}
\begin{leancode}
\leanline{0}{theorem\ paperStep\_pos\ (β\ d\ ε\ :\ ℝ)\ (hβ\ :\ 0\ <\ β)\ (hd\ :\ 0\ <\ d)\ (hε\ :\ 0\ <\ ε)\ :}
\leanline{4}{0\ <\ paperStep\ β\ d\ ε}
\end{leancode}
\leansource{\texttt{paperStep\_\allowbreak{}le\_\allowbreak{}inv\_\allowbreak{}beta} \quad WassersteinMixing/\allowbreak{}Complexity.lean, lines 49–52}
\begin{leancode}
\leanline{0}{theorem\ paperStep\_le\_inv\_beta\ (β\ d\ ε\ :\ ℝ)\ (hβ\ :\ 0\ <\ β)\ (\_hd\ :\ 0\ <\ d)\ (\_hε\ :\ 0\ <\ ε)\ :}
\leanline{4}{paperStep\ β\ d\ ε\ ≤\ 1\ /\ β}
\end{leancode}
\leansource{\texttt{step\_\allowbreak{}mul\_\allowbreak{}iterationPrefactor} \quad WassersteinMixing/\allowbreak{}Complexity.lean, lines 55–61}
\begin{leancode}
\leanline{0}{theorem\ step\_mul\_iterationPrefactor\ (α\ β\ d\ ε\ :\ ℝ)\ (hα\ :\ 0\ <\ α)}
\leanline{4}{(hβ\ :\ 0\ <\ β)\ (hd\ :\ 0\ <\ d)\ (hε\ :\ 0\ <\ ε)\ :}
\leanline{4}{(α\ *\ paperStep\ β\ d\ ε)\ *\ iterationPrefactor\ α\ β\ d\ ε\ =\ 1}
\end{leancode}
\leanlabel{How the Lean proof follows the argument}
\texttt{paper\allowbreak{}Step} and \texttt{iteration\allowbreak{}Prefactor} retain the exact minimum and maximum. \texttt{paper\allowbreak{}Step\_\allowbreak{}eq\_\allowbreak{}inv} checks the two branches explicitly. The positivity, stability and reciprocal-product lemmas clear only nonzero factors established from the positive hypotheses.
\end{leanbox}
\subsection{\texorpdfstring{Allocating half of the accuracy to the bias}{Allocating half of the accuracy to the bias}}
\label{report-step-WassersteinMixing-Complexity-lean-2}


\textbf{Claim.} For \(\alpha>0\), \(\beta,d\ge0\),

\[\sqrt\alpha\sqrt{(\beta/\alpha)\beta d}=\beta\sqrt d.\]


Both sides are nonnegative and their squares agree after cancellation of \(\alpha\), which proves the identity. For the selected step and \(d>0\), it follows that

\[\sqrt\alpha\,6h\sqrt{(\beta/\alpha)\beta d}
=\frac{6\sqrt d}{m}\le\frac\varepsilon2.\]


The last inequality is precisely \(12\sqrt d\le\varepsilon m\), supplied by the second term in the maximum.

\begin{leanbox}{Lean: \texttt{sqrt\_\allowbreak{}curvature\_\allowbreak{}scale}, \texttt{paperStep\_\allowbreak{}bias\_\allowbreak{}le}}
\leansource{\texttt{sqrt\_\allowbreak{}curvature\_\allowbreak{}scale} \quad WassersteinMixing/\allowbreak{}Complexity.lean, lines 64–68}
\begin{leancode}
\leanline{0}{theorem\ sqrt\_curvature\_scale\ (α\ β\ d\ :\ ℝ)\ (hα\ :\ 0\ <\ α)\ (hβ\ :\ 0\ ≤\ β)\ (hd\ :\ 0\ ≤\ d)\ :}
\leanline{4}{Real.sqrt\ α\ *\ Real.sqrt\ ((β\ /\ α)\ *\ β\ *\ d)\ =\ β\ *\ Real.sqrt\ d}
\end{leancode}
\leansource{\texttt{paperStep\_\allowbreak{}bias\_\allowbreak{}le} \quad WassersteinMixing/\allowbreak{}Complexity.lean, lines 71–88}
\begin{leancode}
\leanline{0}{theorem\ paperStep\_bias\_le\ (α\ β\ d\ ε\ :\ ℝ)\ (hα\ :\ 0\ <\ α)\ (hβ\ :\ 0\ <\ β)}
\leanline{4}{(hd\ :\ 0\ <\ d)\ (hε\ :\ 0\ <\ ε)\ :}
\leanline{4}{Real.sqrt\ α\ *\ (6\ *\ paperStep\ β\ d\ ε\ *\ Real.sqrt\ ((β\ /\ α)\ *\ β\ *\ d))\ ≤\ ε\ /\ 2}
\end{leancode}
\leanlabel{How the Lean proof follows the argument}
\texttt{sqrt\_\allowbreak{}curvature\_\allowbreak{}scale} proves equality by equality of squares and nonnegativity. \texttt{paper\allowbreak{}Step\_\allowbreak{}bias\_\allowbreak{}le} substitutes this identity and the reciprocal step formula, then uses the maximum's lower bound to prove the half-budget inequality.
\end{leanbox}
\subsection{\texorpdfstring{The transient estimate and the complete scalar accuracy bound}{The transient estimate and the complete scalar accuracy bound}}
\label{report-step-WassersteinMixing-Complexity-lean-3}


\textbf{Claim.} Let \(0<t\le1\), \(A\ge0\), \(\varepsilon>0\). If \(k\) is a natural number with \(\log(1+2A/\varepsilon)\le tk\), then

\[(1-t)^kA\le\varepsilon/2.\]


\textbf{Proof.} Since \(1-t\le e^{-t}\) and both are nonnegative, \((1-t)^k\le e^{-tk}\). The logarithmic hypothesis gives \(e^{-tk}\le(1+2A/\varepsilon)^{-1}\). Multiplication by \(A\ge0\) yields the result because \(A/(1+2A/\varepsilon)\le\varepsilon/2\). This includes \(A=0\) and \(t=1\).




Take \(t=\alpha h\), \(A=\sqrt\alpha R\), with \(R\ge0\) and \(0<\alpha\le\beta\). The chosen step gives \(t\le1\). If \(k\ge K\log(1+2\sqrt\alpha R/\varepsilon)\), then \(tK=1\) supplies the logarithmic hypothesis. Adding the transient and bias half-budgets proves

\[\sqrt\alpha\left((1-\alpha h)^kR+6h\sqrt{(\beta/\alpha)\beta d}\right)\le\varepsilon.\]\begin{leanbox}{Lean: \texttt{geometric\_\allowbreak{}transient\_\allowbreak{}le}, \texttt{paper\_\allowbreak{}complexity\_\allowbreak{}bound}}
\leansource{\texttt{geometric\_\allowbreak{}transient\_\allowbreak{}le} \quad WassersteinMixing/\allowbreak{}Complexity.lean, lines 92–114}
\begin{leancode}
\leanline{0}{theorem\ geometric\_transient\_le\ (t\ A\ ε\ :\ ℝ)\ (\_ht\ :\ 0\ <\ t)\ (ht1\ :\ t\ ≤\ 1)}
\leanline{4}{(hA\ :\ 0\ ≤\ A)\ (hε\ :\ 0\ <\ ε)\ (k\ :\ ℕ)}
\leanline{4}{(hk\ :\ Real.log\ (1\ +\ 2\ *\ A\ /\ ε)\ ≤\ t\ *\ k)\ :}
\leanline{4}{(1\ -\ t)\ \textasciicircum{}\ k\ *\ A\ ≤\ ε\ /\ 2}
\end{leancode}
\leansource{\texttt{paper\_\allowbreak{}complexity\_\allowbreak{}bound} \quad WassersteinMixing/\allowbreak{}Complexity.lean, lines 118–142}
\begin{leancode}
\leanline{0}{theorem\ paper\_complexity\_bound\ (α\ β\ d\ R\ ε\ :\ ℝ)\ (hα\ :\ 0\ <\ α)\ (hαβ\ :\ α\ ≤\ β)}
\leanline{4}{(hd\ :\ 0\ <\ d)\ (hR\ :\ 0\ ≤\ R)\ (hε\ :\ 0\ <\ ε)\ (k\ :\ ℕ)}
\leanline{4}{(hk\ :\ iterationPrefactor\ α\ β\ d\ ε\ *}
\leanline{6}{Real.log\ (1\ +\ 2\ *\ Real.sqrt\ α\ *\ R\ /\ ε)\ ≤\ (k\ :\ ℝ))\ :}
\leanline{4}{Real.sqrt\ α\ *\ ((1\ -\ α\ *\ paperStep\ β\ d\ ε)\ \textasciicircum{}\ k\ *\ R\ +}
\leanline{6}{6\ *\ paperStep\ β\ d\ ε\ *\ Real.sqrt\ ((β\ /\ α)\ *\ β\ *\ d))\ ≤\ ε}
\end{leancode}
\leanlabel{How the Lean proof follows the argument}
\texttt{geometric\_\allowbreak{}transient\_\allowbreak{}le} bounds the geometric factor by an exponential, uses monotonicity of the exponential to invert the logarithmic bound and checks the final quotient. \texttt{paper\_\allowbreak{}complexity\_\allowbreak{}bound} proves the admissible dimensionless step, invokes the reciprocal-product identity and adds the two half-budgets. No mixing hypothesis is hidden in these scalar statements.
\end{leanbox}
\section{\texorpdfstring{From stationary comparison to the actual invariant-law bounds}{From stationary comparison to the actual invariant-law bounds}}
\label{proof-module-WassersteinMixing-ComparisonTheorems-lean}


These reusable conclusions identify the comparison target with the actual Gibbs law and its chain with actual Gaussian ULA iterates. The constructed invariant law, its proved attraction and the genuine transport triangle give the exact bias, convergence and accuracy bounds once the concrete comparison is supplied.



Earlier results used here: \cite[Section~3.26]{LangevinLean2026ULASC}.
\subsection{\texorpdfstring{Actual law identifications and finite transport}{Actual law identifications and finite transport}}
\label{report-step-WassersteinMixing-ComparisonTheorems-lean-1}


\textbf{Set-up.} Identify the stationary reference law with \(\pi\), the chain law with \(K_h^n\pi\), and the parameter dimension with the physical dimension. The proved stationary comparison then bounds \(W_2(\pi,K_h^n\pi)\) uniformly and proves finiteness. Gaussian ULA iteration preserves probability, and finite second moments imply finite transport from any such iterate to the constructed invariant law \(\pi_h\).



\textbf{Proof.} Substitute the exact law identifications into the comparison coupling. Iteration is the actual Gaussian Markov recursion. Its proved contraction transports the finite initial second-moment distance to the invariant law.

\begin{leanbox}{Lean: \texttt{ULAPoissonComparison}, \texttt{stationary\_\allowbreak{}start\_\allowbreak{}bound}, \texttt{stationary\_\allowbreak{}start\_\allowbreak{}finite} and 2 more}
\leansource{\texttt{ULAPoissonComparison} \quad WassersteinMixing/\allowbreak{}ComparisonTheorems.lean, lines 30–37}
\begin{leancode}
\leanline{0}{structure\ ULAPoissonComparison\ (S\ :\ PoissonProcess\ H\ E\ Q\ P.step\ P.step\_pos)}
\leanline{4}{(score\ :\ H)\ (V\ :\ E\ →\ ℝ)\ extends\ PoissonComparison\ S\ score\ where}
\leanline{2}{dimension\_eq\ :\ P.dimension\ =\ Module.finrank\ ℝ\ E}
\leanline{2}{target\_law\ :\ Q.map\ (reference\ 0)\ =\ ULA.gibbsMeasure\ V}
\leanline{2}{chain\_law\ :\ ∀\ n,\ Q.map\ (chain\ n)\ =}
\leanline{4}{(ULA.ulaLawStep\ V\ P.step)\textasciicircum{}[n]\ (ULA.gibbsMeasure\ V)}
\leanline{0}{\strut}
\leanline{0}{namespace\ ULAPoissonComparison}
\end{leancode}
\leansource{\texttt{stationary\_\allowbreak{}start\_\allowbreak{}bound} \quad WassersteinMixing/\allowbreak{}ComparisonTheorems.lean, lines 43–49}
\begin{leancode}
\leanline{0}{theorem\ stationary\_start\_bound\ (n\ :\ ℕ)\ :}
\leanline{4}{ULA.wasserstein2Real\ (ULA.gibbsMeasure\ V)}
\leanline{6}{((ULA.ulaLawStep\ V\ P.step)\textasciicircum{}[n]\ (ULA.gibbsMeasure\ V))\ ≤\ P.bias}
\end{leancode}
\leansource{\texttt{stationary\_\allowbreak{}start\_\allowbreak{}finite} \quad WassersteinMixing/\allowbreak{}ComparisonTheorems.lean, lines 50–54}
\begin{leancode}
\leanline{0}{theorem\ stationary\_start\_finite\ (n\ :\ ℕ)\ :}
\leanline{4}{ULA.wasserstein2\ (ULA.gibbsMeasure\ V)}
\leanline{6}{((ULA.ulaLawStep\ V\ P.step)\textasciicircum{}[n]\ (ULA.gibbsMeasure\ V))\ ≠\ ⊤}
\end{leancode}
\leansource{\texttt{ula\_\allowbreak{}iterate\_\allowbreak{}isProbability} \quad WassersteinMixing/\allowbreak{}ComparisonTheorems.lean, lines 59–65}
\begin{leancode}
\leanline{0}{theorem\ ula\_iterate\_isProbability\ (P\ :\ Parameters)\ (V\ :\ E\ →\ ℝ)}
\leanline{4}{(hV\ :\ PotentialAssumptions\ P\ V)\ (μ\ :\ Measure\ E)\ [IsProbabilityMeasure\ μ]\ (n\ :\ ℕ)\ :}
\leanline{4}{IsProbabilityMeasure\ ((ULA.ulaLawStep\ V\ P.step)\textasciicircum{}[n]\ μ)}
\end{leancode}
\leansource{\texttt{transport\_\allowbreak{}invariant\_\allowbreak{}finite} \quad WassersteinMixing/\allowbreak{}ComparisonTheorems.lean, lines 69–75}
\begin{leancode}
\leanline{0}{theorem\ transport\_invariant\_finite\ (P\ :\ Parameters)\ (V\ :\ E\ →\ ℝ)}
\leanline{4}{(hV\ :\ PotentialAssumptions\ P\ V)\ (μ\ :\ Measure\ E)\ [IsProbabilityMeasure\ μ]}
\leanline{4}{(hμ\ :\ MemLp\ (id\ :\ E\ →\ E)\ 2\ μ)\ (n\ :\ ℕ)\ :}
\leanline{4}{ULA.wasserstein2\ ((ULA.ulaLawStep\ V\ P.step)\textasciicircum{}[n]\ μ)\ (invariantLaw\ P\ V\ hV)\ ≠\ ⊤}
\end{leancode}
\leanlabel{How the Lean proof follows the argument}
\texttt{ULAPoisson\allowbreak{}Comparison} adds only dimension and actual-law equalities to the generic comparison. The stationary-start theorems substitute these equalities. Probability and finiteness of iterates use the actual Gaussian Markov law and the independently constructed invariant law.
\end{leanbox}


Earlier results used here: \cite[Section~3.17.3]{LangevinLean2026ULASC}; \cite[Section~3.17.4]{LangevinLean2026ULASC}; \cite[Section~3.21.2]{LangevinLean2026ULASC}; \cite[Section~3.26.3]{LangevinLean2026ULASC}; \cite[Section~3.33.11]{LangevinLean2026ULASC}; \cite[Section~3.34.1]{LangevinLean2026ULASC}; \cite[Section~3.36.5]{LangevinLean2026ULASC}; \cite[Section~3.36.7]{LangevinLean2026ULASC}.
\subsection{\texorpdfstring{The exact invariant bias and finite-time convergence}{The exact invariant bias and finite-time convergence}}
\label{report-step-WassersteinMixing-ComparisonTheorems-lean-2}


\textbf{Claim.}

\[W_2(\pi,\pi_h)\le6h\sqrt{\kappa\beta d},\qquad W_2(K_h^n\mu,\pi)\le(1-\alpha h)^nW_2(\mu,\pi_h)+6h\sqrt{\kappa\beta d},\quad\kappa=\beta/\alpha.\]

\textbf{Proof.} Apply the transport triangle through \(K_h^n\pi\). Its first leg has the uniform stationary-start bound; its second tends to zero by actual invariant attraction. Taking this proved limit yields the bias. For the finite-time statement, apply the triangle through \(\pi_h\), bound its first leg by the proved ULA contraction and its second by the bias. All laws and finite distances have already been established.

\begin{leanbox}{Lean: \texttt{bias\_\allowbreak{}from\_\allowbreak{}comparison}, \texttt{convergence\_\allowbreak{}from\_\allowbreak{}comparison}}
\leansource{\texttt{bias\_\allowbreak{}from\_\allowbreak{}comparison} \quad WassersteinMixing/\allowbreak{}ComparisonTheorems.lean, lines 81–101}
\begin{leancode}
\leanline{0}{theorem\ bias\_from\_comparison\ (hV\ :\ PotentialAssumptions\ P\ V)\ (C\ :\ ULAPoissonComparison\ S\ score\ V)\ :}
\leanline{4}{ULA.wasserstein2Real\ (ULA.gibbsMeasure\ V)\ (invariantLaw\ P\ V\ hV)\ ≤}
\leanline{6}{6\ *\ P.step\ *\ Real.sqrt\ (P.condition\ *\ P.beta\ *\ P.dimension)}
\end{leancode}
\leansource{\texttt{convergence\_\allowbreak{}from\_\allowbreak{}comparison} \quad WassersteinMixing/\allowbreak{}ComparisonTheorems.lean, lines 105–124}
\begin{leancode}
\leanline{0}{theorem\ convergence\_from\_comparison\ (hV\ :\ PotentialAssumptions\ P\ V)}
\leanline{4}{(C\ :\ ULAPoissonComparison\ S\ score\ V)\ (μ\ :\ Measure\ E)\ [IsProbabilityMeasure\ μ]}
\leanline{4}{(hμ\ :\ MemLp\ (id\ :\ E\ →\ E)\ 2\ μ)\ (n\ :\ ℕ)\ :}
\leanline{4}{ULA.wasserstein2Real\ ((ULA.ulaLawStep\ V\ P.step)\textasciicircum{}[n]\ μ)\ (ULA.gibbsMeasure\ V)\ ≤}
\leanline{6}{P.contraction\ \textasciicircum{}\ n\ *\ ULA.wasserstein2Real\ μ\ (invariantLaw\ P\ V\ hV)\ +}
\leanline{8}{6\ *\ P.step\ *\ Real.sqrt\ (P.condition\ *\ P.beta\ *\ P.dimension)}
\end{leancode}
\leanlabel{How the Lean proof follows the argument}
\texttt{bias\_\allowbreak{}from\_\allowbreak{}comparison} passes the actual invariant-attraction limit through a uniform triangle bound. \texttt{convergence\_\allowbreak{}from\_\allowbreak{}comparison} combines \texttt{invariant\allowbreak{}Law\_\allowbreak{}mixing}, the proved transport triangle and that bias. Invariant existence, contraction and limit convergence are not comparison fields.
\end{leanbox}


Earlier results used here: \cite[Section~3.17.4]{LangevinLean2026ULASC}; \cite[Section~3.21.2]{LangevinLean2026ULASC}; \cite[Section~3.26.13]{LangevinLean2026ULASC}; \cite[Section~3.26.15]{LangevinLean2026ULASC}; \cite[Section~3.26.3]{LangevinLean2026ULASC}; \cite[Section~3.36.7]{LangevinLean2026ULASC}.
\subsection{\texorpdfstring{The source accuracy scale}{The source accuracy scale}}
\label{report-step-WassersteinMixing-ComparisonTheorems-lean-3}


\textbf{Claim.} With the source step choice and its logarithmic iteration threshold, \(\sqrt\alpha\,W_2(K_h^n\mu,\pi)\le\varepsilon\).



\textbf{Proof.} Multiply the actual convergence bound by \(\sqrt\alpha\). The step choice makes the bias contribution at most \(\varepsilon/2\); the iteration threshold makes the geometric transient at most \(\varepsilon/2\). These exact numerical estimates are proved independently in the complexity calculation. Positive step together with the displayed step formula excludes the zero-dimensional exceptional branch when this accuracy statement is used.

\begin{leanbox}{Lean: \texttt{accuracy\_\allowbreak{}from\_\allowbreak{}comparison}}
\leansource{\texttt{accuracy\_\allowbreak{}from\_\allowbreak{}comparison} \quad WassersteinMixing/\allowbreak{}ComparisonTheorems.lean, lines 128–150}
\begin{leancode}
\leanline{0}{theorem\ accuracy\_from\_comparison\ (hV\ :\ PotentialAssumptions\ P\ V)}
\leanline{4}{(C\ :\ ULAPoissonComparison\ S\ score\ V)\ (μ\ :\ Measure\ E)\ [IsProbabilityMeasure\ μ]}
\leanline{4}{(hμ\ :\ MemLp\ (id\ :\ E\ →\ E)\ 2\ μ)\ (ε\ :\ ℝ)\ (hε\ :\ 0\ <\ ε)}
\leanline{4}{(hstep\ :\ P.step\ =\ paperStep\ P.beta\ P.dimension\ ε)\ (n\ :\ ℕ)}
\leanline{4}{(hn\ :\ iterationPrefactor\ P.alpha\ P.beta\ P.dimension\ ε\ *}
\leanline{6}{Real.log\ (1\ +\ 2\ *\ Real.sqrt\ P.alpha\ *}
\leanline{8}{ULA.wasserstein2Real\ μ\ (invariantLaw\ P\ V\ hV)\ /\ ε)\ ≤\ (n\ :\ ℝ))\ :}
\leanline{4}{Real.sqrt\ P.alpha\ *}
\leanline{6}{ULA.wasserstein2Real\ ((ULA.ulaLawStep\ V\ P.step)\textasciicircum{}[n]\ μ)\ (ULA.gibbsMeasure\ V)\ ≤\ ε}
\end{leancode}
\leanlabel{How the Lean proof follows the argument}
\texttt{accuracy\_\allowbreak{}from\_\allowbreak{}comparison} applies the exact numerical theorem \texttt{paper\_\allowbreak{}complexity\_\allowbreak{}bound} to the actual initial distance from the constructed invariant law. The proof derives the needed positive dimension from the positive step and its prescribed formula.
\end{leanbox}


Earlier results used here: \cite[Section~3.17.4]{LangevinLean2026ULASC}; \cite[Section~3.26.3]{LangevinLean2026ULASC}; \cite[Section~3.36.7]{LangevinLean2026ULASC}.
\section{\texorpdfstring{Assembling the comparison from the actual diffusion and Euler chain}{Assembling the comparison from the actual diffusion and Euler chain}}
\label{proof-module-WassersteinMixing-ConcreteComparison-lean}


The actual reference, Gaussian Euler chain, true score increments and random averaged Hessian supply every primitive datum of the stationary comparison. This reusable constructor takes exact filtration, evaluation and conditional-source identifications. The actual Poisson process and fully realized comparison prove all of them from the potential, so none remains an additional assumption of the source endpoints.

\subsection{\texorpdfstring{Every comparison input is supplied by its actual construction}{Every comparison input is supplied by its actual construction}}
\label{report-step-WassersteinMixing-ConcreteComparison-lean-1}


\textbf{Construction.} Given the Poisson process identified with the actual Brownian filtration and grid evaluations, use the actual stationary diffusion reference, its synchronous Euler chain, the true gradient increments and the actual random segment Hessian.



\textbf{Proof.} Gibbs stationarity and the true Gaussian iteration theorem give the actual target and chain laws. The score moment, compact-time integrability, increment constant, endpoint identity and exact conditional defect mean supply the local estimates. The actual synchronous recursion and averaged-Hessian secant identity give the exact error equation. Actual past adaptation and true curvature symmetry and form bounds supply the predictable geometry. The physical dimension equality is exact. Thus every comparison field is derived from actual objects; no shifted-energy, bias or transport conclusion is supplied as input.

\begin{leanbox}{Lean: \texttt{comparisonFromActualProcess}}
\leansource{\texttt{comparisonFromActualProcess} \quad WassersteinMixing/\allowbreak{}ConcreteComparison.lean, lines 23–68}
\begin{leancode}
\leanline{0}{def\ comparisonFromActualProcess}
\leanline{4}{(S\ :\ PoissonProcess\ (Lp\ (E\ P.dimension)\ 2\ (ULA.gibbsMeasure\ V))}
\leanline{6}{(E\ P.dimension)\ (actualProbability\ P\ V)\ P.step\ P.step\_pos)}
\leanline{4}{(hfiltration\ :\ S.filtration\ =\ actualFiltration\ P\ V\ hV)}
\leanline{4}{(heval\ :\ ∀\ n,\ S.eval\ n\ =\ actualGridEvaluation\ P\ V\ hV\ n)}
\leanline{4}{(hmean\ :\ ∀\ n,}
\leanline{6}{(condExpL2\ (E\ P.dimension)\ ℝ\ (S.filtration.le\ n)\ (actualDefect\ P\ V\ hV\ n)\ :}
\leanline{8}{Lp\ (E\ P.dimension)\ 2\ (actualProbability\ P\ V))\ =}
\leanline{10}{S.eval\ n\ (S.poissonSource\ (actualObservableScore\ P\ V\ hV)))\ :}
\leanline{4}{ULAPoissonComparison\ S\ (actualObservableScore\ P\ V\ hV)\ V\ where}
\leanline{2}{reference\ :=\ actualReference\ P\ V\ hV}
\leanline{2}{chain\ :=\ actualChain\ P\ V\ hV}
\leanline{2}{increment\ :=\ actualScoreIncrement\ P\ V\ hV}
\leanline{2}{curvature\ :=\ actualCurvature\ P\ V\ hV}
\leanline{2}{score\_bound\ :=\ actualObservableScore\_energy\ P\ V\ hV}
\leanline{2}{reference\_stationary\ :=\ fun\ n\ ↦\ (actualReference\_stationary\ P\ V\ hV\ n).trans}
\leanline{4}{(actualReference\_stationary\ P\ V\ hV\ 0).symm}
\leanline{2}{initial\_equal\ :=\ actual\_initial\_equal\ P\ V\ hV}
\leanline{2}{increment\_memLp\ :=\ actualScoreIncrement\_memLp\ P\ V\ hV}
\leanline{2}{increment\_bound\ :=\ fun\ n\ s\ hs\ ↦\ actualScoreIncrement\_energy\ P\ V\ hV\ n\ s\ ⟨hs.1.le,\ hs.2⟩}
\leanline{2}{increment\_endpoint\ :=\ by}
\leanline{4}{intro\ n}
\leanline{4}{rw\ [heval\ n,\ heval\ (n\ +\ 1),\ actualGridEvaluation\_score,\ actualGridEvaluation\_score]}
\leanline{4}{exact\ actualScoreIncrement\_endpoint\ P\ V\ hV\ n}
\leanline{2}{defect\_mean\ :=\ hmean}
\leanline{2}{error\_step\ :=\ actual\_curvature\_error\_step\ P\ V\ hV}
\leanline{2}{error\_adapted\ :=\ by}
\leanline{4}{intro\ n}
\leanline{4}{have\ hm\ :\ (S.filtration\ n\ :\ MeasurableSpace\ (ActualSpace\ P))\ =\ actualPast\ P\ V\ hV\ n\ :=}
\leanline{6}{congrArg\ (fun\ F\ :\ Filtration\ ℕ}
\leanline{8}{(inferInstance\ :\ MeasurableSpace\ (ActualSpace\ P))\ ↦\ F\ n)\ hfiltration}
\leanline{4}{rw\ [hm]}
\leanline{4}{exact\ actualError\_adapted\ P\ V\ hV\ n}
\leanline{2}{curvature\_adapted\ :=\ by}
\leanline{4}{intro\ n\ U\ hU}
\leanline{4}{have\ hm\ :\ (S.filtration\ n\ :\ MeasurableSpace\ (ActualSpace\ P))\ =\ actualPast\ P\ V\ hV\ n\ :=}
\leanline{6}{congrArg\ (fun\ F\ :\ Filtration\ ℕ}
\leanline{8}{(inferInstance\ :\ MeasurableSpace\ (ActualSpace\ P))\ ↦\ F\ n)\ hfiltration}
\leanline{4}{rw\ [hm]\ at\ hU\ ⊢}
\leanline{4}{exact\ actualCurvature\_preserves\_past\ P\ V\ hV\ n\ U\ hU}
\leanline{2}{curvature\_symmetric\ :=\ actualCurvature\_symmetric\ P\ V\ hV}
\leanline{2}{curvature\_lower\ :=\ actualCurvature\_lower\ P\ V\ hV}
\leanline{2}{curvature\_upper\ :=\ actualCurvature\_upper\ P\ V\ hV}
\leanline{2}{dimension\_eq\ :=\ by\ simp\ [E]}
\leanline{2}{target\_law\ :=\ actualReference\_stationary\ P\ V\ hV\ 0}
\leanline{2}{chain\_law\ :=\ actualChain\_law\ P\ V\ hV}
\end{leancode}
\leanlabel{How the Lean proof follows the argument}
\texttt{comparison\allowbreak{}From\allowbreak{}Actual\allowbreak{}Process} fills \texttt{ULAPoisson\allowbreak{}Comparison} using \texttt{actual\allowbreak{}Observable\allowbreak{}Score\_\allowbreak{}energy}, the actual local-defect facts, \texttt{actual\_\allowbreak{}curvature\_\allowbreak{}error\_\allowbreak{}step}, actual adaptation and the true curvature theorems. It transports only the exact filtration/\allowbreak{}evaluation/\allowbreak{}source identifications required by the supplied actual process.
\end{leanbox}


Earlier results used here: \cite[Section~3.26.3]{LangevinLean2026ULASC}.
\section{\texorpdfstring{The fully realized actual-law stationary comparison}{The fully realized actual-law stationary comparison}}
\label{proof-module-WassersteinMixing-RealizedComparison-lean}


The integrated actual conditional defect mean is identified with the constructed Poisson source. This closes the actual-process comparison: every score, path, curvature, conditional mean and law is supplied by its construction from the potential.

\subsection{\texorpdfstring{Closing the exact source identity and all comparison fields}{Closing the exact source identity and all comparison fields}}
\label{report-step-WassersteinMixing-RealizedComparison-lean-1}


\textbf{Claim.} For every real target observable \(f\), the actual process source is

\[(I-P_h)G_hf=\int_0^h(f-P_sf)\,ds.\]

For the actual gradient observable, its evaluated value equals the conditional mean of the genuine local defect. Thus the actual reference and Euler chain form the complete stationary Poisson comparison.



\textbf{Proof.} The semigroup corrector theorem gives the integral source for the extracted real operator. That operator is already identified with the actual real Langevin transition, so the integrands agree. The independently proved conditional defect-mean formula has precisely this integral. Substitute the source identity to obtain the exact comparison mean. The concrete comparison constructor then fills every remaining field from the proved actual scores, increments, curvature, adaptation and laws.

\begin{leanbox}{Lean: \texttt{actualPoissonSource\_\allowbreak{}eq\_\allowbreak{}integral}, \texttt{actualPoissonProcess\_\allowbreak{}defect\_\allowbreak{}mean}, \texttt{actualULAComparison}}
\leansource{\texttt{actualPoissonSource\_\allowbreak{}eq\_\allowbreak{}integral} \quad WassersteinMixing/\allowbreak{}RealizedComparison.lean, lines 22–33}
\begin{leancode}
\leanline{0}{theorem\ actualPoissonSource\_eq\_integral\ (f\ :\ Lp\ (E\ P.dimension)\ 2\ (ULA.gibbsMeasure\ V))\ :}
\leanline{4}{(actualPoissonProcess\ P\ V\ hV).poissonSource\ f\ =}
\leanline{6}{∫\ s\ in\ 0..P.step,\ f\ -\ CanonicalSemigroup.op\ ({\leanfallback 𝕜}\ :=\ ℝ)\ P\ V\ hV\ s\ f}
\end{leancode}
\leansource{\texttt{actualPoissonProcess\_\allowbreak{}defect\_\allowbreak{}mean} \quad WassersteinMixing/\allowbreak{}RealizedComparison.lean, lines 36–42}
\begin{leancode}
\leanline{0}{theorem\ actualPoissonProcess\_defect\_mean\ (n\ :\ ℕ)\ :}
\leanline{4}{(condExpL2\ (E\ P.dimension)\ ℝ\ ((actualPoissonProcess\ P\ V\ hV).filtration.le\ n)}
\leanline{6}{(actualDefect\ P\ V\ hV\ n)\ :\ Lp\ (E\ P.dimension)\ 2\ (actualProbability\ P\ V))\ =}
\leanline{8}{(actualPoissonProcess\ P\ V\ hV).eval\ n}
\leanline{10}{((actualPoissonProcess\ P\ V\ hV).poissonSource\ (actualObservableScore\ P\ V\ hV))}
\end{leancode}
\leansource{\texttt{actualULAComparison} \quad WassersteinMixing/\allowbreak{}RealizedComparison.lean, lines 46–49}
\begin{leancode}
\leanline{0}{def\ actualULAComparison\ :}
\leanline{4}{ULAPoissonComparison\ (actualPoissonProcess\ P\ V\ hV)\ (actualObservableScore\ P\ V\ hV)\ V\ :=}
\leanline{2}{comparisonFromActualProcess\ P\ V\ hV\ (actualPoissonProcess\ P\ V\ hV)\ rfl\ (fun\ \_\ ↦\ rfl)}
\leanline{4}{(actualPoissonProcess\_defect\_mean\ P\ V\ hV)}
\end{leancode}
\leanlabel{How the Lean proof follows the argument}
\texttt{actual\allowbreak{}Poisson\allowbreak{}Source\_\allowbreak{}eq\_\allowbreak{}integral} rewrites the general real corrector source using \texttt{actual\allowbreak{}Real\allowbreak{}Op\_\allowbreak{}eq\_\allowbreak{}op}. \texttt{actual\allowbreak{}Poisson\allowbreak{}Process\_\allowbreak{}defect\_\allowbreak{}mean} matches the independently proved actual defect mean. \texttt{actual\allowbreak{}ULAComparison} invokes the concrete constructor with definitional filtration/\allowbreak{}evaluation identities and this proved mean, adding no realization assumption.
\end{leanbox}


Earlier results used here: \cite[Section~3.26.3]{LangevinLean2026ULASC}.
\section{\texorpdfstring{The original Hessian hypothesis implies the potential contract}{The original Hessian hypothesis implies the potential contract}}
\label{proof-module-WassersteinMixing-CurvatureAssumptions-lean}


Assume \(V\in C^2(H)\) and \(\alpha I\preceq\nabla^2V(x)\preceq\beta I\), with \(0<\alpha\le\beta\). These are the source assumptions. The first-order convexity and smoothness estimates used by the probabilistic construction are consequences, proved here with the same constants.

\subsection{\texorpdfstring{The sharp derivative and gradient bounds}{The sharp derivative and gradient bounds}}
\label{report-step-WassersteinMixing-CurvatureAssumptions-lean-1}


\textbf{Claim.} The quadratic-form bounds imply \(\|\nabla^2V(x)\|\le\beta\), hence \(\|\nabla V(x)-\nabla V(y)\|\le\beta\|x-y\|\).



\textbf{Proof.} Scale the symmetric Hessian by \(\beta^{-1}\). It is a positive contraction in quadratic form, so its squared norm on a vector is bounded by that quadratic form and then by the vector norm squared. Rescaling yields the Hessian norm estimate. The mean-value bound for a differentiable map with uniformly bounded derivative gives the gradient estimate.

\begin{leanbox}{Lean: \texttt{HessianAssumptions}, \texttt{hessian\_\allowbreak{}norm\_\allowbreak{}le\_\allowbreak{}of\_\allowbreak{}forms}, \texttt{gradient\_\allowbreak{}lipschitz\_\allowbreak{}of\_\allowbreak{}hessian}}
\leansource{\texttt{HessianAssumptions} \quad WassersteinMixing/\allowbreak{}CurvatureAssumptions.lean, lines 20–23}
\begin{leancode}
\leanline{0}{structure\ HessianAssumptions\ (P\ :\ Parameters)\ (V\ :\ H\ →\ ℝ)\ :\ Prop\ where}
\leanline{2}{regularity\ :\ ContDiff\ ℝ\ 2\ V}
\leanline{2}{lower\ :\ ∀\ x\ v,\ P.alpha\ *\ ‖v‖\ \textasciicircum{}\ 2\ ≤\ inner\ ℝ\ v\ (hessianOperator\ V\ x\ v)}
\leanline{2}{upper\ :\ ∀\ x\ v,\ inner\ ℝ\ v\ (hessianOperator\ V\ x\ v)\ ≤\ P.beta\ *\ ‖v‖\ \textasciicircum{}\ 2}
\end{leancode}
\leansource{\texttt{hessian\_\allowbreak{}norm\_\allowbreak{}le\_\allowbreak{}of\_\allowbreak{}forms} \quad WassersteinMixing/\allowbreak{}CurvatureAssumptions.lean, lines 26–50}
\begin{leancode}
\leanline{0}{theorem\ hessian\_norm\_le\_of\_forms\ (P\ :\ Parameters)\ (V\ :\ H\ →\ ℝ)}
\leanline{4}{(hV\ :\ HessianAssumptions\ P\ V)\ (x\ :\ H)\ :\ ‖hessianOperator\ V\ x‖\ ≤\ P.beta}
\end{leancode}
\leansource{\texttt{gradient\_\allowbreak{}lipschitz\_\allowbreak{}of\_\allowbreak{}hessian} \quad WassersteinMixing/\allowbreak{}CurvatureAssumptions.lean, lines 53–59}
\begin{leancode}
\leanline{0}{theorem\ gradient\_lipschitz\_of\_hessian\ (P\ :\ Parameters)\ (V\ :\ H\ →\ ℝ)}
\leanline{4}{(hV\ :\ HessianAssumptions\ P\ V)\ :}
\leanline{4}{LipschitzWith\ ⟨P.beta,\ P.beta\_pos.le⟩\ (gradient\ V)}
\end{leancode}
\leanlabel{How the Lean proof follows the argument}
\texttt{Hessian\allowbreak{}Assumptions} records the actual second-derivative inequalities. \texttt{hessian\_\allowbreak{}norm\_\allowbreak{}le\_\allowbreak{}of\_\allowbreak{}forms} applies the proved positive-contraction inequality to the scaled Hessian. \texttt{gradient\_\allowbreak{}lipschitz\_\allowbreak{}of\_\allowbreak{}hessian} applies the derivative mean-value theorem with that sharp bound.
\end{leanbox}
\subsection{\texorpdfstring{Integrating the lower Hessian bound}{Integrating the lower Hessian bound}}
\label{report-step-WassersteinMixing-CurvatureAssumptions-lean-2}


\textbf{Claim.} For all \(x,y\),

\[V(y)\ge V(x)+\langle\nabla V(x),y-x\rangle+\frac\alpha2\|y-x\|^2.\]

\textbf{Proof.} Along \(x+t(y-x)\), the derivative of \(\langle\nabla V(x+t(y-x)),y-x\rangle\) is at least \(\alpha\|y-x\|^2\). Integrating once bounds that scalar gradient; integrating again gives the displayed supporting inequality. Subtracting \(\alpha\|x\|^2/2\) yields supporting planes for a convex function. To verify convexity directly, combine its two supporting inequalities at a convex combination; the linear terms cancel. Thus the original Hessian hypothesis supplies every part of the potential contract.

\begin{leanbox}{Lean: \texttt{strong\_\allowbreak{}support\_\allowbreak{}of\_\allowbreak{}hessian}, \texttt{convexOn\_\allowbreak{}of\_\allowbreak{}supporting\_\allowbreak{}planes}, \texttt{strongConvexity\_\allowbreak{}of\_\allowbreak{}hessian} and 1 more}
\leansource{\texttt{strong\_\allowbreak{}support\_\allowbreak{}of\_\allowbreak{}hessian} \quad WassersteinMixing/\allowbreak{}CurvatureAssumptions.lean, lines 62–103}
\begin{leancode}
\leanline{0}{theorem\ strong\_support\_of\_hessian\ (P\ :\ Parameters)\ (V\ :\ H\ →\ ℝ)}
\leanline{4}{(hV\ :\ HessianAssumptions\ P\ V)\ (x\ y\ :\ H)\ :}
\leanline{4}{V\ x\ +\ inner\ ℝ\ (gradient\ V\ x)\ (y\ -\ x)\ +\ P.alpha\ /\ 2\ *\ ‖y\ -\ x‖\ \textasciicircum{}\ 2\ ≤\ V\ y}
\end{leancode}
\leansource{\texttt{convexOn\_\allowbreak{}of\_\allowbreak{}supporting\_\allowbreak{}planes} \quad WassersteinMixing/\allowbreak{}CurvatureAssumptions.lean, lines 108–122}
\begin{leancode}
\leanline{0}{theorem\ convexOn\_of\_supporting\_planes\ (F\ :\ H\ →\ ℝ)\ (G\ :\ H\ →\ H)}
\leanline{4}{(hplane\ :\ ∀\ x\ y,\ F\ x\ +\ inner\ ℝ\ (G\ x)\ (y\ -\ x)\ ≤\ F\ y)\ :\ ConvexOn\ ℝ\ Set.univ\ F}
\end{leancode}
\leansource{\texttt{strongConvexity\_\allowbreak{}of\_\allowbreak{}hessian} \quad WassersteinMixing/\allowbreak{}CurvatureAssumptions.lean, lines 125–135}
\begin{leancode}
\leanline{0}{theorem\ strongConvexity\_of\_hessian\ (P\ :\ Parameters)\ (V\ :\ H\ →\ ℝ)}
\leanline{4}{(hV\ :\ HessianAssumptions\ P\ V)\ :\ StrongConvexOn\ Set.univ\ P.alpha\ V}
\end{leancode}
\leansource{\texttt{HessianAssumptions.\allowbreak{}toPotentialAssumptions} \quad WassersteinMixing/\allowbreak{}CurvatureAssumptions.lean, lines 138–142}
\begin{leancode}
\leanline{0}{theorem\ HessianAssumptions.toPotentialAssumptions\ (P\ :\ Parameters)\ (V\ :\ H\ →\ ℝ)}
\leanline{4}{(hV\ :\ HessianAssumptions\ P\ V)\ :\ PotentialAssumptions\ P\ V}
\end{leancode}
\leanlabel{How the Lean proof follows the argument}
\texttt{strong\_\allowbreak{}support\_\allowbreak{}of\_\allowbreak{}hessian} uses the actual gradient derivative and interval fundamental theorem. \texttt{convex\allowbreak{}On\_\allowbreak{}of\_\allowbreak{}supporting\_\allowbreak{}planes} proves the convex-combination inequality algebraically. \texttt{strong\allowbreak{}Convexity\_\allowbreak{}of\_\allowbreak{}hessian} and \texttt{Hessian\allowbreak{}Assumptions.\allowbreak{}to\allowbreak{}Potential\allowbreak{}Assumptions} assemble the consequences without additional hypotheses.
\end{leanbox}
\section{\texorpdfstring{The source Wasserstein estimates from the original Hessian bounds}{The source Wasserstein estimates from the original Hessian bounds}}
\label{proof-module-WassersteinMixing-Theorem-lean}


The displayed endpoints concern the actual normalized Gibbs law, the constructed invariant law of ordinary independent-Gaussian ULA and its true finite-time iterates. Their potential input is exactly a twice continuously differentiable potential with the source Hessian bounds; the analytic comparison is constructed internally.

\subsection{\texorpdfstring{The exact invariant bias estimate}{The exact invariant bias estimate}}
\label{report-step-WassersteinMixing-Theorem-lean-1}


\textbf{Claim.} For \(V\in C^2(\mathbb R^d)\), \(\alpha I\preceq\nabla^2V\preceq\beta I\), \(0<\alpha\le\beta\), \(0<h\le1/\beta\), and \(\kappa=\beta/\alpha\),

\[W_2(\pi,\pi_h)\le6h\sqrt{\kappa\beta d}.\]

\textbf{Proof.} The original Hessian bounds imply the first-order potential contract with unchanged constants. From that contract the actual reversible Langevin semigroup, bounded corrector, score and synchronous comparison are constructed. Apply the proved comparison-to-invariant theorem to this constructed comparison and the genuinely constructed Gaussian ULA invariant law. No spectral representation, conditional mean, curvature operator or comparison estimate is an additional input.

\begin{leanbox}{Lean: \texttt{bias\_\allowbreak{}bound}}
\leansource{\texttt{bias\_\allowbreak{}bound} \quad WassersteinMixing/\allowbreak{}Theorem.lean, lines 21–27}
\begin{leancode}
\leanline{0}{theorem\ bias\_bound\ (P\ :\ Parameters)\ (V\ :\ E\ P.dimension\ →\ ℝ)}
\leanline{4}{(hA\ :\ HessianAssumptions\ P\ V)\ :}
\leanline{4}{ULA.wasserstein2Real\ (ULA.gibbsMeasure\ V)}
\leanline{6}{(invariantLaw\ P\ V\ hA.toPotentialAssumptions)\ ≤}
\leanline{8}{6\ *\ P.step\ *\ Real.sqrt\ (P.condition\ *\ P.beta\ *\ P.dimension)}
\end{leancode}
\leanlabel{How the Lean proof follows the argument}
\texttt{bias\_\allowbreak{}bound} invokes \texttt{bias\_\allowbreak{}from\_\allowbreak{}comparison} on \texttt{actual\allowbreak{}ULAComparison}, using only \texttt{to\allowbreak{}Potential\allowbreak{}Assumptions}. Its conclusion is the actual transport distance from \texttt{gibbs\allowbreak{}Measure} to the constructed \texttt{invariant\allowbreak{}Law}.
\end{leanbox}


Earlier results used here: \cite[Section~3.17.4]{LangevinLean2026ULASC}; \cite[Section~3.26.3]{LangevinLean2026ULASC}.
\subsection{\texorpdfstring{The actual finite-time Gaussian ULA estimate}{The actual finite-time Gaussian ULA estimate}}
\label{report-step-WassersteinMixing-Theorem-lean-2}


\textbf{Claim.} For an initial probability law \(\mu\) with finite second moment,

\[W_2(K_h^n\mu,\pi)\le(1-\alpha h)^nW_2(\mu,\pi_h)+6h\sqrt{\kappa\beta d}.\]

\textbf{Proof.} Instantiate the general convergence theorem with the constructed actual comparison. The invariant law exists and attracts Gaussian ULA iterates by its own contraction construction. The true transport triangle, actual ULA contraction and preceding bias estimate yield the displayed two-term bound. The initial moment hypothesis supplies finiteness of every transport distance used in this argument.

\begin{leanbox}{Lean: \texttt{convergence\_\allowbreak{}bound}}
\leansource{\texttt{convergence\_\allowbreak{}bound} \quad WassersteinMixing/\allowbreak{}Theorem.lean, lines 31–39}
\begin{leancode}
\leanline{0}{theorem\ convergence\_bound\ (P\ :\ Parameters)\ (V\ :\ E\ P.dimension\ →\ ℝ)}
\leanline{4}{(hA\ :\ HessianAssumptions\ P\ V)\ (μ\ :\ Measure\ (E\ P.dimension))\ [IsProbabilityMeasure\ μ]}
\leanline{4}{(hμ\ :\ MemLp\ (id\ :\ E\ P.dimension\ →\ E\ P.dimension)\ 2\ μ)\ (n\ :\ ℕ)\ :}
\leanline{4}{ULA.wasserstein2Real\ ((ULA.ulaLawStep\ V\ P.step)\textasciicircum{}[n]\ μ)\ (ULA.gibbsMeasure\ V)\ ≤}
\leanline{6}{P.contraction\ \textasciicircum{}\ n\ *\ ULA.wasserstein2Real\ μ}
\leanline{8}{(invariantLaw\ P\ V\ hA.toPotentialAssumptions)\ +}
\leanline{10}{6\ *\ P.step\ *\ Real.sqrt\ (P.condition\ *\ P.beta\ *\ P.dimension)}
\end{leancode}
\leanlabel{How the Lean proof follows the argument}
\texttt{convergence\_\allowbreak{}bound} supplies the constructed comparison to \texttt{convergence\_\allowbreak{}from\_\allowbreak{}comparison}. Its iterate is the actual independent-Gaussian law step, and its only distributional premise is the source finite-second-moment initial law.
\end{leanbox}


Earlier results used here: \cite[Section~3.17.4]{LangevinLean2026ULASC}; \cite[Section~3.26.3]{LangevinLean2026ULASC}; \cite[Section~3.36.7]{LangevinLean2026ULASC}.
\subsection{\texorpdfstring{The source scale-invariant mixing complexity}{The source scale-invariant mixing complexity}}
\label{report-step-WassersteinMixing-Theorem-lean-3}


\textbf{Claim.} With the source step size

\[h=\beta^{-1}\min\left\{1,\frac{\varepsilon}{12\sqrt d}\right\}\]

and an integer iteration count at least

\[\kappa\max\left\{1,\frac{12\sqrt d}{\varepsilon}\right\}\log\left(1+\frac{2\sqrt\alpha\,W_2(\mu,\pi_h)}{\varepsilon}\right),\]

one has \(\sqrt\alpha\,W_2(K_h^n\mu,\pi)\le\varepsilon\).



\textbf{Proof.} Apply the proved exact complexity arithmetic to the actual finite-time convergence bound. The chosen step reduces the scaled bias to at most half the tolerance, and the logarithmic threshold controls the scaled geometric transient by the other half. Positive step and this formula imply the positive-dimension condition needed for its displayed numerical form.

\begin{leanbox}{Lean: \texttt{accuracy\_\allowbreak{}bound}}
\leansource{\texttt{accuracy\_\allowbreak{}bound} \quad WassersteinMixing/\allowbreak{}Theorem.lean, lines 44–54}
\begin{leancode}
\leanline{0}{theorem\ accuracy\_bound\ (P\ :\ Parameters)\ (V\ :\ E\ P.dimension\ →\ ℝ)}
\leanline{4}{(hA\ :\ HessianAssumptions\ P\ V)\ (μ\ :\ Measure\ (E\ P.dimension))\ [IsProbabilityMeasure\ μ]}
\leanline{4}{(hμ\ :\ MemLp\ (id\ :\ E\ P.dimension\ →\ E\ P.dimension)\ 2\ μ)\ (ε\ :\ ℝ)\ (hε\ :\ 0\ <\ ε)}
\leanline{4}{(hstep\ :\ P.step\ =\ paperStep\ P.beta\ P.dimension\ ε)\ (n\ :\ ℕ)}
\leanline{4}{(hn\ :\ iterationPrefactor\ P.alpha\ P.beta\ P.dimension\ ε\ *}
\leanline{6}{Real.log\ (1\ +\ 2\ *\ Real.sqrt\ P.alpha\ *}
\leanline{8}{ULA.wasserstein2Real\ μ\ (invariantLaw\ P\ V\ hA.toPotentialAssumptions)\ /\ ε)\ ≤\ (n\ :\ ℝ))\ :}
\leanline{4}{Real.sqrt\ P.alpha\ *}
\leanline{6}{ULA.wasserstein2Real\ ((ULA.ulaLawStep\ V\ P.step)\textasciicircum{}[n]\ μ)\ (ULA.gibbsMeasure\ V)\ ≤\ ε}
\end{leancode}
\leanlabel{How the Lean proof follows the argument}
\texttt{accuracy\_\allowbreak{}bound} invokes \texttt{accuracy\_\allowbreak{}from\_\allowbreak{}comparison} with the actual constructed comparison and original Hessian assumptions. The source step and prefactor are the literal functions \texttt{paper\allowbreak{}Step} and \texttt{iteration\allowbreak{}Prefactor}, whose exact arithmetic was proved separately.
\end{leanbox}


Earlier results used here: \cite[Section~3.17.4]{LangevinLean2026ULASC}; \cite[Section~3.26.3]{LangevinLean2026ULASC}; \cite[Section~3.36.7]{LangevinLean2026ULASC}.
\section{\texorpdfstring{Coordinate independence of Gibbs targets, Gaussian ULA and transport}{Coordinate independence of Gibbs targets, Gaussian ULA and transport}}
\label{proof-module-WassersteinMixing-IsometryTransport-lean}


Let \(e:H\to F\) be a linear isometric equivalence of finite-dimensional real inner product spaces with their Borel structures. The normalized Gibbs law for \(V\) pushes forward to the Gibbs law for \(V\circ e^{-1}\); standard Gaussian ULA and all its iterates commute with the same coordinate change; and the exact quadratic transport distance is unchanged. These proved identities identify the IPLA coordinates with the Euclidean coordinates of the sharp theorem.

\subsection{\texorpdfstring{The true gradient in the new coordinates}{The true gradient in the new coordinates}}
\label{report-step-WassersteinMixing-IsometryTransport-lean-1}


\textbf{Claim.} For differentiable \(V\),

\[\nabla(V\circ e^{-1})(y)=e\,\nabla V(e^{-1}y).\]


\textbf{Proof.} The chain rule gives \(D(V\circ e^{-1})(y)[w]=DV(e^{-1}y)[e^{-1}w]\). Identifying derivatives with gradients and using the inner-product identity for the linear isometry gives the inner product of \(w\) with the asserted right side. Equality of all such pairings proves equality of the gradients.

\begin{leanbox}{Lean: \texttt{gradient\_\allowbreak{}comp\_\allowbreak{}isometry}}
\leansource{\texttt{gradient\_\allowbreak{}comp\_\allowbreak{}isometry} \quad WassersteinMixing/\allowbreak{}IsometryTransport.lean, lines 29–40}
\begin{leancode}
\leanline{0}{theorem\ gradient\_comp\_isometry\ (e\ :\ H\ ≃ₗᵢ[ℝ]\ F)\ (V\ :\ H\ →\ ℝ)}
\leanline{4}{(hV\ :\ Differentiable\ ℝ\ V)\ (y\ :\ F)\ :}
\leanline{4}{gradient\ (V\ ∘\ e.symm)\ y\ =\ e\ (gradient\ V\ (e.symm\ y))}
\end{leancode}
\leanlabel{How the Lean proof follows the argument}
\texttt{gradient\_\allowbreak{}comp\_\allowbreak{}isometry} combines the Fréchet derivative chain rule with inner-product extensionality and \texttt{inner\_\allowbreak{}map\_\allowbreak{}map}. Its differentiability hypothesis is exactly what the chain rule needs; no regularity of a separate coordinate gradient is postulated.
\end{leanbox}
\subsection{\texorpdfstring{Volume preservation gives the exact normalized Gibbs law}{Volume preservation gives the exact normalized Gibbs law}}
\label{report-step-WassersteinMixing-IsometryTransport-lean-2}


\textbf{Claim.} Write \(Z_V=\int e^{-V}\) and \(\pi_V=Z_V^{-1}e^{-V}\,dx\), using the library's extended-valued normalization convention. Then

\[Z_{V\circ e^{-1}}=Z_V,\qquad e_\#\pi_V=\pi_{V\circ e^{-1}}.\]


\textbf{Proof.} A linear isometric equivalence preserves finite-dimensional volume. Change variables through \(e\) in the nonnegative integral defining the partition function. For a measurable set, the same change of variables identifies the weighted-volume measures before normalization. The equality of partition functions then identifies their scalar normalizations. These identities hold at the defining measure level; they do not assume that the Gibbs law is invariant or even separately assume a positive finite partition function.

\begin{leanbox}{Lean: \texttt{gibbsPartition\_\allowbreak{}comp\_\allowbreak{}isometry}, \texttt{gibbsMeasure\_\allowbreak{}map\_\allowbreak{}isometry}}
\leansource{\texttt{gibbsPartition\_\allowbreak{}comp\_\allowbreak{}isometry} \quad WassersteinMixing/\allowbreak{}IsometryTransport.lean, lines 44–49}
\begin{leancode}
\leanline{0}{theorem\ gibbsPartition\_comp\_isometry\ (e\ :\ H\ ≃ₗᵢ[ℝ]\ F)\ (V\ :\ H\ →\ ℝ)\ :}
\leanline{4}{ULA.gibbsPartition\ (V\ ∘\ e.symm)\ =\ ULA.gibbsPartition\ V}
\end{leancode}
\leansource{\texttt{gibbsMeasure\_\allowbreak{}map\_\allowbreak{}isometry} \quad WassersteinMixing/\allowbreak{}IsometryTransport.lean, lines 54–66}
\begin{leancode}
\leanline{0}{theorem\ gibbsMeasure\_map\_isometry\ (e\ :\ H\ ≃ₗᵢ[ℝ]\ F)\ (V\ :\ H\ →\ ℝ)\ :}
\leanline{4}{(ULA.gibbsMeasure\ V).map\ e\ =\ ULA.gibbsMeasure\ (V\ ∘\ e.symm)}
\end{leancode}
\leanlabel{How the Lean proof follows the argument}
\texttt{gibbs\allowbreak{}Partition\_\allowbreak{}comp\_\allowbreak{}isometry} uses \texttt{measure\allowbreak{}Preserving} and the nonnegative change-of-variables formula. \texttt{gibbs\allowbreak{}Measure\_\allowbreak{}map\_\allowbreak{}isometry} proves equality on measurable sets after pushing the normalized density forward, applying the same measure-preserving embedding identity. Later potential assumptions establish the probability and moment properties independently.
\end{leanbox}


Earlier results used here: \cite[Section~3.26.1]{LangevinLean2026ULASC}; \cite[Section~3.26.2]{LangevinLean2026ULASC}; \cite[Section~3.26.3]{LangevinLean2026ULASC}.
\subsection{\texorpdfstring{Transporting actual couplings in both directions}{Transporting actual couplings in both directions}}
\label{report-step-WassersteinMixing-IsometryTransport-lean-3}


\textbf{Claim.} Every coupling \(c\) of \(\mu,\nu\) pushes forward through \((x,y)\mapsto(e x,e y)\) to a coupling of \(e_\#\mu,e_\#\nu\) with the same root-mean-square cost. Consequently, including possibly infinite distances,

\[W_2(e_\#\mu,e_\#\nu)=W_2(\mu,\nu).\]


\textbf{Proof.} Composing each coordinate projection with the pair map proves the two marginal identities. Since \(e(x-y)=e x-e y\) and \(e\) preserves norms, the nonnegative quadratic displacement integral and its square root are unchanged. Taking the infimum over original couplings gives one inequality. Apply the same construction using \(e^{-1}\) to the pushed-forward laws to get the reverse inequality. Converting the resulting equality of extended distances to the library's real-distance convention preserves equality; no order argument through an infinite-to-real conversion is used.

\begin{leanbox}{Lean: \texttt{isometryCoupling}, \texttt{isometryCoupling\_\allowbreak{}cost}, \texttt{wasserstein2\_\allowbreak{}map\_\allowbreak{}isometry\_\allowbreak{}le} and 2 more}
\leansource{\texttt{isometryCoupling} \quad WassersteinMixing/\allowbreak{}IsometryTransport.lean, lines 69–81}
\begin{leancode}
\leanline{0}{def\ isometryCoupling\ (e\ :\ H\ ≃ₗᵢ[ℝ]\ F)\ \{μ\ ν\ :\ Measure\ H\}\ (c\ :\ ULA.Coupling\ μ\ ν)\ :}
\leanline{4}{ULA.Coupling\ (μ.map\ e)\ (ν.map\ e)\ where}
\leanline{2}{joint\ :=\ c.joint.map\ (Prod.map\ e\ e)}
\leanline{2}{fst\ :=\ by}
\leanline{4}{rw\ [Measure.map\_map\ measurable\_fst\ (by\ fun\_prop)]}
\leanline{4}{have\ h\ :=\ congrArg\ (Measure.map\ e)\ c.fst}
\leanline{4}{rw\ [Measure.map\_map\ e.continuous.measurable\ measurable\_fst]\ at\ h}
\leanline{4}{exact\ h}
\leanline{2}{snd\ :=\ by}
\leanline{4}{rw\ [Measure.map\_map\ measurable\_snd\ (by\ fun\_prop)]}
\leanline{4}{have\ h\ :=\ congrArg\ (Measure.map\ e)\ c.snd}
\leanline{4}{rw\ [Measure.map\_map\ e.continuous.measurable\ measurable\_snd]\ at\ h}
\leanline{4}{exact\ h}
\end{leancode}
\leansource{\texttt{isometryCoupling\_\allowbreak{}cost} \quad WassersteinMixing/\allowbreak{}IsometryTransport.lean, lines 85–91}
\begin{leancode}
\leanline{0}{theorem\ isometryCoupling\_cost\ (e\ :\ H\ ≃ₗᵢ[ℝ]\ F)\ \{μ\ ν\ :\ Measure\ H\}}
\leanline{4}{(c\ :\ ULA.Coupling\ μ\ ν)\ :}
\leanline{4}{ULA.couplingCost\ (isometryCoupling\ e\ c).joint\ =\ ULA.couplingCost\ c.joint}
\end{leancode}
\leansource{\texttt{wasserstein2\_\allowbreak{}map\_\allowbreak{}isometry\_\allowbreak{}le} \quad WassersteinMixing/\allowbreak{}IsometryTransport.lean, lines 95–100}
\begin{leancode}
\leanline{0}{theorem\ wasserstein2\_map\_isometry\_le\ (e\ :\ H\ ≃ₗᵢ[ℝ]\ F)\ (μ\ ν\ :\ Measure\ H)\ :}
\leanline{4}{ULA.wasserstein2\ (μ.map\ e)\ (ν.map\ e)\ ≤\ ULA.wasserstein2\ μ\ ν}
\end{leancode}
\leansource{\texttt{wasserstein2\_\allowbreak{}map\_\allowbreak{}isometry} \quad WassersteinMixing/\allowbreak{}IsometryTransport.lean, lines 105–110}
\begin{leancode}
\leanline{0}{theorem\ wasserstein2\_map\_isometry\ (e\ :\ H\ ≃ₗᵢ[ℝ]\ F)\ (μ\ ν\ :\ Measure\ H)\ :}
\leanline{4}{ULA.wasserstein2\ (μ.map\ e)\ (ν.map\ e)\ =\ ULA.wasserstein2\ μ\ ν}
\end{leancode}
\leansource{\texttt{wasserstein2Real\_\allowbreak{}map\_\allowbreak{}isometry} \quad WassersteinMixing/\allowbreak{}IsometryTransport.lean, lines 114–116}
\begin{leancode}
\leanline{0}{theorem\ wasserstein2Real\_map\_isometry\ (e\ :\ H\ ≃ₗᵢ[ℝ]\ F)\ (μ\ ν\ :\ Measure\ H)\ :}
\leanline{4}{ULA.wasserstein2Real\ (μ.map\ e)\ (ν.map\ e)\ =\ ULA.wasserstein2Real\ μ\ ν}
\end{leancode}
\leanlabel{How the Lean proof follows the argument}
\texttt{isometry\allowbreak{}Coupling} constructs the pushed joint measure and proves both marginals by composition of measure maps. \texttt{isometry\allowbreak{}Coupling\_\allowbreak{}cost} uses the exact norm equality in the \texttt{e\allowbreak{}Lp\allowbreak{}Norm} change-of-variables formula. \texttt{wasserstein2\_\allowbreak{}map\_\allowbreak{}isometry\_\allowbreak{}le} passes this cost through the coupling infimum; \texttt{wasserstein2\_\allowbreak{}map\_\allowbreak{}isometry} uses the inverse map for equality. \texttt{wasserstein2\allowbreak{}Real\_\allowbreak{}map\_\allowbreak{}isometry} applies the real-value operation only to this established equality.
\end{leanbox}


Earlier results used here: \cite[Section~3.17.1]{LangevinLean2026ULASC}; \cite[Section~3.17.2]{LangevinLean2026ULASC}; \cite[Section~3.17.3]{LangevinLean2026ULASC}; \cite[Section~3.17.4]{LangevinLean2026ULASC}; \cite[Section~3.17.5]{LangevinLean2026ULASC}.
\subsection{\texorpdfstring{The actual independent-Gaussian transition commutes with reindexing}{The actual independent-Gaussian transition commutes with reindexing}}
\label{report-step-WassersteinMixing-IsometryTransport-lean-4}


\textbf{Claim.} If \(V\) is differentiable with a globally Lipschitz gradient and \(\mu\) is an s-finite measure, the actual Gaussian ULA transition satisfies

\[e_\#K_{V,h}\mu=K_{V\circ e^{-1},h}(e_\#\mu).\]


\textbf{Proof.} Its deterministic proposal is \(x-h\nabla V(x)+\sqrt{2h}\,\xi\). The gradient identity and linearity show that applying \(e\) to this proposal equals the proposal in the new coordinates using \(e x,e\xi\). A standard isotropic Gaussian law is preserved by the linear isometry. Therefore the image of the product of the initial law and the independent Gaussian law is exactly the product measure used by the transformed transition. Pushing this identity through the measurable proposal proves the equality of the actual update laws.

\begin{leanbox}{Lean: \texttt{ulaLawStep\_\allowbreak{}map\_\allowbreak{}isometry}}
\leansource{\texttt{ulaLawStep\_\allowbreak{}map\_\allowbreak{}isometry} \quad WassersteinMixing/\allowbreak{}IsometryTransport.lean, lines 120–143}
\begin{leancode}
\leanline{0}{theorem\ ulaLawStep\_map\_isometry\ (e\ :\ H\ ≃ₗᵢ[ℝ]\ F)\ (V\ :\ H\ →\ ℝ)}
\leanline{4}{(hV\ :\ Differentiable\ ℝ\ V)\ (L\ :\ ℝ≥0)\ (hLip\ :\ LipschitzWith\ L\ (gradient\ V))}
\leanline{4}{(h\ :\ ℝ)\ (μ\ :\ Measure\ H)\ [SFinite\ μ]\ :}
\leanline{4}{(ULA.ulaLawStep\ V\ h\ μ).map\ e\ =}
\leanline{6}{ULA.ulaLawStep\ (V\ ∘\ e.symm)\ h\ (μ.map\ e)}
\end{leancode}
\leanlabel{How the Lean proof follows the argument}
\texttt{ula\allowbreak{}Law\allowbreak{}Step\_\allowbreak{}map\_\allowbreak{}isometry} first proves the pointwise proposal identity, then rewrites the standard Gaussian pushforward with \texttt{std\allowbreak{}Gaussian\_\allowbreak{}map}. \texttt{Measure.\allowbreak{}map\_\allowbreak{}prod\_\allowbreak{}map} and composition of measure maps turn it into the desired transition identity. The Lipschitz assumption establishes measurability of the true gradient and proposal.
\end{leanbox}


Earlier results used here: \cite[Section~3.34.1]{LangevinLean2026ULASC}; \cite[Section~3.36.5]{LangevinLean2026ULASC}; \cite[Section~3.36.7]{LangevinLean2026ULASC}.
\subsection{\texorpdfstring{Finite second moments and every Gaussian iterate}{Finite second moments and every Gaussian iterate}}
\label{report-step-WassersteinMixing-IsometryTransport-lean-5}


\textbf{Claim.} A finite second moment is preserved by \(e\). Under the potential contract, for every probability initialization and every nonnegative integer \(k\),

\[e_\#K_{V,h}^{\,k}\mu=K_{V\circ e^{-1},h}^{\,k}(e_\#\mu).\]


\textbf{Proof.} The coordinate map preserves the pointwise norm, so its image has the same square-integrability as the original random vector. For the iterate identity, the case \(k=0\) is the initial law. If the equality holds at \(k\), apply the proved one-step covariance to that iterate and then substitute the induction hypothesis. The already constructed Gaussian transition preserves probability, supplying the measure-finiteness hypothesis needed at every induction step.

\begin{leanbox}{Lean: \texttt{memLp\_\allowbreak{}id\_\allowbreak{}map\_\allowbreak{}isometry}, \texttt{ula\_\allowbreak{}iterate\_\allowbreak{}map\_\allowbreak{}isometry}}
\leansource{\texttt{memLp\_\allowbreak{}id\_\allowbreak{}map\_\allowbreak{}isometry} \quad WassersteinMixing/\allowbreak{}IsometryTransport.lean, lines 147–151}
\begin{leancode}
\leanline{0}{theorem\ memLp\_id\_map\_isometry\ (e\ :\ H\ ≃ₗᵢ[ℝ]\ F)\ (μ\ :\ Measure\ H)}
\leanline{4}{(hμ\ :\ MemLp\ (id\ :\ H\ →\ H)\ 2\ μ)\ :\ MemLp\ (id\ :\ F\ →\ F)\ 2\ (μ.map\ e)}
\end{leancode}
\leansource{\texttt{ula\_\allowbreak{}iterate\_\allowbreak{}map\_\allowbreak{}isometry} \quad WassersteinMixing/\allowbreak{}IsometryTransport.lean, lines 154–165}
\begin{leancode}
\leanline{0}{theorem\ ula\_iterate\_map\_isometry\ (P\ :\ Parameters)\ (e\ :\ H\ ≃ₗᵢ[ℝ]\ F)}
\leanline{4}{(V\ :\ H\ →\ ℝ)\ (hV\ :\ PotentialAssumptions\ P\ V)}
\leanline{4}{(μ\ :\ Measure\ H)\ [IsProbabilityMeasure\ μ]\ (n\ :\ ℕ)\ :}
\leanline{4}{((ULA.ulaLawStep\ V\ P.step)\textasciicircum{}[n]\ μ).map\ e\ =}
\leanline{6}{(ULA.ulaLawStep\ (V\ ∘\ e.symm)\ P.step)\textasciicircum{}[n]\ (μ.map\ e)}
\end{leancode}
\leanlabel{How the Lean proof follows the argument}
\texttt{mem\allowbreak{}Lp\_\allowbreak{}id\_\allowbreak{}map\_\allowbreak{}isometry} uses the mapped-measure characterization of \texttt{Mem\allowbreak{}Lp} and composition with the continuous linear map of the isometry. \texttt{ula\_\allowbreak{}iterate\_\allowbreak{}map\_\allowbreak{}isometry} inducts on the actual update count, supplies the proved probability instance for each iterate, and invokes \texttt{ula\allowbreak{}Law\allowbreak{}Step\_\allowbreak{}map\_\allowbreak{}isometry}. These are distributional identities, so the application does not replace the IPLA recursion with a different sampler.
\end{leanbox}


Earlier results used here: \cite[Section~3.36.7]{LangevinLean2026ULASC}.
\section{\texorpdfstring{The exact IPLA curvature, time and coordinate scaling}{The exact IPLA curvature, time and coordinate scaling}}
\label{proof-module-WassersteinMixing-IPLACurvature-lean}


For \(N>0\) particles, let the joint potential \(U(\theta,x)\) be twice continuously differentiable, \(m\)-strongly convex and have an \(L\)-Lipschitz gradient, with \(0<m\le L\). For a positive IPLA step \(\gamma\) satisfying \(\gamma L\le1\), the scaled state has dimension \(d_N=d_\theta+Nd_x\), potential \(V_N(z)=\sum_{i=1}^NU(\theta,\sqrt N z_i)\), curvature constants \(Nm,NL\), and ULA step \(\gamma/N\). Every lifted potential hypothesis and every coordinate-change identity is proved from these assumptions.



Earlier results used here: \cite[Section~3.1]{LangevinLean2026IPLA}.
\subsection{\texorpdfstring{Deriving the lifted numerical parameters and regularity}{Deriving the lifted numerical parameters and regularity}}
\label{report-step-WassersteinMixing-IPLACurvature-lean-1}


\textbf{Claim.} The scaled IPLA parameters are

\[\alpha_N=Nm,\qquad\beta_N=NL,\qquad h_N=\frac{\gamma}{N},\qquad d_N=d_\theta+Nd_x.\]


They satisfy \(0<\alpha_N\le\beta_N\), \(h_N>0\), and \(h_N\beta_N\le1\). The finite-sum potential \(V_N\) belongs to \(C^2\).




\textbf{Proof.} Positivity of the particle number and of \(m,\gamma\) gives the strict positive parameters. Multiplying \(m\le L\) by \(N\) gives the curvature ordering. Exact cancellation yields \(h_N\beta_N=\gamma L\le1\). Each map \(z\mapsto(\theta,\sqrt N z_i)\) is continuous linear, so composition with \(U\in C^2\) is twice continuously differentiable; a finite sum preserves this regularity.

\begin{leanbox}{Lean: \texttt{iplaParameters}, \texttt{ipla\_\allowbreak{}potential\_\allowbreak{}contDiff}}
\leansource{\texttt{iplaParameters} \quad WassersteinMixing/\allowbreak{}IPLACurvature.lean, lines 19–32}
\begin{leancode}
\leanline{0}{def\ iplaParameters\ (N\ dθ\ dx\ :\ ℕ)\ (m\ L\ γ\ :\ ℝ)\ (hN\ :\ 0\ <\ N)}
\leanline{4}{(hm\ :\ 0\ <\ m)\ (hmL\ :\ m\ ≤\ L)\ (hγ\ :\ 0\ <\ γ)\ (hγL\ :\ γ\ *\ L\ ≤\ 1)\ :\ Parameters\ where}
\leanline{2}{alpha\ :=\ N\ *\ m}
\leanline{2}{beta\ :=\ N\ *\ L}
\leanline{2}{step\ :=\ γ\ /\ N}
\leanline{2}{dimension\ :=\ dθ\ +\ N\ *\ dx}
\leanline{2}{alpha\_pos\ :=\ mul\_pos\ (by\ exact\_mod\_cast\ hN)\ hm}
\leanline{2}{alpha\_le\_beta\ :=\ mul\_le\_mul\_of\_nonneg\_left\ hmL\ (Nat.cast\_nonneg\ N)}
\leanline{2}{step\_pos\ :=\ div\_pos\ hγ\ (by\ exact\_mod\_cast\ hN)}
\leanline{2}{step\_beta\_le\ :=\ by}
\leanline{4}{have\ hn\ :\ (N\ :\ ℝ)\ ≠\ 0\ :=\ ne\_of\_gt\ (by\ exact\_mod\_cast\ hN)}
\leanline{4}{have\ heq\ :\ γ\ /\ (N\ :\ ℝ)\ *\ ((N\ :\ ℝ)\ *\ L)\ =\ γ\ *\ L\ :=\ by\ field\_simp\ [hn]}
\leanline{4}{rw\ [heq]}
\leanline{4}{exact\ hγL}
\end{leancode}
\leansource{\texttt{ipla\_\allowbreak{}potential\_\allowbreak{}contDiff} \quad WassersteinMixing/\allowbreak{}IPLACurvature.lean, lines 36–38}
\begin{leancode}
\leanline{0}{theorem\ ipla\_potential\_contDiff\ \{N\ dθ\ dx\ :\ ℕ\}\ \{U\ :\ IPLA.Joint\ dθ\ dx\ →\ ℝ\}}
\leanline{4}{(hU\ :\ ContDiff\ ℝ\ 2\ U)\ :\ ContDiff\ ℝ\ 2\ (@IPLA.potential\ N\ dθ\ dx\ U)}
\end{leancode}
\leanlabel{How the Lean proof follows the argument}
\texttt{ipla\allowbreak{}Parameters} constructs the existing numerical structure with the displayed constants and proves all its fields, including exact cancellation of the positive particle number. \texttt{ipla\_\allowbreak{}potential\_\allowbreak{}cont\allowbreak{}Diff} composes \texttt{Cont\allowbreak{}Diff} with each actual \texttt{IPLA.\allowbreak{}joint\allowbreak{}CLM} and applies finite-sum regularity.
\end{leanbox}


Earlier results used here: \cite[Section~3.1.1]{LangevinLean2026IPLA}; \cite[Section~3.1.3]{LangevinLean2026IPLA}.
\subsection{\texorpdfstring{Transporting the actual gradient through an orthogonal coordinate map}{Transporting the actual gradient through an orthogonal coordinate map}}
\label{report-step-WassersteinMixing-IPLACurvature-lean-2}


\textbf{Claim.} For a linear isometric equivalence \(e:H\to F\) and a differentiable potential \(V:H\to\mathbb R\),

\[\nabla(V\circ e^{-1})(y)=e\,\nabla V(e^{-1}y).\]


\textbf{Proof.} Pair each side with an arbitrary vector \(w\in F\). The derivative chain rule identifies the pairing on the left as \(DV(e^{-1}y)[e^{-1}w]\). The gradient identity turns it into \(\langle\nabla V(e^{-1}y),e^{-1}w\rangle\). Preservation of the inner product by \(e\) gives the pairing with the right side. Equality against every vector proves the vector identity.

\begin{leanbox}{Lean: \texttt{gradient\_\allowbreak{}reindex}}
\leansource{\texttt{gradient\_\allowbreak{}reindex} \quad WassersteinMixing/\allowbreak{}IPLACurvature.lean, lines 41–51}
\begin{leancode}
\leanline{0}{theorem\ gradient\_reindex}
\leanline{4}{\{H\ F\ :\ Type*\}\ [NormedAddCommGroup\ H]\ [InnerProductSpace\ ℝ\ H]\ [CompleteSpace\ H]}
\leanline{4}{[NormedAddCommGroup\ F]\ [InnerProductSpace\ ℝ\ F]\ [CompleteSpace\ F]}
\leanline{4}{(e\ :\ H\ ≃ₗᵢ[ℝ]\ F)\ (V\ :\ H\ →\ ℝ)\ (hV\ :\ Differentiable\ ℝ\ V)\ (x\ :\ F)\ :}
\leanline{4}{gradient\ (V\ ∘\ e.symm)\ x\ =\ e\ (gradient\ V\ (e.symm\ x))}
\end{leancode}
\leanlabel{How the Lean proof follows the argument}
\texttt{gradient\_\allowbreak{}reindex} uses extensionality by the real inner product, the exact \texttt{Has\allowbreak{}FDeriv\allowbreak{}At.\allowbreak{}comp} chain rule, and \texttt{Linear\allowbreak{}Isometry\allowbreak{}Equiv.\allowbreak{}inner\_\allowbreak{}map\_\allowbreak{}eq\_\allowbreak{}flip}. This proves an identity for the actual gradient rather than introducing a transformed gradient as a separate assumption.
\end{leanbox}
\subsection{\texorpdfstring{Preserving the potential contract under reindexing}{Preserving the potential contract under reindexing}}
\label{report-step-WassersteinMixing-IPLACurvature-lean-3}


\textbf{Claim.} A linear isometric coordinate change preserves \(C^2\) regularity, the same strong-convexity constant \(\alpha\), and the same gradient-Lipschitz constant \(\beta\).




\textbf{Proof.} Smoothness follows by composition with the inverse linear isometry. For strong convexity, apply the original convex-combination inequality to \(e^{-1}x,e^{-1}y\). Linearity preserves their convex combination and the isometry preserves \(\|x-y\|\), giving the same quadratic correction. For smoothness of the gradient, use the proved gradient identity and the isometric distance equalities on both sides of the original \(\beta\)-Lipschitz estimate. None of the numerical constants changes.

\begin{leanbox}{Lean: \texttt{potentialAssumptions\_\allowbreak{}reindex}}
\leansource{\texttt{potentialAssumptions\_\allowbreak{}reindex} \quad WassersteinMixing/\allowbreak{}IPLACurvature.lean, lines 55–71}
\begin{leancode}
\leanline{0}{theorem\ potentialAssumptions\_reindex}
\leanline{4}{\{H\ F\ :\ Type*\}\ [NormedAddCommGroup\ H]\ [InnerProductSpace\ ℝ\ H]\ [CompleteSpace\ H]}
\leanline{4}{[NormedAddCommGroup\ F]\ [InnerProductSpace\ ℝ\ F]\ [CompleteSpace\ F]}
\leanline{4}{(P\ :\ Parameters)\ (e\ :\ H\ ≃ₗᵢ[ℝ]\ F)\ (V\ :\ H\ →\ ℝ)}
\leanline{4}{(hV\ :\ PotentialAssumptions\ P\ V)\ :\ PotentialAssumptions\ P\ (V\ ∘\ e.symm)}
\end{leancode}
\leanlabel{How the Lean proof follows the argument}
\texttt{potential\allowbreak{}Assumptions\_\allowbreak{}reindex} assembles all three fields of \texttt{Potential\allowbreak{}Assumptions}. It transports the original \texttt{Strong\allowbreak{}Convex\allowbreak{}On} inequality directly and rewrites the Lipschitz estimate using \texttt{gradient\_\allowbreak{}reindex}. The result is valid for arbitrary complete real inner product spaces related by the given linear isometric equivalence.
\end{leanbox}
\subsection{\texorpdfstring{Discharging the lifted assumptions from the joint potential}{Discharging the lifted assumptions from the joint potential}}
\label{report-step-WassersteinMixing-IPLACurvature-lean-4}


\textbf{Claim.} The joint assumptions \(U\in C^2\), \(m\)-strong convexity and an \(L\)-Lipschitz gradient imply the entire lifted contract with constants \(Nm,NL\). In particular, the literal joint Hessian bounds

\[mI\preccurlyeq\nabla^2U\preccurlyeq LI\]


suffice, with the same constants.




\textbf{Proof.} The exact identity \(\sum_i\|(\delta\theta,\sqrt N\,\delta z_i)\|^2=N\|\delta z\|^2\) makes the sum of the strong-convexity inequalities have constant \(Nm\). The lifted gradient chain rule, followed by Cauchy–Schwarz over the replicated pair coordinates, gives Lipschitz constant \(NL\). The preceding finite-sum argument supplies \(C^2\) regularity. When joint Hessian bounds are the input, the already proved Hessian-to-potential theorem first derives joint strong convexity and gradient Lipschitzness. Thus no semigroup, invariant measure, stationary process or comparison estimate is an input here.

\begin{leanbox}{Lean: \texttt{ipla\_\allowbreak{}potentialAssumptions}, \texttt{ipla\_\allowbreak{}potentialAssumptions\_\allowbreak{}of\_\allowbreak{}hessian}}
\leansource{\texttt{ipla\_\allowbreak{}potentialAssumptions} \quad WassersteinMixing/\allowbreak{}IPLACurvature.lean, lines 75–86}
\begin{leancode}
\leanline{0}{theorem\ ipla\_potentialAssumptions}
\leanline{4}{\{N\ dθ\ dx\ :\ ℕ\}\ \{U\ :\ IPLA.Joint\ dθ\ dx\ →\ ℝ\}\ \{m\ γ\ :\ ℝ\}\ \{L\ :\ ℝ≥0\}}
\leanline{4}{(hN\ :\ 0\ <\ N)\ (hm\ :\ 0\ <\ m)\ (hmL\ :\ m\ ≤\ L)\ (hγ\ :\ 0\ <\ γ)\ (hγL\ :\ γ\ *\ L\ ≤\ 1)}
\leanline{4}{(hU\ :\ ContDiff\ ℝ\ 2\ U)\ (hsc\ :\ StrongConvexOn\ univ\ m\ U)}
\leanline{4}{(hLip\ :\ LipschitzWith\ L\ (gradient\ U))\ :}
\leanline{4}{PotentialAssumptions\ (iplaParameters\ N\ dθ\ dx\ m\ L\ γ\ hN\ hm\ hmL\ hγ\ hγL)}
\leanline{6}{(@IPLA.potential\ N\ dθ\ dx\ U)}
\end{leancode}
\leansource{\texttt{ipla\_\allowbreak{}potentialAssumptions\_\allowbreak{}of\_\allowbreak{}hessian} \quad WassersteinMixing/\allowbreak{}IPLACurvature.lean, lines 90–100}
\begin{leancode}
\leanline{0}{theorem\ ipla\_potentialAssumptions\_of\_hessian}
\leanline{4}{(Q\ :\ Parameters)\ \{N\ dθ\ dx\ :\ ℕ\}\ \{U\ :\ IPLA.Joint\ dθ\ dx\ →\ ℝ\}}
\leanline{4}{(hU\ :\ HessianAssumptions\ Q\ U)\ (hN\ :\ 0\ <\ N)}
\leanline{4}{\{γ\ :\ ℝ\}\ (hγ\ :\ 0\ <\ γ)\ (hγL\ :\ γ\ *\ Q.beta\ ≤\ 1)\ :}
\leanline{4}{PotentialAssumptions}
\leanline{6}{(iplaParameters\ N\ dθ\ dx\ Q.alpha\ Q.beta\ γ\ hN\ Q.alpha\_pos\ Q.alpha\_le\_beta\ hγ\ hγL)}
\leanline{6}{(@IPLA.potential\ N\ dθ\ dx\ U)}
\end{leancode}
\leanlabel{How the Lean proof follows the argument}
\texttt{ipla\_\allowbreak{}potential\allowbreak{}Assumptions} combines \texttt{ipla\_\allowbreak{}potential\_\allowbreak{}cont\allowbreak{}Diff} with \texttt{IPLA.\allowbreak{}potential\_\allowbreak{}strong\allowbreak{}Convex} and \texttt{IPLA.\allowbreak{}potential\_\allowbreak{}gradient\_\allowbreak{}lipschitz}; these earlier IPLA lemmas prove the sharp factor of the particle number. \texttt{ipla\_\allowbreak{}potential\allowbreak{}Assumptions\_\allowbreak{}of\_\allowbreak{}hessian} applies \texttt{Hessian\allowbreak{}Assumptions.\allowbreak{}to\allowbreak{}Potential\allowbreak{}Assumptions} before this lift. Both return the complete potential contract needed by the constructed Wasserstein comparison.
\end{leanbox}


Earlier results used here: \cite[Section~3.1.1]{LangevinLean2026IPLA}; \cite[Section~3.1.3]{LangevinLean2026IPLA}; \cite[Section~3.1.4]{LangevinLean2026IPLA}.
\section{\texorpdfstring{The linear IPLA discretization term and explicit accuracy arithmetic}{The linear IPLA discretization term and explicit accuracy arithmetic}}
\label{proof-module-WassersteinMixing-IPLARates-lean}


Let \(N>0\), \(0<m\le L\), and \(D=d_x+d_\theta/N\). The lifted sharp bias simplifies exactly to \(B=6\gamma\sqrt{(L/m)LD}\). After multiplication by \(\sqrt m\), it is \(6L\gamma\sqrt D\). These scalar identities yield explicit particle, step and iteration choices for normalized parameter accuracy, and an exact small-step comparison with Project 03. The probabilistic bounds and their assumptions are proved in the application modules; they are not supplied or presumed by these arithmetic lemmas.

\subsection{\texorpdfstring{Exact particle scaling and normalization}{Exact particle scaling and normalization}}
\label{report-step-WassersteinMixing-IPLARates-lean-1}


\textbf{Definitions and claims.} Set \(D=d_x+d_\theta/N\ge0\). The new and Project 03 discretization terms are

\[B(\gamma)=6\gamma\sqrt{\frac Lm\,LD},\qquad B_{03}(\gamma)=\frac{33}{20}\frac Lm\sqrt D\sqrt\gamma.\]


The exact lifted expression satisfies

\[6\frac\gamma N\sqrt{\frac{NL}{Nm}(NL)(d_\theta+Nd_x)}=B(\gamma).\]


Moreover, for \(m>0\), \(L,D\ge0\),

\[\sqrt m\,B(\gamma)=6L\gamma\sqrt D,\qquad \sqrt m\sqrt{\frac{2d_\theta}{Nm}}=\sqrt{\frac{2d_\theta}{N}}.\]


\textbf{Proof.} Cancel the positive particle number in the condition number. Since \(d_\theta+Nd_x=ND\), the expression under the remaining square root is \(N^2(L/m)LD\). Extracting its nonnegative factor \(N\) cancels the denominator outside. For the normalized bias, multiply the radicand by \(m\) to obtain \(L^2D\) and extract \(L\ge0\). For concentration, the same multiplication cancels \(m\) under the square root. All divisions by the particle number and curvature use their proved positivity.

\begin{leanbox}{Lean: \texttt{iplaEffectiveDimension}, \texttt{iplaLinearBias}, \texttt{iplaProject03Bias} and 4 more}
\leansource{\texttt{iplaEffectiveDimension} \quad WassersteinMixing/\allowbreak{}IPLARates.lean, lines 17–18}
\begin{leancode}
\leanline{0}{def\ iplaEffectiveDimension\ (N\ dθ\ dx\ :\ ℕ)\ :\ ℝ\ :=}
\leanline{2}{(dx\ :\ ℝ)\ +\ (dθ\ :\ ℝ)\ /\ N}
\end{leancode}
\leansource{\texttt{iplaLinearBias} \quad WassersteinMixing/\allowbreak{}IPLARates.lean, lines 21–22}
\begin{leancode}
\leanline{0}{def\ iplaLinearBias\ (m\ L\ γ\ D\ :\ ℝ)\ :\ ℝ\ :=}
\leanline{2}{6\ *\ γ\ *\ Real.sqrt\ ((L\ /\ m)\ *\ L\ *\ D)}
\end{leancode}
\leansource{\texttt{iplaProject03Bias} \quad WassersteinMixing/\allowbreak{}IPLARates.lean, lines 25–26}
\begin{leancode}
\leanline{0}{def\ iplaProject03Bias\ (m\ L\ γ\ D\ :\ ℝ)\ :\ ℝ\ :=}
\leanline{2}{(33\ /\ 20\ :\ ℝ)\ *\ (L\ /\ m)\ *\ Real.sqrt\ D\ *\ Real.sqrt\ γ}
\end{leancode}
\leansource{\texttt{iplaEffectiveDimension\_\allowbreak{}nonneg} \quad WassersteinMixing/\allowbreak{}IPLARates.lean, lines 44–47}
\begin{leancode}
\leanline{0}{theorem\ iplaEffectiveDimension\_nonneg\ (N\ dθ\ dx\ :\ ℕ)\ :}
\leanline{4}{0\ ≤\ iplaEffectiveDimension\ N\ dθ\ dx}
\end{leancode}
\leansource{\texttt{ipla\_\allowbreak{}lifted\_\allowbreak{}bias\_\allowbreak{}eq} \quad WassersteinMixing/\allowbreak{}IPLARates.lean, lines 50–65}
\begin{leancode}
\leanline{0}{theorem\ ipla\_lifted\_bias\_eq\ (N\ dθ\ dx\ :\ ℕ)\ (hN\ :\ 0\ <\ N)\ (m\ L\ γ\ :\ ℝ)\ :}
\leanline{4}{6\ *\ (γ\ /\ (N\ :\ ℝ))\ *}
\leanline{6}{Real.sqrt\ ((((N\ :\ ℝ)\ *\ L)\ /\ ((N\ :\ ℝ)\ *\ m))\ *}
\leanline{8}{((N\ :\ ℝ)\ *\ L)\ *\ (dθ\ +\ N\ *\ dx\ :\ ℕ))\ =}
\leanline{6}{iplaLinearBias\ m\ L\ γ\ (iplaEffectiveDimension\ N\ dθ\ dx)}
\end{leancode}
\leansource{\texttt{ipla\_\allowbreak{}linear\_\allowbreak{}bias\_\allowbreak{}scale} \quad WassersteinMixing/\allowbreak{}IPLARates.lean, lines 68–75}
\begin{leancode}
\leanline{0}{theorem\ ipla\_linear\_bias\_scale\ (m\ L\ γ\ D\ :\ ℝ)\ (hm\ :\ 0\ <\ m)}
\leanline{4}{(hL\ :\ 0\ ≤\ L)\ (hD\ :\ 0\ ≤\ D)\ :}
\leanline{4}{Real.sqrt\ m\ *\ iplaLinearBias\ m\ L\ γ\ D\ =\ 6\ *\ L\ *\ γ\ *\ Real.sqrt\ D}
\end{leancode}
\leansource{\texttt{ipla\_\allowbreak{}concentration\_\allowbreak{}scale} \quad WassersteinMixing/\allowbreak{}IPLARates.lean, lines 78–84}
\begin{leancode}
\leanline{0}{theorem\ ipla\_concentration\_scale\ (N\ dθ\ :\ ℕ)\ (hN\ :\ 0\ <\ N)\ (m\ :\ ℝ)\ (hm\ :\ 0\ <\ m)\ :}
\leanline{4}{Real.sqrt\ m\ *\ Real.sqrt\ (2\ *\ (dθ\ :\ ℝ)\ /\ ((N\ :\ ℝ)\ *\ m))\ =}
\leanline{6}{Real.sqrt\ (2\ *\ (dθ\ :\ ℝ)\ /\ N)}
\end{leancode}
\leanlabel{How the Lean proof follows the argument}
The first three definitions retain the literal dimension and the two constants. \texttt{ipla\allowbreak{}Effective\allowbreak{}Dimension\_\allowbreak{}nonneg} proves the dimension sign. \texttt{ipla\_\allowbreak{}lifted\_\allowbreak{}bias\_\allowbreak{}eq} factors the radicand as the square of the particle number times the reduced radicand, then uses the square-root product identity. \texttt{ipla\_\allowbreak{}linear\_\allowbreak{}bias\_\allowbreak{}scale} invokes the proved curvature scaling identity. \texttt{ipla\_\allowbreak{}concentration\_\allowbreak{}scale} combines positive square roots before exact cancellation.
\end{leanbox}
\subsection{\texorpdfstring{A positive step with an exact reciprocal iteration factor}{A positive step with an exact reciprocal iteration factor}}
\label{report-step-WassersteinMixing-IPLARates-lean-2}


\textbf{Definitions.} For \(L>0\) and a desired normalized accuracy \(\varepsilon>0\), set

\[M=\max\!\left\{1,\frac{18\sqrt D}{\varepsilon}\right\},\qquad \gamma_\varepsilon=\frac1{LM},\qquad T_\varepsilon=\frac Lm\,M.\]


\textbf{Claim.} The step is positive, \(\gamma_\varepsilon L\le1\), and \((m\gamma_\varepsilon)T_\varepsilon=1\).




\textbf{Proof.} The maximum is at least one and hence positive, even when \(D=0\). Thus both denominator factors are positive. Direct cancellation gives \(\gamma_\varepsilon L=1/M\le1\). Multiplying the step by its displayed iteration factor cancels the positive constants \(m,L,M\) exactly.

\begin{leanbox}{Lean: \texttt{iplaAccuracyStep}, \texttt{iplaAccuracyPrefactor}, \texttt{iplaAccuracyStep\_\allowbreak{}pos} and 2 more}
\leansource{\texttt{iplaAccuracyStep} \quad WassersteinMixing/\allowbreak{}IPLARates.lean, lines 29–30}
\begin{leancode}
\leanline{0}{def\ iplaAccuracyStep\ (L\ D\ ε\ :\ ℝ)\ :\ ℝ\ :=}
\leanline{2}{1\ /\ (L\ *\ max\ 1\ (18\ *\ Real.sqrt\ D\ /\ ε))}
\end{leancode}
\leansource{\texttt{iplaAccuracyPrefactor} \quad WassersteinMixing/\allowbreak{}IPLARates.lean, lines 33–34}
\begin{leancode}
\leanline{0}{def\ iplaAccuracyPrefactor\ (m\ L\ D\ ε\ :\ ℝ)\ :\ ℝ\ :=}
\leanline{2}{(L\ /\ m)\ *\ max\ 1\ (18\ *\ Real.sqrt\ D\ /\ ε)}
\end{leancode}
\leansource{\texttt{iplaAccuracyStep\_\allowbreak{}pos} \quad WassersteinMixing/\allowbreak{}IPLARates.lean, lines 86–89}
\begin{leancode}
\leanline{0}{theorem\ iplaAccuracyStep\_pos\ (L\ D\ ε\ :\ ℝ)\ (hL\ :\ 0\ <\ L)\ :}
\leanline{4}{0\ <\ iplaAccuracyStep\ L\ D\ ε}
\end{leancode}
\leansource{\texttt{iplaAccuracyStep\_\allowbreak{}mul\_\allowbreak{}L\_\allowbreak{}le} \quad WassersteinMixing/\allowbreak{}IPLARates.lean, lines 91–99}
\begin{leancode}
\leanline{0}{theorem\ iplaAccuracyStep\_mul\_L\_le\ (L\ D\ ε\ :\ ℝ)\ (hL\ :\ 0\ <\ L)\ :}
\leanline{4}{iplaAccuracyStep\ L\ D\ ε\ *\ L\ ≤\ 1}
\end{leancode}
\leansource{\texttt{ipla\_\allowbreak{}step\_\allowbreak{}mul\_\allowbreak{}prefactor} \quad WassersteinMixing/\allowbreak{}IPLARates.lean, lines 102–107}
\begin{leancode}
\leanline{0}{theorem\ ipla\_step\_mul\_prefactor\ (m\ L\ D\ ε\ :\ ℝ)\ (hm\ :\ 0\ <\ m)\ (hL\ :\ 0\ <\ L)\ :}
\leanline{4}{(m\ *\ iplaAccuracyStep\ L\ D\ ε)\ *\ iplaAccuracyPrefactor\ m\ L\ D\ ε\ =\ 1}
\end{leancode}
\leanlabel{How the Lean proof follows the argument}
\texttt{ipla\allowbreak{}Accuracy\allowbreak{}Step} and \texttt{ipla\allowbreak{}Accuracy\allowbreak{}Prefactor} use a maximum rather than dividing by the square root of the dimension; therefore zero dimension is included. \texttt{ipla\allowbreak{}Accuracy\allowbreak{}Step\_\allowbreak{}pos} proves positivity of the denominator. \texttt{ipla\allowbreak{}Accuracy\allowbreak{}Step\_\allowbreak{}mul\_\allowbreak{}L\_\allowbreak{}le} rewrites the scaled step as the reciprocal of that maximum. \texttt{ipla\_\allowbreak{}step\_\allowbreak{}mul\_\allowbreak{}prefactor} establishes exact equality by field cancellation using the positive factors.
\end{leanbox}
\subsection{\texorpdfstring{Three equal budgets give invariant-based normalized accuracy}{Three equal budgets give invariant-based normalized accuracy}}
\label{report-step-WassersteinMixing-IPLARates-lean-3}


\textbf{Claim.} Let \(R\ge0\), \(\varepsilon>0\), \(N>0\), \(0<m\le L\), and \(D\ge0\). With the preceding step, assume

\[N\varepsilon^2\ge18d_\theta,\qquad k\ge\frac Lm\max\!\left\{1,\frac{18\sqrt D}{\varepsilon}\right\}\log\!\left(1+\frac{3\sqrt m\,R}{\varepsilon}\right).\]


Then each of the following three terms is at most \(\varepsilon/3\):

\[\sqrt{\frac{2d_\theta}{N}},\qquad \sqrt m(1-m\gamma_\varepsilon)^kR,\qquad \sqrt m\,B(\gamma_\varepsilon).\]


Consequently their sum is at most \(\varepsilon\).




\textbf{Proof.} Squaring the nonnegative concentration term reduces its budget exactly to the stated particle inequality. The maximum in the step satisfies \(18\sqrt D\le\varepsilon M\), so the normalized bias \(6\sqrt D/M\) is at most \(\varepsilon/3\). Put \(t=m\gamma_\varepsilon\in(0,1]\). Its exact reciprocal factor turns the iteration threshold into \(tk\ge\log(1+3\sqrt mR/\varepsilon)\). The established geometric-transient estimate then bounds \((1-t)^k\sqrt mR\) by \(\varepsilon/3\). Adding the three inequalities gives the conclusion.

\begin{leanbox}{Lean: \texttt{ipla\_\allowbreak{}step\_\allowbreak{}bias\_\allowbreak{}budget}, \texttt{ipla\_\allowbreak{}concentration\_\allowbreak{}budget}, \texttt{ipla\_\allowbreak{}transient\_\allowbreak{}budget} and 1 more}
\leansource{\texttt{ipla\_\allowbreak{}step\_\allowbreak{}bias\_\allowbreak{}budget} \quad WassersteinMixing/\allowbreak{}IPLARates.lean, lines 110–125}
\begin{leancode}
\leanline{0}{theorem\ ipla\_step\_bias\_budget\ (m\ L\ D\ ε\ :\ ℝ)\ (hm\ :\ 0\ <\ m)\ (hL\ :\ 0\ <\ L)}
\leanline{4}{(hD\ :\ 0\ ≤\ D)\ (hε\ :\ 0\ <\ ε)\ :}
\leanline{4}{Real.sqrt\ m\ *\ iplaLinearBias\ m\ L\ (iplaAccuracyStep\ L\ D\ ε)\ D\ ≤\ ε\ /\ 3}
\end{leancode}
\leansource{\texttt{ipla\_\allowbreak{}concentration\_\allowbreak{}budget} \quad WassersteinMixing/\allowbreak{}IPLARates.lean, lines 128–135}
\begin{leancode}
\leanline{0}{theorem\ ipla\_concentration\_budget\ (N\ dθ\ :\ ℕ)\ (hN\ :\ 0\ <\ N)\ (ε\ :\ ℝ)\ (hε\ :\ 0\ <\ ε)}
\leanline{4}{(hparticles\ :\ 18\ *\ (dθ\ :\ ℝ)\ ≤\ (N\ :\ ℝ)\ *\ ε\ \textasciicircum{}\ 2)\ :}
\leanline{4}{Real.sqrt\ (2\ *\ (dθ\ :\ ℝ)\ /\ (N\ :\ ℝ))\ ≤\ ε\ /\ 3}
\end{leancode}
\leansource{\texttt{ipla\_\allowbreak{}transient\_\allowbreak{}budget} \quad WassersteinMixing/\allowbreak{}IPLARates.lean, lines 138–159}
\begin{leancode}
\leanline{0}{theorem\ ipla\_transient\_budget\ (m\ L\ D\ R\ ε\ :\ ℝ)\ (hm\ :\ 0\ <\ m)}
\leanline{4}{(hmL\ :\ m\ ≤\ L)\ (hR\ :\ 0\ ≤\ R)\ (hε\ :\ 0\ <\ ε)\ (k\ :\ ℕ)}
\leanline{4}{(hk\ :\ iplaAccuracyPrefactor\ m\ L\ D\ ε\ *}
\leanline{6}{Real.log\ (1\ +\ 3\ *\ Real.sqrt\ m\ *\ R\ /\ ε)\ ≤\ (k\ :\ ℝ))\ :}
\leanline{4}{Real.sqrt\ m\ *\ ((1\ -\ m\ *\ iplaAccuracyStep\ L\ D\ ε)\ \textasciicircum{}\ k\ *\ R)\ ≤\ ε\ /\ 3}
\end{leancode}
\leansource{\texttt{ipla\_\allowbreak{}accuracy\_\allowbreak{}arithmetic} \quad WassersteinMixing/\allowbreak{}IPLARates.lean, lines 162–173}
\begin{leancode}
\leanline{0}{theorem\ ipla\_accuracy\_arithmetic\ (N\ dθ\ :\ ℕ)\ (hN\ :\ 0\ <\ N)\ (m\ L\ D\ R\ ε\ :\ ℝ)}
\leanline{4}{(hm\ :\ 0\ <\ m)\ (hmL\ :\ m\ ≤\ L)\ (hD\ :\ 0\ ≤\ D)\ (hR\ :\ 0\ ≤\ R)\ (hε\ :\ 0\ <\ ε)}
\leanline{4}{(hparticles\ :\ 18\ *\ (dθ\ :\ ℝ)\ ≤\ (N\ :\ ℝ)\ *\ ε\ \textasciicircum{}\ 2)\ (k\ :\ ℕ)}
\leanline{4}{(hk\ :\ iplaAccuracyPrefactor\ m\ L\ D\ ε\ *}
\leanline{6}{Real.log\ (1\ +\ 3\ *\ Real.sqrt\ m\ *\ R\ /\ ε)\ ≤\ (k\ :\ ℝ))\ :}
\leanline{4}{Real.sqrt\ (2\ *\ (dθ\ :\ ℝ)\ /\ (N\ :\ ℝ))\ +}
\leanline{6}{Real.sqrt\ m\ *\ ((1\ -\ m\ *\ iplaAccuracyStep\ L\ D\ ε)\ \textasciicircum{}\ k\ *\ R\ +}
\leanline{8}{iplaLinearBias\ m\ L\ (iplaAccuracyStep\ L\ D\ ε)\ D)\ ≤\ ε}
\end{leancode}
\leanlabel{How the Lean proof follows the argument}
\texttt{ipla\_\allowbreak{}step\_\allowbreak{}bias\_\allowbreak{}budget} rewrites the normalized bias using \texttt{ipla\_\allowbreak{}linear\_\allowbreak{}bias\_\allowbreak{}scale} and the maximum inequality. \texttt{ipla\_\allowbreak{}concentration\_\allowbreak{}budget} compares squares only after proving both sides nonnegative. \texttt{ipla\_\allowbreak{}transient\_\allowbreak{}budget} applies \texttt{geometric\_\allowbreak{}transient\_\allowbreak{}le} with the reciprocal identity and the proved step range. \texttt{ipla\_\allowbreak{}accuracy\_\allowbreak{}arithmetic} adds these three derived budgets. The application later substitutes the actual nonnegative initial transport distance for the scalar radius.
\end{leanbox}
\subsection{\texorpdfstring{An optional conservative accuracy budget for a doubled bias}{An optional conservative accuracy budget for a doubled bias}}
\label{report-step-WassersteinMixing-IPLARates-lean-4}


\textbf{Definitions and claim.} For the target-based parameter bound, reserve room for the doubled discretization term by choosing

\[\widetilde\gamma_\varepsilon=\frac1{L\max\{1,36\sqrt D/\varepsilon\}},\qquad \widetilde T_\varepsilon=\frac Lm\max\{1,36\sqrt D/\varepsilon\}.\]


If \(N\varepsilon^2\ge18d_\theta\) and \(k\ge\widetilde T_\varepsilon\log(1+3\sqrt mR/\varepsilon)\), then

\[\sqrt{\frac{2d_\theta}{N}}+\sqrt m\left((1-m\widetilde\gamma_\varepsilon)^kR+2B(\widetilde\gamma_\varepsilon)\right)\le\varepsilon.\]


\textbf{Proof.} The new step and prefactor are the preceding definitions with accuracy \(\varepsilon/2\); simplifying their formulas gives the coefficient \(36\). The earlier bias budget at that half accuracy gives \(\sqrt mB\le\varepsilon/6\), so twice the bias costs at most \(\varepsilon/3\). The same reciprocal identity proves the transient budget \(\varepsilon/3\) at the displayed logarithmic threshold. The unchanged concentration budget is at most \(\varepsilon/3\). Add the three terms. This is a valid auxiliary budget for a doubled-bias estimate. The primary application proves a sharper exact-Gibbs initialization bound with a single bias and therefore uses the coefficient \(18\), rather than requiring this conservative coefficient \(36\).

\begin{leanbox}{Lean: \texttt{iplaTargetAccuracyStep}, \texttt{iplaTargetAccuracyPrefactor}, \texttt{iplaTargetAccuracyStep\_\allowbreak{}eq} and 3 more}
\leansource{\texttt{iplaTargetAccuracyStep} \quad WassersteinMixing/\allowbreak{}IPLARates.lean, lines 37–38}
\begin{leancode}
\leanline{0}{def\ iplaTargetAccuracyStep\ (L\ D\ ε\ :\ ℝ)\ :\ ℝ\ :=}
\leanline{2}{iplaAccuracyStep\ L\ D\ (ε\ /\ 2)}
\end{leancode}
\leansource{\texttt{iplaTargetAccuracyPrefactor} \quad WassersteinMixing/\allowbreak{}IPLARates.lean, lines 41–42}
\begin{leancode}
\leanline{0}{def\ iplaTargetAccuracyPrefactor\ (m\ L\ D\ ε\ :\ ℝ)\ :\ ℝ\ :=}
\leanline{2}{iplaAccuracyPrefactor\ m\ L\ D\ (ε\ /\ 2)}
\end{leancode}
\leansource{\texttt{iplaTargetAccuracyStep\_\allowbreak{}eq} \quad WassersteinMixing/\allowbreak{}IPLARates.lean, lines 176–182}
\begin{leancode}
\leanline{0}{theorem\ iplaTargetAccuracyStep\_eq\ (L\ D\ ε\ :\ ℝ)\ :}
\leanline{4}{iplaTargetAccuracyStep\ L\ D\ ε\ =}
\leanline{6}{1\ /\ (L\ *\ max\ 1\ (36\ *\ Real.sqrt\ D\ /\ ε))}
\end{leancode}
\leansource{\texttt{iplaTargetAccuracyPrefactor\_\allowbreak{}eq} \quad WassersteinMixing/\allowbreak{}IPLARates.lean, lines 185–191}
\begin{leancode}
\leanline{0}{theorem\ iplaTargetAccuracyPrefactor\_eq\ (m\ L\ D\ ε\ :\ ℝ)\ :}
\leanline{4}{iplaTargetAccuracyPrefactor\ m\ L\ D\ ε\ =}
\leanline{6}{(L\ /\ m)\ *\ max\ 1\ (36\ *\ Real.sqrt\ D\ /\ ε)}
\end{leancode}
\leansource{\texttt{ipla\_\allowbreak{}target\_\allowbreak{}transient\_\allowbreak{}budget} \quad WassersteinMixing/\allowbreak{}IPLARates.lean, lines 194–217}
\begin{leancode}
\leanline{0}{theorem\ ipla\_target\_transient\_budget\ (m\ L\ D\ R\ ε\ :\ ℝ)\ (hm\ :\ 0\ <\ m)}
\leanline{4}{(hmL\ :\ m\ ≤\ L)\ (hR\ :\ 0\ ≤\ R)\ (hε\ :\ 0\ <\ ε)\ (k\ :\ ℕ)}
\leanline{4}{(hk\ :\ iplaTargetAccuracyPrefactor\ m\ L\ D\ ε\ *}
\leanline{6}{Real.log\ (1\ +\ 3\ *\ Real.sqrt\ m\ *\ R\ /\ ε)\ ≤\ (k\ :\ ℝ))\ :}
\leanline{4}{Real.sqrt\ m\ *\ ((1\ -\ m\ *\ iplaTargetAccuracyStep\ L\ D\ ε)\ \textasciicircum{}\ k\ *\ R)\ ≤\ ε\ /\ 3}
\end{leancode}
\leansource{\texttt{ipla\_\allowbreak{}accuracy\_\allowbreak{}target\_\allowbreak{}arithmetic} \quad WassersteinMixing/\allowbreak{}IPLARates.lean, lines 220–233}
\begin{leancode}
\leanline{0}{theorem\ ipla\_accuracy\_target\_arithmetic\ (N\ dθ\ :\ ℕ)\ (hN\ :\ 0\ <\ N)\ (m\ L\ D\ R\ ε\ :\ ℝ)}
\leanline{4}{(hm\ :\ 0\ <\ m)\ (hmL\ :\ m\ ≤\ L)\ (hD\ :\ 0\ ≤\ D)\ (hR\ :\ 0\ ≤\ R)\ (hε\ :\ 0\ <\ ε)}
\leanline{4}{(hparticles\ :\ 18\ *\ (dθ\ :\ ℝ)\ ≤\ (N\ :\ ℝ)\ *\ ε\ \textasciicircum{}\ 2)\ (k\ :\ ℕ)}
\leanline{4}{(hk\ :\ iplaTargetAccuracyPrefactor\ m\ L\ D\ ε\ *}
\leanline{6}{Real.log\ (1\ +\ 3\ *\ Real.sqrt\ m\ *\ R\ /\ ε)\ ≤\ (k\ :\ ℝ))\ :}
\leanline{4}{Real.sqrt\ (2\ *\ (dθ\ :\ ℝ)\ /\ (N\ :\ ℝ))\ +}
\leanline{6}{Real.sqrt\ m\ *\ ((1\ -\ m\ *\ iplaTargetAccuracyStep\ L\ D\ ε)\ \textasciicircum{}\ k\ *\ R\ +}
\leanline{8}{2\ *\ iplaLinearBias\ m\ L\ (iplaTargetAccuracyStep\ L\ D\ ε)\ D)\ ≤\ ε}
\end{leancode}
\leanlabel{How the Lean proof follows the argument}
\texttt{ipla\allowbreak{}Target\allowbreak{}Accuracy\allowbreak{}Step} and \texttt{ipla\allowbreak{}Target\allowbreak{}Accuracy\allowbreak{}Prefactor} reuse the base definitions at half accuracy, while the two equality lemmas expose the exact coefficient. \texttt{ipla\_\allowbreak{}target\_\allowbreak{}transient\_\allowbreak{}budget} proves the logarithmic bound for that step without changing the target error budget. \texttt{ipla\_\allowbreak{}accuracy\_\allowbreak{}target\_\allowbreak{}arithmetic} combines it with the original concentration budget and twice the half-accuracy bias estimate.
\end{leanbox}
\subsection{\texorpdfstring{An exact regime where the discretization estimate improves Project 03}{An exact regime where the discretization estimate improves Project 03}}
\label{report-step-WassersteinMixing-IPLARates-lean-5}


\textbf{Claim.} For \(m>0\), \(L,D,\gamma\ge0\),

\[m\gamma\le\frac{121}{1600}\ \Longrightarrow\ B(\gamma)\le B_{03}(\gamma),\]
\[\quad m\gamma\le\frac{121}{6400}\ \Longrightarrow\ 2B(\gamma)\le B_{03}(\gamma).\]


\textbf{Proof.} Factor the new term as

\[B(\gamma)=\bigl(6\gamma\sqrt m\bigr)\left(\frac Lm\sqrt D\right).\]


The second factor is nonnegative and is shared with the Project 03 expression. Comparing the first factors reduces to \(6\gamma\sqrt m\le(33/20)\sqrt\gamma\); squaring the nonnegative sides and using the displayed bound on \(m\gamma\) proves it. Replacing \(6\) by \(12\) gives the second threshold and the doubled-bias result. No division by \(\gamma\) or \(D\) is required, so their zero cases are included.




The primary application attains a single bias with either discretized-invariant or exact-Gibbs initialization, so its sharp complete bound uses the first threshold. The doubled statement records a valid conservative alternative. Thus the improved discretization order is linear in the step, within the explicitly stated regularity and step regime, rather than a claim that every Project 03 hypothesis or every allowed step is covered.

\begin{leanbox}{Lean: \texttt{ipla\_\allowbreak{}bias\_\allowbreak{}coefficient\_\allowbreak{}le\_\allowbreak{}old}, \texttt{ipla\_\allowbreak{}doubled\_\allowbreak{}bias\_\allowbreak{}coefficient\_\allowbreak{}le\_\allowbreak{}old}, \texttt{ipla\_\allowbreak{}linear\_\allowbreak{}bias\_\allowbreak{}factorization} and 2 more}
\leansource{\texttt{ipla\_\allowbreak{}bias\_\allowbreak{}coefficient\_\allowbreak{}le\_\allowbreak{}old} \quad WassersteinMixing/\allowbreak{}IPLARates.lean, lines 236–242}
\begin{leancode}
\leanline{0}{theorem\ ipla\_bias\_coefficient\_le\_old\ (m\ γ\ :\ ℝ)\ (hm\ :\ 0\ <\ m)\ (hγ\ :\ 0\ ≤\ γ)}
\leanline{4}{(hsmall\ :\ m\ *\ γ\ ≤\ (121\ /\ 1600\ :\ ℝ))\ :}
\leanline{4}{6\ *\ γ\ *\ Real.sqrt\ m\ ≤\ (33\ /\ 20\ :\ ℝ)\ *\ Real.sqrt\ γ}
\end{leancode}
\leansource{\texttt{ipla\_\allowbreak{}doubled\_\allowbreak{}bias\_\allowbreak{}coefficient\_\allowbreak{}le\_\allowbreak{}old} \quad WassersteinMixing/\allowbreak{}IPLARates.lean, lines 245–251}
\begin{leancode}
\leanline{0}{theorem\ ipla\_doubled\_bias\_coefficient\_le\_old\ (m\ γ\ :\ ℝ)\ (hm\ :\ 0\ <\ m)\ (hγ\ :\ 0\ ≤\ γ)}
\leanline{4}{(hsmall\ :\ m\ *\ γ\ ≤\ (121\ /\ 6400\ :\ ℝ))\ :}
\leanline{4}{12\ *\ γ\ *\ Real.sqrt\ m\ ≤\ (33\ /\ 20\ :\ ℝ)\ *\ Real.sqrt\ γ}
\end{leancode}
\leansource{\texttt{ipla\_\allowbreak{}linear\_\allowbreak{}bias\_\allowbreak{}factorization} \quad WassersteinMixing/\allowbreak{}IPLARates.lean, lines 254–269}
\begin{leancode}
\leanline{0}{theorem\ ipla\_linear\_bias\_factorization\ (m\ L\ γ\ D\ :\ ℝ)\ (hm\ :\ 0\ <\ m)}
\leanline{4}{(hL\ :\ 0\ ≤\ L)\ (hD\ :\ 0\ ≤\ D)\ :}
\leanline{4}{iplaLinearBias\ m\ L\ γ\ D\ =}
\leanline{6}{(6\ *\ γ\ *\ Real.sqrt\ m)\ *\ ((L\ /\ m)\ *\ Real.sqrt\ D)}
\end{leancode}
\leansource{\texttt{ipla\_\allowbreak{}linear\_\allowbreak{}bias\_\allowbreak{}le\_\allowbreak{}project03} \quad WassersteinMixing/\allowbreak{}IPLARates.lean, lines 273–282}
\begin{leancode}
\leanline{0}{theorem\ ipla\_linear\_bias\_le\_project03\ (m\ L\ γ\ D\ :\ ℝ)\ (hm\ :\ 0\ <\ m)}
\leanline{4}{(hL\ :\ 0\ ≤\ L)\ (hγ\ :\ 0\ ≤\ γ)\ (hD\ :\ 0\ ≤\ D)}
\leanline{4}{(hsmall\ :\ m\ *\ γ\ ≤\ (121\ /\ 1600\ :\ ℝ))\ :}
\leanline{4}{iplaLinearBias\ m\ L\ γ\ D\ ≤\ iplaProject03Bias\ m\ L\ γ\ D}
\end{leancode}
\leansource{\texttt{ipla\_\allowbreak{}doubled\_\allowbreak{}bias\_\allowbreak{}le\_\allowbreak{}project03} \quad WassersteinMixing/\allowbreak{}IPLARates.lean, lines 287–296}
\begin{leancode}
\leanline{0}{theorem\ ipla\_doubled\_bias\_le\_project03\ (m\ L\ γ\ D\ :\ ℝ)\ (hm\ :\ 0\ <\ m)}
\leanline{4}{(hL\ :\ 0\ ≤\ L)\ (hγ\ :\ 0\ ≤\ γ)\ (hD\ :\ 0\ ≤\ D)}
\leanline{4}{(hsmall\ :\ m\ *\ γ\ ≤\ (121\ /\ 6400\ :\ ℝ))\ :}
\leanline{4}{2\ *\ iplaLinearBias\ m\ L\ γ\ D\ ≤\ iplaProject03Bias\ m\ L\ γ\ D}
\end{leancode}
\leanlabel{How the Lean proof follows the argument}
\texttt{ipla\_\allowbreak{}bias\_\allowbreak{}coefficient\_\allowbreak{}le\_\allowbreak{}old} and \texttt{ipla\_\allowbreak{}doubled\_\allowbreak{}bias\_\allowbreak{}coefficient\_\allowbreak{}le\_\allowbreak{}old} prove the scalar comparisons by nonnegative-square inequalities. \texttt{ipla\_\allowbreak{}linear\_\allowbreak{}bias\_\allowbreak{}factorization} derives the common-factor identity from the curvature square-root scaling. The final two lemmas multiply the appropriate scalar inequality by the nonnegative common factor and rewrite the exact Project 03 expression.
\end{leanbox}
\section{\texorpdfstring{The sharp actual-law Wasserstein bound in the scaled IPLA coordinates}{The sharp actual-law Wasserstein bound in the scaled IPLA coordinates}}
\label{proof-module-WassersteinMixing-IPLABridge-lean}


The scaled IPLA state is a finite Euclidean space indexed by a parameter block and a particle/\allowbreak{}latent block. A proved linear isometry identifies it with the coordinate convention of the sharp Wasserstein theorem. Pulling the constructed invariant law back through that isometry gives a genuine invariant probability measure with a finite second moment. The actual Gaussian transition, Gibbs target and transport distance all respect this change of coordinates, so the sharp bias and convergence estimates apply to the actual lifted particle law.



Earlier results used here: \cite[Section~3.16]{LangevinLean2026IPLA}.
\subsection{\texorpdfstring{Constructing the invariant law in the original coordinate type}{Constructing the invariant law in the original coordinate type}}
\label{report-step-WassersteinMixing-IPLABridge-lean-1}


\textbf{Construction.} Let \(V\) satisfy the potential contract with parameters \(\alpha,\beta,h,d\), and let \(e:H\to\mathbb R^d\) be a linear isometric equivalence. For \(W=V\circ e^{-1}\), let \(\widehat\pi_W\) be the invariant law already constructed from the actual Gaussian ULA transition. Define

\[\widehat\pi_V=(e^{-1})_\#\widehat\pi_W.\]


\textbf{Claim.} It is a probability law with finite second moment, \(e_\#\widehat\pi_V=\widehat\pi_W\), and \(K_{V,h}\widehat\pi_V=\widehat\pi_V\).




\textbf{Proof.} The measurable inverse coordinate map preserves total mass, and its norm preservation transfers the constructed finite second moment. Composing its pushforward with the forward map gives the exact recovery identity. Gaussian-transition covariance then gives \(e_\#K_{V,h}\widehat\pi_V=K_{W,h}\widehat\pi_W=\widehat\pi_W=e_\#\widehat\pi_V\). Injectivity of pushforward through the measurable equivalence proves invariance in the original coordinates. Thus invariance is derived for the same transition used by the particle algorithm.

\begin{leanbox}{Lean: \texttt{coordinateInvariantLaw}, \texttt{coordinateInvariantLaw\_\allowbreak{}map}, \texttt{coordinateInvariantLaw\_\allowbreak{}isProbability} and 2 more}
\leansource{\texttt{coordinateInvariantLaw} \quad WassersteinMixing/\allowbreak{}IPLABridge.lean, lines 25–28}
\begin{leancode}
\leanline{0}{def\ coordinateInvariantLaw\ (P\ :\ Parameters)}
\leanline{4}{(e\ :\ H\ ≃ₗᵢ[ℝ]\ EuclideanSpace\ ℝ\ (Fin\ P.dimension))}
\leanline{4}{(V\ :\ H\ →\ ℝ)\ (hV\ :\ PotentialAssumptions\ P\ V)\ :\ Measure\ H\ :=}
\leanline{2}{(invariantLaw\ P\ (V\ ∘\ e.symm)\ (potentialAssumptions\_reindex\ P\ e\ V\ hV)).map\ e.symm}
\end{leancode}
\leansource{\texttt{coordinateInvariantLaw\_\allowbreak{}map} \quad WassersteinMixing/\allowbreak{}IPLABridge.lean, lines 33–39}
\begin{leancode}
\leanline{0}{theorem\ coordinateInvariantLaw\_map\ (P\ :\ Parameters)}
\leanline{4}{(e\ :\ H\ ≃ₗᵢ[ℝ]\ EuclideanSpace\ ℝ\ (Fin\ P.dimension))}
\leanline{4}{(V\ :\ H\ →\ ℝ)\ (hV\ :\ PotentialAssumptions\ P\ V)\ :}
\leanline{4}{(coordinateInvariantLaw\ P\ e\ V\ hV).map\ e\ =}
\leanline{6}{invariantLaw\ P\ (V\ ∘\ e.symm)\ (potentialAssumptions\_reindex\ P\ e\ V\ hV)}
\end{leancode}
\leansource{\texttt{coordinateInvariantLaw\_\allowbreak{}isProbability} \quad WassersteinMixing/\allowbreak{}IPLABridge.lean, lines 43–48}
\begin{leancode}
\leanline{0}{theorem\ coordinateInvariantLaw\_isProbability\ (P\ :\ Parameters)}
\leanline{4}{(e\ :\ H\ ≃ₗᵢ[ℝ]\ EuclideanSpace\ ℝ\ (Fin\ P.dimension))}
\leanline{4}{(V\ :\ H\ →\ ℝ)\ (hV\ :\ PotentialAssumptions\ P\ V)\ :}
\leanline{4}{IsProbabilityMeasure\ (coordinateInvariantLaw\ P\ e\ V\ hV)}
\end{leancode}
\leansource{\texttt{coordinateInvariantLaw\_\allowbreak{}memLp} \quad WassersteinMixing/\allowbreak{}IPLABridge.lean, lines 51–56}
\begin{leancode}
\leanline{0}{theorem\ coordinateInvariantLaw\_memLp\ (P\ :\ Parameters)}
\leanline{4}{(e\ :\ H\ ≃ₗᵢ[ℝ]\ EuclideanSpace\ ℝ\ (Fin\ P.dimension))}
\leanline{4}{(V\ :\ H\ →\ ℝ)\ (hV\ :\ PotentialAssumptions\ P\ V)\ :}
\leanline{4}{MemLp\ (id\ :\ H\ →\ H)\ 2\ (coordinateInvariantLaw\ P\ e\ V\ hV)}
\end{leancode}
\leansource{\texttt{coordinateInvariantLaw\_\allowbreak{}invariant} \quad WassersteinMixing/\allowbreak{}IPLABridge.lean, lines 59–70}
\begin{leancode}
\leanline{0}{theorem\ coordinateInvariantLaw\_invariant\ (P\ :\ Parameters)}
\leanline{4}{(e\ :\ H\ ≃ₗᵢ[ℝ]\ EuclideanSpace\ ℝ\ (Fin\ P.dimension))}
\leanline{4}{(V\ :\ H\ →\ ℝ)\ (hV\ :\ PotentialAssumptions\ P\ V)\ :}
\leanline{4}{ULA.ulaLawStep\ V\ P.step\ (coordinateInvariantLaw\ P\ e\ V\ hV)\ =}
\leanline{6}{coordinateInvariantLaw\ P\ e\ V\ hV}
\end{leancode}
\leanlabel{How the Lean proof follows the argument}
\texttt{coordinate\allowbreak{}Invariant\allowbreak{}Law} maps the existing \texttt{invariant\allowbreak{}Law} back through the inverse isometry using \texttt{potential\allowbreak{}Assumptions\_\allowbreak{}reindex}. \texttt{coordinate\allowbreak{}Invariant\allowbreak{}Law\_\allowbreak{}map} proves recovery by composition of measure maps. The probability and moment lemmas transfer the corresponding established properties of \texttt{invariant\allowbreak{}Law}. \texttt{coordinate\allowbreak{}Invariant\allowbreak{}Law\_\allowbreak{}invariant} applies \texttt{ula\allowbreak{}Law\allowbreak{}Step\_\allowbreak{}map\_\allowbreak{}isometry} and the original invariant-law theorem, then uses injectivity of the mapped measure.
\end{leanbox}


Earlier results used here: \cite[Section~3.36.7]{LangevinLean2026ULASC}.
\subsection{\texorpdfstring{Transferring the sharp bias and finite-time comparison}{Transferring the sharp bias and finite-time comparison}}
\label{report-step-WassersteinMixing-IPLABridge-lean-2}


\textbf{Claim.} Set \(\pi_V=Z_V^{-1}e^{-V}\), \(q=1-\alpha h\) and \(B=6h\sqrt{(\beta/\alpha)\beta d}\). Then

\[W_2(\pi_V,\widehat\pi_V)\le B,\qquad W_2(K_{V,h}^{\,k}\mu,\pi_V)\le q^kW_2(\mu,\widehat\pi_V)+B,\]


for probability initializations with finite second moment.




\textbf{Proof.} Push both measures through \(e\). The exact Gibbs pushforward, invariant-law recovery and transport-isometry identities reduce the first inequality to the already constructed sharp bias theorem for \(W\). For the second inequality, the Gaussian-iterate covariance also identifies the pushed-forward iterate with the actual ULA iterate for \(W\). Finite second moments pass to the mapped initialization. The constructed comparison theorem in the Euclidean coordinates then gives precisely the displayed bound, with all constants unchanged.

\begin{leanbox}{Lean: \texttt{coordinate\_\allowbreak{}bias\_\allowbreak{}bound}, \texttt{coordinate\_\allowbreak{}convergence\_\allowbreak{}bound}}
\leansource{\texttt{coordinate\_\allowbreak{}bias\_\allowbreak{}bound} \quad WassersteinMixing/\allowbreak{}IPLABridge.lean, lines 73–80}
\begin{leancode}
\leanline{0}{theorem\ coordinate\_bias\_bound\ (P\ :\ Parameters)}
\leanline{4}{(e\ :\ H\ ≃ₗᵢ[ℝ]\ EuclideanSpace\ ℝ\ (Fin\ P.dimension))}
\leanline{4}{(V\ :\ H\ →\ ℝ)\ (hV\ :\ PotentialAssumptions\ P\ V)\ :}
\leanline{4}{ULA.wasserstein2Real\ (ULA.gibbsMeasure\ V)\ (coordinateInvariantLaw\ P\ e\ V\ hV)\ ≤\ P.bias}
\end{leancode}
\leansource{\texttt{coordinate\_\allowbreak{}convergence\_\allowbreak{}bound} \quad WassersteinMixing/\allowbreak{}IPLABridge.lean, lines 83–94}
\begin{leancode}
\leanline{0}{theorem\ coordinate\_convergence\_bound\ (P\ :\ Parameters)}
\leanline{4}{(e\ :\ H\ ≃ₗᵢ[ℝ]\ EuclideanSpace\ ℝ\ (Fin\ P.dimension))}
\leanline{4}{(V\ :\ H\ →\ ℝ)\ (hV\ :\ PotentialAssumptions\ P\ V)}
\leanline{4}{(μ\ :\ Measure\ H)\ [IsProbabilityMeasure\ μ]\ (hμ\ :\ MemLp\ (id\ :\ H\ →\ H)\ 2\ μ)\ (n\ :\ ℕ)\ :}
\leanline{4}{ULA.wasserstein2Real\ ((ULA.ulaLawStep\ V\ P.step)\textasciicircum{}[n]\ μ)\ (ULA.gibbsMeasure\ V)\ ≤}
\leanline{6}{P.contraction\ \textasciicircum{}\ n\ *\ ULA.wasserstein2Real\ μ\ (coordinateInvariantLaw\ P\ e\ V\ hV)\ +\ P.bias}
\end{leancode}
\leanlabel{How the Lean proof follows the argument}
\texttt{coordinate\_\allowbreak{}bias\_\allowbreak{}bound} rewrites both distances through \texttt{wasserstein2\allowbreak{}Real\_\allowbreak{}map\_\allowbreak{}isometry}, identifies the two laws, and applies \texttt{bias\_\allowbreak{}from\_\allowbreak{}comparison} to \texttt{actual\allowbreak{}ULAComparison}. \texttt{coordinate\_\allowbreak{}convergence\_\allowbreak{}bound} additionally uses \texttt{ula\_\allowbreak{}iterate\_\allowbreak{}map\_\allowbreak{}isometry} and \texttt{mem\allowbreak{}Lp\_\allowbreak{}id\_\allowbreak{}map\_\allowbreak{}isometry} before applying \texttt{convergence\_\allowbreak{}from\_\allowbreak{}comparison}. The actual comparison is constructed from the potential contract rather than passed in as an extra theorem hypothesis.
\end{leanbox}


Earlier results used here: \cite[Section~3.17.4]{LangevinLean2026ULASC}; \cite[Section~3.26.3]{LangevinLean2026ULASC}; \cite[Section~3.36.7]{LangevinLean2026ULASC}.
\subsection{\texorpdfstring{Establishing finite output transport before taking real values}{Establishing finite output transport before taking real values}}
\label{report-step-WassersteinMixing-IPLABridge-lean-3}


\textbf{Claim.} The extended distance \(W_2(K_{V,h}^{\,k}\mu,\pi_V)\) is finite for every finite update count and every finite-second-moment probability initialization under the potential contract.




\textbf{Proof.} The mapped initialization has finite second moment. Positive strong convexity gives the Gibbs probability law and its finite second moment, so their initial transport distance is finite. The established Gaussian ULA output-finiteness theorem gives finite distance from every Euclidean-coordinate iterate to that target. Rewrite its iterate and Gibbs target through the isometry; the exact equality of extended transport distances transfers finiteness back to the original coordinate space. This supplies a genuine finite extended distance before later parameter-projection statements use its ordinary real value.

\begin{leanbox}{Lean: \texttt{coordinate\_\allowbreak{}convergence\_\allowbreak{}finite}}
\leansource{\texttt{coordinate\_\allowbreak{}convergence\_\allowbreak{}finite} \quad WassersteinMixing/\allowbreak{}IPLABridge.lean, lines 97–112}
\begin{leancode}
\leanline{0}{theorem\ coordinate\_convergence\_finite\ (P\ :\ Parameters)}
\leanline{4}{(e\ :\ H\ ≃ₗᵢ[ℝ]\ EuclideanSpace\ ℝ\ (Fin\ P.dimension))}
\leanline{4}{(V\ :\ H\ →\ ℝ)\ (hV\ :\ PotentialAssumptions\ P\ V)}
\leanline{4}{(μ\ :\ Measure\ H)\ [IsProbabilityMeasure\ μ]\ (hμ\ :\ MemLp\ (id\ :\ H\ →\ H)\ 2\ μ)\ (n\ :\ ℕ)\ :}
\leanline{4}{ULA.wasserstein2\ ((ULA.ulaLawStep\ V\ P.step)\textasciicircum{}[n]\ μ)\ (ULA.gibbsMeasure\ V)\ ≠\ ⊤}
\end{leancode}
\leanlabel{How the Lean proof follows the argument}
\texttt{coordinate\_\allowbreak{}convergence\_\allowbreak{}finite} first applies \texttt{ULA.\allowbreak{}wasserstein2\_\allowbreak{}ne\_\allowbreak{}top\_\allowbreak{}of\_\allowbreak{}mem\allowbreak{}Lp} and \texttt{ULA.\allowbreak{}Hilbert.\allowbreak{}dalalyan\_\allowbreak{}output\_\allowbreak{}finite} in the standard coordinates. It then rewrites by the Gaussian iterate, Gibbs pushforward and extended-distance isometry lemmas. This use of the earlier ULA theorem only establishes finiteness; the improved numerical rate still comes from the new sharp comparison.
\end{leanbox}


Earlier results used here: \cite[Section~3.17.3]{LangevinLean2026ULASC}; \cite[Section~3.21.2]{LangevinLean2026ULASC}; \cite[Section~3.26.3]{LangevinLean2026ULASC}; \cite[Section~3.36.7]{LangevinLean2026ULASC}; \cite[Section~3.45.6]{LangevinLean2026ULASC}.
\subsection{\texorpdfstring{Keeping the sharp bias with exact Gibbs-based initialization}{Keeping the sharp bias with exact Gibbs-based initialization}}
\label{report-step-WassersteinMixing-IPLABridge-lean-4}


\textbf{Claim.} The actual Gaussian chain started from the exact Gibbs target has finite transport distance from that target at every update count, and

\[W_2(\pi_V,K_{V,h}^{\,k}\pi_V)\le B.\]


For any finite-second-moment initialization this gives the sharp target-based estimate

\[W_2(K_{V,h}^{\,k}\mu,\pi_V)\le q^kW_2(\mu,\pi_V)+B.\]


\textbf{Proof.} The constructed Poisson comparison already couples a stationary Langevin process to the actual ULA chain started with the same Gibbs state. Its uniform shifted-energy estimate bounds this stationary-start displacement by \(B\), with finite extended transport proved separately. The coordinate isometry transfers both facts exactly to the original space. The contraction theorem for the same Gaussian transition compares its two actual iterates:

\[W_2(K_{V,h}^{\,k}\mu,K_{V,h}^{\,k}\pi_V)\le q^kW_2(\mu,\pi_V).\]


Apply the proved finite real transport triangle inequality through the intermediate law \(K_{V,h}^{\,k}\pi_V\). The displayed contraction bounds the first term and the stationary-start comparison bounds the second by \(B\). This retains the same exact-Gibbs initialization distance as Project 03 with a single copy of the sharp bias. The weaker bound with \((1+q^k)B\) also follows by adding the nonnegative term \(q^kB\); it remains available as a supporting corollary.

\begin{leanbox}{Lean: \texttt{coordinate\_\allowbreak{}stationary\_\allowbreak{}start\_\allowbreak{}bound}, \texttt{coordinate\_\allowbreak{}stationary\_\allowbreak{}start\_\allowbreak{}finite}, \texttt{coordinate\_\allowbreak{}convergence\_\allowbreak{}target\_\allowbreak{}bound\_\allowbreak{}sharp} and 1 more}
\leansource{\texttt{coordinate\_\allowbreak{}stationary\_\allowbreak{}start\_\allowbreak{}bound} \quad WassersteinMixing/\allowbreak{}IPLABridge.lean, lines 116–126}
\begin{leancode}
\leanline{0}{theorem\ coordinate\_stationary\_start\_bound\ (P\ :\ Parameters)}
\leanline{4}{(e\ :\ H\ ≃ₗᵢ[ℝ]\ EuclideanSpace\ ℝ\ (Fin\ P.dimension))}
\leanline{4}{(V\ :\ H\ →\ ℝ)\ (hV\ :\ PotentialAssumptions\ P\ V)\ (n\ :\ ℕ)\ :}
\leanline{4}{ULA.wasserstein2Real\ (ULA.gibbsMeasure\ V)}
\leanline{6}{((ULA.ulaLawStep\ V\ P.step)\textasciicircum{}[n]\ (ULA.gibbsMeasure\ V))\ ≤\ P.bias}
\end{leancode}
\leansource{\texttt{coordinate\_\allowbreak{}stationary\_\allowbreak{}start\_\allowbreak{}finite} \quad WassersteinMixing/\allowbreak{}IPLABridge.lean, lines 130–140}
\begin{leancode}
\leanline{0}{theorem\ coordinate\_stationary\_start\_finite\ (P\ :\ Parameters)}
\leanline{4}{(e\ :\ H\ ≃ₗᵢ[ℝ]\ EuclideanSpace\ ℝ\ (Fin\ P.dimension))}
\leanline{4}{(V\ :\ H\ →\ ℝ)\ (hV\ :\ PotentialAssumptions\ P\ V)\ (n\ :\ ℕ)\ :}
\leanline{4}{ULA.wasserstein2\ (ULA.gibbsMeasure\ V)}
\leanline{6}{((ULA.ulaLawStep\ V\ P.step)\textasciicircum{}[n]\ (ULA.gibbsMeasure\ V))\ ≠\ ⊤}
\end{leancode}
\leansource{\texttt{coordinate\_\allowbreak{}convergence\_\allowbreak{}target\_\allowbreak{}bound\_\allowbreak{}sharp} \quad WassersteinMixing/\allowbreak{}IPLABridge.lean, lines 145–169}
\begin{leancode}
\leanline{0}{theorem\ coordinate\_convergence\_target\_bound\_sharp\ (P\ :\ Parameters)}
\leanline{4}{(e\ :\ H\ ≃ₗᵢ[ℝ]\ EuclideanSpace\ ℝ\ (Fin\ P.dimension))}
\leanline{4}{(V\ :\ H\ →\ ℝ)\ (hV\ :\ PotentialAssumptions\ P\ V)}
\leanline{4}{(μ\ :\ Measure\ H)\ [IsProbabilityMeasure\ μ]\ (hμ\ :\ MemLp\ (id\ :\ H\ →\ H)\ 2\ μ)\ (n\ :\ ℕ)\ :}
\leanline{4}{ULA.wasserstein2Real\ ((ULA.ulaLawStep\ V\ P.step)\textasciicircum{}[n]\ μ)\ (ULA.gibbsMeasure\ V)\ ≤}
\leanline{6}{P.contraction\ \textasciicircum{}\ n\ *\ ULA.wasserstein2Real\ μ\ (ULA.gibbsMeasure\ V)\ +}
\leanline{8}{P.bias}
\end{leancode}
\leansource{\texttt{coordinate\_\allowbreak{}convergence\_\allowbreak{}target\_\allowbreak{}bound} \quad WassersteinMixing/\allowbreak{}IPLABridge.lean, lines 173–182}
\begin{leancode}
\leanline{0}{theorem\ coordinate\_convergence\_target\_bound\ (P\ :\ Parameters)}
\leanline{4}{(e\ :\ H\ ≃ₗᵢ[ℝ]\ EuclideanSpace\ ℝ\ (Fin\ P.dimension))}
\leanline{4}{(V\ :\ H\ →\ ℝ)\ (hV\ :\ PotentialAssumptions\ P\ V)}
\leanline{4}{(μ\ :\ Measure\ H)\ [IsProbabilityMeasure\ μ]\ (hμ\ :\ MemLp\ (id\ :\ H\ →\ H)\ 2\ μ)\ (n\ :\ ℕ)\ :}
\leanline{4}{ULA.wasserstein2Real\ ((ULA.ulaLawStep\ V\ P.step)\textasciicircum{}[n]\ μ)\ (ULA.gibbsMeasure\ V)\ ≤}
\leanline{6}{P.contraction\ \textasciicircum{}\ n\ *\ ULA.wasserstein2Real\ μ\ (ULA.gibbsMeasure\ V)\ +}
\leanline{8}{(1\ +\ P.contraction\ \textasciicircum{}\ n)\ *\ P.bias}
\end{leancode}
\leanlabel{How the Lean proof follows the argument}
\texttt{coordinate\_\allowbreak{}stationary\_\allowbreak{}start\_\allowbreak{}bound} and \texttt{coordinate\_\allowbreak{}stationary\_\allowbreak{}start\_\allowbreak{}finite} transport the corresponding consequences of \texttt{actual\allowbreak{}ULAComparison} through the exact coordinate identities. \texttt{coordinate\_\allowbreak{}convergence\_\allowbreak{}target\_\allowbreak{}bound\_\allowbreak{}sharp} proves initial transport finiteness from moments, supplies probability for both iterates, and combines \texttt{ula\_\allowbreak{}iterate\_\allowbreak{}wasserstein2\allowbreak{}Real\_\allowbreak{}le} with \texttt{wasserstein2\allowbreak{}Real\_\allowbreak{}triangle} and the stationary-start bound. \texttt{coordinate\_\allowbreak{}convergence\_\allowbreak{}target\_\allowbreak{}bound} is an algebraic weakening of that sharp theorem using nonnegativity of the contraction power and the bias. No invariant-triangle loss is needed by the primary target-based result.
\end{leanbox}


Earlier results used here: \cite[Section~3.17.3]{LangevinLean2026ULASC}; \cite[Section~3.17.4]{LangevinLean2026ULASC}; \cite[Section~3.21.2]{LangevinLean2026ULASC}; \cite[Section~3.26.13]{LangevinLean2026ULASC}; \cite[Section~3.26.15]{LangevinLean2026ULASC}; \cite[Section~3.26.3]{LangevinLean2026ULASC}; \cite[Section~3.36.7]{LangevinLean2026ULASC}.
\subsection{\texorpdfstring{Applying the coordinate theorem with the exact IPLA constants}{Applying the coordinate theorem with the exact IPLA constants}}
\label{report-step-WassersteinMixing-IPLABridge-lean-5}


\textbf{Claim.} For \(N>0\), let \(U\in C^2\) be \(m\)-strongly convex with an \(L\)-Lipschitz gradient, \(0<m\le L\), and \(0<\gamma\) with \(\gamma L\le1\). Define \(V_N=\sum_iU(\theta,\sqrt N z_i)\), \(d_N=d_\theta+Nd_x\), and \(D=d_x+d_\theta/N\). Then for the actual lifted Gaussian updates,

\[W_2(K_{V_N,\gamma/N}^{\,k}\mu,\pi_{V_N})\le(1-m\gamma)^kW_2(\mu,\widehat\pi_{V_N})+6\gamma\sqrt{(L/m)LD}.\]


\textbf{Proof.} The index set of the scaled state has exactly \(d_\theta+Nd_x\) elements, so an index equivalence gives a linear isometry to the required standard coordinate space. The joint-potential lifting theorem proves the potential contract with \(\alpha=Nm\), \(\beta=NL\), and \(h=\gamma/N\). Apply the coordinate convergence theorem. Exact cancellation gives \(1-\alpha h=1-m\gamma\), and the proved bias-scaling identity reduces the sharp bias to the displayed linear-step expression. No relation between the parameter and latent dimensions is assumed.

\begin{leanbox}{Lean: \texttt{iplaCoordinates}, \texttt{ipla\_\allowbreak{}wasserstein\_\allowbreak{}bound}}
\leansource{\texttt{iplaCoordinates} \quad WassersteinMixing/\allowbreak{}IPLABridge.lean, lines 186–189}
\begin{leancode}
\leanline{0}{def\ iplaCoordinates\ (N\ dθ\ dx\ :\ ℕ)\ :}
\leanline{4}{IPLA.State\ N\ dθ\ dx\ ≃ₗᵢ[ℝ]\ EuclideanSpace\ ℝ\ (Fin\ (dθ\ +\ N\ *\ dx))\ :=}
\leanline{2}{LinearIsometryEquiv.piLpCongrLeft\ 2\ ℝ\ ℝ}
\leanline{4}{(Fintype.equivFinOfCardEq\ (by\ simp))}
\end{leancode}
\leansource{\texttt{ipla\_\allowbreak{}wasserstein\_\allowbreak{}bound} \quad WassersteinMixing/\allowbreak{}IPLABridge.lean, lines 194–215}
\begin{leancode}
\leanline{0}{theorem\ ipla\_wasserstein\_bound}
\leanline{4}{\{N\ dθ\ dx\ :\ ℕ\}\ \{U\ :\ IPLA.Joint\ dθ\ dx\ →\ ℝ\}\ \{m\ γ\ :\ ℝ\}\ \{L\ :\ ℝ≥0\}}
\leanline{4}{(hN\ :\ 0\ <\ N)\ (hm\ :\ 0\ <\ m)\ (hmL\ :\ m\ ≤\ L)\ (hγ\ :\ 0\ <\ γ)\ (hγL\ :\ γ\ *\ L\ ≤\ 1)}
\leanline{4}{(hU\ :\ ContDiff\ ℝ\ 2\ U)\ (hsc\ :\ StrongConvexOn\ univ\ m\ U)}
\leanline{4}{(hLip\ :\ LipschitzWith\ L\ (gradient\ U))}
\leanline{4}{(μ\ :\ Measure\ (IPLA.State\ N\ dθ\ dx))\ [IsProbabilityMeasure\ μ]}
\leanline{4}{(hμ\ :\ MemLp\ (id\ :\ IPLA.State\ N\ dθ\ dx\ →\ IPLA.State\ N\ dθ\ dx)\ 2\ μ)\ (k\ :\ ℕ)\ :}
\leanline{4}{ULA.wasserstein2Real\ ((ULA.ulaLawStep\ (@IPLA.potential\ N\ dθ\ dx\ U)\ (γ\ /\ N))\textasciicircum{}[k]\ μ)}
\leanline{6}{(ULA.gibbsMeasure\ (@IPLA.potential\ N\ dθ\ dx\ U))\ ≤}
\leanline{6}{(1\ -\ m\ *\ γ)\ \textasciicircum{}\ k\ *\ ULA.wasserstein2Real\ μ}
\leanline{8}{(coordinateInvariantLaw\ (iplaParameters\ N\ dθ\ dx\ m\ L\ γ\ hN\ hm\ hmL\ hγ\ hγL)}
\leanline{10}{(iplaCoordinates\ N\ dθ\ dx)\ (IPLA.potential\ U)}
\leanline{10}{(ipla\_potentialAssumptions\ hN\ hm\ hmL\ hγ\ hγL\ hU\ hsc\ hLip))\ +}
\leanline{8}{iplaLinearBias\ m\ L\ γ\ (iplaEffectiveDimension\ N\ dθ\ dx)}
\end{leancode}
\leanlabel{How the Lean proof follows the argument}
\texttt{ipla\allowbreak{}Coordinates} uses the actual finite index equivalence of the sum/\allowbreak{}product coordinate set. \texttt{ipla\_\allowbreak{}wasserstein\_\allowbreak{}bound} supplies \texttt{ipla\allowbreak{}Parameters} and \texttt{ipla\_\allowbreak{}potential\allowbreak{}Assumptions} to \texttt{coordinate\_\allowbreak{}convergence\_\allowbreak{}bound}, then uses \texttt{ipla\_\allowbreak{}lifted\_\allowbreak{}bias\_\allowbreak{}eq} to simplify the numerical constant. The law is the actual ULA law of the scaled IPLA potential; the final application identifies it with the noise-driven IPLA recursion before parameter projection.
\end{leanbox}


Earlier results used here: \cite[Section~3.17.4]{LangevinLean2026ULASC}; \cite[Section~3.26.3]{LangevinLean2026ULASC}; \cite[Section~3.36.7]{LangevinLean2026ULASC}; \cite[Section~3.1.1]{LangevinLean2026IPLA}.
\section{\texorpdfstring{New linear-step parameter rates for actual Gaussian IPLA}{New linear-step parameter rates for actual Gaussian IPLA}}
\label{proof-module-WassersteinMixing-IPLA-lean}


For \(N>0\) particles and independent parameter and latent dimensions, assume the joint potential \(U\in C^2\) is \(m\)-strongly convex with an \(L\)-Lipschitz gradient, \(0<m\le L\), and \(0<\gamma\le1/L\). For the actual Gaussian IPLA laws, the parameter root-mean-square error is bounded by the equilibrium radius \(\sqrt{2d_\theta/(Nm)}\), a geometric transient, and the linear-step discretization term \(6\gamma\sqrt{(L/m)L(d_x+d_\theta/N)}\). The same sharp term works with initialization measured from either the constructed discretized invariant law or the exact Gibbs target. The resulting explicit accuracy and Project 03 comparison theorems have all structural assumptions displayed.



Earlier results used here: \cite[Section~3.13]{LangevinLean2026IPLA}; \cite[Section~3.4]{LangevinLean2026IPLA}; \cite[Section~3.5]{LangevinLean2026IPLA}.
\subsection{\texorpdfstring{Displaying every structural assumption and deriving the lifted contract}{Displaying every structural assumption and deriving the lifted contract}}
\label{report-step-WassersteinMixing-IPLA-lean-1}


\textbf{Setting.} Let \(N,d_\theta,d_x\) be nonnegative integers with \(N>0\). Assume \(U:\mathbb R^{d_\theta+d_x}\to\mathbb R\) belongs to \(C^2\), is \(m\)-strongly convex, and has an \(L\)-Lipschitz joint gradient, where \(0<m\le L\). Let \(0<\gamma\) satisfy \(\gamma L\le1\).




\textbf{Claim and proof.} For the actual scaled potential \(V_N(z)=\sum_iU(\theta,\sqrt N z_i)\), the previous lifting theorems give \(C^2\) regularity, curvature \(Nm\), gradient Lipschitz constant \(NL\), dimension \(d_\theta+Nd_x\), and ULA step \(\gamma/N\). In particular the required small-step condition reduces exactly to \(\gamma L\le1\). These are consequences of the joint assumptions, not extra analytic hypotheses on a lifted reference process.




The regularity and step range are explicit additions relative to the full scope of Project 03: that theorem needs differentiability and a Lipschitz gradient and permits \(\gamma\le2/(m+L)\). The improved rate is asserted within the regularity class and step interval displayed here.

\begin{leanbox}{Lean: \texttt{IPLASetting}, \texttt{IPLASetting.\allowbreak{}parameters}, \texttt{IPLASetting.\allowbreak{}potentialAssumptions}}
\leansource{\texttt{IPLASetting} \quad WassersteinMixing/\allowbreak{}IPLA.lean, lines 26–41}
\begin{leancode}
\leanline{0}{structure\ IPLASetting\ where}
\leanline{2}{N\ :\ ℕ}
\leanline{2}{dθ\ :\ ℕ}
\leanline{2}{dx\ :\ ℕ}
\leanline{2}{U\ :\ IPLA.Joint\ dθ\ dx\ →\ ℝ}
\leanline{2}{m\ :\ ℝ}
\leanline{2}{L\ :\ ℝ≥0}
\leanline{2}{γ\ :\ ℝ}
\leanline{2}{particles\_pos\ :\ 0\ <\ N}
\leanline{2}{curvature\_pos\ :\ 0\ <\ m}
\leanline{2}{curvature\_le\_smoothness\ :\ m\ ≤\ L}
\leanline{2}{step\_pos\ :\ 0\ <\ γ}
\leanline{2}{step\_small\ :\ γ\ *\ L\ ≤\ 1}
\leanline{2}{regularity\ :\ ContDiff\ ℝ\ 2\ U}
\leanline{2}{strongConvexity\ :\ StrongConvexOn\ univ\ m\ U}
\leanline{2}{gradientLipschitz\ :\ LipschitzWith\ L\ (gradient\ U)}
\end{leancode}
\leansource{\texttt{IPLASetting.\allowbreak{}parameters} \quad WassersteinMixing/\allowbreak{}IPLA.lean, lines 44–46}
\begin{leancode}
\leanline{0}{def\ IPLASetting.parameters\ (A\ :\ IPLASetting)\ :\ Parameters\ :=}
\leanline{2}{iplaParameters\ A.N\ A.dθ\ A.dx\ A.m\ A.L\ A.γ\ A.particles\_pos\ A.curvature\_pos}
\leanline{4}{A.curvature\_le\_smoothness\ A.step\_pos\ A.step\_small}
\end{leancode}
\leansource{\texttt{IPLASetting.\allowbreak{}potentialAssumptions} \quad WassersteinMixing/\allowbreak{}IPLA.lean, lines 49–52}
\begin{leancode}
\leanline{0}{theorem\ IPLASetting.potentialAssumptions\ (A\ :\ IPLASetting)\ :}
\leanline{4}{PotentialAssumptions\ A.parameters\ (@IPLA.potential\ A.N\ A.dθ\ A.dx\ A.U)}
\end{leancode}
\leanlabel{How the Lean proof follows the argument}
\texttt{IPLASetting} contains the joint potential, dimensions, curvature and step constants, and exactly the structural hypotheses just listed. \texttt{IPLASetting.\allowbreak{}parameters} uses the proved \texttt{ipla\allowbreak{}Parameters} construction. \texttt{IPLASetting.\allowbreak{}potential\allowbreak{}Assumptions} invokes \texttt{ipla\_\allowbreak{}potential\allowbreak{}Assumptions} to derive the complete lifted contract. None of these structures contains a comparison bound, invariant law, semigroup or concentration estimate as a premise.
\end{leanbox}


Earlier results used here: \cite[Section~3.1.1]{LangevinLean2026IPLA}.
\subsection{\texorpdfstring{The actual Gaussian algorithm and its constructed statistical objects}{The actual Gaussian algorithm and its constructed statistical objects}}
\label{report-step-WassersteinMixing-IPLA-lean-2}


\textbf{Definitions.} Let \(\pi_N\) be the exact Gibbs law of \(V_N\), and let \(\widehat\pi_{N,\gamma}\) be the genuine Gaussian ULA invariant law pulled back to the scaled particle coordinates. Write \(\mu_k=K_{V_N,\gamma/N}^{\,k}\mu_0\). Set

\[C_N=\sqrt{\frac{2d_\theta}{Nm}},\qquad D=d_x+\frac{d_\theta}{N},\qquad B_N=6\gamma\sqrt{(L/m)LD}.\]


\textbf{Algorithm identification.} Under \(z_i=x_i/\sqrt N\), the kernel is precisely the IPLA recursion

\[\theta_{k+1}=\theta_k-\frac\gamma N\sum_i\nabla_\theta U(\theta_k,x_k^i)+\sqrt{\frac{2\gamma}{N}}\,\xi_{k+1}^{\,0},\qquad x_{k+1}^i=x_k^i-\gamma\nabla_xU(\theta_k,x_k^i)+\sqrt{2\gamma}\,\xi_{k+1}^{\,i},\]


with independent standard Gaussian coordinates at every time and independence from initialization.




\textbf{Proof.} The existing IPLA geometry computes the parameter gradient as the sum of joint parameter gradients and each scaled latent gradient as \(\sqrt N\) times its joint latent gradient. Substitution of the lifted step \(\gamma/N\), followed by rescaling the latent coordinates, gives exactly the displayed updates. The already proved identity of the Gaussian recursion laws equates the laws of every recursion satisfying the independent-innovation specification with these kernel iterates. The invariant law is the one constructed by the actual Gaussian contraction argument, rather than a measure specified only through an assumed invariance property.

\begin{leanbox}{Lean: \texttt{IPLASetting.\allowbreak{}invariantLaw}, \texttt{IPLASetting.\allowbreak{}law}, \texttt{IPLASetting.\allowbreak{}law\_\allowbreak{}eq\_\allowbreak{}gaussian} and 2 more}
\leansource{\texttt{IPLASetting.\allowbreak{}invariantLaw} \quad WassersteinMixing/\allowbreak{}IPLA.lean, lines 55–57}
\begin{leancode}
\leanline{0}{def\ IPLASetting.invariantLaw\ (A\ :\ IPLASetting)\ :\ Measure\ (IPLA.State\ A.N\ A.dθ\ A.dx)\ :=}
\leanline{2}{coordinateInvariantLaw\ A.parameters\ (iplaCoordinates\ A.N\ A.dθ\ A.dx)}
\leanline{4}{(IPLA.potential\ A.U)\ A.potentialAssumptions}
\end{leancode}
\leansource{\texttt{IPLASetting.\allowbreak{}law} \quad WassersteinMixing/\allowbreak{}IPLA.lean, lines 60–62}
\begin{leancode}
\leanline{0}{def\ IPLASetting.law\ (A\ :\ IPLASetting)\ (μ\ :\ Measure\ (IPLA.State\ A.N\ A.dθ\ A.dx))}
\leanline{4}{(k\ :\ ℕ)\ :\ Measure\ (IPLA.State\ A.N\ A.dθ\ A.dx)\ :=}
\leanline{2}{(ULA.ulaLawStep\ (IPLA.potential\ A.U)\ (A.γ\ /\ A.N))\textasciicircum{}[k]\ μ}
\end{leancode}
\leansource{\texttt{IPLASetting.\allowbreak{}law\_\allowbreak{}eq\_\allowbreak{}gaussian} \quad WassersteinMixing/\allowbreak{}IPLA.lean, lines 65–74}
\begin{leancode}
\leanline{0}{theorem\ IPLASetting.law\_eq\_gaussian\ (A\ :\ IPLASetting)}
\leanline{4}{\{Ω\ :\ Type*\}\ [MeasurableSpace\ Ω]\ (P\ :\ Measure\ Ω)\ [IsProbabilityMeasure\ P]}
\leanline{4}{(initial\ :\ Lp\ (IPLA.State\ A.N\ A.dθ\ A.dx)\ 2\ P)}
\leanline{4}{(noise\ :\ ℕ\ →\ Lp\ (IPLA.State\ A.N\ A.dθ\ A.dx)\ 2\ P)}
\leanline{4}{(hGaussian\ :\ IPLA.GaussianInnovations\ P\ initial\ noise)\ (k\ :\ ℕ)\ :}
\leanline{4}{A.law\ (P.map\ initial)\ k\ =\ IPLA.laws\ P\ A.U\ A.γ\ A.L}
\leanline{6}{(A.regularity.differentiable\ (by\ norm\_num))\ A.gradientLipschitz\ initial\ noise\ k}
\end{leancode}
\leansource{\texttt{IPLASetting.\allowbreak{}radius} \quad WassersteinMixing/\allowbreak{}IPLA.lean, lines 77–78}
\begin{leancode}
\leanline{0}{def\ IPLASetting.radius\ (A\ :\ IPLASetting)\ :\ ℝ\ :=}
\leanline{2}{Real.sqrt\ (2\ *\ (A.dθ\ :\ ℝ)\ /\ ((A.N\ :\ ℝ)\ *\ A.m))}
\end{leancode}
\leansource{\texttt{IPLASetting.\allowbreak{}bias} \quad WassersteinMixing/\allowbreak{}IPLA.lean, lines 81–82}
\begin{leancode}
\leanline{0}{def\ IPLASetting.bias\ (A\ :\ IPLASetting)\ :\ ℝ\ :=}
\leanline{2}{iplaLinearBias\ A.m\ A.L\ A.γ\ (iplaEffectiveDimension\ A.N\ A.dθ\ A.dx)}
\end{leancode}
\leanlabel{How the Lean proof follows the argument}
\texttt{IPLASetting.\allowbreak{}invariant\allowbreak{}Law} specializes \texttt{coordinate\allowbreak{}Invariant\allowbreak{}Law} with the actual index isometry. \texttt{IPLASetting.\allowbreak{}law} iterates \texttt{ULA.\allowbreak{}ula\allowbreak{}Law\allowbreak{}Step} for the lifted potential and exact scaled step. \texttt{IPLASetting.\allowbreak{}law\_\allowbreak{}eq\_\allowbreak{}gaussian} invokes \texttt{IPLA.\allowbreak{}laws\_\allowbreak{}eq\_\allowbreak{}ula\allowbreak{}Law\allowbreak{}Step}, whose \texttt{Gaussian\allowbreak{}Innovations} fields encode all coordinate laws and independence requirements. \texttt{radius} and \texttt{bias} define the literal coefficients used in the final statements.
\end{leanbox}


Earlier results used here: \cite[Section~3.36.7]{LangevinLean2026ULASC}; \cite[Section~3.1.1]{LangevinLean2026IPLA}; \cite[Section~3.16.3]{LangevinLean2026IPLA}; \cite[Section~3.3.2]{LangevinLean2026IPLA}; \cite[Section~3.3.3]{LangevinLean2026IPLA}.
\subsection{\texorpdfstring{Projecting the finite transport estimate to the parameter}{Projecting the finite transport estimate to the parameter}}
\label{report-step-WassersteinMixing-IPLA-lean-3}


\textbf{Claim.} Let \(\theta^\star\) maximize the latent-integrated density \(\int e^{-U(\theta,x)}\,dx\). For a finite-second-moment probability initialization,

\[\left(\mathbb E_{\mu_k}\|\theta-\theta^\star\|^2\right)^{1/2}\le W_2(\mu_k,\pi_N)+C_N.\]


\textbf{Proof.} The established IPLA concentration theorem gives \((\mathbb E_{\pi_N}\|\theta-\theta^\star\|^2)^{1/2}\le C_N\) from the same joint strong convexity. The parameter projection has operator norm at most one. For any coupling of \(\mu_k\) and \(\pi_N\), the \(L^2\) triangle inequality bounds the projected distance to \(\theta^\star\) by its coupling displacement plus that target radius. Taking the coupling infimum gives the displayed inequality. The coordinate convergence-finiteness theorem proves finite output transport first; this and concentration prove a finite projected second moment before converting the extended moment to the ordinary square-root integral.

\begin{leanbox}{Lean: \texttt{ipla\_\allowbreak{}parameter\_\allowbreak{}projection}}
\leansource{\texttt{ipla\_\allowbreak{}parameter\_\allowbreak{}projection} \quad WassersteinMixing/\allowbreak{}IPLA.lean, lines 85–104}
\begin{leancode}
\leanline{0}{theorem\ ipla\_parameter\_projection\ (A\ :\ IPLASetting)}
\leanline{4}{(θStar\ :\ EuclideanSpace\ ℝ\ (Fin\ A.dθ))\ (hmax\ :\ IPLA.IsMarginalMaximizer\ A.U\ θStar)}
\leanline{4}{(μ\ :\ Measure\ (IPLA.State\ A.N\ A.dθ\ A.dx))\ [IsProbabilityMeasure\ μ]}
\leanline{4}{(hμ\ :\ MemLp\ (id\ :\ IPLA.State\ A.N\ A.dθ\ A.dx\ →\ IPLA.State\ A.N\ A.dθ\ A.dx)\ 2\ μ)}
\leanline{4}{(k\ :\ ℕ)\ :}
\leanline{4}{Real.sqrt\ (∫\ z,\ ‖IPLA.parameterCLM\ z\ -\ θStar‖\ \textasciicircum{}\ 2\ ∂A.law\ μ\ k)\ ≤}
\leanline{6}{ULA.wasserstein2Real\ (A.law\ μ\ k)\ (@IPLA.gibbsTarget\ A.N\ A.dθ\ A.dx\ A.U)\ +}
\leanline{8}{A.radius}
\end{leancode}
\leanlabel{How the Lean proof follows the argument}
\texttt{ipla\_\allowbreak{}parameter\_\allowbreak{}projection} obtains finite transport from \texttt{coordinate\_\allowbreak{}convergence\_\allowbreak{}finite}, obtains the sharp statistical radius from \texttt{IPLA.\allowbreak{}gibbs\allowbreak{}Target\_\allowbreak{}parameter\_\allowbreak{}concentration}, and applies the proved one-Lipschitz parameter-projection theorem. \texttt{projected\allowbreak{}RMS\_\allowbreak{}ne\_\allowbreak{}top} establishes finiteness before \texttt{projected\allowbreak{}RMSReal\_\allowbreak{}eq\_\allowbreak{}sqrt\_\allowbreak{}integral} rewrites the final expression. The maximizer premise identifies the statistical center; the concentration estimate itself is derived.
\end{leanbox}


Earlier results used here: \cite[Section~3.17.4]{LangevinLean2026ULASC}; \cite[Section~3.1.1]{LangevinLean2026IPLA}; \cite[Section~3.1.6]{LangevinLean2026IPLA}; \cite[Section~3.13.3]{LangevinLean2026IPLA}; \cite[Section~3.3.1]{LangevinLean2026IPLA}; \cite[Section~3.3.2]{LangevinLean2026IPLA}.
\subsection{\texorpdfstring{The linear-step parameter bounds with either initialization distance}{The linear-step parameter bounds with either initialization distance}}
\label{report-step-WassersteinMixing-IPLA-lean-4}


\textbf{Claim.} Define \(q=1-m\gamma\in[0,1)\), \(R_h=W_2(\mu_0,\widehat\pi_{N,\gamma})\), and \(R_0=W_2(\mu_0,\pi_N)\). Then the actual IPLA parameter error satisfies both

\[\left(\mathbb E_{\mu_k}\|\theta-\theta^\star\|^2\right)^{1/2}\le C_N+q^kR_h+B_N,\qquad \left(\mathbb E_{\mu_k}\|\theta-\theta^\star\|^2\right)^{1/2}\le C_N+q^kR_0+B_N.\]


\textbf{Proof.} For the invariant-based form, apply the fully constructed lifted Wasserstein comparison and then the parameter-projection inequality. For the target-based form, compare the actual iterates from \(\mu_0\) and \(\pi_N\) by the Gaussian contraction theorem, yielding \(q^kR_0\). The stationary-start Poisson comparison bounds the distance of the second iterate from \(\pi_N\) by \(B_N\). Their transport triangle and the same projection bound give the second formula with only a single bias. The lifted constants simplify exactly to \(q\) and \(B_N\) by the established particle-scaling identities. The supporting uniform corollary with \(2B_N\) follows by weakening this bound using \(B_N\ge0\); it is not required by the sharper target statement.

\begin{leanbox}{Lean: \texttt{ipla\_\allowbreak{}parameter\_\allowbreak{}error}, \texttt{ipla\_\allowbreak{}parameter\_\allowbreak{}error\_\allowbreak{}target}, \texttt{ipla\_\allowbreak{}parameter\_\allowbreak{}error\_\allowbreak{}target\_\allowbreak{}uniform}}
\leansource{\texttt{ipla\_\allowbreak{}parameter\_\allowbreak{}error} \quad WassersteinMixing/\allowbreak{}IPLA.lean, lines 107–120}
\begin{leancode}
\leanline{0}{theorem\ ipla\_parameter\_error\ (A\ :\ IPLASetting)}
\leanline{4}{(θStar\ :\ EuclideanSpace\ ℝ\ (Fin\ A.dθ))\ (hmax\ :\ IPLA.IsMarginalMaximizer\ A.U\ θStar)}
\leanline{4}{(μ\ :\ Measure\ (IPLA.State\ A.N\ A.dθ\ A.dx))\ [IsProbabilityMeasure\ μ]}
\leanline{4}{(hμ\ :\ MemLp\ (id\ :\ IPLA.State\ A.N\ A.dθ\ A.dx\ →\ IPLA.State\ A.N\ A.dθ\ A.dx)\ 2\ μ)}
\leanline{4}{(k\ :\ ℕ)\ :}
\leanline{4}{Real.sqrt\ (∫\ z,\ ‖IPLA.parameterCLM\ z\ -\ θStar‖\ \textasciicircum{}\ 2\ ∂A.law\ μ\ k)\ ≤}
\leanline{6}{A.radius\ +\ (1\ -\ A.m\ *\ A.γ)\ \textasciicircum{}\ k\ *\ ULA.wasserstein2Real\ μ\ A.invariantLaw\ +\ A.bias}
\end{leancode}
\leansource{\texttt{ipla\_\allowbreak{}parameter\_\allowbreak{}error\_\allowbreak{}target} \quad WassersteinMixing/\allowbreak{}IPLA.lean, lines 123–147}
\begin{leancode}
\leanline{0}{theorem\ ipla\_parameter\_error\_target\ (A\ :\ IPLASetting)}
\leanline{4}{(θStar\ :\ EuclideanSpace\ ℝ\ (Fin\ A.dθ))\ (hmax\ :\ IPLA.IsMarginalMaximizer\ A.U\ θStar)}
\leanline{4}{(μ\ :\ Measure\ (IPLA.State\ A.N\ A.dθ\ A.dx))\ [IsProbabilityMeasure\ μ]}
\leanline{4}{(hμ\ :\ MemLp\ (id\ :\ IPLA.State\ A.N\ A.dθ\ A.dx\ →\ IPLA.State\ A.N\ A.dθ\ A.dx)\ 2\ μ)}
\leanline{4}{(k\ :\ ℕ)\ :}
\leanline{4}{Real.sqrt\ (∫\ z,\ ‖IPLA.parameterCLM\ z\ -\ θStar‖\ \textasciicircum{}\ 2\ ∂A.law\ μ\ k)\ ≤}
\leanline{6}{A.radius\ +\ (1\ -\ A.m\ *\ A.γ)\ \textasciicircum{}\ k\ *}
\leanline{8}{ULA.wasserstein2Real\ μ\ (@IPLA.gibbsTarget\ A.N\ A.dθ\ A.dx\ A.U)\ +}
\leanline{8}{A.bias}
\end{leancode}
\leansource{\texttt{ipla\_\allowbreak{}parameter\_\allowbreak{}error\_\allowbreak{}target\_\allowbreak{}uniform} \quad WassersteinMixing/\allowbreak{}IPLA.lean, lines 150–161}
\begin{leancode}
\leanline{0}{theorem\ ipla\_parameter\_error\_target\_uniform\ (A\ :\ IPLASetting)}
\leanline{4}{(θStar\ :\ EuclideanSpace\ ℝ\ (Fin\ A.dθ))\ (hmax\ :\ IPLA.IsMarginalMaximizer\ A.U\ θStar)}
\leanline{4}{(μ\ :\ Measure\ (IPLA.State\ A.N\ A.dθ\ A.dx))\ [IsProbabilityMeasure\ μ]}
\leanline{4}{(hμ\ :\ MemLp\ (id\ :\ IPLA.State\ A.N\ A.dθ\ A.dx\ →\ IPLA.State\ A.N\ A.dθ\ A.dx)\ 2\ μ)}
\leanline{4}{(k\ :\ ℕ)\ :}
\leanline{4}{Real.sqrt\ (∫\ z,\ ‖IPLA.parameterCLM\ z\ -\ θStar‖\ \textasciicircum{}\ 2\ ∂A.law\ μ\ k)\ ≤}
\leanline{6}{A.radius\ +\ (1\ -\ A.m\ *\ A.γ)\ \textasciicircum{}\ k\ *}
\leanline{8}{ULA.wasserstein2Real\ μ\ (@IPLA.gibbsTarget\ A.N\ A.dθ\ A.dx\ A.U)\ +\ 2\ *\ A.bias}
\end{leancode}
\leanlabel{How the Lean proof follows the argument}
\texttt{ipla\_\allowbreak{}parameter\_\allowbreak{}error} combines \texttt{ipla\_\allowbreak{}wasserstein\_\allowbreak{}bound} with \texttt{ipla\_\allowbreak{}parameter\_\allowbreak{}projection}. \texttt{ipla\_\allowbreak{}parameter\_\allowbreak{}error\_\allowbreak{}target} uses \texttt{coordinate\_\allowbreak{}convergence\_\allowbreak{}target\_\allowbreak{}bound\_\allowbreak{}sharp}, rewrites the lifted contraction and bias by their exact identities, and adds the derived concentration radius. \texttt{ipla\_\allowbreak{}parameter\_\allowbreak{}error\_\allowbreak{}target\_\allowbreak{}uniform} retains a coarser supporting estimate. The final statements use the actual law and require no supplied transport comparison or invariant-law construction.
\end{leanbox}


Earlier results used here: \cite[Section~3.17.4]{LangevinLean2026ULASC}; \cite[Section~3.1.1]{LangevinLean2026IPLA}; \cite[Section~3.1.6]{LangevinLean2026IPLA}; \cite[Section~3.3.1]{LangevinLean2026IPLA}; \cite[Section~3.3.2]{LangevinLean2026IPLA}; \cite[Section~3.5.1]{LangevinLean2026IPLA}.
\subsection{\texorpdfstring{A direct comparison with the complete Project 03 bound}{A direct comparison with the complete Project 03 bound}}
\label{report-step-WassersteinMixing-IPLA-lean-5}


\textbf{Claim.} Under the preceding structural assumptions and the additional comparison regime \(m\gamma\le121/1600\), the new target-based right side is no greater than

\[C_N+q^kW_2(\mu_0,\pi_N)+\frac{33}{20}\frac Lm\sqrt D\sqrt\gamma,\]


the right side used by Project 03 for the same dimensions and initialization.




\textbf{Proof.} The two bounds have the same statistical radius and exactly the same Gibbs-based transient. The only change is the discretization term. The established coefficient inequality \(B_N\le(33/20)(L/m)\sqrt D\sqrt\gamma\) holds throughout the stated regime. Add the unchanged two terms to that inequality and combine it with the new sharp parameter bound. Thus the linear-step estimate gives an explicit improvement for sufficiently small steps; it does not claim to cover the larger step interval or weaker regularity assumptions of every case of Project 03.

\begin{leanbox}{Lean: \texttt{ipla\_\allowbreak{}parameter\_\allowbreak{}error\_\allowbreak{}improves\_\allowbreak{}project03}}
\leansource{\texttt{ipla\_\allowbreak{}parameter\_\allowbreak{}error\_\allowbreak{}improves\_\allowbreak{}project03} \quad WassersteinMixing/\allowbreak{}IPLA.lean, lines 165–178}
\begin{leancode}
\leanline{0}{theorem\ ipla\_parameter\_error\_improves\_project03\ (A\ :\ IPLASetting)}
\leanline{4}{(θStar\ :\ EuclideanSpace\ ℝ\ (Fin\ A.dθ))\ (hmax\ :\ IPLA.IsMarginalMaximizer\ A.U\ θStar)}
\leanline{4}{(μ\ :\ Measure\ (IPLA.State\ A.N\ A.dθ\ A.dx))\ [IsProbabilityMeasure\ μ]}
\leanline{4}{(hμ\ :\ MemLp\ (id\ :\ IPLA.State\ A.N\ A.dθ\ A.dx\ →\ IPLA.State\ A.N\ A.dθ\ A.dx)\ 2\ μ)}
\leanline{4}{(k\ :\ ℕ)\ (hsmall\ :\ A.m\ *\ A.γ\ ≤\ (121\ /\ 1600\ :\ ℝ))\ :}
\leanline{4}{Real.sqrt\ (∫\ z,\ ‖IPLA.parameterCLM\ z\ -\ θStar‖\ \textasciicircum{}\ 2\ ∂A.law\ μ\ k)\ ≤}
\leanline{6}{A.radius\ +\ (1\ -\ A.m\ *\ A.γ)\ \textasciicircum{}\ k\ *}
\leanline{8}{ULA.wasserstein2Real\ μ\ (@IPLA.gibbsTarget\ A.N\ A.dθ\ A.dx\ A.U)\ +}
\leanline{8}{iplaProject03Bias\ A.m\ A.L\ A.γ\ (iplaEffectiveDimension\ A.N\ A.dθ\ A.dx)}
\end{leancode}
\leanlabel{How the Lean proof follows the argument}
\texttt{ipla\_\allowbreak{}parameter\_\allowbreak{}error\_\allowbreak{}improves\_\allowbreak{}project03} applies \texttt{ipla\_\allowbreak{}parameter\_\allowbreak{}error\_\allowbreak{}target} and \texttt{ipla\_\allowbreak{}linear\_\allowbreak{}bias\_\allowbreak{}le\_\allowbreak{}project03} with the proved nonnegative effective dimension. The conclusion is an actual parameter error inequality with Project 03's complete numerical right side. The sharper linear-step inequality has already been proved separately, so this theorem records the exact regime in which its right side is at least as good.
\end{leanbox}


Earlier results used here: \cite[Section~3.17.4]{LangevinLean2026ULASC}; \cite[Section~3.1.1]{LangevinLean2026IPLA}; \cite[Section~3.1.6]{LangevinLean2026IPLA}; \cite[Section~3.3.1]{LangevinLean2026IPLA}; \cite[Section~3.3.2]{LangevinLean2026IPLA}.
\subsection{\texorpdfstring{Explicit normalized accuracy and its iteration count}{Explicit normalized accuracy and its iteration count}}
\label{report-step-WassersteinMixing-IPLA-lean-6}


\textbf{Claim.} Fix \(\varepsilon>0\). Let \(R\) be either \(R_h\) or \(R_0\), matching the desired initialization form. Choose particles and step so that

\[N\varepsilon^2\ge18d_\theta,\qquad \gamma=\frac1{L\max\{1,18\sqrt D/\varepsilon\}}.\]


For every nonnegative integer \(k\) satisfying

\[k\ge\frac Lm\max\!\left\{1,\frac{18\sqrt D}{\varepsilon}\right\}\log\!\left(1+\frac{3\sqrt m\,R}{\varepsilon}\right),\]


the actual parameter law has normalized root-mean-square error at most \(\varepsilon\):

\[\sqrt m\left(\mathbb E_{\mu_k}\|\theta-\theta^\star\|^2\right)^{1/2}\le\varepsilon.\]


\textbf{Proof.} Multiply the corresponding sharp parameter bound by \(\sqrt m\). The concentration scaling identity gives \(\sqrt{2d_\theta/N}\). The scalar accuracy theorem allocates at most \(\varepsilon/3\) to that radius, the logarithmically controlled transient, and the normalized linear bias \(6L\gamma\sqrt D\). The displayed particle budget, step and iteration count are precisely its hypotheses, so their sum is at most \(\varepsilon\). The maximum keeps the step positive even when the effective dimension is zero.

\begin{leanbox}{Lean: \texttt{ipla\_\allowbreak{}parameter\_\allowbreak{}accuracy}, \texttt{ipla\_\allowbreak{}parameter\_\allowbreak{}accuracy\_\allowbreak{}target}}
\leansource{\texttt{ipla\_\allowbreak{}parameter\_\allowbreak{}accuracy} \quad WassersteinMixing/\allowbreak{}IPLA.lean, lines 181–202}
\begin{leancode}
\leanline{0}{theorem\ ipla\_parameter\_accuracy\ (A\ :\ IPLASetting)}
\leanline{4}{(θStar\ :\ EuclideanSpace\ ℝ\ (Fin\ A.dθ))\ (hmax\ :\ IPLA.IsMarginalMaximizer\ A.U\ θStar)}
\leanline{4}{(μ\ :\ Measure\ (IPLA.State\ A.N\ A.dθ\ A.dx))\ [IsProbabilityMeasure\ μ]}
\leanline{4}{(hμ\ :\ MemLp\ (id\ :\ IPLA.State\ A.N\ A.dθ\ A.dx\ →\ IPLA.State\ A.N\ A.dθ\ A.dx)\ 2\ μ)}
\leanline{4}{(ε\ :\ ℝ)\ (hε\ :\ 0\ <\ ε)}
\leanline{4}{(hparticles\ :\ 18\ *\ (A.dθ\ :\ ℝ)\ ≤\ (A.N\ :\ ℝ)\ *\ ε\ \textasciicircum{}\ 2)}
\leanline{4}{(hstep\ :\ A.γ\ =\ iplaAccuracyStep\ A.L\ (iplaEffectiveDimension\ A.N\ A.dθ\ A.dx)\ ε)}
\leanline{4}{(k\ :\ ℕ)}
\leanline{4}{(hk\ :\ iplaAccuracyPrefactor\ A.m\ A.L\ (iplaEffectiveDimension\ A.N\ A.dθ\ A.dx)\ ε\ *}
\leanline{6}{Real.log\ (1\ +\ 3\ *\ Real.sqrt\ A.m\ *\ ULA.wasserstein2Real\ μ\ A.invariantLaw\ /\ ε)\ ≤\ k)\ :}
\leanline{4}{Real.sqrt\ A.m\ *}
\leanline{6}{Real.sqrt\ (∫\ z,\ ‖IPLA.parameterCLM\ z\ -\ θStar‖\ \textasciicircum{}\ 2\ ∂A.law\ μ\ k)\ ≤\ ε}
\end{leancode}
\leansource{\texttt{ipla\_\allowbreak{}parameter\_\allowbreak{}accuracy\_\allowbreak{}target} \quad WassersteinMixing/\allowbreak{}IPLA.lean, lines 205–228}
\begin{leancode}
\leanline{0}{theorem\ ipla\_parameter\_accuracy\_target\ (A\ :\ IPLASetting)}
\leanline{4}{(θStar\ :\ EuclideanSpace\ ℝ\ (Fin\ A.dθ))\ (hmax\ :\ IPLA.IsMarginalMaximizer\ A.U\ θStar)}
\leanline{4}{(μ\ :\ Measure\ (IPLA.State\ A.N\ A.dθ\ A.dx))\ [IsProbabilityMeasure\ μ]}
\leanline{4}{(hμ\ :\ MemLp\ (id\ :\ IPLA.State\ A.N\ A.dθ\ A.dx\ →\ IPLA.State\ A.N\ A.dθ\ A.dx)\ 2\ μ)}
\leanline{4}{(ε\ :\ ℝ)\ (hε\ :\ 0\ <\ ε)}
\leanline{4}{(hparticles\ :\ 18\ *\ (A.dθ\ :\ ℝ)\ ≤\ (A.N\ :\ ℝ)\ *\ ε\ \textasciicircum{}\ 2)}
\leanline{4}{(hstep\ :\ A.γ\ =\ iplaAccuracyStep\ A.L\ (iplaEffectiveDimension\ A.N\ A.dθ\ A.dx)\ ε)}
\leanline{4}{(k\ :\ ℕ)}
\leanline{4}{(hk\ :\ iplaAccuracyPrefactor\ A.m\ A.L\ (iplaEffectiveDimension\ A.N\ A.dθ\ A.dx)\ ε\ *}
\leanline{6}{Real.log\ (1\ +\ 3\ *\ Real.sqrt\ A.m\ *}
\leanline{8}{ULA.wasserstein2Real\ μ\ (@IPLA.gibbsTarget\ A.N\ A.dθ\ A.dx\ A.U)\ /\ ε)\ ≤\ k)\ :}
\leanline{4}{Real.sqrt\ A.m\ *}
\leanline{6}{Real.sqrt\ (∫\ z,\ ‖IPLA.parameterCLM\ z\ -\ θStar‖\ \textasciicircum{}\ 2\ ∂A.law\ μ\ k)\ ≤\ ε}
\end{leancode}
\leanlabel{How the Lean proof follows the argument}
\texttt{ipla\_\allowbreak{}parameter\_\allowbreak{}accuracy} scales \texttt{ipla\_\allowbreak{}parameter\_\allowbreak{}error} and applies \texttt{ipla\_\allowbreak{}accuracy\_\allowbreak{}arithmetic} with the actual invariant-law distance. \texttt{ipla\_\allowbreak{}parameter\_\allowbreak{}accuracy\_\allowbreak{}target} uses the same arithmetic with \texttt{ipla\_\allowbreak{}parameter\_\allowbreak{}error\_\allowbreak{}target} and the actual Gibbs-based initial distance. \texttt{ipla\_\allowbreak{}concentration\_\allowbreak{}scale} converts the statistical radius exactly; nonnegativity of real Wasserstein distances supplies the scalar initialization sign. Both conclusions concern the actual Gaussian IPLA law.
\end{leanbox}


Earlier results used here: \cite[Section~3.17.4]{LangevinLean2026ULASC}; \cite[Section~3.1.1]{LangevinLean2026IPLA}; \cite[Section~3.1.6]{LangevinLean2026IPLA}; \cite[Section~3.3.1]{LangevinLean2026IPLA}; \cite[Section~3.3.2]{LangevinLean2026IPLA}.
\subsection{\texorpdfstring{The sharp rate for the original independent Gaussian recursion}{The sharp rate for the original independent Gaussian recursion}}
\label{report-step-WassersteinMixing-IPLA-lean-7}


\textbf{Claim.} Take a square-integrable scaled initial state on an actual probability space and an innovation sequence satisfying IPLA's original coordinatewise standard-Gaussian laws, independence across time and coordinates, and independence from initialization. Under the displayed \(C^2\), joint curvature and step assumptions, the original noise-driven IPLA laws satisfy

\[\left(\mathbb E\|\theta_k-\theta^\star\|^2\right)^{1/2}\le\sqrt{\frac{2d_\theta}{Nm}}+(1-m\gamma)^kW_2(\operatorname{Law}(Z_0),\pi_N)+6\gamma\sqrt{(L/m)L(d_x+d_\theta/N)}.\]


\textbf{Proof.} Assemble the structural setting from the displayed potential, dimensions and numerical constants. The pushforward law of the square-integrable initial state has finite second moment by the mapped-measure characterization of square integrability. Apply the already proved sharp target-based parameter theorem to this initial law. The exact identity of the Gaussian recursion laws replaces its kernel-iterate law by the law of the original recursion. This proves the same improved bound in the original algorithm, with all potential assumptions and all distributional requirements displayed explicitly.

\begin{leanbox}{Lean: \texttt{ipla\_\allowbreak{}gaussian\_\allowbreak{}parameter\_\allowbreak{}error}}
\leansource{\texttt{ipla\_\allowbreak{}gaussian\_\allowbreak{}parameter\_\allowbreak{}error} \quad WassersteinMixing/\allowbreak{}IPLA.lean, lines 233–261}
\begin{leancode}
\leanline{0}{theorem\ ipla\_gaussian\_parameter\_error}
\leanline{4}{\{Ω\ :\ Type*\}\ [MeasurableSpace\ Ω]\ (P\ :\ Measure\ Ω)\ [IsProbabilityMeasure\ P]}
\leanline{4}{(N\ dθ\ dx\ :\ ℕ)\ (hN\ :\ 0\ <\ N)\ (U\ :\ IPLA.Joint\ dθ\ dx\ →\ ℝ)}
\leanline{4}{(m\ γ\ :\ ℝ)\ (L\ :\ ℝ≥0)\ (hm\ :\ 0\ <\ m)\ (hmL\ :\ m\ ≤\ L)}
\leanline{4}{(hγ\ :\ 0\ <\ γ)\ (hγL\ :\ γ\ *\ L\ ≤\ 1)}
\leanline{4}{(hU\ :\ ContDiff\ ℝ\ 2\ U)\ (hsc\ :\ StrongConvexOn\ univ\ m\ U)}
\leanline{4}{(hLip\ :\ LipschitzWith\ L\ (gradient\ U))}
\leanline{4}{(θStar\ :\ EuclideanSpace\ ℝ\ (Fin\ dθ))\ (hmax\ :\ IPLA.IsMarginalMaximizer\ U\ θStar)}
\leanline{4}{(initial\ :\ Lp\ (IPLA.State\ N\ dθ\ dx)\ 2\ P)}
\leanline{4}{(noise\ :\ ℕ\ →\ Lp\ (IPLA.State\ N\ dθ\ dx)\ 2\ P)}
\leanline{4}{(hGaussian\ :\ IPLA.GaussianInnovations\ P\ initial\ noise)\ (k\ :\ ℕ)\ :}
\leanline{4}{Real.sqrt\ (∫\ z,\ ‖IPLA.parameterCLM\ z\ -\ θStar‖\ \textasciicircum{}\ 2}
\leanline{6}{∂IPLA.laws\ P\ U\ γ\ L\ (hU.differentiable\ (by\ norm\_num))\ hLip\ initial\ noise\ k)\ ≤}
\leanline{6}{Real.sqrt\ (2\ *\ (dθ\ :\ ℝ)\ /\ ((N\ :\ ℝ)\ *\ m))\ +}
\leanline{8}{(1\ -\ m\ *\ γ)\ \textasciicircum{}\ k\ *\ ULA.wasserstein2Real\ (P.map\ initial)}
\leanline{10}{(@IPLA.gibbsTarget\ N\ dθ\ dx\ U)\ +}
\leanline{8}{iplaLinearBias\ m\ L\ γ\ (iplaEffectiveDimension\ N\ dθ\ dx)}
\end{leancode}
\leanlabel{How the Lean proof follows the argument}
\texttt{ipla\_\allowbreak{}gaussian\_\allowbreak{}parameter\_\allowbreak{}error} has the full original-style probability-space, initialization and innovation binders. It constructs \texttt{IPLASetting} from those structural inputs, derives \texttt{Mem\allowbreak{}Lp id 2} of the initial pushforward using \texttt{Lp.\allowbreak{}mem\allowbreak{}Lp}, and applies \texttt{ipla\_\allowbreak{}parameter\_\allowbreak{}error\_\allowbreak{}target}. Finally \texttt{IPLASetting.\allowbreak{}law\_\allowbreak{}eq\_\allowbreak{}gaussian} rewrites the actual output law. Neither an auxiliary reference process nor an assumed transport estimate appears among the theorem inputs.
\end{leanbox}


Earlier results used here: \cite[Section~3.17.4]{LangevinLean2026ULASC}; \cite[Section~3.1.1]{LangevinLean2026IPLA}; \cite[Section~3.1.6]{LangevinLean2026IPLA}; \cite[Section~3.3.1]{LangevinLean2026IPLA}; \cite[Section~3.3.2]{LangevinLean2026IPLA}; \cite[Section~3.3.3]{LangevinLean2026IPLA}.
\section{\texorpdfstring{Explicit linear-step IPLA bounds}{Explicit linear-step IPLA bounds}}
\label{proof-module-WassersteinMixing-IPLATheorem-lean}


Fix an observation \(y\). Let the joint potential \(U_y\) be of class \(C^2\), \(m\)-strongly convex with \(m>0\), and have an \(L\)-Lipschitz gradient with \(L\ge0\). For \(N\ge1\) particles, \(\gamma>0\), \(\gamma L\le1\), a square-integrable scaled initial law \(\mu_0\) and independent standard Gaussian innovations, the actual IPLA recursion satisfies

\[\bigl(\mathbb E\|\theta_k-\theta_\star\|^2\bigr)^{1/2}\le\sqrt{\frac{2d_\theta}{Nm}}+(1-m\gamma)^kW_2(\mu_0,\Pi_N)+\frac{6L}{\sqrt m}\sqrt{d_x+\frac{d_\theta}N}\,\gamma\]

for every \(k\in\mathbb N\). The assumptions, notation and first two terms are those of Project 03; only the regularity, the step range and the last term differ. The same bound holds with any square-integrable stationary law of the recursion in place of \(\Pi_N\), and explicit particle, step and iteration choices give normalized accuracy \(\varepsilon\). On a nonzero joint space, the ordering \(m\le L\) follows from the potential assumptions, so the step restriction is equivalently \(\gamma\le1/L\). In zero joint dimension the state space is a point and the parameter error is zero, including when \(L=0\).



Earlier results used here: \cite[Section~3.18]{LangevinLean2026IPLA}.
\subsection{\texorpdfstring{Writing the discretization term explicitly}{Writing the discretization term explicitly}}
\label{report-step-WassersteinMixing-IPLATheorem-lean-1}


\textbf{Claim.} Let \(m>0\) and \(L,D\ge0\). Then

\[6\gamma\sqrt{\frac Lm\,LD}=\frac{6L}{\sqrt m}\sqrt D\,\gamma.\]

If also \(\gamma\ge0\), then

\[\frac{6L}{\sqrt m}\sqrt D\,\gamma=\frac{40}{11}\sqrt{m\gamma}\cdot\frac{33}{20}\frac Lm\sqrt D\sqrt\gamma,\]

and \(m\gamma\le121/1600\) implies \(\frac{6L}{\sqrt m}\sqrt D\,\gamma\le\frac{33}{20}\frac Lm\sqrt D\sqrt\gamma\).



\textbf{Proof.} Multiplying the left side of the first identity by \(\sqrt m\) gives \(6L\gamma\sqrt D\) by the earlier curvature scaling identity; divide by \(\sqrt m>0\). For the second identity write \(\sqrt{m\gamma}=\sqrt m\sqrt\gamma\), use \((\sqrt m)^2=m\), \((\sqrt\gamma)^2=\gamma\) and \(\frac{40}{11}\cdot\frac{33}{20}=6\). The inequality is the earlier factorized comparison with Project 03, rewritten by the first identity. Equivalently, \(\frac{40}{11}\sqrt{m\gamma}\le1\) exactly when \(m\gamma\le121/1600\).

\begin{leanbox}{Lean: \texttt{iplaLinearBias\_\allowbreak{}eq\_\allowbreak{}explicit}, \texttt{ipla\_\allowbreak{}linear\_\allowbreak{}step\_\allowbreak{}term\_\allowbreak{}eq\_\allowbreak{}project03\_\allowbreak{}mul}, \texttt{ipla\_\allowbreak{}linear\_\allowbreak{}step\_\allowbreak{}term\_\allowbreak{}le\_\allowbreak{}project03}}
\leansource{\texttt{iplaLinearBias\_\allowbreak{}eq\_\allowbreak{}explicit} \quad WassersteinMixing/\allowbreak{}IPLATheorem.lean, lines 22–30}
\begin{leancode}
\leanline{0}{theorem\ iplaLinearBias\_eq\_explicit\ (m\ L\ γ\ D\ :\ ℝ)\ (hm\ :\ 0\ <\ m)\ (hL\ :\ 0\ ≤\ L)\ (hD\ :\ 0\ ≤\ D)\ :}
\leanline{4}{iplaLinearBias\ m\ L\ γ\ D\ =\ 6\ *\ (L\ /\ Real.sqrt\ m)\ *\ Real.sqrt\ D\ *\ γ}
\end{leancode}
\leansource{\texttt{ipla\_\allowbreak{}linear\_\allowbreak{}step\_\allowbreak{}term\_\allowbreak{}eq\_\allowbreak{}project03\_\allowbreak{}mul} \quad WassersteinMixing/\allowbreak{}IPLATheorem.lean, lines 33–41}
\begin{leancode}
\leanline{0}{theorem\ ipla\_linear\_step\_term\_eq\_project03\_mul\ (m\ L\ γ\ D\ :\ ℝ)\ (hm\ :\ 0\ <\ m)\ (hγ\ :\ 0\ ≤\ γ)\ :}
\leanline{4}{6\ *\ (L\ /\ Real.sqrt\ m)\ *\ Real.sqrt\ D\ *\ γ\ =}
\leanline{6}{(40\ /\ 11\ :\ ℝ)\ *\ Real.sqrt\ (m\ *\ γ)\ *}
\leanline{8}{((33\ /\ 20\ :\ ℝ)\ *\ (L\ /\ m)\ *\ Real.sqrt\ D\ *\ Real.sqrt\ γ)}
\end{leancode}
\leansource{\texttt{ipla\_\allowbreak{}linear\_\allowbreak{}step\_\allowbreak{}term\_\allowbreak{}le\_\allowbreak{}project03} \quad WassersteinMixing/\allowbreak{}IPLATheorem.lean, lines 44–50}
\begin{leancode}
\leanline{0}{theorem\ ipla\_linear\_step\_term\_le\_project03\ (m\ L\ γ\ D\ :\ ℝ)\ (hm\ :\ 0\ <\ m)\ (hL\ :\ 0\ ≤\ L)}
\leanline{4}{(hγ\ :\ 0\ ≤\ γ)\ (hD\ :\ 0\ ≤\ D)\ (hsmall\ :\ m\ *\ γ\ ≤\ 121\ /\ 1600)\ :}
\leanline{4}{6\ *\ (L\ /\ Real.sqrt\ m)\ *\ Real.sqrt\ D\ *\ γ\ ≤}
\leanline{6}{(33\ /\ 20\ :\ ℝ)\ *\ (L\ /\ m)\ *\ Real.sqrt\ D\ *\ Real.sqrt\ γ}
\end{leancode}
\leanlabel{How the Lean proof follows the argument}
\texttt{ipla\allowbreak{}Linear\allowbreak{}Bias\_\allowbreak{}eq\_\allowbreak{}explicit} rewrites \texttt{ipla\allowbreak{}Linear\allowbreak{}Bias} through \texttt{ipla\_\allowbreak{}linear\_\allowbreak{}bias\_\allowbreak{}scale} and field cancellation by the positive square root. \texttt{ipla\_\allowbreak{}linear\_\allowbreak{}step\_\allowbreak{}term\_\allowbreak{}eq\_\allowbreak{}project03\_\allowbreak{}mul} splits \texttt{Real.\allowbreak{}sqrt (m * \ensuremath{\gamma})}, clears denominators and substitutes the two squared roots. \texttt{ipla\_\allowbreak{}linear\_\allowbreak{}step\_\allowbreak{}term\_\allowbreak{}le\_\allowbreak{}project03} rewrites \texttt{ipla\_\allowbreak{}linear\_\allowbreak{}bias\_\allowbreak{}le\_\allowbreak{}project03}; its right side is literally Project 03's discretization term.
\end{leanbox}
\subsection{\texorpdfstring{Ordering the curvature constants and the one-point space}{Ordering the curvature constants and the one-point space}}
\label{report-step-WassersteinMixing-IPLATheorem-lean-2}


\textbf{Claim.} If \(d_\theta+d_x>0\) and a differentiable joint potential \(U\) is \(m\)-strongly convex with an \(L\)-Lipschitz gradient, then \(m\le L\). On a one-point measurable space, any two probability measures are equal.



\textbf{Proof.} Take a unit coordinate vector \(v\). Strong monotonicity of the gradient, the Cauchy–Schwarz inequality and the Lipschitz bound give

\[m=m\|v\|^2\le\langle\nabla U(v)-\nabla U(0),v\rangle\le\|\nabla U(v)-\nabla U(0)\|\le L.\]

On a one-point space every set is empty or the whole space, where both probability measures take the values \(0\) and \(1\).

\begin{leanbox}{Lean: \texttt{ipla\_\allowbreak{}curvature\_\allowbreak{}le\_\allowbreak{}smoothness}, \texttt{probability\_\allowbreak{}eq\_\allowbreak{}of\_\allowbreak{}subsingleton}}
\leansource{\texttt{ipla\_\allowbreak{}curvature\_\allowbreak{}le\_\allowbreak{}smoothness} \quad WassersteinMixing/\allowbreak{}IPLATheorem.lean, lines 53–65}
\begin{leancode}
\leanline{0}{theorem\ ipla\_curvature\_le\_smoothness\ \{dθ\ dx\ :\ ℕ\}\ [Nonempty\ (Fin\ dθ\ ⊕\ Fin\ dx)]}
\leanline{4}{\{U\ :\ IPLA.Joint\ dθ\ dx\ →\ ℝ\}\ \{m\ :\ ℝ\}\ \{L\ :\ ℝ≥0\}}
\leanline{4}{(hdiff\ :\ Differentiable\ ℝ\ U)\ (hsc\ :\ StrongConvexOn\ univ\ m\ U)}
\leanline{4}{(hLip\ :\ LipschitzWith\ L\ (gradient\ U))\ :\ m\ ≤\ L}
\end{leancode}
\leansource{\texttt{probability\_\allowbreak{}eq\_\allowbreak{}of\_\allowbreak{}subsingleton} \quad WassersteinMixing/\allowbreak{}IPLATheorem.lean, lines 68–74}
\begin{leancode}
\leanline{0}{theorem\ probability\_eq\_of\_subsingleton\ \{α\ :\ Type*\}\ [MeasurableSpace\ α]\ [Subsingleton\ α]}
\leanline{4}{(μ\ ν\ :\ Measure\ α)\ [IsProbabilityMeasure\ μ]\ [IsProbabilityMeasure\ ν]\ :\ μ\ =\ ν}
\end{leancode}
\leanlabel{How the Lean proof follows the argument}
\texttt{ipla\_\allowbreak{}curvature\_\allowbreak{}le\_\allowbreak{}smoothness} uses \texttt{Euclidean\allowbreak{}Space.\allowbreak{}single i 1} for an index supplied by the nonempty coordinate type, then combines \texttt{ULA.\allowbreak{}strong\_\allowbreak{}gradient\_\allowbreak{}monotone}, \texttt{real\_\allowbreak{}inner\_\allowbreak{}le\_\allowbreak{}norm} and \texttt{Lipschitz\allowbreak{}With.\allowbreak{}dist\_\allowbreak{}le\_\allowbreak{}mul}. \texttt{probability\_\allowbreak{}eq\_\allowbreak{}of\_\allowbreak{}subsingleton} proves measure extensionality by splitting each set into the empty and nonempty cases.
\end{leanbox}


Earlier results used here: \cite[Section~3.2.3]{LangevinLean2026ULASC}; \cite[Section~3.1.1]{LangevinLean2026IPLA}.
\subsection{\texorpdfstring{Stationary laws of the Gaussian IPLA transition}{Stationary laws of the Gaussian IPLA transition}}
\label{report-step-WassersteinMixing-IPLATheorem-lean-3}


\textbf{Claim.} Let \(T\) be the Gaussian IPLA transition on scaled-state laws, under the structural assumptions \(U_y\in C^2\), \(m\)-strong convexity with \(m>0\), an \(L\)-Lipschitz gradient, \(\gamma>0\) and \(\gamma L\le1\). Then \(T\) has a stationary probability law with finite second moment. If \(\widehat\Pi\) is the constructed invariant law, \(\Pi'\) is any stationary law with finite second moment, and \(\mu\) has finite second moment, then

\[W_2(\mu,\widehat\Pi)\le W_2(\mu,\Pi').\]

\textbf{Proof.} When \(d_\theta+d_x>0\), the previous step gives \(L\ge m>0\), and the invariant law constructed for the lifted chain is a stationary probability law with finite second moment. When \(d_\theta=d_x=0\), the state space is a point; its Dirac law is stationary because both it and its image are probability measures on a point. In this case every probability law is that same point mass, so all distances in the displayed comparison are zero.



For the remaining comparison suppose \(d_\theta+d_x>0\). Both stationary laws have finite second moments, so \(W_2(\Pi',\widehat\Pi)<\infty\). One step of synchronous contraction gives

\[W_2(\Pi',\widehat\Pi)=W_2(T\Pi',T\widehat\Pi)\le(1-m\gamma)\,W_2(\Pi',\widehat\Pi).\]

Here \(0\le1-m\gamma<1\) because \(m\gamma\le L\gamma\le1\). This forces \(W_2(\Pi',\widehat\Pi)=0\). The triangle inequality then gives \(W_2(\mu,\widehat\Pi)\le W_2(\mu,\Pi')+0\).

\begin{leanbox}{Lean: \texttt{IPLASetting.\allowbreak{}wasserstein2Real\_\allowbreak{}invariantLaw\_\allowbreak{}le}, \texttt{ipla\_\allowbreak{}stationary\_\allowbreak{}law\_\allowbreak{}exists}}
\leansource{\texttt{IPLASetting.\allowbreak{}wasserstein2Real\_\allowbreak{}invariantLaw\_\allowbreak{}le} \quad WassersteinMixing/\allowbreak{}IPLATheorem.lean, lines 78–109}
\begin{leancode}
\leanline{0}{theorem\ IPLASetting.wasserstein2Real\_invariantLaw\_le\ (A\ :\ IPLASetting)}
\leanline{4}{(μ\ stationary\ :\ Measure\ (IPLA.State\ A.N\ A.dθ\ A.dx))}
\leanline{4}{[IsProbabilityMeasure\ μ]\ [IsProbabilityMeasure\ stationary]}
\leanline{4}{(hμ\ :\ MemLp\ (id\ :\ IPLA.State\ A.N\ A.dθ\ A.dx\ →\ IPLA.State\ A.N\ A.dθ\ A.dx)\ 2\ μ)}
\leanline{4}{(hmoment\ :\ MemLp\ (id\ :\ IPLA.State\ A.N\ A.dθ\ A.dx\ →\ IPLA.State\ A.N\ A.dθ\ A.dx)\ 2\ stationary)}
\leanline{4}{(hstationary\ :\ ULA.ulaLawStep\ (IPLA.potential\ A.U)\ (A.γ\ /\ A.N)\ stationary\ =\ stationary)\ :}
\leanline{4}{ULA.wasserstein2Real\ μ\ A.invariantLaw\ ≤\ ULA.wasserstein2Real\ μ\ stationary}
\end{leancode}
\leansource{\texttt{ipla\_\allowbreak{}stationary\_\allowbreak{}law\_\allowbreak{}exists} \quad WassersteinMixing/\allowbreak{}IPLATheorem.lean, lines 112–148}
\begin{leancode}
\leanline{0}{theorem\ ipla\_stationary\_law\_exists}
\leanline{4}{(N\ dθ\ dx\ :\ ℕ)\ (hN\ :\ 0\ <\ N)}
\leanline{4}{\{Data\ :\ Type*\}\ (U\ :\ Data\ →\ IPLA.Joint\ dθ\ dx\ →\ ℝ)\ (y\ :\ Data)}
\leanline{4}{(m\ γ\ :\ ℝ)\ (L\ :\ ℝ≥0)}
\leanline{4}{(hU\ :\ ContDiff\ ℝ\ 2\ (U\ y))\ (hsc\ :\ StrongConvexOn\ univ\ m\ (U\ y))}
\leanline{4}{(hLip\ :\ LipschitzWith\ L\ (gradient\ (U\ y)))}
\leanline{4}{(hm\ :\ 0\ <\ m)\ (hγ\ :\ 0\ <\ γ)\ (hγL\ :\ γ\ *\ L\ ≤\ 1)\ :}
\leanline{4}{∃\ stationary\ :\ Measure\ (IPLA.State\ N\ dθ\ dx),\ IsProbabilityMeasure\ stationary\ ∧}
\leanline{6}{MemLp\ (id\ :\ IPLA.State\ N\ dθ\ dx\ →\ IPLA.State\ N\ dθ\ dx)\ 2\ stationary\ ∧}
\leanline{6}{ULA.ulaLawStep\ (IPLA.potential\ (U\ y))\ (γ\ /\ N)\ stationary\ =\ stationary}
\end{leancode}
\leanlabel{How the Lean proof follows the argument}
\texttt{IPLASetting.\allowbreak{}wasserstein2\allowbreak{}Real\_\allowbreak{}invariant\allowbreak{}Law\_\allowbreak{}le} applies \texttt{ula\_\allowbreak{}iterate\_\allowbreak{}wasserstein2\allowbreak{}Real\_\allowbreak{}le} with one step, rewrites both laws by stationarity and \texttt{coordinate\allowbreak{}Invariant\allowbreak{}Law\_\allowbreak{}invariant}, and concludes with \texttt{wasserstein2\allowbreak{}Real\_\allowbreak{}triangle}; finiteness comes from \texttt{ULA.\allowbreak{}wasserstein2\_\allowbreak{}ne\_\allowbreak{}top\_\allowbreak{}of\_\allowbreak{}mem\allowbreak{}Lp}. \texttt{ipla\_\allowbreak{}stationary\_\allowbreak{}law\_\allowbreak{}exists} returns \texttt{IPLASetting.\allowbreak{}invariant\allowbreak{}Law} with its proved probability, moment and invariance properties, and uses \texttt{Measure.\allowbreak{}dirac 0} in zero dimensions.
\end{leanbox}


Earlier results used here: \cite[Section~3.17.4]{LangevinLean2026ULASC}; \cite[Section~3.21.2]{LangevinLean2026ULASC}; \cite[Section~3.36.7]{LangevinLean2026ULASC}; \cite[Section~3.1.1]{LangevinLean2026IPLA}.
\subsection{\texorpdfstring{The parameter bounds in Project 03's form}{The parameter bounds in Project 03's form}}
\label{report-step-WassersteinMixing-IPLATheorem-lean-4}


\textbf{Claim.} Under the assumptions in the summary, for every \(k\in\mathbb N\),

\[\bigl(\mathbb E\|\theta_k-\theta_\star\|^2\bigr)^{1/2}\le\sqrt{\frac{2d_\theta}{Nm}}+(1-m\gamma)^kW_2(\mu_0,\Pi_N)+\frac{6L}{\sqrt m}\sqrt{d_x+\frac{d_\theta}N}\,\gamma\]

The same bound holds with \(W_2(\mu_0,\Pi_N)\) replaced by \(W_2(\mu_0,\Pi')\) for any stationary law \(\Pi'\) of the recursion with finite second moment.



\textbf{Proof.} Suppose \(d_\theta+d_x>0\). Then \(m\le L\), so the earlier application gives the sharp target-based bound with discretization term \(6\gamma\sqrt{(L/m)L(d_x+d_\theta/N)}\), which equals the displayed term by the first step. For the stationary form, apply the invariant-law bound of the application and replace \(W_2(\mu_0,\widehat\Pi)\) by the larger \(W_2(\mu_0,\Pi')\), using the previous step and \((1-m\gamma)^k\ge0\), which holds because \(m\gamma\le L\gamma\le1\). If \(d_\theta=d_x=0\), the parameter error vanishes and every transport distance on the one-point state space is zero, so both sides vanish.

\begin{leanbox}{Lean: \texttt{ipla\_\allowbreak{}linear\_\allowbreak{}step\_\allowbreak{}bound}, \texttt{ipla\_\allowbreak{}linear\_\allowbreak{}step\_\allowbreak{}bound\_\allowbreak{}stationary}}
\leansource{\texttt{ipla\_\allowbreak{}linear\_\allowbreak{}step\_\allowbreak{}bound} \quad WassersteinMixing/\allowbreak{}IPLATheorem.lean, lines 153–186}
\begin{leancode}
\leanline{0}{theorem\ ipla\_linear\_step\_bound}
\leanline{4}{\{Ω\ :\ Type*\}\ [MeasurableSpace\ Ω]\ (P\ :\ Measure\ Ω)\ [IsProbabilityMeasure\ P]}
\leanline{4}{(N\ dθ\ dx\ :\ ℕ)\ (hN\ :\ 0\ <\ N)}
\leanline{4}{\{Data\ :\ Type*\}\ (U\ :\ Data\ →\ IPLA.Joint\ dθ\ dx\ →\ ℝ)\ (y\ :\ Data)}
\leanline{4}{(m\ γ\ :\ ℝ)\ (L\ :\ ℝ≥0)}
\leanline{4}{(hU\ :\ ContDiff\ ℝ\ 2\ (U\ y))\ (hsc\ :\ StrongConvexOn\ univ\ m\ (U\ y))}
\leanline{4}{(hLip\ :\ LipschitzWith\ L\ (gradient\ (U\ y)))}
\leanline{4}{(hm\ :\ 0\ <\ m)}
\leanline{4}{(hγ\ :\ 0\ <\ γ)\ (hγL\ :\ γ\ *\ L\ ≤\ 1)}
\leanline{4}{(θStar\ :\ EuclideanSpace\ ℝ\ (Fin\ dθ))\ (hMaximizer\ :\ IPLA.IsMarginalMaximizer\ (U\ y)\ θStar)}
\leanline{4}{(initial\ :\ Lp\ (IPLA.State\ N\ dθ\ dx)\ 2\ P)\ (noise\ :\ ℕ\ →\ Lp\ (IPLA.State\ N\ dθ\ dx)\ 2\ P)}
\leanline{4}{(hGaussian\ :\ IPLA.GaussianInnovations\ P\ initial\ noise)}
\leanline{4}{(k\ :\ ℕ)\ :}
\leanline{4}{Real.sqrt\ (∫\ z,\ ‖IPLA.parameterCLM\ z\ -\ θStar‖\ \textasciicircum{}\ 2}
\leanline{6}{∂(IPLA.laws\ P\ (U\ y)\ γ\ L\ (hU.differentiable\ (by\ norm\_num))\ hLip\ initial\ noise\ k))\ ≤}
\leanline{6}{Real.sqrt\ (2\ *\ (dθ\ :\ ℝ)\ /\ ((N\ :\ ℝ)\ *\ m))\ +}
\leanline{8}{(1\ -\ m\ *\ γ)\ \textasciicircum{}\ k\ *\ ULA.wasserstein2Real\ (P.map\ initial)\ (IPLA.gibbsTarget\ (U\ y))\ +}
\leanline{8}{6\ *\ ((L\ :\ ℝ)\ /\ Real.sqrt\ m)\ *}
\leanline{10}{Real.sqrt\ ((dx\ :\ ℝ)\ +\ (dθ\ :\ ℝ)\ /\ (N\ :\ ℝ))\ *\ γ}
\end{leancode}
\leansource{\texttt{ipla\_\allowbreak{}linear\_\allowbreak{}step\_\allowbreak{}bound\_\allowbreak{}stationary} \quad WassersteinMixing/\allowbreak{}IPLATheorem.lean, lines 190–245}
\begin{leancode}
\leanline{0}{theorem\ ipla\_linear\_step\_bound\_stationary}
\leanline{4}{\{Ω\ :\ Type*\}\ [MeasurableSpace\ Ω]\ (P\ :\ Measure\ Ω)\ [IsProbabilityMeasure\ P]}
\leanline{4}{(N\ dθ\ dx\ :\ ℕ)\ (hN\ :\ 0\ <\ N)}
\leanline{4}{\{Data\ :\ Type*\}\ (U\ :\ Data\ →\ IPLA.Joint\ dθ\ dx\ →\ ℝ)\ (y\ :\ Data)}
\leanline{4}{(m\ γ\ :\ ℝ)\ (L\ :\ ℝ≥0)}
\leanline{4}{(hU\ :\ ContDiff\ ℝ\ 2\ (U\ y))\ (hsc\ :\ StrongConvexOn\ univ\ m\ (U\ y))}
\leanline{4}{(hLip\ :\ LipschitzWith\ L\ (gradient\ (U\ y)))}
\leanline{4}{(hm\ :\ 0\ <\ m)}
\leanline{4}{(hγ\ :\ 0\ <\ γ)\ (hγL\ :\ γ\ *\ L\ ≤\ 1)}
\leanline{4}{(θStar\ :\ EuclideanSpace\ ℝ\ (Fin\ dθ))\ (hMaximizer\ :\ IPLA.IsMarginalMaximizer\ (U\ y)\ θStar)}
\leanline{4}{(stationary\ :\ Measure\ (IPLA.State\ N\ dθ\ dx))\ [IsProbabilityMeasure\ stationary]}
\leanline{4}{(hmoment\ :\ MemLp\ (id\ :\ IPLA.State\ N\ dθ\ dx\ →\ IPLA.State\ N\ dθ\ dx)\ 2\ stationary)}
\leanline{4}{(hstationary\ :\ ULA.ulaLawStep\ (IPLA.potential\ (U\ y))\ (γ\ /\ N)\ stationary\ =\ stationary)}
\leanline{4}{(initial\ :\ Lp\ (IPLA.State\ N\ dθ\ dx)\ 2\ P)\ (noise\ :\ ℕ\ →\ Lp\ (IPLA.State\ N\ dθ\ dx)\ 2\ P)}
\leanline{4}{(hGaussian\ :\ IPLA.GaussianInnovations\ P\ initial\ noise)}
\leanline{4}{(k\ :\ ℕ)\ :}
\leanline{4}{Real.sqrt\ (∫\ z,\ ‖IPLA.parameterCLM\ z\ -\ θStar‖\ \textasciicircum{}\ 2}
\leanline{6}{∂(IPLA.laws\ P\ (U\ y)\ γ\ L\ (hU.differentiable\ (by\ norm\_num))\ hLip\ initial\ noise\ k))\ ≤}
\leanline{6}{Real.sqrt\ (2\ *\ (dθ\ :\ ℝ)\ /\ ((N\ :\ ℝ)\ *\ m))\ +}
\leanline{8}{(1\ -\ m\ *\ γ)\ \textasciicircum{}\ k\ *\ ULA.wasserstein2Real\ (P.map\ initial)\ stationary\ +}
\leanline{8}{6\ *\ ((L\ :\ ℝ)\ /\ Real.sqrt\ m)\ *}
\leanline{10}{Real.sqrt\ ((dx\ :\ ℝ)\ +\ (dθ\ :\ ℝ)\ /\ (N\ :\ ℝ))\ *\ γ}
\end{leancode}
\leanlabel{How the Lean proof follows the argument}
Both theorems have the binders of \texttt{IPLA.\allowbreak{}theorem1}, with \texttt{h\allowbreak{}U : Cont\allowbreak{}Diff \ensuremath{\mathbb R} 2 (U y)} in place of differentiability and \texttt{h\ensuremath{\gamma}\allowbreak{}L : \ensuremath{\gamma} * L \ensuremath{\le} 1} in place of \texttt{\ensuremath{\gamma} \ensuremath{\le} 2 /\allowbreak{} (m + L)}; the stationary form adds the stationary law with its moment and invariance hypotheses. In nonzero dimension, \texttt{ipla\_\allowbreak{}linear\_\allowbreak{}step\_\allowbreak{}bound} applies \texttt{ipla\_\allowbreak{}gaussian\_\allowbreak{}parameter\_\allowbreak{}error} and rewrites by \texttt{ipla\allowbreak{}Linear\allowbreak{}Bias\_\allowbreak{}eq\_\allowbreak{}explicit}; \texttt{ipla\_\allowbreak{}linear\_\allowbreak{}step\_\allowbreak{}bound\_\allowbreak{}stationary} applies \texttt{ipla\_\allowbreak{}parameter\_\allowbreak{}error}, \texttt{IPLASetting.\allowbreak{}law\_\allowbreak{}eq\_\allowbreak{}gaussian} and \texttt{IPLASetting.\allowbreak{}wasserstein2\allowbreak{}Real\_\allowbreak{}invariant\allowbreak{}Law\_\allowbreak{}le}. The zero-dimensional branch uses \texttt{IPLA.\allowbreak{}parameter\allowbreak{}RMS\_\allowbreak{}zero\_\allowbreak{}dimension} and \texttt{ULA.\allowbreak{}wasserstein2\allowbreak{}Real\_\allowbreak{}eq\_\allowbreak{}zero\_\allowbreak{}of\_\allowbreak{}subsingleton}, as in Project 03.
\end{leanbox}


Earlier results used here: \cite[Section~3.17.4]{LangevinLean2026ULASC}; \cite[Section~3.36.7]{LangevinLean2026ULASC}; \cite[Section~3.51.3]{LangevinLean2026ULASC}; \cite[Section~3.1.1]{LangevinLean2026IPLA}; \cite[Section~3.1.6]{LangevinLean2026IPLA}; \cite[Section~3.18.1]{LangevinLean2026IPLA}; \cite[Section~3.3.1]{LangevinLean2026IPLA}; \cite[Section~3.3.2]{LangevinLean2026IPLA}; \cite[Section~3.3.3]{LangevinLean2026IPLA}; \cite[Section~3.5.1]{LangevinLean2026IPLA}.
\subsection{\texorpdfstring{Explicit accuracy with either initialization distance}{Explicit accuracy with either initialization distance}}
\label{report-step-WassersteinMixing-IPLATheorem-lean-5}


\textbf{Claim.} Let \(\varepsilon>0\), \(N\varepsilon^2\ge18d_\theta\), and \(D_N=d_x+d_\theta/N\). If \(D_N=0\), both dimensions are zero and the normalized parameter error is zero for every iteration and every admissible positive step; no accuracy-driven step or iteration count is needed. Suppose \(D_N>0\). Then \(L\ge m>0\), and choose

\[\gamma=\frac1{L\max\bigl(1,18\sqrt{D_N}/\varepsilon\bigr)}.\]

If \(R\) is \(W_2(\mu_0,\Pi_N)\), or \(W_2(\mu_0,\Pi')\) for a stationary law \(\Pi'\) with finite second moment, then every \(k\in\mathbb N\) with

\[k\ge\frac Lm\max\Bigl(1,\frac{18\sqrt{D_N}}\varepsilon\Bigr)\log\Bigl(1+\frac{3\sqrt m\,R}\varepsilon\Bigr)\]

satisfies \(\sqrt m\,\bigl(\mathbb E\|\theta_k-\theta_\star\|^2\bigr)^{1/2}\le\varepsilon\).



\textbf{Proof.} The zero-dimensional conclusion follows directly from the zero parameter error. In nonzero dimension, the displayed step is positive and satisfies \(\gamma L\le1\). The earlier accuracy theorems then apply: they bound the normalized concentration, transient and discretization terms by \(\varepsilon/3\) each. For a stationary law \(\Pi'\), the previous comparison \(W_2(\mu_0,\widehat\Pi)\le W_2(\mu_0,\Pi')\), monotonicity of the logarithm and nonnegativity of the prefactor show that the displayed count also satisfies the condition stated with the constructed law \(\widehat\Pi\).

\begin{leanbox}{Lean: \texttt{ipla\_\allowbreak{}linear\_\allowbreak{}step\_\allowbreak{}accuracy}, \texttt{ipla\_\allowbreak{}linear\_\allowbreak{}step\_\allowbreak{}accuracy\_\allowbreak{}stationary}}
\leansource{\texttt{ipla\_\allowbreak{}linear\_\allowbreak{}step\_\allowbreak{}accuracy} \quad WassersteinMixing/\allowbreak{}IPLATheorem.lean, lines 251–294}
\begin{leancode}
\leanline{0}{theorem\ ipla\_linear\_step\_accuracy}
\leanline{4}{\{Ω\ :\ Type*\}\ [MeasurableSpace\ Ω]\ (P\ :\ Measure\ Ω)\ [IsProbabilityMeasure\ P]}
\leanline{4}{(N\ dθ\ dx\ :\ ℕ)\ (hN\ :\ 0\ <\ N)}
\leanline{4}{\{Data\ :\ Type*\}\ (U\ :\ Data\ →\ IPLA.Joint\ dθ\ dx\ →\ ℝ)\ (y\ :\ Data)}
\leanline{4}{(m\ γ\ :\ ℝ)\ (L\ :\ ℝ≥0)}
\leanline{4}{(hU\ :\ ContDiff\ ℝ\ 2\ (U\ y))\ (hsc\ :\ StrongConvexOn\ univ\ m\ (U\ y))}
\leanline{4}{(hLip\ :\ LipschitzWith\ L\ (gradient\ (U\ y)))}
\leanline{4}{(hm\ :\ 0\ <\ m)}
\leanline{4}{(θStar\ :\ EuclideanSpace\ ℝ\ (Fin\ dθ))\ (hMaximizer\ :\ IPLA.IsMarginalMaximizer\ (U\ y)\ θStar)}
\leanline{4}{(initial\ :\ Lp\ (IPLA.State\ N\ dθ\ dx)\ 2\ P)\ (noise\ :\ ℕ\ →\ Lp\ (IPLA.State\ N\ dθ\ dx)\ 2\ P)}
\leanline{4}{(hGaussian\ :\ IPLA.GaussianInnovations\ P\ initial\ noise)}
\leanline{4}{(ε\ :\ ℝ)\ (hε\ :\ 0\ <\ ε)\ (hparticles\ :\ 18\ *\ (dθ\ :\ ℝ)\ ≤\ (N\ :\ ℝ)\ *\ ε\ \textasciicircum{}\ 2)}
\leanline{4}{(hstep\ :\ γ\ =\ 1\ /\ ((L\ :\ ℝ)\ *}
\leanline{6}{max\ 1\ (18\ *\ Real.sqrt\ ((dx\ :\ ℝ)\ +\ (dθ\ :\ ℝ)\ /\ (N\ :\ ℝ))\ /\ ε)))}
\leanline{4}{(k\ :\ ℕ)}
\leanline{4}{(hk\ :\ (L\ :\ ℝ)\ /\ m\ *\ max\ 1\ (18\ *\ Real.sqrt\ ((dx\ :\ ℝ)\ +\ (dθ\ :\ ℝ)\ /\ (N\ :\ ℝ))\ /\ ε)\ *}
\leanline{6}{Real.log\ (1\ +\ 3\ *\ Real.sqrt\ m\ *}
\leanline{8}{ULA.wasserstein2Real\ (P.map\ initial)\ (IPLA.gibbsTarget\ (U\ y))\ /\ ε)\ ≤\ k)\ :}
\leanline{4}{Real.sqrt\ m\ *\ Real.sqrt\ (∫\ z,\ ‖IPLA.parameterCLM\ z\ -\ θStar‖\ \textasciicircum{}\ 2}
\leanline{6}{∂(IPLA.laws\ P\ (U\ y)\ γ\ L\ (hU.differentiable\ (by\ norm\_num))\ hLip\ initial\ noise\ k))\ ≤\ ε}
\end{leancode}
\leansource{\texttt{ipla\_\allowbreak{}linear\_\allowbreak{}step\_\allowbreak{}accuracy\_\allowbreak{}stationary} \quad WassersteinMixing/\allowbreak{}IPLATheorem.lean, lines 298–359}
\begin{leancode}
\leanline{0}{theorem\ ipla\_linear\_step\_accuracy\_stationary}
\leanline{4}{\{Ω\ :\ Type*\}\ [MeasurableSpace\ Ω]\ (P\ :\ Measure\ Ω)\ [IsProbabilityMeasure\ P]}
\leanline{4}{(N\ dθ\ dx\ :\ ℕ)\ (hN\ :\ 0\ <\ N)}
\leanline{4}{\{Data\ :\ Type*\}\ (U\ :\ Data\ →\ IPLA.Joint\ dθ\ dx\ →\ ℝ)\ (y\ :\ Data)}
\leanline{4}{(m\ γ\ :\ ℝ)\ (L\ :\ ℝ≥0)}
\leanline{4}{(hU\ :\ ContDiff\ ℝ\ 2\ (U\ y))\ (hsc\ :\ StrongConvexOn\ univ\ m\ (U\ y))}
\leanline{4}{(hLip\ :\ LipschitzWith\ L\ (gradient\ (U\ y)))}
\leanline{4}{(hm\ :\ 0\ <\ m)}
\leanline{4}{(θStar\ :\ EuclideanSpace\ ℝ\ (Fin\ dθ))\ (hMaximizer\ :\ IPLA.IsMarginalMaximizer\ (U\ y)\ θStar)}
\leanline{4}{(stationary\ :\ Measure\ (IPLA.State\ N\ dθ\ dx))\ [IsProbabilityMeasure\ stationary]}
\leanline{4}{(hmoment\ :\ MemLp\ (id\ :\ IPLA.State\ N\ dθ\ dx\ →\ IPLA.State\ N\ dθ\ dx)\ 2\ stationary)}
\leanline{4}{(hstationary\ :\ ULA.ulaLawStep\ (IPLA.potential\ (U\ y))\ (γ\ /\ N)\ stationary\ =\ stationary)}
\leanline{4}{(initial\ :\ Lp\ (IPLA.State\ N\ dθ\ dx)\ 2\ P)\ (noise\ :\ ℕ\ →\ Lp\ (IPLA.State\ N\ dθ\ dx)\ 2\ P)}
\leanline{4}{(hGaussian\ :\ IPLA.GaussianInnovations\ P\ initial\ noise)}
\leanline{4}{(ε\ :\ ℝ)\ (hε\ :\ 0\ <\ ε)\ (hparticles\ :\ 18\ *\ (dθ\ :\ ℝ)\ ≤\ (N\ :\ ℝ)\ *\ ε\ \textasciicircum{}\ 2)}
\leanline{4}{(hstep\ :\ γ\ =\ 1\ /\ ((L\ :\ ℝ)\ *}
\leanline{6}{max\ 1\ (18\ *\ Real.sqrt\ ((dx\ :\ ℝ)\ +\ (dθ\ :\ ℝ)\ /\ (N\ :\ ℝ))\ /\ ε)))}
\leanline{4}{(k\ :\ ℕ)}
\leanline{4}{(hk\ :\ (L\ :\ ℝ)\ /\ m\ *\ max\ 1\ (18\ *\ Real.sqrt\ ((dx\ :\ ℝ)\ +\ (dθ\ :\ ℝ)\ /\ (N\ :\ ℝ))\ /\ ε)\ *}
\leanline{6}{Real.log\ (1\ +\ 3\ *\ Real.sqrt\ m\ *\ ULA.wasserstein2Real\ (P.map\ initial)\ stationary\ /\ ε)\ ≤\ k)\ :}
\leanline{4}{Real.sqrt\ m\ *\ Real.sqrt\ (∫\ z,\ ‖IPLA.parameterCLM\ z\ -\ θStar‖\ \textasciicircum{}\ 2}
\leanline{6}{∂(IPLA.laws\ P\ (U\ y)\ γ\ L\ (hU.differentiable\ (by\ norm\_num))\ hLip\ initial\ noise\ k))\ ≤\ ε}
\end{leancode}
\leanlabel{How the Lean proof follows the argument}
In the nonzero-dimensional branch, the step hypothesis \texttt{hstep} determines \(\gamma\); \texttt{ipla\allowbreak{}Accuracy\allowbreak{}Step\_\allowbreak{}pos} and \texttt{ipla\allowbreak{}Accuracy\allowbreak{}Step\_\allowbreak{}mul\_\allowbreak{}L\_\allowbreak{}le} supply the structure fields. \texttt{ipla\_\allowbreak{}linear\_\allowbreak{}step\_\allowbreak{}accuracy} applies \texttt{ipla\_\allowbreak{}parameter\_\allowbreak{}accuracy\_\allowbreak{}target}, whose step and prefactor are definitionally the displayed formulas. \texttt{ipla\_\allowbreak{}linear\_\allowbreak{}step\_\allowbreak{}accuracy\_\allowbreak{}stationary} proves the logarithmic comparison with \texttt{Real.\allowbreak{}log\_\allowbreak{}le\_\allowbreak{}log} and \texttt{gcongr}, then applies \texttt{ipla\_\allowbreak{}parameter\_\allowbreak{}accuracy}. Both rewrite the output law by \texttt{IPLASetting.\allowbreak{}law\_\allowbreak{}eq\_\allowbreak{}gaussian}; their zero-dimensional branches conclude directly from the zero parameter error.
\end{leanbox}


Earlier results used here: \cite[Section~3.17.4]{LangevinLean2026ULASC}; \cite[Section~3.36.7]{LangevinLean2026ULASC}; \cite[Section~3.1.1]{LangevinLean2026IPLA}; \cite[Section~3.1.6]{LangevinLean2026IPLA}; \cite[Section~3.18.1]{LangevinLean2026IPLA}; \cite[Section~3.3.1]{LangevinLean2026IPLA}; \cite[Section~3.3.2]{LangevinLean2026IPLA}; \cite[Section~3.3.3]{LangevinLean2026IPLA}.
\section{\texorpdfstring{The Wasserstein mixing theorem and its IPLA application}{The Wasserstein mixing theorem and its IPLA application}}
\label{proof-module-WassersteinMixing-lean}


This import bundle collects the actual Brownian Gibbs diffusion, its proved reversible strongly continuous contraction semigroup, the bounded Poisson corrector, the actual synchronous gradient comparison and the constructed Gaussian ULA invariant law. The source bias, convergence and accuracy estimates follow from the original twice continuously differentiable Hessian bounds, with no extra realization or smoothing assumption.



The additional IPLA application proves linear-step parameter rates for the actual Gaussian particle algorithm, both with invariant-law and exact-Gibbs initialization, together with explicit accuracy choices and an exact comparison to Project 03. Its joint \(C^2\) regularity and small-step assumptions are stated separately from the scope of Project 03.

\section{\texorpdfstring{Kernel audits of the source mixing results and the IPLA application}{Kernel audits of the source mixing results and the IPLA application}}
\label{proof-module-WassersteinMixing-Check-lean}


This audit exposes the dependencies of the actual curvature, Brownian and Markov constructions, reversible semigroup, bounded corrector, shifted-energy argument, invariant law and source endpoints. It also audits the IPLA lifting, coordinate covariance, constructed invariant law, sharp target-based comparison, identity of the Gaussian recursion laws, parameter error, Project 03 comparison and both normalized accuracy endpoints. The permitted logical axioms are standard propositional extensionality, choice and quotient soundness. No added mathematical axiom or proof placeholder is permitted.


\begin{thebibliography}{ULA-LSI}
\bibitem[ULA-SC]{LangevinLean2026ULASC}
LangevinLean Project. \emph{Unadjusted Langevin under Strong Convexity}. LL-01, 2026. Draft report, content revision d34ead06987c. \url{https://odakyildiz.github.io/LangevinLean/docs/pdf/UnadjustedLangevin/FullProof.pdf}.
\bibitem[IPLA-SC]{LangevinLean2026IPLA}
LangevinLean Project. \emph{Interacting Particle Langevin under Strong Convexity}. LL-03, 2026. Draft report, content revision ec0201df2207. \url{https://odakyildiz.github.io/LangevinLean/docs/pdf/IPLA/FullProof.pdf}.
\bibitem[Kinetic]{LangevinLean2026Kinetic}
LangevinLean Project. \emph{Kinetic Langevin Monte Carlo}. LL-06, 2026. Draft report, content revision b2fee42a9593. \url{https://odakyildiz.github.io/LangevinLean/docs/pdf/KineticLangevin/FullProof.pdf}.
\end{thebibliography}

\end{document}
